\section{B. Good reduction and inertia characters}
\subsection{Explicit counterexamples to the mod-$p$ good-reduction criterion}

\begin{theorem}[B-RD1: a counterexample for every odd prime]\label{thm:B-RD1}
\status{Proved.}
For every prime $p\geq3$, there is a complete discretely valued field $K$
of characteristic zero, residue field $\overline{\mathbf F}_p$, and
absolute ramification index
\[
 e(K)=\frac{p^2-1}{2},
\]
containing $\mu_p$, and an elliptic curve $E/K$ whose full $p$-torsion
is $K$-rational and whose reduction is additive, potentially good and
supersingular. Its semistability defect is exactly $2$. Its inertia
character is $-2$ on the nontrivial element of the finite inertia image.
\end{theorem}
\begin{proof}
We give the field and the curve by an explicit algebraic construction.
Put $k=\overline{\mathbf F}_p$, $R=W(k)$ and $F=\operatorname{Frac}R$.
Let $m=(p-1)/2$ and choose any root $\lambda\in k$ of
\[
 H_p(T)=\sum_{i=0}^{m}\binom mi^2T^i.
\]
Such a root exists since $m\geq1$. It is neither $0$ nor $1$:
$H_p(0)=1$ and
$H_p(1)=\binom{2m}{m}\equiv(-1)^m\pmod p$.
The coefficient calculation for Frobenius on the one-dimensional space
$H^1(\mathcal O)$ of a Legendre elliptic curve gives $(-1)^mH_p(\lambda)$
as its Hasse--Witt scalar. Thus
\[
 E_0/R:\quad y^2=x(x-1)(x-[\lambda])
\]
is smooth and has supersingular special fibre. Here $[\lambda]$ is the
Teichmuller representative. The discriminant is
$16[\lambda]^2(1-[\lambda])^2$, an $R$-unit.

Let $L=F(E_0[p])$. We now verify its ramification without invoking a
Breuil--Kisin dictionary. Every nonzero geometric $p$-torsion point
specializes to the identity, since the special fibre is supersingular.
For a local parameter $T$ of the formal group at that identity, its
height is two. Consequently $[p](T)/T$ has constant coefficient $p$,
its first coefficient which is a unit modulo $p$ is that of
$T^{p^2-1}$, and all lower nonconstant coefficients are divisible by $p$.
Weierstrass preparation gives a polynomial of degree $p^2-1$ which is
Eisenstein over $R$. Thus every nonzero $P\in E_0[p]$ generates a
completely ramified extension $F(P)/F$ of degree $p^2-1$. This degree is
prime to $p$, so the extension is tame. Since $F$ has algebraically
closed residue field, its unique tame extension of that degree is
Galois cyclic and contains every such point. Hence
\[
 \operatorname{Gal}(L/F)\simeq C_{p^2-1}.
\]
Its action on $E_0[p]\setminus\{0\}$ is free and transitive. The unique
involution in this cyclic group acts as $-1$: an involution in
$\mathrm{GL}_2(\mathbf F_p)$ with any eigenvalue $+1$ would fix a
nonzero vector.

Let $H$ be this order-two subgroup and put $K=L^H$. The Weil pairing
shows $\mu_p\subset L$, and $H$ acts on it by the determinant of $-1$,
which is $1$. Thus $\mu_p\subset K$. We have $e(K)=(p^2-1)/2$.
Choose any nonzero $P\in E_0[p]$, and put $\delta=y(P)^2\in K^\times$.
A nonzero odd-order point has $y(P)\ne0$, so $L=K(\sqrt\delta)$.
Write $E_0$ as $y^2=x^3+a_2x^2+a_4x+a_6$ and define
\[
 E:\quad y^2=x^3+\delta a_2x^2+\delta^2a_4x+\delta^3a_6.
\]
Over $L$, the map $(x,y)\mapsto(\delta x,\delta\sqrt\delta\,y)$
identifies $E_0$ with $E$. For every $Q\in E_0[p]$, both $x(Q)$ and
$\sqrt\delta\,y(Q)$ are fixed by $H$. Hence every point of $E[p]$
is $K$-rational. Equivalently, its representation is the tensor product
of the scalar quadratic representation on $E_0[p]$ and its inverse.

The curve $E_0$ has good reduction over $K$, so for every prime
$\ell\ne p$ inertia acts trivially on $T_\ell(E_0)$. On $T_\ell(E)$ it
acts by the nontrivial ramified quadratic character of $L/K$. The
Neron--Ogg--Shafarevich criterion\cite[Theorem~1]{SerreTate1968Good} therefore excludes good reduction.
It has good supersingular reduction over $L$, so its bad reduction is
additive and its minimal good-reduction inertia quotient has order $2$.
The trace on its covariant Dieudonne module is $-2$ at the involution,
since that involution acts as $[-1]$ on the special elliptic curve.
\end{proof}

\begin{theorem}[B-RD2: full $3$-torsion already over $\mathbf Q_3(\zeta_3)$]\label{thm:B-RD2}
\status{Proved, with an exact computation.}
Let $K=\mathbf Q_3(\zeta)$, where $\zeta^2+\zeta+1=0$. Then $K$ has
absolute ramification index $2$ and residue field $\mathbf F_3$. The curve
\[
 E/K:\quad y^2+y=x^3
\]
has $E[3]\subset E(K)$, additive potentially supersingular reduction,
and semistability defect exactly $6$. Its inertia character takes the
value $-2$ on the element of order $2$ of its finite inertia image.
\end{theorem}
\begin{proof}
The discriminant and $j$-invariant are $\Delta=-27$ and $j=0$.
The third division polynomial is
\[
 \psi_3(x)=3x(x^3+1).
\]
Besides the identity, the following are the eight distinct $3$-torsion
points:
\[
 (0,0),\ (0,-1),\qquad
 (x,y)\quad\text{with }x\in\{-1,-\zeta,-\zeta^2\},\quad
 y\in\{\zeta,\zeta^2\}.
\]
They exhaust $E[3]$, which has degree nine in characteristic zero.
The exact program \texttt{verify\_p3\_full\_torsion.py} checks all these
points, their orders, all $81$ sums, and that $(0,0),(-1,\zeta)$ form a
basis. Its recorded output is part of the computation certificate.

The displayed equation has $v_K(\Delta)=6$. Discriminants of two
Weierstrass equations of the same curve differ by a twelfth power.
Thus no equation over $K$ can have discriminant valuation zero, and
$E$ does not have good reduction.

We construct good reduction after ramification degree six. Make first
the unramified quadratic extension $K'=K(i)$, with $i^2=-1$, and choose
$\sqrt3\in K'$. Adjoin $r$ with $r^3=-1/4$, and then $u$ with
$u^2=r\sqrt3$. The polynomial for $s=r+1$ is
\[
 s^3-3s^2+3s-3/4.
\]
Its Newton polygon has the single slope $-2/3$ for the valuation
normalized on $K'$, so $K'(r)/K'$ is completely ramified of degree three.
The element $r\sqrt3$ has odd valuation in that field, so adjoining $u$
gives a completely ramified quadratic extension. Put $L=K'(r,u)$.
After completing the square and setting
\[
 x=r+u^2X,\qquad y+\tfrac12=u^3Y,
\]
the equation becomes
\[
 Y^2=X^3+\frac{3r}{u^2}X^2+X.
\]
The middle coefficient has positive valuation and the special fibre is
$Y^2=X^3+X$, which is smooth and supersingular in characteristic three.
Thus $E$ has good reduction over an extension of ramification degree six.

We verify minimality, including possible unramified extensions. Let
$F/K$ have good reduction. The discriminant valuation implies
$2\mid e(F/K)$. In a good integral Weierstrass equation over $F$, complete
the square; this is integral since $2$ is a unit. The resulting equation
has the form $Y^2=X^3+b_2X^2+b_4X+b_6$. Comparison with
$y^2=x^3+1/4$ uses $x=r_0+u_0^2X$, $y=u_0^3Y$, and forces
$v_3(u_0)=1/4$. Integrality of $b_6$ gives
\[
 v_3(r_0^3+1/4)\geq3/2.
\]
It follows that $r_0$ is a unit with residue $-1$. Writing $s_0=r_0+1$,
if $v_3(s_0)<1/3$ the term $s_0^3$ uniquely has lowest valuation in
$s_0^3-3s_0^2+3s_0-3/4$; if $v_3(s_0)>1/3$ the constant term uniquely
has lowest valuation. Both alternatives give valuation at most one.
Hence $v_3(s_0)=1/3$, so $3\mid e(F/\mathbf Q_3)$, and therefore
$3\mid e(F/K)$. We conclude $6\mid e(F/K)$, proving the defect is six.

Over the unramified extension $K'$, the extension just constructed is
Galois: $r\mapsto\zeta r$, $u\mapsto\zeta^2u$ has order three and
commutes with $u\mapsto-u$. Its inertia group is cyclic of order six.
The latter involution induces $Y\mapsto-Y$ on the good special fibre,
that is $[-1]$, so its trace is $-2$.
\end{proof}

\begin{proposition}[B-RD3: the ordinary case of R3(a)]\label{prop:B-RD3}
\status{Proved.}
Let $K$ be complete discretely valued of characteristic zero, with
perfect residue field of characteristic $p\geq3$, and assume
$\mu_p\subset K$. If an abelian variety $A/K$ has potentially good
ordinary reduction and $A[p]$ is unramified, then $A$ has good reduction.
There is no restriction on $\dim A$.
\end{proposition}
\begin{proof}
We may pass to the completed maximal unramified extension: good
reduction is equivalent before and after this passage by the
Neron--Ogg--Shafarevich criterion. The residue field is now algebraically
closed. Choose a finite Galois extension $L/K$ giving good reduction,
and let $\mathcal A/\mathcal O_L$ be the resulting abelian scheme.
For its ordinary special fibre, the connected subgroup
$\mathcal A[p]^0$ is of multiplicative type and its etale quotient has
rank $p^g$. Indeed the Cartier dual of the connected subgroup has
etale special fibre, hence is etale by Nakayama's lemma applied to its
relative differentials. Over the strictly henselian ring
$\mathcal O_L$, it is constant; thus
$\mathcal A[p]^0\simeq\mu_p^g$.

Every inertia element acts semilinearly on this subgroup by its action
on the base, composed with a matrix in $\mathrm{GL}_g(\mathbf F_p)$.
The base action fixes $\mu_p$, and its generic action on $A[p]$ is
trivial. The matrix is therefore the identity. The action on the etale
quotient is trivial for the same reason. Over the algebraically closed
residue field the connected--etale sequence splits, and there are no
nonzero homomorphisms between the connected and etale factors in either
direction. Thus inertia acts identically on $\mathcal A_k[p]$.

If $\alpha$ is one of the resulting finite-order automorphisms of
$\mathcal A_k$, then $\alpha-1$ factors through $[p]$, so
$\alpha=1+p\beta$ for some $\beta\in\operatorname{End}(\mathcal A_k)$.
The kernel of reduction modulo $p$ in the units of this endomorphism
ring has no nontrivial finite-order elements for $p\geq3$. To see this,
pass to its finite free $\mathbf Z_p$-completion. If
$\alpha=1+p^a\gamma$ with $a\geq1$ and $\gamma\notin p\operatorname{End}$,
the binomial expansion shows that a prime-to-$p$ power has the same
leading valuation, and the $p$-th power has leading valuation $a+1$;
neither can be one. This proves the assertion for any finite order.
Hence inertia is trivial on the entire special abelian variety.
Proper smooth base change now makes $T_\ell(A)$ unramified for
$\ell\ne p$, and the same good-reduction criterion concludes.
\end{proof}

\begin{corollary}[B-RD6: the manuscript's cas test is negative]\label{cor:B-RD6}
\status{Refuted, for every odd prime.}
For every $p\geq3$ there is a field $K$ as in Theorem~\ref{thm:B-RD1}
and two elliptic curves $E,E'/K$ with potentially good reduction and
isomorphic $p$-torsion, for which the inertia traces differ on an element
of order prime to $p$.
\end{corollary}
\begin{proof}
Use $E$ and $K$ from Theorem~\ref{thm:B-RD1}. Choose an ordinary
elliptic curve $C/\overline{\mathbf F}_p$; for example choose a Legendre
parameter outside the finite root set of $H_p$ and outside $\{0,1\}$.
The ordinary Serre--Tate coordinates of
\cite[Theorem~2.1]{Katz1981SerreTate} give its canonical deformation
over $W(\overline{\mathbf F}_p)$, whose $p$-divisible group is
\[
 \mu_{p^\infty}\oplus\mathbf Q_p/\mathbf Z_p.
\]
This formal elliptic scheme algebraizes: its identity section provides
the compatible relatively ample bundles $\mathcal O(3[0])$, so the
formal existence theorem applies. Let $E'$ be its generic fibre after
extension to $K$. Then $E'$ has good reduction, and $\mu_p\subset K$
gives $E'[p]\simeq(\mathbf Z/p)^2\simeq E[p]$ over $K$.
The good curve has inertia trace $2$; the bad quadratic twist has trace
$-2$ at the inertia involution. That involution has order prime to $p$.
\end{proof}

\subsection{The first torsion level: a sharp tame ramification range}

\begin{theorem}[B-RD7: the first-level tame range]\label{thm:B-RD7}
\status{Proved; independently audited by the parent agent.}
Let $p\geq5$, and permit arbitrary perfect residue fields of characteristic
$p$. There exists a complete discretely valued $K$ of characteristic zero
and absolute ramification index $e$, and an elliptic curve $E/K$ with
$E[p]\subset E(K)$, additive potentially good reduction, and semistability
defect $m$, if and only if
\[
 m\in\{2,3,4,6\},\qquad m\mid(p+1),\qquad
 (p-1)\mid e,\qquad e\geq\frac{p^2-1}{m}.
\]
In the existence direction one may take the residue field to be
$\overline{\mathbf F}_p$. In particular the set of possible $e$ is
\[
 \left\{e\in(p-1)\mathbf Z_{>0}:
 e\geq\frac{p^2-1}{m_{\max}(p)}\right\},
\]
where $m_{\max}(p)=2,6,4,6$ for $p\equiv1,5,7,11\pmod {12}$,
respectively. For $p=3$, the same assertion with a specified tame defect
$m=2$ or $4$ holds: defect $2$ occurs exactly at even $e\geq4$, and
defect $4$ occurs at every positive even $e$. No assertion about wild
3-primary inertia is included in the latter statement.
\end{theorem}

\begin{proof}
We first prove necessity. Completion of the maximal unramified extension
does not change $e$ or the inertia image. The ordinary case is excluded
by Proposition~\ref{prop:B-RD3}. For $p\geq5$, potential-good inertia is
tame by \cite[Corollary~2 of Theorem~2]{SerreTate1968Good}. Thus the
minimal good-reduction extension $L/K$ is cyclic of order $m$, with
supersingular reduction $C$. Over an algebraically closed field of
characteristic greater than three, writing $C$ in short Weierstrass form
shows that its origin-preserving automorphisms have the form
$(x,y)\mapsto(u^2x,u^3y)$, with $u^4a=a$ and $u^6b=b$.
The possibilities are $m=2,3,4,6$. The curves with $j=0$ are supersingular
exactly when $p\equiv-1\pmod3$, and those with $j=1728$ exactly when
$p\equiv-1\pmod4$. Both assertions follow directly by extracting the
coefficient of $x^{p-1}$ from $(x^3+b)^{(p-1)/2}$ and
$(x^3+ax)^{(p-1)/2}$, respectively. Consequently $m\mid p+1$.
The same discussion for tame inertia in characteristic three gives
$m=2$ or $4$.

Normalize the valuation $v_L$ by $v_L(L^\times)=\mathbf Z$.
An average over the tame cyclic action provides a formal parameter $t$
at the identity of the good model which is an exact eigenvector.
If its eigenvalue is a primitive $m$th root $\zeta^q$, then every
$K$-rational nonzero $p$-torsion point $P$ satisfies
\[
 v_L(t(P))\equiv q\pmod m,\qquad \gcd(q,m)=1.
\]
All these valuations are positive, since $C$ is supersingular. Choose a
point $P$ whose parameter has maximum valuation $A$, and a point $Q$
whose parameter has minimum valuation $B$. If $A>B$ they are independent;
if $A=B$, choose any basis. The $p-1$ nonzero multiples of $P$ have
valuation $A$. The other $p(p-1)$ points have valuation $B$: if $A>B$,
this follows from the linear term of the formal group law; if $A=B$,
it follows from the assumed common maximum and minimum.
The constant term of $[p](t)/t$ is $p$, and its distinguished degree is
$p^2-1$. Weierstrass preparation therefore gives
\[
 m e=v_L(p)=(p-1)(A+pB).                         \tag{B.7.1}
\]
Here $A\geq B\geq1$ and $A\equiv B\equiv q\pmod m$.
Since $p\equiv-1\pmod m$, the displayed identity implies both
$(p-1)\mid e$ and $me\geq p^2-1$.

We give a deformation-theoretic converse with the equivariance made
explicit. Put $k=\overline{\mathbf F}_p$. Choose a supersingular elliptic
curve $C/k$ with an automorphism $\alpha$ of order $m$. For $m=3,6$
choose $j=0$, for $m=4$ choose $j=1728$ (the unique supersingular
$j$ in characteristic three is also zero), and for $m=2$ choose any
supersingular curve. The derivative of $\alpha$ is a primitive $m$th
root; a finite automorphism of order prime to $p$ with derivative one
is the identity, by averaging a local coordinate.

The framed deformation functor of $C$, with full Drinfeld $p$-level
structure, is represented by a complete regular local ring $R$ of
dimension two. If $P,Q$ are the universal labelled basis, their formal
coordinates form a regular system of parameters. These are the
precise statements of \cite[Theorem~2.1.2 and
Proposition~2.2.3(i)]{CarterStrauch2021LubinTate}; Serre--Tate identifies
the height-two formal-group deformation with the elliptic deformation
\cite[Theorem~1.2.1]{Katz1981SerreTate}. The universal formal elliptic
curve algebraizes: the identity gives the compatible ample line
bundles $\mathcal O(3[0])$, so the projective formal existence theorem
applies over the complete noetherian ring $R$.

Changing the framing by $\alpha$ gives an action of $C_m$ on this
functor, hence a semilinear action on the universal elliptic curve.
It preserves the labels $P,Q$. For any initial formal coordinate $t_0$,
replace it by
\[
 t=\frac1m\sum_{i=0}^{m-1}\zeta^{-i}\alpha^{i*}t_0,
\]
where $\zeta$ is the derivative eigenvalue. The linear coefficient is
a unit, so $t$ remains a coordinate and $\alpha^*t=\zeta t$.
Thus $x=t(P)$ and $y=t(Q)$ are regular parameters satisfying exactly
$\alpha(x)=\zeta x$, $\alpha(y)=\zeta y$. The action fixes $W(k)$.
Cohen's presentation theorem gives an equivariant surjection
\[
 W(k)[[X,Y]]\longrightarrow R,\qquad X\mapsto x,\quad Y\mapsto y,
\]
where $\alpha$ multiplies $X,Y$ by $\zeta$.

We detail its single relation, to avoid assuming that an arbitrary
monomial specialization has the intended ramification. The full-level
factorization gives, in $R$,
\[
 p=u\prod_{(a,b)\in\mathbf F_p^2\setminus\{0\}}
          t([a]P+[b]Q),\qquad u\in R^\times.
\]
Lift $u$ to a unit $U$ and each factor to a power series $T_{a,b}$
in $W(k)[[X,Y]]$ by using the multiplication and addition power
series of the formal group, lifting their coefficients. These lifts
have zero constant term and linear term $aX+bY$ (using integral
representatives for $a,b$); crucially, $T_{a,0}$ is divisible by $X$
with unit quotient when $a\ne0$. Then
\[
 F_0=p-U\prod_{(a,b)\ne0}T_{a,b}
\]
lies in the kernel and generates it: it has coefficient one on $p$
modulo the square of the maximal ideal of $W(k)[[X,Y]]$.
Replace $F_0$ by its $C_m$-average $F$. This still generates the kernel,
has the same linear term $p$, and is invariant.

Write $t=e/(p-1)$ and choose
\[
 B=1,\qquad A=1+m\left(t-\frac{p+1}{m}\right).
\]
These are positive integers, $A\geq B$, and $A\equiv B\equiv1\pmod m$.
Reduce the relation modulo $p$, and give $X,Y$ weights $A,B$. If $A>B$, the weighted initial term of the
product in $F_0$ is a nonzero scalar times
$X^{p-1}Y^{p(p-1)}$, of weight
$d=(p-1)(A+pB)$. If $A=B$, the weighted initial term is the nonzero
homogeneous polynomial $\prod_{(a,b)\ne0}(aX+bY)$ of degree $p^2-1$.
In either case that initial form is invariant, since its total degree
is $p^2-1$, divisible by $m$. Averaging therefore preserves it.
Choose a Teichmuller unit $c$ arbitrarily when $A>B$, and with
$\bar c\notin\mathbf F_p$ when $A=B$. Then
\[
 R_L=W(k)[[\pi]]/(F(\pi^A,c\pi^B))
\]
is a complete discrete valuation ring with uniformizer $\pi$ and
$v_L(p)=d$: modulo $p$ its defining equation has first nonzero term
of degree $d$, while its coefficient on $p$ is a unit. This also
proves directly that the specified specialization exists.

The action $\pi\mapsto\zeta\pi$ preserves this equation, by invariance
of $F$ and $A\equiv B\equiv1\pmod m$. It has order $m$, acts trivially
on the residue field, and makes the map $R\to R_L$ equivariant.
Let $L=\operatorname{Frac}R_L$ and $K=L^{C_m}$. Then $L/K$ is tame
of ramification degree $m$, and
\[
 e(K)=d/m=e.
\]
The generic elliptic curve descends to $K$. Since the labels of the
level structure were preserved, its full $p$-torsion is rational over
$K$. The inertia action on the good special fibre is the faithful
order-$m$ action $\alpha$. The Neron--Ogg--Shafarevich criterion gives
exactly semistability defect $m$, so reduction is additive. Finally the
Weil pairing implies $\mu_p\subset K$. This proves every sufficiency
assertion made in the theorem, including $p=3$ and $m=2,4$.
\end{proof}

\subsection{A missing wild semistability defect in characteristic three}

\begin{theorem}[B-RD8: defect twelve with full rational $3$-torsion]\label{thm:B-RD8}
\status{Proved; exact computation recorded.}
The list $2,3,4,6$ in B.2 is incomplete in residue characteristic three.
Put
\[
 K=\mathbf Q_3(i,\zeta),\qquad i^2=-1,\quad
 \zeta^2+\zeta+1=0,\qquad \pi=\zeta-1,\quad t=1+i\pi.
\]
Then $K$ has absolute ramification index $2$ and residue field
$\mathbf F_9$. The elliptic curve
\[
 E:\quad y^2=x^3-27t(t^3+8)x+54(t^6-20t^3-8)
\]
has full $3$-torsion over $K$ and additive potentially supersingular
reduction, with semistability defect exactly $12$.
\end{theorem}
\begin{proof}
Write $a=-27t(t^3+8)$ and $b=54(t^6-20t^3-8)$.
Set $\delta=1+2\zeta=3+2\pi$, so $\delta^2=-3$.
The eight nonzero $3$-torsion points are
\[
 \left(-9t^2,\ \pm12\delta(t^3-1)\right)
\]
and, for $\epsilon\in\{1,\zeta,\zeta^2\}$,
\[
 \left(3t^2+12\epsilon t+12\epsilon^2,
       \ \pm36(\epsilon t^2+\epsilon^2t+1)\right).
\]
Indeed their four $x$-coordinates give the factorization
\[
 3x^4+6ax^2+12bx-a^2
 =3(x+9t^2)\prod_{\epsilon^3=1}
       (x-3t^2-12\epsilon t-12\epsilon^2).
\]
Substitution verifies their $y$-coordinates; the points are distinct
and exhaust the eight nonzero geometric $3$-torsion points.
The exact certificate \texttt{verify\_p3\_defect12.py} performs these
calculations in
$\mathbf Q[i,\pi]/(i^2+1,\pi^2+3\pi+3)$, verifies the curve equation
and the tripling identity for every point, and verifies independence
of points with different $x$-coordinates. No numerical approximation
is used.

Normalize $v_K(\pi)=1$. Direct calculation gives
\[
 \Delta=2^{12}3^9(t^3-1)^3,\quad
 v_K(t^3-1)=v_K(t^3+8)=3,\quad v_K(\Delta)=27.
\]
In particular $E$ is not good over $K$, since the valuation of a
Weierstrass discriminant can change only by a multiple of $12$.
We next both construct good reduction and prove the required minimality.

In the $2$-division polynomial $x^3+ax+b$, substitute
$x=3\pi^2 Z$. Dividing by $(3\pi^2)^3$ gives
\[
 Z^3+A Z+B,
 \qquad A=4i\pi+3\pi+6,\quad B=-12i+20\pi+30.
\]
Both $A$ and $B$ have valuation one. The polynomial is Eisenstein,
so a root $z$ generates a totally ramified cubic extension $K_1/K$,
and $v_K(z)=1/3$. Set $x_0=3\pi^2 z$. Then
\[
 v_K(x_0)=13/3,\qquad v_K(3x_0^2+a)=9.
\]
The second equality can also be read from the derivative
$3z^2+A$, whose valuation is one. Adjoin $u$ satisfying
\[
 u^4=3x_0^2+a.
\]
In the normalized valuation of $K_1$, the right side has value $27$,
prime to four. Thus this is a totally ramified degree-four extension;
for $L=K_1(u)$ we have $e(L/K)=12$ and $v_K(u)=9/4$.
The substitution $x=x_0+u^2X$, $y=u^3Y$ gives
\[
 Y^2=X^3+\frac{3x_0}{u^2}X^2+X.
\]
The middle coefficient has valuation $11/6>0$. The special fibre is
$Y^2=X^3+X$, smooth and supersingular in characteristic three.
Therefore good reduction is obtained with ramification degree twelve.

Conversely, let $F/K$ be any finite extension over which $E$ has good
reduction. Its discriminant valuation implies
$12\mid27e(F/K)$, hence $4\mid e(F/K)$. Since $2$ is invertible,
the $2$-torsion of a good elliptic scheme is finite etale; all its
points are defined over the maximal unramified extension of $F$.
The Eisenstein $2$-division root has $K$-valuation $1/3$ after the
indicated scaling, so $3\mid e(F/K)$. Consequently
$12\mid e(F/K)$. Together with the constructed extension, this proves
that the semistability defect is exactly twelve.
\end{proof}

\subsection{The complete first-level range in residue characteristic three}

\begin{lemma}[B-RD9: a variable-ramification Hesse family]\label{lem:B-RD9}
\status{Proved.}
Let $K$ have algebraically closed residue field of characteristic three,
absolute ramification index $e$, and a uniformizer $\varpi$. Suppose
$\mu_3\subset K$. Choose an integer $r$ with
$0<2r<e$ and $3\nmid e+r$, and put $t=1+\varpi^r$. The Hesse
Weierstrass curve of Theorem~\ref{thm:B-RD8} has full rational
$3$-torsion and semistability defect
\[
 m=\frac{12}{\gcd(4,e+r)}.
\]
\end{lemma}
\begin{proof}
The torsion formulas in Theorem~\ref{thm:B-RD8} are identities in $t$,
so apply unchanged. Write $w=\varpi^r$ and first scale the $2$-division
polynomial by $x=3X$. After translating $X=w^2+Z$, the resulting cubic is
\[
 Z^3+3w^2Z^2-(12w^3+18w^2+36w+27)Z
       +12w^4-36w^3-117w^2-108w-54.
\]
Its coefficients, from constant through quadratic, have valuations
$e+4r,e+3r,e+2r$. The inequalities $2r<e$ make the displayed leading
terms unique. Its Newton polygon has the single slope
$-(e+4r)/3$. Because $3\nmid e+r$, adjoining a root $z$ gives a
completely ramified cubic extension $K_1/K$. Put $x_0=3(w^2+z)$.
Then
\[
 v_K(x_0)=e+2r,\qquad v_K(3x_0^2+a)=3(e+r).
\]
The derivative computation giving the second equality has a unique
lowest term, of valuation $e+3r$, before multiplying by $3^2$.

Choose $u$ with $u^4=3x_0^2+a$. Since the residue field is algebraically
closed and $4$ is invertible, adjoining this root has ramification
degree $4/\gcd(4,e+r)$ over $K_1$. The equation obtained from
$x=x_0+u^2X$, $y=u^3Y$ is
\[
 Y^2=X^3+(3x_0/u^2)X^2+X,
\]
and the middle coefficient has valuation $(e+r)/2>0$. Thus it is good.
Conversely $v_K(\Delta)=9(e+r)$, which forces the tame factor
$4/\gcd(4,e+r)$ in any good-reduction ramification degree. The
$2$-division cubic forces the additional factor three, exactly as in
Theorem~\ref{thm:B-RD8}. This proves minimality.
\end{proof}

\begin{theorem}[B-RD10: all first-level defects at $p=3$]\label{thm:B-RD10}
\status{Proved; independently audited by the parent agent.}
Allowing arbitrary perfect residue fields, a potentially good additive
elliptic curve with full rational $3$-torsion exists at every positive
even absolute ramification index. The precise possible defects are
\[
\begin{array}{c|c}
\text{defect}&\text{absolute ramification indices}\\\hline
2& e\in2\mathbf Z,\ e\geq4\cr
3& e\in2\mathbf Z,\ e\geq4\cr
4& e\in2\mathbf Z,\ e\geq2\cr
6& e\in2\mathbf Z,\ e\geq2\cr
12& e\in2\mathbf Z,\ e\geq2.
\end{array}
\]
Over the particular field $\mathbf Q_3(\zeta_3)$, Theorem~\ref{thm:B-RD2}
already gives an example of defect six.
\end{theorem}
\begin{proof}
The special fibre is supersingular by Proposition~\ref{prop:B-RD3}.
Over an algebraically closed field of characteristic three it is
isomorphic to $y^2=x^3-x$. Its automorphisms are
\[
 (x,y)\longmapsto(u^2x+r,u^3y),\qquad
 u^4=1,\quad r^3-r=0.
\]
Thus its automorphism group has order twelve, is $C_3\rtimes C_4$,
and the only possible nontrivial subgroup orders are $2,3,4,6,12$.
The inertia image embeds in this group.

Over a good-reduction extension of ramification degree $m$, all eight
nonzero $3$-torsion parameters have positive integral valuations.
Their sum is $v_L(3)=me$, by the constant term of $[3](T)/T$.
Consequently $me\geq8$, which rules out $m=2,3$ at $e=2$.
The Weil pairing gives $\mu_3\subset K$, hence $e$ is even.
These are all the necessary restrictions in the table.

The tame cases $m=2,4$ are Theorem~\ref{thm:B-RD7}.
We explain the remaining constructions explicitly. For any desired
even $e$, a convenient base with algebraically closed residue field is
\[
 K_e=\operatorname{Frac}\left(
 W(\overline{\mathbf F}_3)[\varpi]/
 (\varpi^e+3\varpi^{e/2}+3)\right).
\]
The polynomial is Eisenstein, and $1+\varpi^{e/2}$ is a primitive
cube root of unity. In all the cases below the variable-ramification
family is taken over $K_e$.

For defect twelve and $e\geq4$, use $r=1$ unless $3\mid e+1$, in
which case use $r=3$. The second case first occurs at $e=8$, so always
$0<2r<e$. The integer $e+r$ is odd and prime to three. The missing
case $e=2$ is Theorem~\ref{thm:B-RD8}.

For defect six with $e\equiv0\pmod4$, use $r=2$, replacing it by
$r=6$ when $3\mid e+2$. These choices work for every such $e\geq8$.
For $e\equiv2\pmod4$, use $r=4$, replacing it by $r=8$ when
$3\mid e+4$. These choices work except at $e=2,6,14$.
In each applicable case, $e+r\equiv2\pmod4$ and $3\nmid e+r$,
so Lemma~\ref{lem:B-RD9} applies. The case $e=2$ is
Theorem~\ref{thm:B-RD2}. The case $e=4$ is obtained from the
order-twelve example by the tame quadratic extension: its inertia
image then has order six. The case $e=14$ follows from the order-six
example at $e=2$ by a tame degree-seven extension.
For $e=6$, start with the order-six example over a base of algebraically
closed residue field, and adjoin a cube root of a uniformizer.
Its good-reduction cubic subfield is obtained by adjoining a cube root
of $-1/4$, whose Kummer class has valuation zero; the uniformizer class
has valuation one. These cyclic cubic extensions are distinct and
linearly disjoint. Thus the order-six inertia image remains unchanged.
All finite extensions here are totally ramified.

For defect three with $e\equiv0\pmod4$, use $r=4$, replacing it by
$r=8$ when $3\mid e+4$. These choices work for every such $e\geq12$.
For $e\equiv2\pmod4$ and $e\geq6$, use $r=2$, replacing it by
$r=6$ when $3\mid e+2$; only $e=10$ is exceptional.
The examples at $e=4,8$ come from the order-six example at $e=2$
by tame extensions of degree two and four: they remove its order-two
inertia and retain the order-three subgroup.
For $e=10$, work over $K_{10}$ and put $t=1+\varpi^2$.
In the $2$-division cubic scaled by $x=3X$, translate
$X=\varpi^4+\varpi^6+Z$. Its constant, linear and quadratic
coefficients have valuations $22,16,14$, respectively. Its Newton
polygon therefore has the single slope $-22/3$, forcing a ramified
cubic extension. Here $v_K(\Delta)=108$, and the good-model
construction has $v_K(u^4)=36$; its fourth root requires no additional
ramification over the cubic extension. Hence the defect is exactly
three. The exact certificate \texttt{verify\_p3\_ranges.py} verifies these
coefficient valuations by reducing polynomials modulo
$\varpi^{10}+3\varpi^5+3$, and separately audits all the index recipes
for every even $e\leq10000$.

For the assertions about tame base change, one can work throughout
over completed maximal unramified bases. The order-six example has
cyclic inertia $C_6$. The order-twelve example has the full group
$C_3\rtimes C_4$, whose tame quotient is $C_4$; a tame quadratic
extension leaves the order-six subgroup. These observations also
verify that the stated defects are not merely upper bounds.
\end{proof}

\begin{corollary}[B-RD11: the cyclotomic field]\label{cor:B-RD11}
\status{Proved.}
For $p\geq5$, no elliptic curve over $\mathbf Q_p(\zeta_p)$ can have
full rational $p$-torsion and additive potentially good reduction.
For $p=3$, such a curve exists.
\end{corollary}
\begin{proof}
For $p\geq7$, the necessary inequality in
Theorem~\ref{thm:B-RD7} gives
$e\geq(p^2-1)/6>p-1$, contradicting $e(\mathbf Q_p(\zeta_p))=p-1$.
For $p=5$, the only possible defect is six. We justify that the good
extension may be chosen totally ramified, so its residue field really
is $\mathbf F_5$. Choose a minimal integral short Weierstrass equation
$y^2=x^3+ax+b$, and put $\delta=v_K(\Delta)$. Since the reduction is
potentially good, $v_K(j)\geq0$, giving $3v_K(a)\geq\delta$.
The discriminant identity then gives $2v_K(b)\geq\delta$ as well.
Put $m'=12/\gcd(12,\delta)$. Over the totally ramified extension
$K(\varpi^{1/m'})$, scaling by $u=\varpi^{\delta/12}$ gives integral
coefficients and unit discriminant. Conversely the discriminant
forces $m'$ to divide every good-reduction ramification degree.
Thus $m'=6$, and this construction gives the required degree-six
extension with unchanged residue field. Equality now holds in the
valuation bound: over this good-reduction extension every nonzero
$5$-torsion formal parameter has valuation one. Dividing those
parameters by a uniformizer and reducing gives an injective
$\mathbf F_5$-linear map
\[
 E[5](K)\longrightarrow k_L=k_K=\mathbf F_5.
\]
Additivity follows from the linear term of the formal group law;
injectivity follows because every nonzero parameter has valuation
one. Its source would have dimension two and its target dimension
one, a contradiction. The claim for $p=3$ is
Theorem~\ref{thm:B-RD2}.
\end{proof}

\subsection{All torsion levels for tame potentially good elliptic curves}

\begin{theorem}[B-RD12: the sharp tame higher-level criterion]\label{thm:B-RD12}
\status{Proved; independently audited by the parent agent.}
Fix $p\geq5$, $n\geq1$, and $m\in\{2,3,4,6\}$ with $m\mid p+1$.
Assume $e$ is a positive multiple of $p-1$, and put
$s=v_p(e/(p-1))$. There is a complete discretely valued $K$ of absolute
ramification index $e$ and an elliptic curve with full rational
$p^n$-torsion and semistability defect $m$ if and only if
\[
 s\geq n-1,\qquad
 e\geq\frac{p^2-1}{m}\,
        p^{\max(2n-2,\,4n-4-s)}.                  \tag{B.12.1}
\]
Arbitrary perfect residue fields are allowed; the constructions use
$\overline{\mathbf F}_p$. The same assertion holds for $p=3$ with a
specified tame defect $m=2$ or $4$. In that assertion no restriction on
wild defects is claimed here; their sharp higher-level criterion is proved separately in Theorem~\ref{thm:B-RD26-wild}.
\end{theorem}
\begin{proof}
We first obtain an exact integer criterion. Pass to completed maximal
unramified bases and a tame cyclic degree-$m$ good-reduction extension
$L/K$. Use the eigenparameter from Theorem~\ref{thm:B-RD7}.
Let $A$ and $B$ be the maximum and minimum valuations of a primitive
$p^n$-torsion point. Points $P,Q$ achieving them form a basis if $A>B$:
otherwise $Q=aP+pR$ with $a$ a unit; multiplication by $p$ strictly
increases every positive formal valuation, contradicting minimality
of $B$. If $A=B$, choose any basis. We have
\[
 A\geq B\geq1,\qquad A\equiv B\equiv q\pmod m,
 \qquad \gcd(q,m)=1.                              \tag{B.12.2}
\]

The full-level deformation ring has the factorization of
\cite[Proposition~2.2.3(vii)]{CarterStrauch2021LubinTate}.
Iterated $j$ times, it expresses the coordinate of $[p^j]Q$ as a unit
times a product of $p^j$ primitive level-$n$ coordinates, each raised
to $p^j$. The relevant lines have representatives $(a,1)$ with
$a\equiv0\pmod {p^{n-j}}$. All these coordinates have valuation $B$,
because the contribution $[a]P$ has valuation strictly greater than
$B$. Consequently
\[
 v_L(t([p^j]Q))=p^{2j}B\quad(0\leq j\leq n-1).   \tag{B.12.3}
\]
This factorization applies to the actual good elliptic curve by its
map from the universal framed deformation ring with full level.

Primitive vectors modulo scalar units are the lines of
$\mathbf P^1(\mathbf Z/p^n)$. There are $p^n$ representatives $(a,1)$,
all of valuation $B$: when $A=B$, a cancellation would contradict the
chosen maximum $A$. There is the line $(1,0)$, of valuation $A$.
For $1\leq j<n$, there are $(p-1)p^{n-j-1}$ representatives $(1,b)$
with $v_p(b)=j$, each of valuation $\min(A,p^{2j}B)$. Again equality
cannot cause cancellation beyond $A$, since $A$ is the maximum among
primitive points. The series $[p^n](T)/[p^{n-1}](T)$ has constant term
$p$, distinguished degree $(p^2-1)p^{2n-2}$, and precisely the
primitive $p^n$-torsion roots. Taking its constant term gives
\[
 me=(p-1)p^{n-1}\left(
 A+p^nB+\sum_{j=1}^{n-1}(p-1)p^{n-j-1}
                \min(A,p^{2j}B)\right).          \tag{B.12.4}
\]

We prove that every pair of integers satisfying (B.12.2) and (B.12.4)
can be realized; for the stated range it will suffice to use $B=q=1$.
Use the full Drinfeld $p^n$-level deformation ring $R_n$, whose regular
parameters $x,y$ are the universal basis coordinates
\cite[Theorem~2.1.2 and Proposition~2.2.3(i)]{CarterStrauch2021LubinTate}.
The eigenparameter averaging and equivariant Cohen presentation of
Theorem~\ref{thm:B-RD7} apply verbatim to this level.
We justify the weighted specialization, including its possible
cancellations, as follows.

Work modulo $p$ in the power-series presentation and give $X,Y$
weights $A,B$. Lift the full-level factorization (B.12.3), rather than
an arbitrary expression for $[p^j]y$. Each factor with representative
$(a,1)$ and $a\equiv0\pmod {p^{n-j}}$ has weighted initial form $Y$,
up to a nonzero constant. Therefore a lift of $[p^j]y$ has weighted
initial form $c_jY^{p^{2j}}$, with $c_j\in k^\times$. It follows that
for a primitive line $(1,b)$, $v_p(b)=j$, one can choose a lift whose
initial form is the lower-weight term of
\[
 X+c_bY^{p^{2j}},\qquad c_b\ne0,
\]
or that entire binomial when the two weights agree. The other
primitive lines are handled by their linear terms. These choices are
legitimate lifts because the factorization identities hold in $R_n$;
changes by the defining relation do not alter their images.

The constant-term identity for $[p^n](T)/[p^{n-1}](T)$ gives a generator
of the Cohen kernel of the form $p$ minus a unit times the product of
all primitive coordinates. Average this generator under $C_m$.
Its weighted initial form modulo $p$ is preserved: every initial
factor has scalar character $\zeta$, since
$p^{2j}\equiv1\pmod m$, and the total number
$(p^2-1)p^{2n-2}$ of factors is divisible by $m$.
Now specialize $X=\pi^A$, $Y=c\pi^B$, choosing the residue of $c$
to avoid the finitely many roots of the displayed initial forms.
The resulting relation has linear term $p$ and first nonzero term
modulo $p$ of degree $d$, where $d$ is the right side of (B.12.4).
It defines a complete DVR with $v_L(p)=d$.
Every primitive-coordinate weight is strictly smaller than $d$, so
terms in its chosen lift divisible by $p$ cannot change that weight.
Thus all the asserted valuations hold after specialization, with no
circular assumption on $v_L(p)$.

For $A\equiv B\equiv1\pmod m$, the action $\pi\mapsto\zeta\pi$
preserves the relation and the specialization. Tame descent gives a
curve over the invariant field, of absolute ramification $d/m$,
with full rational $p^n$-torsion and exact defect $m$, as in
Theorem~\ref{thm:B-RD7}. For a general common residue $q$ in
(B.12.2), replace the action on $\pi$ by multiplication by
$\zeta^{q^{-1}}$, where the inverse is taken modulo $m$.

It remains to simplify (B.12.4), which also proves that $B=1$ suffices.
Choose $r\in\{0,\ldots,n-1\}$ so that
$p^{2r}B\leq A\leq p^{2r+2}B$, with the upper bound omitted for
$r=n-1$. Summing the geometric series in (B.12.4) gives
\[
 me=(p-1)p^{2n-r-2}(A+p^{2r+1}B).                \tag{B.12.5}
\]
It implies the two conditions
\[
 p^{2n-r-2}(p-1)\mid e,\qquad
 e\geq \frac{p^2-1}{m}p^{2n+r-2}.                \tag{B.12.6}
\]
Conversely suppose these conditions hold for some $r$. Increase $r$
until it is $n-1$, or until the second inequality fails for $r+1$;
the divisibility remains true. Put
\[
 B=1,\qquad
 A=\frac{me}{(p-1)p^{2n-r-2}}-p^{2r+1}.
\]
Then $A$ is an integer, $A\equiv1\pmod m$, and
$p^{2r}\leq A<p^{2r+2}$ unless $r=n-1$, when only the lower bound
is needed. Hence it realizes (B.12.4) by the construction above.
Finally the least permissible $r$ is
$\max(0,2n-2-s)$; such an $r\leq n-1$ exists exactly when
$s\geq n-1$. Substitution gives (B.12.1).
\end{proof}

\begin{corollary}[B-RD13: the least good-forcing level in the tame case]
\label{cor:B-RD13}
\status{Proved; independently audited by the parent agent.}
Let $p\geq5$, $e\in(p-1)\mathbf Z_{>0}$, and let $m_{\max}(p)$ be
as in Theorem~\ref{thm:B-RD7}. If $m_{\max}(p)e<p^2-1$, put $N=0$.
Otherwise put
\[
 s=v_p(e/(p-1)),\qquad
 L=\left\lfloor\log_p\frac{m_{\max}(p)e}{p^2-1}\right\rfloor,
 \qquad
 N=1+\min\left(s,\left\lfloor\frac L2\right\rfloor,
                     \left\lfloor\frac{L+s}{4}\right\rfloor\right).
\]
Among potentially good elliptic curves over all such fields of
absolute ramification $e$, the least positive level $n$ at which
full rational $p^n$-torsion forces good reduction is $N+1$.
For $p=3$, the same formula with $m_{\max}=4$ gives the least level
within the class of potentially good curves with tame inertia image.
\end{corollary}
\begin{proof}
For a fixed level, the largest allowed defect gives the weakest
inequality in Theorem~\ref{thm:B-RD12}. Its inequalities are equivalent
to $n-1\leq s$, $2n-2\leq L$, and $4n-4-s\leq L$.
Solving them gives exactly the displayed largest bad level $N$.
The construction realizes it, and its lower levels, while every
potentially good curve outside the bad cases is good.
\end{proof}

\begin{remark}[The scope of ``forces good reduction'']\label{auto:B:b9f011fa8019}
\status{Proved under stated hypotheses}

The potentially good hypothesis is essential. If $K$ contains
$\mu_{p^n}$, a split Tate curve with parameter $q=a^{p^n}$,
$0<|a|<1$, has full rational $p^n$-torsion and split multiplicative
reduction. Thus the preceding least-level assertion is not a
criterion for all elliptic curves without that hypothesis.
\end{remark}

\subsection{Cyclotomic bounds in every dimension and a character lower bound}

\begin{lemma}[B-RD14: full torsion forces the cyclotomic extension]
\label{lem:B-RD14}
\status{Proved.}
Let $A/K$ be an abelian variety of positive dimension over a
characteristic-zero field. If $A[p^n]$ is rational over $K$, then
$\mu_{p^n}\subset K$. A principal polarization, or a polarization of
degree prime to $p$, is not needed for this implication.
\end{lemma}
\begin{proof}
Choose any polarization defined over $K$. Its Weil pairing gives a
nonzero $G_K$-equivariant integral alternating form
\[
 b:T_pA\times T_pA\longrightarrow\mathbf Z_p(1).
\]
Divide all its coefficients by their largest common power of $p$.
The resulting form is still integral and equivariant, and at least
one coefficient is a unit. Since $G_K$ acts trivially on $T_pA/p^n$,
equivariance yields
$(\chi_p(\sigma)-1)b(x,y)\equiv0\pmod {p^n}$ for every $x,y$.
Choose a pair with unit value; then
$\chi_p(\sigma)\equiv1\pmod {p^n}$ for every $\sigma\in G_K$.
This is precisely the claimed cyclotomic containment.
\end{proof}

\begin{proposition}[B-RD15: the unconditional reading of the least-level question]
\label{prop:B-RD15}
\status{Proved; the final implication is explicitly vacuous.}
Fix $p\geq3$, a positive multiple $e$ of $p-1$, and $g\geq1$, and put
$s=v_p(e/(p-1))$. Among \emph{all} abelian varieties of dimension $g$
over all complete discretely valued fields with perfect residue field
and absolute ramification $e$, the least positive $n$ for which full
rational $p^n$-torsion implies good reduction is $s+2$.
At this level full rational $p^n$-torsion cannot occur at all.
Consequently a substantive good-reduction threshold must specify the
potentially good class, as in Corollary~\ref{cor:B-RD13}.
\end{proposition}
\begin{proof}
Lemma~\ref{lem:B-RD14} and the ramification degree
$[\mathbf Q_p(\mu_{p^n}):\mathbf Q_p]=p^{n-1}(p-1)$ imply
$p^{n-1}(p-1)\mid e$. Thus full rational $p^n$-torsion is impossible
for $n\geq s+2$.

Conversely, for each $1\leq n\leq s+1$, choose a complete field $K$
with residue field $\overline{\mathbf F}_p$, absolute ramification $e$,
and $\mu_{p^n}\subset K$. Such a field is obtained by taking the
cyclotomic extension of $\operatorname{Frac}W(\overline{\mathbf F}_p)$
and a further totally ramified extension of the required degree.
By \cite[Theorem~1, p.6 of the author version]{Tate1995Uniformization},
a split Tate curve with parameter $q=a^{p^n}$, where $a$ is a
uniformizer, has all its $p^n$-torsion represented by
$\zeta_{p^n}^i a^j$ in $K^\times/q^{\mathbf Z}$, and has split
multiplicative reduction. For dimension $g$, take its product with
$g-1$ canonical ordinary elliptic lifts. Their $p$-divisible groups
are $\mu_{p^\infty}\oplus\mathbf Q_p/\mathbf Z_p$, so they too have
full rational $p^n$-torsion. The product still has bad reduction.
This proves the asserted least level and explains its vacuity.
\end{proof}

\begin{proposition}[B-RD16: necessary lower bounds for character determination]
\label{prop:B-RD16}
\status{Proved.}
Let $h\geq2$. For $p\geq5$ and absolute ramification $e$, let $N$
be the largest bad elliptic level from Corollary~\ref{cor:B-RD13}.
If $N\geq1$, any universal least index $\nu$ in B.4 must satisfy
$\nu\geq N$. More precisely, for each $1\leq n\leq N$ there are two
height-$h$ BT groups satisfying the manuscript's monodromy condition,
with isomorphic $p^n$-torsion, whose characters on inertia elements of
order prime to $p$ differ. For $p=3$, the same lower bound holds using
the tame value $m_{\max}=4$ in that corollary.
\end{proposition}
\begin{proof}
Take the elliptic curve $E/K$ supplied by Theorem~\ref{thm:B-RD12}
at level $n$, and a canonical ordinary elliptic lift $E'/K$.
Lemma~\ref{lem:B-RD14} gives $\mu_{p^n}\subset K$, so both
$E[p^n]$ and $E'[p^n]$ are constant of rank $p^{2n}$.
Their inertia characters differ: the ordinary good curve has trace
two, whereas the nontrivial tame automorphism of the supersingular
good reduction has trace different from two. Append the same constant
etale BT group of height $h-2$ to both. Their characters each increase
by $h-2$, preserving the difference, and their truncated groups remain
isomorphic. These groups are potentially good, hence satisfy the
stated monodromy condition without using any extension argument.
Finally the manuscript's $M(\nu)$ is $M[p^{\nu+1}]$; failure at
$n=N$ gives $\nu\geq N$.
\end{proof}

\subsection{The monodromy condition and the Fontaine comparison}

\begin{proposition}[B-RD4: Lead B' needs an additional hypothesis]\label{prop:B-RD4}
\status{Refuted as stated.}
For every odd prime $p$, there is a height-two Barsotti--Tate group over
$K=\mathbf Q_p(\zeta_p)$ satisfying the manuscript's monodromy condition
whose rational Tate module is not potentially Hodge--Tate, and therefore
is not potentially semistable. Its manuscript inertia character is the
constant function $2$. Thus the monodromy condition in 9:11 cannot be
identified with potential semistability, and Fontaine's usual
Weil--Deligne representation need not be defined for its Tate module.
\end{proposition}
\begin{proof}
Write $\chi:G_K\to\mathbf Z_p^\times$ for the cyclotomic character.
Since $\mu_p\subset K$, it takes its values in $1+p\mathbf Z_p$. Put
\[
 \rho(g)=\begin{pmatrix}1&p^{-1}\log\chi(g)\\0&1\end{pmatrix}.
\]
This is a continuous representation on $T=\mathbf Z_p^2$, because the
logarithm is additive on $1+p\mathbf Z_p$ and takes values in
$p\mathbf Z_p$. Over a characteristic-zero field, the equivalence
between finite etale group schemes and finite continuous Galois modules
applied to $T/p^nT$ gives compatible finite etale group schemes $M[p^n]$.
The transition maps are induced by multiplication by $p$ on the lattices,
and the resulting inductive system is a height-two Barsotti--Tate group.
The invariant first coordinate and the quotient both give the constant
group $\mathbf Q_p/\mathbf Z_p$. They extend over $\mathcal O_K$, so the
specified two-step filtration satisfies the manuscript condition and
its character is $1+1=2$.

For $V=T\otimes\mathbf Q_p$, the Sen operator is
\[
 \Theta_V=\frac1p\begin{pmatrix}0&1\\0&0\end{pmatrix}.
\]
Indeed the logarithm of a unipotent matrix $1+aN$, with $N^2=0$, is
$aN$, and division by $\log\chi(g)$ gives the displayed operator.
It remains nonzero and nilpotent after restriction to any finite
extension. Sen's criterion says that a representation is Hodge--Tate
exactly when this operator is diagonalizable with integral eigenvalues;
the identical unscaled example and the criterion are explicitly given
in \cite[\S II.1.2, p.~9 of the author's version]{Berger2004Introduction}.
The operator is not diagonalizable. Potentially semistable
representations are potentially Hodge--Tate, so no finite extension
makes this representation semistable.
\end{proof}

\begin{proposition}[The corrected Fontaine dictionary]
\label{prop:B-RD4-dictionary}
\status{Proved with the additional potential semistability hypothesis.}
Let $M$ satisfy the manuscript monodromy condition, and assume in
addition that $V=T_p(M)\otimes_{\mathbf Z_p}\mathbf Q_p$ is potentially
semistable. For arbitrary perfect residue field, its manuscript
character is the inertia trace of Fontaine's $p$-module of Deligne
$D_{\mathrm{pst}}(V)$. When the residue field is finite, this is the
inertia character of the usual Weil--Deligne representation. With
\[
 D_{\mathrm{cov}}(G)=D_{\mathrm{contra}}(G^D),
\]
the precise comparison for a good graded group $G$ over a finite
extension $L/K$ is
\[
 D_{\mathrm{cov}}(G_{k_L})[1/p]
       \simeq D_{\cris,L}(T_pG\otimes\mathbf Q_p)(-1).
                                                        \tag{B-RD4.1}
\]
The Tate twist in this formula does not change the inertia trace.
\end{proposition}
\begin{proof}
Choose a finite Galois extension $L/K$ over which $V$ is semistable
and every graded group has good reduction. All subrepresentations
and quotients in the filtration are semistable over $L$.
Crystalline comparison for a BT group, with the indicated
covariant-versus-contravariant convention, gives (B-RD4.1)
functorially and compatibly with descent; see
\cite[Lemma~2.2.4 and Corollary~2.2.6, author pp.26--27; Appendix~A.2
and~A.6]{Kisin2006Crystalline}. One can check the convention on
$\mathbf Q_p/\mathbf Z_p$ and $\mu_{p^\infty}$: their covariant
Frobenius slopes are respectively one and zero, whereas the
crystalline Tate-module slopes are zero and minus one.

The semistable period functor is exact on semistable
representations. Thus its filtration has graded terms the
crystalline period modules of the good graded groups, and traces
of every finite descent-inertia element add. Exactness and the
passage from the completed maximal unramified extension back to
$K$ are treated in \cite[\S\S5.1.5--5.1.7,
pp.156--158]{Fontaine1994Semi}; alternatively the exact
$D_{\mathrm{pst}}$ functor is stated in
\cite[\S2.3.4, p.337]{Fontaine1994WD}.
The crystalline Tate twist is unramified in this period-module
sense, so the right side of (B-RD4.1) has the same inertia trace as
its untwisted version. Summing proves the assertion.

Over an arbitrary perfect residue field the natural object is the
$p$-module of Deligne, with semilinear Galois action, Frobenius and
monodromy; its definition is
\cite[\S\S1.1.1--1.1.2, pp.322--323]{Fontaine1994WD}.
If the residue field is finite, the corresponding linear
Weil--Deligne action is
$\rho_0(w)=\varphi^{-a(w)}\rho(w)$, as in
\cite[\S1.3.5, p.329; \S2.3.7, pp.339--340]{Fontaine1994WD}.
For inertia $a(w)=0$, so the trace just computed is exactly the
Weil--Deligne inertia trace. No finite-residue-field hypothesis is
needed for the preceding Deligne-module statement.
\end{proof}

\begin{proposition}[B-RD5: the height-one case of B.4]\label{prop:B-RD5}
\status{Proved.}
Suppose $\mu_p\subset K$. For height $h=1$ and every prime $p$, the manuscript
character on prime-to-$p$ inertia elements is determined by $M(0)$.
The least index in its indexing is therefore $\nu=0$.
\end{proposition}
\begin{proof}
For $p=2$, the finite special inertia action is contained in
$(\mathbf Q_2^\times)_{\mathrm{tors}}=\{1,-1\}$. Every odd-order
inertia element therefore acts trivially. Its requested character is
identically one, so $\nu=0$ suffices. We now assume $p\geq3$.

A height-one group satisfying the monodromy condition itself acquires
good reduction. Its good model has dimension zero or one. In dimension
zero it is etale, and in dimension one its Cartier dual is etale.
Over the completed maximal unramified extension of a good-reduction
field its model is respectively $\mathbf Q_p/\mathbf Z_p$ or
$\mu_{p^\infty}$. Its finite descent inertia character therefore takes
values in the torsion subgroup of $\mathbf Z_p^\times$, namely
$\mu_{p-1}$. In dimension zero, reduction modulo $p$ gives precisely
this character. In dimension one, the Tate module acquires additionally
the cyclotomic factor, whose reduction is trivial under $\mu_p\subset K$.
Thus in both cases the finite character is the unique Teichmuller lift
of the one-dimensional mod-$p$ inertia character. The same conclusion
holds after summing characters of height-one potentially good graded
pieces: the trace on a prime-to-$p$ element is the Brauer character of
the mod-$p$ representation, and is independent of the chosen filtration.
\end{proof}


\subsection{A finite bound for the manuscript character in every height}

\begin{lemma}[Comparison after a common good-reduction extension]
\label{lem:B-RD17-common}
\status{Proved under stated hypotheses}
Let $M,N$ be height-$h$ BT groups over $K$, each equipped with a
filtration by BT subgroups whose graded pieces have good reduction
over the same finite Galois extension $L/K$. Let $c$ be an integer
such that $p^c$ times every generic-fibre homomorphism between finite
flat commutative $p$-group schemes over $\mathcal O_L$ extends over
$\mathcal O_L$. If
\[
 M[p^n]\simeq N[p^n],\qquad n>2ch^2,
                                                        \tag{B-RD17.1}
\]
then their manuscript characters agree on all prime-to-$p$ inertia
elements. No good-reduction assertion about the extensions between
the graded pieces is assumed.
\end{lemma}
\begin{proof}
Use the given isomorphism to identify the two generic finite groups
with $X$. The two BT filtrations induce filtrations
\[
 0=X_0\subset\cdots\subset X_r=X,
 \qquad 0=Y_0\subset\cdots\subset Y_s=X.
\]
Write $h_i$ and $h'_j$ for the heights of their respective graded BT
groups. Thus $\sum_i h_i=\sum_j h'_j=h$.
The Zassenhaus refinement has common generic pieces
\[
 Q_{ij}=
 \frac{X_i\cap Y_j}
 {(X_{i-1}\cap Y_j)+(X_i\cap Y_{j-1})}.
                                                        \tag{B-RD17.2}
\]
In the good model of the $i$th graded group at level $p^n$, take
schematic closures of the subgroup filtration induced by the
$Y_j$. Consecutive quotients give finite flat models $G_{ij}$ of
$Q_{ij}$. Refine the good model of the $j$th graded group from $N$
in the same way, obtaining another model $H_{ij}$. These quotients
need not be truncated BT groups; only their finite flatness is used.
All closures and quotients are stable under the descent action.
The descent actions on the good BT models exist uniquely by Tate's
extension theorem for homomorphisms over a normal base.

Let $k_L$ be the residue field, and extend it algebraically if
necessary. The special Dieudonne module of $G_{ij}$ is a subquotient
of the free $W(k_L)/p^n$-module belonging to the $i$th graded BT
model. Consequently it can be generated by at most $h_i$ elements
over the DVR $W(k_L)$. Similarly the module for $H_{ij}$ has at most
$h'_j$ generators. Here exactness of Dieudonne theory over a perfect
field is applied only to special fibres of exact finite flat
sequences.

The generic identification of $G_{ij}$ and $H_{ij}$, and its inverse,
extend after multiplication by $p^c$. The two composites are
$p^{2c}$, because this equality holds generically and the source
coordinate algebras are flat. Uniqueness of these extensions makes
them compatible with inertia. Fix a prime-to-$p$ element $\gamma$
of inertia, of order $m$. After extending the perfect residue field
so that all $m$th roots of unity are present, the integral
idempotents
\[
 \frac1m\sum_{a=0}^{m-1}\zeta^{-a}\gamma^a
\]
split every special Dieudonne module into $\zeta$-eigensummands.
On a fixed eigensummand, the kernel of either comparison map is
annihilated by $p^{2c}$, and its cokernel is also annihilated by
$p^{2c}$. A module with at most $d$ generators over a DVR has
$p^{2c}$-torsion of length at most $2cd$, and any quotient killed
by $p^{2c}$ has length at most $2cd$. It follows that the two
$\zeta$-eigenlengths differ by at most
\[
 2c\max(h_i,h'_j)\leq2ch_i h'_j.
                                                        \tag{B-RD17.3}
\]
(The last inequality uses that all nonzero graded heights are
positive integers.)

Exactness along each refinement adds these lengths. A good BT group
of height $d$ has a free integral Dieudonne module of rank $d$, and
its $p^n$-torsion contributes $n$ times the multiplicity of each
$\gamma$-eigenvalue in that module. Thus, if $a_\zeta,b_\zeta$ are
the multiplicities in the sums of graded isocrystals for $M,N$,
\[
 \bigl|n(a_\zeta-b_\zeta)\bigr|
 \leq 2c\sum_{i,j}h_i h'_j=2ch^2.
\]
Under (B-RD17.1), the integer $a_\zeta-b_\zeta$ is zero. This holds
for every eigenvalue and every prime-to-$p$ inertia element, proving
the assertion. In particular, applying the same argument to two
filtrations of a single $M$, with arbitrarily large $n$, proves
independence of the prime-to-$p$ character from the chosen filtration.
\end{proof}

\begin{lemma}[A uniform ramification degree for the graded pieces]
\label{lem:B-RD17-degree}
\status{Proved under stated hypotheses}
Put $d=4h^2$ and
\[
 D_{p,h}=
 \begin{cases}
  |\GL_d(\mathbf F_p)|,&p>2,\\
  |\GL_d(\mathbf Z/4\mathbf Z)|
     =2^{d^2}|\GL_d(\mathbf F_2)|,&p=2.
 \end{cases}
\]
For two height-$h$ groups satisfying the monodromy condition, their
graded pieces have simultaneous good reduction over a finite Galois
extension $L/K$ with ramification index at most $D_{p,h}$.
\end{lemma}
\begin{proof}
Choose any finite Galois good-reduction extension $L_0/K$ for all
the graded pieces. Let $I_0$ be its inertia group. Their direct sum
has height $2h$; denote its special Dieudonne isocrystal by $D$.
The image $H$ of $I_0$ lies in
$\End_F(D)^\times$, where $F$ denotes semilinear Frobenius.
This endomorphism algebra has $\mathbf Q_p$-dimension at most
$(2h)^2=d$. Here is an elementary justification not requiring a
classification of isocrystals. In any semilinear Frobenius module,
linearly independent fixed vectors over the fixed field
$W(k_{L_0})[1/p]^{F}=\mathbf Q_p$ remain linearly independent over
$W(k_{L_0})[1/p]$: normalize a shortest relation to have one
coefficient $1$, apply Frobenius, and subtract. A shorter relation
would result unless every coefficient were fixed. Apply this to
conjugation by Frobenius on the endomorphism module, of rank $d$.

Left multiplication embeds $H$ in a general linear group over
$\mathbf Q_p$ of degree at most $d$. A finite group preserves a
$\mathbf Z_p$-lattice, obtained by summing the translates of any
lattice. For $p>2$, the kernel of reduction modulo $p$ in
$\GL_d(\mathbf Z_p)$ is torsion free. For $p=2$, the kernel of
reduction modulo $4$ is torsion free. The $p$-adic logarithm, or the
binomial valuation argument, proves both assertions. Thus
$|H|\leq D_{p,h}$.

Let $J$ be the kernel of $I_0\to H$. It is normal in
$\operatorname{Gal}(L_0/K)$, and $L=L_0^J$ is Galois over $K$ with
$e(L/K)=|H|$. Each rational graded Tate module is crystalline over
$L$. Indeed crystalline comparison identifies its special
isocrystal, with its descent action, with the crystalline period
module, up to the dual or Tate twist dictated by the convention.
The twist has trivial inertia action. The group $J$ acts trivially,
and $L_0,L$ have the same residue field. Hence
\[
 D_{\cris,L}(V)=D_{\cris,L_0}(V)^J
\]
has full dimension. The period-invariant and unramified-completion
statements used here are also explicitly treated in
\cite[\S\S5.1.5--5.1.7, pp.156--158]{Fontaine1994Semi}.
The required comparison, including functoriality,
is given in \cite[Lemma~2.2.4 and the proof of Corollary~2.2.6]{Kisin2006Crystalline}.
Corollary~2.2.6 of that paper, valid for arbitrary perfect residue
field and every prime, now gives a good BT group with this rational
Tate module.

For completeness, the actual graded lattice, rather than just its
rational representation, also has a good model. If $T_0$ is the Tate
lattice of one good model, rescale the desired lattice $T$ so that
$p^aT_0\subset T\subset T_0$. Take the schematic closure of
$T/p^aT_0$ in the finite flat group of $p^a$-torsion, and quotient the
good BT group by this finite flat subgroup. Its Tate lattice is
$p^{-a}T$, which is isomorphic to $T$. Finite flat quotients of BT
groups give BT groups, level by level. This argument uses only the
rational existence theorem at $p=2$, and does not assume an integral
$p=2$ equivalence in Kisin's Theorem~2.2.7.
\end{proof}

\begin{theorem}[B-RD17: a finite universal index exists]
\label{thm:B-RD17}
\status{Proved; independent proof and primary-source audit passed.}
For every prime $p$, every absolute ramification index $e$, and every
height $h\geq1$, a finite universal character-determination index in
B.4 exists. With $D_{p,h}$ as above, set
\[
 q=\left\lceil\log_p(eD_{p,h})\right\rceil,
 \qquad C=(2q+2)(q+2).
\]
An explicit sufficient index in the manuscript's convention is
\[
                       \nu=2Ch^2.                    \tag{B-RD17.4}
\]
In other words, the prime-to-$p$ manuscript character is determined
by $M[p^{\nu+1}]$ at this index. This bound is not asserted to be
sharp. For every prime and $h=1$, Proposition~\ref{prop:B-RD5} gives
the sharp value $\nu=0$; Proposition~\ref{prop:B-RD18} supplies
nonzero lower bounds for every admissible ramification in higher
height.
\end{theorem}
\begin{proof}
Choose $L$ as in Lemma~\ref{lem:B-RD17-degree}. Its absolute
ramification $e_L$ satisfies $e_L\leq eD_{p,h}$.
By \cite[Corollary~2 and the bound following Theorem~1, p.2 of the
author version]{VasiuZink2012Boundedness}, the generic Hom extension
exponent can be chosen to satisfy
\[
 c\leq(\log_p e_L+v_p(e_L)+2)(v_p(e_L)+2)
       \leq(2q+2)(q+2)=C.
\]
At level $n=2Ch^2+1$, Lemma~\ref{lem:B-RD17-common} applies.
Since $M(\nu)=M[p^{\nu+1}]$, this is precisely (B-RD17.4).
\end{proof}

\begin{remark}[The remaining sharp problem]\label{auto:B:5338697ffdba}
\status{Proved under stated hypotheses}

Theorem~\ref{thm:B-RD17} excludes the alternative that no finite
index exists. It does not determine the least index in height
$h\geq2$, and does not turn the manuscript monodromy condition into
potential semistability. The logarithmic cyclotomic extension in
Proposition~\ref{prop:B-RD4} remains a counterexample to that claim.
The refinement proof above works directly with good graded pieces
and therefore includes that extension and analogous nongood
extensions.
\end{remark}


\subsection{Stronger lower bounds from formal groups}

\begin{proposition}[B-RD18: every admissible ramification has a height-two counterexample]
\label{prop:B-RD18}
\status{Proved; independent construction and determinant audit passed.}
Let $p$ be any prime, let $e$ be a positive multiple of $p-1$, and
put $s=v_p(e/(p-1))$. Suppose $n\geq1$ satisfies
\[
 s\geq n-1,\qquad
 e\geq(p-1)p^{\max(2n-2,\,4n-4-s)}.                \tag{B-RD18.1}
\]
Then there are a complete discretely valued field $K$ with
algebraically closed residue field of characteristic $p$, absolute
ramification $e$, and $\mu_p\subset K$, and two height-two BT groups
$M,N$ over $K$ with
\[
 M[p^n]\simeq N[p^n],
\]
whose characters on prime-to-$p$ inertia differ. Both groups
individually acquire good reduction. The group $M$ can be taken to
have tame good-reduction defect $p+1$, and $N$ to be constant etale.
The same examples in every height $h\geq2$ are obtained by adding
constant etale groups.
\end{proposition}
\begin{proof}
Choose a height-two Honda formal group over
$k=\overline{\mathbf F}_p$. Its endomorphism ring contains
$W(\mathbf F_{p^2})$, so the Teichmuller units give a tame
order-$p+1$ automorphism $\alpha$ whose tangent eigenvalue $\zeta$
is primitive of that order. The two eigenvalues on its isocrystal
are $\zeta,\zeta^p=\zeta^{-1}$.

Apply the full-level Drinfeld deformation-ring construction in the
proof of Theorem~\ref{thm:B-RD12}, with $m=p+1$. Every ingredient
before the elliptic algebraization step involves only the formal
height-two group: the regular full-level ring, the two torsion
coordinates, averaging for the tame action, the product formula
for $p$, its valuation regions, and the equivariant discrete
valuation specialization. The condition $m\mid p+1$ holds, and its
arithmetic existence condition is exactly (B-RD18.1), because
$(p^2-1)/m=p-1$.

Stop at that formal-group construction and descend its generic BT
group along the cyclic tame extension. Over a characteristic-zero
field this is ordinary Galois descent for the compatible finite
etale groups at all levels. The full level structure is fixed by
the descent, so $M[p^n]$ is constant over $K$. The special descent
action is $\alpha$, hence is faithful of order $p+1$. Its trace at a
generator is $\zeta+\zeta^{-1}\ne2$.
The constant etale height-two group $N$ has the same constant
$p^n$-torsion and has inertia trace two. Adding constant etale
height $h-2$ adds $h-2$ to both traces and preserves their difference.

The proof also applies at $p=2$: the chosen order is three, so every
averaging denominator remains a unit; all Drinfeld product and
valuation formulas used in the construction hold for that prime.
Finally the cyclotomic containment can be checked directly, without
assuming a polarization of $M$. Over the good extension, the
rank-one crystalline representation $\det(T_pM)\otimes\mathbf Q_p$
has the single Hodge--Tate weight one. Kisin's
Corollary~2.2.6 realizes it by a height-one good group, which over
this strictly henselian complete base is multiplicative. Thus
$\det(T_pM)\chi_p^{-1}$ is a finite character. Its inertia type is
$\det(\alpha)=\zeta^{p+1}=1$, so it is unramified, hence trivial
because the residue field is algebraically closed. Therefore
$\det(T_pM)=\chi_p$. The full rational $p^n$-torsion of $M$ forces
$\chi_p\equiv1\pmod{p^n}$, proving $\mu_{p^n}\subset K$.
\end{proof}

\begin{corollary}[An explicit general-height lower bound]
\label{cor:B-RD18}
\status{Proved under stated hypotheses}
Let $h\geq2$, $\mu_p\subset K$, and define
\[
 L=\left\lfloor\log_p\!\left(\frac e{p-1}\right)\right\rfloor,
 \qquad
 N=1+\min\left\{s,\left\lfloor\frac L2\right\rfloor,
                         \left\lfloor\frac{L+s}4\right\rfloor\right\}.
\]
Any universal least index in B.4 satisfies $\nu\geq N$.
In particular $N\geq1$, so for every admissible $e$ and every
$h\geq2$, the first truncation $M(0)$ does not determine the
requested character. This is stronger than restricting the lower
bound to elliptic examples: the tame automorphism of a formal
height-two group need not arise from an automorphism of an
elliptic curve.
\end{corollary}
\begin{proof}
The integer equivalence between (B-RD18.1) and $n\leq N$ is the
same calculation as in Corollary~\ref{cor:B-RD13}, with
$m=p+1$. Failure at level $p^N=M(N-1)$ gives $\nu\geq N$.
\end{proof}

\begin{remark}[Finite arithmetic certificate]\label{auto:B:c9c8e4df645d}
\status{Proved under stated hypotheses}

The proof above is symbolic. As a separate finite check,
\path{code/verify_honda_bounds.py} enumerates the exact integer-pair
formula for $p=2,3,5,7,11,13,17,19$, levels $n=1,2,3$, and
$e\leq1200$, always with $m=p+1$. It verifies all 24 range tables,
3803 constructed pairs, and 11088 least-level formula comparisons
through $e\leq5000$. All checks pass; the recorded run took
$0.508$ seconds. This certificate verifies the integer arithmetic,
not the formal deformation theory.
\end{remark}


\subsection{A wild restriction and a sharp minimum in characteristic three}

\begin{lemma}[B-RD19: additive full odd torsion is potentially good]
\label{lem:B-RD19}
\status{Proved.}
Let $p$ be odd and let an elliptic curve $E/K$ have full rational
$p$-torsion. If its reduction is additive, it is potentially good.
Thus the additive existence question in B.3 is the potentially good
one, although its unrestricted good-forcing question also includes
multiplicative Tate curves.
\end{lemma}
\begin{proof}
The Weil pairing gives $\mu_p\subset K$. Suppose instead that
$v_K(j(E))<0$. The integral power series for $j(q)^{-1}$ starts with
$q$, so its compositional inverse gives $q\in K$, $v_K(q)>0$, with
$j(E_q)=j(E)$, where $E_q$ is the Tate curve
\cite[Theorem~1 and the $j$-expansion, pp.4--6]{Tate1995Uniformization}.
This value of $j$ is neither zero nor $1728$. Elliptic curves with
that same $j$ are quadratic twists, so $E=E_q\otimes\epsilon$ for
a quadratic character $\epsilon$. Tate uniformization gives
\[
       E[p]^{\mathrm{ss}}
        \simeq\overline\epsilon\,\overline\chi_p
                      \oplus\overline\epsilon.
\]
The left side is trivial and $\overline\chi_p=1$, so
$\overline\epsilon=1$. As $p$ is odd, reduction modulo $p$ is
injective on $\{1,-1\}$, hence $\epsilon=1$. The curve is therefore
split multiplicative, contrary to additive reduction. Thus $j(E)$
is integral, which is the elliptic criterion for potential good
reduction.
\end{proof}

\begin{lemma}[The maximum formal valuation is divisible by three]
\label{lem:B-RD20-wild}
\status{Proved; independent proof and source audits by the parent and obstructions/crystals agents passed.}
Suppose $p=3$, $E[3^n]\subset E(K)$, and $E$ has additive reduction
with wild semistability defect $m\in\{3,6,12\}$.
After extending the residue field algebraically, let $L/K$ be the
minimal good-reduction extension, with its normalized valuation
$v_L$. The maximum valuation $A$ of a primitive $3^n$-torsion
parameter in the good formal group is divisible by three.
\end{lemma}
\begin{proof}
First the discriminant of any Weierstrass equation for $E$ is a
cube in $K^\times$. The curve $E_0:y^2+y=x^3$ has discriminant
$-27=(-3)^3$ and full rational $3$-torsion over $K\supset\mu_3$,
as checked in Theorem~\ref{thm:B-RD2}. Choose a symplectic
isomorphism $E_0[3]\simeq E[3]$. The universal property of the
Hesse pencil makes $E$ $K$-isomorphic to a member of the Hesse
pencil of $E_0$ \cite[Theorem~1.1, p.2; proof, pp.7--8]
{AnemaTopTuijp2018Hesse}. Completing the square gives the short
equation $y^2=x^3+1/4$ with the same discriminant. Proposition~5.1
of that source, pp.4--5, computes the discriminant of every finite
smooth member as $\Delta(E_0)(\det A_t)^3$; the member at infinity
is $E_0$ itself. Thus every such discriminant is a cube.
A change of Weierstrass coordinates changes the discriminant by a
twelfth power. Hence $\delta=v_K(\Delta)$ is divisible by three.

Write the relation between a $K$-rational invariant differential and
an integral generator on the good model over $L$ as
$\omega_L=u\omega_K$. Since the good discriminant is a unit,
\[
                  v_L(u)=m\delta/12.
\]
This is an integer divisible by three, for each
$m=3,6,12$. Let $\sigma$ generate the wild subgroup $C_3$, and set
$F=L^{C_3}$. Choose $b_0\in F^\times$ with
$v_L(b_0)=v_L(u)$. Replacing the integral differential by
$(b_0/u)\omega_L$ makes it $C_3$-invariant without changing its
integrality, since $b_0/u$ is a unit. Choose a formal parameter with
this linear differential and odd under $[-1]$, which is possible
because two is a unit. Thus the wild descent series has linear
coefficient exactly one and only odd powers.

The special fibre is supersingular by Proposition~\ref{prop:B-RD3}.
On $y^2=x^3-x$ its order-three automorphisms are translations
$x\mapsto x+r$, $r\in\mathbf F_3^\times$. In the odd parameter
$t=-x/y$, their first nonidentity term has degree three and nonzero
coefficient. Therefore the wild descent series, or its inverse
according to the coordinate convention, has the form
\[
                   t\longmapsto t+c t^3+O(t^5),
                   \qquad c\in\mathcal O_L^\times.
\]
For a $K$-rational nonzero torsion point with parameter $x$ this gives
\[
                     v_L(\sigma x-x)=3v_L(x).       \tag{B-RD20.1}
\]

Let $b=v_L(\sigma\pi-\pi)-1$ be the positive lower ramification
break, for a uniformizer $\pi$ of $L$. The standard uniformizer
expansion, with coefficients in the fixed Cohen ring $W(k)$, gives
\[
 v_L(\sigma x-x)\geq v_L(x)+b,
 \quad\text{with equality if }3\nmid v_L(x).       \tag{B-RD20.2}
\]
Indeed the term of lowest valuation in $x$ gives equality when its
exponent is prime to three, and every higher term has strictly
larger valuation in its difference. If $B$ is the minimum primitive
valuation, (B-RD20.1)--(B-RD20.2) imply $b\leq2B$. If $A>B$ and
$3\nmid A$, the same equations give $b=2A>2B$, a contradiction.

It remains to exclude $A=B$ with $3\nmid A$. In that case $b=2A$.
For every primitive torsion parameter $x$, write
$x=a_x\pi^A+O(\pi^{A+1})$, with $a_x\in k^\times$. The first
coefficient of $\sigma x-x$, divided by $\pi^{3A}$, is a fixed
nonzero scalar times $a_x$; the first coefficient of the right
side of the descent equation is $\bar c a_x^3$. Thus all these
coefficients solve one equation
\[
                           \bar c X^3=dX,
                           \qquad d\ne0.
\]
Here $d=A\overline{(\sigma\pi-\pi)/\pi^{b+1}}\ne0$, since
$3\nmid A$. Its roots have dimension at most one over
$\mathbf F_3$. Choose a
primitive basis $P,Q$ of $E[3^n]$. Sending $(a,b)\in\mathbf F_3^2$
to the leading coefficient of $[\widetilde a]P+[\widetilde b]Q$
is well defined and linear: changing an integer lift by a multiple
of three does not affect the leading coefficient, because formal
addition is addition to first order and the residue characteristic
is three. This map is injective, since a nonzero combination is
primitive and cancellation would give valuation greater than the
maximum $A$. This requires dimension two, a contradiction. Hence
$3\mid A$ also in this case.
\end{proof}

\begin{proposition}[B-RD20: a stricter wild bound and an exact minimum]
\label{prop:B-RD20}
\status{Proved; independent proof and source audits by the parent and obstructions/crystals agents passed.}
For $p=3$ write $e=2u$ and $s=v_3(u)$. If full rational $3^n$-torsion
occurs with wild defect $m\in\{3,6,12\}$, necessarily
\[
 s\geq n-1,\qquad
 e\geq\frac8m\,3^{\max(2n-2,\,4n-4-s)}.        \tag{B-RD20.3}
\]
Allowing every perfect residue field and every defect, the least
absolute ramification index admitting an additive elliptic curve
with full rational $3^n$-torsion is exactly
\[
                            e_{\min}(n)=2\cdot3^{2n-2}.
                                                        \tag{B-RD20.4}
\]
In particular $e_{\min}(2)=18$: absolute indices six and twelve,
although compatible with the cyclotomic necessary condition,
do not admit an additive curve with full rational $9$-torsion.
\end{proposition}
\begin{proof}
The cyclotomic necessary condition gives $s\geq n-1$. The
valuation identity in the proof of Theorem~\ref{thm:B-RD12} uses
only the good height-two formal group, not the tame descent
congruences. In a valuation region
$3^{2r}B\leq A\leq3^{2r+2}B$, $0\leq r\leq n-1$, it therefore
still gives
\[
                  me=2\cdot3^{2n-r-2}(A+3^{2r+1}B).
                                                        \tag{B-RD20.5}
\]
The upper bound for $A$ is omitted when $r=n-1$.
By Lemma~\ref{lem:B-RD20-wild}, the parenthesis is divisible by
three. Since $v_3(m)=1$, this forces
$s\geq2n-r-2$, or $r\geq\max(0,2n-2-s)$.
The lower bound on $A$ and $B\geq1$ give
$me\geq8\cdot3^{2n+r-2}$. Substitute the least possible $r$ to
obtain (B-RD20.3).

To prove (B-RD20.4), tame defects satisfy the lower bound
$e\geq2\cdot3^{2n-2}$ by Theorem~\ref{thm:B-RD12} with $m\leq4$.
For a wild defect use $m\leq12$ in (B-RD20.3). If
$s\geq2n-2$, divisibility of $e$ by $2\cdot3^s$ already gives
the claimed bound. Otherwise $s\leq2n-3$, and the second exponent
in (B-RD20.3) is at least $2n-1$. Consequently
\[
 e\geq\frac23\,3^{2n-1}=2\cdot3^{2n-2}.
\]
Finally the tame construction with $m=4$ and
$e=2\cdot3^{2n-2}$ meets the criterion of
Theorem~\ref{thm:B-RD12}, so it attains this bound for every $n$.
Lemma~\ref{lem:B-RD19} covers every additive curve in the question.
\end{proof}

\begin{corollary}[B-RD21: the remaining wild question is one level]
\label{cor:B-RD21}
\status{Proved; independently audited by the parent agent; arithmetic certificate supplied.}
Fix $p=3$ and a positive even absolute ramification index $e$.
Let $N(e)$ be the largest $n\geq1$ for which an additive elliptic
curve with full rational $3^n$-torsion exists over some such field
with perfect residue field. Put
\[
 s=v_3(e/2),\qquad L=\lfloor\log_3(e/2)\rfloor,
\]
and define
\begin{align*}
 N_{\mathrm t}
   &=1+\min\left(s,\left\lfloor\frac L2\right\rfloor,
                         \left\lfloor\frac{L+s}{4}\right\rfloor\right),\\
 N_{\mathrm w}
   &=1+\min\left(s,\left\lfloor\frac{L+1}{2}\right\rfloor,
                         \left\lfloor\frac{L+s+1}{4}\right\rfloor\right).
\end{align*}
Then
\[
             N_{\mathrm t}\leq N(e)\leq N_{\mathrm w}
                             \leq N_{\mathrm t}+1.       \tag{B-RD21.1}
\]
In particular the full answer is already determined whenever
$N_{\mathrm t}=N_{\mathrm w}$. A gap occurs precisely when there is
an integer $k\geq1$ such that
\[
             k\leq s\leq2k-1,\qquad L=4k-s-1.          \tag{B-RD21.2}
\]
Then $N_{\mathrm t}=k$ and $N_{\mathrm w}=k+1$, so only existence
at level $n=k+1$ remains undecided. For full $9$-torsion, the only
indices not decided by these results are
\[
                              e=24,30,42,48.            \tag{B-RD21.3}
\]
At $e=24,30$ an example, if one exists, must have defect twelve;
at $e=42,48$ its possible wild defects are six and twelve.
\end{corollary}
\begin{proof}
The sharp tame criterion with defect four gives the lower bound
$N_{\mathrm t}$. The bound in Proposition~\ref{prop:B-RD20} is
weakest for defect twelve; its logarithmic inequalities are
$n-1\leq s$, $2n-2\leq L+1$, and $4n-4-s\leq L+1$.
They give the upper bound $N_{\mathrm w}$. Increasing $L$ by one
increases each displayed minimum by at most one.

Since $s\leq L$, a level $k+1$ admitted by the wild inequalities
but excluded by the tame ones must satisfy
\[
 k\leq s,\qquad \max(2k,4k-s)=L+1.
\]
The maximum cannot be $2k$ with $s\geq2k$, because that would
give $L=2k-1<s$. Thus $s\leq2k-1$ and $L=4k-s-1$.
Conversely these conditions make the two minima respectively
$k-1$ and $k$, which proves (B-RD21.2).
For $k=1$ we have $s=1,L=2$, so $e/2$ is a multiple of three
but not nine in the interval $[9,27)$; the four possibilities
are $12,15,21,24$. Substitute them into (B-RD20.3) for the
defect assertions. The finite integer certificate
\path{code/verify_p3_higher_wild_bounds.py} checks these formulae,
their extrema, and their compatibility with the unsimplified
valuation identity in the stated finite range.
\end{proof}


\subsection{An unramified height-two obstruction beyond the first level}

\begin{proposition}[B-RD22: four-torsion does not suffice at $p=2,e=1$]
\label{prop:B-RD22}
\status{Proved; exact finite certificate and independent parent and obstructions/crystals audits passed.}
For $p=2$, $e=1$, and every height $h\geq2$, the least index in
B.4 is at least two. In the manuscript's indexing, even
$M(1)=M[4]$ does not determine the character on inertia elements of
odd order. This already occurs over the unramified quadratic field
$K=\mathbf Q_2(\zeta_3)$.
\end{proposition}
\begin{proof}
Let
\[
 E/K:\quad y^2=x(x-1)(x+3),\qquad
 M=E[2^\infty],\qquad
 N=\mu_{2^\infty}\oplus\mathbf Q_2/\mathbf Z_2.
\]
The group $N$ extends to a split Barsotti--Tate group over
$\mathcal O_K$. Put $a=1+2\zeta_3$, so that $a^2=-3$, and choose
$i^2=-1$. The points
\[
       P=(a,ia(a-1)),\qquad Q=(3,6)
\]
satisfy $2P=(0,0)$ and $2Q=(1,0)$, and hence form a basis of
$E[4]$. Every element of $G_K$ fixes $a$ and sends $i$ to either
$i$ or $-i$. Consequently it sends $P$ to $\chi_2 P$ modulo four
and fixes $Q$. The equivalence between finite etale group schemes
over the characteristic-zero field $K$ and finite $G_K$-sets gives
an isomorphism of group schemes
\[
                     M[4]\simeq\mu_4\oplus\mathbf Z/4
                              =N[4].                 \tag{B-RD22.1}
\]

For the displayed integral Weierstrass equation the exact invariants
are
\[
 \Delta=2^8\,3^2,\qquad c_4=2^4\,13,\qquad
 j=\frac{2^4\,13^3}{9}.
\]
The discriminant valuation is eight, which is minimal since every
integral change to another Weierstrass equation changes it by a
multiple of twelve and a discriminant of an integral equation has
nonnegative valuation. In particular the curve is bad. Its integral
$j$ gives potential good reduction
\cite[Section~2, p.496]{SerreTate1968Good}. Since $v_K(j)=4>0$,
the good special fibre has $j=0$. In characteristic two its equation
has the form $y^2+a_3y=x^3+a_4x+a_6$, with $a_3\ne0$
\cite[Section~9.1, p.8]{AnemaTopTuijp2018Hesse}. Negation sends
$(x,y)$ to $(x,y+a_3)$, so there is no nonzero geometric
two-torsion point. Thus the good fibre is supersingular.

Pass to the completed maximal unramified extension and let $L/K$
be a minimal good-reduction extension, with inertia image of order
$m$. A change to a unit-discriminant equation over $L$ shows that
\[
                            12\mid8m,
\]
so $3\mid m$. The inertia image therefore contains an automorphism
of order three. On the supersingular special elliptic curve its
action on the rational Dieudonne module has eigenvalues $\zeta_3$
and $\zeta_3^{-1}$, hence trace $-1$: equivalently, use the standard
supersingular equation $y^2+y=x^3$ and its order-three automorphism
$(x,y)\mapsto(\zeta_3x,y)$; the two eigenvalues are interchanged by
Frobenius. The order-three trace of $N$, which has good reduction,
is $2$. Thus (B-RD22.1) does not determine the requested character.
Both groups satisfy the manuscript's monodromy condition.

For height $h>2$, add $h-2$ copies of the constant etale group
$\mathbf Q_2/\mathbf Z_2$ to both sides. Their equal traces add the
same constant, so the discrepancy remains. The entire finite
four-torsion calculation is checked with exact rational arithmetic
in $\mathbf Q(\sqrt{-3},i)$ by
\path{code/verify_p2_character_level2.py}, including all sixteen
points, all $256$ additions, and the Galois action.
\end{proof}


\subsection{The complete full-nine-torsion range}

\begin{theorem}[B-RD23: the four remaining full-nine cases exist]
\label{thm:B-RD23}
\status{Proved; independent parent and obstructions/crystals proof/source audits passed.}
For $p=3$, an additive elliptic curve with full rational $9$-torsion
exists over a complete discretely valued field of absolute
ramification index $e$ and perfect residue field if and only if
\[
                              6\mid e,\qquad e\geq18.
                                                        \tag{B-RD23.1}
\]
The examples can all be taken over finite extensions of $\mathbf Q_3$.
In particular all four indices $24,30,42,48$ left by
Corollary~\ref{cor:B-RD21} occur, with semistability defect twelve.
\end{theorem}
\begin{proof}
Necessity follows from the cyclotomic condition and
Proposition~\ref{prop:B-RD20}. We first make the example at $e=18$
explicit over a finite local field. Let $K_0=\mathbf Q_3(i)$,
$i^2=-1$, and $E_0:y^2=x^3-x$. This has good supersingular
reduction and the endomorphism
$[i]:(x,y)\mapsto(-x,iy)$, defined over $K_0$. Put
$L=K_0(E_0[9])$, the finite Galois division field.

Every primitive nine-torsion parameter $t_0=-x/y$ has
$v_3(t_0)=1/72$. Indeed, the unramified height-two formal group's
$[3]$-series has linear coefficient three, its first unit
coefficient is in degree nine, and every intervening coefficient
is divisible by three. A nonzero three-torsion parameter therefore
has valuation $1/8$. A preimage under $[3]$ must have smaller
valuation, on which the degree-nine term has strictly least
valuation; hence its valuation is $1/72$.

On the other hand, $E_0[9]$ is free of rank one over
$\mathbf Z_3[i]/9$. To see this directly, reduction modulo three
has dimension two over $\mathbf F_3$ and the operator $i$ has
irreducible polynomial $X^2+1$; it is therefore one-dimensional
over $\mathbf F_9$. Nakayama's lemma and the equal cardinalities
then give the rank-one assertion modulo nine. Galois commutes
with $i$, so its faithful image has order at most
$|(\mathbf Z_3[i]/9)^\times|=72$. The preceding valuation
denominator forces ramification degree at least $72$.
Consequently $L/K_0$ is totally ramified of degree $72$ and this
Galois image is the full unit group.

Let $\sigma\in\operatorname{Gal}(L/K_0)$ act on $E_0[9]$ by $[i]$,
and let $K=L^{\langle\sigma\rangle}$. Choose a primitive point
$Q\in E_0[9]$ and put $u=-x(Q)/y(Q)$. It is a uniformizer of
$L$, and the exact coordinate formula for $[i]$ gives
$\sigma u=iu$. Thus $a=u^4\in K$ has $v_K(a)=1$. The curve
\[
                         E:y^2=x^3-ax
\]
has full rational nine-torsion: its points obtained from
$P\in E_0[9]$ are $(u^2x(P),u^3y(P))$, and both coordinates are
fixed by $\sigma$. The absolute ramification index of $K$ is
$72/4=18$. The curve becomes $E_0$ over $L$, while
$v_K(\Delta(E))=v_K(64a^3)=3$ forces exact defect four. This gives
the promised finite local-field example at the endpoint.

We next give a uniform explicit construction for every multiple
$e=6d>18$.

Choose $i^2=-1$, a primitive ninth root of unity $\eta$, and put
$\zeta=\eta^3$. Let
\[
 K=\mathbf Q_3(i,\eta,\pi),\qquad \pi^d=\eta-1.
                                                        \tag{B-RD23.2}
\]
The field $\mathbf Q_3(i)$ is the unramified quadratic extension,
and $\eta-1$ is a uniformizer of its degree-six cyclotomic
extension. The final polynomial in (B-RD23.2) is Eisenstein.
Thus $K$ has absolute ramification $e$, residue field $\mathbf F_9$,
and uniformizer $\pi$. Write $v(\pi)=1$.

Put
\[
 h=1+\pi^2,\quad D=3(4h^3-1),\quad
 y=\frac{-3+\sqrt D}{2(1-h^3)}.                         \tag{B-RD23.3}
\]
Here $v(h^3-1)=6$, $v(4h^3-1)=6$, and $v(D)=e+6$.
Moreover $3/\pi^e\equiv-1\pmod\pi$, by the Eisenstein
cyclotomic equation for $\eta-1$. Consequently
$D/\pi^{e+6}\equiv-1\pmod\pi$. This unit has a square root in
$K$ lifting either $i$ or $-i$, by ordinary Hensel lifting, since
two is a unit. This specifies the square root in (B-RD23.3).
The quadratic formula gives
\[
       h^3=1+\frac3y+\frac3{y^2},\qquad
       r:=v(y)=\frac e2-3.                             \tag{B-RD23.4}
\]
In particular $r>0$. Define
\[
 \lambda=-\eta^2,\qquad c=\frac{\lambda}{\zeta-1},\qquad
 d_0=(1+y)^3-1=y^3h^3,\qquad
 t=1+\frac{1-\zeta}{d_0}.                             \tag{B-RD23.5}
\]
Since $(\zeta-1)^2=-3\zeta$, one has
$c^3=(1-\zeta)/9$. Therefore
\[
 \frac{t-1}{9}=\left(\frac{c}{yh}\right)^3,
 \qquad
 \frac{t-\zeta}{9}
   =\left(\frac{c(1+y)}{yh}\right)^3.                 \tag{B-RD23.6}
\]

We show that $(t-\zeta^2)/9$ is also a cube. Its ratio with the
first expression in (B-RD23.6) is
\[
              a=1-\zeta^2d_0=1+\lambda^3d_0.
\]
Use $\beta_0=1+\lambda y$ as a first approximation to a cube root.
Then
\[
 a-\beta_0^3
      =3(\lambda^3-\lambda)y
                   +3(\lambda^3-\lambda^2)y^2.
                                                        \tag{B-RD23.7}
\]
The first coefficient in parentheses has valuation $e/6$ and the
second is a unit. Since $r>e/6$, the valuation in (B-RD23.7) is
\[
             e+q,\qquad q=e/6+r=2e/3-3>e/2.
\]
The polynomial
\[
       \frac{(\beta_0+\pi^q Z)^3-a}{3\pi^q}
\]
has integral coefficients. Its linear coefficient is the unit
$\beta_0^2$, its quadratic coefficient is $\beta_0\pi^q$, and
its cubic coefficient has valuation $2q-e>0$. Its reduction is
therefore a degree-one polynomial, which has a simple root over
$\mathbf F_9$. Hensel lifting gives $\beta\in K$ with
$\beta^3=a$. This argument uses the displayed rescaling; it does
not invoke the generally inapplicable stronger inequality
$v(a-\beta_0^3)>2v(3)$.

Now use the original full-level modular-curve equations of
Fisher \cite[Lemma~4.2(i), p.8; Theorem~4.3, equation~(9), p.9]
{Fisher2015Nine}. They state that the family over $Y(3)$ is
\[
             E_t:\quad y_0^2+3tx_0y_0+(t^3-1)y_0=x_0^3,
                                                        \tag{B-RD23.8}
\]
and the cover $X(9)\to X(3)$ is given by the three equations
\[
     t-\theta=9(u+v\theta+w\theta^2)^3,
                     \qquad \theta=1,\zeta,\zeta^2.    \tag{B-RD23.9}
\]
The three cube roots supplied above are respectively
$c/(yh)$, $c(1+y)/(yh)$, and $c\beta/(yh)$. An invertible Fourier
matrix over $K$ gives $u,v,w$ with precisely these three values.
Thus (B-RD23.9) gives a $K$-point of $Y(9)$. By the moduli
interpretation in Section~2, p.4 of that source, it gives
$E_t[9]\simeq\mu_9\oplus\mathbf Z/9$. Since $\mu_9\subset K$,
the entire $9$-torsion is rational. The parameter is not a cusp,
because the valuation calculation below gives $v(t)<0$.

Finally, (B-RD23.4)--(B-RD23.5) give $v(t)=9-e<0$. The exact
invariants of (B-RD23.8) are
\[
 c_4=9t(t^3+8),\qquad \Delta=27(t^3-1)^3,\qquad
 j=\frac{27t^3(t^3+8)^3}{(t^3-1)^3}.
\]
Hence
\[
                   v(j)=27,\qquad v(\Delta)=81-6e
                                             \equiv9\pmod {12}.
\]
The integral $j$ gives potential good reduction
\cite[Section~2, p.496]{SerreTate1968Good}, and the discriminant
congruence excludes good reduction. Thus the reduction is additive.
An explicit integral equation is obtained by the substitution
$x_0=\pi^{2(9-e)}X$, $y_0=\pi^{3(9-e)}Y$: its two nonzero
coefficients have valuations $e$ and zero, and its discriminant
has valuation $6e-27$.

At any of $e=24,30,42,48$, the discriminant congruence forces the
minimal good inertia order $m$ to be divisible by four. Its possible
orders in characteristic three are therefore four and twelve.
The sharp tame full-nine criterion excludes four at these indices,
so $m=12$. All displayed finite polynomial identities, the Fourier
inversion, and the valuation inequalities at the four previously
undecided indices are checked in
\path{code/verify_full9_construction.py}. The square-root and cubic
Hensel steps are justified above, rather than treated as numerical
approximations.
\end{proof}

\begin{remark}[Passage to all higher levels]\label{auto:B:b02333f8ed15}
\status{Proved under stated hypotheses}

This first settles level $n=2$ of the wild higher-torsion question,
including every exceptional band with $k=1$ in Corollary~\ref{cor:B-RD21}.
The first higher case beyond this construction is full $27$-torsion
at $e=216$. It is subsequently realized by the wild trace-invariant
construction of Theorem~\ref{thm:B-RD26-wild}, which settles every
higher wild level and yields Corollary~\ref{cor:B-RD27-exact-three}.
\end{remark}


\subsection{An explicit finite algebraic test for elliptic torsion existence}

\begin{proposition}[A bounded test; the output at the first wild gap is not yet evaluated]
\label{prop:B-RD24-finite-test}
\status{Proved; complete proof and original-source audit independently passed by the parent, with an independent rerun of the exact polynomial certificate.  No numerical output for $(p,e,n)=(3,216,3)$ is asserted.}
For every odd prime $p$ and positive integers $e,n$, the existence of a complete discretely valued field of absolute ramification $e$, perfect residue characteristic $p$, and an additive elliptic curve with full rational $p^n$-torsion is equivalent to the finite polynomial test specified below.  The test uses only arithmetic in finite-length Witt vectors, polynomial division and discriminants over $\mathbf Q$, and effective elimination over $\overline{\mathbf F}_p$.  If the test is affirmative, an example exists over a finite extension of $\mathbf Q_p$.
\end{proposition}

\par\smallskip\noindent\textbf{Constants and the division polynomial.}
Put $m=p^n$, $N=(m^2-1)/2$, and $a=ne$.  Write
\[
 \psi_m(X)=mX^N+\sum_{i=0}^{N-1}d_iX^i
\]
for the usual odd division polynomial (whose recurrences and torsion-root interpretation are recorded in \cite[section~1.1.3, pp.~8--9]{Burhanuddin2007}) of
\[
 y^2+a_1xy+a_3y=x^3+a_2x^2+a_4x+a_6.
\]
The long-equation version used here can be specified without any division in its coefficient ring.  Put
\[
\begin{gathered}
 b_2=a_1^2+4a_2,\quad b_4=a_1a_3+2a_4,\quad
 b_6=a_3^2+4a_6,\\
 b_8=a_1^2a_6+4a_2a_6-a_1a_3a_4+a_2a_3^2-a_4^2,\qquad
 B(X)=4X^3+b_2X^2+2b_4X+b_6.
\end{gathered}
\]
For odd $j$ let $P_j=\psi_j$, and for even $j$ let
$P_j=\psi_j/\psi_2$.  Their initial values are
\[
\begin{aligned}
 P_0&=0,\quad P_1=P_2=1,\\
 P_3&=3X^4+b_2X^3+3b_4X^2+3b_6X+b_8,\\
 P_4&=2X^6+b_2X^5+5b_4X^4+10b_6X^3+10b_8X^2\\
    &\hspace{7mm}+(b_2b_8-b_4b_6)X+b_4b_8-b_6^2.
\end{aligned}
\]
For subsequent indices the recurrences are
\[
 P_{2r}=P_r(P_{r+2}P_{r-1}^2-P_{r-2}P_{r+1}^2)
\]
and
\[
 P_{2r+1}=
 \begin{cases}
 P_{r+2}P_r^3-B^2P_{r-1}P_{r+1}^3,&r\text{ odd},\\
 B^2P_{r+2}P_r^3-P_{r-1}P_{r+1}^3,&r\text{ even}.
 \end{cases}
\]
They follow from the usual recurrences and the identity $\psi_2^2=B(X)$.
Completing the square and then translating $X$ to remove its quadratic coefficient over $\mathbf Q$ identifies these initial values and recurrences with the short-equation division polynomials in the cited source; thus it also gives their stated torsion-root interpretation.  All the displayed operations are integral polynomial operations.  In particular the $d_i$ can be calculated by this finite recursion.  Define
\[
 c_m=\frac{\operatorname{disc}\bigl(\psi_m(X;0,0,0,-1,0)\bigr)}
                  {64^{\,N(N-1)/6}}\ \in\mathbf Q^\times,
 \qquad t_m=v_p(c_m).
\]
This definition is an exact rational computation, not a conjectured closed formula for the discriminant constant.
Let
\[
 D_0=\begin{cases}6e+11,&p=3,\\11,&p\geq5,\end{cases}
 \qquad
 D_F=a(N-1)(N-2)+e\,t_m+\frac{N(N-1)}6D_0,
\]
and take the deliberately generous positive integer
\[
 D=2|D_F|+a(4N+3)+ND_0+1,\qquad
 H=\max\{2D+1,\,3e+7,\,D_0+1\},\qquad q=\lceil H/e\rceil.
\]
The proof below shows that $D_F$ is in fact nonnegative and bounds the valuation of the monic polynomial discriminant used in the test.  Absolute values in the definition of $D$ merely make positivity manifest before this verification.

\par\smallskip\noindent\textbf{The finite coefficient ring.}
All quantified residue-field variables range over $k=\overline{\mathbf F}_p$.
Choose $e$ Witt vectors $c_0,\ldots,c_{e-1}\in W_q(k)$ with the zeroth Witt coordinate of $c_0$ nonzero, and form the finite-length ring
\[
 R=W_q(k)[\pi]\big/\left(\pi^e+pC(\pi)\right),
 \qquad C(\pi)=\sum_{i=0}^{e-1}c_i\pi^i.
\]
This ring is free of rank $e$ over $W_q(k)$, so all its operations are explicit polynomial operations in $eq$ residue-field coordinates.  Its maximal ideal is $(\pi)$ and its length is $eq\geq H$.  Membership in $\pi^rR$ is the finite polynomial condition
\[
                  z\in\pi^rR\quad\Longleftrightarrow\quad
                  (\exists w\in R)\ z=\pi^rw.
\]
Thus every condition below is a finite first-order polynomial formula over $k$.

\par\smallskip\noindent\textbf{The test.}
Ask whether there exist such $c_i$, coefficients $a_1,a_2,a_3,a_4,a_6\in R$, an integer $1\leq\delta\leq D_0$, and $N$ pairs $(X_i,Y_i)\in R^2$ satisfying the following conditions.

\begin{enumerate}
\item The universal invariants satisfy
\[
 \Delta\in\pi^\delta R\setminus\pi^{\delta+1}R,
 \qquad c_4^3\in\pi^\delta R.
\]
\item There are no $S,R_1,T_1\in R$ for which all five expressions
\[
\begin{aligned}
 U_1&=a_1+2S,\\
 U_2&=a_2-Sa_1+3r-S^2,\\
 U_3&=a_3+ra_1+2t,\\
 U_4&=a_4-Sa_3+2ra_2-(t+rS)a_1+3r^2-2St,\\
 U_6&=a_6+ra_4+r^2a_2+r^3-ta_3-rta_1-t^2,
\end{aligned}
\qquad r=\pi^{-e}R_1,\quad t=\pi^{-e}T_1
\]
satisfy $\pi^{3e}U_j\in\pi^{3e+j}R$, for $j=1,2,3,4,6$.
This notation contains no inversion of $\pi$ in $R$: first expand each $\pi^{3e}U_j$ and substitute $r=\pi^{-e}R_1,t=\pi^{-e}T_1$; every resulting exponent of $\pi$ is nonnegative.
\item Form the two integral polynomials
\[
\begin{aligned}
 F(X)&=X^N+\sum_{i=0}^{N-1}
       (-C(\pi))^n\pi^{a(N-i-1)}d_iX^i,\\
 G(X,Y)&=Y^2+(a_1\pi^aX+a_3\pi^{2a})Y\\
 &\hspace{9mm}-\bigl(\pi^aX^3+a_2\pi^{2a}X^2
                   +a_4\pi^{3a}X+a_6\pi^{4a}\bigr).
\end{aligned}
\]
For every $i$ require
\[
 F(X_i),G(X_i,Y_i)\in\pi^HR,\qquad
 F'(X_i)G_Y(X_i,Y_i)\notin\pi^{D+1}R.
\]
For $i\ne j$ require $X_i-X_j\notin\pi^{D+1}R$.
\end{enumerate}
Here $R_1$ in condition~2 is a variable, not the coefficient ring $R$.
The test returns \emph{yes} exactly when this finite formula is true in an algebraically closed field of characteristic $p$.

\begin{proof}
We give the bounds and the exact approximation argument, since these distinguish the test from an enumeration of unknown local fields or unknown infinite lifts.
Lemma~\ref{lem:B-RD19} proves that full rational odd-prime torsion excludes additive potentially multiplicative reduction.  Consequently every additive curve in the statement is potentially good, so the following integral-$j$ bounds cover the whole stated problem.

\emph{Minimal discriminants.}
Let an integral minimal equation have potentially good reduction, and put $\delta=v(\Delta)$.  Since $j=c_4^3/\Delta$ is integral, $v(c_4)\geq\delta/3$.  The identity $c_4^3-c_6^2=1728\Delta$ then gives $v(c_6)\geq\delta/2$.  Its short equation has coefficients $-c_4/48$ and $-c_6/864$.  If $p=3$ and $\delta\geq6e+12$, their valuations are at least $4$ and $6$; scaling $x=\pi^2x',y=\pi^3y'$ gives an integral equation with discriminant valuation $\delta-12$, a contradiction.  For $p\geq5$ the same argument applies to $\delta\geq12$.  Hence $\delta\leq D_0$.

An integral equation is nonminimal exactly when a change of variables with scaling factor $\pi$ produces another integral equation.  Indeed, a change with factor $\pi^b$, $b\geq1$, also gives one with factor $\pi$ using the same translations; unit scaling does not affect integrality.  The standard change-of-variables formulas are precisely the $U_j$ above.  The first condition gives $s\in\OO$ because $2$ is a unit.  The second gives $v(r)\geq-v(3)\geq-e$; the third then gives $v(t)\geq-e$.  Conversely these are all the possible translations.  Clearing denominators by $\pi^{3e}$ proves that condition~2 is exactly minimality and depends only on the displayed finite precision.

\emph{The two discriminant bounds.}
Over $\mathbf Q[a_1,a_2,a_3,a_4,a_6,\Delta^{-1}]$, the nonzero $m$-torsion modulo the free involution $P\mapsto-P$ is finite \'etale, with coordinate algebra defined by $\psi_m/m$.  Consequently $\operatorname{disc}(\psi_m)$ is a unit after inverting $\Delta$.  The universal Weierstrass discriminant is irreducible over $\mathbf Q$: after the invertible change to a short equation this is the irreducibility of $4A^3+27B^2$.  Its weighted degree is $12$, while the weighted degree of $\operatorname{disc}(\psi_m)$ is $2N(N-1)$.  It follows that
\[
           \operatorname{disc}(\psi_m)=c_m\Delta^{N(N-1)/6};
\]
specializing to $y^2=x^3-x$ proves exactly the formula defining $c_m$.

Over a lift to a complete DVR, $p/\pi^e=-C(\pi)^{-1}$, so $F(X)$ is exactly
\[
                 \frac{\pi^{aN}}{m}\psi_m(\pi^{-a}X).
\]
It is monic integral, and therefore
\[
 v(\operatorname{disc}F)
   =a(N-1)(N-2)+e\,t_m+\frac{N(N-1)}6\delta
   \leq D_F.
\]
The algebra
\[
                 \mathcal A=\OO[X,Y]/(F,G)
\]
is finite free of rank $2N=m^2-1$, with \'etale generic fibre equal to the coordinate algebra of $E[m]\setminus\{0\}$.
For completeness the needed second resultant has an elementary derivation.  Write
\[
 B(x)=4x^3+b_2x^2+2b_4x+b_6.
\]
After completing the square, at a two-torsion abscissa $\alpha$ the division-polynomial recurrences give
\[
 \psi_m(\alpha)=(-1)^{(m-1)/2}
         f'(\alpha)^{(m^2-1)/4},\qquad
 f(x)=x^3+(b_2/4)x^2+(b_4/2)x+b_6/4.
\]
The base cases are $\psi_1=1$ and $\psi_3(\alpha)=-f'(\alpha)^2$.
For $m=2r+1$ one of the two terms in the odd recurrence vanishes at $\alpha$, and the other proves the formula by induction.  Multiplying over the three roots of $f$ gives
\[
                 \operatorname{Res}(\psi_m,B)
                     =(-1)^{(m-1)/2}\Delta^{N/2}.
\]
The discriminant of the quadratic $G$ is $\pi^{4a}B(\pi^{-a}X)$.
The discriminant formula for a tower of finite free algebras, or its direct determinant proof after splitting $F$, now gives
\[
 v(\operatorname{disc}\mathcal A)
 =2v(\operatorname{disc}F)+a(4N-3)+(N/2)\delta<D.
\]
This proves the required uniform bound without using a formula for $c_m$ that has not been independently established.

\emph{Why approximate roots are exact enough.}
If $E[m]$ is rational, all $N$ roots of the monic $F$ and both roots of the corresponding monic quadratics $G$ lie in $\OO$.  At each point the Jacobian determinant of $(F,G)$ is $F'G_Y$.  The product of these determinants over all $2N$ points is, up to sign, the discriminant of $\mathcal A$.  All the factors are integral, so each has valuation less than $D$.  Distinct roots $X_i,X_j$ have $v(X_i-X_j)\leq v(\operatorname{disc}F)/2<D$, by the Vandermonde product.  Thus reduction modulo the stated precision supplies all the conditions.

Conversely, lift approximate pairs to $\OO$.  Their equation errors have valuation at least $H\geq2D+1$, whereas their Jacobian determinants have valuation at most $D$.  The following elementary multivariate Hensel argument applies to each pair.  If a polynomial map $f$ has integral Jacobian $J$ at $z$, determinant valuation $b\leq D$, and error valuation at least $H$, substitute $z'=z+\pi^{H-b}w$ and consider
\[
                \pi^{-H}\operatorname{adj}(J)
                              f(z+\pi^{H-b}w).
\]
This is an integral polynomial system.  Its linear Jacobian modulo $\pi$ is $(\det J/\pi^b)I$, which is invertible; all nonlinear coefficients are divisible by $\pi$, since $H-2b\geq1$.  The reduction is therefore a linear system with a unique solution, and successive coefficient lifting gives an exact zero.  Since $J$ is invertible over the fraction field, this is also a zero of the original system.  Its difference from $z$ has valuation at least $H-b\geq D+1$.  The distinctness conditions consequently retain $N$ distinct exact abscissas.  Each exact pair is a nonzero $m$-torsion point; its inverse is also rational.  These $2N$ points exhaust $E[m]\setminus\{0\}$.

The first two conditions give a minimal integral equation with positive discriminant valuation and integral $j$.  It is therefore additive and potentially good, and the torsion conclusion follows.

\emph{Why the fields and the decision procedure are bounded.}
By the Cohen structure theorem \cite[Theorem~10.160.8]{Stacks032A}, a mixed-characteristic complete DVR with perfect residue field contains its Witt coefficient ring.  It is generated over that ring by a uniformizer: reduction modulo $p$ and completeness give the basis $1,\pi,\ldots,\pi^{e-1}$, and the relation is Eisenstein of degree $e$.  Conversely every Eisenstein relation used in the test lifts to such a DVR, and $R$ is its quotient modulo $p^q$.  All conditions above are stable under arbitrary lifts because $eq\geq H$.

For the necessity direction, start with its arbitrary perfect residue field $k_0$ and pass to the completed maximal unramified extension.  The absolute ramification remains $e$, rational torsion remains rational, and the new residue field is an algebraic closure $\overline{k_0}$.  The identification of the absolute Galois group of this completed extension with the original inertia group is stated explicitly in \cite[section~5.1.5, pp.~156--157]{Fontaine1994Semi}.  Serre--Tate's good-reduction criterion \cite[Theorem~1, p.~493]{SerreTate1968Good} therefore shows that a bad curve remains bad.  Its integral $j$ remains integral, so it remains additive and potentially good.  Choose a new minimal integral equation after this extension; no invariance claim about the old equation's minimality is needed.

The field $\overline{k_0}$ need not be algebraic over $\mathbf F_p$.  The passage from that field to the field $k=\overline{\mathbf F}_p$ in the test is instead a logical transfer of a finite formula over $\mathbf F_p$.  Addition and multiplication on $W_q$ are universal polynomial operations in $q$ residue-field coordinates, and reduction by the monic Eisenstein polynomial gives polynomial operations on $R$.  Thus the entire displayed test, including its existential and universal quantifiers, is a finite first-order polynomial formula with coefficients in $\mathbf F_p$.  Elimination in algebraically closed fields makes its truth independent of the algebraically closed field of characteristic $p$ used to evaluate it.  In particular a witness over $\overline{k_0}$ gives an affirmative test over $\overline{\mathbf F}_p$, without identifying these two residue fields.

One can describe that elimination procedure directly.  Replace inequations by inverses, express Boolean formulas as constructible sets, and compute existential images by elimination of polynomial ideals.  On each irreducible component of the closure of an image, a Gr\"obner basis over its function field supplies a nonzero denominator whose complement lies in the image; repeat on the remaining closed subset.  Dimension strictly decreases, so this is finite.  Complements handle universal quantifiers.  All coefficients lie in $\mathbf F_p$, so a nonempty constructible output has a point over $\overline{\mathbf F}_p$ with finitely many coordinates in one finite field.  Lifting their Witt coordinates to that finite unramified extension of $\mathbf Q_p$ and the Eisenstein coefficients to finite Witt polynomials produces the asserted finite local-field example.
\end{proof}

\par\smallskip\noindent\textbf{Exact computational certificate.}
The program \path{code/verify_finite_elliptic_test.py} reconstructs the division polynomials and computes their exact discriminants and two-torsion resultants for nine specializations, including the degree-$364$ polynomial for $m=27$ on $y^2=x^3-x$.  It also checks $2000$ integer-bound instances.  The original run took $74.857496652$ seconds, and the parent's independent rerun took $74.239452977$ seconds; both passed.  The retained output records the large precision bound at $(3,216,3)$ and explicitly records that its elimination test was not evaluated.  These finite checks support the identities and arithmetic used in the proof; they do not substitute for the universal argument. The formerly open index $216$ is instead decided by the independent construction in Theorem~\ref{thm:B-RD26-wild}.

\par\smallskip\noindent\textbf{Scope.}
This bounded test itself has not been evaluated at $(e,n)=(216,3)$. The simplified arithmetic characterization is now proved separately in Theorem~\ref{thm:B-RD26-wild} and Corollary~\ref{cor:B-RD27-exact-three}, without executing this elimination test.  It also does not decide Part~B.4: arbitrary truncated BT data with a finite wild descent action are not thereby supplied with an infinite equivariant BT lift.


\subsection{Why a ramification-nine base change of the full-nine examples does not solve the next gap}

\begin{proposition}\label{prop:B-RD25-basechange-obstruction}
\status{Proved; independent complete proof audit and exact certificate rerun by the parent passed.}
Let $K$ have absolute ramification index $24$, residue characteristic $3$, and perfect residue field.  If an additive elliptic curve $E/K$ has full rational $9$-torsion, then no extension of $K$ of ramification index $9$ makes its full $27$-torsion rational.  More generally, over its minimal good-reduction field after completed unramified base change, the next full division field has ramification degree at least $27$.

This rules out the proposed construction of an $(e,n)=(216,3)$ example by a ramification-nine base change of an $e=24$ full-nine example.  It does not rule out a new curve defined directly over a field of absolute index $216$; such a curve is constructed in Theorem~\ref{thm:B-RD26-wild}.
\end{proposition}
\begin{proof}
Pass first to completed unramified bases.  The proved tame criterion and the wild bound of Proposition~\ref{prop:B-RD20} force the semistability defect to be $12$.  Let $L/K$ be its minimal good-reduction field, so $v_L(3)=288$.  The good special fibre is supersingular.  Write $t$ for an integral formal parameter of the good model.

The universal valuation identity (B.12.4), together with the wild divisibility of the maximum valuation, gives the unique primitive-nine extrema
\[
                            A=21,\qquad B=1.
\]
Explicitly the only possible region is $r=1$, and
$288=6(A+27B)$ with $A\geq9B$, $3\mid A$, $B\geq1$
gives $(A,B)=(21,1)$; region $r=0$ would require
$A+3B=16$ with $3\mid A$, and is impossible.
Choose a basis $P,Q$ of $E[9]$ with $v_L(t(P))=21$ and $v_L(t(Q))=1$.

Formula (B.12.3) gives the minimum nonzero three-torsion valuation $9$.
The points of larger valuation form a line $C\subset E[3]$ of order three: two independent such points would make every nonzero three-torsion point have larger valuation.  The constant-term product for $[3](T)/T$ gives their common valuation $c$:
\[
                         288=6\cdot9+2c,\qquad c=117.
\]
All these torsion points are $L$-rational.

Let $R$ be a point with $3R=Q$, and let $S$ be any point with $3S=P$.
The Weierstrass factorization of $[3](T)$ is a unit times the product of $T-t(U)$ over $U\in E[3]$.  A point of valuation at least $9$ has image under $[3]$ of valuation at least $81$.  A point of smaller positive valuation $v$ has image of valuation $9v$.  Consequently
\[
                      v_L(t(R))=\frac19,\qquad
                      v_L(t(S))=\frac73.
\]
The field $F=L(R)$ has degree at most $9$, because $E[3]$ is rational; its ramification degree is at least $9$ by the valuation of $t(R)$.  Hence it is a totally ramified Galois extension of degree $9$.
The map $\sigma\mapsto\sigma R-R$ identifies
$\Gamma=\operatorname{Gal}(F/L)$ with $E[3]\simeq(\mathbf Z/3)^2$.
Put $z=t(R)$; it is a uniformizer of $F$.
Integral formal addition gives
\[
 v_F(\sigma z-z)=
 \begin{cases}
 1053,&0\ne\sigma R-R\in C,\\
 81,&\sigma R-R\notin C.
 \end{cases}
\]
Here the valuation of a formal difference equals that of the translating torsion point: the coefficient of that point in the difference is a unit.

Let $H\subset\Gamma$ correspond to $C$.  For every $0\ne\sigma\in H$ the lower ramification number is $1052$.  The elementary ramification estimate proved in Lemma~\ref{lem:B-RD20-wild} says
\[
 v_F(\sigma x-x)\geq v_F(x)+1052\qquad(x\in\mathfrak m_F).
\]
If $S$ belonged to $E(F)$, then $v_F(t(S))=21$; a nonzero translation of $S$ by three-torsion would have difference valuation either $81$ or $1053$, both less than $21+1052=1073$.  Thus every element of $H$ fixes $S$.  Set $F_1=F^H$; then $S$ would be $F_1$-rational, and $v_{F_1}(t(S))=7$.

We compute the ramification number of $F_1/L$ directly, without an appeal to a normalization of the upper numbering.  Its uniformizer
\[
                             w=N_{F/F_1}(z)
\]
satisfies $w=z^3+\epsilon$ with $v_F(\epsilon)\geq1055$.
For $\tau\notin H$, put $\delta=\tau z-z$, so $v_F(\delta)=81$.
Since $v_F(3)=2592$,
\[
 v_F\bigl((z+\delta)^3-z^3\bigr)
   =\min\{2592+2+81,\ 2592+1+162,\ 243\}=243.
\]
The correction $\tau\epsilon-\epsilon$ has valuation at least $1055$ and cannot cancel this term.  Hence
$v_{F_1}(\tau w-w)=81$, so the nontrivial automorphisms of $F_1/L$ have lower ramification number $80$.

For an element of valuation prime to $3$, the elementary ramification estimate is an equality: expand in the uniformizer with Teichm\"uller coefficients; the first term has coefficient equal to its exponent modulo $3$, which is nonzero.  Applied to $t(S)$ of valuation $7$, it gives
\[
                         v_{F_1}(\tau t(S)-t(S))=7+80=87.
\]
But $\tau S-S$ is a nonzero three-torsion point, so its formal coordinate has valuation $3\cdot9=27$ or $3\cdot117=351$ in $F_1$.  Neither equals $87$, a contradiction.
It follows that $F$ does not contain full $27$-torsion.  The full division field over $L$ is a $3$-extension, since the action is trivial modulo $9$; it contains $F$, and its degree is therefore at least $27$.

Finally, an extension $K'/K$ of ramification index $9$, followed by completed unramified base change, has degree $9$.  The compositum with $L$ has degree at most $9$ over $L$.  If it contained all $27$-torsion it would contradict the preceding argument.  This proves the asserted obstruction.
\end{proof}

\begin{remark}[Finite arithmetic certificate]\label{auto:B:c83d5f781371}
\status{Proved under stated hypotheses}

The exact certificate \path{code/verify_level27_basechange_obstruction.py}
enumerates every feasible $(A,B,r)$, checks the three-torsion and
next-division valuations, and verifies both ramification contradictions.
Its original run returned PASS in $0.006479513998783659$ seconds;
the parent independently reran it with PASS in $0.033169171001645736$ seconds.
These checks supplement the formal-group and field arguments above.
\end{remark}


\subsection{Wild equivariant arcs and the exact characteristic-three range}

\begin{lemma}[A universal wild trace invariant]
\label{lem:B-RD26-trace-arc}
\status{Proved; independent complete proof and primary-source audit by the parent passed, with an independent exact certificate rerun.}
Let $n\geq1$, and let $A\geq1$ be divisible by three. Put
\[
 d=2\cdot3^{n-1}\left(A+3^n+
        \sum_{j=1}^{n-1}2\cdot3^{n-j-1}\min(A,3^{2j})\right).
                                                        \tag{B-RD26.1}
\]
There is a good supersingular elliptic curve over a complete DVR
$R$ with algebraically closed residue field of characteristic three,
absolute ramification index $d$, and full rational $3^n$-torsion,
whose primitive torsion parameters have extrema $A,1$.
Moreover $R$ and this marked curve admit faithful semilinear actions
of the following groups, fixing every marked torsion point:
\[
 C_3;\qquad C_6\ \text{if }A\text{ is odd};\qquad
 C_3\rtimes C_4\ \text{if }A\equiv1\pmod4.
\]
The action on the special fibre is the usual faithful subgroup of
$\operatorname{Aut}(Y^2=X^3-X)$, and all these actions are trivial
on the residue field. The order-four generator acts on $C_3$ by
inversion.
\end{lemma}

\begin{proof}
Let $k$ be an algebraic closure of $\mathbf F_3$, and work over
$W=W(k)$. The elliptic curve
\[
       \mathcal E_u:\quad Y^2=X^3+uX^2-X
                                                        \tag{B-RD26.2}
\]
over $W[[u]]$ is a universal framed deformation of its special
fibre $E_0$. Here the framing identifies the special fibre with
$E_0$, and isomorphisms respect that identification. To check the
parameter directly, complete the square in any deformation, move
the simple root lifting $X=0$ to zero, and take the unique fourth
root reducing to one to make the coefficient of $X$ equal to $-1$.
The only remaining coefficient is $u$. The simple-root and
fourth-root operations are unique because two and four are units.
This also proves the required uniqueness for coordinate changes
reducing to the identity. Its discriminant is $16(u^2+4)$, a unit.
Serre--Tate identifies this deformation with that of the height-two
formal group; the precise nilpotent-base theorem is
\cite[Theorem~1.2.1]{Katz1981SerreTate}.

Use the odd formal parameter $T=X^{-1/2}$ at infinity, with the
choice of sign for which
\[
 Y=-T^{-3}(1+uT^2-T^4)^{1/2},\qquad
 \omega=\frac{dX}{2Y}=
                  \frac{dT}{(1+uT^2-T^4)^{1/2}}.
\]
The coefficient of $T^3$ in $[3](T)$ is exactly $4u$.
Indeed the formal logarithm begins $T-uT^3/6$, so substitution in
$\log([3](T))=3\log(T)$ gives this coefficient. This identity is
an identity of integral coefficients, although it was calculated
after inverting three.

Let $R_n$ be the full Drinfeld-level deformation ring, and let
$x,y$ be the parameters of a universal basis of $3^n$-torsion.
It is a complete regular local ring with these two regular
parameters, by \cite[Theorem~2.1.2 and Proposition~2.2.3(i)]
{CarterStrauch2021LubinTate}. The universal family in (B-RD26.2)
pulls back to this ring. Changing the special-fibre framing gives
an action on $R_n$ and its family of
\[
 G=\langle\sigma,\tau\rangle=C_3\rtimes C_4,
\]
where on $E_0$ the first generator translates $X$ by one and the
second is the order-four automorphism with $T\mapsto iT$,
$i^2=-1$. This is the functorial action by precomposition of the
framing; its isomorphisms preserve the universal level structure.
In the explicit family, the second action is $u\mapsto-u$ and
$T\mapsto iT$ exactly. The first action reduces to
\[
              h(T)=\frac{T}{(1+T^2)^{1/2}}.
\]
The isomorphisms preserve the origin. Consequently, for $j=0,1,2$,
the difference between the series for $\sigma^j(T)$ and
\[
                 h_j(T)=\frac{T}{(1+jT^2)^{1/2}}
\]
belongs to $(3,u)T W[[u,T]]$. The factor $T$ here is essential:
it follows from preservation of the zero section, in addition to
agreement of the special-fibre series.

We record two bounds in $R_n$ needed to make the forthcoming
rational expressions integral on a chart. Set
\[
               M=3^{2n-2},\qquad H=6M.
\]
Every three-torsion coordinate belongs to
$\mathfrak m_{R_n}^{M}$, by iterating the prime factorization in
\cite[Proposition~2.2.3(vii)]{CarterStrauch2021LubinTate}.
Weierstrass preparation factors $[3](T)$ as a unit times
$T\prod_{0\ne Q\in E[3]}(T-t(Q))$. Its coefficient of $T^3$
is a sum of terms containing at least six of these coordinates.
The formula $[3]_{T^3}=4u$ therefore proves
\[
                    u\in\mathfrak m_{R_n}^{H}.          \tag{B-RD26.3}
\]
Likewise $3\in\mathfrak m_{R_n}^{8M}$. In particular $H\geq6$.
The special formal series is $[3](T)=-T^9$. For an explicit
algebraic check, the three-division polynomial of (B-RD26.2) is
$3X^4+4uX^3-6X^2-1$, and its tripling numerator reduces to $X^9$
modulo $(3,u)$. Substitution $X=T^{-2}$ with the square-root
branch whose leading coefficient is three gives the asserted
series. The sign is also checked directly by the characteristic-three
addition law: it gives $Y([3]P)=-Y^9$, using the polynomial identity
\[
 ((X^3-X)^2-1)(X-X^9)-(X^3-X)+(X^3-X)^5=0
                                      \quad\text{in }\mathbf F_3[X].
\]
Hence the successive coordinates $[3^j]y$ have weighted
initial forms $(-1)^jY^{3^{2j}}$ whenever $j<n$: the error ideal
$(3,u)$ has weight at least $H>M$, while these weights are at
most $M$.

Here is a precise weighted chart construction, including the
possible coefficients divisible by three in a Cohen presentation.
Give $x,y,3$ weights $A,1,d$. The factorization argument in the
proof of Theorem~\ref{thm:B-RD12}, which does not use tame descent
at this stage, gives lifts of all primitive torsion coordinates
with the prescribed weighted initial forms: a nonzero multiple
of $Y$, or the lower-weight term of $X+cY^{3^{2j}}$, or that
binomial at an equality of weights. Choose $c\in k^\times$ so
none of these finitely many initial forms vanishes at $(c,1)$.
In fact one may choose $c=i$ uniformly. All nonzero scalar
coefficients of the boundary binomials are in $\mathbf F_3^\times$
by the preceding explicit series and the linear terms of formal
addition and multiplication by units. The only excluded values
at a boundary are therefore $1,-1$, whereas $i^2=-1$ makes
$i\ne\pm1$. This optional choice fixes $\lambda=-1$ below.
The relation between three and the product of all primitive
coordinates supplies a generator of the Cohen kernel
\[
          R_n\simeq W[[X_0,Y_0]]/(3-P(X_0,Y_0)).        \tag{B-RD26.4}
\]
The series $P$ is a unit times the product of the chosen lifts.
It has weight at least $d$, and its weight-$d$ initial form,
evaluated at $(c,1)$, is a nonzero element $v_0\in k$.
Terms of a lifted factor divisible by three have weight at least
$d$, strictly greater than the weight of that factor. In the
whole product they consequently have weight strictly greater
than $d$. Thus the weight-$d$ part is independent of the variable
used below for $3/\pi^d$. This explicitly deals with low ordinary
degree monomials whose coefficients are divisible by three.
The generator in (B-RD26.4) has coefficient one at its linear
term $3$, so it generates the Cohen kernel, rather than merely
an ideal contained in it.

Adjoin the chart variables
\[
          x=\pi^A Z,\qquad y=\pi,\qquad3=\pi^d V,
                                                        \tag{B-RD26.5}
\]
where $Z,V$ are units, and complete at $Z=c,V=v_0,\pi=0$.
This chart has the two relations
\[
           3=\pi^d V,\qquad
           V=P(\pi^A Z,\pi)/\pi^d.                    \tag{B-RD26.6}
\]
The division in the second relation is legitimate: expand each
coefficient of $P$ in Teichmuller digits and replace each factor
$3^b$ by $\pi^{db}V^b$ using the first relation. The preceding
weighted bound gives a convergent series with nonnegative powers
of $\pi$. Its reduction is a polynomial in $Z$ alone, with value
$v_0$ at $c$. The second equation therefore has derivative one
with respect to $V$ modulo the maximal ideal. Equivalently, this
is the completed chart obtained by adjoining $x/y^A$ and $3/y^d$
to $R_n$ and removing the exceptional torsion. Before imposing
the next equation it is a complete regular local ring of
dimension two, with parameters $\pi,Z-c$. Formula (B-RD26.3)
gives $u\in\pi^H$ on this chart, while $3\in\pi^d$ and
$d\geq8M>H$.

The $C_3$ action already extends to this chart, before any trace
is formed. Indeed $\sigma x/x$ and $\sigma y/y$ are units of
$R_n$, and
\[
 \sigma Z=(\sigma x/x)(\sigma y/y)^{-A}Z,\qquad
 \sigma V=(\sigma y/y)^{-d}V.
\]
Both displayed multipliers reduce to one. The action therefore
preserves the centre $Z=c,V=v_0$ and extends continuously to its
completion. These formulas come from the actual ratios in the
fraction field, so they respect both chart relations.

In its fraction ring put
\[
 z_x=x^{-2},\quad z_y=y^{-2},\quad
 D_y=\operatorname{Tr}_{\langle\sigma\rangle}(z_y^2),
 \quad
 I=\frac{\operatorname{Tr}_{\langle\sigma\rangle}
                      (z_y^2(z_x-z_y))}{D_y},\quad
 U_y=N_{\langle\sigma\rangle}(y)^{-2}.                \tag{B-RD26.7}
\]
These are exact $C_3$-invariants. We verify all divisions and
error estimates, interpreting $O(\pi^b)$ inside the fraction
ring of this chart. The origin-preserving series above give
\[
 (\sigma^j x)^{-2}=x^{-2}+j+O(\pi^{H-2A}),\qquad
 (\sigma^j y)^{-2}=y^{-2}+j+O(\pi^{H-2}).
\]
It follows that
\[
 \begin{split}
 D_y&=3y^{-4}+6y^{-2}+5+O(\pi^{H-4})
                         =-1+O(\pi^{H-4}),\\
 I&=z_x-z_y+O(\pi^{H-2A-4}),\\
 U_y&=y^{-6}-y^{-2}+O(\pi^{H-6}).
 \end{split}                                           \tag{B-RD26.8}
\]
For the middle equality, subtract $(z_x-z_y)D_y$ before
estimating the numerator: each remaining factor has valuation
at least $H-2A-4$. For the last equality use
$\prod_{j=0}^2(y^{-2}+j)=y^{-6}+3y^{-4}+2y^{-2}$;
its difference from $y^{-6}-y^{-2}$ has a factor three.
Since $H>4$, $D_y$ is an actual unit of the chart. These
calculations also show that all the normalized expressions used
next are convergent integral series, not just valuation
estimates after an unspecified specialization.

Put $a=A/3$, and take $\lambda\in W^\times$ with residue $c^{-2}$.
Impose the exact invariant equation $I=\lambda U_y^a$.
After multiplication by $x^2$ this is an integral equation on
the chart, and (B-RD26.8) gives its explicit form
\[
 x^2(I-\lambda U_y^a)
   =1-\pi^{2A-2}Z^2-\lambda Z^2(1-\pi^4)^a
                                  +O(\pi^{H-4}).       \tag{B-RD26.9}
\]
Its reduction is $1-\bar\lambda Z^2$. At $Z=c$ its derivative
is $-2\bar\lambda c\ne0$. The formal implicit-function theorem,
applied together with the second equation of (B-RD26.6), now
eliminates $Z$ and $V$. The resulting ring $R$ is a complete DVR
with uniformizer $\pi$ and $3=\pi^d$ times a unit. This is a
simple-root argument over a complete local ring; no division by
three or averaging by the wild group order was used.
The choice of $c$ guarantees that every primitive coordinate has
exactly its assigned valuation, so the extrema are $A,1$.

Finally the chart and this branch are stable under the asserted
groups. Since every coordinate isomorphism preserves the
origin, $\sigma x/x$ and $\sigma y/y$ are units. They are both
one modulo $\pi$; hence $\sigma$ fixes the residue of $Z=x/y^A$
and of $V=3/y^d$. The rational equation is exactly
$C_3$-invariant, so its unique simple branch is stable.
For $\tau$, $x,y$ have eigenvalue $i$ exactly. Conjugation
permutes the three factors in the trace and norm, giving
\[
             \tau I=-I,\qquad \tau U_y=-U_y.
\]
If $A\equiv1\pmod4$, then $a=A/3$ is odd, the equation is
stable, and $\tau(Z)\equiv i^{1-A}c=c$. Also $4\mid d$ in
this case. For the order-two element $\tau^2$, stability of the
equation is automatic; its branch is stable when $A$ is odd,
and then $2\mid d$. The wild action is faithful on $R$:
\[
                 \sigma\pi=\pi-\tfrac12\pi^3
                                      +O(\pi^5),
\]
because the deformation error has order at least $H+1>3$.
The tame action has the stated faithful eigenvalue on $\pi$.
All actions fix $W$ and the residue field. The generic marked
elliptic curve and its level structure descend over
$\operatorname{Frac}(R)^G$ (or the indicated subgroup's fixed
field). This last assertion concerns the generic fibre; the
integral good model over $R$ is used for its faithful special
inertia action. This proves every assertion of the lemma.
\end{proof}

\begin{remark}[Exact finite certificate]\label{auto:B:10f3741cc5dc}
\status{Proved under stated hypotheses}

The accompanying program \path{code/verify_wild_trace_arcs.py}
checks the trace and norm identities, the special tripling
identity, the logarithmic coefficient $[3]_{T^3}=4u$, the
equivariance of the residual curve, and the exact integer
construction and least-index formula over a finite range.
It checks $66$ symbolic equivariance identities, $109727$ constructed
integer arcs, and $25000$ least-index comparisons through $e=50000$
and $n=8$. The final sign-expanded author run passed in
$2.5154165040003136$ seconds; the independent parent and
obstructions/crystals audits passed in $2.827041984997777$ and
$2.6579352160006238$ seconds respectively. Both independent
audits also checked the complete proof and original sources.
It does not computationally instantiate the universal formal
deformation or its descent; those are proved above.
\end{remark}

\begin{theorem}[B-RD26: the sharp wild higher-torsion criterion]
\label{thm:B-RD26-wild}
\status{Proved; independent parent proof, source, and exact arithmetic audits passed.}
Let $p=3$, $e$ be a positive even integer, $s=v_3(e/2)$, and
$m\in\{3,6,12\}$. An additive elliptic curve with full rational
$3^n$-torsion and exact semistability defect $m$ exists over a
complete DVR of absolute ramification index $e$ and perfect
residue field if and only if
\[
 s\geq n-1,\qquad
 e\geq\frac8m\,3^{\max(2n-2,\,4n-4-s)}.              \tag{B-RD26.10}
\]
The fields constructed here have algebraically closed residue
field. No finite-residue-field assertion is needed for this
existence criterion.
\end{theorem}

\begin{proof}
Necessity is Proposition~\ref{prop:B-RD20}. Suppose the two
inequalities hold. Start at $r_0=\max(0,2n-2-s)$ and increase
$r$ until $r=n-1$ or the inequality
\[
                 me\geq8\cdot3^{2n+r-2}
\]
would fail on increasing $r$ by one. Define
\[
       A=\frac{me}{2\cdot3^{2n-r-2}}-3^{2r+1},
       \qquad B=1.                                    \tag{B-RD26.11}
\]
Then $3\mid A$, $A\geq3^{2r}$, and, unless $r=n-1$,
$A<3^{2r+2}$. The divisibility by three follows from
$v_3(me/2)=s+1$ and $2n-r-2\leq s$.
The geometric-series identity (B-RD20.5) therefore says that
(B-RD26.1) equals $me$. For $m=6$, $A$ is odd; for $m=12$,
$A\equiv1\pmod4$. Apply Lemma~\ref{lem:B-RD26-trace-arc}
and the subgroup of order $m$. Its fixed field has absolute
index $me/m=e$. Galois descent is effective for the elliptic
curve, for instance using the invariant divisor $3[0]$.
The marked torsion points descend and are all rational.
The special-fibre action is faithful of order $m$, so the
good-reduction criterion gives exact semistability defect $m$.
In particular reduction is additive and potentially good.
\end{proof}

\begin{corollary}[B-RD27: the complete higher elliptic range]
\label{cor:B-RD27-exact-three}
\status{Proved; the full wild construction and integer criterion have passed independent audit.}
For $p=3$, write $e=2u$, $s=v_3(u)$ and
$L=\lfloor\log_3u\rfloor$. The largest level admitting an additive
elliptic curve with full rational $3^n$-torsion is
\[
 N(e)=1+\min\left(s,\left\lfloor\frac{L+1}{2}\right\rfloor,
                         \left\lfloor\frac{L+s+1}{4}\right\rfloor\right).
                                                        \tag{B-RD27.1}
\]
Equivalently, existence at level $n$ is exactly
\[
  v_3(e/2)\geq n-1,\qquad
  e\geq\frac23\,3^{\max(2n-2,4n-4-v_3(e/2))}.          \tag{B-RD27.2}
\]
Among potentially good elliptic curves, the least positive full
torsion level forcing good reduction is $N(e)+1$.
In particular full $27$-torsion with additive reduction exists
at $e=216$, with exact defect twelve. Together with the already
proved tame criterion, these formulas decide the elliptic
higher-torsion existence question at every odd prime.
\end{corollary}

\begin{proof}
Defect twelve gives the weakest wild inequality and is always
available. Its realized maximum is the upper bound $N_{\mathrm w}$
already calculated in Corollary~\ref{cor:B-RD21}. The tame maximum
is no larger. The final index statement uses the potentially good
hypothesis; multiplicative Tate examples retain the unrestricted
qualification in Proposition~\ref{prop:B-RD15}.
\end{proof}


\subsection{A sharper wild lower bound for the character problem}

\begin{proposition}[B-RD28: wild character ambiguity in every higher height]
\label{prop:B-RD28-wild-character}
\status{Proved; independent parent proof and indexing audit passed.}
Let $p=3$, let $e$ be a positive even integer, and put
\[
 s=v_3(e/2),\qquad L=\lfloor\log_3(e/2)\rfloor,\qquad
 N_{\mathrm w}=1+\min\left\{s,
     \left\lfloor\frac{L+1}{2}\right\rfloor,
     \left\lfloor\frac{L+s+1}{4}\right\rfloor\right\}.
\]
For every height $h\geq2$ there are two height-$h$ BT groups
$M,N$ over one complete discretely valued field $K$ of absolute
ramification $e$, with perfect residue field and $\mu_3\subset K$,
such that both groups acquire good reduction,
\[
                 M[3^{N_{\mathrm w}}]\simeq
                 N[3^{N_{\mathrm w}}],
\]
and their manuscript characters differ on a prime-to-three
inertia element. Consequently the least universal B.4 index,
if denoted by $\nu_{\min}(3,e,h)$, satisfies
\[
                       \nu_{\min}(3,e,h)\geq N_{\mathrm w}.
                                                        \tag{B-RD28.1}
\]
In particular the lower bounds at $(e,h)=(24,2)$ and $(216,2)$
are respectively two and three.
\end{proposition}

\begin{proof}
Apply Theorem~\ref{thm:B-RD26-wild} with $n=N_{\mathrm w}$ and
$m=12$. The construction gives an elliptic curve $E/K$ with
full rational $3^n$-torsion and faithful special inertia group
$C_3\rtimes C_4$ on its supersingular good reduction
$Y^2=X^3-X$. By the Weil pairing, $\mu_{3^n}\subset K$.
Set
\[
 M=E[3^\infty]\times(\mathbf Q_3/\mathbf Z_3)^{h-2},
 \qquad N=(\mathbf Q_3/\mathbf Z_3)^h.
\]
These generic BT groups have isomorphic constant $3^n$-torsion.
The second group has good reduction already over $\mathcal O_K$;
the first acquires good reduction over the degree-twelve field
of the construction. Thus both satisfy the manuscript monodromy
condition with a single good-after-extension graded piece.

For the order-four automorphism $\tau$ of the supersingular
elliptic reduction, $\tau^2=[-1]$. Its two eigenvalues on the
rational covariant Dieudonne module are $i,-i$; one may see the
opposite pair either from the degree-one determinant or from
Frobenius interchanging the two primitive fourth roots in
characteristic three. Hence its trace is zero. On $M$ the trace
is $h-2$, whereas on $N$ it is $h$. The period twist in the
corrected Fontaine dictionary has trivial finite inertia action,
so these are exactly the manuscript character values.
The characters therefore differ. Failure at
$3^n=M(n-1)$ gives $\nu_{\min}\geq n$, with precisely the indexing
specified in the manuscript.
\end{proof}

\begin{remark}[Scope of the lower bound]\label{auto:B:688dec0e6255}
\status{Proved under stated hypotheses}

This improves the tame Honda bound in
Corollary~\ref{cor:B-RD18} in the wild exceptional ranges.
It does not assert that the displayed lower bound is the least
universal index: the latter concerns arbitrary pairs of groups
with isomorphic truncations, and the comparison group need not
be constant. The independent characteristic-two example in
Proposition~\ref{prop:B-RD22} already illustrates that distinction.
The finite upper bound in Theorem~\ref{thm:B-RD17} continues to
apply to every pair satisfying the full filtration condition.
The numerical formula here is included among the $25000$ exact
least-index comparisons in the certificate for
Theorem~\ref{thm:B-RD26-wild}; no new unexecuted computation is
being used.
\end{remark}


\subsection{A level-nine character ambiguity at the smallest odd ramification}

\begin{lemma}[Cyclic isogenies and a Tate comparison]
\label{lem:B-RD29-isogeny-Tate}
\status{Proved under stated hypotheses}
Let $K$ be a complete discretely valued field of characteristic zero,
let $r\geq2$, and let $E/K$ have a $K$-rational point $P$ of exact
order $r$. Put $B=E/\langle P\rangle$. There is a split Tate curve
$T_q/K$, with a parameter $q\in K^\times$ of positive valuation,
such that $B[r]\simeq T_q[r]$ as finite group schemes over $K$.
The generic $p$-divisible group of a Tate curve satisfies the
manuscript monodromy condition: its two graded pieces are
$\mu_{p^\infty}$ and $\mathbf Q_p/\mathbf Z_p$.
\end{lemma}
\begin{proof}
The dual of the quotient isogeny is a morphism
$\widehat\phi:B\to E$ with kernel canonically Cartier dual to
$\langle P\rangle\simeq\mathbf Z/r$. Its restriction gives an
exact sequence of finite etale $K$-group schemes
\[
 0\longrightarrow\mu_r\longrightarrow B[r]
   \xrightarrow{\widehat\phi}\langle P\rangle
   \longrightarrow0.                                      \tag{B-RD29.1}
\]
It is exact on the right: an inverse image of $P$ under
$\widehat\phi$ is killed by $r$, since
$[r]=\phi\widehat\phi$ on $B$ and $\phi(P)=0$.
Identify the last group with $\mathbf Z/r$ using its generator $P$.
The sequence is killed by $r$, so it splits as a sequence of
abstract $\mathbf Z/r$-modules. Its Galois descent class lies in
\[
 H^1(K,\mu_r)=K^\times/(K^\times)^r.
\]
Choose a representative $q_0$ and multiply it by an $r$th power
of a sufficiently high power of a uniformizer to obtain $v_K(q)>0$.
The finite group associated to this Kummer class is exactly
$T_q[r]$. This is the explicit finite group scheme
$G_{K,q,r}$ in \cite[``Points of finite order'', pp.18--19]{Tate1995Uniformization};
its component over $a\in\mathbf Z/r$ is given by $z^r=q^a$.
That source proves both its extension-class description and its
identification with the torsion of the Tate curve. The latter
curve exists by its Theorem~1, p.6.

For $r=p^n$ these descriptions at every level yield
\[
 0\longrightarrow\mu_{p^\infty}\longrightarrow T_q[p^\infty]
   \longrightarrow\mathbf Q_p/\mathbf Z_p\longrightarrow0.
\]
Both graded groups have good models over $\mathcal O_K$.
Their finite inertia characters are trivial, so the manuscript
character of $T_q[p^\infty]$ is constantly two on inertia.
\end{proof}

\begin{proposition}[B-RD29: nine-torsion fails already at $p=3,e=2$]
\label{prop:B-RD29-p3-level-two}
\status{Proved; independent parent complete proof, original-source and exact-certificate audits passed.}
Let $K=\mathbf Q_3(\zeta_3)$, whose absolute ramification is two
and whose residue field is $\mathbf F_3$. Define
\[
 B:\quad y^2-xy+y=x^3-x^2-28x-53,
 \qquad q=2^5\,3^6=23328,
\]
and let $T_q/K$ be the split Tate curve with this parameter.
Then
\[
 B[9]\simeq T_q[9],                                      \tag{B-RD29.2}
\]
but the manuscript characters of $B[3^\infty]$ and
$T_q[3^\infty]$ differ on an inertia element of order two in
the finite inertia image. The first group is potentially good;
the second satisfies the filtration condition with good rank-one
graded pieces. Consequently
\[
                 \nu_{\min}(3,2,h)\geq2\qquad(h\geq2).
                                                        \tag{B-RD29.3}
\]
The same lower bound holds for every $e\equiv2\pmod4$.
\end{proposition}

\begin{proof}
We give the quotient, the Kummer class and the inertia calculation
explicitly. Start with
\[
 E:\quad y^2-3xy-12y=x^3-12x^2,
 \qquad P=(0,0).
\]
Exact addition gives
\[
 \begin{array}{c|rrrrrrrr}
 a&1&2&3&4&5&6&7&8\\\hline
 x(aP)&0&12&4&6&6&4&12&0\\
 y(aP)&0&48&8&18&12&16&0&12
 \end{array}
\]
and $9P=0$. In particular $P$ has exact order nine.
The quotient by this subgroup is
\[
 E':\quad y^2-3xy-12y=x^3-12x^2-2310x-36190.
                                                        \tag{B-RD29.4}
\]
Here is a direct verification of the quotient formula, so that
no equation is imported without checking its isogeny. In short
coordinates $X=x-13/4$, $Y=y-(3x+12)/2$, put
\[
 A=-219/16,\quad C=827/32,\quad F(X)=X^3+AX+C,
 \quad\mathcal R=\{-13/4,35/4,3/4,11/4\}.
\]
Define
\[
 R(X)=X+\sum_{a\in\mathcal R}
   \left(\frac{2(3a^2+A)}{X-a}
              +\frac{4F(a)}{(X-a)^2}\right).
\]
Exact polynomial arithmetic gives
\[
 F(X)R'(X)^2=R(X)^3-\frac{37179}{16}R(X)
                         -\frac{1397493}{32}.
                                                        \tag{B-RD29.5}
\]
The reduced numerator and denominator of $R$ have degrees nine
and eight, and the denominator has double zeros precisely at
$\mathcal R$. Thus $(X,Y)\mapsto(R(X),YR'(X))$ extends to a
pointed degree-nine morphism of smooth elliptic curves, hence an
isogeny, and its kernel consists of the nine displayed points.
The target short equation is exactly (B-RD29.4) in the original
long coordinates. Finally $x=9x_1+1$, $y=27y_1+21$ changes
(B-RD29.4) to the displayed equation for $B$. This identifies
$B$ with $E/\langle P\rangle$ over $\mathbf Q$.

To compute the extension class in (B-RD29.1), consider
\[
 f(x,y)=\frac{(x^2-y)^3(y-3x+4)}{(x-4)^3}.
                                                        \tag{B-RD29.6}
\]
The tangent and chord construction gives
\[
 \operatorname{div}\left(\frac{x^2-y}{x-4}\right)
       =3[P]-[3P]-2[0],
 \qquad
 \operatorname{div}(y-3x+4)=3[3P]-3[0].
\]
For the first equality, the usual two addition lines are $y=0$
and $y=4x$, and the equation of $E$ gives
$y(y-4x)=(x-12)(x^2-y)$. For the second, substitution of
$y=3x-4$ into the equation gives $-(x-4)^3$.
Consequently $\operatorname{div}(f)=9[P]-9[0]$.

Let $\delta:E(K)\to K^\times/(K^\times)^9$ be the connecting
homomorphism of the dual isogeny. The precise divisor description
of this map is \cite[Appendix~A.1, Lemmas~A.1--A.2, p.22]{AntieauAuel2021Descent}.
There is a constant $c$ such that $cf$ is normalized for that
map. Since $5P$ and $4P$ avoid both its zero and its pole,
homomorphicity cancels this constant and gives
\[
 \delta(P)=\frac{f(5P)}{f(4P)}
       =\frac{-3456}{2916}
       =-\frac{32}{27}
       \equiv2^5\,3^6\pmod{(K^\times)^9}.               \tag{B-RD29.7}
\]
The last ratio is the ninth power $(-3)^9$. One can also compare
this calculation directly with the $N=9$ family and Figure~9 on
pp.26--27 of that source, specialized at $\lambda=1$.
Thus (B-RD29.1) is the Kummer extension with the specified
$q=23328$, whose valuation over $K$ is twelve. The preceding
lemma proves (B-RD29.2) with this explicit Tate parameter.

The exact invariants of $B$ are
\[
 c_4=17\,3^4,\qquad \Delta=-2\,3^9,
 \qquad j=-\frac{17^3\,3^3}{2}.                          \tag{B-RD29.8}
\]
The integral $j$ gives potential good reduction by
\cite[Section~2, p.496]{SerreTate1968Good}.
Over the completed maximal unramified extension of $K$, take
an extension of minimal good-reduction degree $m$. Its inertia
acts faithfully on the good special elliptic curve.
The displayed equation has discriminant valuation eighteen.
A change to a unit-discriminant good equation over the extension
therefore implies
\[
                            12\mid18m,
\]
so $m$ is even. The inertia image contains an element of order
two. In residue characteristic three the only nontrivial
order-two origin-preserving elliptic automorphism is $[-1]$;
it acts as $-1$ on the rational Dieudonne module and has trace
$-2$. The Tate comparison has trace $2$, by its two good graded
pieces. Hence the requested characters differ.
The inertia calculation uses the actual finite descent action
of the good model and therefore has the conventions of
Proposition~\ref{prop:B-RD4-dictionary}.

For $e\equiv2\pmod4$, take an extension of $K$ generated by a
root of $X^{e/2}-(\zeta_3-1)$. This polynomial is Eisenstein,
so its absolute ramification is $e$ and residue field remains
$\mathbf F_3$. The isomorphism (B-RD29.2) survives base change.
The raw discriminant valuation becomes $9e$, and
$12\mid9em$ still forces $m$ even. Thus the same trace
comparison holds. Adding $h-2$ identical constant etale factors
to both groups gives every $h\geq2$. Finally level nine is
$M(1)$, so its failure proves $\nu_{\min}\geq2$.
\end{proof}

\begin{remark}[Exact certificate and remaining question]\label{auto:B:7f506124b28c}
\status{Proved under stated hypotheses}

The program \path{code/verify_p3_character_level2.py} verifies all
nine subgroup points and all $81$ additions, the degree-nine
isogeny identity, the change to $B$, both discriminants and
$j$-invariants, the two divisor identities, the Kummer ratio and
its ninth-power adjustment. Its recorded exact run passed in
$4.882327191000513$ seconds; the independent parent rerun passed
in $4.293367760001274$ seconds. The program verifies finite
algebraic assertions; the Tate extension dictionary and inertia
argument are proved above.
This proves that the full-torsion bound $N_{\mathrm w}=1$ at
$e=2$ is strictly smaller than the least B.4 index. Whether
$\nu=2$ already suffices for arbitrary height-two groups at
$(p,e)=(3,2)$ is settled affirmatively in Theorem~\ref{thm:B-RD42}.
\end{remark}


\subsection{A logarithmic universal upper bound}

\begin{proposition}[B-RD30: the sharper generic-Hom exponent]
\label{prop:B-RD30}
\status{Proved; independent proof and primary-source audit passed.}
Use the integer $D_{p,h}$ of Lemma~\ref{lem:B-RD17-degree} and put
\[
 c_{p,e,h}=\left\lfloor\log_p\left(\frac{peD_{p,h}}{p-1}\right)
                  \right\rfloor.
\]
Then a sufficient universal index in B.4 is
\[
                    \nu=2h^2c_{p,e,h}.                 \tag{B-RD30.1}
\]
In particular, the previously established quadratic logarithmic
bound in Theorem~\ref{thm:B-RD17} can be replaced by this bound.
This is an upper bound, not a formula for the least index.
\end{proposition}
\begin{proof}
Bondarko's original Theorems~A and~B
\cite[p.661]{Bondarko2006Generic} assert the following, for every
prime, over a complete discrete valuation ring of absolute
ramification $a$: every homomorphism between generic fibres of
finite flat commutative group schemes extends after multiplication
by $p^c$, where
\[
                  c=\lfloor\log_p(pa/(p-1))\rfloor.
\]
These theorems apply to arbitrary finite flat models, not just
truncated BT groups. The proof is in \S4, pp.681--684; the passage
to arbitrary residue fields is stated on pp.662 and 681. The
stronger value of this exponent is also explicitly confirmed in
\cite[p.2 of the author version]{VasiuZink2012Boundedness}.

Choose the common good-reduction field $L$ from
Lemma~\ref{lem:B-RD17-degree}. Then $e_L\leq eD_{p,h}$, so the
exponent $c$ in Lemma~\ref{lem:B-RD17-common} can be taken at most
$c_{p,e,h}$. Apply that lemma at
$n=2h^2c_{p,e,h}+1$. Its Zassenhaus pieces are finite flat and
therefore satisfy exactly the hypotheses just stated. Finally,
$M(\nu)=M[p^{\nu+1}]$ gives (B-RD30.1).
\end{proof}

\begin{remark}[A useful stronger comparison for actual formal BT groups]\label{auto:B:cb3f24f10664}
\status{Proved under stated hypotheses}

Bondarko's Theorem~E(1), pp.663--664, proved from Proposition~18
on p.685 of \cite{Bondarko2006Generic}, has a different conclusion:
if two finite-height formal groups over $\mathcal O_L$ have
isomorphic generic $p^{\ell+u}$-torsion and
$e_L\leq(p-1)p^u$, then their $p^\ell$-torsion is isomorphic over
$\mathcal O_L$, compatibly with the original generic
identification. Uniqueness of a generic-fibre extension makes this
comparison equivariant for any given descent actions. This
stronger statement will be used only for actual formal BT groups.
It is not being applied to the possibly non-BT Zassenhaus pieces.
\end{remark}


\subsection{The smallest odd-prime case reduces to two levels}

\begin{lemma}[Finite inertia in a height-two formal group at $3$]
\label{lem:B-RD31-finite}
\status{Proved under stated hypotheses}
Let $G/K$ have potentially good reduction to a one-dimensional
formal BT group of height two, with residue characteristic $3$.
After passage to the completed maximal unramified extension, it
has good reduction over a Galois extension of degree at most $12$.
Its faithful special inertia image is either cyclic of order
dividing $8$, or has a normal subgroup of order $3$ and order at
most $12$; in the latter case its tame quotient has order $1,2$,
or $4$. Two such groups have simultaneous good reduction over an
extension of ramification degree at most $36$.
\end{lemma}
\begin{proof}
We work with algebraically closed residue field. Lazard's
classification theorem \cite[Theorem~IV, pp.270--271]{Lazard1955Formal}
identifies the special formal group with the height-two Honda
group. Its rational Dieudonne module has a basis in which
\[
 F=\begin{pmatrix}0&3\\1&0\end{pmatrix}\sigma.
\]
This is also obtained directly from the Dieudonne module of the
supersingular elliptic curve $Y^2=X^3-X$; a change of basis over
$W(\bar k)[1/3]$ absorbs the sign in $F^2=-3$.
The endomorphisms commuting with $F$ are precisely
\[
 A(a,c)=\begin{pmatrix}a&3\sigma(c)\\c&\sigma(a)\end{pmatrix},
       \qquad a,c\in\mathbf Q_{3^2}.
                                                        \tag{B-RD31.1}
\]
Here $\mathbf Q_{3^2}=W(\mathbf F_9)[1/3]$ and $\sigma$ is its
nontrivial automorphism. The determinant is
$N(a)-3N(c)$, whose two nonzero terms have different parity of
$3$-adic valuation. Thus this algebra is a quaternion division
algebra. Its valuation ring consists of the matrices with
$a,c\in W(\mathbf F_9)$; its residue field is $\mathbf F_9$, and
the kernel of reduction on units is pro-$3$. These facts also
follow by writing the algebra as
$\mathbf Q_{3^2}\langle\Pi\rangle$, with
$\Pi^2=3$ and $\Pi a=\sigma(a)\Pi$.

Let $H$ be any finite subgroup of its units. Its reduction has
cyclic image in $\mathbf F_9^\times$. Its kernel $P$ is a finite
$3$-group. If $P\ne1$, choose a central element $z$ of order $3$.
Its centralizer is the quadratic field $\mathbf Q_3(z)$, so all
of $P$ lies in that field. Its finite $3$-power roots of unity
form $\mu_3$: an element of order $9$ would have minimal
polynomial of degree $6$, whereas every element of the
quaternion algebra satisfies a quadratic polynomial. Hence
$P=C_3$. The centralizer of $P$ in $H$ lies in the roots of unity
of $\mathbf Q_3(\zeta_3)$, namely $\mu_6$. Conjugation embeds
$H/C_H(P)$ in $\operatorname{Aut}(C_3)=C_2$. Consequently
$|H|\leq12$. Its tame quotient has order dividing $8$ and at
most $4$, as asserted. If $P=1$, reduction embeds $H$ in the
cyclic group of order $8$.

The faithful special inertia image gives a good-reduction
extension of that degree by exactly the descent and lattice
argument of Lemma~\ref{lem:B-RD17-degree}. In particular, the
good model remains a one-dimensional formal group, since its
special isocrystal remains of slope $1/2$.

For two groups, their joint wild image embeds in $P_1\times P_2$,
and the joint tame quotient is cyclic. If both wild images are
nontrivial, its order is at most $9\cdot4=36$. If only one is
nontrivial, it is at most $3\cdot8=24$. If neither is nontrivial,
it is at most $8$. Applying the same descent argument to the
joint image proves the last assertion.
\end{proof}

\begin{lemma}[A point bound for supersingular formal groups]
\label{lem:B-RD31-point}
\status{Proved under stated hypotheses}
Let $F$ be a one-dimensional formal group over a complete DVR
of mixed characteristic $(0,3)$, of absolute ramification $a$,
whose special fibre has height at least two. If $F$ has a
rational point of exact order $3^n$, then
\[
                         a\geq3^n.                     \tag{B-RD31.2}
\]
\end{lemma}
\begin{proof}
Normalize the valuation to give a uniformizer value $1$.
Write $[3](T)=3T+\sum_{j\geq2}b_jT^j$. For an invariant
differential $u(T)dT$, with $u(T)$ an integral unit series, the
identity $u([3](T))[3]'(T)=3u(T)$ shows that $j b_j$ is divisible by $3$.
Thus $b_2\in3\mathcal O$. Since the special fibre has height
at least two, $b_3,\ldots,b_8$ lie in the maximal ideal.
For $x$ of positive integral valuation $v$ this gives
\[
 v([3]x)\geq\min\{a+v,\,1+3v,\,9v\}.
                                                        \tag{B-RD31.3}
\]
A nonzero point $y$ of order $3$ satisfies
$v(y)\leq(a-1)/2$. Indeed, if $v(y)>(a-1)/2$, the linear term
$3y$ has strictly smaller valuation than all terms of degree
$2$ through $8$. It also has smaller valuation than all terms
of degree at least $9$: if $a\geq2$, then
$8v(y)>4(a-1)\geq a$, with strict inequality also at $a=2$;
if $a=1$, the integral inequality $v(y)\geq1$ gives it
directly. This would contradict $[3]y=0$.

For a point of exact order $3^n$, write
$v_i=v([3]^i x)$ for $0\leq i<n$. The final point satisfies
$v_{n-1}\leq(a-1)/2$. Each earlier $v_i$ also lies below this
bound: (B-RD31.3) strictly increases every positive valuation,
so a value already above that bound could never reach the
final one. Consequently $a\geq2v_i+1$, and
(B-RD31.3) gives $v_{i+1}\geq3v_i+1$.
Starting from $v_0\geq1$, induction yields
\[
 v_{n-1}\geq\frac{3^n-1}{2}.
\]
Combining the two inequalities gives $a\geq3^n$.
\end{proof}

\begin{proposition}[The mixed formal and filtered case is detected by $27$-torsion]
\label{prop:B-RD31-mixed}
\status{Proved under stated hypotheses}
Suppose $p=3$, $e_K=2$, and $\mu_3\subset K$.
Let $M$ be a height-two potentially good group with supersingular
special fibre. Let $N$ admit a filtration with height-one graded
pieces and with sum of their dimensions equal to $1$.
Then
\[
                         M[27]\not\simeq N[27].
                                                        \tag{B-RD31.4}
\]
\end{proposition}
\begin{proof}
Pass to the completed maximal unramified extension, which
preserves $e_K=2$ and the given torsion comparison. A potentially
good height-one character is then $\chi_3^b\epsilon$, where
$b\in\{0,1\}$ and $\epsilon$ has values in $\{1,-1\}$.
This is the height-one classification used in
Proposition~\ref{prop:B-RD5}. The determinant of $T_3M$ is
$\chi_3\epsilon_0$, with $\epsilon_0$ likewise quadratic;
this follows by taking the determinant in crystalline comparison,
as in Proposition~\ref{prop:B-RD18}.

An asserted isomorphism transports the first height-one
subgroup of $N[27]$ to a cyclic direct summand $C\subset M[27]$
whose character is $\chi_3^b\epsilon$. If $b=0$, twist $M$ by
$\epsilon$ to obtain a rational point of order $27$. If $b=1$,
form the BT quotient $M'=M/C$. The image of $M[27]$ in $M'[27]$
is cyclic of order $27$ and has character
$\epsilon_0\epsilon^{-1}$. Twisting $M'$ by this quadratic
character again gives a rational point of exact order $27$.
The quotient and twist remain potentially good, one-dimensional,
and supersingular. A finite flat closure of the kernel in a good
model justifies the assertion for the quotient; the twist becomes
trivial over a quadratic extension.

By Lemma~\ref{lem:B-RD31-finite}, the resulting group acquires
good formal reduction over an extension of ramification degree
at most $12$. Its absolute ramification is therefore at most
$24$. Its rational point extends to a point of this formal
group: finite flat torsion is proper over the DVR, and its
connected special fibre has only the identity geometric point.
Lemma~\ref{lem:B-RD31-point} would require absolute ramification
at least $27$, a contradiction.
\end{proof}

\begin{theorem}[B-RD31: three possibilities in the smallest odd-prime case]
\label{thm:B-RD31}
\status{Proved; two independent complete proof audits and primary-source checks passed.}
In B.4 one has
\[
                         2\leq\nu_{\min}(3,2,2)\leq4.
                                                        \tag{B-RD31.5}
\]
More precisely, $M[243]$ determines the prime-to-$3$ character
of every height-two group satisfying the manuscript monodromy
condition. Any counterexample to determination by $M[27]$ or
$M[81]$ must involve two groups which themselves have
potentially good supersingular reduction. Thus deciding the
two levels $27$ and $81$ in that restricted case decides
$\nu_{\min}(3,2,2)$.
\end{theorem}
\begin{proof}
The lower bound is Proposition~\ref{prop:B-RD29-p3-level-two}.
For the upper bound, let $M[243]\simeq N[243]$ and work after
completed maximal unramified base change. For a height-two
monodromy filtration, let $d\in\{0,1,2\}$ be the sum of the
dimensions of its good graded models. Its determinant character
on inertia is $\chi_3^d\epsilon$, with $\epsilon$ finite
quadratic. The image of $\chi_3$ modulo $9$ has order $3$:
$\mu_3\subset K$, whereas $e_K=2$ excludes $\mu_9\subset K$.
Equality of determinants modulo $9$ therefore forces equal $d$
and equal $\epsilon$ for $M$ and $N$.

If $d=0$, all good graded models are etale. Their special
inertia character is the Brauer lift of the semisimplification
of the generic representation modulo $3$, evaluated on
prime-to-$3$ elements. It is therefore determined already by
$M[3]$. If $d=2$, Cartier duality reduces to $d=0$; the
cyclotomic twist has trivial finite descent inertia. This
also treats extensions between the graded pieces, since
semisimplification and the character are additive.

Suppose $d=1$. Each group either has a filtration by height-one
pieces, or is itself potentially good and supersingular.
Indeed a potentially ordinary height-two graded piece has a
canonical connected--etale filtration on its good model;
uniqueness and descent of this filtration produce the claimed
height-one filtration over $K$.
For two groups with height-one filtrations, the graded generic
characters are $\chi_3^{b_i}\epsilon_i$. Since
$\chi_3\equiv1\pmod3$ on inertia, their finite characters
$\epsilon_i$ and hence their sum are determined by the
semisimplified $3$-torsion. A pair of different kinds is excluded
already at level $27$ by Proposition~\ref{prop:B-RD31-mixed}.

It remains to compare two potentially supersingular groups.
Lemma~\ref{lem:B-RD31-finite} gives a common good formal
extension $L$ with $e_L\leq72$. Since
\[
                         72\leq2\cdot3^4,
\]
Bondarko's Theorem~E(1), with $u=4$ and $\ell=1$, extends the
generic $3^5$-torsion comparison to an isomorphism between the
special $3$-torsion models, equivariantly for descent. The
special Dieudonne modules modulo $3$ consequently have the
same prime-to-$3$ Brauer character. A finite-order matrix of
order prime to $3$ on a free Witt module has eigenvalues which
are Teichmuller lifts of its residual eigenvalues. Thus this is
the required equality of characteristic-zero inertia
characters. The same reductions, without the final application
of Theorem~E, give the assertion about possible counterexamples
at levels $27$ and $81$.
\end{proof}


\begin{remark}[Exact auxiliary certificate and remaining gate]\label{auto:B:a26975327b26}
\status{Proved under stated hypotheses}

The program \path{code/verify_height_two_small_bounds.py} checks the
quaternion matrix identities, all $49$ joint-image numerical cases,
$16000$ formal-valuation inequalities and the $36$ determinant
comparisons used above. It passed in $0.05041721599991433$ seconds;
its JSON output is retained. These finite checks certify the
explicit algebra and inequalities, not the existence or
nonexistence of a counterexample at either remaining level.
At this intermediate stage neither $27$ nor $81$ was decided for arbitrary pairs
of potentially supersingular groups over fields of absolute
ramification $2$.
\end{remark}

\par\smallskip\noindent\textbf{Subsequent exact determination.}
The remaining small-height alternatives isolated here are settled in
Theorem~\ref{thm:B-RD42}: the exact manuscript index is
$\nu_{\min}(3,2,2)=2$. The intermediate bounds and conditional
reductions above are retained as steps in that proof.


\subsection{The remaining $81$-torsion comparison has only two inertia pairs}

Theorem~\ref{thm:B-RD34} below excludes both of the level-$81$
pairs isolated here and strengthens the level-$27$ restriction to
joint inertia order $18$ or $36$. The present proposition records
the group-theoretic reduction used in that proof.

\begin{lemma}[Norm and prime-to-$3$ characters in the quaternion algebra]
\label{lem:B-RD32-norm}
\status{Proved under stated hypotheses}
For the finite subgroups $H$ in Lemma~\ref{lem:B-RD31-finite},
the reduced norm is trivial whenever $3\mid |H|$.
If it is nontrivial, $H$ is cyclic of order $8$.
If it is trivial, the possible orders are
\[
                           1,2,3,4,6,12.
                                                        \tag{B-RD32.1}
\]
In the norm-one case the prime-to-$3$ character depends only on
the tame quotient order $t\in\{1,2,4\}$, and on a prime-to-$3$
element whose image in that quotient has order $r$ it is
\[
       r=1:\ 2,\qquad r=2:\ -2,\qquad r=4:\ 0.
                                                        \tag{B-RD32.2}
\]
Thus two norm-one special inertia images with the same tame
quotient order have identical prime-to-$3$ characters as
functions on inertia.
\end{lemma}
\begin{proof}
Retain the explicit quaternion algebra (B-RD31.1). A noncentral
element has characteristic polynomial of degree two, whose
constant term is its reduced norm. There is only one element
of order two, namely $-1$, because the algebra is a division
algebra and $(x-1)(x+1)=0$ forces $x=1$ or $-1$.

If the wild subgroup is nontrivial, it is $P=C_3$. Its
centralizer in $H$ is contained in
$\mu_6\subset\mathbf Q_3(\zeta_3)$, on which the reduced norm
is $1$. If $H$ is not equal to this centralizer, its quotient
by the centralizer has order two. The tame quotient of order
two cannot give a nonabelian group of order six, since its
involutions would be noncentral. Therefore the nonabelian
case has order twelve. Its Sylow $2$-subgroup injects under
reduction into $\mathbf F_9^\times$, so it is cyclic of order
four. A generator $j$ satisfies $j^2=-1$ and has reduced
characteristic polynomial $X^2+1$, hence reduced norm $1$.
The centralizer and $j$ generate $H$. This proves norm one
for every wild $H$.

If $H$ is tame, reduction embeds it in $\mathbf F_9^\times$,
so it is cyclic of order dividing eight. A generator lies in
a quadratic commutative subfield. For an element of order
eight this is the unramified quadratic field, and its reduced
norm is $\zeta_8^{1+3}=\zeta_8^4=-1$. For elements of order
dividing four the norm is $1$. This proves (B-RD32.1) and
the assertion about nontrivial norm.

In a norm-one subgroup, a prime-to-$3$ element has order
$1,2$, or $4$. Its trace is respectively $2,-2$, or $0$,
the last two facts following from $-1$ and $X^2+1$.
Its cyclic subgroup maps injectively to the tame quotient,
so its order equals the order of its image there.
Over the completed maximal unramified base, the absolute
tame inertia is procyclic and has a unique quotient of each
of the orders $1,2,4$. The order of the image of an element
in that quotient is independent of the choice of a generator.
Formula (B-RD32.2) therefore proves equality as functions on
inertia, even when the two wild good-reduction fields differ.
\end{proof}

\begin{proposition}[B-RD32: exact restriction of the two undecided tests]
\label{prop:B-RD32}
\status{Proved; parent and independent second proof audits passed.}
Keep $p=3$, $e_K=2$, $\mu_3\subset K$, and height two.
For a pair with isomorphic $27$-torsion and different
prime-to-$3$ characters, both groups must be potentially
supersingular, both finite special inertia images must have
trivial reduced norm, their tame quotient orders must differ,
and their joint faithful good-reduction image must have order
strictly greater than $9$.

For a pair with isomorphic $81$-torsion and different
prime-to-$3$ characters, after interchanging the pair their
individual faithful special inertia orders must be
\[
                             (12,3)\quad\hbox{or}\quad(12,6),
                                                        \tag{B-RD32.3}
\]
and their joint faithful special inertia image must have
order exactly $36$. In particular its wild subgroup has
order $9$ and the two wild order-three quotients are
independent.

Consequently $\nu_{\min}(3,2,2)\leq3$ if and only if
no $81$-torsion comparison of the two specified kinds has
different prime-to-$3$ characters. This is a restricted
remaining test, not a claim that either kind exists.
\end{proposition}
\begin{proof}
Theorem~\ref{thm:B-RD31} reduces either test to two whole
potentially supersingular groups. Its determinant argument
shows that their finite determinant characters agree.
Under crystalline comparison this finite determinant is
the reduced norm of the special quaternionic inertia
action: the comparison duality or period Tate twist does
not change finite descent inertia.

If this determinant is nontrivial, Lemma~\ref{lem:B-RD32-norm}
makes both special images cyclic of order eight and tame.
Their common good extension then has ramification at most
eight, hence absolute ramification at most sixteen.
Bondarko's Theorem~E(1), with $u=2$ and $\ell=1$, applies
since $16\leq2\cdot3^2$. Already an isomorphism on
$27$-torsion makes their special $3$-torsion equivariantly
isomorphic and their prime-to-$3$ characters equal.
Thus any counterexample has norm one on both sides.

Their tame quotient orders must differ by
Lemma~\ref{lem:B-RD32-norm}. Write $m$ for the order of
the joint faithful image. Its wild subgroup has order
$1,3$, or $9$, and its tame quotient has order $1,2$, or
$4$. In particular
\[
                m\in\{1,2,3,4,6,9,12,18,36\}.
                                                        \tag{B-RD32.4}
\]
For $m\leq9$, the common absolute ramification is at most
$18=2\cdot3^2$, so the same $u=2$ argument settles the
$27$-torsion comparison. This proves the first assertion.

At level $81$, take $u=3$ and $\ell=1$. The comparison is
settled whenever $2m\leq2\cdot3^3$, or $m\leq27$.
The only larger value in (B-RD32.4) is $36$. This forces
both wild images to have order three and the joint tame
quotient to have order four. One individual image must
therefore have order twelve. If both had that order,
their tame quotient orders and hence their prime-to-$3$
characters would agree. The other image instead has
tame quotient order one or two, and hence total order
three or six. Joint wild order nine is exactly the
asserted independence. Every failure of the unrestricted
$81$ test has now been reduced to (B-RD32.3); the converse
implication about the least index is immediate because
these specified pairs are themselves allowed in B.4.
\end{proof}

\begin{remark}[Exact auxiliary certificate]\label{auto:B:72e408a14989}
\status{Proved under stated hypotheses}

The retained program
\path{code/verify_height_two_character_gate.py}
checks the eight cyclotomic reduced norms and traces, all twenty-seven
numerical joint-image possibilities, and the two remaining level-$81$
pairs. Its exact run passed in $0.050076540999725694$ seconds.
The numerical level-$27$ list is an upper list of possibilities; no
existence claim is made for a listed group-theoretic combination.
\end{remark}

\par\smallskip\noindent\textbf{Subsequent exact determination.}
The remaining small-height alternatives isolated here are settled in
Theorem~\ref{thm:B-RD42}: the exact manuscript index is
$\nu_{\min}(3,2,2)=2$. The intermediate bounds and conditional
reductions above are retained as steps in that proof.


\subsection{A hypothetical comparison has a prescribed Frobenius graph}

\begin{proposition}[B-RD33: the graph at every smaller level]
\label{prop:B-RD33}
\status{Proved. Independent root and obstruction/crystal proof and primary-source
audits passed. The exact graph is a necessary condition and does not by itself
improve the numerical upper bound.}
Keep the completed maximal unramified base used in
Proposition~\ref{prop:B-RD32}. Let $M,N$ be height-two,
one-dimensional, potentially supersingular $3$-divisible groups
with trivial finite determinant characters and different tame
special-inertia quotient orders. Let $L/K$ be a common finite
Galois good-reduction extension, with good models
$\mathcal M,\mathcal N$ over $R=\mathcal O_L$ and algebraically
closed residue field $\bar k$. Write $M_0,N_0$ for their special
fibres. Suppose
\[
             f:M[3^n]\xrightarrow{\sim}N[3^n]
                                                        \tag{B-RD33.1}
\]
is an isomorphism over $K$.
For $1\le r\le n$, let $H_r$ be the scheme-theoretic closure of
the graph of $f|_{M[3^r]}$ in
$\mathcal M[3^r]\times_R\mathcal N[3^r]$.
Then, as closed subgroups of the indicated product,
\[
 (H_r)_{\bar k}
       =\ker(F_{M_0}^{r})\times\ker(F_{N_0}^{r}),
                                                        \tag{B-RD33.2}
\]
where $F_{M_0}^{r}:M_0\to M_0^{(3^r)}$ is the iterated
relative Frobenius, and likewise for $N_0$.
In particular a remaining comparison at level $27$ has graph
special fibre $\ker(F_{M_0}^{3})\times\ker(F_{N_0}^{3})$;
at level $81$ it has special fibre $M_0[9]\times N_0[9]$.

Put $H_0=0$ and $Q_r=H_r/H_{r-1}$. Every $Q_r$ is finite flat
of rank nine, its special fibre is $\alpha_3\times\alpha_3$,
and its generic fibre is the common representation $M[3]$.
Multiplication by $3$ induces maps
\[
                    Q_{r+1}\longrightarrow Q_r
                      \qquad(1\le r<n)
                                                        \tag{B-RD33.3}
\]
that are isomorphisms on generic fibres and zero on special
fibres. These statements are necessary conditions. No lift
of this prescribed chain with the required descent actions
is asserted here.
\end{proposition}
\begin{proof}
\par\smallskip\noindent\textbf{Perfect pairings on the full good models.}
The determinant argument of Theorem~\ref{thm:B-RD31}, with
trivial finite determinant and after the completed maximal
unramified base change, gives
\[
             \det T_3(M)=\mathbf Z_3(1),\qquad
             \det T_3(N)=\mathbf Z_3(1).
\]
Here equality means an isomorphism of rank-one Galois lattices.
Indeed the weight-zero quotient of the determinant by the
cyclotomic character has trivial finite inertia and is
unramified, hence trivial on this base. A choice of each
isomorphism supplies a perfect alternating pairing of Tate
lattices into $\mathbf Z_3(1)$, and therefore an isomorphism
of the generic $3$-divisible group with its Cartier dual.
Tate's integral full-faithfulness theorem extends this map
and its inverse over $R$; their composites remain identity
by generic-fibre faithfulness. Thus both good models have
perfect alternating pairings at every torsion level.
Alternation and Galois-descent compatibility also extend
uniquely from the generic fibre. The exact source is
\cite[Theorem~4 and Corollaries~1--2, printed p.~181]
{Tate1967PDivisible}; this argument uses full $3$-divisible
groups for extension, not finite-level full faithfulness.

Since a linear isomorphism of free rank-two
$\mathbf Z/3^n$-modules multiplies an alternating pairing by
its determinant, there is a unit
$\lambda\in(\mathbf Z/3^n)^\times$ such that
\[
            e_N(fx,fy)=e_M(x,y)^\lambda.
\]
Choose any lift of $\lambda$ to $\mathbf Z_3^\times$.
On each product $P_r=\mathcal M[3^r]\times\mathcal N[3^r]$
use the perfect alternating product pairing
$e_M\,e_N^{-\lambda^{-1}}$. The graph is maximal isotropic
on the generic fibre.

\par\smallskip\noindent\textbf{Closure and orthogonality commute with reduction.}
Each graph closure $H_r$ is finite flat over the DVR:
its coordinate algebra is the image of that of $P_r$ in
the generic graph algebra, and is finite and torsion-free.
It is a closed subgroup, of rank $3^{2r}$, stable under
$\operatorname{Gal}(L/K)$. The restricted pairing on
$H_r\times_RH_r$ is trivial, because this product is flat
and the pairing is trivial generically.

Finite-flat Cartier duality is exact, so the orthogonal
$H_r^\perp$ is finite flat of rank $3^{2r}$ and its formation
commutes with base change. There is a closed immersion
$H_r\subset H_r^\perp$. The coordinate-algebra surjection
is between free $R$-modules of equal rank; it is therefore
an isomorphism. Thus $H_r=H_r^\perp$, also on the special
fibre. Exact finite-flat duality and the rank formula are
recalled in \cite[\S1.2--1.3, printed pp.~159--160]
{Tate1967PDivisible}.

\par\smallskip\noindent\textbf{Tame idempotents split the graph.}
The norm-one possibilities in Lemma~\ref{lem:B-RD32-norm}
have tame orders $1,2,4$. Choose an element $\tau$ in a
Sylow $2$-subgroup of the joint finite descent image whose
image generates its tame quotient; Schur--Zassenhaus, or
the cyclic tame quotient and the odd wild kernel, supplies
such a lift. Its action on the special fibre is
$\bar k$-linear.

If one tame order is four and the other is one or two,
then $\tau^2$ acts as $[-1]$ on the first special group
and as $[1]$ on the second. Indeed the unique quaternionic
involution is $-1$, and rational Dieudonne faithfulness
implies the same equality of integral group endomorphisms.
Consequently $(1-\tau^2)/2$ is the projection onto the
first factor. If the two orders are one and two, use
$(1-\tau)/2$ instead. These are well-defined endomorphisms
of the $3$-primary group schemes. They preserve
$(H_r)_{\bar k}$ because that subgroup is stable under
$\tau$ and under scalar multiplication. Hence
\[
                 (H_r)_{\bar k}=J_M\times J_N
                                                        \tag{B-RD33.4}
\]
for closed subgroups $J_M\subset M_0[3^r]$ and
$J_N\subset N_0[3^r]$.

A one-dimensional connected formal group over $\bar k$
has a unique finite closed subgroup of any specified
$p$-power rank: its completed coordinate ring is
$\bar k[[T]]$, and the only ideal of colength $3^a$ is
$(T^{3^a})$. When it is a subgroup this is the kernel of
relative $3^a$-Frobenius. Thus
$J_M=\ker F_{M_0}^a$ and $J_N=\ker F_{N_0}^b$ for some
$a,b$, with $a+b=2r$.

The orthogonal of $\ker F_{M_0}^a$ in the perfectly
self-dual group $M_0[3^r]$ has rank $3^{2r-a}$.
Uniqueness makes it $\ker F_{M_0}^{2r-a}$, and similarly
for $N_0$. Since the product in (B-RD33.4) is its own
orthogonal, $a=r$ and $b=r$. This proves (B-RD33.2).
For $r=4$, uniqueness of the rank-$3^4$ subgroup identifies
$\ker F^4$ with the $9$-torsion of a height-two formal group;
no choice identifying a Frobenius twist with the group is
needed.

\par\smallskip\noindent\textbf{The successive finite-flat models.}
The graph closure in $P_r$ is also its closure in $P_n$,
since $P_r$ is a closed subgroup of $P_n$. Hence the $H_r$
form nested finite-flat closed subgroups, and their finite
flat quotients exist. Base change of the quotient sequence
is exact. Each special quotient
$\ker F^r/\ker F^{r-1}$ of a supersingular one-dimensional
group is $\alpha_3$: it is a rank-three subgroup of the
appropriate height-two Frobenius twist, whose only
rank-three subgroup is its Frobenius kernel and is local
with local Cartier dual. Thus $(Q_r)_{\bar k}=\alpha_3^2$.
The generic quotient is the usual quotient
$M[3^r]/M[3^{r-1}]\simeq M[3]$.

Multiplication by $3$ maps $H_{r+1}$ to $H_r$, by the
same generic-fibre closure argument, and maps $H_r$ to
$H_{r-1}$. It gives (B-RD33.3), generically an isomorphism.
On a height-two special formal group, multiplication by
$3$ maps $\ker F^{r+1}$ into $\ker F^{r-1}$: its kernel
has rank $3^2$ and its image in the displayed subgroup
has rank at most $3^{r-1}$, so uniqueness gives the
inclusion (with the case $r=1$ giving zero). The induced
map on the rank-three quotients is therefore zero.
This proves the last assertion.
\end{proof}

\begin{remark}[What the graph gate does not decide]\label{auto:B:18f97bffbcfd}
\status{Proved under stated hypotheses}

Bondarko's kernel theorem
\cite[Theorem~5, printed pp.~681--682]{Bondarko2006Generic}
says that the kernel of the special fibre of a generically
injective finite-flat map is killed by $F^u$ when
$e_L<2\cdot3^u$ (and also at equality for connected source).
Applied to the graph projections, (B-RD33.2) forces
$n\le u$. This recovers the proved numerical exclusions.
The additional content here is the exact special graph
and the entire chain (B-RD33.3). The refined logarithm argument
of Theorem~\ref{thm:B-RD34} subsequently excludes the
level-$81$ case. At level $27$, existence of a compatible
chain with joint inertia order $18$ or $36$ is not excluded by this proposition alone.
The remaining cases are excluded in Theorem~\ref{thm:B-RD42}.
\end{remark}


\subsection{A supersingular logarithm estimate excludes the remaining level $81$}

\begin{lemma}[A coefficientwise estimate for a product-type special fibre]
\label{lem:B-RD34-log}
\status{Proved under stated hypotheses}
Let $R$ be a complete discrete valuation ring of mixed characteristic
$(0,3)$, with perfect residue field, uniformizer $\pi$, and normalized
valuation $v(3)=a$. Let $G$ be a commutative formal group of finite height
and dimension $d$ over $R$. Suppose the coordinates of $[3]_{G_k}$ have
no monomial of total degree less than nine. Define
\[
 q_0=1,\qquad q_{i+1}=\min\{a+q_i,\,1+3q_i,\,9q_i\}.
                                                        \tag{B-RD34.1}
\]
If $q_u>a/2$, then for every $h\in tR[[t]]^d$ the normalized logarithm
satisfies
\[
       3^u\log_G(\pi h)\in\pi^{q_u}R[[t]]^d.
                                                        \tag{B-RD34.2}
\]
The hypothesis holds for a special fibre isomorphic to a product of
one-dimensional height-two formal groups, and is invariant under an
arbitrary integral change of formal coordinates.
\end{lemma}
\begin{proof}
Write the matrix of normalized invariant differentials as $U(X)$;
its entries are in $R[[X]]$ and $U(0)=1$. Multiplication by three gives
the exact differential identity
\[
     U([3]_G X)\,J_{[3]_G}(X)=3U(X).
                                                        \tag{B-RD34.3}
\]
Since $U([3]_GX)$ is invertible over $R[[X]]$, every entry of the
Jacobian $J_{[3]_G}$ is divisible by $3$. If a monomial $X^\alpha$ has
an exponent not divisible by three, its coefficient in $[3]_G$ is
therefore divisible by three. In particular this holds for every
monomial of total degree two. The linear part is $3X$. The special-fibre
hypothesis makes every coefficient of total degree three through eight
divisible by $\pi$, and all higher coefficients are integral. Thus an
augmentation vector whose coefficients have valuation at least $q\geq1$
is sent by $[3]_G$ to one whose coefficients have valuation at least
\[
                        \min\{a+q,1+3q,9q\}.
\]
Iteration starting with $\pi h$ proves the bound $q_i$ for
$[3]^i_G(\pi h)$.

For completeness this also controls the multivariable logarithm.
Write a coefficient of $X^\alpha$ in any coordinate of $\log_G$ as
$c_\alpha$, and put $m=|\alpha|\geq1$. Since $d\log_G=U(X)dX$,
\[
 v(c_\alpha)\geq
       -a\min_{j:\alpha_j>0}v_3(\alpha_j)
       \geq-a\lfloor\log_3m\rfloor.
                                                        \tag{B-RD34.4}
\]
For every integer $r\geq0$, $3^r-1\geq2r$. Consequently
$\lfloor\log_3m\rfloor\leq(m-1)/2$. If a vector $z$ has coefficient
valuations at least $q>a/2$, every coefficient of
$c_\alpha z^\alpha$ has valuation at least
$mq-a\lfloor\log_3m\rfloor\geq q$. Formal substitution is legitimate
because $z$ is an augmentation vector; each coefficient involves only
finitely many terms. Hence $\log_G(z)\in\pi^qR[[t]]^d$.
Apply this to $z=[3]^u_G(\pi h)$ and use
$\log_G([3]^u_G(\pi h))=3^u\log_G(\pi h)$.

For a one-dimensional height-two group in characteristic three,
multiplication by three starts in degree nine. This persists under
products. An integral coordinate change and its inverse have zero
constant terms and invertible linear parts. Substitution in either
direction preserves the ideal generated by monomials of total degree
at least nine, so the hypothesis is coordinate independent. No
coordinate choice on the mixed-characteristic quotient is assumed.
\end{proof}

\begin{lemma}[A refined kernel estimate]
\label{lem:B-RD34-kernel}
\status{Proved under stated hypotheses}
Keep $R,a$ as above. Let $f:F\to G$ be a homomorphism of finite-height
formal groups having no nonzero finite flat subgroup scheme in its
kernel, and with finite special-fibre kernel. This is absence of a
flat kernel; it does not require finite-level closed immersions and
does not preclude a nonzero kernel after reduction. Suppose $G_k$
satisfies the hypothesis of
Lemma~\ref{lem:B-RD34-log}. If $q_u>a/2$, then
\[
                         F^u\ker(f_k)=0,
                                                        \tag{B-RD34.5}
\]
where $F^u$ denotes iterated relative Frobenius. In particular one can
take $u=2$ when $a\leq24$, and $u=3$ when $a=72$.
\end{lemma}
\begin{proof}
We give the change to the original Cartier-module proof explicitly.
Use Bondarko's invariant module $D_G\subset L[[\Delta]]^d$, where
$L=\operatorname{Frac}(R)$, and put
\[
 \mathcal R=R[[\Delta]][1/3],\qquad
 J=\pi^{\lceil a/2\rceil}R.
\]
Here $\mathcal R$ is precisely the ring of series with coefficient
valuations bounded below. Let $G_\pi$ be the formal group law
$\pi^{-1}G(\pi X,\pi Y)$, and let $\mathcal C_G$ denote Bondarko's
completed module $M_G$. We use the exact identities
\[
 \mathcal C_G=\Omega\,\pi D_{G_\pi},\qquad
 J[[\Delta]]^d\subset D_G,\qquad
 E(G_k)=D_G/(D_G\cap\mathcal R^d).
                                                        \tag{B-RD34.6}
\]
They are respectively Proposition~10(4), Proposition~9, and Lemma~3
of \cite[printed pp.~674--675]{Bondarko2006Generic}. The canonical
embedding $D_G\to\mathcal C_G\subset L\{\{\Delta\}\}^d$ agrees
with the literal series on every bounded coefficient projection;
this is Proposition~11(1)--(2), printed p.~675.

By the definition of $D_{G_\pi}$, an element of $\pi D_{G_\pi}$,
viewed as a $3$-typical series in $t$, has the form
$\log_G(\pi h)$ for $h\in tR[[t]]^d$.
Lemma~\ref{lem:B-RD34-log} therefore gives
\[
             3^u\mathcal C_G
                  \subset\pi^{q_u}R\{\{\Delta\}\}^d
                  \subset J\{\{\Delta\}\}^d.
                                                        \tag{B-RD34.7}
\]
Indeed $\Omega$ consists of two-sided series with integral Witt
coefficients, tending to zero towards negative powers. Its action
preserves every coefficient valuation ideal; the formula for this
action is displayed immediately before Proposition~10 in the cited
source. Thus the estimate on $\pi D_{G_\pi}$ passes to its
$\Omega$-span. This replaces the general logarithm bound in
Proposition~11(5).

Here is the needed replacement for Proposition~12(1)(ii). If
$z\in D_G$, $A$ is an integral matrix, and $Az\in\mathcal R^b$, then
\[
                     \mathbf f^uAz\in J[[\Delta]]^b.
                                                        \tag{B-RD34.8}
\]
To see this, (B-RD34.7) and the bounded-projection property give
$3^uAz\in J[[\Delta]]^b$. The Cartier Frobenius satisfies
\[
        \mathbf fz=3(z-z(0))\Delta^{-1},\qquad z(0)\in R^d.
\]
Each $\mathbf f^jz$ again lies in $D_G$, so its constant term is
integral. Induction gives
\[
   \mathbf f^uAz\equiv3^uAz\Delta^{-u}
                        \pmod{3R((\Delta))^b}.
\]
Since $3R\subset J$, the right side and the error lie in
$J((\Delta))^b$. The left side has no negative powers, proving
(B-RD34.8). This is equation~(8) and the proof of Proposition~12,
printed p.~676, with the improved bound inserted.

We now spell out the final step of Proposition~16, printed p.~681.
Let $B=df(0)$, $I'=B(L^{\dim F})$, and $I=I'\cap R^d$.
Choose an integral direct-sum complement $R^d=I\oplus C$ and
the associated projection $\operatorname{pr}:R^d\to C$.
Propositions~6 and~8 identify
\[
                     B D_F=D_G\cap I'[[\Delta]].
                                                        \tag{B-RD34.9}
\]
Indeed the formal group belonging to the right side factors $f$;
the resulting map from $F$ is an isogeny, and generic injectivity
forces its finite flat kernel to vanish, so it is an isomorphism.
The same reduction identity used in Proposition~16 then embeds
$E(F_k)$ in $E(G_k)$.

Take $z\in D_G$ whose reduction $\bar z$ belongs to
$\mathbf Q_3E(F_k)\cap E(G_k)$. By (B-RD34.6), after multiplying
by a power of three, $z$ differs from an element of $BD_F$ by a
bounded-coefficient series. Hence
$\operatorname{pr}(z)\in C\otimes_R\mathcal R$.
Formula (B-RD34.8) shows
\[
 c=\mathbf f^u\operatorname{pr}(z)\in J C[[\Delta]]
                                                \subset D_G.
\]
This element is bounded, so reduces to zero in $E(G_k)$. Moreover
$\mathbf f^uz-c$ lies in the right side of (B-RD34.9). Therefore
$\mathbf f^u\bar z\in E(F_k)$. The saturation quotient
\[
 (\mathbf Q_3E(F_k)\cap E(G_k))/E(F_k)
\]
is the Dieudonne module of $\ker(f_k)$ by Theorem~2(8), equation~(6),
printed p.~668 of the same source. The operator $\mathbf f$
corresponds to relative Frobenius, as stated immediately after that
equation. Full faithfulness proves (B-RD34.5).

For the numerical assertions, if $a\leq2$ then $q_2=2a+1$;
if $3\leq a\leq9$ then $q_2=a+4$; and if $9\leq a\leq24$
then $q_2=13$. Each is greater than $a/2$. At $a=72$ the successive
values are $1,4,13,40$, and $40>36$.
\end{proof}

\begin{theorem}[B-RD34: the height-two least index has only two possibilities]
\label{thm:B-RD34}
\status{Proved. Independent parent and obstruction/crystal proof and original-source
audits passed; exact arithmetic checks also passed.}
For the character question B.4 with $p=3$, absolute ramification
$e_K=2$, $\mu_3\subset K$, and height two,
\[
                         2\leq\nu_{\min}(3,2,2)\leq3.
                                                        \tag{B-RD34.10}
\]
Thus an isomorphism of $81$-torsion always determines the requested
prime-to-$3$ character. A failure at level $27$, if one exists,
must compare two whole potentially supersingular groups with
norm-one finite special inertia and different tame quotient orders;
the order of their joint faithful good-reduction image is $18$ or
$36$.
\end{theorem}
\begin{proof}
The lower bound is the explicit level-nine comparison in
Proposition~\ref{prop:B-RD29-p3-level-two}. The reductions in
Theorem~\ref{thm:B-RD31} and Proposition~\ref{prop:B-RD32} show
that only two whole potentially supersingular groups need be
considered. At level $81$, their joint image would have order
$36$, so their common good-reduction DVR $R$ has $a=72$.
After completed unramified base change, form the graph closure
$H_n\subset P=\mathcal M\times\mathcal N$ at level $3^n$.
Proposition~\ref{prop:B-RD33} gives
\[
       (H_n)_k=\ker F_{M_0}^n\times\ker F_{N_0}^n.
                                                        \tag{B-RD34.11}
\]
The quotient $Q=P/H_n$ is a $3$-divisible group over $R$.
Its special fibre is the product of the corresponding Frobenius
twists of $M_0,N_0$, because quotient formation commutes with base
change. In particular it is a product of one-dimensional
height-two formal groups. Use the homomorphism of formal groups
\[
 g:P\longrightarrow Q\times\mathcal M,
                 \qquad z\longmapsto(qz,\operatorname{pr}_{\mathcal M}z).
                                                        \tag{B-RD34.12}
\]
The generic graph of an isomorphism meets $0\times N$ trivially,
so $g$ has no nonzero finite flat subgroup scheme in its kernel.
By (B-RD34.11), its special kernel is
$0\times\ker F_{N_0}^{n}$. Its target special fibre is a product
of three height-two groups, so Lemma~\ref{lem:B-RD34-kernel}
applies. For $n=4,a=72$ it says that $F^3$ kills
$\ker F_{N_0}^4$. This is impossible: on a one-dimensional formal
group that Frobenius image has rank three. This excludes every
remaining level-$81$ case and proves the upper bound.

For a hypothetical level-$27$ comparison, the same construction
with $n=3$ and $a\leq24$ would force $F^2\ker F_{N_0}^3=0$,
again contradicting its rank-three image. The norm-one joint-image
list in (B-RD32.4) therefore leaves only $18$ and $36$. No
existence claim for a remaining pair, and no decision between
least indices two and three, is made by this intermediate bound.
The final decision $\nu_{\min}(3,2,2)=2$ is proved in
Theorem~\ref{thm:B-RD42}.
\end{proof}

\begin{remark}[Exact arithmetic certificate]\label{auto:B:2b832ed8e5d4}
\status{Proved under stated hypotheses}

The retained program \path{code/verify_supersingular_log_bound.py}
checks $282$ multivariable monomials, $18\,432$ coefficient-valuation
cases, and $40\,000$ logarithm-degree cases. It also verifies every
ramification value $1\leq a\leq24$, the sequence $(1,4,13,40)$ at
$a=72$, the remaining joint orders $18,36$, and both rank-three
Frobenius images. Its exact author run passed in
$0.05011515400110511$ seconds. These are arithmetic checks of the
proof; no existence of a remaining level-$27$ comparison is asserted.
\end{remark}


\subsection{The remaining character comparison can be realized by elliptic curves}

\begin{lemma}[Integral conjugacy of the finite norm-one action]
\label{lem:B-RD35-conjugacy}
\status{Proved under stated hypotheses}
Let $k$ be algebraically closed of characteristic three, let
$E_0/k$ be $y^2=x^3-x$, and write $G_0=E_0[3^\infty]$.
Every finite reduced-norm-one subgroup of $\operatorname{Aut}(G_0)$
is conjugate by an automorphism of $G_0$ into
$\operatorname{Aut}(E_0)$. In particular this is integral conjugacy,
not only conjugacy by a quasi-isogeny.
\end{lemma}
\begin{proof}
Use the quaternion division algebra $D$ and its valuation ring
$\mathcal O_D=\End(G_0)$ from (B-RD31.1), normalizing the division
valuation by $v_D(3)=2$. Write
\[
 \zeta:(x,y)\longmapsto(x+1,y),\qquad
 j:(x,y)\longmapsto(-x,iy),\quad i^2=-1.
\]
These automorphisms have
\[
 \zeta^3=1,\quad j^2=-1,\quad
 j\zeta j^{-1}=\zeta^{-1},\qquad
 H_0=\langle\zeta,j\rangle\quad\hbox{of order twelve}.
                                                        \tag{B-RD35.1}
\]
Here $-1$ denotes the elliptic involution, not the identity.
Every origin-preserving automorphism of this curve has the form
$(x,y)\mapsto(u^2x+r,u^3y)$, with $u^4=1$ and $r^3-r=0$:
the pole orders of $x,y$ give the usual Weierstrass change of
variables, and comparison of the $y$-terms eliminates additions
to $y$ since $2$ is invertible. Substitution then gives the two
displayed conditions. Thus $H_0=\operatorname{Aut}(E_0)$.
Put $E=\mathbf Q_3(\zeta)\subset D$. We have
\[
 v_D(\zeta-1)=1,\qquad
 (\zeta-1)j(\zeta-1)^{-1}=-\zeta j.
                                                        \tag{B-RD35.2}
\]
The first equality follows from
$(\zeta-1)(\zeta^{-1}-1)=3$. Consequently $\zeta-1$
normalizes $H_0$ and permits adjustment of a conjugator's valuation
by any integer.

By Lemma~\ref{lem:B-RD32-norm}, a finite norm-one group $H$ is
cyclic of order $1,2,3,4$, or $6$, or is the nonabelian group of
order twelve described there. We give conjugators rather than
appealing only to rational Skolem--Noether.

If $z\in H$ has order three, the element $d=z+\zeta+1$
satisfies $dz=\zeta d$. If $d=0$, then $z=\zeta^{-1}$,
and $j$ instead conjugates $z$ to $\zeta$. Otherwise $d$ is
invertible because $D$ is a division algebra. Left multiplication
by a power of $\zeta-1$, which commutes with $\zeta$, makes $d$
a valuation-zero element. It is then an element of
$\mathcal O_D^\times=\operatorname{Aut}(G_0)$.
This handles the cyclic groups of order three and six; the unique
involution is $-1$.

For a generator $t$ of a cyclic group of order four, $t^2=-1$.
Unless $t=-j$, the nonzero element $d=t+j$ satisfies $dt=jd$.
Multiplication on the left by a suitable power of $3$ makes its
valuation either zero or one. In the latter case multiply once by
$q=\zeta-\zeta^{-1}$, which has valuation one and conjugates
$j$ to $-j$, and then divide by $3$. The resulting conjugator is
a unit and sends $\langle t\rangle$ onto $\langle j\rangle$.
The case $t=-j$ already has the desired subgroup.

Finally suppose $|H|=12$. First use the integral conjugacy just
proved to put its order-three subgroup inside $\langle\zeta\rangle$.
Its order-four generator $t$ conjugates $\zeta$ to $\zeta^{-1}$,
so $t=bj$ for some $b\in E^\times$. The equality $t^2=-1$ gives
$b\bar b=1$. If $b\ne-1$, take $c=1+b$; then $c/\bar c=b$.
If $b=-1$, take $c=\zeta-\zeta^{-1}$ instead, with the same
identity. Thus $c^{-1}tc=j$, while conjugation by $c$ fixes
$\zeta$. Multiplying this rational conjugator on the left by a
power of $\zeta-1$ makes it a unit without changing the target
group $H_0$, by (B-RD35.2). This proves integral conjugacy in
every case.
\end{proof}

\begin{proposition}[B-RD35: elliptic realization with its exact Tate lattice]
\label{prop:B-RD35}
\status{Proved. Independent parent and obstruction/crystal proof, original-source,
and exact-code audits passed.}
Let $K$ be a complete discretely valued field of mixed characteristic
$(0,3)$ with algebraically closed residue field. Let $M/K$ be a
height-two $3$-divisible group that has good reduction to a
one-dimensional supersingular formal group over a finite extension.
Assume its finite special-inertia action has reduced norm one.
Then there is an elliptic curve $E/K$ and an isomorphism of
$3$-divisible groups
\[
                             E[3^\infty]\simeq M.
                                                        \tag{B-RD35.3}
\]
Thus the comparison is integral on the full Tate lattice. It is not
merely an equality of rational representations or of finite inertia
characters.

In particular, after the completed unramified base change already
used in B-RD31--B-RD34, every remaining level-$27$ pair in the
$p=3,e_K=2,h=2$ problem is a pair of elliptic curves with
isomorphic $27$-torsion. The two curves have different tame special
inertia orders, norm-one special inertia, and joint faithful image
of order $18$ or $36$.
\end{proposition}
\begin{proof}
Choose a finite Galois good-reduction extension $L/K$, with
$R=\mathcal O_L$ and group $\Gamma$. Its residue extension is
trivial. Let $\mathcal G/R$ be the good model of $M_L$.
Tate's integral full-faithfulness theorem extends the generic
semilinear descent isomorphisms, and their inverses, to
$\mathcal G$; uniqueness preserves the cocycle relations
\cite[Theorem~4 and Corollaries~1--2, p.~181]{Tate1967PDivisible}.
Reduction gives the specified finite action of $\Gamma$ on
$\mathcal G_k$.

Lazard's height classification identifies $\mathcal G_k$ with
$G_0=E_0[3^\infty]$
\cite[Theorem~IV and its proof, pp.~270--271]{Lazard1955Formal}.
By Lemma~\ref{lem:B-RD35-conjugacy}, this identification can be
chosen so that the reduced action is the $3$-divisible-group action
of a homomorphism $\Gamma\to\operatorname{Aut}(E_0)$.
The homomorphism, including its labelling by elements of $\Gamma$,
is retained; no alteration of the inertia representation is made.

For $R_n=R/\pi^nR$, the prime three is nilpotent. Serre--Tate
equivalence gives a compatible elliptic deformation $\mathcal E_n/R_n$
of $E_0$, together with an isomorphism
$\mathcal E_n[3^\infty]\simeq\mathcal G_{R_n}$ respecting the
chosen special-fibre identification
\cite[Theorem~1.2.1, printed p.~143]{Katz1981SerreTate}.
Compatibility follows from uniqueness in that equivalence.
Each identity section is an effective relative Cartier divisor,
and the line bundles $\mathcal O_{\mathcal E_n}(3\cdot0)$ are
compatible and relatively ample. Grothendieck existence therefore
algebraizes the formal system to a projective curve
$\mathcal E/R$; the identity and group operations algebraize by
formal full faithfulness. The formal flatness and smoothness
criteria give smoothness along the closed fibre. Since $R$ is
local and $\mathcal E$ is proper, the closed nonsmooth locus
cannot avoid the closed fibre unless it is empty. Thus
$\mathcal E$ is an elliptic scheme over $R$. The exact existence
and full-faithfulness passages used here are the same as in C.1,
namely \cite[Theorems~5.4.1 and~5.4.5]{EGAIII1}.
For every $r$, the compatible finite-level Hopf algebra
isomorphisms modulo $\pi^n$ pass to an isomorphism over $R$,
because the coordinate algebras of $\mathcal G[3^r]$ and
$\mathcal E[3^r]$ are finite free and $R$ is complete.
They commute with the transitions in $r$. Hence
\[
                    \mathcal E[3^\infty]\simeq\mathcal G
                                                        \tag{B-RD35.4}
\]
integrally over $R$.

For every $\gamma\in\Gamma$, the specified special automorphism
of $E_0$ and the corresponding semilinear isomorphism of
$\mathcal G$ lift, by Serre--Tate full faithfulness on each $R_n$,
to a semilinear isomorphism of $\mathcal E_n$. The lifts are
compatible in $n$, and their cocycle relations hold by the same
uniqueness. Proper formal full faithfulness algebraizes these
isomorphisms and relations over $R$.

Restrict this descent datum to the generic elliptic curve
$\mathcal E_L$. Finite Galois descent is effective: the identity
section and its canonical ample line bundle
$\mathcal O(3\cdot0)$ are preserved, giving polarized projective
descent. We obtain an elliptic curve $E/K$. The isomorphism
(B-RD35.4) respects the descent data, so its generic restriction
descends to (B-RD35.3). Only the generic elliptic curve is
descended here; the ramified good model is not claimed to descend
to an elliptic scheme over $\mathcal O_K$.

The last assertion now follows from Theorem~\ref{thm:B-RD34}.
Conversely, such elliptic pairs are a subclass of the original
$3$-divisible-group question. Thus existence of a remaining
comparison after this unramified completion is equivalent to
existence in the elliptic subclass, with the same exact torsion
isomorphism and the same finite inertia characters.
\end{proof}

\begin{remark}[An explicit bounded-dimensional parameterization]\label{auto:B:665ea52b93eb}
\status{Proved under stated hypotheses}

Since $2$ is invertible, every curve in this reduction has an
integral Weierstrass equation
\[
        y^2=x^3+a_2x^2+a_4x+a_6,\qquad a_2,a_4,a_6\in\mathcal O_K,
                                                        \tag{B-RD35.5}
\]
with nonzero discriminant, after scaling. Hence the remaining
question concerns two triples of integral coefficients and an
isomorphism of their rank-$3^6$ generic group schemes of
$27$-torsion. The proposition justifies this reduction for every
allowed potentially supersingular BT pair; it does not evaluate
the resulting congruence problem or assert that such a pair exists.
\end{remark}

\begin{remark}[Exact quaternion checks]\label{auto:B:4074b39ba0d6}
\status{Proved under stated hypotheses}

The retained program \path{code/verify_elliptic_inertia_conjugacy.py}
works in the rational quaternion algebra
$\mathbf Q(\zeta_3)\oplus\mathbf Q(\zeta_3)j$.
It verifies the displayed normalizer and conjugator identities,
all $144$ products of the twelve automorphisms, their element-order
distribution, and the integral-normalization algorithm on eighty
exact rational conjugates. The author run passed in
$0.15318646800005808$ seconds. These finite checks illustrate the
explicit algebraic proof of all cases; they do not replace its
classification or its valuation argument.
\end{remark}


\subsection{A coefficient and discriminant form of the remaining level-$27$ test}

\begin{proposition}[B-RD36: the tame character is read from the discriminant]
\label{prop:B-RD36}
\status{Proved. Complete root and independent proof, original-source, and exact-certificate audits passed.}
Let $K$ have algebraically closed residue field of characteristic three,
and let $E/K$ be an elliptic curve with potentially supersingular
reduction. For any Weierstrass equation put
$\delta=v_K(\Delta)$. The order $t$ of its tame special-inertia
quotient is
\[
                  t=\frac4{\gcd(4,\delta)}.
                                                        \tag{B-RD36.1}
\]
Here $\gcd(4,0)=4$, and changing the equation changes $\delta$ by
a multiple of twelve. Hence the three possibilities are
\[
 \begin{array}{c|ccc}
 \delta\pmod4&0&2&1\text{ or }3\\ \hline
 t&1&2&4.
 \end{array}
                                                        \tag{B-RD36.2}
\]
Together with (B-RD32.2), this determines the requested
prime-to-$3$ character as a function on inertia.

For absolute ramification $e_K=2$, the unresolved alternative
$\nu_{\min}(3,2,2)=3$ is consequently equivalent to the existence,
for some algebraically closed $k$ of characteristic three, of two
triples $(a_2,a_4,a_6)$ in
\[
 R_k=W(k)[\pi]/(\pi^2-3),\qquad K_k=\operatorname{Frac}(R_k),
                                                        \tag{B-RD36.3}
\]
such that the integral equations
$y^2=x^3+a_2x^2+a_4x+a_6$ have nonzero discriminants, satisfy
\[
        0\leq v(\Delta)\leq23,\qquad
        3v(c_4)>v(\Delta),
                                                        \tag{B-RD36.4}
\]
have isomorphic generic $27$-torsion group schemes, and their
discriminant valuations belong to different columns of
(B-RD36.2). The convention is $v(0)=+\infty$.
This is a six-coefficient elliptic test. It is not asserted here
to be a finite residue-field elimination that has been evaluated.
\end{proposition}
\begin{proof}
Take a minimal faithful Galois good-reduction extension $L/K$.
Its group has order $m=3^wt$, with $w\in\{0,1\}$ and
$t\in\{1,2,4\}$, by the norm-one classification and the
canonical alternating pairing of an elliptic curve. A change
from the given Weierstrass equation to a good equation over
$\mathcal O_L$ has scaling factor $u$, with
\[
                  r=v_L(u)=m\delta/12\in\mathbf Z.
                                                        \tag{B-RD36.5}
\]
Choose a tame element $\tau$ of order $t$. Its action on a
uniformizer has residual multiplier a primitive $t$-th root
$\xi$. The residual multiplier of $\tau(u)/u$ is therefore
$\xi^r$. Up to inversion, it is the multiplier on the invariant
differential of the good special elliptic curve. By the explicit
automorphisms (B-RD35.1), the tangent character has precisely the
wild order-three group as its kernel. Thus the multiplier of
$\tau$ has order exactly $t$, and $\gcd(r,t)=1$.

If $t=1$, equation (B-RD36.5) gives $4\mid\delta$.
If $t>1$, $r$ is odd; taking two-adic valuations in
$12r=3^wt\delta$ gives
$2=v_2(t)+v_2(\delta)$. This proves (B-RD36.1)--(B-RD36.2).
The character identification follows from
Lemma~\ref{lem:B-RD32-norm}.

Every complete mixed-characteristic DVR of absolute ramification
two with algebraically closed residue field $k$ contains its Witt
coefficient ring. Write $3=\varpi^2v$ for a uniformizer $\varpi$
and a unit $v$. Since $k$ is algebraically closed and $2$ is
invertible, Hensel lifting gives a square root of $v$. Replacing
$\varpi$ by $\varpi\sqrt v$ gives (B-RD36.3); the degree-two
coefficient-ring description follows from the Cohen structure
argument used in Proposition~\ref{prop:B-RD24-finite-test}.
Such a field contains $\mu_3$, since its coefficient ring contains
a square root of $-1$.

By Proposition~\ref{prop:B-RD35}, every remaining BT comparison is
an elliptic comparison over one of these fields. Choose a minimal
integral Weierstrass equation and complete the square, which is an
integral change since $2$ is a unit. The discriminant bound in
Proposition~\ref{prop:B-RD24-finite-test}, with $p=3,e=2$,
gives $v(\Delta)\leq6e+11=23$.

We justify the strict $j$-condition in (B-RD36.4). Potential good
reduction is equivalent to integral $j$
\cite[Section~2, p.~496]{SerreTate1968Good}. On a good equation
in characteristic three the Cartier operator on
$\omega=dx/y$, for $y^2=x^3+a_2x^2+a_4x+a_6$, is
\[
                 C(\omega)=a_2^{1/3}\omega.
\]
Indeed $\omega=(x^3+a_2x^2+a_4x+a_6)\,dx/y^3$, and Cartier
selects the $x^2dx$ coefficient. Also $c_4\equiv a_2^2\pmod3$.
Thus a good elliptic special fibre is supersingular exactly when
its $j$-invariant is zero. Since $j=c_4^3/\Delta$, potential
supersingularity is exactly $3v(c_4)>v(\Delta)$, including $j=0$.

Conversely any pair in the stated coefficient test is potentially
supersingular and gives two allowed height-two BT groups. Its
different discriminant columns give different prime-to-$3$
characters, while its isomorphic $27$-torsion gives a failure at
manuscript index two. The proved upper bound in
Theorem~\ref{thm:B-RD34} then makes the least index three.
If no such pair exists, B-RD35 and B-RD34 instead make the least
index two. No existence or nonexistence of such a pair is inferred
from this reformulation.
\end{proof}

\par\smallskip\noindent\textbf{Exact finite certificate.}
The retained program
\path{code/verify_discriminant_tame_character.py} checks all twenty-four
displayed discriminant valuations against the necessary scaling equation
$12r=m\delta$ and the tame coprimality condition. It also checks 4,221
changes $\delta\mapsto\delta+12z$, 977 instances of the preceding crude
minimal-discriminant bound, and the exact polynomial identities
$c_4^3-c_6^2=1728\Delta$, $c_4\bmod3=a_2^2$, and the Cartier coefficient
$a_2$. The author run passed in $0.04732804800005397$ seconds;
\path{code/verify_discriminant_tame_character.output.json} retains its
full output. These finite checks certify the displayed arithmetic and
identities. They assert neither existence of all necessary scaling cases
nor existence of a level-$27$ comparison pair. Independent reruns by the
root reviewer and the second reviewer passed in respectively
$0.04524370399667532$ and $0.061048945997754345$ seconds; the author output
was preserved.

\par\smallskip\noindent\textbf{Subsequent exact determination.}
The remaining small-height alternatives isolated here are settled in
Theorem~\ref{thm:B-RD42}: the exact manuscript index is
$\nu_{\min}(3,2,2)=2$. The intermediate bounds and conditional
reductions above are retained as steps in that proof.


\subsection{Scalar residual representations are excluded by a unit obstruction}

\begin{lemma}[The five Hesse valuation patterns at absolute ramification two]
\label{lem:B-RD37-patterns}
\status{Proved. Complete root and independent proof, original-source, and exact-certificate audits passed.}
Let $K$ have algebraically closed residue field of characteristic three,
absolute ramification index two, and normalized valuation $v$.
Fix a primitive cube root $\zeta\in K$. For
\[
 E_t:\quad y^2+3txy+(t^3-1)y=x^3,\qquad t^3\ne1,
                                                        \tag{B-RD37.1}
\]
suppose $E_t$ is potentially supersingular. The ordered vector
$\boldsymbol a(t)=(v(t-1),v(t-\zeta),v(t-\zeta^2))$ has one
of the following five forms; the indicated permutations are allowed.
\[
\begin{array}{c|c|c}
\text{form}&\boldsymbol a(t)&\text{tame special-inertia order}\\\hline
\mathrm A&(-1,-1,-1)&4\\
\mathrm B&(0,0,0)&2\\
\mathrm C&(1,1,1)&4\\
\mathrm D&\text{a permutation of }(2,1,1)&2\\
\mathrm E&\text{a permutation of }(3,1,1)&4
\end{array}                                             \tag{B-RD37.2}
\]
The two sets of residue vectors modulo three in the last-column
classes two and four are disjoint, and remain disjoint after negating
either vector and permuting its coordinates.
\end{lemma}
\begin{proof}
The exact invariants are
\[
 c_4=9t(t^3+8),\qquad \Delta(E_t)=27(t^3-1)^3,\qquad
 j(E_t)=\frac{27t^3(t^3+8)^3}{(t^3-1)^3}.
                                                        \tag{B-RD37.3}
\]
If $v(t)=-a<0$, then $v(j)=6-3a$. Potential supersingularity
requires $v(j)>0$, so $a=1$, giving form A. If $t$ is integral
and its residue is not one, all three entries are zero, giving B.
This includes $t=0$, with the usual convention $v(0)=+\infty$.

It remains that $\overline t=1$. The three roots of unity have
pairwise difference valuation one, since
$(\zeta-1)^2=-3\zeta$. Therefore every entry is at least one,
and at most one is greater than one. Write the vector, after
permutation, as $(r,1,1)$ with $r\ge1$. Then
\[
 v(t^3-1)=r+2,
 \qquad t^3+8=(t^3-1)+9,
 \qquad v(9)=4.
\]
For $r=1$ the numerator factor also has valuation three and
$v(j)=6$. For $r=2$ its valuation is at least four and
$v(j)\ge6$, including $j=0$. For $r\ge3$ its valuation is
exactly four, and $v(j)=12-3r$; positivity permits only $r=3$.
This proves the list.

Put $S=\sum a_i(t)$. The displayed discriminant has valuation
$\delta=6+3S$, respectively $-3,6,15,18,21$ in the five forms.
Proposition~\ref{prop:B-RD36} gives the last column; negative raw
discriminant valuations are permitted since the same proposition
applies to any Weierstrass equation. Modulo three the order-two
vectors are $(0,0,0)$ and permutations of $(2,1,1)$. The
order-four vectors are $(2,2,2)$, $(1,1,1)$ and permutations of
$(0,1,1)$. Negating the second order-two form gives permutations
of $(1,2,2)$, which occur in neither order-four list. The asserted
disjointness follows, also for arbitrary coordinate permutations.
\end{proof}

\begin{lemma}[Cubes close to one at absolute ramification two]
\label{lem:B-RD37-units}
\status{Proved. Complete root and independent proof, original-source, and exact-certificate audits passed.}
If $u\in(K^\times)^3$ and $v(u-1)>0$, then $v(u-1)\ge3$.
\end{lemma}
\begin{proof}
Write $u=x^3$. Then $x$ is a unit and its residue has cube one.
In characteristic three this forces $\overline x=1$. Write
$x=1+w$, with $v(w)\ge1$. Each term of
$3w+3w^2+w^3$ has valuation at least three, proving the claim.
\end{proof}

\begin{proposition}[B-RD37: scalar residual representations are excluded]
\label{prop:B-RD37}
\status{Proved. Complete root and independent proof, original-source, and exact-certificate audits passed.}
Let $K$ be as in Lemma~\ref{lem:B-RD37-patterns}. Let $E,E'/K$
be potentially supersingular elliptic curves with
$E[9]\simeq E'[9]$. If their common residual representation on
three-torsion is scalar, then their finite special-inertia characters
agree on all inertia elements of order prime to three.

Consequently the common residual image in the remaining level-$27$
comparison of Theorem~\ref{thm:B-RD34} is non-scalar. After a
simultaneous quadratic twist it is either the nontrivial unipotent
cyclic group $C_3$ or the irreducible tame cyclic group $C_4$.
This excludes every scalar residual representation from both
remaining joint-inertia orders eighteen and thirty-six.
\end{proposition}
\begin{proof}
First suppose the three-torsion is rational. The full-level moduli
description of $Y(3)$ and its elliptic family, in
\citet[Section~2, p.~4; Lemma~4.2(i), p.~8]{Fisher2015Nine},
allows us to write $E$ as (B-RD37.1), with parameter $t_E\in K$.
An arbitrary isomorphism of nine-torsion has a Weil-pairing multiplier
in $(\mathbf Z/9)^\times$. Every such multiplier is plus or minus
a square. Composing with a scalar automorphism of the target torsion
therefore makes the isomorphism either symplectic or antisymplectic.
These are exactly the direct and reverse nine-congruences in
\citet[Introduction, p.~1; Section~1, first paragraph, p.~2]
{Fisher2015Nine}.

The precise original equations of
\citet[Theorem~4.3, equations~(10)--(11), pp.~9--10]
{Fisher2015Nine}, with its parameter $\Delta=1$, give a parameter
$t\in K$ representing $E'$ by the same family (B-RD37.1).
For every $\theta\in\{1,\zeta,\zeta^2\}$ they impose respectively
\[
 \frac{t-\theta}{t_E-\theta}\in(K^\times)^3
       \quad\text{in the direct case},
 \qquad
 \frac{(t-\theta^{-1})(t_E-\theta)}3\in(K^\times)^3
       \quad\text{in the reverse case}.
                                                        \tag{B-RD37.4}
\]
These are necessary equations for the actual nine-congruence, not
an inference from equality of the mod-three representations.
In particular the two valuation vectors agree modulo three in the
direct case. In the reverse case their coordinate sums are two
modulo three, with the permutation $\theta\mapsto\theta^{-1}$:
\[
       a_i(t)+a_{i^{-1}}(t_E)\equiv v(3)=2\pmod3.
                                                        \tag{B-RD37.5}
\]
The factor three in the reverse equation is essential; it is
present in the original printed equation (11).
Lemma~\ref{lem:B-RD37-patterns} excludes different tame columns
in the direct case. Equal tame quotient orders give equal trace
functions by Lemma~\ref{lem:B-RD32-norm}.

Suppose now that the congruence is reverse and the tame orders
are different. Interchanging the curves if necessary, take $E$
to have tame order two. Comparing the five forms in (B-RD37.5)
permits only B/A or D/E, with the first letter denoting the
source $t_E$. We exclude both by their unit parts; the valuation
congruence alone does not exclude them.

In case B/A, $v(t_E-1)=0$ and $v(t)=-1$. Put
\[
 R_\theta=\frac{(t-\theta^{-1})(t_E-\theta)}3.
\]
Every $R_\theta$ is a cube by (B-RD37.4), whereas
\[
 \frac{R_\zeta}{R_1}
  =\left(1+\frac{1-\zeta^2}{t-1}\right)
   \left(1+\frac{1-\zeta}{t_E-1}\right).
                                                        \tag{B-RD37.6}
\]
The two deviations have valuations two and one. The deviation of
this product from one therefore has valuation exactly one, because
its three summands have distinct valuations $2,1,3$. This contradicts
Lemma~\ref{lem:B-RD37-units}.

In case D/E, let $\theta_0$ be the source root with valuation two.
Equation (B-RD37.5) makes $\theta_0^{-1}$ the target root with
valuation three. Replace $t_E,t,\theta$ algebraically by
$t_E/\theta_0,t\theta_0,\theta/\theta_0$; the products in
(B-RD37.4) are unchanged. Thus the exceptional root is one.
Write $t_E=1+a$, $t=1+b$, with $v(a)=2$, $v(b)=3$, and put
$\lambda=\zeta-1$, $\lambda_2=\zeta^2-1$. Their exact product
is $\lambda\lambda_2=3$. The reverse cube condition at $\zeta$
becomes
\[
 R_\zeta=\frac{(a-\lambda)(b-\lambda_2)}3
      =1-\frac a\lambda-\frac b{\lambda_2}+\frac{ab}3.
                                                        \tag{B-RD37.7}
\]
The three nonconstant terms have valuations $1,2,3$, respectively.
Thus $v(R_\zeta-1)=1$, again contradicting the cube lemma.
No reverse opposite-character case remains. This proves the
proposition when the three-torsion is rational, for arbitrary
Weil-pairing multiplier.

For a scalar residual representation, write it as
$\varepsilon\operatorname{id}$ with
$\varepsilon:G_K\to\mathbf F_3^\times=\{\pm1\}$.
Use the same quadratic character to twist both elliptic curves.
Their three-torsion becomes rational and the nine-torsion
isomorphism persists. Potential supersingularity also persists.
Over a common good-reduction field the new descent action is the
old action multiplied by the scalar automorphism
$[\varepsilon(\sigma)]$. Hence both special-inertia characters
are multiplied by the same function $\varepsilon$; equality after
twisting is equivalent to equality before twisting. This proves
the assertion in the scalar case as well. The Weil-pairing
multiplier is preserved by the simultaneous twist, although the
preceding argument treats both its square classes.

For completeness, after the completed maximal unramified extension
the entire Galois group is inertia. The common residual image lies
in $\operatorname{SL}_2(\mathbf F_3)$ because $\mu_3\subset K$.
Its wild subgroup is a normal three-subgroup and its quotient is
cyclic of order prime to three. If the wild subgroup is nontrivial,
it has order three and its normalizer in
$\operatorname{SL}_2(\mathbf F_3)$ is $C_6$; the image is $C_3$
or $C_6$. A simultaneous quadratic twist removes the central
minus sign in the latter case. If the wild subgroup is trivial,
the image is cyclic of order $1,2$ or $4$. The first two are
scalar and have just been excluded. An order-four generator has
irreducible characteristic polynomial $X^2+1$ over $\mathbf F_3$.
These elementary subgroup assertions are also exhaustively checked
by the retained finite certificate. Thus the two stated
residual forms remain as possible tests; their existence as
opposite-character level-$27$ pairs is not asserted.
\end{proof}

\begin{remark}[Compatibility with the known level-nine ambiguity]
\label{rem:B-RD37-known-example}
\status{Proved under stated hypotheses}
The pair in Proposition~\ref{prop:B-RD29-p3-level-two} consists of
a potentially supersingular elliptic group and a Tate group with
good rank-one graded pieces. Moreover its common three-torsion is
non-scalar. Indeed its Kummer extension parameter is
$q=2^5 3^6$, whose cube class is nontrivial even after completed
unramified base change. If $2$ were a cube, so would be $-2$.
A putative unit $x$ with $x^3=-2$ has residue one; writing
$x=1+w$ gives $v(w)\ge1$ and
\[
 v\bigl((1+w)^3-1\bigr)
   =v(3w+3w^2+w^3)\ge3,
\]
whereas $v(-2-1)=v(-3)=2$. This is impossible.
Thus the example has nontrivial unipotent residual image $C_3$,
and supplies no exception to Proposition~\ref{prop:B-RD37}.
\end{remark}

\par\smallskip\noindent\textbf{Exact finite certificate and correction record.}
The program \path{code/verify_scalar_residual_patterns.py} checks
the nine ordered valuation patterns, all 240 opposite-column
comparisons allowing arbitrary coordinate permutations, every
nine-torsion multiplier, both displayed Weierstrass invariant
identities, both reverse-unit identities and their distinct valuation
terms, and all fifteen subgroups of
$\operatorname{SL}_2(\mathbf F_3)$. Its run passed in
$0.0822893020012998$ seconds, with the complete output retained
beside the program. In particular it uses the reverse valuation
shift \emph{two}. An earlier unpromoted argument omitted the
factor three in Fisher's original equation (11); a visual source
check caught this error. That argument is retained as withdrawn
in \path{failed_candidates/B37_scalar_exclusion_missing_factor3.withdrawn.tex}.
The current proof adds the exact unit obstructions (B-RD37.6)--
(B-RD37.7), which exclude the two reverse valuation patterns left
by the corrected equation. It thereby excludes scalar residual
level-$27$ pairs, while leaving the non-scalar cases and the least
index unresolved. Independent author-output-preserving reruns by the root
reviewer and a second reviewer passed in $0.06197141099983128$ and
$0.09311267000157386$ seconds, respectively.


\subsection{Irreducible tame residual image leaves only joint inertia eighteen}

\begin{proposition}[B-RD38: the cyclic-four residual case]
\label{prop:B-RD38}
\status{Proved. Complete root and independent proof, original-source, and exact-certificate audits passed.}
Let $K$ be a complete field of mixed characteristic $(0,3)$ with
algebraically closed residue field and absolute ramification index two.
Suppose an elliptic curve $E/K$ has residual image
$\rho_{E,3}(G_K)=C_4\subset\operatorname{SL}_2(\mathbf F_3)$.
Then
\[
                  v_K(j(E))=6,
       \qquad     6\mid v_K(\Delta(E)),
                                                        \tag{B-RD38.1}
\]
where the discriminant is that of any Weierstrass equation.
In particular $E$ is potentially supersingular and its tame
special-inertia quotient has order one or two, never four.

Consequently, in the remaining level-$27$ test of
Theorem~\ref{thm:B-RD34}, a common irreducible $C_4$ residual image
can occur with different prime-to-three characters only if the two
individual faithful special-inertia orders are $3$ and $6$, and the
joint faithful order is $18$. Every possible failure with joint order
$36$ has nontrivial unipotent residual image $C_3$ after the common
quadratic twist in Proposition~\ref{prop:B-RD37}.
\end{proposition}
\begin{proof}
Write $L=K(E[3])$. Since its Galois group has order four, $L/K$
is totally tamely ramified and cyclic. Let $F$ be its quadratic
subfield and let $\sigma$ be the nontrivial element of
$\operatorname{Gal}(F/K)$. The residual image of $G_F$ is the
scalar subgroup $\{\pm I\}$, with nontrivial character given by
$L/F$.

We first identify the projective level-three parameter carefully.
The full-level modular curve $X(3)$ carries the action of
$\operatorname{SL}_2(\mathbf F_3)/\{\pm I\}=A_4$, with quotient
the $j$-line, as in \citet[Section~2, p.~4]{Fisher2015Nine}.
Use the parameter and universal elliptic family in
\citet[Lemma~4.2(i), p.~8]{Fisher2015Nine}:
\[
 E_t:\quad y^2+3txy+(t^3-1)y=x^3,
 \qquad
 j(t)=\frac{27t^3(t^3+8)^3}{(t^3-1)^3}.
                                                        \tag{B-RD38.2}
\]
The nontrivial projective residual image is a double transposition
in $A_4$. After changing the symplectic basis, it acts by
\[
                  h(t)=\frac{t+2}{t-1}.
                                                        \tag{B-RD38.3}
\]
Here no transformation formula is being assumed without a check:
$h^2=1$, $j\circ h=j$, and $h$ interchanges the cusp pairs
$\{\infty,1\}$ and $\{\zeta,\zeta^2\}$, where $\zeta\in K$
is a primitive cube root of unity. Thus $h$ is one of the three
double transpositions, which are conjugate under $A_4$.

A symplectic basis of $E[3]$ over $L$ gives a noncuspidal point
of $X(3)$. The central scalar $-I$ acts trivially on this modular
curve, so its coordinate $t$ belongs to $F$, and
\[
                         \sigma(t)=h(t).
                                                        \tag{B-RD38.4}
\]
This permits $t\in K$ when the point is fixed by $h$; the ensuing
norm identity still holds. Put $z=t-1$. Equation (B-RD38.4)
becomes $\sigma(z)=3/z$, whence $N_{F/K}(z)=3$.

Choose $\pi\in K$ with $\pi^2=3$, as in
Proposition~\ref{prop:B-RD36}. With normalized valuations,
$v_F(3)=4$ and $v_F(\pi)=2$. Since $F/K$ is totally ramified,
$v_K(N_{F/K}(z))=v_F(z)$, so $v_F(z)=2$.
The unit $u=z/\pi$ has norm one. The residue action of
$\sigma$ is trivial, and therefore $\overline u^2=1$:
\[
                         \overline u=1\quad\hbox{or}\quad-1.
                                                        \tag{B-RD38.5}
\]
On the other hand, $(\zeta-1)^2=-3\zeta$ gives
\[
             \left(\frac{\zeta-1}{\pi}\right)^2=-\zeta.
\]
Its residual square is $-1$, so its residue is neither $1$
nor $-1$. The same holds for $(\zeta^2-1)/\pi$.
It follows that the three quantities $t-1,t-\zeta,t-\zeta^2$
all have $F$-valuation two. Thus
\[
 v_F(t^3-1)=6,\qquad v_F(t^3+8)=6,
 \qquad v_F(\Delta(E_t))=30,
 \qquad v_F(j(t))=12.                                  \tag{B-RD38.6}
\]
For the second equality use $t^3+8=(t^3-1)+9$ and $v_F(9)=8$.
Since $j(t)=j(E)\in K$, the last equality gives $v_K(j(E))=6$.

It remains to compare the discriminants of the actual elliptic
curves, rather than just their common modular point. Over $F$,
twist $E_F$ by the quadratic character of $L/F$. This kills its
scalar residual character and produces a full rational level-three
basis. Its point in $X(3)$ is exactly the point $t$ above: the
quadratic twist multiplies the basis descent by a central sign,
which is invisible in $X(3)$. The full-level moduli description
therefore identifies this twist with $E_t$ over $F$.

The quadratic extension $L/F$ is tame and ramified. It has the
form $F(\sqrt d)$ with $v_F(d)$ odd; every unit is a square by
Hensel lifting over the algebraically closed residue field.
A quadratic twist changes the discriminant by $d^6$, up to a
twelfth power from a Weierstrass coordinate change. Hence, writing
$\delta=v_K(\Delta(E))$, equation (B-RD38.6) gives
\[
                   30\equiv2\delta+6\pmod {12},
       \qquad\text{and thus}\qquad 6\mid\delta.
                                                        \tag{B-RD38.7}
\]
This comparison also covers the special values of $j$: the
full-level structure eliminates elliptic automorphisms fixing the
basis, so the identified twist is the same marked object, not
merely a curve with the same $j$-invariant.

The strict positivity of $v_K(j)$ gives potential supersingularity
by Proposition~\ref{prop:B-RD36}, whose discriminant formula now
makes the tame special-inertia order one or two. In a remaining
level-$27$ failure, Theorem~\ref{thm:B-RD34} permits only joint
orders eighteen or thirty-six. Joint order thirty-six would require
tame order four on one side, so it is excluded here. Joint order
eighteen with distinct tame characters forces tame orders one and
two and independent wild order-three images. The individual
faithful orders are therefore three and six. Finally apply the two
residual alternatives in Proposition~\ref{prop:B-RD37}.
\end{proof}

\par\smallskip\noindent\textbf{Exact finite certificate.}
The retained program \path{code/verify_c4_residual_gate.py} checks
the modular involution, its cusp permutation, the exact $c_4$ and
discriminant transformation factors, the two normalized residues in
$\mathbf F_9$, the quadratic-twist discriminant identity, and the
valuation congruence (B-RD38.7). The author run passed in
$0.07729700100026093$ seconds; its complete output is retained
beside the program.
These checks establish the finite identities used above; neither
existence nor nonexistence of a remaining $C_3$ or $C_4$ level-$27$
comparison is inferred. Independent author-output-preserving reruns by
the root and second reviewers passed in $0.1112131780027994$ and
$0.10795447899727151$ seconds, respectively.

\par\smallskip\noindent\textbf{Subsequent exact determination.}
The remaining small-height alternatives isolated here are settled in
Theorem~\ref{thm:B-RD42}: the exact manuscript index is
$\nu_{\min}(3,2,2)=2$. The intermediate bounds and conditional
reductions above are retained as steps in that proof.


\subsection{Unipotent residual image leaves only the order-twelve versus order-six pair}

\begin{lemma}[The two unit depths over absolute ramification two]
\label{lem:B-RD39-units}
\status{Proved. Complete root and independent proof, original-source, and exact-certificate audits passed.}
Let $K$ be complete of mixed characteristic $(0,3)$, with algebraically
closed residue field, and normalize $v_K(3)=2$. Choose $\pi^2=3$.
Every unit in $1+\pi^3\mathcal O_K$ is a cube. A noncube unit $D$
therefore has a well-defined cube-approximation depth $s\in\{1,2\}$:
for every unit $t$ with $v_K(t^3-D)>0$, one has
$v_K(t^3-D)=s$. The inverse unit $D^{-1}$ has the same depth.

If $L/K$ is totally ramified of degree three, and a cube
$u\in(L^\times)^3$ satisfies $0<v_L(u-1)<9$, then
$v_L(u-1)$ is either three or six.
\end{lemma}
\begin{proof}
For $b\in\mathcal O_K$, the equation
\[
       (1+\pi w)^3=1+\pi^3b
       \quad\Longleftrightarrow\quad
                     w^3+w+\pi w^2=b
                                                        \tag{B-RD39.1}
\]
has a solution in $\mathcal O_K$: its reduction has a root over the
algebraically closed residue field and its derivative is a unit.
This proves the first assertion. Choose a unit $t$ whose residue is
the unique cube root of $\overline D$. The positive valuation of
$t^3-D$ must be one or two, since valuation at least three would
make $D/t^3$, and hence $D$, a cube. Any two such $t$ differ by a
multiple of $\pi$, and their cubes differ by a multiple of $\pi^3$.
The positive valuation is therefore independent of $t$.
Replacing $t$ and $D$ by their inverses preserves it.

For the last assertion write $u=x^3$. Its residue is one, so
$x=1+w$ with $r=v_L(w)\ge1$. Since $v_L(3)=6$, the three terms in
$u-1=3w+3w^2+w^3$ have valuations $6+r,6+2r,3r$.
For $r=1,2$ their minimum is uniquely $3r$. For $r\ge3$ every term
has valuation at least nine. This proves the claim.
\end{proof}

\begin{proposition}[B-RD39: the cyclic-three residual case]
\label{prop:B-RD39}
\status{Proved. Complete root and independent proof, original-source, and exact-certificate audits passed.}
Let $K$ be as in Lemma~\ref{lem:B-RD39-units}. Suppose two potentially
supersingular elliptic curves $E,E'/K$ have isomorphic nine-torsion,
common nontrivial unipotent residual image $C_3$, and different
prime-to-three special-inertia characters. Then their common
three-division field is generated by a cube root of a noncube
\emph{unit} of $K$. Their tame special-inertia quotient orders are
two and four, and their $j$-valuations, in some order, are
\[
                            3\quad\hbox{and}\quad b\ge6,
                                                        \tag{B-RD39.2}
\]
where $b=+\infty$ is allowed.

Consequently every remaining level-$27$ comparison with this residual
image has individual faithful special-inertia orders $12$ and $6$,
joint faithful order $36$, and independent wild order-three images.
Combined with Proposition~\ref{prop:B-RD38}, after the simultaneous
quadratic normalization of Proposition~\ref{prop:B-RD37}, the entire
remaining level-$27$ test has the following necessary alternatives:
\[
\begin{array}{c|c|c|c}
\text{residual image}&\text{individual special orders}&
 \text{joint order}&\text{$j$-valuations}\\\hline
C_4&(3,6)&18&(6,6)\\
C_3&(12,6)&36&(3,b),\ b\ge6.
\end{array}                                             \tag{B-RD39.3}
\]
The table asserts necessary restrictions, not existence of either kind
of comparison. Undoing the common quadratic twist may interchange
the tame orders one and two; it leaves order four and both
$j$-invariants unchanged.
\end{proposition}
\begin{proof}
A nontrivial unipotent residual image fixes a nonzero three-torsion
point. By \citet[Lemmas~4.1--4.2, p.~8]{Fisher2015Nine}, write
\[
 E=E_{t_E,D}:\quad y^2+3t_Exy+(t_E^3-D)y=x^3,
 \qquad L=K(E[3])=K(\sqrt[3]D).
                                                        \tag{B-RD39.4}
\]
Concretely, in the rational-three-point form
$y^2+a_1xy+a_3y=x^3$, take $t_E=a_1/3$ and
$D=t_E^3-a_3$. Its actual discriminant is $(3a_3)^3D$,
so its cube class is precisely the parameter class $D$.
The element $D$ is not a cube. A Weierstrass scaling replaces
$(t_E,D)$ by $(t_E/c,D/c^3)$, so we normalize
$d=v_K(D)\in\{0,1,2\}$. The common residual field $L/K$ is
cyclic and totally ramified of degree three. Fix $\theta\in L$
with $\theta^3=D$; then $v_L(\theta)=d$ and
$v_L(\zeta-1)=3$.

An arbitrary nine-congruence can be made direct or reverse by a scalar
change, as in Proposition~\ref{prop:B-RD37}. Fisher's exact
Theorem~4.3, equations~(10)--(11), pp.~9--10, supplies a parameter
$t\in K$ for $E'$ and requires
\[
\begin{array}{ll}
 \displaystyle\frac{t-\theta}{t_E-\theta}\in(L^\times)^3,
   &E'=E_{t,D}\quad\text{(direct)},\\[6pt]
 \displaystyle\frac{(t-\theta^{-1})(t_E-\theta)}3
                 \in(L^\times)^3,
   &E'=E_{t,D^{-1}}\quad\text{(reverse)}.
\end{array}                                             \tag{B-RD39.5}
\]
Each condition also holds with $\theta$ replaced by $\zeta\theta$
or $\zeta^2\theta$. The reverse denominator three is retained.

Here are all the required valuations, with no minimal-model
assumption. For $r=v_K(t)$ and $S=v_K(t^3-D)$, the exact invariants are
\[
 c_4(E_{t,D})=9t(t^3+8D),\qquad
 \Delta(E_{t,D})=27D(t^3-D)^3.
                                                        \tag{B-RD39.6}
\]
If $d=1$ or $2$, the valuations $3r$ and $d$ are unequal.
Potential supersingularity requires $r\ge-1$, and
\[
             S=\min(3r,d),\quad
             v_K(j)=6+3r-d,\quad
             \delta=6+d+3S.
                                                        \tag{B-RD39.7}
\]
The value $t=0$ belongs to the last row below and has $j=0$.
By Proposition~\ref{prop:B-RD36}, the tame quotient is read from
$\delta$ modulo four. The full direct-source tables are
\[
\begin{array}{c|ccc|ccc}
 &\multicolumn{3}{c|}{d=1}&\multicolumn{3}{c}{d=2}\\
r&S&\delta&\text{tame}&S&\delta&\text{tame}\\\hline
-1&-3&-2&2&-3&-1&4\\
0&0&7&4&0&8&1\\
\ge1&1&10&2&2&14&2
\end{array}                                             \tag{B-RD39.8}
\]
For the reverse target, whose parameter is $D^{-1}$, put
$S'=v_K(t^3-D^{-1})$. The analogous complete table is
\[
\begin{array}{c|ccc|ccc}
 &\multicolumn{3}{c|}{d=1}&\multicolumn{3}{c}{d=2}\\
r&S'&\delta'&\text{tame}&S'&\delta'&\text{tame}\\\hline
-2&-6&-13&4&-6&-14&2\\
-1&-3&-4&1&-3&-5&4\\
\ge0&-1&2&2&-2&-2&2
\end{array}                                             \tag{B-RD39.9}
\]

Take norms in (B-RD39.5). Since a norm of a cube is a cube in $K$,
the necessary congruences are
\[
    S-S_E\equiv0\pmod3\quad\text{(direct)},\qquad
    S'+S_E\equiv0\pmod3\quad\text{(reverse)}.
                                                        \tag{B-RD39.10}
\]
In the reverse case the norm of the denominator is $27$, whose
$K$-valuation is six, hence zero modulo three.

Suppose first $d=1$ or $2$. In the direct case, the only
opposite-tame pairs compatible with (B-RD39.10) have
$\{S,S_E\}=\{0,-3\}$. Interchange the pair if needed so that
$S_E=-3$, $S=0$. The quotient of the cube conditions at
$\zeta\theta$ and $\theta$ is
\[
 \frac{1+(1-\zeta)\theta/(t-\theta)}
      {1+(1-\zeta)\theta/(t_E-\theta)}.
                                                        \tag{B-RD39.11}
\]
The two displayed deviations have $L$-valuations $3+d$ and $6+d$.
The quotient is congruent to one with deviation valuation exactly
$3+d$, which is four or five. It cannot be a cube by
Lemma~\ref{lem:B-RD39-units}.

In the reverse case, compare the two tables using (B-RD39.10).
All opposite-tame pairs and the two deviation valuations in their
cube quotient are as follows:
\[
\begin{array}{c|rr|rr}
d&S_E&S'&3-d-S'&3+d-S_E\\\hline
1&-3&-6&8&7\\
1&-3&-3&5&7\\
1&0&-3&5&4\\
2&-3&-6&7&8\\
2&0&-6&7&5\\
2&0&-3&4&5
\end{array}                                             \tag{B-RD39.12}
\]
Indeed, on dividing the condition at $\zeta\theta$ by that at
$\theta$, one obtains the product
\[
 \left(1+\frac{(1-\zeta^{-1})\theta^{-1}}
                      {t-\theta^{-1}}\right)
 \left(1+\frac{(1-\zeta)\theta}{t_E-\theta}\right).
                                                        \tag{B-RD39.13}
\]
The valuations of its two deviations are the last two columns.
They are unequal; their minimum is four, five, or seven. The cross
term has strictly larger valuation. Thus the deviation of this
required cube from one has valuation four, five, or seven, again
contrary to Lemma~\ref{lem:B-RD39-units}. This excludes $d=1,2$
for either orientation.

We are left with the noncube unit $D$, of depth $s=1$ or $2$ in
Lemma~\ref{lem:B-RD39-units}. Its inverse has the same depth.
For either of the parameters $D,D^{-1}$, potential
supersingularity and the unit-depth lemma give exactly the following
possibilities:
\[
\begin{array}{c|c|c|c}
S&\delta&\text{tame}&v_K(j)\\\hline
-3&-3&4&3\\
0&6&2&\ge6\\
1\ (s=1)&9&4&6\\
2\ (s=2)&12&1&6
\end{array}                                             \tag{B-RD39.14}
\]
For $S=-3$ one has $v_K(t)=-1$. For $S=0$ the parameter is
integral and its residue does not equal the cube root of the unit
parameter; the zero value of $t$ is allowed. If $S>0$, $t$ is a
unit and $v_K(t^3+8D)=S$, since $S< v_K(9D)=4$.
These observations and (B-RD39.6) verify the table.

The direct norm congruence in (B-RD39.10) allows different tame
columns only for $\{S,S_E\}=\{-3,0\}$. The reverse norm
congruence likewise allows different columns only for that pair:
both curves have the same positive depth $s$, and $2s$ is not
divisible by three. Formula (B-RD39.14) now proves (B-RD39.2) and
the tame orders two and four.

At level $27$, Theorem~\ref{thm:B-RD34} restricts the joint faithful
order to eighteen or thirty-six. A tame quotient of order four
rules out eighteen, so the joint order is thirty-six. Its wild
subgroup has order nine; the two individual wild quotients have
order three and are independent. Their distinct tame orders four
and two therefore give individual orders twelve and six. The
other row of (B-RD39.3) is Proposition~\ref{prop:B-RD38}.
\end{proof}

\par\smallskip\noindent\textbf{Exact certificate and scope.}
The program \path{code/verify_c3_residual_gate.py} checks the
Weierstrass invariants, the exact Hensel polynomial, both source
and reverse-target tables, every opposite-tame norm-congruence
case, and all the required forbidden cube deviations. Its output
records the author runtime $0.053551534001599066$ seconds. Independent
author-output-preserving reruns by the root and second reviewers passed
in $0.045766205999825615$ and $0.05912446099682711$ seconds.
These calculations narrow the remaining test;
they do not construct or exclude either final row of (B-RD39.3).


\subsection{The remaining cyclic-four comparison must reverse the pairing}

\begin{proposition}[B-RD40: no direct nine-congruence in the cyclic-four branch]
\label{prop:B-RD40}
\status{Proved. Complete root and independent proof, original-source, and exact-certificate audits passed.}
Let $K$ be as in Proposition~\ref{prop:B-RD38}, and let $E,E'/K$
have common irreducible cyclic-four residual representation. If there
is a symplectic isomorphism $E[9]\simeq E'[9]$, their requested
prime-to-three special-inertia characters agree.

Consequently any opposite-character level-$27$ comparison in the
$C_4$ row of (B-RD39.3) has pairing multiplier congruent to $-1$
modulo three. Every such comparison can be made antisymplectic by a
scalar change; none can be made symplectic. This is an orientation
restriction, not an exclusion of the entire $C_4$ row.
\end{proposition}
\begin{proof}
Put $L=K(E[3])=K(E'[3])$ and let $F$ be its quadratic subfield.
Use the given symplectic identification to choose symplectic bases of
both three-torsion groups with the same actual $C_4$ Galois matrix.
As in Proposition~\ref{prop:B-RD38}, choose the bases so that its
projective action on $X(3)$ is
\[
 h(t)=\frac{t+2}{t-1},\qquad \sigma(t)=h(t),
                                                        \tag{B-RD40.1}
\]
where $\sigma\in\operatorname{Gal}(F/K)$ is nontrivial.
Write $t_E,t'$ for the two parameters, $z=t_E-1$, $z'=t'-1$,
and choose $\pi^2=3$ in $K$. Thus
\[
 u=z/\pi,\quad u'=z'/\pi,\qquad
 N_{F/K}(u)=N_{F/K}(u')=1,\qquad
 \overline u,\overline{u'}\in\{1,-1\}.
                                                        \tag{B-RD40.2}
\]

\par\smallskip\noindent\textbf{The actual marked differential cocycle.}
We determine the lift of (B-RD40.1), rather than using only its
projective action. Let $\zeta$ be a primitive cube root and
$c=\zeta-\zeta^2$, so $c^2=-3$. On the Hesse pencil
\[
 F_t(x,y,z)=x^3+y^3+z^3-3txyz
\]
use the matrix
\[
 M=\begin{pmatrix}1&1&1\\1&\zeta&\zeta^2\\
                         1&\zeta^2&\zeta\end{pmatrix}.
\]
Direct polynomial calculation gives
\[
 \det M=-3c,\qquad
 F_t(MX)=3(1-t)F_{h(t)}(X).
                                                        \tag{B-RD40.3}
\]
For the Poincare residue differential
$\eta_t=\operatorname{Res}(\Omega/F_t)$, where
$\Omega=x\,dy\wedge dz-y\,dx\wedge dz+z\,dx\wedge dy$, this gives
\[
                  M^*\eta_t=\kappa(t)\eta_{h(t)},
        \qquad    \kappa(t)=\frac c{t-1},
        \qquad    \kappa(t)\kappa(h(t))=-1.
                                                        \tag{B-RD40.4}
\]
The two Hesse translation generators are the cyclic coordinate
permutation and $\operatorname{diag}(1,\zeta,\zeta^2)$.
Neither generator has a fixed point on a smooth member of the
pencil: the diagonal generator fixes only coordinate vertices, and
the eigenvectors of the cyclic permutation lie on the cubic only
when $t=1$. Thus both are translations; they generate all nine
three-torsion translations. Conjugation by $M$ interchanges them
with one inverted, so its marked action is an order-four symplectic matrix, whose square is
$-I$. Moving the image flex back to the chosen origin is a
translation and does not change the invariant differential.
The possible central sign is fixed by the actual common marked
matrix on $E[3]$ and $E'[3]$, hence is the same on both sides.
The use of these marked Hesse translations agrees with the full-level
moduli convention in \citet[Section~2, pp.~4--5]{Fisher2015Nine}.

We also check that passing to the Weierstrass family used in
Fisher's cube equations introduces no parameter-dependent factor.
On the plane cubic set
\[
              X=z,\qquad Y=x-y,\qquad Z=x+y+tz.
\]
Its equation becomes
\[
 3ZY^2=4(t^3-1)X^3-9t^2X^2Z+6tXZ^2-Z^3.
                                                        \tag{B-RD40.5}
\]
Put $a=4(t^3-1)/3$ and take affine coordinates $aX/Z,aY/Z$,
followed by the usual $x$-translation to remove the square term.
The resulting short Weierstrass coefficients are
\[
 A_H=-\frac{t(t^3+8)}3,\qquad
 B_H=\frac{2(t^6-20t^3-8)}{27}.
\]
The displayed family $E_t:y^2+3txy+(t^3-1)y=x^3$ instead has
\[
 A_E=-\frac{3t(t^3+8)}{16},\qquad
 B_E=-\frac{t^6-20t^3-8}{32}.
\]
They are related by the constant Weierstrass scaling $2/c$.
The determinant of the preceding plane coordinate change is two;
tracking the residue differential through (B-RD40.5) gives
\[
                         \eta_t=\frac c3\,\omega_t,
                                                        \tag{B-RD40.6}
\]
where $\omega_t$ is the invariant differential of $E_t$ under this
identification. Thus (B-RD40.4) is also the exact marked differential
factor for this Weierstrass family. This computation supplies the
normalization behind the constant twist mentioned in
\citet[Lemma~4.2(i), p.~8]{Fisher2015Nine}.

\par\smallskip\noindent\textbf{Different tame characters force opposite unit residues.}
Take marked isomorphisms $E_L\simeq E_{t_E,L}$ and
$E'_L\simeq E_{t',L}$, and let their pullbacks of the universal
invariant differentials be $\alpha\omega_E$ and
$\alpha'\omega_{E'}$, respectively. For a generator
$\tau\in\operatorname{Gal}(L/K)$, equation (B-RD40.4) and the
common actual marked matrix show that
\[
       \frac{\tau(\alpha'/\alpha)}{\alpha'/\alpha}
                =\frac{z'}z
          \quad\hbox{or its inverse},
                                                        \tag{B-RD40.7}
\]
according to the simultaneous choice of generator and direction of
the universal isomorphism. Either choice has the same residue for
the argument below. A common central minus sign cancels in the ratio.

Let $\delta,\delta'$ be discriminant valuations over $K$, and
put $n=v_L(\alpha)$, $n'=v_L(\alpha')$. The invariant discriminant
section $\Delta\omega^{12}$ and the valuation
$v_L(\Delta(E_t))=60$ from Proposition~\ref{prop:B-RD38} give
\[
                 4\delta=60+12n,\qquad
                 4\delta'=60+12n'.
                                                        \tag{B-RD40.8}
\]
If the tame characters differ, their orders are one and two, so
$\delta'-\delta\equiv2\pmod4$. Hence
$n'-n\equiv2\pmod4$. In a tame cyclic extension of degree four,
the residue of $\tau(\beta)/\beta$ for an element of valuation
$m$ is $\xi^m$, where $\xi$ is a primitive fourth root.
Apply this to $\beta=\alpha'/\alpha$. Equation (B-RD40.7) forces
\[
                           \overline{u'}=-\overline u.
                                                        \tag{B-RD40.9}
\]

\par\smallskip\noindent\textbf{The direct nine-congruence contradicts a norm cube.}
Over $F$, twist both curves by the same quadratic character of $L/F$.
Their three-torsion becomes rational, they become the marked curves
$E_{t_E},E_{t'}$, and the symplectic nine-congruence persists.
Fisher's direct equation (10) therefore makes
\[
       \frac{t'-\theta}{t_E-\theta}\in(F^\times)^3
          \qquad(\theta=1,\zeta,\zeta^2).
                                                        \tag{B-RD40.10}
\]
At $\theta=1$ this says $u'/u$ is a cube. Its negative is a cube
congruent to one by (B-RD40.9). Since $v_F(3)=4$, the expansion
$(1+w)^3-1=3w+3w^2+w^3$ shows that any such cube has deviation
valuation at least three. Thus $v_F(u'+u)\ge3$.

Set
\[
 T=z+3/z=\pi(u+u^{-1}),\qquad
 T'=z'+3/z'=\pi(u'+u'^{-1})\in K.
\]
Both $T$ and $T'$ have $K$-valuation one. Moreover
\[
       T'+T=\pi(u'+u)\bigl(1+(u'u)^{-1}\bigr).
\]
The second factor has positive $F$-valuation by (B-RD40.9), and
$v_F(\pi)=2$, so $v_K(T'+T)\ge3$.
For $\lambda=\zeta-1$, the exact relation
$\lambda^2=-3\zeta$ gives
\[
 N_{F/K}(t_E-\zeta)
       =(z-\lambda)(3/z-\lambda)=-\lambda(T+3).
\]
Taking the norm of the cube in (B-RD40.10) at $\theta=\zeta$
makes
\[
                 q=\frac{T'+3}{T+3}\in(K^\times)^3.
\]
But
\[
              -q-1=-\frac{T'+T+6}{T+3}
\]
has $K$-valuation exactly one: its numerator has valuation two
and its denominator valuation one. This contradicts
Lemma~\ref{lem:B-RD37-units}, since $-q$ is a cube congruent to one.
The assumed different tame characters are impossible.

Finally, a unit of $\mathbf Z/27$ congruent to one modulo three is
a square, by Hensel lifting for $X^2-a$. A torsion isomorphism with
such a multiplier could be rescaled to be symplectic, contradicting
the result just proved on nine-torsion. Every remaining multiplier
is congruent to $-1$, and its negative is a square, giving the stated
antisymplectic normalization.
\end{proof}

\par\smallskip\noindent\textbf{Exact certificate and scope.}
The program \path{code/verify_c4_orientation_gate.py} checks all
Fourier, cubic, differential-scaling and short-Weierstrass identities,
the exact norm identity, the valuation inequalities, and every
unit multiplier modulo twenty-seven. Its retained output records
the author runtime $0.1905790330056334$ seconds. An initial
certificate used structural equality for two differently factored
SymPy expressions; that failed version and traceback are retained in
\path{failed_candidates/B40_certificate_structural_equality_v1.py} and
its adjacent output. The corrected test proves their exact difference
is zero. No antisymplectic level-$27$ pair is constructed or
excluded by this proposition alone; it is excluded in
Theorem~\ref{thm:B-RD42}. Independent author-output-preserving reruns by the
root and second reviewers passed in $0.1678152059976128$ and
$0.1795537689977209$ seconds, respectively.


\subsection{A three-isogeny eliminates the unipotent residual branch}

\begin{lemma}[The normalized three-isogeny and its invariants]
\label{lem:B-RD41-isogeny}
\status{Proved; independently checked by the root and obstruction agents, including the complete proof and exact certificate.}
Over a characteristic-zero field let
\[
 E:\ y^2+a_1xy+a_3y=x^3,\qquad a_3(a_1^3-27a_3)\ne0.
\]
The quotient by the subgroup generated by $(0,0)$ is
\[
 E_1:\ y^2+a_1xy+a_3y
       =x^3-5a_1a_3x-a_3(a_1^3+7a_3).
                                                        \tag{B-RD41.1}
\]
Its invariants are
\[
 c_4(E_1)=a_1(a_1^3+216a_3),\qquad
 \Delta(E_1)=a_3(a_1^3-27a_3)^3.
                                                        \tag{B-RD41.2}
\]
For the Hesse-form coefficients $a_1=3t$, $a_3=t^3-D$, this gives
\[
          j(E_1)=\frac{27t^3(9t^3-8D)^3}{D^3(t^3-D)}.
                                                        \tag{B-RD41.3}
\]
\end{lemma}
\begin{proof}
The point $(0,0)$ has order three, and its inverse is $(0,-a_3)$,
as in \citet[Lemma~4.1, p.~8]{Fisher2015Nine}. Here the quotient
can be checked directly. On the open set $x\ne0$ set
\[
\begin{split}
 X&=x+\frac{a_1a_3}{x}+\frac{a_3^2}{x^2},\\
 Y&=y-\frac{a_3(a_1x+2a_3)y}{x^3}
        -\frac{a_3(a_1x+a_3)^2}{x^3}.
\end{split}                                             \tag{B-RD41.4}
\]
Substitution into (B-RD41.1), followed by reduction using the source
equation, gives zero identically. Also
\[
          2Y+a_1X+a_3=(2y+a_1x+a_3)\frac{dX}{dx},
\]
so the invariant differential is normalized. These are universal
polynomial identities after clearing denominators, checked by the
retained exact certificate.

The rational map extends to the smooth projective curves. The
function $X$ has pole divisor exactly
$2(O+(0,0)+(0,-a_3))$: at the two finite points $x$ is a uniformizer
and the term $a_3^2/x^2$ has nonzero leading coefficient. At $O$,
the leading term is $x$. Therefore the morphism has degree three
and sends all three displayed points to the origin. An origin-preserving
morphism of elliptic curves is a group homomorphism (apply rigidity
to the difference of the two addition maps). Its kernel is exactly
the indicated subgroup, so this is the quotient isogeny.
The standard Weierstrass invariant polynomials applied to
(B-RD41.1) give (B-RD41.2); substituting the Hesse-form coefficients
gives (B-RD41.3).
\end{proof}

\begin{proposition}[B-RD41: no level-twenty-seven failure with unipotent residual image]
\label{prop:B-RD41}
\status{Proved; independently checked by the root and obstruction agents, including the complete proof and exact certificate.}
Let $K$ be complete of mixed characteristic $(0,3)$, with algebraically
closed residue field and absolute ramification index two. No two
potentially supersingular elliptic curves with isomorphic $27$-torsion,
different prime-to-three special-inertia characters, and common
residual image $C_3$ or $C_6$ exist.

Thus any still-hypothetical height-two comparison in B.4 with
$p=3,e_K=2$ is restricted to the cyclic-four row:
\[
 \rho_3(G_K)=C_4,\qquad
 (|I_E|,|I_{E'}|)=(3,6),\qquad
 |I_{E,E'}|=18,\qquad
 v_K(j(E))=v_K(j(E'))=6.
                                                        \tag{B-RD41.5}
\]
Here $I_E,I_{E'}$ denote faithful special-inertia images, and the
pair is unordered. This proposition alone leaves least indices two
and three; Theorem~\ref{thm:B-RD42} subsequently excludes
(B-RD41.5) and proves that the least index is two.
\end{proposition}
\begin{proof}
If the common residual image is $C_6$, apply the same quadratic
twist to both curves to remove its central sign. The $27$-torsion
isomorphism persists, both $j$-invariants are unchanged, and
inequality of the two requested characters persists because both
are multiplied by the same quadratic character. We may therefore
assume the common image is the nontrivial unipotent group $C_3$.

By Proposition~\ref{prop:B-RD39}, the given comparison on nine-torsion
has the unit-parameter Hesse forms described there. After possibly
interchanging the curves, the first has parameter valuation $r=-1$
and $S=v_K(t^3-D)=-3$, and the second has $r\ge0$ and $S=0$.
The parameters $D$ and, in the reverse case, $D^{-1}$ are units.
The tame special-inertia orders are four and two.

Let $C\subset E[3]$ and $C'\subset E'[3]$ be their unique invariant
lines. These are the subgroups generated by the rational points
$(0,0)$ in the respective Hesse forms. Uniqueness follows because
a nontrivial unipotent matrix in dimension two has a unique invariant
line. The $27$-torsion isomorphism $f$ therefore carries $C$ onto $C'$.
Write $\varphi:E\to E_1=E/C$ and
$\varphi':E'\to E'_1=E'/C'$.

We spell out the level change. The inverse image of $E_1[9]$ under
$\varphi$ is
\[
 T=\{P\in E[27]:[9]P\in C\}.
\]
Indeed $[9]\varphi(P)=0$ is equivalent to $[9]P\in C$, which
also forces $[27]P=0$. Thus $E_1[9]=T/C$ as finite group schemes.
The same construction on the other side gives $T'/C'$, and $f$
induces an isomorphism
\[
                             E_1[9]\simeq E'_1[9]
                                                        \tag{B-RD41.6}
\]
over $K$. Both quotient curves are potentially supersingular.
An isogeny induces an isomorphism of rational Tate modules and of
their potential good-reduction Dieudonne modules. In particular the
special-inertia characters and their tame orders remain four and two.

Formula (B-RD41.3) gives a different pair of $j$-valuations. If
$r=-1$ and $D$ is a unit, then
$v_K(9t^3)=1$ and $v_K(8D)=0$, so
\[
                        v_K(j(E_1))=6.
\]
For the other curve, $r\ge0$ and $S=0$, while its unit parameter
is either $D$ or $D^{-1}$. Again $9t^3-8D$ (or its inverse-parameter
version) is a unit, so
\[
                       v_K(j(E'_1))=6+3r\ge6,
                                                        \tag{B-RD41.7}
\]
with $+\infty$ allowed for $t=0$.

Now classify the common residual representation in (B-RD41.6).
Since $\mu_3\subset K$, it lies in
$\operatorname{SL}_2(\mathbf F_3)$; the elementary inertia
classification in Proposition~\ref{prop:B-RD37} gives scalar image,
cyclic-four image, or cyclic-three/six image. A scalar image is
impossible by Proposition~\ref{prop:B-RD37}, since the characters
are different. Cyclic-four image is impossible by
Proposition~\ref{prop:B-RD38}, since one tame special-inertia order
is four. In the remaining case, a common quadratic twist gives
nontrivial unipotent residual image. It leaves both valuations in
(B-RD41.7) unchanged and preserves inequality of the characters.
Proposition~\ref{prop:B-RD39} would then require one $j$-valuation
to be three, contradicting (B-RD41.7). This excludes (B-RD41.6)
and hence the proposed original $27$-torsion comparison.

The reductions in Theorem~\ref{thm:B-RD34} and
Propositions~\ref{prop:B-RD35}, \ref{prop:B-RD37} identify every
remaining BT comparison with one of these elliptic residual cases.
Only the cyclic-four case survives, and
Proposition~\ref{prop:B-RD38} gives exactly (B-RD41.5).
\end{proof}

\par\smallskip\noindent\textbf{Exact certificate.}
The program \path{code/verify_c3_isogeny_exclusion.py} checks the
universal rational-map substitution and normalized differential,
both target invariant polynomials, the specialized quotient
$j$-formula, and all required valuation cases. It also models the
finite module passage from $27$-torsion to quotient nine-torsion,
with the actual invariant line and subgroup ranks. The retained
output records the author runtime $0.08725505400070688$ seconds.
It does not evaluate the remaining
cyclic-four antisymplectic branch.


\subsection{The reverse cyclic-four case and the exact smallest height-two index}

\begin{lemma}[The two cyclic-four wild families and their cubic fields]
\label{lem:B-RD42-normal}
\status{Proved. Complete root and independent original-source, proof and exact-certificate audits passed.}
Let $K$ be complete of mixed characteristic $(0,3)$, with algebraically
closed residue field and $v_K(3)=2$. Fix $\pi^2=3$ and $i^2=-1$ in $K$.
A potentially supersingular elliptic curve with irreducible cyclic-four
three-torsion representation and faithful special inertia of order six
is $K$-isomorphic to
\[
 M_b:\ y^2=x^3-3x+3b,\qquad b\in\OO_K^\times.
                                                        \tag{B-RD42.1}
\]
A curve with the same residual condition and faithful special inertia
of order three is the quadratic $\pi$-twist $M_b^{(\pi)}$ of such a curve.
The minimal discriminant valuations are respectively six and twelve.

Put $L_b=K(r)$, where $r^3-3r+3b=0$. This is a cyclic cubic extension,
and $M_b^{(\pi)}$ has good reduction over $L_b$, with faithful special
inertia of order three. The curve $M_b$ has good reduction over
$L_b(\sqrt\pi)$, with faithful special inertia of order six. Moreover
\[
 L_b=L_{b'}\quad\Longleftrightarrow\quad
                   \overline {b'}=\overline b\ \hbox{or}\ -\overline b.
                                                        \tag{B-RD42.2}
\]
\end{lemma}
\begin{proof}
The discriminant bound in Proposition~\ref{prop:B-RD36} and
Proposition~\ref{prop:B-RD38} give
$0\le\delta\le23$, $6\mid\delta$, and $v_K(j)=6$.
Thus $v_K(c_4)=(\delta+6)/3$. Passing to the short equation
$y^2=x^3+Ax+B$, with $A=-c_4/48$ and $B=-c_6/864$, gives an
integral equation when $\delta=6,12,18$. Indeed
$c_6^2=c_4^3-1728\Delta$ gives respectively
$v_K(c_6)\ge6,9,12$, so the pairs of coefficient valuations satisfy
\[
 (v_K(A),v_K(B))=(2,\ge0),\ (4,\ge3),\ (6,\ge6).
\]
The last equation scales down by $x=\pi^2X,y=\pi^3Y$, contradicting
minimality. The tame-order formula of Proposition~\ref{prop:B-RD36}
therefore forces $\delta=6$ for inertia order six and $\delta=12$
for order three; $\delta=0$ would mean good reduction.

In the first case write $A=3a$, $a$ a unit. The four three-torsion
lines have $x$-coordinates given by
\[
                     x^4+6ax^2+4Bx-3a^2=0.
                                                        \tag{B-RD42.3}
\]
If $v_K(B)=0$, its Newton polygon has one root of valuation two
and three of valuation zero. If $v_K(B)=1$, it has one root of
valuation one and three of valuation $1/3$. In each case the unique
root valuation gives a Galois-stable three-torsion line, contrary
to irreducibility of the cyclic-four representation. Hence
$v_K(B)\ge2$. If $v_K(B)\ge3$, the $\pi$-twist has coefficient
valuations at least $(4,6)$; scaling by $\pi$ makes its discriminant
a unit. The original curve would acquire good reduction over a
quadratic extension, contrary to its wild inertia of order three.
Thus $B=3b_0$ with $b_0$ a unit. A fourth root of $-a$ exists in $K$
by Hensel's lemma and normalizes $a$ to $-1$, giving (B-RD42.1).

For the order-three curve, twist its integral short equation by
$\pi$ and scale by $\pi$. The result has discriminant valuation six,
coefficient valuations $(2,\ge0)$, cyclic-four residual image, and
faithful special inertia of order six. (The unique nontrivial tame
quadratic character multiplies its central tame action by $-1$.)
The preceding argument applies. Twisting back proves the assertion.

For completeness the good fields can be exhibited directly.
The Newton polygon of $r^3-3r+3b$ gives $v_K(r)=2/3$, so its
root field has degree three. Its discriminant is
$27(4-9b^2)$, a square in $K$; it is therefore cyclic. All three
pairwise root differences have valuation one, since Galois acts
transitively on the three unordered pairs and the discriminant has
valuation six. The $\pi$-twist of (B-RD42.1) has equation
$y^2=x^3-9x+\pi^5b$. The substitution
$x=\pi^2X+\pi r$, $y=\pi^3Y$ gives
\[
              Y^2=X^3+\pi rX^2+(r^2-1)X.
                                                        \tag{B-RD42.4}
\]
Its reduction is $Y^2=X^3-X$, a good supersingular curve.
A nontrivial $\sigma\in\operatorname{Gal}(L_b/K)$ changes $X$ by
$(r-\sigma r)/\pi$, whose residue is nonzero, and fixes $Y$.
Thus its action on the special curve is nontrivial of order three.
Untwisting requires $\sqrt\pi$; the added quadratic action is
$[-1]$, proving both faithful inertia assertions. In particular
these displayed good fields are the minimal good fields.

One can also check the ramification directly: $r^2/\pi$ is a
uniformizer of $L_b$, and
\[
 v_{L_b}\bigl(\sigma(r^2/\pi)-r^2/\pi\bigr)=2.
\]
The cubic lower break is one. This observation is not needed for
the equality criterion below.

Choose the square root $v_b=\sqrt{1-9b^2/4}$ congruent to one, and put
\[
                 D_b=v_b+\frac{3i}{2}b.
                                                        \tag{B-RD42.5}
\]
If $u^3=iD_b$, then $u^3+u^{-3}=-3b$, so $u+u^{-1}$ is a root
of the cubic. A cube root of $i$ lies in $K$, and
\[
                    D_b\equiv1+\frac{3i}{2}b\pmod{\pi^3}.
                                                        \tag{B-RD42.6}
\]
This element is not a cube: its deviation from one has valuation
two, whereas a cube congruent to one has deviation at least three.
Consequently $K(\sqrt[3]{D_b})=L_b$. By
Lemma~\ref{lem:B-RD39-units}, every element of $1+\pi^3\OO_K$ is a
cube. The Kummer classes of $D_b$ and $D_{b'}$ therefore generate
the same line in $K^\times/(K^\times)^3$ exactly when
$\overline{b'}=\pm\overline b$. This proves (B-RD42.2).
\end{proof}

\begin{lemma}[The reverse marked cocycle identifies the cubic good fields]
\label{lem:B-RD42-reverse}
\status{Proved. Complete root and independent original-source, proof and exact-certificate audits passed.}
Suppose that $M_b[9]\simeq M_{b'}^{(\pi)}[9]$ is antisymplectic,
and that their common residual representation is irreducible cyclic
of order four. Then $L_b=L_{b'}$.
\end{lemma}
\begin{proof}
Let $L=K(M_b[3])$ and let $F$ be its quadratic subfield. The residual
isomorphism gives the same fields on the other side. Choose symplectic
three-torsion markings with a fixed Weil-pairing value $\zeta$. Since
the given isomorphism reverses the pairing, its matrix can be taken
to be $R=\operatorname{diag}(1,-1)$ between these markings. Choose
the cyclic-four generator as
$S=\left(\begin{smallmatrix}0&-1\\1&0\end{smallmatrix}\right)$
on the first side. On the second side its matrix is
$RSR^{-1}=S^{-1}=-S$. This is precisely Fisher's reverse-marking
involution $\iota=\operatorname{diag}(1,-1)$; his inverse-transpose
twist dictionary is \citet[Section~3, p.~7, and Lemma~3.2, pp.~7--8]
{Fisher2015Nine}.

The projective action on the two Hesse parameters is the same
involution $h(t)=(t+2)/(t-1)$. Write $z=t-1$, $z'=t'-1$ and
$u=z/\pi$, $u'=z'/\pi$. As in Proposition~\ref{prop:B-RD40},
\[
 N_{F/K}(z)=N_{F/K}(z')=3,\qquad
 N_{F/K}(u)=N_{F/K}(u')=1,\qquad
 \overline u,\overline{u'}\in\{1,-1\}.
                                                        \tag{B-RD42.7}
\]
The actual Fourier lift for $S$ has differential multiplier
$\kappa(t)=c/(t-1)$, where $c=\zeta-\zeta^2$. The lift for $S^{-1}$
is its central negative: its multiplier is $-\kappa(t)$.
Indeed $S^{-1}=-S$, the central element is elliptic inversion,
and the exact normalization and identity
$\kappa(t)\kappa(h(t))=-1$ were proved in
Proposition~\ref{prop:B-RD40}. Thus no change of plane-residue
normalization or parameter-dependent factor is being suppressed.

Let $\alpha,\alpha'$ compare the respective universal Hesse
invariant differentials with invariant differentials over $K$.
The analogue of (B-RD40.7), now with opposite actual lifts, is
\[
       \frac{\tau(\alpha'/\alpha)}{\alpha'/\alpha}
                 =-\frac{z'}z\quad\hbox{or its inverse},
                                                        \tag{B-RD42.8}
\]
for a generator $\tau$ of $\operatorname{Gal}(L/K)$.
The discriminant valuations are six and twelve, so the invariant
section $\Delta\omega^{12}$ gives
$v_L(\alpha')-v_L(\alpha)=2$. Tame leading coefficients therefore
make the residue of the left side equal to $-1$. It follows that
\[
                            \overline{u'}=\overline u.
                                                        \tag{B-RD42.9}
\]

Over $F$ twist both curves by the same quadratic character of $L/F$.
Their three-torsion becomes rational, the Hesse curves are $E_t,E_{t'}$,
and the reverse nine-congruence persists. Fisher's reverse equation
(11), with its essential numerator three, gives at $\theta=1$
\[
                          u'u=\frac{(t'-1)(t-1)}3
                                    \in(F^\times)^3.
                                                        \tag{B-RD42.10}
\]
Here we use exactly the full-level marking convention of
\citet[Section~2, pp.~4--5; Theorem~4.3 and its proof, pp.~9--10]
{Fisher2015Nine}. By (B-RD42.9), this cube is congruent to one.
Since $v_F(3)=4$, expansion of $(1+w)^3$ shows
$v_F(u'u-1)\ge3$. Set
\[
 T=z+3/z=\pi(u+u^{-1}),\qquad
 T'=z'+3/z'=\pi(u'+u'^{-1})\in K.
\]
The exact identity
\[
       T'-T=\pi(u'-u^{-1})\left(1-\frac u{u'}\right)
                                                        \tag{B-RD42.11}
\]
and (B-RD42.9) give $v_F(T'-T)\ge2+3+1=6$, hence
\[
                              v_K(T'-T)\ge3.
                                                        \tag{B-RD42.12}
\]

We recover the residues of the normal-form parameters from these
traces. Eliminating $z^2-Tz+3=0$ from the Hesse $j$-formula gives
\[
 j(E_t)=27\left(\frac{T(T+4)}{T+3}\right)^3,
 \qquad
 j(M_b)=\frac{6912}{4-9b^2}.
                                                        \tag{B-RD42.13}
\]
These are exact rational identities, and twisting does not change
$j$. Put $H=T(T+4)/(T+3)$. If $\overline u=\epsilon\in\{1,-1\}$,
then the identity
\[
 u+u^{-1}-2\epsilon=\frac{(u-\epsilon)^2}{u}
\]
and $v_F(\pi)=2$ give $v_F(T-2\epsilon\pi)\ge4$. Since
$T\in K$, we may write $T=2\epsilon\pi+\pi^2a$ with
$a\in\OO_K$. Also
\[
 H-4=\frac{T^2-12}{T+3}\in\pi^2\OO_K,\qquad
 H^3=\frac{64}{1-9b^2/4}.
\]
Writing $H=4+\pi^2d$ and comparing the latter identity modulo
$\pi^5$ gives $d\equiv b^2\pmod\pi$. The former identity gives
$d\equiv2a\pmod\pi$. Consequently
\[
                T\equiv2\epsilon\pi+\frac{\pi^2}{2}b^2
                                      \pmod{\pi^3}.
                                                        \tag{B-RD42.14}
\]
The same formula holds for $T'$ and $b'$ with the same $\epsilon$
by (B-RD42.9). Equation (B-RD42.12) therefore gives
$\overline{b'}^{\,2}=\overline b^{\,2}$. In characteristic three
this says $\overline{b'}=\pm\overline b$, and
Lemma~\ref{lem:B-RD42-normal} proves $L_b=L_{b'}$.
\end{proof}

\begin{theorem}[B-RD42: the exact least height-two index at $p=3,e=2$]
\label{thm:B-RD42}
\status{Proved. Complete root and independent original-source, proof and exact-certificate audits passed.}
For the full class of height-two groups in B.4 satisfying the manuscript
monodromy hypothesis,
\[
                              \boxed{\nu_{\min}(3,2,2)=2}.
                                                        \tag{B-RD42.15}
\]
Equivalently, $27$-torsion determines the requested character on
prime-to-three inertia elements, while nine-torsion does not always do so.
This is a sharp small-height case; no general least-index formula is
asserted for other heights or ramification indices.
\end{theorem}
\begin{proof}
The lower bound is the explicit nine-torsion example of
Proposition~\ref{prop:B-RD29-p3-level-two}. The reduction and bounds in
Theorem~\ref{thm:B-RD34} and Proposition~\ref{prop:B-RD35} reduce
any putative failure at $27$ to two potentially supersingular elliptic
curves after completed maximal unramified base change. These steps
already cover arbitrary perfect residue fields and all the allowed
height-two monodromy filtrations. The completed maximal unramified
base change preserves inertia and absolute ramification; its algebraically
closed residue field makes every unit a square, so $\pi^2=3$ and
$\mu_3\subset K$ can be arranged without a further ramified extension.

Propositions~\ref{prop:B-RD37} and \ref{prop:B-RD41} exclude the
scalar and unipotent residual cases. By
Propositions~\ref{prop:B-RD38} and \ref{prop:B-RD40}, the sole remaining
case is an antisymplectic comparison with cyclic-four residual image,
individual faithful special-inertia orders three and six, and joint
faithful inertia order eighteen. Lemma~\ref{lem:B-RD42-normal}
writes this pair, after interchanging it, as $M_b,M_{b'}^{(\pi)}$.
Restriction of the normalized antisymplectic $27$-isomorphism to
nine-torsion satisfies Lemma~\ref{lem:B-RD42-reverse}; hence
$L_b=L_{b'}$. Both curves then have good reduction over the same
extension $L_b(\sqrt\pi)$ of degree six. Their joint faithful special
inertia has order at most six, contradicting the required order
eighteen. This excludes the last case and proves (B-RD42.15).
\end{proof}

\par\smallskip\noindent\textbf{Exact certificate and scope.}
The program \path{code/verify_c4_final_exclusion.py} checks the universal
normal-form and good-model invariants, the Cardano and Kummer identities,
the trace elimination and truncated expansions, the opposite marked
matrices, and the integer valuation cases used in the proof. Its output
is retained separately, with author runtime $0.1415409550027107$ seconds.
The initial structural-equality check and traceback are retained in
\path{failed_candidates/B42_certificate_structural_equality_v1.py} and
its adjacent output; the corrected certificate tests exact polynomial
differences. The preceding proof uses the algebraic-closure
residue hypothesis only in the established reduction and the explicitly
stated Hensel and unit-cube steps. The finite-elimination procedure of
B-RD24 remains unevaluated; it is not used to obtain this exact index.

\par\smallskip\noindent\textbf{Independent audits.}
The complete proof and exact certificate were independently checked by
two reviewers. Their author-output-preserving reruns passed in
$0.11300166400178568$ and $0.10002119700220646$ seconds, respectively.
The reverse-marking convention was checked against Fisher's original
Section~3 and the full proof of Theorem~4.3. This establishes only the
explicit small instance (B-RD42.15); the general least-index problem
in B.4 remains open.


\subsection{An explicit eight-torsion ambiguity in the unramified dyadic case}

\begin{proposition}[B-RD43: even eight-torsion does not suffice at $p=2,e=1$]
\label{prop:B-RD43}
\status{Proved. Complete root and independent original-source, proof and exact-certificate audits passed.}
Let $K=\mathbf Q_2(\omega)$, where $\omega^2+\omega+1=0$, and put
$\lambda=1+2\omega$, so $\lambda^2=-3$. Define
\[
 E:\ y^2+(5-\lambda)xy-4\lambda y=x^3-2\lambda x^2,
 \qquad P=(0,0),\qquad B=E/\langle P\rangle,
 \qquad q=-\frac{4\lambda}{3}.
                                                        \tag{B-RD43.1}
\]
Then $P$ has exact order eight, $v_K(q)=2>0$, and the split Tate
curve $T_q$ satisfies
\[
                             B[8]\simeq T_q[8].
                                                        \tag{B-RD43.2}
\]
The requested characters of $B[2^\infty]$ and $T_q[2^\infty]$
differ on an inertia element whose image has order three in the
finite special-inertia quotient. Consequently
\[
                       \nu_{\min}(2,1,h)\ge3\qquad(h\ge2).
                                                        \tag{B-RD43.3}
\]
The same lower bound holds for every absolute ramification index
$e$ not divisible by three. No assertion of sufficiency of
sixteen-torsion is made.
\end{proposition}
\begin{proof}
The field $K$ is unramified quadratic over $\mathbf Q_2$, since
$X^2+X+1$ is separable and irreducible over $\mathbf F_2$.
The displayed equation is integral. Exact addition gives
\[
\begin{array}{c|c|c}
 n&x(nP)&y(nP)\\\hline
 1&0&0\\
 2&2\lambda&-6-6\lambda\\
 3&-6+2\lambda&12\\
 4&(-3+\lambda)/2&3\\
 5&-6+2\lambda&12-12\lambda\\
 6&2\lambda&0\\
 7&0&4\lambda
\end{array}
\qquad 8P=0.
                                                        \tag{B-RD43.4}
\]
In particular $P$ has exact order eight. This curve arises from
the $X_1(8)$ Tate normal form in
\citet[Appendix~A.2, pp.~26--27, Figure~8]{AntieauAuel2021Descent}
by taking its parameter to be $\lambda$ and scaling
$x=X/4$, $y=Y/8$. The displayed addition certificate verifies
its torsion without relying on that parametrization.

We compute the exact extension class in
Lemma~\ref{lem:B-RD29-isogeny-Tate}, so that the Tate parameter
in (B-RD43.1) is specified rather than merely asserted to exist.
The tangent at $P$ is $y=0$. The tangent at $2P$ is
\[
 \ell_2=y-\frac{-9+\lambda}{2}x-(-3+3\lambda)=0.
\]
Put $x_2=2\lambda$ and $x_4=(-3+\lambda)/2$.
The Miller doubling identity gives a function with divisor
$8[P]-8[0]$:
\[
 f=\frac{y^4\ell_2^2}{(x-x_2)^4(x-x_4)}.
\]
Here the last tangent is the vertical line at the two-torsion point
$4P$. Reducing by the equation of $E$ cancels the denominator and gives
\[
\begin{split}
 f(x,y)={}&x^4-4\lambda x^3-8x^2
            +(2\lambda-6)x^2y\\
          &+(12\lambda+4)xy+16y.
\end{split}                                             \tag{B-RD43.5}
\]
This polynomial identity and all the displayed point additions
are verified exactly in $\mathbf Q(\lambda)$.
The divisor description of the dual-isogeny connecting homomorphism
is \citet[Appendix~A.1, Lemmas~A.1--A.2, p.~22]
{AntieauAuel2021Descent}. Its unspecified normalizing constant
cancels in the ratio at $3P$ and $2P$, since both avoid the zero
and pole of $f$. Thus
\[
 \delta(P)=\frac{f(3P)}{f(2P)}
           =\frac{576\lambda}{-432}
           =-\frac{4\lambda}{3}
                   \quad\hbox{in }K^\times/(K^\times)^8.
                                                        \tag{B-RD43.6}
\]
The quotient $B=E/\langle P\rangle$ therefore fits into the
Kummer extension
$0\to\mu_8\to B[8]\to\mathbf Z/8\to0$ with the indicated
class. Since $v_K(q)=2$, Lemma~\ref{lem:B-RD29-isogeny-Tate}
and Tate's explicit finite-order construction identify it with
$T_q[8]$. As an additional check, the Figure~8 representative
$\lambda^7(\lambda-1)^6(\lambda+1)^4$ divided by (B-RD43.6)
is exactly $(-3-\lambda)^8$.

The invariant calculation gives
\[
 c_4=-200-312\lambda,\qquad
 c_6=15488+576\lambda,\qquad
 \Delta=20736(\lambda-2),
\]
so
\[
                     v_K(c_4)=4,\quad v_K(\Delta)=8,
                     \quad v_K(j)=4.
                                                        \tag{B-RD43.7}
\]
For example, $c_4=-16(32+39\omega)$ and
$\lambda-2=-1+2\omega$ are respectively sixteen times a unit
and a unit. The displayed discriminant is minimal, since eight
is less than twelve and integral discriminant valuations are
nonnegative. Integral $j$ gives potential good reduction by
\citet[Section~2, p.~496]{SerreTate1968Good}; its positive valuation
gives supersingular special reduction in characteristic two,
as in the already verified proof of Proposition~\ref{prop:B-RD22}.

Over the completed maximal unramified field, let $m$ be a minimal
good-reduction degree. The change to a unit-discriminant equation
implies $12\mid8m$, hence $3\mid m$. Faithful special inertia
therefore contains an element of order three. Its trace on the
rank-two rational Dieudonne module is $-1$, as in
Proposition~\ref{prop:B-RD22}. The isogeny $E\to B$ preserves this
rational module and its finite inertia action. The Tate group's
two good rank-one graded pieces have trace $2$ on that element.
The order-three assertion concerns the finite special-inertia image;
no finite-order lift in the absolute inertia group is asserted.
Both groups satisfy the manuscript monodromy condition, and the
characters differ despite (B-RD43.2).

Finally, take an Eisenstein extension of $K$ of degree $e$.
The torsion isomorphism persists and the raw discriminant valuation
becomes $8e$. If $3\nmid e$, the relation $12\mid8em$ still
forces $3\mid m$, so the same character discrepancy survives.
Adding $h-2$ identical constant etale factors proves every height
$h\ge2$. Since eight-torsion is $M(2)$, its failure proves
(B-RD43.3).
\end{proof}

\par\smallskip\noindent\textbf{Exact certificate and scope.}
The program \path{code/verify_p2_character_level3.py} checks all eight
subgroup points and all $64$ additions, the exact invariant
polynomials, both tangent lines, the affine Miller-function identity,
its two evaluations and the eighth-power parameter comparison.
Its retained author output records $0.5159330129972659$ seconds.
The finite-scheme extension dictionary and the special-inertia
argument are proved above using the stated primary sources.
This counterexample raises the lower bound; it does not compute
the least dyadic index.

Independent reruns took $0.5299817019986222$ and
$0.6142895730008604$ seconds, preserving the author output. Both reviewers
checked the complete proof and the original descent and Tate sources.


\subsection{The dyadic sixteen-torsion problem reduces to two supersingular curves}

\begin{lemma}[Finite dyadic inertia and integral elliptic conjugacy]
\label{lem:B-RD44-conjugacy}
\status{Proved; complete independent parent and obstructions/crystals proof, source and exact-code audits passed.}
Let $k$ be algebraically closed of characteristic two and let
$G_0=E_0[2^\infty]$, where $E_0:y^2+y=x^3$.
Every finite subgroup of $\operatorname{Aut}(G_0)$ is conjugate,
by an element of $\operatorname{Aut}(G_0)$ itself, into
$\operatorname{Aut}(E_0)$. The possible groups are
\[
 1,\ C_2,\ C_3,\ C_4,\ C_6,\ Q_8,\ \operatorname{SL}_2(\mathbf F_3).
 \tag{B-RD44.1}
\]
Their reduced norms are all one.
\end{lemma}
\begin{proof}
The height-two Dieudonne calculation of Lemma~\ref{lem:B-RD31-finite},
with $3$ replaced by $2$, identifies the rational endomorphism algebra
with the quaternion division algebra $D/\mathbf Q_2$, and
$\operatorname{End}(G_0)$ with its valuation ring $\mathcal O_D$.
Its residue field is $\mathbf F_4$ and the reduction kernel on units
is pro-$2$. Normalize $v_D$ by $v_D(2)=2$.
For a finite subgroup $H$, its reduction kernel $P$ is consequently
a normal Sylow two-subgroup, and $H/P$ has order one or three.

Every element of $D$ satisfies a quadratic polynomial over $\mathbf Q_2$.
Thus no element of $P$ has order eight, since
$[\mathbf Q_2(\zeta_8):\mathbf Q_2]=4$; every element has order at
most four. The only involution in a division ring of characteristic
zero is $-1$, by $(x-1)(x+1)=0$.
If two order-four elements $i,j$ generate different cyclic subgroups,
then $ij$ is neither $1$ nor $-1$, so $(ij)^2=-1$ and $ij=-ji$.
They generate $Q_8$. Any further order-four element outside this
quaternion group would anticommute with both $i,j$. In the basis
$1,i,j,ij$ its simultaneous anticommutant is $\mathbf Q_2ij$;
squaring then forces the element to be $ij$ or $-ij$, a contradiction.
This proves $P=1,C_2,C_4$, or $Q_8$.

If $H/P$ has order three, Cauchy's theorem gives an element $z$ of
order three and $H=P\rtimes\langle z\rangle$.
For $P=1,C_2$ this gives $C_3,C_6$.
For $P=C_4$ the action is trivial, since
$\operatorname{Aut}(C_4)=C_2$; this would give an element of order
twelve, impossible because $\mathbf Q_2(\zeta_{12})$ has degree four.
For $P=Q_8$ the action must be nontrivial for the same reason.
It cyclically permutes the three order-four subgroups, giving the
binary tetrahedral group $Q_8\rtimes C_3=\operatorname{SL}_2(\mathbf F_3)$.
This recovers the degree-two case of Hewett's general classification
\cite[Theorem~1.4 and Corollary~1.5, p.~520]{Hewett1995Division}.
The argument above proves the particular classification used here
without importing a conjugacy assertion from that theorem.

We now prove integral conjugacy. Choose $\omega\in\mathbf F_4$
with $\omega^2+\omega+1=0$. The maps
\[
 (x,y)\longmapsto
 (u^2x+s^2,\ y+u^2sx+t),\qquad
 u\in\mathbf F_4^\times,\ s\in\mathbf F_4,\ t^2+t=s^3,
 \tag{B-RD44.2}
\]
are twenty-four distinct automorphisms fixing the origin.
Substitution verifies their equations and closure under composition.
The eight with $u=1$ form $Q_8$: the six with $s\ne0$ square
to $(x,y)\mapsto(x,y+1)=[-1]$.
Choose generators $i,j$ of this quaternion subgroup.
Its rational span is $D$, with $i^2=j^2=-1$ and $ij=-ji$.
The eight order-three automorphisms in (B-RD44.2) induce the eight
order-three automorphisms of $Q_8$; this follows either by composing
the displayed maps, or by choosing cyclic images of the three axes
and the two independent signs. Denote this group of twenty-four
automorphisms by $H_0$.

An order-four element $a$ can be conjugated to $i$ by $a+i$,
unless $a=-i$, which already generates the desired subgroup:
$(a+i)a=i(a+i)$.
The centralizer of $i$ contains $1+i$, of norm two and valuation one.
Multiplication of this conjugator on the left by any power of $1+i$
does not change its conjugation of $a$, and makes its valuation zero.
Thus the cyclic order-four conjugacy is integral.

Put $\zeta=(-1+i+j+ij)/2$.
Its conjugation cyclically permutes the quaternion axes and has
order three. Choose the order-three member of $H_0$ inducing that
same automorphism of $Q_8$. Their quotient is central in $D$,
and its cube is one, so it is one; thus $\zeta\in H_0$.
For any order-three $z$, the element $d=z+\zeta+1$ satisfies
$dz=\zeta d$. If $d=0$, then $z=\zeta^{-1}$ and its cyclic
subgroup is already the required one. Otherwise $d$ is invertible.
The element $i-j$ has norm two, valuation one, and conjugates
$\zeta$ to $\zeta^{-1}$. Hence multiplying $d$ on the left by
a suitable integral power of $i-j$ makes it a unit while still
conjugating $\langle z\rangle$ to $\langle\zeta\rangle$.
This also treats $C_6$, since the involution is $-1$.

For a quaternion subgroup, first put one generator equal to $i$
by the preceding integral conjugacy. Its other generator is $cj$,
where $c\in\mathbf Q_2(i)$ and $c\bar c=1$.
If $c\ne-1$, take $b=1+c$; if $c=-1$, take $b=i$.
In both cases $b/\bar b=c$, so $b^{-1}(cj)b=j$ while $b$ commutes
with $i$. The element $1+i$ normalizes the target $Q_8$, sending
$j$ to $ij$, and again has valuation one. Left multiplication
of $b^{-1}$ by a suitable power of $1+i$ therefore makes the
conjugator integral without changing the target subgroup.

Finally put the quaternion subgroup of an order-twenty-four $H$
inside $H_0$ by this integral conjugacy. An element $z\in H$ of
order three induces the same automorphism of $Q_8$ as some
order-three $z_0\in H_0$. Then $zz_0^{-1}$ centralizes the
rational span $D$ of $Q_8$, so is a scalar in $\mathbf Q_2^\times$.
Its cube is one, forcing it to be one. Hence $H\subset H_0$.
The trivial and order-two cases require no conjugacy.
All correcting elements have valuation zero at the end and
therefore belong to $\mathcal O_D^\times$, as required.
Finally order-four generators have minimal polynomial $T^2+1$,
and order-three generators have $T^2+T+1$; their norms, and that
of $-1$, are one. They generate every group in (B-RD44.1).
\end{proof}

\begin{proposition}[Exact dyadic elliptic realization]
\label{prop:B-RD44-realization}
\status{Proved; complete independent parent and obstructions/crystals proof, source and exact-code audits passed.}
Let $K$ be a complete discretely valued field of mixed characteristic
$(0,2)$ with algebraically closed residue field. Every height-two
$2$-divisible group $M/K$ which has potentially good one-dimensional
supersingular reduction is isomorphic to $E[2^\infty]$ for an
elliptic curve $E/K$. The isomorphism preserves the full integral
Tate lattice and the labelled finite special-inertia action.
\end{proposition}
\begin{proof}
Choose a finite Galois good extension $L/K$ and the good model
$\mathcal G/\mathcal O_L$. Tate's integral full faithfulness
extends its generic semilinear descent isomorphisms and their
inverses uniquely, with their cocycle relations
\cite[Theorem~4 and Corollaries~1--2, p.~181]{Tate1967PDivisible}.
Lazard's height classification
\cite[Theorem~IV and proof, pp.~270--271]{Lazard1955Formal}
identifies the special formal group with $G_0$.
Lemma~\ref{lem:B-RD44-conjugacy} makes the entire labelled finite
action an action by actual automorphisms of $E_0$.

The labelled Serre--Tate construction and algebraization in the
proof of Proposition~\ref{prop:B-RD35} now apply with $2$ in place
of $3$: their deformation bases are $\mathcal O_L/\pi^n$, on
which $2$ is nilpotent, and their precise inputs impose no
odd-prime restriction. Namely Katz's Theorem~1.2.1
\cite[p.~143 and proof pp.~143--146]{Katz1981SerreTate}
produces compatible elliptic deformations with the chosen BT
isomorphisms and semilinear action. The compatible ample line
bundles $\mathcal O(3\cdot0)$ and EGA~III$_1$, Theorems~5.4.1
and~5.4.5, algebraize the elliptic scheme and its operations
\cite[pp.~156--158]{EGAIII1}.
For each $r$, completeness of the finite free Hopf algebras
algebraizes the isomorphism on $2^r$-torsion, compatibly in $r$.
Serre--Tate uniqueness and proper formal full faithfulness lift
and algebraize the labelled descent datum and all cocycles.
Effective generic polarized Galois descent yields $E/K$ and
$E[2^\infty]\simeq M$. As in B-RD35, only the generic elliptic
curve is descended across the ramified extension; its good
integral model is not asserted to descend.
\end{proof}

\begin{lemma}[No rational sixteen-point in the unramified supersingular case]
\label{lem:B-RD44-point}
\status{Proved; complete independent parent and obstructions/crystals proof, source and exact-code audits passed.}
Let $K$ have perfect residue field of characteristic two and
absolute ramification index one. If an elliptic curve $E/K$ has
potentially supersingular good reduction, then
$E(K)[2^\infty]$ has exponent at most eight.
In particular it has no point of exact order sixteen.
\end{lemma}
\begin{proof}
Pass to the completed maximal unramified extension if necessary;
this preserves the absolute ramification and any asserted torsion
point. Choose a minimal integral Weierstrass equation.
Its coefficient $a_1$ lies in $2\mathcal O_K$.
For good supersingular reduction this is the vanishing of the
characteristic-two Hasse coefficient. For bad reduction, potential
goodness gives $v(j)\ge0$ and $v(\Delta)>0$, whence $v(c_4)>0$;
the identity $c_4\equiv a_1^4\pmod2$ proves the same assertion.

The formal parameter $T=-x/y$ identifies the reduction kernel
$E_1(K)$ with the maximal ideal and has integral multiplication
series
\[
             [2](T)=2T-a_1T^2+\sum_{j\ge3}b_jT^j.
 \tag{B-RD44.3}
\]
These integral formal-group assertions are Tate's Theorems~4.1--4.2
and the expansions in his Section~4
\cite[pp.~41--44]{Tate1975Algorithm}.
If $v(T)\ge1$, the first term of (B-RD44.3) has value $1+v(T)$,
strictly less than $1+2v(T)$ and than every $jv(T)$ for $j\ge3$.
Consequently $[2]T\ne0$ for every nonzero formal point, and
$E_1(K)[2^\infty]=0$.

If reduction is good supersingular, its special curve has no
nonidentity geometric two-primary torsion. Every rational
two-primary point lies in $E_1(K)$ and is therefore zero.
Otherwise reduction is additive: potential good reduction excludes
the multiplicative case. The smooth connected reduction group is
$\mathbf G_a$, so $E_0(K)/E_1(K)$ has exponent two. The geometric
component group's two-primary subgroup has exponent at most four.
Both assertions hold in residue characteristic two: in Tate's
Section~6 table, p.~46, these are the component-group and
identity-component rows above the explicit boundary marking the
later assumptions $\operatorname{char}k\ne2,3$
\cite[Section~6, pp.~45--46]{Tate1975Algorithm}.
Thus any rational two-primary point $P$ has $[4]P\in E_0(K)$
and $[8]P\in E_1(K)$. The latter group is two-primary torsion-free,
so $[8]P=0$.
\end{proof}

\begin{proposition}[B-RD44: the mixed case is excluded by sixteen-torsion]
\label{prop:B-RD44}
\status{Proved; complete independent parent and obstructions/crystals proof, source and exact-code audits passed.}
Suppose $p=2$, $e_K=1$, and $M,N$ have height two and satisfy
the manuscript monodromy condition. If
$M[16]\simeq N[16]$ but their prime-to-two special-inertia
characters differ, then both groups themselves have potentially
good supersingular reduction. After completed unramified base
change they are the exact BT groups of two elliptic curves.
Their faithful finite inertia images have respective orders
\[
       m_1\in\{1,2,4,8\},\qquad m_2\in\{3,6,24\},
 \tag{B-RD44.4}
\]
in some order; their joint good-reduction degree belongs to
$\{3,6,12,24,48,96,192\}$.
No assertion of sixteen-torsion sufficiency for these remaining
pairs is made.
\end{proposition}
\begin{proof}
Work after completed maximal unramified base change. Write $d_M,d_N$
for the sums of the dimensions of the good graded models.
They belong to $\{0,1,2\}$. The determinant characters are
$\chi_2^{d_M}\epsilon_M,\chi_2^{d_N}\epsilon_N$, where
$\epsilon_M,\epsilon_N$ are finite quadratic characters, by the
height-one determinant classification already used in B-RD18 and
B-RD31. The cyclotomic image modulo sixteen is all
$(\mathbf Z/16)^\times$: the polynomial $\Phi_{16}(1+T)$ is
Eisenstein over $\mathcal O_K$ because $v_K(2)=1$, so adjoining
$\zeta_{16}$ still has degree eight, also for a transcendental
perfect residue field.
Equality of determinants modulo sixteen forces $d_M=d_N$:
if their difference is $\pm1$, its cyclotomic character has
order four and cannot be a sign character; if it is $\pm2$,
an inertia element with cyclotomic value $3$ has squared value
$9$, which is neither $1$ nor $-1$ modulo sixteen.

For common dimension zero, all good graded groups are etale,
and the prime-to-two Brauer character is read directly from the
semisimplified two-torsion. Dimension two reduces to this by
Cartier duality. For common dimension one, each group either
has a filtration by height-one pieces, or is itself potentially
good supersingular. A potentially ordinary good graded piece
has its canonical connected--etale filtration, which descends
uniquely and gives the first alternative.
Height-one finite characters have values in
$(\mathbf Q_2^\times)_{\mathrm{tors}}=\{1,-1\}$; hence every
height-one-filtered group has prime-to-two special character two.

It remains to exclude a comparison between a potentially
supersingular $M$ and a height-one-filtered $N$.
Proposition~\ref{prop:B-RD44-realization} gives $M=E[2^\infty]$
with its exact Tate lattice. The first height-one subgroup of
$N[16]$ transports to a cyclic direct summand $C\subset E[16]$
with character $\chi_2^b\epsilon$, where $b=0$ or $1$ and
$\epsilon$ is quadratic. If $b=0$, the quadratic twist by
$\epsilon$ has a rational point of order sixteen.
If $b=1$, use the elliptic quotient $E'=E/C$.
The image of $E[16]$ is the cyclic subgroup $E[16]/C$ of
$E'[16]$; the Weil pairing, or the determinant $\det E[16]=\chi_2$,
gives this quotient the character $\epsilon^{-1}$.
The corresponding quadratic twist of $E'$ again has a rational
point of order sixteen. Isogenies and quadratic twists preserve
potentially supersingular reduction. Both conclusions contradict
Lemma~\ref{lem:B-RD44-point}.

Both remaining groups are therefore potentially supersingular,
and Proposition~\ref{prop:B-RD44-realization} applies to each.
In the list (B-RD44.1), every odd-order element is either the
identity, of trace two, or has order three, of trace minus one.
If both faithful inertia orders are divisible by three, their
characters agree as functions on prime-to-two inertia: the tame
procyclic quotient has a unique quotient of order three, so the
two nontrivial cubic quotients have the same kernel. If neither
order is divisible by three, both characters are two. Different
characters therefore require exactly one order divisible by three,
which gives (B-RD44.4). The joint wild image embeds in the product
of two groups of orders at most eight, and its tame quotient has
order three. Its order is consequently $3\cdot2^r$, $0\le r\le6$.
The faithful-image descent of Lemma~\ref{lem:B-RD17-degree}
makes this a simultaneous good-reduction extension.
\end{proof}

\begin{remark}[Scope and finite verification]\label{auto:B:e906eee05992}
\status{Proved under stated hypotheses}

Proposition~\ref{prop:B-RD43} proves that eight-torsion can fail
at $(p,e,h)=(2,1,2)$. The preceding result removes the mixed
case at the next level, but leaves the all-supersingular
sixteen-torsion comparison open. No finite elimination has been
executed for those pairs.
The certificate \path{code/verify_p2_height_two_mixed_gate.py}
checks the twenty-four special automorphisms, their group orders
and quaternion relations, the explicit conjugator and normalizer
identities, the formal valuation inequalities, and the determinant
dimension separation modulo sixteen. Its exact scope and runtime
are retained in the adjacent JSON output. The author run passed
in $0.07470255399675807$ seconds; these computations do
not decide the remaining pairs.
\end{remark}

\par\smallskip\noindent\textbf{Independent verification.}
The root rerun took
$0.06556921700394014$ seconds (actual audit record
2026-11-19 at 13:57:27.954980 UTC); the obstructions/crystals
agent reran it in $0.06939755400526337$ seconds.
Both independently read the full proof and the specified original
sources. The author output was preserved. These assertions exclude
the mixed sixteen-torsion comparison and give an exact elliptic
realization; they do not decide the remaining comparison between
two potentially supersingular objects.


\subsection{A bounded list of dyadic supersingular comparison levels}

\begin{lemma}[Dyadic graph and logarithm estimates]
\label{lem:B-RD45}
\status{Proved; complete independent parent and obstructions/crystals proof, original-source and exact-code audits passed.}
Let $K$ be a complete discretely valued field of mixed characteristic
$(0,2)$ with algebraically closed residue field.
Suppose $M,N$ are height-two potentially supersingular groups with
different prime-to-two special-inertia characters and
$M[2^n]\simeq N[2^n]$. Let $L/K$ be a common finite Galois
good-reduction extension, put $R=\mathcal O_L$ and $a=v_L(2)$,
and write $M_0,N_0$ for the special good groups.
Define
\[
 q_0=1,\qquad q_{r+1}=\min\{a+q_r,1+2q_r,4q_r\}.
 \tag{B-RD45.1}
\]
Then $q_{n-1}\le a$.
More precisely, the special graph at every level $2^r$, $r\le n$,
is $\ker F_{M_0}^r\times\ker F_{N_0}^r$; if $q_u>a$, then
$u\ge n$.
\end{lemma}
\begin{proof}
We first supply the dyadic replacement for the tame idempotent
in Proposition~\ref{prop:B-RD33}.
By Lemma~\ref{lem:B-RD44-conjugacy}, exactly one of the two
faithful finite inertia orders is divisible by three.
Choose an element $\tau$ of order three in the joint finite image.
On one special group it is the identity; on the other it has
minimal polynomial $T^2+T+1$ in the quaternion endomorphism algebra.
Consequently
\[
                       e_0=(1+\tau+\tau^2)/3
 \tag{B-RD45.2}
\]
is the integral projection onto the first factor. Division by
three is valid on every two-primary group scheme.

By Proposition~\ref{prop:B-RD44-realization} the generic groups
are exact elliptic BT groups and have perfect alternating
pairings into $\mu_{2^\infty}$. Alternatively their finite
determinants are one and the determinant argument in B-RD33
gives the same pairings. Tate's integral full faithfulness
extends each pairing and its inverse to the good models.
At level $2^r$, rescale one factor by the unit multiplier of
the asserted rank-two isomorphism. Its graph is then maximal
isotropic. Its finite flat closure $H_r$ remains maximal
isotropic after reduction: restricted pairings extend by
flatness, finite-flat duality commutes with reduction, and
the closed inclusion into its orthogonal has equal free rank.
These are the exact arguments of B-RD33, with no division by two.

The idempotent (B-RD45.2) preserves $(H_r)_k$ and splits it
as a product of subgroups of $M_0,N_0$.
In a one-dimensional connected formal group the unique subgroup
of rank $2^b$ is $\ker F^b$, since its defining ideal in
$k[[T]]$ must be $(T^{2^b})$.
The two ranks have exponents summing to $2r$; self-orthogonality
in each factor makes each exponent $r$.
Thus
\[
                 (H_r)_k=\ker F_{M_0}^r\times\ker F_{N_0}^r.
 \tag{B-RD45.3}
\]

We give the logarithm argument in the required prime-two form.
Let $G/R$ be any finite-height formal group whose special
two-series has no monomial of total degree less than four.
This condition holds for a product of one-dimensional
height-two groups and is invariant under integral changes
of formal coordinates. Its linear term is $2X$; coefficients
of total degrees two and three are in $\pi R$, and all higher
coefficients are integral. Therefore $[2]_G$ maps coefficient
valuation at least $q$ to valuation at least
$\min(a+q,1+2q,4q)$.
For any augmentation vector $h\in tR[[t]]^{\dim G}$,
the coefficients of $[2]^u_G(\pi h)$ have valuations at
least $q_u$.

Write $c_\alpha$ for a coefficient of total degree $m$ of
the normalized logarithm. Its differential is the integral
matrix of invariant differentials, so
\[
 v(c_\alpha)\ge
 -a\min_{j:\alpha_j>0}v_2(\alpha_j)
 \ge-a\lfloor\log_2m\rfloor.
\]
Since $\lfloor\log_2m\rfloor\le m-1$, substitution of a
vector of coefficient valuation $q>a$ gives logarithm
coefficient valuation at least $q$. All substitutions are
coefficientwise finite because the vectors have zero constant
term. Consequently
\[
        2^u\log_G(\pi h)\in\pi^{q_u}R[[t]]^{\dim G}
                  \quad\text{if }q_u>a.
 \tag{B-RD45.4}
\]

The Cartier-module proof of Lemma~\ref{lem:B-RD34-kernel}
now applies literally with prime two, with
\[
 \mathcal R=R[[\Delta]][1/2],\qquad J=\pi^aR=2R.
\]
To specify the prime-two point completely, the published original
explicitly imposes no $p>2$ condition
\cite[p.~663]{Bondarko2006Generic}.
Its Propositions~9--11 and Lemma~3, pp.~674--675, give
$J[[\Delta]]^d\subset D_G$, the identity
$\mathcal C_G=\Omega\pi D_{G_\pi}$, and the bounded-projection
embedding. The integral Witt-coefficient action of $\Omega$
preserves coefficient valuation ideals, so (B-RD45.4) gives
$2^u\mathcal C_G\subset J\{\{\Delta\}\}^d$.
Equation~(8), p.~676, gives
\[
 \mathbf f^uAz\equiv2^uAz\Delta^{-u}
                              \pmod{2R((\Delta))^b}.
\]
Because $2R=J$, an integral matrix $A$ and a bounded projection
$Az$ therefore satisfy $\mathbf f^uAz\in J[[\Delta]]^b$.
The saturation argument of Proposition~16, p.~681, and
Theorem~2(8), p.~668, then show that $F^u$ kills the special
kernel of any formal homomorphism having no nonzero finite
flat subgroup in its kernel, finite special kernel, and target
$G$ of this type. This is exactly the proof written in
B-RD34, including its distinction between absence of a flat
kernel and a finite-level closed immersion; no additional
prime-two full-faithfulness claim is inserted.

Finally set $P=\mathcal M\times\mathcal N$, $Q=P/H_n$, and
consider
\[
        P\longrightarrow Q\times\mathcal M,
        \qquad z\longmapsto(qz,\operatorname{pr}_{\mathcal M}z).
\]
It has no nonzero finite flat kernel, since generically the
graph of an isomorphism meets $0\times N$ trivially.
By (B-RD45.3) its special kernel is
$0\times\ker F_{N_0}^n$.
The target special fibre is a product of three height-two
formal groups: the quotient special group is the product of
the two Frobenius twists, because finite-flat quotient formation
commutes with base change. If $q_u>a$, the preceding estimate
would make $F^u$ annihilate this special kernel.
For $u<n$ its image has rank $2^{n-u}>1$, a contradiction.
Taking $u=n-1$ proves the asserted inequality.
\end{proof}

\begin{theorem}[B-RD45: a finite set of possible dyadic least indices]
\label{thm:B-RD45}
\status{Proved; complete independent parent and obstructions/crystals proof, original-source and exact-code audits passed.}
For the manuscript question with $(p,e,h)=(2,1,2)$,
\[
                         3\le\nu_{\min}(2,1,2)\le7.
 \tag{B-RD45.5}
\]
More precisely a failure at the indicated torsion level must
compare two potentially supersingular elliptic BT groups after
completed unramified base change, with joint faithful finite
inertia degree in the corresponding row:
\[
\begin{array}{c|c|c}
\text{torsion level}&\text{manuscript index}&\text{remaining joint degrees}\\
\hline
16&3&24,48,96,192\\
32&4&48,96,192\\
64&5&96,192\\
128&6&192\\
256&7&\text{none}.
\end{array}
 \tag{B-RD45.6}
\]
The table is a necessary-condition table. It asserts no
counterexample in any remaining row, and does not decide
which of the five possible least indices is attained.
\end{theorem}
\begin{proof}
The lower bound is Proposition~\ref{prop:B-RD43}.
Proposition~\ref{prop:B-RD44} excludes every mixed case from
level sixteen onward and gives joint degrees $m=3\cdot2^r$,
$0\le r\le6$. Since $e_K=1$ and the residue field is algebraically
closed, a faithful joint good-reduction extension has
$a=v_L(2)=m$.
For $q\ge1$ the recurrence in (B-RD45.1) is simply
$q_{i+1}=\min(a+q_i,1+2q_i)$.
Starting from one, its first crossing of $a$ occurs at the
first integer $u$ with $2^{u+1}-1>a$.
For $a=3\cdot2^r$ with $r\ge1$, this is $u=r+1$, since
$2^{r+1}-1<a<2^{r+2}-1$. For $r=0$, the first values are
$1,3,6$, so the crossing occurs at $u=2$; this degree is already
excluded at every level considered in the table.
Lemma~\ref{lem:B-RD45} therefore requires $n\le r+1$ for
any remaining $r\ge1$ at level $2^n$.
For $n=4,5,6,7,8$ this gives exactly the displayed rows,
including impossibility at $n=8$ because $r\le6$.
The manuscript convention is $M(\nu)=M[2^{\nu+1}]$, which
gives the stated upper index seven.
\end{proof}

\begin{remark}[Scope of the arithmetic certificate]\label{auto:B:a2534788021e}
\status{Proved under stated hypotheses}

The program \path{code/verify_p2_height_two_log_gate.py}
checks the cubic projector in the quaternion algebra, the
recurrence and first-crossing formula, the logarithm degree
inequalities, and every joint-degree entry of (B-RD45.6).
Its retained output records the exact scope and runtime; the
author run passed in $0.0776823170017451$ seconds, with
$192000$ logarithm-degree cases and $24576$ coefficient-valuation
cases. The first draft's crossing formula omitted the exceptional
case $a=3$, caught during the proof review and retained in
\path{failed_candidates/B45_first_crossing_r0_v1.tex}.
Its correction did not change any row of (B-RD45.6).
The full Cartier saturation and graph arguments are proved
above, not inferred from the finite checks. In particular
no finite decision procedure for sixteen-congruences has
been executed.
\end{remark}

\par\smallskip\noindent\textbf{Independent audits.}
The parent independently reread the entire argument and the published
Bondarko original at pp.663, 668, 674--676 and 681, then reran the
complete certificate in $0.23222537100082263$ seconds, with the
audit recorded at 2026-11-19 14:14:46.950471 UTC.
The obstructions/crystals agent completed the same full proof,
source and code audit, rerunning it in $0.05902020299981814$ seconds.
Both preserved the author output. The bound and necessary-condition
table are proved. The sharper exact value
$\nu_{\min}(2,1,2)=3$ is subsequently established in
Theorem~\ref{thm:B-RD48} by the marked-congruence analysis;
the present theorem retains its original independent scope.


\subsection{Excluding the trivial residual dyadic comparison}

\begin{lemma}[B-RD46: the cubic part of dyadic elliptic inertia]
\label{lem:B-RD46-discriminant}
\status{Proved; complete independent parent and obstructions/crystals proof and exact-code audits passed.}
Let $K$ be a complete discretely valued field of characteristic zero
with perfect residue field of characteristic two. Suppose that an
elliptic curve $E/K$ is potentially supersingular. If $\delta$ is the
valuation of the discriminant of any Weierstrass equation over $K$,
then the order of the finite special-inertia image is divisible by
three if and only if $3\nmid\delta$. Consequently its character on
elements of odd order is determined by whether $3\mid\delta$.
Here ``odd order'' refers to the image in a common finite
special-inertia quotient.
\end{lemma}
\begin{proof}
Pass to the completed maximal unramified extension, which preserves
the valuation and the special-inertia representation. Let $L/K$ be
the minimal Galois extension over which the curve has good reduction,
with faithful special action $H$ and $m=[L:K]$. The explicit
supersingular automorphism group and finite-subgroup classification
in Lemma~\ref{lem:B-RD44-conjugacy} show that
$m\in\{1,2,3,4,6,8,24\}$.
If $\omega$ is the invariant differential of the equation over $K$
and $\omega_L$ a good-model invariant differential, write
$\omega=c\omega_L$. For a coordinate scaling $x=u^2x_0$, $y=u^3y_0$, one has
$\Delta=u^{12}\Delta_0$ and $\omega=u^{-1}\omega_0$.
Thus the discriminant transformation rule gives
\[
                         12v_L(c)=-m\delta.             \tag{B-RD46.1}
\]
This identity is valid even if the original equation is not integral.
If $3\nmid m$, it immediately gives $3\mid\delta$.

Suppose $3\mid m$ and choose $\tau\in H$ of order three.
Its action on a uniformizer satisfies
$\overline{\tau(\pi_L)/\pi_L}=\zeta$, a primitive cube root:
the kernel of the tame cotangent character is a two-group.
Since the residue action is trivial, the multiplier of $\tau$ on
the good special differential is $\zeta^{-v_L(c)}$.
This multiplier is nontrivial. Indeed an automorphism of order
three with derivative one at the origin can be linearized by
averaging its local parameter, using that three is invertible;
it would then be the identity on the formal completion and hence
on the elliptic curve. Faithfulness excludes that possibility.
Thus $3\nmid v_L(c)$, and (B-RD46.1), with $v_3(m)=1$, gives
$3\nmid\delta$.

The only odd orders in $H$ are one and three. An order-three
automorphism has eigenvalues $\zeta_3,\zeta_3^{-1}$ on the
rational special Dieudonne module and trace $-1$; the identity
has trace two. If the cubic quotient is nontrivial, its kernel
on tame inertia is the unique subgroup of index three.
Hence this determines the character as a function on inertia,
not merely its abstract character table.
\end{proof}

\begin{proposition}[B-RD46: rational two-torsion excludes the ambiguity]
\label{prop:B-RD46}
\status{Proved; complete independent parent and obstructions/crystals proof and exact-code audits passed.}
Let $K$ be a complete discretely valued field of mixed characteristic
$(0,2)$, with perfect residue field and $v_K(2)=1$. Let $E,E'/K$ be
potentially supersingular elliptic curves with rational full
two-torsion. If $E[4]\simeq E'[4]$ over $K$, their special
inertia characters agree on elements of odd order.

In particular, in the remaining height-two, ramification-one
sixteen-torsion comparison of Proposition~\ref{prop:B-RD44},
the common residual inertia image cannot be trivial. After
completed maximal unramified base change, it is either $C_2$
or $C_3$. This proposition decides neither remaining case.
\end{proposition}
\begin{proof}
Every curve with rational full two-torsion is a quadratic twist
of a Legendre curve
\[
              E_\lambda:\quad y^2=x(x-1)(x-\lambda),
              \qquad \lambda\in K\setminus\{0,1\}.
\]
Permuting the three nonzero two-torsion points permits the
normalization
\[
 v_K(\lambda)=q\ge0,\qquad v_K(1-\lambda)=0.
                                                               \tag{B-RD46.2}
\]
If $q=0$, both differences are units. If $q>0$, the formula
\[
 j(E_\lambda)=2^8\frac{(1-\lambda+\lambda^2)^3}
                         {\lambda^2(1-\lambda)^2}
\]
gives $v_K(j)=8-2q$. Potential supersingularity requires
$v_K(j)>0$, so $q\in\{1,2,3\}$ in this case.
A twist by $d\in K^\times$ has discriminant
$2^4d^6\lambda^2(1-\lambda)^2$, and therefore
\[
 \delta\equiv4+2q\pmod3.
 \quad
 \begin{array}{c|rrrr}
 q&0&1&2&3\\\hline
 \delta\bmod3&1&0&2&1
 \end{array}                                             \tag{B-RD46.3}
\]
The preceding lemma shows that different odd-order characters
would require one normalized parameter to have $q=1$ and the
other to have $q'=0,2$, or $3$.

We compute the complete projective four-torsion action. Choose
\[
 i^2=-1,\qquad a^2=\lambda,\qquad b^2=1-\lambda,
\]
and take the basis
\[
 P=(a,ia(a-1)),\qquad Q=(1+b,b(1+b)).
                                                               \tag{B-RD46.4}
\]
The chord-and-tangent law gives
$2P=(0,0)$ and $2Q=(1,0)$.
Changing the sign of $i$ negates $P$; changing the sign of $a$
sends $P$ to $P+2Q$; changing the sign of $b$ sends $Q$ to
$Q+2P$. These statements remain valid when the three quadratic
characters are dependent: their simultaneous effects commute
modulo four. Write $\chi_z:G_K\to\mathbf F_2$ for the
quadratic Kummer character of $z$, and set
\[
 N=\begin{pmatrix}0&0\\1&0\end{pmatrix},\quad
 U=\begin{pmatrix}0&1\\0&0\end{pmatrix},\quad
 T=\begin{pmatrix}1&0\\0&0\end{pmatrix}.
\]
For the quadratic twist by $d$ the exact action is
\[
 \rho_4(\sigma)=I+2A(\sigma),\qquad
 A(\sigma)=\chi_d(\sigma)I+\chi_{-1}(\sigma)T
             +\chi_\lambda(\sigma)N
             +\chi_{1-\lambda}(\sigma)U.              \tag{B-RD46.5}
\]
Thus the scalar term accounts for every possible quadratic twist.
An isomorphism of the two four-torsion modules conjugates their
matrices $A$ by one fixed $C\in\GL_2(\mathbf F_2)$.

The valuation modulo two is a linear functional on
$K^\times/K^{\times2}$, zero on unit classes and nonzero on
the class of two. Kummer theory therefore supplies an element
$\sigma_0\in G_K$ fixing all unit square roots needed in both
displays and changing the sign of $\sqrt2$.
It is enough to apply this to the finite-dimensional space of
classes occurring in the two equations and twists; there is no
assumption that the full square-class group is finite.
For a normalized parameter of valuation $q$, evaluation at
$\sigma_0$ gives
\[
                       A(\sigma_0)=\epsilon I+(q\bmod2)N.
                                                               \tag{B-RD46.6}
\]
If $q=1$ and $q'=0$ or $2$, one of these matrices is nonscalar
and the other scalar, so they cannot be conjugate.

It remains to compare $q=1$ and $q'=3$.
Equation (B-RD46.6) forces
$CNC^{-1}=N+(\epsilon+\epsilon')I$.
The determinant distinguishes $N$ from $N+I$, so
$\epsilon=\epsilon'$ and $C$ centralizes $N$.
The centralizer in $\GL_2(\mathbf F_2)$ is exactly
$\{I,I+N\}$. Conjugation by either matrix preserves the
upper-right entry of every matrix. Taking that entry in
(B-RD46.5), for every $\sigma$, now forces
\[
                         \chi_{1-\lambda}
                              =\chi_{1-\lambda'}.       \tag{B-RD46.7}
\]
The right side is zero: since $v_K(\lambda')=3$,
Hensel's lemma applied to $X^2-(1-\lambda')$ at $X=1$
gives a square root, using $3>2v_K(2)=2$.
The left side is nonzero. If $(1+w)^2=1-\lambda$, reduction
forces $v_K(w)\ge1$, whereas
$v_K(2w+w^2)\ge2$ contradicts $v_K(\lambda)=1$.
This contradiction proves the proposition.

For its final assertion, use the exact elliptic realizations
of Proposition~\ref{prop:B-RD44-realization} and restrict the
given sixteen-torsion isomorphism to four-torsion. A trivial
residual inertia action gives rational full two-torsion after
completed maximal unramified base change, which has just been
excluded. Finally the wild image in $\GL_2(\mathbf F_2)=S_3$
is a normal two-subgroup and its tame quotient is cyclic of
odd order. This excludes $S_3$ itself and leaves, among the
nontrivial images, precisely $C_2$ and $C_3$.
\end{proof}

\par\smallskip\noindent\textbf{Exact certificate and scope.}
The retained program
\path{code/verify_p2_residual_trivial_gate.py}
checks the displayed halving and translation identities over an
exact polynomial quotient, the six conjugating matrices, their
centralizer and upper-right-entry property, and the valuation
table. It does not enumerate elliptic curves or decide the
remaining residual $C_2$ and $C_3$ cases. Both new assertions have passed complete independent audits.

\par\smallskip\noindent\textbf{Self-check correction.}
The initial written draft put the opposite sign in (B-RD46.1).
The coordinate-scaling check above corrected it before independent
audit; the original is retained at
\path{failed_candidates/B46_differential_sign_v1.tex}.
Only nondivisibility by three is used, so this correction changes
no conclusion. The exact certificate passed in
$0.12948673400387634$ seconds.

\par\smallskip\noindent\textbf{Independent audits.}
The parent independently checked the entire proof and code, with a
preserved-output rerun of $0.35755470999720274$ seconds recorded
at 2026-11-19 14:16:15.762965 UTC. The obstructions/crystals agent
completed the same full audit, rerunning the certificate in
$0.3844016500006546$ seconds. The prior B44 original-source
verification supplies the finite special-inertia classification;
the remaining new steps are proved explicitly here.

\par\smallskip\noindent\textbf{Subsequent closure.}
The remaining residual branches are settled in
Theorem~\ref{thm:B-RD48}, which gives the exact value
$\nu_{\min}(2,1,2)=3$. The assertions and finite certificate
in this subsection retain their stated narrower scope.


\subsection{The residual order-two dyadic branch}

\begin{lemma}[B-RD47: the two quadratic fields of a two-isogeny]
\label{lem:B-RD47}
\status{Proved; complete independent parent and monodromy proof and exact-code audits passed.}
Let $K$ be complete, with algebraically closed residue field of
characteristic two and $v_K(2)=1$. Suppose $E/K$ is potentially
supersingular, and both $E$ and its quotient by a rational
order-two subgroup have nontrivial residual image $C_2$.
Write
\[
 E:\quad y^2=x(x^2+a x+b),\qquad D=a^2-4b.
\]
The two quadratic two-torsion fields are $K(\sqrt D)$ and
$K(\sqrt b)$. The special-inertia image has order divisible by
three if and only if $v_K(D)$ and $v_K(b)$ have different parities.
\end{lemma}
\begin{proof}
Completing the square and translating the rational two-torsion
point gives the displayed equation, with $bD\ne0$.
The normalized quotient is
\[
 E_1:\quad Y^2=X(X^2-2aX+D),\qquad
 X=x+a+b/x,\quad Y=y(1-b/x^2).                         \tag{B-RD47.1}
\]
Direct substitution verifies the equation and
$dX/Y=dx/y$. The map has degree two, sends the origin and
$(0,0)$ to the origin, and hence is the quotient isogeny.
Its source two-torsion field is $K(\sqrt D)$; the target
quadratic discriminant is $16b$, giving $K(\sqrt b)$.
The hypotheses say that both $D$ and $b$ are nonsquares.

We shall use
\[
                         1+4\mathcal O_K
                              \subset K^{\times2}.     \tag{B-RD47.2}
\]
Indeed the equation $z^2+z=u$ has a residue solution for every
$u\in\mathcal O_K$, and that solution lifts by Hensel's lemma
since $2z+1$ is a unit. Then $(1+2z)^2=1+4u$.

A quadratic twist replaces $(a,b,D)$ by $(ta,t^2b,t^2D)$.
It preserves the two quadratic fields, their valuation
parities, and the discriminant valuation modulo three.
Choose such a twist to normalize
$\beta=v_K(b)\in\{0,1\}$, and write $\alpha=v_K(a)$,
allowing $\alpha=+\infty$ when $a=0$.
The exact invariants are
\[
 \Delta=16b^2D,\qquad c_4=16(a^2-3b),\qquad
 v_K(j)=8+3v_K(a^2-3b)-2\beta-v_K(D).                 \tag{B-RD47.3}
\]

If $\beta=0$ and $\alpha\le0$, then
$D=a^2(1-4b/a^2)$ is a square by (B-RD47.2), contrary to the
hypothesis. Thus $\alpha\ge1$. For $\alpha\ge2$ one has
$v_K(D)=2$ and $v_K(a^2-3b)=0$, so
$(v_K(\Delta),v_K(j))=(6,6)$.
For $\alpha=1$, put $a=2c$ with $c$ a unit and
$r=v_K(c^2-b)\ge0$.
If $r=0$, the same pair is $(6,6)$.
If $r\ge2$, the identity
$b=c^2(1-(c^2-b)/c^2)$ and (B-RD47.2) make $b$ a square,
again contrary to the hypothesis. Hence the only other case
is $r=1$, when $v_K(D)=3$, $v_K(a^2-3b)=0$, and the pair
is $(7,5)$.

If $\beta=1$ and $\alpha\le0$, the same factorization
$D=a^2(1-4b/a^2)$ makes $D$ a square. Thus $\alpha\ge1$.
For $\alpha=1$, its terms have valuations two and three,
so $v_K(D)=2$; also $v_K(a^2-3b)=1$. The invariant pair
is $(8,7)$. For $\alpha\ge2$, one instead has
$v_K(D)=3$, $v_K(a^2-3b)=1$, giving $(9,6)$.
We have exhausted all possibilities under the hypotheses:
\[
\begin{array}{c|c|c|c|c|c}
 \beta&\text{case}&v_K(D)&v_K(\Delta)&v_K(j)
       &3\mid |H|\\\hline
 0&\alpha\ge2\text{ or }(\alpha,r)=(1,0)&2&6&6&\text{no}\\
 0&(\alpha,r)=(1,1)&3&7&5&\text{yes}\\
 1&\alpha=1&2&8&7&\text{yes}\\
 1&\alpha\ge2&3&9&6&\text{no}
\end{array}                                                   \tag{B-RD47.4}
\]
The last column is Lemma~\ref{lem:B-RD46-discriminant}.
It is precisely the assertion about the two valuation parities.
\end{proof}

\begin{proposition}[B-RD47: order-two residual inertia is excluded]
\label{prop:B-RD47}
\status{Proved; complete independent parent and monodromy proof and exact-code audits passed.}
Let $K$ be a complete discretely valued field of mixed characteristic
$(0,2)$, with perfect residue field and $v_K(2)=1$. Suppose two
potentially supersingular elliptic curves have isomorphic
eight-torsion and their common residual inertia image is $C_2$.
Their special-inertia characters agree on elements of odd
order in a common finite special-inertia quotient.

Consequently every hypothetical height-two, ramification-one
sixteen-torsion ambiguity left by
Proposition~\ref{prop:B-RD44} has common residual inertia image
exactly $C_3$. No existence or nonexistence in that final branch
is asserted here.
\end{proposition}
\begin{proof}
Pass to the completed maximal unramified extension, preserving
the given isomorphism, the valuation, and the special character.
The unique nonzero invariant vector in the common residual
$C_2$-module identifies rational order-two subgroups
$C\subset E[2]$ and $C'\subset E'[2]$ under the given
isomorphism. Put $E_1=E/C$ and $E'_1=E'/C'$.
There is an exact scheme-theoretic construction
\[
 T=\{P\in E[8]:[4]P\in C\},\qquad
                         T/C\simeq E_1[4].             \tag{B-RD47.5}
\]
Indeed the preimage of $E_1[4]$ under the quotient isogeny is
exactly $T$; if $[4]P\in C$ then $[8]P=0$.
The given eight-torsion isomorphism therefore induces
$E_1[4]\simeq E'_1[4]$. In particular their two-torsion
representations agree.

Each quotient has a rational two-torsion point, namely the
dual-isogeny kernel, so its residual image is either trivial
or $C_2$. If it is trivial, both quotient curves have rational
full two-torsion, and Proposition~\ref{prop:B-RD46} applies to
their four-torsion isomorphism. Isogeny preserves the rational
special-inertia character, proving the assertion in this case.

Otherwise both quotient residual images are $C_2$. The
original two-torsion representations identify the nontrivial
quadratic fields $K(\sqrt D)$ and $K(\sqrt {D'})$, while the
quotient two-torsion representations identify
$K(\sqrt b)$ and $K(\sqrt {b'})$ in the equations of
Lemma~\ref{lem:B-RD47}. Equality of quadratic extensions
is equality of their nonzero Kummer classes, so it preserves
the valuation parity of each class. That lemma gives equality
of the cubic special-inertia conditions, hence equality of the
requested characters by Lemma~\ref{lem:B-RD46-discriminant}.

For the final statement, use the exact elliptic realizations
and mixed exclusion of Proposition~\ref{prop:B-RD44}, restrict
the sixteen-torsion isomorphism to eight-torsion, and combine
this proposition with the trivial-residual exclusion of
Proposition~\ref{prop:B-RD46}.
\end{proof}

\par\smallskip\noindent\textbf{Exact verification.}
The retained program
\path{code/verify_p2_residual_c2_gate.py}
verifies the quotient equation and normalized differential,
the two discriminants and $j$ formulas, and every row and
excluded-square case of the valuation table. The proof of
the finite-group-scheme quotient in (B-RD47.5) is given
above. No search of the remaining residual $C_3$ branch is
represented by that certificate. Both statements have passed complete independent audits.

The author certificate passed in $0.20210709299863083$ seconds.

\par\smallskip\noindent\textbf{Independent audits.}
The parent independently checked the full proof and exact code,
rerunning the certificate in $0.047221027001796756$ seconds
at 2026-11-19 14:16:15.762965 UTC. The monodromy agent completed
its own full proof and code audit, with an author-preserving
rerun of $0.15887224199832417$ seconds. Both checked the
finite-group-scheme quotient and the entire valuation table.

\par\smallskip\noindent\textbf{Subsequent closure.}
The remaining residual branches are settled in
Theorem~\ref{thm:B-RD48}, which gives the exact value
$\nu_{\min}(2,1,2)=3$. The assertions and finite certificate
in this subsection retain their stated narrower scope.


\subsection{The final dyadic residual branch and its exact least index}

\begin{lemma}[B-RD48: the marked level-eight square conditions]
\label{lem:B-RD48-models}
\status{Proved; complete independent proof, source and code audits passed.}
Let $E:y^2=x^3+A x+B$ be an elliptic curve over a field $K$ of
characteristic zero, let $\theta$ be a root of its cubic, and put
\[
 \mathcal D=-4A^3-27B^2,\qquad
 \gamma=-(3\theta^2+A),\qquad
 q_\theta(t)=t^2-\frac{2\theta}{3}t-
                         \frac{2\theta^2+A}{9}.          \tag{B-RD48.1}
\]
If an elliptic curve $E'/K$ has an isomorphism $E[8]\simeq E'[8]$
of Weil-pairing multiplier $r\in\{1,3,5,7\}$, its induced
marked level-four point has a coordinate $t\in\mathbf P^1(K)$.
For finite $t$, the following element is a square in $K(\theta)$:
\[
 \alpha_r q_\theta(t),\qquad
 \begin{array}{c|rrrr}
 r&1&5&3&7\\\hline
 \alpha_r&1&\mathcal D&\mathcal D/\gamma&\gamma.
 \end{array}                                             \tag{B-RD48.2}
\]
For the level-four family $E_t$ with $E_\infty=E$, let
\[
 \Phi_E(t)=\operatorname{Norm}_{K(\theta)/K}q_\theta(t)
                                                               \tag{B-RD48.3}
\]
when the cubic is irreducible. Then
\[
                  \Delta(E_t)=\Delta(E)\Phi_E(t)^4.      \tag{B-RD48.4}
\]
The curve $E'$ is $E_t$ for $r=1,5$, and the quadratic twist
$E_t^{\Delta(E)}$ for $r=3,7$, in the compatible families.
At $t=\infty$ these are respectively $E$ and $E^{\Delta(E)}$.
In particular the discriminant valuations of $E'$ and $E_t$
are congruent modulo three over any discretely valued $K$.
\end{lemma}
\begin{proof}
Use Chen's marked moduli models
\cite[Theorems~1.1--1.4, pp.2--3; Section~3, pp.5--7;
Sections~5--6, pp.10--16; Appendix, p.21]{Chen2014Mod8}.
Here we spell out both the field scope and the normalization.
The finite-level moduli problem records an actual isomorphism
of eight-torsion with prescribed Weil-pairing multiplier.
Its forgetful map records the induced four-torsion marking.
Chen constructs its covers of the level-four parameter line
by square roots and computes the twists by cocycles in the
finite group $\GL_2(\mathbf Z/8)$.
The calculation is performed on generators of that finite
group (Lemmas~5.1--5.5 and 6.1--6.3), so it can be restricted
to the image of $G_K$ for any characteristic-zero field $K$.
Equivalently, these are identities for the universal marked
elliptic curve over
$\mathbf Q[A,B,(4A^3+27B^2)^{-1}]$, followed by base change.
They do not require that $A,B$ be algebraic over $\mathbf Q$,
and they apply at $j=0$ and $1728$ as well. Nontrivial elliptic
automorphisms cause no extra quadratic ambiguity: a full
level-four marking has no nonidentity automorphism fixing it.

The level-four completeness and field scope also follow directly
from Fisher's Theorem~13.2, stated for arbitrary $K$, with its
proof and the discussion of the special $j$ values
\cite[Theorem~13.2 and proof, preprint pp.29--30;
discussion and reverse-four twist, p.31]{Fisher2010Hessian}.
In Chen's coordinate change of Section~3, put
$T(u,v)=uv(16u^4-v^4)$.
Its matrix is
$\left(\begin{smallmatrix}u&-T_v\\v&T_u\end{smallmatrix}\right)$,
whose determinant is $uT_u+vT_v=6T$.
The original elliptic discriminant is nonzero, so $T\ne0$.
Thus this coordinate change is invertible also at $j=0,1728$;
the exceptional-coordinate restriction of Fisher's
Corollary~13.4 is not being used.

For clarity, the equation for multiplier one, after comparing
coefficients of $1,\theta,\theta^2$, is exactly
\[
                  (a_0+a_1\theta+a_2\theta^2)^2
                                =q_\theta(t).
\]
The other three displayed models give the scaling factors in
(B-RD48.2). In particular the coefficient of $t^2$ in the
multiplier-three model is
\[
                    6A\theta^2-9B\theta+4A^2
                              =\mathcal D/\gamma,
\]
using $\theta^3+A\theta+B=0$; its remaining coefficients
are those of $\mathcal D q_\theta/\gamma$.
The multiplier-seven model similarly gives $\gamma q_\theta$.
This proves the necessary square conditions, including the
actual reversed markings. The quadratic twist for the
reversed level-four family is Chen's Section~3 and Lemma~3.1.

Here is an independent polynomial verification of the
discriminant normalization. Write $c=-A/27$, $d=-B/54$.
The norm in (B-RD48.3) is
\[
 \Phi_E(t)=t^6-15ct^4-40dt^3-45c^2t^2-24cdt+27c^3-32d^2.
                                                               \tag{B-RD48.5}
\]
Let $a_E(t),b_E(t)$ be the two polynomials in Chen's Appendix,
so that $E_t$ has coefficients $-27a_E(t),-54b_E(t)$.
Expansion gives the universal identity
\[
       a_E(t)^3-b_E(t)^2=(c^3-d^2)\Phi_E(t)^4,
\]
which is (B-RD48.4). All its coefficients, and the coefficient
comparisons above, are checked exactly in the retained certificate.
This is independently the degree-four instance of Fisher's
Hesse identity
\cite[Lemma~8.4, equation~(8.1), and Remark~8.6,
preprint pp.19--20]{Fisher2010Hessian}.
The leading terms of the two Appendix polynomials are
$ct^8,dt^{12}$, so the family at infinity is $E$ after the
coordinate scaling by $t^2$.

\emph{Source sign check.} The minus sign in (B-RD48.1) is
the one in Chen's Theorems~1.1--1.4 and agrees with the
Appendix discriminant identity just verified. The cusp-sum
formula in printed Lemma~3.2 has the opposite sign.
We do not import that sign into the computation. Replacing
$t$ by $-t$ changes no valuation used below, but the displayed
models and the certificate consistently use (B-RD48.1).
\end{proof}

\begin{proposition}[B-RD48: the residual cubic branch is excluded]
\label{prop:B-RD48}
\status{Proved; complete independent proof, source and code audits passed.}
Let $K$ be a complete discretely valued field of mixed
characteristic $(0,2)$ with perfect residue field and $v_K(2)=1$.
If two potentially supersingular elliptic curves have
isomorphic eight-torsion and common residual inertia image
$C_3$, their special-inertia characters agree on the elements
of odd order in a common finite special-inertia quotient.
\end{proposition}
\begin{proof}
Pass to the completed maximal unramified extension.
Its residue field is algebraically closed, every unit is a cube
by Hensel's lemma, and it contains $\zeta_3$.
The common two-torsion field is the unique tame cubic extension
\[
                   F=K(\alpha),\qquad \alpha^3=2,
                   \qquad v_F(2)=3.
\]
Suppose for a contradiction that the characters differ, and
name $E$ as the curve whose finite special-inertia image has
order prime to three. This is possible by
Lemma~\ref{lem:B-RD46-discriminant} and its character conclusion.

Complete the square and remove the quadratic term in a
Weierstrass equation. Its roots lie in $F$, have trace zero,
and form a three-cycle. Thus one root has the form
$a\alpha+b\alpha^2$ with $a,b\in K$.
The terms, when nonzero, have distinct valuations modulo three.
Direct multiplication of the three roots gives
\[
 A=-6ab,\qquad B=-2a^3-4b^3,\qquad
 \Delta=-1728(a^3-2b^3)^2.                             \tag{B-RD48.6}
\]
If $r_0=\min\{3v_K(a)+1,3v_K(b)+2\}$, the same formula gives
$v_K(\Delta)=4+2r_0$.
This valuation is divisible by three, by our choice of $E$.
Consequently the first term is the smaller one, $a\ne0$,
and $v_K(b)\ge v_K(a)$.
Twist both curves and their eight-torsion isomorphism by
$a^{-1}$. Such a common quadratic twist does not change the
characters on odd-order elements. The first curve becomes
\[
 E_u:\quad y^2=x^3-6u x-2-4u^3,
 \qquad u=b/a\in\mathcal O_K,
 \qquad \theta=\alpha+u\alpha^2.                       \tag{B-RD48.7}
\]
Here $v_F(\theta)=1$, $v_K(\Delta(E_u))=6$, and
$v_F(\mathcal D)=6$. Also
\[
             v_F(\gamma)=v_F(3\theta^2+A)=2,
\]
because the other term has valuation at least three.
The four scaling factors in (B-RD48.2) therefore have
valuations $0,6,4,2$, all even.

Let $t$ be the marked level-four coordinate of the second
curve supplied by Lemma~\ref{lem:B-RD48-models}.
The point at infinity gives $E_u$ or its discriminant twist,
which has the same odd-order character, so it is already
inconsistent with our supposition. Assume $t\in K$ and put
$k=v_K(t)$ when $t\ne0$.
The four terms of $q_\theta(t)$ in (B-RD48.1) have valuations
\[
                6k,\qquad4+3k,\qquad5,
                         \qquad3+3v_K(u),              \tag{B-RD48.8}
\]
with the last term omitted when $u=0$.
If $k\le0$, the first is the unique minimum, so
$v_F(q_\theta(t))=6k$.
If $k\ge1$, the unique minimum is three when $u$ is a unit,
and five when $u$ is not a unit. The same assertion holds
at $t=0$ after omitting its first two terms.
In the latter cases every
$\alpha_rq_\theta(t)$ has odd valuation in $F$ and hence
cannot be a square. This contradicts the actual level-eight
marking, for every possible multiplier $r$.

Thus $k\le0$. Since $F/K$ is totally ramified with residue
degree one,
\[
          v_K\!\left(\operatorname{Norm}_{F/K}q_\theta(t)\right)
                              =v_F(q_\theta(t))=6k.
\]
The exact discriminant identity gives
\[
                         v_K(\Delta(E_t))=6+24k,
\]
which is divisible by three. For a reversed multiplier the
additional discriminant twist changes this by a multiple of
six, hence preserves divisibility by three.
Lemma~\ref{lem:B-RD46-discriminant} says that the second
potentially supersingular curve has no cubic special inertia
either. This is the desired contradiction.
\end{proof}

\begin{theorem}[B-RD48: the sharp dyadic height-two index]
\label{thm:B-RD48}
\status{Proved; complete independent proof, source and code audits passed.}
In the notation and monodromy class of B.4,
\[
                            \nu_{\min}(2,1,2)=3.
\]
Thus sixteen-torsion determines the requested prime-to-two
special-inertia character of every permitted height-two group
over absolute ramification index one, while eight-torsion does not.
For the subclass consisting of two potentially supersingular
elliptic BT groups, eight-torsion already determines that character.
\end{theorem}
\begin{proof}
For two potentially supersingular elliptic curves, the residual
inertia images over the completed maximal unramified field are
only $1,C_2,C_3$. The trivial case follows from
Proposition~\ref{prop:B-RD46}, the order-two case from
Proposition~\ref{prop:B-RD47}, and the cubic case from the
preceding proposition. This proves the stated subclass result.

For arbitrary permitted height-two groups with isomorphic
sixteen-torsion, Proposition~\ref{prop:B-RD44} already excludes
every discrepancy except a comparison of two exact potentially
supersingular elliptic realizations. Apply the subclass result
to their eight-torsion. This proves the upper bound three in
the manuscript's indexing. Proposition~\ref{prop:B-RD43}
supplies the explicit ramification-one eight-torsion comparison
with unequal characters, proving the lower bound three.
No assertion about the least index for greater height or
greater absolute ramification follows from this small instance.
\end{proof}

\par\smallskip\noindent\textbf{Exact certificate and current status.}
The complete retained certificate
\path{code/verify_p2_residual_c3_gate.py}
checks the full universal discriminant polynomial, all needed
coefficient and scaling identities, the trace-zero normal form,
and $20\,301$ valuation cases together with $3\,721$ normalization
cases. It passed in $0.6555589939962374$ seconds.
The finite arithmetic is additional verification of the symbolic
proof; no finite sample replaces the parameter or moduli arguments.
All three assertions have passed complete independent proof,
source and code audits. The parent reran the complete certificate
in $0.8482220709993271$ seconds, with its audit recorded at
2026-11-19 14:30:51.582619 UTC. The obstructions/crystals agent
independently reread the original Chen and Fisher passages and
reran the full certificate in $0.5984987490010099$ seconds.
Both preserved the author output. The source-sign discrepancy
and its exact resolution remain recorded above; no failed
normalization has been silently discarded.


\subsection{Tame height two at minimal cyclotomic ramification}

\begin{lemma}[A rational second-level point forces a wild absolute index]
\label{lem:B-RD49-point}
\status{Proved; complete root and obstruction/crystal proof, original-source, and exact-code audits passed.}
Let $R$ be a complete discrete valuation ring of mixed characteristic
$(0,p)$ with algebraically closed residue field, and normalize its
valuation by $v(R^\times)=0$ and $v(\pi)=1$. Put $a=v(p)$.
Let $F/R$ be a one-dimensional formal group whose special fibre has
height two. If
\[
             a\leq(p-1)(p^2-1),\qquad p\nmid a,
                                                        \tag{B-RD49.1}
\]
then $F(\mathfrak m_R)$ has no point of exact order $p^2$.
\end{lemma}
\begin{proof}
Choose an identification of the special fibre with the height-two
Honda group, using \cite[Theorem~IV, pp.270--271]{Lazard1955Formal}.
The $p$-typical universal deformation gives an integral coordinate
and $u\in\mathfrak m_R$ in which
\[
 \begin{gathered}
 [p](T)=pT+c_pT^p+\sum_{j>p}c_jT^j,\qquad
 c_p=(1-p^{p-1})u,\\
 c_j\in(p,u)R\ (p<j<p^2),\qquad c_{p^2}\in R^\times.
 \end{gathered}
                                                        \tag{B-RD49.2}
\]
All coefficients are integral. Here the precise source is
\cite[Proposition~5.7 and its displayed congruence; equations
(12.3)--(12.5) and Proposition~12.10 with proof]{HopkinsGross1994Bundles}.
The latter proposition explicitly allows a complete local
$\mathbf Z_p$-algebra whose residue field extends $\mathbf F_p$;
there is no finite-residue-field restriction on $R$.
For clarity, the logarithm starts
$f(T)=T+(u/p)T^p+O(T^{p^2})$. Comparing
$f([p]T)=pf(T)$ gives the zero coefficients in degrees $2$ through
$p-1$ and the displayed exact value of $c_p$.
Proposition~5.7, with its second coefficient specialized to $1$,
gives the remaining coefficient ideal and the unit coefficient.
An integral invertible change of parameter preserves positive
valuations, so it does not change the question about rational
formal points.

Every nonzero point $z$ of order $p$ satisfies
\[
                           v(z)\leq a/(p-1).            \tag{B-RD49.3}
\]
Indeed, if this inequality failed, $pz$ would have strictly smaller
valuation than every other term of (B-RD49.2), an impossibility.
Now suppose $x$ has exact order $p^2$, and write
$A=v(x)\geq1$ and $Z=v([p]x)$. By (B-RD49.1) and (B-RD49.3),
\[
              Z\leq a/(p-1)\leq p^2-1,
              \qquad a+A>Z,\qquad p^2A>Z.             \tag{B-RD49.4}
\]
Put $b=v(u)$, allowing $b=+\infty$ when $u=0$.
If $b\geq a$, the linear term and every term of degree less than
$p^2$ have valuation at least $a+A$, while all later terms have
valuation at least $p^2A$. This contradicts (B-RD49.4).
Thus $1\leq b<a$. Every term strictly between degrees $p$ and
$p^2$ has valuation at least $b+jA>b+pA$; every term beyond degree
$p^2$ has valuation greater than $p^2A$. Therefore (B-RD49.4)
forces the unique contributing minimum to be the degree-$p$ term:
\[
                              Z=b+pA.                  \tag{B-RD49.5}
\]
At the point $[p]x$ of order $p$, one has
$(p^2-p)Z>b$, so the degree-$p^2$ term and all subsequent terms
have greater valuation than the degree-$p$ term. The intermediate
terms also have greater valuation. Cancellation consequently
requires the linear and degree-$p$ terms to have equal valuation:
\[
            a+Z=b+pZ,\qquad
            a=b+(p-1)Z=p\bigl(b+(p-1)A\bigr).
                                                        \tag{B-RD49.6}
\]
This contradicts $p\nmid a$.
\end{proof}

\begin{lemma}[Tame formal inertia for $p\geq5$]
\label{lem:B-RD49-inertia}
\status{Proved; complete root and obstruction/crystal proof, original-source, and exact-code audits passed.}
Let $p\geq5$ and let $K$ be complete discretely valued of mixed
characteristic $(0,p)$ with algebraically closed residue field.
A height-two BT group which acquires one-dimensional formal good
reduction has a faithful good-reduction extension which is tame
cyclic of degree dividing $p^2-1$. Two such groups have simultaneous
good formal reduction over a tame cyclic extension of degree
$m\mid p^2-1$.
\end{lemma}
\begin{proof}
The endomorphism algebra of the special height-two Honda group is
\[
 D=\mathbf Q_{p^2}\langle\Pi\rangle,
 \qquad \Pi^2=p,\qquad \Pi a=\sigma(a)\Pi,
 \quad a\in\mathbf Q_{p^2}.
                                                        \tag{B-RD49.7}
\]
Here $\sigma$ is the nontrivial unramified field automorphism. This is the same direct Dieudonne calculation
as (B-RD31.1), with $3$ replaced by $p$.
The maximal order has residue field $\mathbf F_{p^2}$, and the
kernel of reduction on its units is pro-$p$.
Every element of $D$ satisfies its reduced quadratic polynomial.
An element of order $p$ would generate
$\mathbf Q_p(\zeta_p)$ of degree $p-1>2$, so cannot occur.
Cauchy's theorem therefore makes every finite subgroup of the
unit group prime to $p$, and reduction embeds it in the cyclic
group $\mathbf F_{p^2}^{\times}$.

The kernel of the finite special inertia action can be removed
from a good-reduction extension by the crystalline descent and
actual-lattice argument of Lemma~\ref{lem:B-RD17-degree}.
Thus the faithful quotient of order dividing $p^2-1$ gives good
reduction for the original lattice, and its special fibre is still
of slope $1/2$. Over this strictly henselian complete base such an
extension is tame cyclic. The tame inertia group is procyclic,
so the common extension of two of them has degree equal to the
least common multiple of their degrees, again dividing $p^2-1$.
\end{proof}

\begin{proposition}[The mixed case is already detected at $p^2$]
\label{prop:B-RD49-mixed}
\status{Proved; complete root and obstruction/crystal proof, original-source, and exact-code audits passed.}
Let $p\geq5$, let $K$ be a complete discretely valued field of mixed
characteristic $(0,p)$ with perfect residue field, and suppose
$e_K=p-1$ and $\mu_p\subset K$.
Let $M$ be a height-two group which acquires supersingular good
reduction. Let $N$ admit a filtration by two potentially good
height-one groups whose dimensions sum to one. Then
\[
                            M[p^2]\not\simeq N[p^2].    \tag{B-RD49.8}
\]
\end{proposition}
\begin{proof}
Pass to the completed maximal unramified extension. Its absolute
ramification is unchanged, its residue field is algebraically
closed, and the torsion comparison would remain an isomorphism.
A potentially good height-one character is
$\chi_p^b\epsilon$, where $b\in\{0,1\}$ and
$\epsilon$ takes values in $\mu_{p-1}\subset\mathbf Z_p^\times$.
The determinant of $T_pM$ is $\chi_p\epsilon_0$ with
$\epsilon_0$ another such finite character. These are the same
height-one and determinant arguments used in
Propositions~\ref{prop:B-RD5} and~\ref{prop:B-RD18}.

Transport the first filtered subgroup of $N[p^2]$ to a cyclic
direct summand $C\subset M[p^2]$, with character
$\chi_p^b\epsilon$. If $b=0$, twist $M$ by $\epsilon^{-1}$;
the resulting height-two group has a rational point of order
$p^2$. If $b=1$, first form the BT quotient $M'=M/C$.
The image of $M[p^2]$ in $M'[p^2]$ is cyclic of order $p^2$
with character $\epsilon_0\epsilon^{-1}$; twisting by its
inverse again produces a rational point of exact order $p^2$.
The schematic closure of $C$ in a good model is finite flat,
and the quotient good BT group remains isogenous to the
original one. These quotients and finite twists therefore remain
potentially good, of dimension one and special slope $1/2$.

Lemma~\ref{lem:B-RD49-inertia} gives a good formal model over
an extension of ramification degree $m\mid p^2-1$.
Its absolute ramification is $a=(p-1)m$, which satisfies
(B-RD49.1). The rational torsion point extends by properness of
the finite flat torsion scheme, and reduces to the identity
because its special fibre is connected. It is consequently a
formal point forbidden by Lemma~\ref{lem:B-RD49-point}.
\end{proof}

\begin{theorem}[B-RD49: a one-step least-index interval in tame height two]
\label{thm:B-RD49}
\status{Proved; complete root and obstruction/crystal proof, original-source, and exact-code audits passed.}
For every prime $p\geq5$, in the setting of B.4 with
$e=p-1$, $\mu_p\subset K$, and height two,
\[
                         1\leq\nu_{\min}(p,p-1,2)\leq2.
                                                        \tag{B-RD49.9}
\]
More precisely, $p^3$-torsion determines the requested character
for every group satisfying the manuscript monodromy condition.
Any failure at $p^2$ must involve two whole potentially
supersingular groups. After completed maximal unramified base
change their finite determinant characters agree, and their
joint faithful tame inertia has degree
\[
                            m\mid p^2-1,\qquad m>p.    \tag{B-RD49.10}
\]
No existence of one of these remaining pairs is asserted.
\end{theorem}
\begin{proof}
The lower bound is Proposition~\ref{prop:B-RD18} with $n=1$.
Suppose two allowed groups have isomorphic $p^2$-torsion and
pass to the completed maximal unramified base. For either
monodromy filtration let $d\in\{0,1,2\}$ be the sum of the
dimensions of its good graded pieces. The generic determinant
on inertia is $\chi_p^d\epsilon$, with $\epsilon$ a finite
$\mu_{p-1}$-valued character. The image of $\chi_p$ modulo
$p^2$ has order $p$: the field contains $\mu_p$ but cannot
contain $\mu_{p^2}$, since that would require absolute
ramification divisible by $p(p-1)$. Equality of the two
determinants modulo $p^2$ thus gives equality of $d$, because
$|d_M-d_N|\leq2<p$, and then equality of $\epsilon$.

For $d=0$, the special character is the Brauer lift of the
semisimplified generic mod-$p$ inertia representation;
additivity makes this true for any extensions between the
graded pieces. For $d=2$, Cartier duality reduces to $d=0$.
For $d=1$, a potentially ordinary height-two graded group has
a canonical connected--etale filtration, which descends to
$K$, so each group is either filtered by height-one groups or
is itself potentially supersingular. Two filtered groups have
finite graded characters determined by the semisimplified
mod-$p$ representation, since $\chi_p\equiv1\pmod p$ on
inertia. The mixed case is excluded by
Proposition~\ref{prop:B-RD49-mixed}.

It remains to compare two whole potentially supersingular
groups. Lemma~\ref{lem:B-RD49-inertia} supplies a common good
extension with $e_L=(p-1)m$, $m\mid p^2-1$. As
\[
                     e_L\leq(p-1)(p^2-1)<(p-1)p^2,
\]
\cite[Theorem~E(1), pp.663--664; Proposition~18 and the full
proof of Theorem~E, p.685]{Bondarko2006Generic}, with $u=2$
and $\ell=1$, extends an isomorphism of generic $p^3$-torsion
to an isomorphism of the special $p$-torsion models. The
extension is compatible with the prescribed generic map and
therefore with descent. The resulting equal mod-$p$ special
Dieudonne representations have the same prime-to-$p$ Brauer
character, hence the same characteristic-zero finite inertia
character. This proves the upper bound. If $m\leq p$, the
same theorem with $u=1$, $\ell=1$ works already from
$p^2$-torsion and excludes a counterexample. The finite
norm characters agree by the determinant argument above,
which proves every restriction in (B-RD49.10).
\end{proof}

\begin{remark}[Scope and exact arithmetic check]
\label{rem:B-RD49-scope}
\status{Proved under stated hypotheses}
This argument includes nontrivial finite reduced norms and does
not replace arbitrary formal groups by elliptic curves. The
least value within (B-RD49.9) remains open here. The coefficient
normalization and universal deformation assertions used above
are read from the original Hopkins--Gross paper; the direct PDF
download and screenshot endpoint failed, while its indexed
original full text was successfully opened and read.
The retained program \path{code/reduction_descent/verify_tame_height_two_gate.py}
checks the exact coefficient identity, the cancellation
arithmetic and the tame degree and determinant restrictions.
It reconstructed the $p$-series through degree $p^2$ for
$p=2,3,5,7,11$, checked $136220$ cancellation configurations,
$6984$ joint tame-degree cases and $1320$ determinant cases.
The author run passed in $1.3878417579981033$ seconds; independent
root and obstruction/crystal runs passed in $0.6928875380035606$
and $0.7361230790047557$ seconds, respectively. Their runs preserved
the author output. No existence claim for the remaining pairs is
inferred from these finite checks.
\end{remark}


% Post-CP17 refinement; full root and obstruction/crystal audits passed.
\subsection{The logarithm equality boundary and square finite determinants}

\begin{lemma}[The closed logarithm boundary]
\label{lem:B-RD50-boundary}
\status{Proved; complete root and obstruction/crystal proof, original-source, and exact-code audits passed.}
Let $R$ be a complete discrete valuation ring of mixed characteristic
$(0,p)$ with perfect residue field, and put $a=v(p)$ for its normalized
valuation. Let $f:F\to G$ be a homomorphism of finite-height formal
groups with no nonzero finite flat subgroup in its kernel and with
finite special-fibre kernel. Suppose $G_k$ is isomorphic to a product
of one-dimensional height-two formal groups. If
\[
                              a\leq p^2-1,
                                                        \tag{B-RD50.1}
\]
then relative Frobenius annihilates $\ker f_k$.
\end{lemma}
\begin{proof}
We keep the actual Cartier-module argument of
Lemma~\ref{lem:B-RD34-kernel}, and verify its coefficient bound at
an equality previously left unused. The normalized invariant
differential matrix gives
\[
                  U([p]X)J_{[p]}(X)=pU(X).
\]
Every monomial coefficient with an exponent not divisible by $p$
is consequently divisible by $p$. In particular all nonlinear
terms of total degree less than $p$ are divisible by $p$.
The product-height-two special-fibre condition says that every
coefficient of total degree $p$ through $p^2-1$ lies in the
maximal ideal. It is invariant under integral changes of formal
coordinates, exactly as proved in Lemma~\ref{lem:B-RD34-log}.
Thus for any augmentation vector $h\in tR[[t]]^{\dim G}$,
\[
        [p]_G(\pi h)\in\pi^qR[[t]]^{\dim G},\qquad
        q=\min\{a+1,p+1,p^2\}\geq a/(p-1).
                                                        \tag{B-RD50.2}
\]
The last inequality follows from (B-RD50.1).

For a coefficient $c_\alpha$ of the normalized logarithm, of
positive total degree $n=|\alpha|$, integration of the invariant
differentials gives
\[
 v(c_\alpha)\geq-a\lfloor\log_p n\rfloor,
 \qquad \lfloor\log_p n\rfloor\leq\frac{n-1}{p-1}.
\]
It follows that if all coefficients of an augmentation vector
$z$ have valuation at least $q\geq a/(p-1)$, then every
coefficient of $\log_G(z)$ has valuation at least $q$.
Equality in the bound causes no problem: formal substitution
in the augmentation variable $t$ involves finitely many terms
in each coefficient. No analytic-convergence assertion at the
boundary is used. Apply this to (B-RD50.2) to obtain
\[
                      p\log_G(\pi h)\in\pi^qR[[t]]^{\dim G}.
                                                        \tag{B-RD50.3}
\]

For precision, put $J=\pi^{\lceil a/(p-1)\rceil}R$ and use
Bondarko's modules $D_G,\mathcal C_G$ and his bounded-series ring
$\mathcal R=R[[\Delta]][1/p]$. His exact identities are
\[
 \mathcal C_G=\Omega\pi D_{G_\pi},\quad
 J[[\Delta]]^{\dim G}\subset D_G,\quad
 E(G_k)=D_G/(D_G\cap\mathcal R^{\dim G}).
\]
The ideal in his Proposition~9 is
$\mathfrak m^{-[a/(1-p)]}=J$, including the integer equality
case. The precise source is
\cite[Propositions~9--12 and Lemma~3, pp.~674--676;
Proposition~16 and its full proof, p.~681]{Bondarko2006Generic}.
The integral Witt-coefficient action of $\Omega$ preserves each
coefficient-valuation ideal. Equation (B-RD50.3) therefore gives
$p\mathcal C_G\subset J\{\{\Delta\}\}^{\dim G}$.
The bounded-projection identity in Proposition~11(2), together
with $\mathbf fz=p(z-z(0))\Delta^{-1}$ and $pR\subset J$,
then gives
\[
 z\in D_G,\ Az\in\mathcal R^b
       \quad\Longrightarrow\quad
                 \mathbf fAz\in J[[\Delta]]^b
                                                        \tag{B-RD50.4}
\]
for every integral matrix $A$. This is the exact replacement
for (B-RD34.8), with one Cartier Frobenius.

The remaining saturation argument is unchanged, but its target
is explicitly $G$. Put $I'=df(0)(\operatorname{Frac}(R)^{\dim F})$,
choose an integral complement to $I'\cap R^{\dim G}$, and let
$\operatorname{pr}$ be the associated projection. Absence of a
finite flat kernel gives
\[
 df(0)D_F=D_G\cap I'[[\Delta]].
\]
For a class in the saturation of $E(F_k)$ inside $E(G_k)$,
choose $z\in D_G$. Then $\operatorname{pr}z$ is bounded.
By (B-RD50.4), $c=\mathbf f\operatorname{pr}z$ lies in
$J[[\Delta]]$, hence in $D_G$, and is bounded. Thus
$\mathbf fz-c\in df(0)D_F$, while $c$ reduces to zero.
The saturation quotient is annihilated by $\mathbf f$.
By \cite[Theorem~2(8), equation~(6), p.~668]{Bondarko2006Generic},
this is the Dieudonne module of $\ker f_k$, and $\mathbf f$
corresponds to relative Frobenius. This proves the assertion.
\end{proof}

\begin{proposition}[The square-determinant tame comparison]
\label{prop:B-RD50-square}
\status{Proved; complete root and obstruction/crystal proof, original-source, and exact-code audits passed.}
Let $p\geq5$ and let $K$ be complete discretely valued of mixed
characteristic $(0,p)$, with perfect residue field,
$e_K=p-1$, and $\mu_p\subset K$.
Suppose two height-two groups satisfying the manuscript monodromy
condition have isomorphic $p^2$-torsion and different requested
characters. After completed maximal unramified base change,
Theorem~\ref{thm:B-RD49} makes both groups potentially
supersingular, and their determinants have the common form
$\chi_p\psi$, with $\psi:I_K\to\mu_{p-1}$ finite.
Then $\psi$ has no square root in
$\operatorname{Hom}_{\mathrm{cont}}(I_K,\mu_{p-1})$.
In particular the joint faithful tame degree cannot equal $p+1$.
\end{proposition}
\begin{proof}
Assume instead that $\psi=\epsilon^2$ for a finite
$\mu_{p-1}$-valued character. Twist both groups by
$\epsilon^{-1}$. This preserves the torsion isomorphism and
multiplies both special characters by the same nonzero scalar
character, so it preserves inequality. Both new determinants
are cyclotomic. Their special finite actions have reduced norm
one, by the determinant comparison already used in
Proposition~\ref{prop:B-RD32}. The height-two division algebra
calculation in Lemma~\ref{lem:B-RD49-inertia} shows that their
finite images are cyclic of orders dividing $p+1$.
Indeed they are prime-to-$p$ roots of unity in an unramified
quadratic field, and the reduced norm is the $(p+1)$st power.
Their common faithful good field $L$ is therefore tame cyclic
of degree $m\mid p+1$, and $a=e_L=(p-1)m\leq p^2-1$.

Let $\mathcal M,\mathcal N$ be their good formal models and
let $\tau$ generate $\operatorname{Gal}(L/K)$.
Different norm-one characters mean that the two unordered
special eigenvalue sets $\{\zeta,\zeta^{-1}\}$ are disjoint.
Their minimal polynomials over $\mathbf Z_p$ have relatively
prime reductions: the roots are distinct Teichmuller roots
of unity, with their $p$th-power conjugate equal to their
inverse. The Chinese remainder theorem consequently gives
an integral polynomial in $\tau$ which is the projection
onto one factor of $M_0\times N_0$.

Let $H$ be the finite flat closure of the graph of the given
isomorphism in $\mathcal M[p^2]\times\mathcal N[p^2]$.
The perfect alternating pairings needed here are actual
integral BT pairings: cyclotomic determinants give them on
the full generic Tate lattices, and Tate's theorem extends
both the duality map and its inverse over $\mathcal O_L$.
This is exactly the argument in Proposition~\ref{prop:B-RD33},
using \cite[Theorem~4 and Corollaries~1--2, p.~181]
{Tate1967PDivisible}. Scale one pairing by the unit multiplier
of the torsion isomorphism. The graph closure is then its own
orthogonal, including after reduction, by flatness and equal
rank as in that proposition.

The integral tame projector preserves $H_k$, so
$H_k=J_M\times J_N$. A one-dimensional connected formal group
over an algebraically closed field has a unique subgroup of
rank $p^r$, namely its relative $F^r$-kernel. Orthogonality in
the perfectly paired $p^2$-torsion makes both exponents equal
to two. Thus
\[
                    H_k=M_0[p]\times N_0[p].            \tag{B-RD50.5}
\]
Set $Q=(\mathcal M\times\mathcal N)/H$. Its special fibre is
$M_0^{(p^2)}\times N_0^{(p^2)}$, a product of height-two groups.
The homomorphism
\[
 \mathcal M\times\mathcal N\longrightarrow Q\times\mathcal M,
              \qquad z\longmapsto(qz,\operatorname{pr}_{\mathcal M}z)
\]
has no nonzero finite flat subgroup in its kernel, because its
generic graph meets $0\times N$ trivially. Its special kernel
is $0\times N_0[p]=0\times\ker F_{N_0}^2$.
Lemma~\ref{lem:B-RD50-boundary} says that relative Frobenius
annihilates that kernel. Its Frobenius image instead has
rank $p$, a contradiction.

Finally, if the original joint degree were $p+1$, both original
special norm characters would already be trivial, so the same
argument would apply without twisting.
\end{proof}

\begin{corollary}[An explicit normalized remaining character class]
\label{cor:B-RD50-remaining}
\status{Proved; complete root and obstruction/crystal proof, original-source, and exact-code audits passed.}
In the remaining $p^2$ test of Theorem~\ref{thm:B-RD49}, the
common finite determinant must have nonsquare value on a tame
generator. After a simultaneous finite height-one twist one
may require that this value be primitive of order $p-1$.
In that normalization the joint faithful tame degree has the form
\[
          m=\frac{p^2-1}{d},\qquad
          d\mid\frac{p+1}{2^{v_2(p+1)}}.
                                                        \tag{B-RD50.6}
\]
For $p=5$ this leaves $m=8,24$; for $p=7$ it leaves only
$m=48$. These are necessary conditions, with no assertion that
any remaining pair exists. The least-index interval remains
$\{1,2\}$.
\end{corollary}
\begin{proof}
The tame inertia group is procyclic. A character into the cyclic
group $\mu_{p-1}$ has a square root precisely when its value on
a generator is a square. Choose a character $\eta$ with primitive
value. For nonsquare $\psi$, the character $\psi\eta^{-1}$ is
a square, so a simultaneous twist makes the determinant character
$\eta$.

A special tame eigenvalue is a Teichmuller root in
$\mu_{p^2-1}$, say $\zeta^c$ for a fixed primitive root.
Its reduced norm is $(\zeta^{p+1})^c$. Primitivity of the finite
determinant implies $\gcd(c,p-1)=1$ (the inverse norm convention,
if used, has the same consequence). Therefore $c$ is odd and
$\gcd(c,p^2-1)$ is an odd divisor of $p+1$, since
$\gcd(p-1,p+1)=2$. Each individual faithful order is consequently
$(p^2-1)/d_i$ with $d_i$ an odd divisor of $p+1$. Their joint
order is the least common multiple, obtained with
$d=\gcd(d_1,d_2)$, proving (B-RD50.6).
\end{proof}

\begin{remark}[Boundary scope and exact certificate]
\label{rem:B-RD50-scope}
\status{Proved under stated hypotheses.}
The new assertion uses the closed coefficientwise logarithm
bound. It neither assumes analytic convergence at equality nor
transfers Hasse-coefficient thickness from $\mathcal M,\mathcal N$
to the quotient $Q$. Only the product shape of $Q_k$ is used.
The earlier strict logarithm lemmas remain correct within their
stated ranges; this argument supplies an additional boundary case.
The retained exact program
\path{code/reduction_descent/verify_tame_norm_square_boundary.py}
checks 8,243,000 coefficientwise logarithm inequalities, 4,128
coprime tame-spectrum pairs, and 2,908 primitive-norm exponents.
The author run passed in 1.6194844739948167 seconds; the
independent root and obstruction/crystal runs passed in
1.491409575995931 and 1.5930737649978255 seconds, respectively.
Both independent runs preserved the author output. The actual
Cartier-module proof and original published Bondarko passages
were independently reread in full. These finite checks do not
assert the existence of any remaining torsion-congruent pair.
\end{remark}


% Separate post-CP17 refinement. B50 remains unchanged during its audit.
\subsection{The exact finite-character list after primitive-norm normalization}

\begin{lemma}[Primitive norm fibres and Frobenius orbits]
\label{lem:B-RD51-orbits}
\status{Proved; complete root and obstruction/crystal proof and exact-code audits passed.}
Let $p\geq5$, and put $m=(p^2-1)/d$, where $d$ is an odd
divisor of $p+1$. Write
\[
                      q=(p+1)/d,\qquad m=(p-1)q.
                                                        \tag{B-RD51.1}
\]
For any fixed primitive $\eta\in\mu_{p-1}$, the fibre
\[
             \{\lambda\in\mu_m:\lambda^{p+1}=\eta\}
                                                        \tag{B-RD51.2}
\]
has $q$ elements and is the disjoint union of $q/2$ orbits
of the involution $\lambda\mapsto\lambda^p$.
Consequently two distinct rank-two tame characters with that
primitive norm and joint faithful order $m$ exist as abstract
characters if and only if $q\geq4$.

In the normalization of Corollary~\ref{cor:B-RD50-remaining},
any unresolved $p^2$ comparison must therefore satisfy
\[
 m=(p-1)q,\qquad q\mid p+1,\qquad
 v_2(q)=v_2(p+1),\qquad q\geq4.                         \tag{B-RD51.3}
\]
The existence assertion in the first paragraph is about finite
characters only; it does not construct Barsotti--Tate groups
or a torsion congruence.
\end{lemma}
\begin{proof}
Since $d$ is odd and divides $p+1$, it is relatively prime to
$p-1$. Hence
\[
             \gcd(m,p+1)=q,\qquad
             |\operatorname{im}(\lambda\mapsto\lambda^{p+1})|
                      =m/q=p-1.
\]
This proves the fibre cardinality. Frobenius preserves the
fibre and has square equal to the identity there, since
$m\mid p^2-1$. A fixed point would belong to $\mu_{p-1}$;
its norm would then be $\lambda^2$, which cannot be primitive
in the even-order group $\mu_{p-1}$. Thus every orbit has
exactly two elements. These orbits are precisely the two
eigenvalues of the rank-two tame characters under discussion.

If $q=2$, the entire norm fibre is a single orbit and no two
different characters occur. Conversely, suppose $q\geq4$.
There is a primitive $m$th root in any specified primitive-norm
fibre. To see this without assuming a preferred generator,
the reduction map $(\mathbf Z/m\mathbf Z)^\times\to
(\mathbf Z/(p-1)\mathbf Z)^\times$ is surjective: lift a unit
modulo $p-1$, and avoid the finitely many additional primes
dividing $m$ by the Chinese remainder theorem. A primitive
$m$th root can therefore be chosen with norm $\eta$.
Choose an element of the fibre outside its Frobenius orbit.
The two resulting characters are distinct, and their joint
order is $m$ because the first one already has order $m$.
Equation (B-RD51.3) now follows from the normalized necessary
list in Corollary~\ref{cor:B-RD50-remaining}.
\end{proof}

\begin{corollary}[Two finite-character patterns in the first remaining prime]
\label{cor:B-RD51-p5}
\status{Proved; complete root and obstruction/crystal proof and exact-code audits passed.}
For the unresolved test $\nu_{\min}(5,4,2)\in\{1,2\}$,
any failure at $25$-torsion has, after the common normalization
above, joint faithful tame degree $24$.
Choose a primitive $24$th root $\zeta$ with
$\zeta^6=\eta$, the prescribed primitive norm. The only
individual character types are
\[
 \begin{array}{c|c|c}
 \text{type}&\text{eigenvalue exponents in }\zeta&
              \text{faithful order}\\ \hline
 A&\{1,5\}&24\\
 B&\{9,21\}&8\\
 C&\{13,17\}&24.
 \end{array}                                             \tag{B-RD51.4}
\]
Up to interchange and simultaneous quadratic twist, there
are only two distinct-type patterns: $\{A,B\}$ and $\{A,C\}$.
The first has individual orders $24,8$; the second consists
of the two opposite quadratic twists of a faithful order-$24$
character. This is an explicit remaining character list, not
an assertion that either pattern is realized by a pair of groups
with isomorphic $25$-torsion.
\end{corollary}
\begin{proof}
The odd divisors of $6$ are $1,3$. The latter gives $q=2$,
which Lemma~\ref{lem:B-RD51-orbits} excludes; hence $m=24$.
The norm fibre consists of the six exponents congruent to
$1$ modulo $4$. Multiplication by $5$ modulo $24$ gives
exactly the three pairs displayed in (B-RD51.4).
A quadratic twist adds $12$ to both exponents, interchanging
$A,C$ and fixing $B$. Therefore the pairs $\{A,B\}$ and
$\{B,C\}$ are equivalent under simultaneous quadratic twist,
while $\{A,C\}$ is the other possibility.
\end{proof}

\begin{remark}[Retained finite calculation and scope]
\label{rem:B-RD51-scope}
\status{Proved under stated hypotheses.}
The exact program
\path{code/reduction_descent/verify_tame_primitive_norm_orbits.py}
enumerates the primitive-norm fibres, their Frobenius orbits,
and the joint orders of distinct pairs for the listed small
primes. It also reproduces (B-RD51.4) and its simultaneous
quadratic-twist orbits. The general assertion is proved above;
the finite calculation is a retained check of the character
arithmetic and makes no group-existence inference. It checked
1,394 primitive-norm fibres and 292,614 distinct-character pairs
for the primes from $5$ through $101$. The author run took
0.406761682999786 seconds; independent root and
obstruction/crystal runs took 0.31523789799393853 and
0.3501194319978822 seconds. Both independent runs preserved
the author output.
\end{remark}


% Separate post-CP17 candidate. No unrestricted tame least-index conclusion.
\subsection{The remaining tame ambiguity cannot have two CM formal models}

Here a potentially supersingular height-two group is called
\emph{potentially CM} if, over some finite extension on which it
has good formal reduction, its rational endomorphism algebra
contains a quadratic field over $\mathbf Q_p$. This requires
endomorphisms defined over a finite extension; it is not the
endomorphism algebra of the constant etale group over an
algebraic closure of the generic field.

\begin{lemma}[Quadratic endomorphisms and a nonsquare determinant]
\label{lem:B-RD52-CM}
\status{Proved; complete root and obstruction/crystal proof, original-source and exact-code audits passed.}
Let $p\geq5$, and let $K$ be complete discretely valued of mixed
characteristic $(0,p)$, with algebraically closed residue field,
$e_K=p-1$, and $\mu_p\subset K$. Let $M$ be a potentially CM,
potentially supersingular height-two group. Write its determinant
as $\chi_p\psi$, where $\psi$ is its finite special determinant.
If $\psi$ is nonsquare as a $\mu_{p-1}$-valued character, then
its CM field is the unramified quadratic field $E=\mathbf Q_{p^2}$,
all its finite-extension endomorphisms are defined over $K$, and its Tate lattice
is free of rank one over $\mathcal O_E$.
\end{lemma}
\begin{proof}
Pass to a faithful tame good field $L/K$, as in
Lemma~\ref{lem:B-RD49-inertia}. Tate's full faithfulness
\cite[Theorem~4 and Corollaries~1--2, p.~181]{Tate1967PDivisible}
extends a generic finite-extension endomorphism to the good
Barsotti--Tate model; for a quasi-endomorphism, first clear a
power of $p$. Reduction of homomorphisms between finite-height
one-dimensional formal groups is injective, as in
\cite[Proposition~4.2 and its full proof]{HopkinsGross1994Bundles}.
Indeed, if a nonzero homomorphism had all coefficients divisible
by $\pi$, divide by their minimum positive valuation, say $\pi^r$.
Reducing the identity $f\circ[p]=[p]\circ f$ modulo
$\pi^{r+1}$ would give a nonzero series composed with the
nonzero special $[p]$ on the left and zero on the right.
This is impossible. The assertion also holds rationally.

Any endomorphism defined after a further finite extension of
$L$ therefore descends to $L$: after a common finite Galois
extension, its Galois conjugates have the same reduction, since
the residue field is algebraically closed and the model over
$L$ has good reduction. Injectivity makes the conjugates equal.
The rational endomorphism algebra over $L$ embeds both in the
height-two division algebra $D$ of the special fibre and in
$M_2(\mathbf Q_p)$ on the Tate module. It is a division algebra,
and the two-dimensional Tate module is a vector space over it;
its dimension over $\mathbf Q_p$ is consequently at most two.
Under the CM hypothesis it is a quadratic field $E$.

The cyclic finite special inertia group $\Gamma\subset\mathcal O_D^\times$
normalizes $E$. If $E$ is ramified, its prime-to-$p$ roots of
unity have order dividing $p-1$, and
$N_D(E)/E^\times$ has order two. Hence $|\Gamma|\mid2(p-1)$.
If $p\equiv3\pmod4$, every such root in
$\mathbf F_{p^2}^\times$ has square reduced norm: its
$(p^2-1)/2$ power is one. This makes $\psi$ square.
If $p\equiv1\pmod4$, use the generic normalizer instead.
On the one-dimensional $E$-space $V_p(M)$, its elements are
$a$ or $a\iota$, where $\iota$ is the nontrivial field
automorphism; their determinants are respectively
$N_{E/\mathbf Q_p}(a)$ and $-N_{E/\mathbf Q_p}(a)$.
A Galois element preserves the Tate lattice, so this determinant
is a unit and $a\in\mathcal O_E^\times$. The residue of its
norm is a square for a ramified quadratic field, and $-1$ is
also a square when $p\equiv1\pmod4$. Since
$\chi_p\equiv1\pmod p$, the residue of $\psi$ is again
always square. Thus ramified $E$ is impossible.

For unramified $E$, write
$D=E\oplus E\Pi$, with $\Pi^2=p$ and
$\Pi a=\sigma(a)\Pi$, where $\sigma$ is the nontrivial
automorphism of $E/\mathbf Q_p$. The noncentral coset of $N_D(E)$ consists
of $E^\times\Pi$, whose elements have odd $D$-valuation.
It has no units. Therefore $\Gamma$ centralizes $E$,
so the whole generic Galois action centralizes $E$.

Choose a Galois element $g$ for which $\psi(g)$ is nonsquare.
Then $\rho(g)\in E^\times$ has unit norm with nonsquare
residue, and hence its residue is not in $\mathbf F_p$.
It follows that $\mathbf Z_p[\rho(g)]=\mathcal O_E$:
its quadratic characteristic polynomial has irreducible,
separable reduction. The Tate lattice, stable under $\rho(g)$,
is thus stable under $\mathcal O_E$, and is free of rank one
over this discrete valuation ring.
\end{proof}

\begin{lemma}[The integral CM model and its full first wild layer]
\label{lem:B-RD52-LT}
\status{Proved; complete root and obstruction/crystal proof, original-source and exact-code audits passed.}
In the nonsquare case of Lemma~\ref{lem:B-RD52-CM}, choose the
orientation of $E$ on the tangent line. Let $F_0$ be the
Lubin--Tate formal $\mathcal O_E$-module with
$[p](T)=pT+T^{p^2}$, and let $\theta$ denote its Tate
character after base change to $K$. Then
\[
                  T_p(M)\simeq\mathcal O_E(\theta\epsilon)
                                                        \tag{B-RD52.1}
\]
as a $\mathbf Z_p$-representation, for a finite tame character
$\epsilon:G_K\to\mu_{p^2-1}$. Over every common faithful
tame good field, the image of wild inertia under $\theta$
modulo $p^2$ contains all of
$(1+p\mathcal O_E)/(1+p^2\mathcal O_E)$.
\end{lemma}
\begin{proof}
Tate's full faithfulness extends the integral
$\mathcal O_E$-action on the Tate lattice to the good formal
model; see \cite[Theorem~4 and Corollaries~1--2, p.~181]
{Tate1967PDivisible}. We include the uniqueness argument for
this model over an arbitrary algebraically closed perfect
residue field, rather than silently identifying that field
with $\overline{\mathbf F}_p$.

Here the integral conjugacy can be checked without passing
through a rational conjugator. Put $N=p^2-1$ and let $v$
be the image of a primitive $N$th root under the special
CM action. After possibly interchanging the two embeddings
of $E$, its residue in $\mathcal O_D/(\Pi)$ equals that of
the corresponding standard root $\lambda$. Then
\[
 x=\sum_{i=0}^{N-1}v^i\lambda^{-i},\qquad
 vx=x\lambda,\qquad x\equiv N\pmod\Pi.
\]
Thus $x$ is an integral unit and conjugates the whole
$\mathcal O_E$-action to the standard action on the
reduction of $F_0$.

The height-two universal deformation ring is $W(k)[[u]]$;
the arbitrary-residue-field assertion and the degree-$p$
coefficient are \cite[Proposition~5.7 and equations~(5.3)--(5.8);
Proposition~12.10 and its full proof]{HopkinsGross1994Bundles}.
Let $\lambda$ be a primitive $(p^2-1)$st root. Its special
automorphism induces an action on this ring whose tangent
eigenvalue is $a=\lambda^{p-1}\ne1$; this follows directly
by conjugating the degree-$p$ term by $T\mapsto\lambda T$.
Choose the parameter to vanish at the canonical deformation
$F_0$. Averaging it against this finite prime-to-$p$ action
gives a parameter $u_*$ on which the action is exactly
$u_*\mapsto a u_*$. An $\mathcal O_E$-linear deformation
is fixed by this action, so its classifying homomorphism
annihilates $u_*$, since $a-1$ is a unit. It is therefore
the base change of $F_0$. This proves the required integral
uniqueness. Changing the chosen orientation conjugates both
$\theta$ and $\epsilon$ by the unramified field automorphism,
and does not change the underlying rank-two special character.

The Tate character of the good model is consequently
$\theta|_{G_L}$. The quotient by $\theta$ over $K$ is an
$\mathcal O_E^\times$-valued character trivial on $G_L$,
and hence a finite tame character $\epsilon$. This gives
(B-RD52.1).

For the wild assertion put $q=p^2$ and $K_0=W(k)[1/p]$.
The elementary proof of
\cite[Theorem~2(b),(d),(e) and full proof, pp.~383--384]
{LubinTate1965FormalCM} applies after this residue extension.
Indeed, the primitive $p^2$-torsion polynomial is
\[
                 \Phi_2(T)=(pT+T^q)^{q-1}+p.           \tag{B-RD52.2}
\]
It remains Eisenstein over $W(k)$, of degree $q(q-1)$.
Every primitive root generates its entire division field:
the other torsion points are its $\mathcal O_E/p^2$ multiples.
The resulting Galois group is exactly
$(\mathcal O_E/p^2)^\times$. Its wild Sylow subgroup is
$(1+p\mathcal O_E)/(1+p^2\mathcal O_E)$.
Since $L/K_0$ is tame, its intersection with this division
field is fixed by that Sylow subgroup. The subgroup therefore
survives unchanged after base change to $L$.
\end{proof}

\begin{theorem}[Exact threshold in the potentially CM formal subclass]
\label{thm:B-RD52}
\status{Proved; complete root and obstruction/crystal proof, original-source and exact-code audits passed.}
Let $K$ be complete discretely valued of mixed characteristic
$(0,p)$ with perfect residue field, where $p\geq5$,
$e_K=p-1$, and $\mu_p\subset K$. Two
potentially CM, potentially supersingular height-two groups
with isomorphic $p^2$-torsion have equal requested special
characters. The level $p$ does not suffice within this same
subclass. Thus the least manuscript index in this subclass
is exactly one.

In particular any remaining pair in Theorem~\ref{thm:B-RD49}
with unequal characters and isomorphic $p^2$-torsion must
have at least one formal model without quadratic
finite-extension endomorphisms. This does not exclude
non-CM deformations in either of the two patterns of
Corollary~\ref{cor:B-RD51-p5}.
\end{theorem}
\begin{proof}
Pass to the completed maximal unramified base. The determinant
comparison in Theorem~\ref{thm:B-RD49} makes the two finite
determinants equal. If that common character is square,
Proposition~\ref{prop:B-RD50-square} already proves equality.
In the nonsquare case, the preceding lemmas give integral
identifications $T_p(M_i)=\mathcal O_E(\theta\epsilon_i)$,
with the same oriented $\theta$ after changing an
$\mathcal O_E$-orientation if needed.

Let $A$ be an intertwiner modulo $p^2$. On wild inertia in
a common tame good field, both representations contain
every multiplication operator $1+pb$, $b\in\mathcal O_E$.
Hence the reduction $\bar A$ commutes with all multiplication
operators from $\mathbf F_{p^2}$, and is itself multiplication
by a nonzero element of that field. Reducing the intertwining
identity for any $g\in G_K$ modulo $p$ now gives
$\bar\epsilon_1(g)=\bar\epsilon_2(g)$, since $\bar A$
commutes with all field scalars and $\bar\theta(g)\ne0$.
Reduction is injective on these prime-to-$p$ roots of unity;
therefore $\epsilon_1=\epsilon_2$. The special characters
are equal.

For the sharpness assertion, take $k=\overline{\mathbf F}_p$,
$K_0=W(k)[1/p]$, and $K=K_0(\mu_p)$. Choose a uniformizer
$\pi\in K$ with $\pi^{p-1}=-p$: the cyclotomic extension
has tame degree $p-1$, and over the strictly henselian
$K_0$ the unique tame extension of that degree is
$K_0((-p)^{1/(p-1)})$. Equivalently, remove the unit in a
uniformizer's $(p-1)$st power by Hensel's lemma.
Write $\omega$ for
the Teichmuller tame character given by a root
$\pi^{1/(p^2-1)}$. The first torsion equation for $F_0$ is
$T^{p^2-1}=-p$, so
$\bar\theta=\bar\omega^{p-1}$.
Twist $F_0$ generically by the two finite
$\mathcal O_E$-linear characters
\[
            \epsilon_1=\omega,\qquad
            \epsilon_2=\omega^{p^2-p+1}.               \tag{B-RD52.3}
\]
These twists are effective on the generic p-divisible groups
by finite Galois descent, and have good CM reduction after
a tame extension. Their mod-$p$ representations are,
respectively, multiplication by $\bar\omega^p$ and by
$\bar\omega$, which are isomorphic over $\mathbf F_p$
via field Frobenius.
Their finite special eigenvalue exponents are respectively
$\{1,p\}$ and $\{2-p,2p-1\}$ modulo $p^2-1$.
Here the exponents use the covariant convention in
Proposition~\ref{prop:B-RD4-dictionary}:
$D_{\mathrm{cov}}(G)[1/p]=D_{\cris}(V_pG)(-1)$.
The good Lubin--Tate factor has trivial finite special inertia,
the finite twist contributes $\epsilon_i$, and the period
Tate twist has trivial finite descent inertia. Thus these
are the actual displayed covariant exponents, not merely
an assertion up to inversion.
The second first exponent is neither $1$ nor $p$ modulo
$p^2-1$ for $p\ge5$, so the special characters differ.
Both finite determinants have the same primitive norm.
This proves failure at $p$ and the exact restricted index.
\end{proof}

\begin{remark}[Certificate and exact restriction]
\label{rem:B-RD52-scope}
\status{Proved under the stated CM restriction; no unrestricted least-index claim.}
The retained exact program
\path{code/reduction_descent/verify_tame_CM_gate.py}
checks the symbolic quadratic-field centralizer, all
$260{,}070$ two-by-two matrices in its six small-prime cases,
$75{,}984$ finite-field norm cases, $2{,}264$ ramified-norm
cases, $1{,}014$ Eisenstein coefficients, and the displayed
lower character pair for every prime from $5$ through $101$.
It also verifies $75{,}960$ terms in the unit-conjugacy telescoping identity.
Its final successful author runtime was $0.3032921849953709$ seconds.
Independent root and obstruction/crystal reruns took
$0.33256562399765244$ and $0.3663983289952739$ seconds,
respectively, and preserved the author output.
The first run used structural equality of unexpanded symbolic
matrix entries; that failed version and output are retained,
and the corrected comparison expands each difference exactly.
These computations verify finite algebra. The integral
deformation, endomorphism and division-field arguments are
the proofs above. They do not imply that every remaining
tame formal model is potentially CM, nor do they replace
the unresolved non-CM cases by elliptic curves.
\end{remark}


% Separate post-CP18 candidate. No unrestricted least-index conclusion.
\subsection{Residual irreducibility and a cycle of tame character types}

In this subsection ``irreducible residual inertia'' means irreducible
over $\mathbf F_p$ after passing to the completed maximal unramified
base. It does not mean absolutely irreducible over
$\overline{\mathbf F}_p$. The distinction matters: the irreducible
two-dimensional residual representations occurring here are tame
and split over $\mathbf F_{p^2}$.

\begin{lemma}[The leading-coordinate character]
\label{lem:B-RD53-leading}
\status{Proved; full independent proof and exact-code audits passed.}
Let $p\ge5$, and let $K$ be complete discretely valued of mixed
characteristic $(0,p)$ with algebraically closed residue field,
$e_K=p-1$, and $\mu_p\subset K$. Put $N=p^2-1$.
Fix a uniformizer of $K$, the associated compatible tame
fundamental characters, and a generator $\tau$ of tame inertia
whose level-two value is a primitive root $\zeta\in\mu_N$.
Let $M$ be potentially supersingular of height two, and assume
that $M[p]$ is irreducible over $\mathbf F_p$.
If the special tangent character takes the value
$\zeta^c$ at $\tau$, then the residual eigenvalues are
\[
              \bigl\{\zeta^{c+p-1},
                       \zeta^{p(c+p-1)}\bigr\}.       \tag{B-RD53.1}
\]
Here tangent means the action on the one-dimensional Lie space
of the special formal group. The special covariant Dieudonne
character has eigenvalues $\{\zeta^c,\zeta^{pc}\}$.
\end{lemma}
\begin{proof}
The image of wild inertia in $\GL(M[p])$ is a normal finite
$p$-group. Such a group has a nonzero invariant vector: its
orbits on the finite set $M[p]$ show that the number of fixed
vectors is divisible by $p$, and zero is fixed. Its invariant
subspace is stable under all inertia, so irreducibility makes
wild inertia trivial. The remaining image is cyclic. An
irreducible cyclic subgroup of $\GL_2(\mathbf F_p)$ lies in
$\mathbf F_{p^2}^{\times}$: the minimal polynomial of a
generator is irreducible of degree two. Consequently the
$p$-division field is tame of degree dividing $N$.

By Lemma~\ref{lem:B-RD49-inertia}, a faithful good field for
the actual lattice is also tame of degree dividing $N$.
Let $P$ be their compositum and put $t=[P:K]$, so $t\mid N$.
Choose a uniformizer $\varpi$ of $P$ with
$\varpi^t=\pi_K$, and choose the generator $\sigma$ induced
by $\tau$, so that
\[
                  \sigma(\varpi)=\zeta^{N/t}\varpi.
\]
All $p$-torsion points extend to the good finite flat model,
and lie in its one-dimensional formal group because its
special fibre is connected. Tate full faithfulness extends
the generic descent maps and their inverses to this model
\cite[Theorem~4 and Corollaries~1--2, p.~181]
{Tate1967PDivisible}. Thus the Galois action preserves the
valuation of every nonzero formal point: its coordinate
change has a unit linear term and higher terms have larger
valuation.

The points whose coordinate has valuation at least any
fixed bound form an inertia-stable $\mathbf F_p$-subspace.
If the nonzero points had different valuations, one of
these would be a proper nonzero subspace. They therefore
all have a common positive integer valuation $A$.
The Weierstrass polynomial of $[p](T)$ is monic of degree
$p^2$, has precisely these $p^2$ torsion coordinates as
roots, and differs from $[p](T)$ by an integral unit.
Comparing its linear coefficient with $[p]'(0)=p$ gives
\[
                         NA=e_P=(p-1)t,
 \qquad A=\frac{t}{p+1}.                              \tag{B-RD53.2}
\]
In particular $p+1$ divides $t$.
This is the same one-level product-of-torsion-coordinates
calculation used in Theorem~\ref{thm:B-RD12}; it does not
assume a choice of a CM lift.

Taking the leading coefficient of $T(Q)/\varpi^A$ gives
an injective $\mathbf F_p$-linear map $M[p]\to k$.
The formal group addition has ordinary addition as its
leading term. Under $\sigma$, the coordinate's field action
contributes $\zeta^{(N/t)A}=\zeta^{p-1}$ and the descent
map's linear term contributes $\zeta^c$. Therefore this
two-dimensional $\mathbf F_p$-space is acted on by
multiplication by $\alpha=\zeta^{c+p-1}$.
Irreducibility says $\alpha\notin\mathbf F_p$; its two
eigenvalues are $\alpha,\alpha^p$, proving (B-RD53.1).
The last special-character statement is the two embeddings
of the unramified quadratic subfield of the height-two
division algebra, with the covariant convention of
Proposition~\ref{prop:B-RD4-dictionary}.
\end{proof}

\begin{proposition}[Maximal joint order and adjacent character types]
\label{prop:B-RD53-cycle}
\status{Proved; full independent proof and exact-code audits passed.}
In the remaining $p^2$ comparison of Theorem~\ref{thm:B-RD49},
assume the common residual inertia representation is irreducible
over $\mathbf F_p$. Make the common finite height-one twist
that identifies the common finite determinant with the
level-one tame fundamental character of $K$.
If the two requested special characters differ, their joint
faithful good degree is exactly $N=p^2-1$.

The normalized finite types are the Frobenius orbits of
\[
 c_j=1+(p-1)j\quad\text{in }\mathbf Z/N\mathbf Z,
            \qquad j\in\mathbf Z/(p+1)\mathbf Z.
\]
Label each orbit by its unique even representative $j=2k$,
where $k\in\mathbf Z/n\mathbf Z$ and $n=(p+1)/2$.
The two types in any different-character comparison must
be adjacent vertices in the cycle
\[
                    k\longleftrightarrow k+1
                 \quad\text{on }\mathbf Z/n\mathbf Z. \tag{B-RD53.3}
\]
This is a necessary type condition. The statement does not
construct a pair of groups realizing any edge at level $p^2$.
\end{proposition}
\begin{proof}
The common determinant is nonsquare by
Proposition~\ref{prop:B-RD50-square}. Its quotient by the
level-one tame fundamental character is a square finite
$\mu_{p-1}$-valued character, so the indicated common twist
exists. This character is defined from a uniformizer of $K$;
it is not the mod-$p$ cyclotomic character, which is trivial
because $\mu_p\subset K$. A scalar twist preserves residual
irreducibility.

Let $c_1,c_2\in\mathbf Z/N\mathbf Z$ be the actual tangent
exponents of the two good models. The normalization gives
$c_i\equiv1\pmod{p-1}$. Their isomorphic residual modules
and Lemma~\ref{lem:B-RD53-leading} give either $c_2=c_1$ or
\[
                     c_2=pc_1+(p-1)^2\pmod N.          \tag{B-RD53.4}
\]
The first case has equal special characters and is excluded.
Put $d=\gcd(N,c_1,c_2)$, so the joint faithful order is $N/d$.
Equation (B-RD53.4) implies $d\mid(p-1)^2$, whereas
$c_1\equiv1\pmod{p-1}$ gives $\gcd(d,p-1)=1$.
Thus $d=1$ and the joint order is $N$.

On the $j$-indices, ordinary Frobenius is
$j\mapsto1-j$ modulo $p+1$. Because $p+1$ is even, every
orbit has exactly one even index. The affine involution
(B-RD53.4) is $j\mapsto-1-j$. Applied to the even index
$2k$, it gives an odd index whose Frobenius partner is
$2k+2$. Applied to the other representative of the same
orbit, it gives $2k-2$. Hence the possible distinct-type
relations are exactly the adjacent edges in (B-RD53.3).
No orbit equals its neighbor, since $n\ge3$.
The equality of residual eigenvalue pairs for an abstract
edge follows from the same formula; it does not provide a
Barsotti--Tate realization or a $p^2$-torsion congruence.
\end{proof}

\begin{corollary}[Excluded quadratic pairs and the first new restriction]
\label{cor:B-RD53-quadratic}
\status{Proved; full independent proof and exact-code audits passed.}
For $p\equiv3\pmod4$, $p\ge7$, two different opposite
quadratic special types in the remaining $p^2$ test must
have reducible residual inertia over $\mathbf F_p$.
For $p\equiv1\pmod4$, exactly one unordered quadratic pair
among the normalized abstract finite types is an edge in
(B-RD53.3). Also every remaining normalized joint degree
strictly smaller than $p^2-1$ requires reducible residual
inertia.

For example, at $p=7$ the irreducible-residual possibilities
are the adjacent pairs in the following four-cycle:
\[
 \{1,7\},\quad\{13,43\},\quad\{25,31\},\quad\{19,37\}
                 \quad\text{modulo }48,                \tag{B-RD53.5}
\]
including the edge joining the last pair back to the first.
Its two opposite pairs are excluded in this residual case.
At $p=5$ the cycle is a triangle, so this corollary does
not remove either of the two patterns in
Corollary~\ref{cor:B-RD51-p5}.
\end{corollary}
\begin{proof}
A quadratic twist adds $N/2$ to $c$, hence adds
$n=(p+1)/2$ to $j$. If $n$ is even it acts on cycle
indices by $k\mapsto k+n/2$, which is not an edge for
$n\ge4$. If $n$ is odd it acts by the reflection
\[
                    k\longmapsto(1-n)/2-k\pmod n.
\]
Exactly one unordered exchanged pair is adjacent: solving
$2k=(1-n)/2\mp1$ gives the two ends of the same edge,
because two is invertible modulo odd $n$.
The nonmaximal-order assertion is the preceding proposition.
The displayed four-cycle follows by substituting
$p=7$ and the even indices $0,2,4,6$.
\end{proof}

\begin{remark}[Remaining scope]
\label{rem:B-RD53-scope}
\status{Proved scope and limitation; no least-index or pair-existence assertion.}
The new restriction separates irreducible and reducible
residual inertia over $\mathbf F_p$. It does not assert that
the residual representation of every potentially supersingular
formal group is irreducible. In particular, all remaining
non-CM deformations with reducible residual inertia remain
in the problem. The general interval in
Theorem~\ref{thm:B-RD49} is unchanged.
The retained exact program
\path{code/reduction_descent/verify_tame_residual_cycle.py}
checks $9{,}495$ distinct finite-type pairs and $740$
rational valuation cases for every prime from $5$ through
$101$, including the exact four-cycle at $p=7$, the
triangle at $p=5$, and the quadratic-twist parity assertions.
Its author run passed in $0.04175966899492778$ seconds.
The full independent root and companion audits passed, with
author-output-preserving reruns in $0.04199664699990535$ and
$0.047423224001249764$ seconds, respectively.
The finite calculation supplements the proofs; it does not
instantiate a formal group or assert a torsion-congruent pair.
\end{remark}


% Separate post-CP18 candidate. No unrestricted least-index conclusion.
\subsection{A finite-level neighborhood of each canonical CM lift}

\begin{lemma}[Equivariant continuity of formal torsion]
\label{lem:B-RD54-continuity}
\status{Proved; full independent proof and original-source audits passed.}
Let $L$ be complete discretely valued of mixed characteristic $(0,p)$,
with algebraically closed residue field $k$, and write $a=v_L(p)$.
Let $\mathcal F_u$ be a one-dimensional height-two universal formal
deformation over $W(k)[[u]]$, with $\mathcal F_0$ the canonical
unramified quadratic Lubin--Tate formal group. Fix an integer $n\ge1$.
If $x\in\mathfrak m_L$ satisfies
\[
                          v_L(x)>2na,                 \tag{B-RD54.1}
\]
then there is a canonical isomorphism of finite Galois modules
\[
               \mathcal F_0[p^n](\overline L)
                   \simeq\mathcal F_x[p^n](\overline L).
                                                        \tag{B-RD54.2}
\]
It matches roots whose coordinate difference has valuation greater
than $na$.

Suppose moreover that a finite group $\Gamma$ acts on $L$ and on
the universal family, that $u=0$ is fixed, and that the point $u=x$
is equivariant for these actions. The universal family action is
understood to include the isomorphisms of formal groups over the
induced automorphisms of the parameter ring. Then (B-RD54.2)
commutes with the induced semilinear descent actions as well.
\end{lemma}
\begin{proof}
Write $R=W(k)[[u]]$. Weierstrass preparation over the complete
local ring $R$ gives
\[
                  [p^n]_{\mathcal F_u}(T)=U_u(T)W_u(T),
\]
where $U_u(T)$ is an integral unit series and $W_u(T)$ is monic
of degree $p^{2n}$. All its nonleading coefficients are in
$(p,u)$. Uniqueness of preparation shows that its specialization
at $u=0$ is the prepared polynomial for $\mathcal F_0$; hence
every coefficient of $W_u-W_0$ is divisible by $u$.
The same divisibility holds for the difference of the two formal
group laws.

Let $E/L$ be the splitting field of $\mathcal F_0[p^n]$, and
extend $v_L$ to $E$ without changing its normalization. For
every torsion coordinate $\alpha$ in this field,
\[
                       v_L\bigl(W_0'(\alpha)\bigr)=na.
                                                        \tag{B-RD54.3}
\]
Indeed, if $w(T)\,dT$ is an invariant differential with $w(0)=1$,
then
\[
 w([p^n](T))[p^n]'(T)=p^n w(T).
\]
At a torsion point all factors except $p^n$ are units, and the
prepared unit $U_0(\alpha)$ is also a unit. This proves
(B-RD54.3), including at $\alpha=0$.

Put $b=v_L(x)$. Coefficientwise divisibility gives
\[
 v_L(W_x(\alpha))\ge b,
 \qquad v_L(W_x'(\alpha)-W_0'(\alpha))\ge b.
\]
Since $b>2na$, the derivative still has valuation $na$.
The elementary Newton--Hensel argument therefore gives a root
$\beta\in E$ of $W_x$ with
\[
                       v_L(\beta-\alpha)\ge b-na>na.
                                                        \tag{B-RD54.4}
\]
For completeness, after one correction the error valuation
is at least $2(b-na)$; subsequent corrections strictly improve
this bound. Completeness gives a root. In a ball whose difference
valuations exceed $na$, the linear term of the Taylor expansion
strictly dominates every higher term, so this root is unique.
This argument works with the extended rational-valued valuation
on $E$ and introduces no further field extension.

Distinct torsion points of $\mathcal F_0[p^n]$ have coordinate
differences of valuation at most $a/(p^2-1)$. To see this,
translate their difference using the integral formal group law;
its valuation is unchanged by translation. If its exact order
is $p^r$, the canonical Lubin--Tate equation gives valuation
\[
                      \frac{a}{(p^2-1)p^{2(r-1)}}.
\]
In particular this is at most $a/(p^2-1)<na$.
Thus the balls in (B-RD54.4) are disjoint. There are $p^{2n}$
of them, and they give all roots of the separable polynomial
$W_x$.

The matching respects addition. If $\alpha_1,\alpha_2$ match
$\beta_1,\beta_2$, integrality of the formal laws and the
coefficientwise difference divisible by $x$ show that
\[
 \mathcal F_x(\beta_1,\beta_2)
       -\mathcal F_0(\alpha_1,\alpha_2)
\]
has valuation greater than $na$. The first coordinate is again
a $p^n$-torsion point. Uniqueness in the matching ball proves
additivity, and hence (B-RD54.2). The same uniqueness proves
$G_L$-equivariance.

Finally let $g\in\Gamma$. In universal coordinates its formal
isomorphism specializes at $u=0$ to the corresponding automorphism
of $\mathcal F_0$. Their difference has every coefficient
divisible by $u$. Evaluating at the equivariant point $x$,
and using $v_L(gx)=v_L(x)$, shows that applying the two descent
maps to a matched pair again gives a pair with difference
valuation greater than $na$. Uniqueness proves compatibility
with descent. This uses the actual isomorphisms and their
cocycle, not merely the action on isomorphism classes. For the
universal height-two family these are supplied uniquely by
\cite[Propositions~14.7 and~14.9, including (14.8)--(14.10)]
{HopkinsGross1994Bundles}.
\end{proof}

\begin{proposition}[An excluded open neighborhood in the remaining test]
\label{prop:B-RD54-neighborhood}
\status{Proved; full independent proof and original-source audits passed.}
Use the normalized remaining $p^2$ comparison of
Theorem~\ref{thm:B-RD49}, with $p\ge5$ and $e_K=p-1$.
Thus the residue field is algebraically closed, $\mu_p\subset K$,
and the two potentially supersingular height-two groups have
the same nonsquare finite determinant character.

For each group choose a faithful good tame field $L_i$ and
identify its special formal group with the Honda formal group
so that its finite descent subgroup is contained in
$\mu_{p^2-1}\subset\mathbf Z_{p^2}^{\times}$.
Let $x_i\in\mathfrak m_{L_i}$ be its parameter in the universal
deformation based at the canonical unramified quadratic CM lift.
If
\[
                   v_{L_i}(x_i)>4e_{L_i}
                         \qquad(i=1,2),                \tag{B-RD54.5}
\]
where $v(0)=+\infty$, then isomorphic $p^2$-torsion forces
the two requested special characters to be equal.

Consequently every remaining different-character pair has,
for any such two choices of normalized framings, at least one
parameter satisfying
\[
                       0<v_{L_i}(x_i)\le4e_{L_i}.
                                                        \tag{B-RD54.6}
\]
In particular, its normalized deformation parameter cannot
lie arbitrarily deeply in the canonical CM neighborhood on
both sides of the comparison.
\end{proposition}
\begin{proof}
The faithful finite subgroups and integral conjugacy are those
of Lemma~\ref{lem:B-RD49-inertia} and the integral conjugacy
argument in Lemma~\ref{lem:B-RD52-CM}. The latter averaging
argument applies to a finite cyclic subgroup of the special
division order whether or not the given lift has CM.
The canonical lift is fixed by this subgroup. The universal
deformation action, with its unique marked isomorphisms, is
described in
\cite[Propositions~12.10, 14.7, 14.9 and~14.13]
{HopkinsGross1994Bundles}.

Tate full faithfulness extends the generic descent maps to the
actual good models, so their marked universal parameters are
equivariant. At $x_i=0$, the same descent datum gives a
potentially CM generic group $M_i^0/K$: its good formal group
is the canonical unramified quadratic lift, and its descent
maps are the corresponding commuting roots of unity. Generic
finite-level Galois descent is effective and compatible with
the transition maps, so this constructs a height-two
Barsotti--Tate group over $K$. Its finite special character is
the original one.

Lemma~\ref{lem:B-RD54-continuity}, with $n=2$, now gives
\[
                      M_i[p^2]\simeq M_i^0[p^2]
\]
over $K$, including the full descent action. An isomorphism
between the original two truncations therefore gives one
between the two potentially CM truncations. Apply
Theorem~\ref{thm:B-RD52}. This proves (B-RD54.5) and its
contrapositive (B-RD54.6).
\end{proof}

\begin{remark}[Scope of the neighborhood restriction]
\label{rem:B-RD54-scope}
\status{Proved scope and limitation; the remaining finite interval is unresolved.}
The parameter is chosen after specifying the special Honda
framing and its finite descent subgroup; it is not an invariant
of the unframed formal group. The result is quantified for
any such fixed choices and does not identify different
coordinate neighborhoods without their framings.
For a fixed framed CM origin, another integral parameter is
$u_{\mathrm{new}}=u h(u)$ with $h(0)$ a unit; hence its value
has the same valuation. An integral change of formal-group
coordinate also preserves the matching, by the same integral
Lipschitz and coefficient-divisibility argument as in the proof.
It imposes an explicit upper bound on at least one nonzero
parameter in a hypothetical counterexample. It neither
constructs a counterexample within this interval nor proves
the unrestricted $p^2$ character criterion.
No Hasse-thickness estimate is transferred to an isogenous
quotient in this argument. Full independent root and companion
proof and original-source audits passed; no numerical
experiment is asserted for this general continuity lemma.
\end{remark}


% Post-CP19 theorem; complete root and obstruction/crystal audits passed.
\subsection{The smallest wild layer cannot hide two tame characters}

\begin{lemma}[Primitive second-level valuations]
\label{lem:B-RD55-valuations}
\status{Proved and independently checked.}
Let $p\ge5$, let $K$ be complete discretely valued of mixed
characteristic $(0,p)$ with algebraically closed residue field,
$e_K=p-1$, and $\mu_p\subset K$. Let $M$ be a potentially
supersingular height-two group whose residual inertia is
irreducible over $\mathbf F_p$. Put $q=p^2$ and $N=q-1$.
Let $L/K$ be a tame extension of degree $t\mid N$ over which
$M$ has a good formal model and all its $p$-torsion is rational.
Write $W$ for the wild image on $M[p^2]$, and let $F/L$ be
its full $p^2$-division field. Then
\[
       \operatorname{Gal}(F/L)=W\subset1+pM_2(\mathbf F_p),
       \qquad |W|\in\{p^2,p^3,p^4\}.                 \tag{B-RD55.1}
\]
Here the displayed inclusion means the kernel of reduction in
$\GL_2(\mathbf Z/p^2)$, with its usual additive matrix description.
Every nonzero $p$-torsion coordinate has $v_L$-valuation
$a/N$, where $a=(p-1)t$, and every primitive $p^2$-torsion
coordinate has valuation $a/(qN)$.
If $|W|=q$, the latter coordinates have integer $v_F$-valuation
\[
                         A=\frac{t}{p+1},\qquad p\nmid A.
                                                        \tag{B-RD55.2}
\]
\end{lemma}
\begin{proof}
The first-level valuation filtration and product argument are
those of Lemma~\ref{lem:B-RD53-leading}. They give the common
nonzero first-level valuation $\alpha=a/N$.
The negative Newton polygon of the formal $p$-series is therefore
the segment from $(1,a)$ to $(q,0)$: these are exactly its
nonzero roots of positive valuation. In particular its degree-$j$
coefficient has valuation at least $(q-j)\alpha$ for $1\le j<q$.
For a nonzero first-level point $P$, subtract its coordinate
from this series. The constant term now has valuation $\alpha$;
every intermediate coefficient is strictly above the line
from $(0,\alpha)$ to $(q,0)$, since
\[
               (q-j)\alpha>(1-j/q)\alpha\quad(0<j<q).
\]
All $q$ preimages of $P$ consequently have valuation
$\alpha/q$. These are precisely all primitive second-level points.

Wild inertia is trivial on $M[p]$ by residual irreducibility,
as proved in Lemma~\ref{lem:B-RD53-leading}. Its second-level
image is therefore an elementary abelian subgroup of the
four-dimensional additive group $M_2(\mathbf F_p)$.
The remaining quotient of the finite image is tame and is
exactly its residual image: a prime-to-$p$ quotient cannot
have a nontrivial subgroup in this reduction kernel.
Since $L$ contains the residual division field and has tame
degree, adjoining it leaves precisely $W$ as division group.
There is no residue extension, so $[F:L]=|W|$ is the
ramification degree.

The primitive coordinate valuation, normalized integrally on
$F$, is
\[
             \frac{(p-1)t|W|}{(q-1)q}
                  =\frac{t|W|}{(p+1)p^2}.
\]
It is an integer. Since $p\nmid t(p+1)$, this forces
$p^2\mid |W|$, proving (B-RD55.1). If $|W|=p^2$, the formula
is (B-RD55.2); its nonzero integer $A$ divides the prime-to-$p$
integer $t$ and is therefore itself prime to $p$.
\end{proof}

\begin{proposition}[Exclusion of the minimal wild image]
\label{prop:B-RD55-minimal-wild}
\status{Proved and independently checked.}
In the normalized remaining $p^2$ comparison of
Theorem~\ref{thm:B-RD49}, assume that the common residual
inertia representation is irreducible over $\mathbf F_p$.
If the wild image on the common $p^2$-torsion has order $p^2$,
the two requested special characters are equal.
Consequently any different-character pair in this residual
case has wild image of order $p^3$ or $p^4$.

At $p=5$, this excludes order $25$ in both normalized patterns
of Corollary~\ref{cor:B-RD51-p5}; any remaining pair with
irreducible residual inertia has wild order $125$ or $625$.
There is no CM assumption in this assertion.
\end{proposition}
\begin{proof}
Assume that the characters differ. Lemma~\ref{lem:B-RD53-leading}
and Proposition~\ref{prop:B-RD53-cycle} apply. Take the compositum
$L$ of the two faithful good tame fields and the common residual
division field. It has degree $N=p^2-1$ over $K$; in particular
it contains every tame subextension of the common second-level
division field. Put $F=L(M_1[p^2])=L(M_2[p^2])$.
By hypothesis $\operatorname{Gal}(F/L)=W$ has order $q=p^2$.
Choose a uniformizer $\varpi$ of $F$. By the preceding lemma,
every primitive second-level coordinate in either model has
integer valuation $A=p-1$, while every nonzero first-level
coordinate has valuation $qA$.

Choose Honda coordinates on both special formal groups, so
their special $p$-series is $T^q$. For the actual second-level
module $R_i=M_i[p^2](F)$ put $V_i=R_i/pR_i$. Taking the leading
coefficient of $T_i(P)/\varpi^A$ induces an injective
$\mathbf F_p$-linear map
\[
                         \ell_i:V_i\longrightarrow k.
                                                        \tag{B-RD55.3}
\]
Indeed, two primitive points have the same leading coefficient
if and only if their difference has valuation larger than $A$,
which is equivalent to their difference lying in $pR_i$.
Formal addition has ordinary addition as its leading term.
The leading coefficient of $T_i(pP)/\varpi^{qA}$ is
$\ell_i(P)^q$: the degree-$q$ term of $[p]$ is the unique
lowest-valuation term at a primitive coordinate, by the Newton
polygon calculation above, and its residue coefficient is one.
The first-level leading space is a scalar multiple of
$\mathbf F_q$, by residual irreducibility and
Lemma~\ref{lem:B-RD53-leading}. Hence the same holds for
the image in (B-RD55.3). Write it as $d_i\mathbf F_q$,
with $d_i\ne0$.

Let $g\in W$ be nontrivial, and define its lower break $b_g$ by
\[
 v_F(g\varpi-\varpi)=b_g+1,\qquad
 d_g=\overline{(g\varpi-\varpi)/\varpi^{b_g+1}}\ne0.
\]
Choose a coefficient ring $\Lambda\subset\mathcal O_K$ by
\citet[Theorem~10.160.8 and proof]{Stacks032A}, and fix a
representative $[c]\in\Lambda$ for each residue $c\in k$.
Every $g$ fixes these representatives. No canonical embedding
of the Witt ring is required. For every element $x$ of valuation
$A$, its uniformizer expansion with these fixed coefficients gives
\[
 \begin{split}
 v_F(gx-x)&=A+b_g,\\
 \overline{(gx-x)/\varpi^{A+b_g}}
       &=\overline A\,d_g\,\overline{x/\varpi^A}.
 \end{split}                                            \tag{B-RD55.4}
\]
To check the possible unit error explicitly, expand
$x=\sum_{j\ge A}[a_j]\varpi^j$. The degree-$A$ term changes
by its nonzero linear binomial term because $p\nmid A$.
Every higher degree-$j$ term changes with valuation at least
$j+b_g>A+b_g$. Thus none can cancel the displayed term.
Equivalently, a unit changes with valuation at least $b_g+1$.

Write the action of $g$ on $R_i$ as $1+pB_i(g)$, where
$B_i(g)\in\operatorname{End}_{\mathbf F_p}(V_i)$.
For every primitive point its displacement is a first-level
point. Equation (B-RD55.4) says that this displacement is
never zero. Comparing valuations gives
\[
                    b_g=(q-1)A\quad(g\ne1).             \tag{B-RD55.5}
\]
Formal subtraction has the same leading coefficient as ordinary
coordinate subtraction, since their quotient is an integral
unit with residue one. Using the $p$-series leading term now gives
\[
          \ell_i(B_i(g)v)^q=\overline A\,d_g\,\ell_i(v).
                                                        \tag{B-RD55.6}
\]
Identify $V_i$ with $\mathbf F_q$ by $\ell_i(v)=d_i z$.
Since the $q$th power fixes $\mathbf F_q$, equation (B-RD55.6)
says that $B_i(g)$ is multiplication by
\[
           \lambda_i(g)=\kappa_i d_g\in\mathbf F_q,
           \qquad \kappa_i=\overline A\,d_i^{1-q}.
                                                        \tag{B-RD55.7}
\]
The map $g\mapsto B_i(g)$ is injective on $W$, so
$\lambda_i(W)=\mathbf F_q$; define $d_g=0$ when $g=1$. Thus both models
give the same $q$-element leading ramification space, multiplied
by fixed nonzero constants $\kappa_i$.

The isomorphism of the second-level modules induces the same
linear isomorphism on $V_i$ and on $pR_i$ under multiplication
by $p$. In the two leading-coordinate identifications these
are related by $z\mapsto z^q$, which is the identity on
$\mathbf F_q$: the first-level leading coefficient is normalized
by $d_i^q$. Thus the leading-coordinate convention and tame
generator are exactly those of Lemma~\ref{lem:B-RD53-leading}.
Because the special characters differ,
Lemma~\ref{lem:B-RD53-leading} says that the residual intertwiner
must be Frobenius-semilinear, of the form
$z\mapsto c z^p$ with $c\in\mathbf F_q^\times$; the linear
case would identify the actual tangent exponents and hence
the special characters.
Commuting with every $g\in W$ consequently requires
\[
                 \lambda_2(g)=\lambda_1(g)^p.
\]
But (B-RD55.7) also gives
$\lambda_2(g)=(\kappa_2/\kappa_1)\lambda_1(g)$.
Since $\lambda_1(W)=\mathbf F_q$, this would make the
nontrivial Frobenius automorphism of $\mathbf F_q/\mathbf F_p$
a scalar multiplication on the entire field. Evaluating at
one forces that scalar to be one, and an element outside
$\mathbf F_p$ gives a contradiction.

This excludes $|W|=p^2$. The preceding lemma leaves exactly
$p^3$ or $p^4$, with the stated numerical specialization at five.
\end{proof}

\begin{remark}[Remaining scope]
\label{rem:B-RD55-scope}
\status{Proved scope statement; neither remaining wild order is excluded.}
The new restriction concerns the actual wild image on the
common second-level Galois module. It applies to non-CM as
well as CM deformations. It does not identify a group solely
from its abstract image order, and it does not assert existence
of a different-character pair of either larger order.
Reducible residual representations, and the two larger wild
orders in the irreducible case, remain unresolved here.
The exact finite supplement
\path{code/reduction_descent/verify_tame_minimal_wild_layer.py}
checks $3{,}700$ valuation cases, $166$ minimal integral cases,
$75{,}984$ field elements and separately $75{,}984$ scalar
candidates, and $1{,}180$
distinct-type orientation pairs for the primes from $5$ through
$101$. The author run passed in $0.5875725139994756$ seconds.
The independent root and obstruction/crystal runs passed in
$0.6879981279998901$ and $0.6387307759941905$ seconds,
respectively, without changing the author output. All three
statements also received complete independent mathematical audits.
The calculation does not instantiate a formal group or a congruent pair.
\end{remark}


% Post-CP20 theorem; complete independent root and obstruction/crystal audits passed.
\subsection{A cubic wild layer and the residual projective order}

\begin{lemma}[The anti-Cartan layer and exact tame coordinates]
\label{lem:B-RD56-anti-layer}
\status{Proved; complete independent proof and exact-code audits passed.}
Use the normalized remaining $p^2$ comparison of
Theorem~\ref{thm:B-RD49}, with $p\ge5$, over a complete discretely
valued mixed-characteristic $(0,p)$ field $K$ with algebraically
closed residue field, $e_K=p-1$, and $\mu_p\subset K$.
Assume that the common residual representation is irreducible over
$\mathbf F_p$, that its projective image has order $h>2$, and that
the two special characters differ. Put $q=p^2$ and $N=q-1$.
If its common second-level wild image $W$ has order $p^3$, then
$W$ contains a distinguished two-dimensional subspace $U$ of
anti-Cartan matrices. Conjugation by tame inertia acts irreducibly
on $U$. All its nonzero elements have the same lower break $b$.

Take the common tame good field $L/K$ of degree $N$ and the common
second-level division field $F/L$. One can choose a tame complement
$\langle\sigma\rangle\subset\operatorname{Gal}(F/K)$, a uniformizer
$\varpi$ of $F$, and formal coordinates $T_i$ for the two good
models such that, for a primitive $N$th root $\zeta$ and their
actual tangent exponents $c_i$,
\[
 \begin{gathered}
 \sigma\varpi=\eta\varpi,\qquad \eta=\zeta^p,\qquad
 T_i(\sigma P)=\zeta^{c_i}\sigma(T_i(P)),\\
 A=v_F(T_i(P))=p(p-1)\quad(P\text{ primitive of order }p^2),\\
 v_F(T_i(Q))=qA\quad(0\ne Q\in M_i[p]),\\
                         e_F=NqA.                       \tag{B-RD56.1}
 \end{gathered}
\]
The special coordinates remain Honda coordinates with $[p](T)=T^q$.
\end{lemma}
\begin{proof}
Identify either residual space with $\mathbf F_q$, so its tame
image acts by multiplication by $\lambda\in\mathbf F_q^\times$.
Every endomorphism is uniquely
\[
                         z\longmapsto az+bz^p.
\]
Conjugation fixes the two-dimensional Cartan space of $a$'s and
acts on the two-dimensional anti-Cartan space of $b$'s by
$\lambda^{1-p}$. This element has order $h>2$; since its order
divides $p+1$, it is not in $\mathbf F_p$. Thus the anti-Cartan
space is irreducible over $\mathbf F_p$. The conjugation action
has prime-to-$p$ order and is semisimple by averaging. Its invariant
three-dimensional subspace $W$ must consequently be the direct
sum of this anti-Cartan space $U$ and a Cartan line.
The lower ramification filtration is invariant under conjugation.
Its intersections with the irreducible space $U$ are either all
of $U$ or zero, giving a single break $b$ for its nonzero elements.

Proposition~\ref{prop:B-RD53-cycle} gives the common tame good
field of degree $N$. The argument of Lemma~\ref{lem:B-RD55-valuations}
gives $\operatorname{Gal}(F/L)=W$ and the displayed valuations.
The extension of $C_N$ by the abelian group $W$ has a complement:
if a lift $s$ of a generator has $s^N=w\in W$, then $w$ is fixed
by $s$, and replacing $s$ by $s(-N^{-1}w)$ makes its $N$th power
one. Here $W$ is written additively and $N$ is invertible modulo $p$.
Choose the generator $\sigma$ to act on a tame uniformizer of $L$
by $\zeta$. Its multiplier $\eta$ on a uniformizer of $F$ satisfies
$\eta^{p^3}=\zeta$, hence $\eta=\zeta^p$. Averaging a uniformizer
against this character makes $\sigma\varpi=\eta\varpi$ exact,
without changing its leading term.

The actual semilinear descent isomorphisms extend to the good models
by Tate full faithfulness, as in Lemma~\ref{lem:B-RD53-leading}.
Choose the special Honda framings in which the finite tame actions
are $T\mapsto\zeta^{c_i}T$. Averaging any integral lift of that
coordinate against this character over the cyclic group of order
$N$ gives the exact semilinear eigen-coordinate in (B-RD56.1).
Its special reduction is unchanged, so its special $p$-series is
still $T^q$. This uses the actual descent cocycle; it is not an
assumption that an arbitrary deformation has CM.
\end{proof}

\begin{proposition}[Exclusion of order $p^3$ when the projective order exceeds two]
\label{prop:B-RD56-cubic-wild}
\status{Proved; complete independent proof and exact-code audits passed.}
Under the normalized hypotheses of the preceding lemma, if the
common residual projective image has order greater than two,
isomorphic $p^2$-torsion with different special characters must
have wild image of order $p^4$. In particular wild order $p^3$
is excluded, in addition to the order $p^2$ already excluded by
Proposition~\ref{prop:B-RD55-minimal-wild}.
The same exclusion holds when the residual projective order is two,
provided that $W$ contains the entire anti-Cartan plane and all its
nonzero elements have one common lower break.
\end{proposition}
\begin{proof}
Suppose instead that $|W|=p^3$ and first use the preceding lemma.
For the additional projective-order-two assertion, take the full
anti-Cartan plane $U$ assumed in the statement. The same common field,
valuation and exact-coordinate construction applies: only the lemma's
automatic existence and single-break assertion used $h>2$.
Fix a coefficient ring in $\mathcal O_K$ and residue representatives
as in (B-RD55.4). Since $e_F=NqA>qA$, reduction modulo
$\varpi^{qA+1}$ identifies these coefficients with the field $k$:
all errors in their addition and multiplication are divisible by
$p$ and vanish in this quotient. We may therefore calculate with
characteristic-$p$ power series to this precision.

Let $R_i=M_i[p^2](F)$ and $V_i=R_i/pR_i$. Primitive leading
coefficients give $\ell_i:V_i\hookrightarrow k$ with image
$d_i\mathbf F_q$; multiplication by $p$ raises leading coefficients
to the $q$th power, exactly as in the proof of
Proposition~\ref{prop:B-RD55-minimal-wild}. Write $\ell_i(v)=d_i z$.
Changing a representative of $v$ adds a first-level point, so it
changes its formal coordinate only in valuation at least $qA$.
Thus, for all $j<qA$, the coefficients in
\[
 T_i(P)=d_i z\varpi^A+
                \sum_{A<j<qA} f_{i,j}(z)\varpi^j
                     \pmod{\varpi^{qA}}                \tag{B-RD56.2}
\]
are well-defined functions of $z\in\mathbf F_q$.
No additivity of every coefficient function is assumed.

For $g\in U\setminus\{0\}$, put
\[
 d_g=\overline{(g\varpi-\varpi)/\varpi^{b+1}}.
\]
Composing two such expansions adds their leading coefficients;
the kernel is $U\cap G_{b+1}=0$. Thus $g\mapsto d_g$ is
injective and additive on $U$. Applying
$\sigma g\sigma^{-1}$ to $\varpi$ shows directly that tame
conjugation multiplies the leading coefficient by $\eta^b$.
This realizes the two-dimensional conjugation module $U$ inside $k$.
When $h>2$ that module is irreducible; when $h=2$ conjugation acts
on the entire anti-Cartan plane by the scalar $-1$. In either case
its action is nontrivial, so
\[
                         \eta^b\ne1.                  \tag{B-RD56.3}
\]
Put $b_0=N(p-1)$. Equation (B-RD56.3) implies $b\ne b_0$.
We use $b_0$ only as this numerical comparison value; no identification
of a lower break with a cyclotomic upper break is made.

Every nonzero anti-Cartan matrix $z\mapsto\beta z^p$ is invertible.
Consequently $gP-P$ is a nonzero first-level point for every primitive
$P$. Formal subtraction differs from ordinary coordinate subtraction
by a unit with residue one, so
\[
                        v_F(gT_i(P)-T_i(P))=qA.          \tag{B-RD56.4}
\]
For a monomial of degree $j$, if $p^s$ is the exact power of $p$
dividing $j$, the first term of its change has degree
$j+p^s b$ and coefficient $(j/p^s)d_g^{p^s}$ modulo $p$.
This follows by writing $(\varpi+d_g\varpi^{b+1}+
\text{higher terms})^j$ in characteristic $p$.
Terms divisible by $p$ in mixed characteristic have valuation
strictly larger than $qA$ and cannot affect this calculation.

Suppose first that $b<b_0$. The degree-$A$ monomial in (B-RD56.2)
first changes in degree $A+pb<qA$, with nonzero coefficient
$(p-1)d_i z d_g^p$. Every other monomial whose degree is divisible
by $p$ first changes strictly later. Among monomials of degree
prime to $p$, take the first with a nonzero coefficient. If its
degree were less than
\[
                           j_0=A+(p-1)b,
\]
its change would be the unique earlier nonzero term, contradicting
(B-RD56.4). A later degree cannot cancel the degree-$A$ term.
Hence $p\nmid j_0$, and cancellation forces, for every $z$,
\[
             j_0 f_{i,j_0}(z)
                     =-(p-1)d_i d_g^{p-1}z.            \tag{B-RD56.5}
\]
Fix any $z\ne0$. The left side of (B-RD56.5) is independent of
$g$, whereas $d_i z\ne0$. Hence $d_g^{p-1}$ would have to be the
same nonzero constant for all $g\in U\setminus\{0\}$.
The injective map $g\mapsto d_g$ has $q-1>p-1$ distinct nonzero
values, but a polynomial $X^{p-1}-c$ has at most $p-1$ roots.
This contradiction excludes $b<b_0$. It needs neither a CM assumption
nor a choice of a one-dimensional character inside $U$.

Suppose now that $b>b_0$. Every monomial of degree divisible
by $p$ changes in degree at least $A+pb>qA$.
Equation (B-RD56.4) therefore forces the first nonzero coefficient
of prime-to-$p$ degree to have
\[
                         j=qA-b,\qquad p\nmid j.
\]
Here $A<j<qA$ is necessary; if this fails, (B-RD56.4) is
already impossible. Its leading displacement is
$j f_{i,j}(z)d_g$.
Write the matrix of $g$ on $R_i$ as $1+pB_i(g)$.
On $U$ it has $B_i(g)(z)=\beta_i(g)z^p$, and
$\beta_i:U\to\mathbf F_q$ is an additive isomorphism.
The first-level leading coefficient comparison gives
\[
               d_i^q\beta_i(g)z^p
                           =j f_{i,j}(z)d_g.            \tag{B-RD56.6}
\]
Fixing one nonzero $g$ shows that $f_{i,j}(z)$ is a fixed
nonzero multiple of $z^p$. Applying the equation to every
$g\in U$ then gives constants $\kappa_i\ne0$ with
\[
                          \beta_i(g)=\kappa_i d_g.
                                                        \tag{B-RD56.7}
\]

Because the special characters differ, the common residual
intertwiner in these leading coordinates is $z\mapsto c z^p$,
by Lemma~\ref{lem:B-RD53-leading}. Multiplication by $p$
identifies the induced maps on $V_i$ and $pR_i$, since the
normalized leading-coordinate map is $z\mapsto z^q=z$.
Conjugation of the anti-Cartan matrices consequently gives
\[
                      \beta_2(g)=c^{1-p}\beta_1(g)^p.
\]
Together with (B-RD56.7), this makes Frobenius a fixed scalar
multiplication on $\beta_1(U)=\mathbf F_q$, an impossibility.
This excludes $b>b_0$ and proves the proposition.
\end{proof}

\begin{corollary}[A smaller wild list at five]
\label{cor:B-RD56-p5}
\status{Proved; complete independent proof and exact-code audits passed.}
In the $p=5,e_K=4$ remaining test with irreducible residual
inertia, the pattern $\{A,B\}$ of Corollary~\ref{cor:B-RD51-p5}
requires wild order $625$. The pattern $\{A,C\}$ still has
only the necessary wild list $125$ or $625$ from B55.
Neither remaining case is asserted to exist.
\end{corollary}
\begin{proof}
For $\{A,B\}$, the actual tangent exponents giving equal
residual characters are $1,21$, up to interchange and simultaneous quadratic twist. The residual eigenvalues have order $24$
and their projective ratio has order $6>2$.
For $\{A,C\}$, the equal-residual tangent exponents are $5,17$
or their simultaneous variants; the residual eigenvalues have
order $8$ and projective ratio order $2$.
Apply the proposition in the first case and retain the exact
scope of Proposition~\ref{prop:B-RD55-minimal-wild} in the second.
\end{proof}

\begin{remark}[Scope and numbering caution]
\label{rem:B-RD56-scope}
\status{Proved scope and limitation.}
The argument concerns the actual irreducible residual representation
and the common second-level wild Galois module. It does not assume
that the groups have CM or that one is a full Tate-module twist of
the other. Reducible residual representations and wild order $p^4$ are not
excluded here. At projective order two, the new assertion excludes
only an entire anti-Cartan plane with a single common lower break;
a one-dimensional anti-Cartan part or two distinct anti-Cartan
breaks are not settled by this statement. The intermediate outline's
suggestion to identify a Cartan lower break with a cyclotomic upper
break is not used: the proof works entirely with lower breaks and the
explicit comparison number $b_0=(p^2-1)(p-1)$.
The exact finite supplement
\path{code/reduction_descent/verify_tame_cubic_wild_layer.py}
checks $2{,}498$ monomial first terms, $2{,}494$ lower-break cases,
$2{,}336$ tame-weight cases, and $45{,}864$ explicit anti-Cartan
conjugacy cases. It reproduces all six oriented normalized cases
at $p=5$. The original author run passed in $0.4074130669978331$ seconds.
Its source and output are retained before the additional low-break
cardinality checks. The expanded run passed in
$0.4129977519987733$ seconds, additionally checking the nonconstancy
of $d^{p-1}$ on $360$ nonzero elements in four quadratic finite fields.
The independent root and companion expanded runs passed in
$0.40707763800310204$ and $0.37780446800024947$ seconds,
respectively, preserving the author output. Both agents read and
rederived the complete strengthened proof. This is finite arithmetic
support, not a formal-group construction or a pair-existence computation.
\end{remark}


% Proved after full root and companion audits; exact candidate retained under notes/.
\subsection{Exhausting the cubic wild image in the irreducible residual case}

Retain the normalized hypothetical different-character comparison of
Theorem~\ref{thm:B-RD49}. The field is complete discretely valued of
mixed characteristic $(0,p)$, with algebraically closed residue field,
$p\ge5$, $e_K=p-1$, and $\mu_p\subset K$. Throughout this subsection
the common residual inertia is irreducible over $\mathbf F_p$ and the
common second-level wild image has order $p^3$. Use the common fields
and exact coordinates constructed in B56, and put
\[
 q=p^2,\quad N=q-1,\quad A=p(p-1),\quad b_0=N(p-1),
 \qquad e_F=NqA.
\]
On the residual space identified with $\mathbf F_q$, let $Z$ denote
the Cartan plane of maps $z\mapsto az$, and let $U$ denote the
anti-Cartan plane of maps $z\mapsto az^p$. Both planes are intrinsic
once the nonsplit Cartan field is specified; the Frobenius change
between the two leading-coordinate orientations preserves them.

\begin{lemma}[The first degree prime to $p$]
\label{lem:B-RD57-first-term}
\status{Proved; full independent root and companion proof/source/code audits passed.}
Fix one of the two models, and let $g\in W$ have a nonzero matrix
$B(g)$ belonging either to $Z$ or to $U$, in the notation
$g=1+pB(g)$. Let $b$ be its lower break.
For each primitive leading coordinate $d z$, $z\in\mathbf F_q^\times$,
expand its point coordinate in the chosen uniformizer of $F$.
If $b<b_0$, the first nonzero term of degree prime to $p$ below
$qA$ has degree
\[
                     j=A+(p-1)b,
\]
and its coefficient is a fixed nonzero multiple of $z$.
In particular $p\nmid b$.
If $b>b_0$, that first degree is
\[
                         j=qA-b.
\]
Its coefficient is a nonzero multiple of $z$ when $B(g)\in Z$,
and a nonzero multiple of $z^p$ when $B(g)\in U$.
Each asserted first degree and coefficient formula holds for every
$z\ne0$, in the same fixed coordinate expansion.
\end{lemma}
\begin{proof}
Every nonzero matrix in either plane is invertible. Thus $gP-P$
is a nonzero first-level point for every primitive $P$, and its
coordinate displacement has valuation exactly $qA$.
Work modulo $\varpi^{qA+1}$; this is characteristic $p$ because
$e_F=NqA>qA$. A monomial of degree $j$ first changes in degree
$j+p^{v_p(j)}b$. The leading degree is $A$, with $v_p(A)=1$;
all other degrees divisible by $p$ first change strictly later
than $A+pb$.

If $b<b_0$, the leading change in degree $A+pb<qA$ must cancel.
The first prime-to-$p$ degree can be neither earlier nor later
than $j=A+(p-1)b$. Its coefficient satisfies exactly
(B-RD56.5), for this fixed $g$, and is therefore a nonzero
multiple of $z$. Since $j\equiv-b\pmod p$, it also gives
$p\nmid b$.

If $b>b_0$, all terms of degree divisible by $p$ change after
$qA$. The first prime-to-$p$ term must therefore have degree
$qA-b$. Its coefficient is determined by the leading displacement
of $pB(g)P$. Multiplication by $p$ raises primitive leading
coefficients to the $q$th power. After the fixed normalizations,
this preserves the coordinate $z\in\mathbf F_q$. Consequently
that coefficient is proportional to $az$ for $B(g)(z)=az$,
and to $az^p$ for $B(g)(z)=az^p$, proving the claim.
These are assertions about the same first nonzero term, not
about independently changeable formal coordinates for each $g$.
\end{proof}

\begin{lemma}[A full anti-Cartan plane cannot have two breaks]
\label{lem:B-RD57-anti}
\status{Proved; full independent root and companion proof/source/code audits passed.}
In residual projective order two, if $U\subset W$, all nonzero
elements of $U$ have one common lower break. Hence this case is
excluded for a different-character pair by
Proposition~\ref{prop:B-RD56-cubic-wild}.
\end{lemma}
\begin{proof}
The tame complement acts on $U$ by $-1$. If a nonzero element
has break $b$, its leading ramification coefficient transforms
by $\eta^b$, where $\eta$ is the primitive $N$th root in
(B-RD56.1). Additivity of the leading coefficient gives
\[
                         \eta^b=-1,
              \qquad b\equiv N/2\pmod N.                \tag{B-RD57.1}
\]
In particular no anti-Cartan break equals $b_0$.

Suppose that $g_1,g_2\in U\setminus\{0\}$ have breaks
$b_1<b_2$. If $b_1<b_0$, the first lemma supplies the first
prime-to-$p$ degree $j_1=A+(p-1)b_1$, with coefficient a
nonzero multiple of $z$. For the action of $g_2$, that term
changes in degree $j_1+b_2$, strictly before all terms whose
degree is divisible by $p$, since
\[
                         j_1+b_2<A+pb_2.
\]
It is the unique first change. Its required valuation $qA$
would make the leading matrix $B(g_2)$ a nonzero Cartan-linear
map, contradicting $B(g_2)\in U$. The functions $z$ and $z^p$
on $\mathbf F_q$ are linearly independent over $k$.
If instead $b_1>b_0$, the first prime-to-$p$ degree is
$j_1=qA-b_1$. Under $g_2$ its change and every later such
term have valuation greater than $qA$; every degree divisible
by $p$ does too. This contradicts the nonzero first-level
displacement. These cases exhaust $b_1\ne b_0$ and prove
that only one break occurs. B56 then excludes the pair.
\end{proof}

\begin{lemma}[The remaining one-dimensional anti-Cartan part]
\label{lem:B-RD57-Cartan}
\status{Proved; full independent root and companion proof/source/code audits passed.}
In residual projective order two, the only remaining possible
three-dimensional $W$ has the form
\[
                             W=Z\oplus U_1,
\]
where $U_1$ is a line in $U$. If $b$ is the break of
$U_1\setminus\{0\}$, then
\[
 \begin{gathered}
 0<b<b_0,\qquad b\equiv N/2\pmod N,\qquad p\nmid b,\\
 c=p b_0-(p-1)b>b_0                                      \tag{B-RD57.2}
 \end{gathered}
\]
is the single common lower break of all nonzero elements of $Z$.
At $p=5$ the necessary pairs $(b,c)$ are
\[
                       (12,432),\quad(36,336),\quad(84,144).
                                                        \tag{B-RD57.3}
\]
These are necessary ramification patterns for one model under the
stated image assumptions, not existence assertions.
\end{lemma}
\begin{proof}
Tame conjugation has eigenspaces $Z$ and $U$, with eigenvalues
$1$ and $-1$. Since $p$ is odd it splits every invariant subspace.
A three-dimensional $W$ therefore either contains the entire $U$,
or is $Z\oplus U_1$. The former has just been excluded.
All nonzero elements of the one-dimensional $U_1$ have the same
break $b$, and (B-RD57.1) applies. A nonzero Cartan element of
break $c$ instead has $\eta^c=1$, so $N\mid c$. In particular
no Cartan break equals $b$.

We first show $b<b_0$. Suppose $b>b_0$. The anti-Cartan element
makes the first prime-to-$p$ degree $j=qA-b$, with coefficient
a nonzero multiple of $z^p$. Choose any nonzero Cartan element
and call its break $c$. If $c<b_0$, the first lemma forces
that same first coefficient to be a nonzero multiple of $z$,
a contradiction. If $b_0\le c<b$, the degree-$j$ term changes
before $qA$, whereas all degrees divisible by $p$ change at or
after $qA$. This contradicts the prescribed displacement.
If $c>b$, all changes have valuation greater than $qA$.
The excluded equality $c=b$ leaves no possibility for $c$.
Thus $b>b_0$ is impossible, and $b=b_0$ was already excluded.

For $b<b_0$, the first prime-to-$p$ degree is
$j=A+(p-1)b$, with coefficient a nonzero multiple of $z$.
A Cartan break $c<b$ would force a strictly earlier such degree,
contradicting uniqueness of this first term. Equality is excluded
by tame characters. For $c>b$, the degree-$j$ term changes
strictly before every term of degree divisible by $p$.
Its required displacement valuation $qA$ therefore forces
\[
                   c=qA-j=p b_0-(p-1)b>b_0.
\]
This holds for every nonzero Cartan element, proving the single-break
assertion. The first lemma also gives $p\nmid b$.
For $p=5$, the positive integers below $96$ congruent to $12$
modulo $24$ are $12,36,60,84$. The condition $5\nmid b$
removes $60$; substitution gives precisely (B-RD57.3).
\end{proof}

\begin{theorem}[Only the full wild matrix space remains]
\label{thm:B-RD57}
\status{Proved; full independent root and companion proof/source/code audits passed.}
In the normalized remaining $p^2$ comparison, assume the common
residual inertia is irreducible over $\mathbf F_p$.
If the special characters differ, the wild image on the common
$p^2$-torsion has order $p^4$, and hence equals the entire
additive matrix space $M_2(\mathbf F_p)$.
In particular at $p=5$ both normalized different-type patterns
require wild order $625$; neither $25$ nor $125$ remains.
\end{theorem}
\begin{proof}
B55 excludes order $p^2$, and B56 excludes order $p^3$ at residual
projective order greater than two. An irreducible residual cyclic
image cannot have projective order one: that would make it scalar
over $\mathbf F_p$. It remains to exclude the projective-order-two
case with $|W|=p^3$.

By the preceding lemmas its only possible form is $Z\oplus U_1$,
with the full Cartan plane $Z$ having one break $c>b_0$.
For the common uniformizer of $F$ write
\[
 d_g=\overline{(g\varpi-\varpi)/\varpi^{c+1}}
                         \quad(g\in Z).
\]
This is an injective additive map on $Z$. In the leading coordinates
of model $i$, write $B_i(g)(z)=\beta_i(g)z$.
The map $\beta_i:Z\to\mathbf F_q$ is an additive isomorphism.
The high-break coefficient calculation at $j=qA-c$ gives
\[
                     d_i^q\beta_i(g)z
                              =j f_{i,j}(z)d_g.          \tag{B-RD57.4}
\]
Fixing one nonzero $g$ identifies $f_{i,j}(z)$ as a fixed nonzero
multiple of $z$. Varying $g$ then gives constants $\kappa_i\ne0$
with $\beta_i(g)=\kappa_i d_g$ for every $g\in Z$.

The different special characters force the common residual
intertwiner to have the form $z\mapsto a z^p$ by B53.
The same map on the first-level submodule is obtained from
multiplication by $p$, with normalized coordinates $z^q=z$.
Conjugating scalar Cartan multiplication therefore gives
\[
                            \beta_2(g)=\beta_1(g)^p.
\]
Together with their two fixed multiples of $d_g$, this would make
Frobenius scalar-linear on $\beta_1(Z)=\mathbf F_q$, impossible.
Thus order $p^3$ is excluded. Lemma~\ref{lem:B-RD55-valuations}
leaves only order $p^4$. A subgroup of that order in the
four-dimensional additive space is the whole space.
\end{proof}

\begin{remark}[Remaining exact gate]
\label{rem:B-RD57-scope}
\status{Proved scope and limitation; no full-wild-layer or pair-existence assertion.}
This excludes an image size, not all second-level ambiguities.
The full four-dimensional wild image and all reducible residual
representations remain unresolved. No potentially CM hypothesis,
quotient-thickness transfer, or identification of lower and upper
ramification numbering is used. The p=5 bound on the manuscript
least index therefore remains $\nu_{\min}(5,4,2)\in\{1,2\}$.
All coordinate comparisons are performed in the same actual division
field and in fixed exact tame coordinates; the functions $z$ and
$z^p$ are never identified merely because they span the same abstract
residual representation.
The exact finite supplement
\path{code/reduction_descent/verify_tame_cubic_wild_exhaustion.py}
checks $942$ first-term cases, $100{,}962$ anti-plane break comparisons,
$214{,}860$ Cartan/anti break comparisons, $45{,}864$ finite Cartan
conjugacies, and $364$ scalar/Frobenius comparisons. It also enumerates
all $156$ hyperplanes in the four-dimensional matrix space at $p=5$
and identifies the twelve invariant ones of the two stated shapes.
The expanded author run passed in $0.19211209900095128$ seconds;
the independent root run passed in $0.20058958700246876$ seconds, and
the independent companion run in $0.18477749200246762$ seconds.
The exact prepromotion proof and earlier first program version and output
are retained. These are finite
checks of the written proof, not model or pair constructions.
\end{remark}


% Separate post-CP20 candidate. Full audits required before promotion.
\subsection{The final ramification group in the full wild image}

Throughout this subsection $K$ is complete discretely valued of mixed
characteristic $(0,p)$, $p\ge5$, with algebraically closed residue field,
$e_K=p-1$ and $\mu_p\subset K$. We use the normalized hypothetical
different-character comparison of Theorem~\ref{thm:B-RD49}, and assume
that its common residual inertia is irreducible over $\mathbf F_p$.
Set $q=p^2$, $N=q-1$ and $b_0=N(p-1)$. By
Theorem~\ref{thm:B-RD57}, its remaining wild image is
$W=M_2(\mathbf F_p)$, written additively through $1+pB$.
The common tame good field $L/K$ has degree $N$, and the common
division field $F/L$ has group $W$.

\begin{lemma}[Point fields and the final Cartan plane]
\label{lem:B-RD58-final-Cartan}
\status{Proved; full independent root and companion proof, original-source, and exact-code audits passed.}
For every primitive second-level point $P$, the field $L(P)/L$ is
Galois of degree $q$ and has a single lower and upper ramification
break $b_0$. The last nonzero upper group of $F/L$ is
\[
                         W^{b_0}=Z,                    \tag{B-RD58.1}
\]
where $Z$ is the full nonsplit Cartan plane. Consequently all nonzero
Cartan elements have a single common lower break $c>b_0$, and every
element outside $Z$ has lower break strictly less than $c$.
\end{lemma}
\begin{proof}
Fix a primitive $P$ and let $v$ be its nonzero image in
$M[p^2]/pM[p^2]$. Its stabilizer in $W$ is
\[
                   H_v=\{B:Bv=0\},
\]
a two-dimensional subspace. Since $W$ is abelian, its fixed field
$E=L(P)$ is Galois over $L$, of degree $q$.
Lemma~\ref{lem:B-RD55-valuations} gives
\[
 v_E(T(P))=p-1,\qquad
 v_E(T(Q))=q(p-1)\quad(0\ne Q\in M[p]).
\]
Every nonidentity automorphism of $E/L$ displaces $P$ by a nonzero
first-level point. Formal subtraction has the same valuation as
ordinary coordinate subtraction. Expand $T(P)$ in a uniformizer
of $E$, using residue representatives fixed by the Galois group as
in (B-RD55.4). Its valuation $p-1$ is prime to $p$, so an
automorphism of lower break $b$ changes it in valuation $p-1+b$.
It follows that $b=q(p-1)-(p-1)=b_0$ for every nonidentity
automorphism. A single lower break is also the single upper break.

We now use the quotient compatibility of upper ramification groups,
in exactly the totally ramified, perfect-residue-field setting of
\cite[\S3.1, Proposition~3, p.~128]{Serre1961LocalClosed}. The full norm-polynomial proof and Herbrand passage are also given in
\cite[Proposition~3.7, Theorem~3.10 and Proposition~3.13,
pp.~14--16]{ColmezCorpsLocaux}.
For every $v\ne0$, the image of $W^{b_0}$ in $W/H_v$ is all of
$W/H_v$, and its image is zero at every upper index greater than
$b_0$. Since $\bigcap_{v\ne0}H_v=0$, there is no nonzero upper
group after $b_0$. Put $S=W^{b_0}$. Thus
\[
                 S\longrightarrow\mathbf F_q,
                         B\longmapsto Bv
               \quad\hbox{is surjective for every }v\ne0.       \tag{B-RD58.2}
\]

The determinant on $W$ is $1+p\operatorname{tr}(B)$.
The determinant of the good height-two, dimension-one group over
$L$ is cyclotomic, by the determinant argument already proved in
B18 and B49; its finite tame determinant has become trivial over $L$.
Thus its degree-$p$ quotient field is $E_{\rm cyc}=L(\mu_{p^2})$.
This field has degree $p$: the cyclotomic extension over the
unramified coefficient field has ramification degree $p(p-1)$,
whereas $e_L=(p-1)N$ is prime to $p$.
For $x=\zeta_{p^2}-1$ one has
\[
 v_{E_{\rm cyc}}(x)=N,
 \qquad v_{E_{\rm cyc}}(g x-x)=pN\quad(g\ne1).
\]
The same prime-to-$p$ valuation calculation gives its single break
$pN-N=b_0$. Quotient compatibility therefore shows that
$\operatorname{tr}:S\to\mathbf F_p$ is nonzero.

Choose a tame complement of order $N$ and an eigen-uniformizer
$\varpi$ of $F$. As in B56, its multiplier is a primitive $N$th
root $\eta$; here $\eta=\zeta$, since $p^4\equiv1\pmod N$.
If $c$ is the last lower break, the leading ramification map
\[
 S\hookrightarrow k,\qquad
 g\longmapsto d_g=
    \overline{(g\varpi-\varpi)/\varpi^{c+1}}
\]
is injective and additive, with tame action multiplication by
$\eta^c$. Since trace on $S$ is nonzero and tame-invariant,
averaging supplies a nonzero tame-fixed vector of $S$.
Its nonzero image in $k$ forces $\eta^c=1$.
Injectivity then says that tame conjugation fixes all of $S$.
The centralizer of the irreducible tame residual generator is
precisely $Z=\mathbf F_q$. Thus $S\subset Z$.
Surjectivity in (B-RD58.2) forces $\dim_{\mathbf F_p}S\ge2$,
so $S=Z$.

This is the final lower group, and it becomes trivial immediately
after $c$; hence every nonzero element of $Z$ has break $c$.
The increasing Herbrand function for $F/L$ satisfies
$\varphi(c)=b_0$. Since $S$ is a proper subgroup of $W$,
its slope is strictly less than one on a nonempty interval before
$c$, so $c>b_0$. All other elements disappear before this last
group, proving the final assertion.
\end{proof}

\begin{lemma}[The earlier irreducible anti-Cartan layer]
\label{lem:B-RD58-anti-signature}
\status{Proved; full independent root and companion proof, original-source, and exact-code audits passed.}
Suppose the common residual projective order is $h>2$.
The anti-Cartan plane $U$ has one lower break $b<c$, and
$0<b<b_0$. In an exact Honda-framed coordinate, write the leading
coordinate of a primitive point as $d z\varpi^A$, where
\[
                     A=q(p-1),\qquad z\in\mathbf F_q.
\]
For each $z\ne0$, the first nonzero term of degree prime to $p$
below $qA$ has degree
\[
                     K=A+Nb,                            \tag{B-RD58.3}
\]
and its coefficient is one fixed nonzero multiple of $z$.
Every nonzero term of degree $J$ with $v_p(J)=1$ satisfies
\[
                     J>A+(q-p)b.                        \tag{B-RD58.4}
\]
\end{lemma}
\begin{proof}
The action on $U$ is the irreducible two-dimensional action by
$\lambda^{1-p}$ from B56. Thus every nonzero element of $U$ has
the same lower break $b$, and $b<c$ by the preceding lemma.
Its leading ramification coefficient space
$D_U=\{d_g:g\in U\}\subset k$ is a two-dimensional
$\mathbf F_p$-space stable under multiplication by $\eta^b$.
This scalar has degree two over $\mathbf F_p$, so
\[
                 D_U=d_U\mathbf F_q\quad(d_U\ne0).
\]
In particular $\eta^b\ne1$, so $b\ne b_0$.

Primitive coordinates have valuation $A=q(p-1)$ and nonzero
first-level coordinates have valuation $qA$, by B55. Moreover
$e_F=NqA>qA$, so the fixed-coefficient expansions modulo
$\varpi^{qA+1}$ are characteristic-$p$ expansions. Their coefficients
below $qA$ depend only on the primitive residual class $z$, exactly
as in B56. Every nonzero anti-Cartan matrix is invertible, so
\[
                      v_F(gT(P)-T(P))=qA\quad(g\in U\setminus0).
                                                               \tag{B-RD58.5}
\]
For a monomial of degree $j$, its first change has degree
$j+p^{v_p(j)}b$. The leading degree $A$ has exact $p$-valuation
two, and changes with coefficient $(p-1)d z d_g^q$ in degree
$A+qb$. Every other degree divisible by $q$ changes strictly later.
Write $J$ for the first nonzero degree with $v_p(J)=1$, and $K$
for the first nonzero degree prime to $p$; allow either to be
infinite if absent below $qA$.

If $b>b_0$, then $A+qb>qA$. Any earlier change could only come
from the degree-$J$ and degree-$K$ terms. If their first changes
coincide, their coefficient is a polynomial
$a d_g^p+b' d_g$ with coefficients independent of $g$.
A nonzero polynomial of degree at most $p$ cannot vanish at all
$q$ elements of $D_U$. Therefore (B-RD58.5) forces
$J+pb\ge qA$ and $K+b\ge qA$.
For a nonzero Cartan element of break $c>b$, all these changes
then occur strictly after $qA$, as do those of every degree
divisible by $q$. This contradicts its nonzero first-level
displacement. Hence $b>b_0$ is impossible.

For $b<b_0$, the leading degree changes before $qA$.
There cannot be an even earlier change from $J$ and $K$, by
the same degree-$p$ polynomial argument. Cancellation in degree
$A+qb$ must therefore involve a polynomial
\[
              (p-1)d z X^q+a X^p+b'X
\]
which vanishes on $D_U=d_U\mathbf F_q$.
The monic polynomial with that root space is
$X^q-d_U^{q-1}X$. Comparing coefficients shows that the
$X^p$ coefficient must be zero and the linear coefficient
must cancel the nonzero $X^q$ coefficient.
Thus $J+pb>A+qb$, proving (B-RD58.4), while
$K+b=A+qb$, proving (B-RD58.3). Its coefficient equation is
\[
 K f_K(z)=-(p-1)d\,d_U^{q-1}z.                           \tag{B-RD58.6}
\]
In particular $p\nmid K$ and this first coefficient is the same
nonzero scalar multiple of $z$ for every $z\ne0$.
No linearity of all the other coordinate coefficients has been
assumed.
\end{proof}

\begin{theorem}[Only residual projective order two remains]
\label{thm:B-RD58}
\status{Proved; full independent root and companion proof, original-source, and exact-code audits passed.}
In the normalized $p^2$ comparison with irreducible residual
inertia, different special characters require residual projective
order exactly two. In particular, if $p\equiv3\pmod4$, isomorphic
$p^2$-torsion in this irreducible residual case determines the
requested special character. If $p\equiv1\pmod4$, the only
remaining irreducible-residual type pair is the unique opposite
quadratic pair from Corollary~\ref{cor:B-RD53-quadratic}, and its
wild image must still be the full $M_2(\mathbf F_p)$.
At $p=5$, the pattern $\{A,B\}$ is excluded in the irreducible
residual case; only $\{A,C\}$ with wild order $625$ remains there.
\end{theorem}
\begin{proof}
Assume $h>2$. For a nonzero Cartan element of break $c>b$, the
degree-$K$ term from the preceding lemma changes before every
term of degree divisible by $p$:
\[
 \begin{split}
 J+pc-(K+c)&>(p-1)(c-b)>0,\\
 A+qc-(K+c)&=N(c-b)>0.
 \end{split}                                             \tag{B-RD58.7}
\]
It is also the first prime-to-$p$ term, so it gives the unique
first change. The required displacement valuation is $qA$.
Consequently $K+c=qA$ and the leading-coefficient equation on
the full final Cartan plane $Z$ is
\[
                   d_i^q\beta_i(g)z=K f_{i,K}(z)d_g.
                                                               \tag{B-RD58.8}
\]
Equation (B-RD58.6) yields constants $\kappa_i\ne0$ with
$\beta_i(g)=\kappa_i d_g$ for every $g\in Z$.
Since $Z$ is the full Cartan plane, $\beta_i(Z)=\mathbf F_q$.
But the different-type residual intertwiner is Frobenius-semilinear
by B53, and multiplication by $p$ uses the same orientation.
Thus $\beta_2(g)=\beta_1(g)^p$, just as in (B-RD57.4).
The two fixed scalar multiples of $d_g$ would make Frobenius
scalar-linear on all of $\mathbf F_q$, a contradiction.
Together with B57, this excludes every possible wild image at
$h>2$. Irreducibility excludes $h=1$.

For $h=2$, put $c_1\equiv1\pmod{p-1}$ for the first actual
tangent exponent. The residual eigenvalue is
$\lambda=\zeta^{c_1+p-1}$ and its projective order is the order
of $\lambda^{p-1}$. Order two means
\[
              (p-1)(c_1+p-1)\equiv N/2\pmod N.
\]
The different-type relation (B-RD53.4) then gives
$c_2-c_1\equiv N/2$, precisely an opposite quadratic pair.
Also $c_1+p-1$ is odd. For its class modulo $p+1$ to have
order two, $(p+1)/2$ must be odd, so $p\equiv1\pmod4$.
The cycle calculation of B53 gives exactly one such unordered
pair. The three-type table at $p=5$ gives the stated specialization.
\end{proof}

\begin{remark}[Scope and numbering]
\label{rem:B-RD58-scope}
\status{Proved scope and limitation; no pair-existence or least-index assertion.}
This argument uses upper quotient compatibility explicitly and only
then identifies the last lower group. It never sets a lower break
equal to the cyclotomic upper break without this passage.
Reducible residual representations remain outside its scope.
At $p=5$ the unrestricted least-index interval is still
$\nu_{\min}(5,4,2)\in\{1,2\}$, since the quadratic irreducible
case and the reducible residual case remain open.
The second lemma requires the irreducible tame action on the entire
anti-Cartan plane. Its scalar-$\mathbf F_q$ leading space is not
available automatically at projective order two, where one must
retain a general additive quadratic polynomial and possibly two
anti-Cartan breaks.
The retained exact supplement
\path{code/reduction_descent/verify_tame_full_wild_projective.py}
checks $24$ ramification-index cases, $5{,}574{,}772$ low-break
inequality cases, $1{,}180$ different-type cases, $92{,}456$
additive-polynomial coefficient cases, $364$ scalar/Frobenius cases,
$130{,}248$ primitive-point stabilizer cases and $3{,}026$ centralizer
cases. Its author run passed in $7.383174480994057$ seconds.
The complete independent root and companion audits passed; their
author-output-preserving exact reruns took $7.658215595998627$
and $7.812601506004285$ seconds, respectively.
These are finite supplements to the proof; no formal group,
division extension, or congruent pair is instantiated.
\end{remark}


% Separate post-CP20 candidate. No promotion before full independent audits.
\subsection{The two possible anti-Cartan break patterns}

Retain the complete mixed-characteristic $(0,p)$, algebraically closed
residue, $p\ge5$, $e_K=p-1$, $\mu_p\subset K$ hypotheses of B58.
We consider the hypothetical different-character $p^2$ comparison
with irreducible residual inertia. By B57 and B58 only the full
wild image and residual projective order two remain. Write
\[
 W=Z\oplus U=M_2(\mathbf F_p),\quad q=p^2,\quad N=q-1,
 \quad A=q(p-1),\quad b_0=N(p-1).
\]
Here tame conjugation is $+1$ on the Cartan plane $Z$ and $-1$ on
the anti-Cartan plane $U$. Lemma~\ref{lem:B-RD58-final-Cartan}
identifies the last lower group as all of $Z$, with one break
$c>b_0$. Every anti-Cartan break is strictly smaller than $c$.

\begin{lemma}[The final-plane coefficient bound]
\label{lem:B-RD59-final-bound}
\status{Proved; full independent root and companion proof and exact-code audits passed.}
Fix one of the two models and a primitive residual coordinate
$z\ne0$. In its fixed uniformizer expansion, let $J$ be the first
nonzero degree of exact $p$-valuation one, and let $K$ be the first
nonzero degree prime to $p$, both below $qA$. An absent degree is
interpreted as $+\infty$. Then
\[
          J+pc\ge qA,\qquad K+c\ge qA,\qquad
                       \min(J+pc,K+c)=qA.              \tag{B-RD59.1}
\]
All degrees divisible by $q$ first change after $qA$ under a
nonzero element of $Z$.
\end{lemma}
\begin{proof}
The leading degree is $A$, with exact $p$-valuation two, and its
first change has degree $A+qc>qA$ because $c>b_0$.
Every higher degree divisible by $q$ changes later still.
The other two possible earliest changes have degrees $J+pc$ and
$K+c$. If their minimum were below $qA$, it could not occur
uniquely, since every nonzero Cartan matrix displaces every primitive
point by a nonzero first-level point. If both degrees attained that
minimum, their coefficient would be
\[
                           a d_g^p+b d_g,
\]
with at least one coefficient nonzero. The leading ramification
space $\{d_g:g\in Z\}$ has $q$ elements, because $Z$ is the final
lower group. A nonzero polynomial of degree at most $p$ cannot
vanish on this space. Thus both first changes are at least $qA$.
At least one must equal $qA$, as otherwise every change is later
than the required first-level displacement. This proves (B-RD59.1).
\end{proof}

\begin{lemma}[Both anti-Cartan patterns force an ordinary first term]
\label{lem:B-RD59-anti-patterns}
\status{Proved; full independent root and companion proof and exact-code audits passed.}
For each model the first nonzero prime-to-$p$ degree $K$ is
independent of $z\ne0$, has coefficient a fixed nonzero multiple
of $z$, and satisfies
\[
                           K+c=qA.                     \tag{B-RD59.2}
\]
Its change under every nonzero final Cartan element is the unique
first change. More precisely, if the anti-Cartan plane has a single
break $b$, then
\[
            0<b<b_0,\qquad K=A+Nb,\qquad
                           J\ge A+(q-p)b.              \tag{B-RD59.3}
\]
If it has two breaks $b_1<b_2$, then
\[
       J=A+(q-p)b_1,\qquad
       K=J+(p-1)b_2,\qquad b_1<b_0.                    \tag{B-RD59.4}
\]
These are conditional coefficient formulas, not existence assertions
for a formal group with either ramification pattern.
\end{lemma}
\begin{proof}
The lower filtration on the two-dimensional plane $U$ has either
one or two breaks. Its nonzero elements are all invertible as
endomorphisms of the residual space. Thus their displacement has
valuation $qA$ for every $z\ne0$.

Suppose first that all nonzero elements of $U$ have one break $b$.
Its leading coefficient space $D$ is an arbitrary two-dimensional
$\mathbf F_p$-subspace of $k$, not assumed to be a scalar multiple
of $\mathbf F_q$. Its root polynomial has the form
\[
                    P_D(X)=X^q+sX^p+tX,\qquad t\ne0.   \tag{B-RD59.5}
\]
Indeed, choose a basis $a,b$ of $D$ and put
$Y=X^p-a^{p-1}X$. Then
\[
 P_D(X)=Y^p-(b^p-a^{p-1}b)^{p-1}Y.
\]
The two factors giving its linear coefficient are nonzero by
linear independence. This proves (B-RD59.5) directly.
Tame conjugation on $U$ is $-1$, so $b\ne b_0$.
If $b>b_0$, the degree-$A$ change is already after $qA$.
Cancellation of an earlier change between the $J$ and $K$ terms
would require a polynomial of degree at most $p$ to vanish on
all of $D$, which is impossible. Thus $J+pb\ge qA$ and
$K+b\ge qA$, contradicting (B-RD59.1) because $c>b$.
Hence $b<b_0$.

As in the preceding argument, no change from $J$ or $K$ can occur
before $A+qb$: any such cancellation has degree at most $p$ in
the leading ramification coefficient. In degree $A+qb$, cancellation
must therefore be a nonzero scalar multiple of the full root
polynomial (B-RD59.5), with leading coefficient $(p-1)d_i z$.
Its nonzero linear coefficient forces $K+b=A+qb$, with
$Kf_{i,K}(z)$ a fixed nonzero multiple of $z$.
The possible $X^p$ term gives $J+pb\ge A+qb$; equality is
permitted and is not discarded. These are (B-RD59.3).
At the later Cartan break $c$ the prime-to-$p$ term is strictly
earlier than the other two classes, since
\[
 J+pc-(K+c)\ge(p-1)(c-b)>0,\qquad
 A+qc-(K+c)=N(c-b)>0.
\]
Equation (B-RD59.1) proves (B-RD59.2) and uniqueness.

Now suppose the two anti-Cartan breaks are $b_1<b_2<c$.
For the fixed coordinate $z$, consider the three possible first-change
degrees as functions of a lower break $b$:
\[
             A+qb,\qquad J+pb,\qquad K+b.               \tag{B-RD59.6}
\]
By (B-RD59.1), at least one of the last two degrees is strictly
less than $qA$ at $b_2$. Thus the minimum in (B-RD59.6) at
$b_2$ is below $qA$ and must occur at least twice.
It cannot include the first degree $A+qb_2$: if that degree
were at most both others at $b_2$, its strictly larger slope
$q$ would make it uniquely smallest at $b_1$, still below
$qA$. This would give an uncancelled earlier displacement.
Consequently
\[
             J+pb_2=K+b_2<A+qb_2,\qquad
                              J+pb_2<qA.               \tag{B-RD59.7}
\]
At $b_1$, the degree $J+pb_1$ is strictly smaller than $K+b_1$.
It is below $qA$. Cancellation now forces
$A+qb_1=J+pb_1$. These two equalities give (B-RD59.4), in
particular $b_1<b_0$.

Choose fixed nonzero anti-Cartan elements $g_1,g_2$ with these
two breaks. The cancellation at $b_1$ is between the leading
degree $A$ and degree $J$, so its coefficient equation makes
$f_{i,J}(z)$ a fixed nonzero multiple of $z$.
The cancellation at $b_2$ is then between degrees $J$ and $K$,
so it makes $f_{i,K}(z)$ a fixed nonzero multiple of
$f_{i,J}(z)$ and hence of $z$. The equations involve the same
fixed elements $g_1,g_2$ for every $z\ne0$. Thus the degrees
and scalar multiples do not depend on $z$.
Finally
\[
                  J+pc-(K+c)=(p-1)(c-b_2)>0.
\]
Together with (B-RD59.1), this gives (B-RD59.2); all degrees
divisible by $q$ are later by the first lemma. The prime-to-$p$
term is therefore uniquely first, proving the lemma in both cases.
\end{proof}

\begin{theorem}[Character recovery for every irreducible residual comparison]
\label{thm:B-RD59}
\status{Proved; full independent root and companion proof and exact-code audits passed.}
For $p\ge5$ and $e_K=p-1$, with $\mu_p\subset K$ and algebraically
closed residue field, isomorphic $p^2$-torsion determines the
requested special tame-order character for the potentially
supersingular height-two groups in the normalized comparison,
provided their common residual inertia representation is irreducible
over $\mathbf F_p$.
Consequently every remaining different-character pair in the general
minimal-ramification test has reducible residual inertia.
In particular this applies to any remaining failure at $25$-torsion
for $(p,e,h)=(5,4,2)$.
\end{theorem}
\begin{proof}
B57 and B58 leave only the projective-order-two full-wild case.
On its final Cartan plane the preceding lemma gives, for both models,
\[
                   d_i^q\beta_i(g)z=K f_{i,K}(z)d_g,
                     \qquad g\in Z.                   \tag{B-RD59.8}
\]
Since $f_{i,K}(z)$ is a fixed nonzero multiple of $z$, this
forces $\beta_i(g)=\kappa_i d_g$ on all of $Z$.
But $\beta_i:Z\to\mathbf F_q$ is bijective, while the
different-type residual intertwiner is Frobenius-semilinear and
therefore requires $\beta_2(g)=\beta_1(g)^p$.
As in B55 and B57, this would make Frobenius scalar-linear on
all of $\mathbf F_q$, a contradiction.
Thus no irreducible-residual different-character comparison remains.
\end{proof}

\begin{remark}[The unresolved residual case]
\label{rem:B-RD59-scope}
\status{Proved scope and limitation; reducible residual comparisons remain open.}
The theorem does not assume CM and does not replace a torsion module
by a different integral lattice. It uses the actual common division
field and the upper-ramification final-Cartan lemma of B58.
It leaves reducible residual representations entirely open, so the
unrestricted least-index interval at $(5,4,2)$ is still $\{1,2\}$.
The two-break argument compares all three slopes $q,p,1$; it is
not an extrapolation of the earlier cubic-wild two-slope argument.
The exact finite supplement
\path{code/reduction_descent/verify_tame_full_wild_quadratic.py}
checks 252 single-break and 124,788 two-break arithmetic cases,
27,648 three-slope candidates (432 satisfying the conditional ties),
and all 221 two-dimensional planes over $\mathbf F_5$, $\mathbf F_7$
and $\mathbf F_{11}$ in their cubic extensions. All these planes
have a nonzero middle coefficient in (B-RD59.5); the program checks
806 directly multiplied polynomial coefficients and 200,449 exact
root tests, as well as 364 scalar/Frobenius comparisons. The author
run passed in 1.1602564200002234 seconds. The complete independent
root and companion audits passed; their author-output-preserving
reruns took 1.100486028000887 and 1.1509096779991523 seconds,
respectively. The first version's
three-slope subtest had an empty admissible range; that source and
output are retained, and the revised range has the explicit
non-vacuity assertion just reported. These are finite checks of
the proof's identities, not constructions of a formal group or pair.
\end{remark}


% Separate post-CP21 candidate. All source/finite checks retained.
\subsection{The Newton filtration in the remaining reducible case}

Throughout B60--B61, $p\ge5$, $K$ is complete discretely valued of
mixed characteristic $(0,p)$ with algebraically closed residue field,
$e_K=p-1$, and $\mu_p\subset K$. We work in the normalized remaining
comparison of B49--B50: both groups are potentially supersingular of
height two, and their common finite determinant has been twisted to
the level-one tame fundamental character. Thus, putting
\[
              q=p^2,\qquad N=q-1,
\]
we may take the common good field $L/K$ to be tame cyclic of degree
$N$, and its absolute index is $e_L=N(p-1)$. Choose compatible tame
uniformizers and a primitive $N$th root $\zeta$ so that the generator
$\sigma$ acts on a uniformizer of $L$ by $\zeta$. An actual special
tangent exponent $c$ then satisfies
\[
                    c\equiv1\pmod{p-1}.                \tag{B-RD60.1}
\]
The special covariant character is the pair of exponents $c,pc$.
The normalization is a common finite height-one twist, as in B53;
it is not an identification of the mod-$p$ cyclotomic character
with the fundamental character.

\begin{lemma}[An elementary lower-break bound]
\label{lem:B-RD60-break}
\status{Proved; full independent root and companion proof, original-source, and exact-code audits passed.}
Let $E/D$ be a finite totally ramified Galois extension of complete
discretely valued fields with residue characteristic $p$, and let
$g$ be a nonidentity wild inertia element. Normalize $v_E$ to have
value group $\mathbf Z$, and put
\[
       b(g)=v_E(g\varpi-\varpi)-1
\]
for a uniformizer $\varpi$ of $E$. For every $x\in\mathfrak m_E$,
\[
 v_E(gx-x)\ge v_E(x)+b(g).                              \tag{B-RD60.2}
\]
Equality holds if $p\nmid v_E(x)$; if $p\mid v_E(x)$ the inequality
is strict. The statements allow $gx=x$, with valuation $+\infty$.
\end{lemma}
\begin{proof}
Choose representatives in $\mathcal O_D$ of every residue-field
element. Completeness gives an expansion
$x=\sum_{j\ge A}a_j\varpi^j$, with $A=v_E(x)$, all $a_j$ among
these fixed representatives and $a_A$ a unit. Write
$g\varpi=\varpi(1+d\varpi^{b(g)}+\cdots)$, with $d$ a unit.
The change of the degree-$j$ term has valuation at least $j+b(g)$.
If $p\nmid A$, its degree-$A$ term has exactly that valuation and
nonzero leading coefficient $Aa_A\bar d$; all higher terms change
later. If $p\mid A$, the linear binomial coefficient is in the
maximal ideal, while every higher binomial term has strictly larger
valuation because $b(g)\ge1$. The degree-$A$ change and all changes
from $j>A$ are therefore strictly later than $A+b(g)$. This proves
all assertions. No identification of lower and upper numbering is
used here.
\end{proof}

\begin{lemma}[The canonical high line and its two valuations]
\label{lem:B-RD60-Newton}
\status{Proved; full independent root and companion proof, original-source, and exact-code audits passed.}
Let $M$ be one normalized potentially supersingular group, and
suppose its residual inertia representation $M[p]$ is reducible over
$\mathbf F_p$. Its nonzero $p$-torsion points have exactly two
valuations in the good formal coordinate over $L$. There is a unique
high-valuation line $H\subset M[p]$ and an integer $B$ such that
\[
 \begin{gathered}
 0<B<p(p-1),\qquad
 v_L(T(Q))=\alpha=N-B\quad(0\ne Q\in H),\\
 v_L(T(Q))=\beta=B/p\quad(Q\in M[p]\setminus H),\qquad
 \alpha>\beta.                                        \tag{B-RD60.3}
 \end{gathered}
\]
The line $H$ is inertia-stable and is rational over $L$.
In the exact tame coordinate, its character on $\sigma$ is
\[
                  \zeta^{c-B},\qquad
                         B\equiv c\pmod{p+1}.         \tag{B-RD60.4}
\]
For every point $P$ of order $p^2$ with $pP\notin H$,
\[
                       v_L(T(P))=B/p^3.               \tag{B-RD60.5}
\]
All these valuations are unchanged by an integral invertible
formal-coordinate change.
\end{lemma}
\begin{proof}
Use the height-two universal coordinate and coefficient estimate
in Lemma~\ref{lem:B-RD49-point}, read there from the original
Hopkins--Gross Propositions~5.7 and~12.10
\cite[Proposition~5.7, equations (12.3)--(12.5), and
Proposition~12.10 with proof]{HopkinsGross1994Bundles}. Put $e=e_L$ and
$b=v_L(u)$, allowing $+\infty$. The only possible negative-slope
vertices for nonzero $p$-torsion are
\[
                        (1,e),\quad(p,b),\quad(q,0).
\]
Indeed coefficients in degrees $2,\ldots,p-1$ vanish in that
coordinate, coefficients strictly between $p$ and $q$ lie in
$(p,u)$, and the degree-$q$ coefficient is a unit.
If $b\ge ep/(p+1)$, all nonzero points have one valuation $e/N$.
Their leading-coordinate map is injective: equal leading
coefficients would give a nonzero difference of larger valuation,
again a $p$-torsion point. Wild inertia acts trivially on leading
coefficients after good reduction and hence on this entire residual
module. The remaining cyclic tame representation, acting as scalar
multiplication on its leading-coordinate plane in the residue field,
is either irreducible over $\mathbf F_p$ or scalar over
$\mathbf F_p$. The former contradicts the hypothesis; the latter
would give a square residual determinant. Our normalized determinant
is the nonsquare level-one character, so this case is impossible.

Consequently $0<b<ep/(p+1)$. The Newton polygon gives $p-1$ points
of valuation $(e-b)/(p-1)$ and $q-p$ points of valuation
$b/(p(p-1))$. The former, together with zero, form a line $H$:
formal addition and inversion preserve the condition of having at
least that valuation. Descent preserves valuations, by the same
integral good-model and Tate full-faithfulness argument as in B53,
so $H$ is inertia-stable. Its character has values in
$\mathbf F_p^\times$; it is tame of order dividing $p-1$ and is
therefore killed over $L$. Its coordinate valuation is an integer.
Writing that integer as $N-B$ gives $b=(p-1)B$, with $B$ an integer
in the range of (B-RD60.3).

For the character calculation use tame averaging, as in B56, to
choose an integral parameter in which the descent action on points
is $T\mapsto\zeta^c\sigma(T)$. In this parameter the coefficient
$a_p$ of degree $p$ in $[p](T)$ still has valuation $b<e$.
To check this last assertion without assuming a $p$-typical averaged
coordinate, reduce the coordinate change modulo $p$: the first
nonzero degree of $[p](T)$ is $p$, and its coefficient is changed
by a unit $(p-1)$st power. Terms from $pT$ lie in $p\mathcal O_L$.
Thus the valuation $b<e$ is unchanged. Compatibility with $[p]$
gives
\[
                  \sigma(a_p)=\zeta^{(p-1)c}a_p.
\]
Taking its leading coefficient in the tame uniformizer shows
$b\equiv(p-1)c\pmod N$, hence $B\equiv c\pmod{p+1}$.
For $0\ne Q\in H$, the coordinate valuation is $N-B$;
its field action contributes $\zeta^{N-B}$ and its descent action
contributes $\zeta^c$. This proves (B-RD60.4).

Finally $pP$ has valuation $\beta=B/p<p-1$. In the equation
$[p](T)=T(pP)$, its constant term and unit degree-$q$ term give a
single Newton segment: the linear term has coefficient valuation
$e>\beta$, the degree-$p$ coefficient has valuation
$b=(p-1)B>\beta$, and all intermediate coefficients are at least
as large as $b$. Every one of its $q$ roots thus has valuation
$\beta/q=B/p^3$. This is (B-RD60.5).
\end{proof}

\begin{proposition}[The first residual wild break]
\label{prop:B-RD60-residual-break}
\status{Proved; full independent root and companion proof, original-source, and exact-code audits passed.}
Let $R=L(M[p])$. Then $R=L$ or $R/L$ is cyclic of degree $p$.
In the former case $p\mid B$. In the latter case $H$ is the unique
line fixed by wild inertia and, writing $r$ for the single lower
break of $R/L$,
\[
 \begin{cases}
 r=pN-(p+1)B,&p\nmid B,\\
 r<pN-(p+1)B,&p\mid B.
 \end{cases}                                          \tag{B-RD60.6}
\]
If two normalized groups have isomorphic residual modules with
nontrivial wild image, then their high lines agree under that
isomorphism. If both associated integers $B_i$ are prime to $p$,
or both are divisible by $p$, their special characters already
agree. In the only remaining case, after indexing so that
$p\nmid B_1$ and $p\mid B_2$, one has $B_1>B_2$.
\end{proposition}
\begin{proof}
A reducible subgroup of $\GL_2(\mathbf F_p)$ is triangular. Its
wild image has order at most $p$, and its two tame diagonal
characters have order dividing $p-1$, killed over $L$. This proves
the assertion about $R$. If $R=L$, the low valuation $B/p$ is
integral, so $p\mid B$.

Otherwise the nontrivial unipotent group has one fixed line,
necessarily the invariant high line $H$. Every nonzero $g$ sends
a point $Q\notin H$ to $Q+Q_0$ for some $0\ne Q_0\in H$.
With $v_R$ normalized, its coordinate has valuation $B$, while
its difference has valuation $p(N-B)$. Ordinary coordinate
difference and formal-group subtraction have the same valuation:
the latter is the former times an integral unit. Applying
Lemma~\ref{lem:B-RD60-break} proves (B-RD60.6).

The common residual representation identifies the two unique high
lines and their characters. Thus
\[
                   c_1-B_1\equiv c_2-B_2\pmod N.       \tag{B-RD60.7}
\]
If both $B_i$ are prime to $p$, their common break gives $B_1=B_2$,
and (B-RD60.7) gives $c_1=c_2$ modulo $N$.
If both are divisible by $p$, write $B_i=pb_i$.
The range in (B-RD60.3) gives $1\le b_i\le p-2$.
Reducing (B-RD60.7) modulo $p-1$, and using (B-RD60.1), gives
$b_1\equiv b_2\pmod{p-1}$; the range forces equality, and again
the actual tangent exponents agree.
In the remaining case, equality of the common break and the strict
inequality on the second side give
\[
      pN-(p+1)B_1=r<pN-(p+1)B_2,
\]
so $B_1>B_2$ as asserted. No pair is constructed by these necessary
relations.
\end{proof}


% Separate post-CP21 candidate; numerical least-index claim needs full audits.
\subsection{Second-level point fields close the reducible branch}

Retain the exact mixed-characteristic, residue, determinant and
valuation conventions of B60.

\begin{lemma}[The elementary-abelian point field]
\label{lem:B-RD61-point-field}
\status{Proved; full independent root and companion proof, original-source, and exact-code audits passed.}
Let $R$ contain the full $p$-torsion of a height-two group $M$, and
let $P$ have order $p^2$. Then $E=R(P)$ is Galois over $R$, with
elementary abelian group of order at most $p^2$. If its degree is
$p^2$, the map
\[
       \operatorname{Gal}(E/R)\longrightarrow M[p],
                       \qquad g\longmapsto gP-P        \tag{B-RD61.1}
\]
is an isomorphism of $\mathbf F_p$-spaces. An isomorphism of the
actual generic $p^2$-torsion modules identifies the point fields
and this map for corresponding points.
\end{lemma}
\begin{proof}
Over $R$, the image on $M[p^2]\simeq(\mathbf Z/p^2)^2$ lies
in $1+pM_2(\mathbf F_p)$, an elementary abelian group.
Consequently every subfield of the full $p^2$-division field is
Galois over $R$. All conjugates of $P$ lie in the fibre of $[p]$
over the rational point $pP$, a set of $p^2$ points. This bounds
the degree of $R(P)$. The displayed map is an injective homomorphism
because all translations are $R$-rational; at degree $p^2$ it is
surjective. Finally an equivariant torsion-module isomorphism gives
exactly the same stabilizer for corresponding points and translates,
which proves the last assertion. The coordinate field of a point
in the good formal group equals its field of definition: an
integral parameter is a formal coordinate, so fixing it fixes the
point, and conversely.
\end{proof}

\begin{proposition}[No remaining wild reducible comparison]
\label{prop:B-RD61-wild}
\status{Proved; full independent root and companion proof, original-source, and exact-code audits passed.}
Two normalized potentially supersingular height-two groups with
isomorphic $p^2$-torsion and reducible residual inertia of nontrivial
wild image have the same special character.
\end{proposition}
\begin{proof}
By Proposition~\ref{prop:B-RD60-residual-break}, only the case
$p\nmid B_1$, $p\mid B_2$ and $B_1>B_2$ could have different
characters. Identify their torsion modules, and let $R/L$ be their
common full $p$-division field, cyclic of degree $p$.
Their unique high lines coincide; call the common line $H$.
Choose corresponding points $P_i$ of order $p^2$ with
$pP_i\notin H$. Formula (B-RD60.5) gives
\[
                 v_R(T_i(P_i))=B_i/p^2.
\]
On the first side this fraction has denominator $p^2$. The
ramification index of its point field is therefore divisible by
$p^2$. Lemma~\ref{lem:B-RD61-point-field} bounds its degree by
$p^2$, so the common point field $E=R(P_i)$ is Galois of degree
exactly $p^2$. In particular
\[
       v_E(T_i(P_i))=B_i,
       \qquad [E:L]=p^3.                              \tag{B-RD61.2}
\]
Every $p$-torsion translation occurs in (B-RD61.1).
Choose $g$ whose nonzero translation lies outside $H$.
Its coordinate difference in model $i$ has valuation
\[
       [E:L]\,\frac{B_i}{p}=p^2 B_i.                 \tag{B-RD61.3}
\]
Since $p\nmid B_1$, Lemma~\ref{lem:B-RD60-break} identifies the
actual lower break of this particular $g$ on $E/R$ as
\[
                          b_E(g)=NB_1.
\]
Applying just the inequality of that lemma to the second model
then gives
\[
                         p^2B_2\ge B_2+NB_1,
\]
which contradicts $B_1>B_2$. This excludes the final wild case.
The argument uses an actual common point field, not an abstract
field with the same order or a substituted integral lattice.
\end{proof}

\begin{proposition}[No remaining tame reducible comparison]
\label{prop:B-RD61-tame}
\status{Proved; full independent root and companion proof, original-source, and exact-code audits passed.}
Two normalized potentially supersingular height-two groups with
isomorphic $p^2$-torsion and tame reducible residual inertia have
the same special character.
\end{proposition}
\begin{proof}
Here the common $p$-division field over $L$ is $L$ itself.
Write $B_i=pb_i$, with $1\le b_i\le p-2$, by B60.
Their high lines $H_1,H_2$ are not initially identified.
Choose a common nonzero residual vector outside their union,
and lift it to corresponding points $P_i$ of order $p^2$.
This is possible because the union of two lines is a proper
subset of $\mathbf F_p^2$. The low primitive valuations are
\[
                          v_L(T_i(P_i))=b_i/p^2.
\]
As $p\nmid b_i$, their common point field $E=L(P_i)$ has degree
$p^2$ by Lemma~\ref{lem:B-RD61-point-field} and the valuation
denominator. Hence $v_E(T_i(P_i))=b_i$ and all translations in
$M_i[p]$ occur.

Choose a translation outside both high lines. Its coordinate
difference in model $i$ has valuation $p^2b_i$.
The equality part of Lemma~\ref{lem:B-RD60-break} therefore gives
one and the same lower break
\[
                       b_E(g)=Nb_1=Nb_2.
\]
Thus $b_1=b_2=b$. If the high lines were different, choose a
nonzero translation lying in $H_1$ and outside $H_2$.
The corresponding difference valuations would be, respectively,
\[
                   p^2(N-pb)\quad\hbox{and}\quad p^2b.
\]
Both coordinates of $P_i$ have the same prime-to-$p$ valuation
$b$, so their differences under this same $g$ must have the same
valuation by the equality part of the lower-break lemma. That
would force $N-pb=b$, or $b=p-1$, outside its allowed range.
Consequently $H_1=H_2$. Their common character
$\zeta^{c_i-pb}$ in (B-RD60.4) now gives $c_1=c_2$ modulo $N$,
as required.
\end{proof}

\begin{theorem}[The least index at minimal cyclotomic ramification]
\label{thm:B-RD61}
\status{Proved; full independent root and companion proof, original-source, and exact-code audits passed.}
For every prime $p\ge5$, under the manuscript's monodromy condition,
height two and $e_K=p-1$, with $\mu_p\subset K$ and perfect residue
field, the exact least index in B.4 is
\[
                         \nu_{\min}(p,p-1,2)=1.         \tag{B-RD61.4}
\]
Thus $p^2$-torsion always determines the requested special
prime-to-$p$ inertia character, while $p$-torsion need not.
In particular $\nu_{\min}(5,4,2)=1$.
\end{theorem}
\begin{proof}
The lower bound is the actual first-level counterexample of B18,
as used in Theorem~\ref{thm:B-RD49}. For the upper bound pass to
the completed maximal unramified extension, without changing
absolute ramification or the character question on inertia.
The dimension and mixed-case reductions of B49 leave only two
whole potentially supersingular groups with the same finite
determinant. Proposition~\ref{prop:B-RD50-square} handles its
square case. In the nonsquare case the common finite height-one
twist normalizes the determinant exactly as in B53 and B60.
The common residual module is either irreducible over $\mathbf F_p$,
handled by Theorem~\ref{thm:B-RD59}, or reducible. Its wild image
is then trivial or cyclic of order $p$; the two preceding
propositions handle these cases. The special characters agree
in every case, proving the upper bound and equality.
\end{proof}

\begin{remark}[Scope of the new equality]
\label{rem:B-RD61-scope}
\status{Proved scope and limitation; other ramification indices and heights remain separate.}
The equality treats minimal cyclotomic ramification and height two
for $p\ge5$, including non-CM deformations and reducible residual
representations. It does not determine the general B.4 least index
at other ramification indices or heights. All point-field degrees
used above follow from the actual torsion action and reduced
valuation denominators. No existence conclusion is drawn from
an enumeration of possible Newton slopes.
The exact finite supplement
\path{code/reduction_descent/verify_tame_reducible_recovery.py}
checks 74,828 Newton/tame-character records for every prime from
5 through 101, 73,696 prime-to-$p$ wild records, 1,132 divisible
wild comparisons, 36,848 mixed point-field inequalities, 1,698
tame residual pairs (including 566 distinct-type exclusions), and
5,184 exact integer-binomial bounds. The author run passed in
1.8442117410013452 seconds. The complete independent root and
companion audits passed; their author-output-preserving exact reruns
took 1.8523498740032665 and 1.721134293002251 seconds, respectively. Its separate exploratory $p=5$ program
and full output are retained: the 20 Newton records give eight
necessary wild rows and two necessary tame rows before the
point-field contradiction. Neither finite list is used as a
substitute for the general proof.
\end{remark}


\subsection{Cyclic level structures and new filtered height-two counterexamples}

\begin{lemma}[A tame equivariant arc with a rational cyclic level structure]
\label{lem:B-RD62-arc}
\status{Proved; independent audits passed.}
Let $p\ge5$, let $n\ge2$, and let $e$ be a positive multiple of $p-1$
such that
\[
                    p^{2n-3}<e<p^{2n-2}.                 \tag{B-RD62.1}
\]
There is a complete discretely valued field $K$ of mixed characteristic
$(0,p)$, residue field $\overline{\mathbf F}_p$, and absolute ramification
$e$, containing $\mu_p$, and a height-two BT group $G/K$ with a
$K$-rational point of exact order $p^n$.  It has a good one-dimensional
formal model over a cyclic tame extension $L/K$ of degree $p^2-1$.
On its covariant special isocrystal, a generator of this descent group
has eigenvalues $\zeta^{-1},\zeta^{-p}$, where $\zeta$ is a primitive
$(p^2-1)$st root of unity.
\end{lemma}
\begin{proof}
Put $k=\overline{\mathbf F}_p$, $W=W(k)$, $q=p^2$, $N=q-1$, and
\[
             d=Ne,\qquad b=d-(p-1)q^{n-1}.              \tag{B-RD62.2}
\]
The inequalities in (B-RD62.1) imply
\[
 \begin{gathered}
 b>(q-p)q^{n-2}>0,\qquad q^{n-1}<d<Nq^{n-1},\\
 b\equiv-(p-1)\pmod N.
 \end{gathered}                                         \tag{B-RD62.3}
\]
For the first inequality, subtracting its right side from $b$ gives
$N(e-p^{2n-3})$.  The last congruence uses $q\equiv1\pmod N$.

Use the integral $p$-typical group $F_u$ over $W[[u]]$ obtained by
specializing $v_1=u$, $v_2=1$, and $v_j=0$ for $j>2$ in
\cite[Proposition~5.7, equations (5.3)--(5.5), and (12.3)--(12.5)]
{HopkinsGross1994Bundles}.
Its logarithm is
\[
 f_u(T)=\sum_{j\ge0}l_j(u)T^{p^j},\qquad l_0=1,
 \qquad p l_j=l_{j-1}u^{p^{j-1}}+l_{j-2},               \tag{B-RD62.4}
\]
where $l_{-1}=0$.  This is precisely the specialized recursion (5.5),
not a claim that the multiplication series itself is a polynomial.
The same source gives
\[
 \begin{gathered}
 [p]_u(T)=pT+(1-p^{p-1})uT^p+\sum_{j>p}c_j(u)T^j,\\
 c_j(u)\in(p,u)W[[u]]\quad(p<j<q),\\
 c_q(u)\equiv1\pmod{(p,u)}.
 \end{gathered}                                         \tag{B-RD62.5}
\]
The coefficients in degrees $2,\ldots,p-1$ vanish.  The special group
$F_0\otimes k$ is the height-two Honda group.

Define integral power series in $W[[X]]$ by
\[
 \begin{split}
 Y_0(X)&=X,\\
 Y_{j+1}(X)&=[p]_{X^b}\bigl(Y_j(X)\bigr),\\
 \Phi_u(T)&=[p]_u(T)/T,\\
 P(X)&=\Phi_{X^b}\bigl(Y_{n-1}(X)\bigr).
 \end{split}                                           \tag{B-RD62.6}
\]
All compositions converge in the indicated formal power-series ring.
Modulo $p$, induction using (B-RD62.3)--(B-RD62.5) gives
\[
                  Y_j(X)=X^{q^j}+O(X^{q^j+1})
                  \quad(0\le j\le n-1).                \tag{B-RD62.7}
\]
Indeed, at an input of order $q^{j-1}$ the $q$th-degree term is
strictly earlier than the $p$th-degree term because
$b>(q-p)q^{j-1}$; intermediate coefficients are multiples of $u$
modulo $p$, and all terms of degree greater than $q$ are later.
In the last series $P$, its $uY_{n-1}^{p-1}$ term has order
$b+(p-1)q^{n-1}=d$.  Every intermediate term is later, and its
$Y_{n-1}^{q-1}$ term has order $Nq^{n-1}>d$.  Consequently
\[
                    P(0)=p,\qquad
                    P(X)\equiv X^d+O(X^{d+1})\pmod p.  \tag{B-RD62.8}
\]
Weierstrass preparation therefore identifies
$R=W[[X]]/(P)$ with the ring defined by a degree-$d$ Eisenstein
polynomial.  It is a complete DVR, the image $x$ of $X$ is a
uniformizer, and $v_R(p)=d$.

The group $F_{x^b}$ has good height-two reduction over $R$.
The congruences (B-RD62.7) give $v_R(Y_j(x))=q^j$: any term whose
coefficient is divisible by $p$ has valuation at least $d>q^{n-1}$,
and the other terms are strictly later.  In particular
$Y_{n-1}(x)\ne0$.  Since $P(x)=0$,
\[
 [p]^n(x)=Y_{n-1}(x)\Phi_{x^b}(Y_{n-1}(x))=0.
\]
Thus $x$ is a point of exact order $p^n$ of this good formal group.

It remains to descend both the group and this particular point.
Choose $\zeta\in W^\times$ of order $N$ and put $a=\zeta^{-1}$.
The logarithm recursion proves, coefficient by coefficient,
\[
 l_j(a^{p-1}u)=a^{p^j-1}l_j(u),\qquad
 f_{a^{p-1}u}(T)=a^{-1}f_u(aT).                         \tag{B-RD62.9}
\]
For the $l_{j-2}$ term the exponents agree modulo $N$ because
$N\mid p^j-p^{j-2}$.  Hence multiplication of the coordinate by $a$
is an integral isomorphism $F_{a^{p-1}u}\to F_u$ and
\[
 [p]_{a^{p-1}u}(a^{-1}T)=a^{-1}[p]_u(T),\qquad
 \Phi_{a^{p-1}u}(a^{-1}T)=\Phi_u(T).                    \tag{B-RD62.10}
\]
The substitution $X\mapsto\zeta X$ sends $X^b$ to $a^{p-1}X^b$.
Induction gives $Y_j(\zeta X)=a^{-1}Y_j(X)$, and thus the exact
identity $P(\zeta X)=P(X)$.  It defines an automorphism $\sigma$ of
$R$ of order $N$, fixing $W$ and sending $x$ to $\zeta x$.

Let $L=\operatorname{Frac}(R)$ and $K=L^{\langle\sigma\rangle}$.
The extension $L/K$ is cyclic tame of degree $N$, with unchanged
residue field.  It follows that $v_K(p)=d/N=e$.
The isomorphism $\sigma^*F_{x^b}\to F_{x^b}$ with coordinate $aT$
is a descent datum: its $N$fold composite is the identity.
Moreover it fixes the marked point, since $a\sigma(x)=x$.
Effective Galois descent at every generic finite level gives the
BT group $G/K$ and its rational point.  Only the generic group is
descended; no good integral model over $\mathcal O_K$ is asserted.

On the Honda special fibre the descent is $T\mapsto aT$.
Its covariant isocrystal eigenvalues are $a,a^p$ in the conventions
of Proposition~\ref{prop:B-RD4-dictionary}; equivalently these are
the two embeddings of the unramified quadratic Teichmuller unit.
Finally $K$ contains $\mu_p$.  To see this without a polarization,
write $-p=\pi_K^e u$ with $u$ a unit.  Since $k$ is algebraically
closed, Hensel's lemma gives a $(p-1)$st root of $u$.  As $p-1\mid e$,
the field $K$ contains a $(p-1)$st root of $-p$.  Over
$\operatorname{Frac}W(k)$ this generates the unique tame extension
of degree $p-1$, which is the cyclotomic extension generated by
$\mu_p$.  No assertion that $\mu_{p^n}\subset K$ is used.
\end{proof}

\begin{lemma}[Lifting a finite rank-one extension]
\label{lem:B-RD62-extension}
\status{Proved; independent audits passed.}
Let $K$ have characteristic zero, and let
$\psi:G_K\to\mathbf Z_p^\times$ be a finite character of order
prime to $p$.  Every extension of free $\mathbf Z/p^n$-modules with
continuous Galois action
\[
 0\longrightarrow\mu_{p^n}\longrightarrow V_n
       \longrightarrow(\mathbf Z/p^n)(\psi^{-1})
       \longrightarrow0                               \tag{B-RD62.11}
\]
is the $p^n$-torsion of a height-two BT group $H/K$ admitting a
filtration with graded groups $\mu_{p^\infty}$ and
$(\mathbf Q_p/\mathbf Z_p)(\psi^{-1})$.
\end{lemma}
\begin{proof}
An underlying module splitting identifies the extension class with
$H^1(K,T/p^nT)$ for $T=\mathbf Z_p(1)\otimes\psi$.
Let $K'/K$ kill $\psi$, with group $\Delta$ of order prime to $p$.
Restriction identifies this cohomology with the $\Delta$-invariants
of $H^1(K',\mu_{p^n})\otimes\psi$, by averaging and
inflation--restriction.  The Kummer descriptions
\[
 H^1(K',\mu_{p^m})=K'^\times/(K'^\times)^{p^m}
\]
show that the transition map from level $m+1$ to level $m$ is
surjective: it sends the class of an element to the class of the
same element.  Averaging a lift by $|\Delta|^{-1}\sum_{\delta\in\Delta}\delta$
proves surjectivity also on the invariant parts.  Thus the class
at level $n$ lifts through all higher levels.

Choose continuous cocycles representing these compatible classes.
If a chosen cocycle at the next level reduces only up to a
coboundary, lift its correcting zero-cochain and subtract that
coboundary.  The cocycles can therefore be made exactly compatible.
Their inverse limit is a continuous $T$-valued cocycle and defines
an extension of free $\mathbf Z_p$-modules
\[
 0\longrightarrow\mathbf Z_p(1)\longrightarrow T_H
       \longrightarrow\mathbf Z_p(\psi^{-1})
       \longrightarrow0.
\]
The corresponding compatible finite etale group schemes over $K$
give the required BT group.  Its two rank-one graded groups acquire
good reduction after killing $\psi$, so it satisfies exactly the
monodromy condition used in B.4.  The entire extension is not claimed
to acquire good reduction.
\end{proof}

\begin{theorem}[Character ambiguity beyond minimal cyclotomic ramification]
\label{thm:B-RD62}
\status{Proved; independent audits passed.}
For every $p\ge5$, $n\ge2$, and multiple $e$ of $p-1$ satisfying
(B-RD62.1), there are height-two BT groups $M,H$ over a field $K$
with the properties in Lemma~\ref{lem:B-RD62-arc}, such that
\[
                            M[p^n]\simeq H[p^n],        \tag{B-RD62.12}
\]
both satisfy the monodromy condition, and their special characters
on prime-to-$p$ inertia differ.  The group $M$ individually acquires
good supersingular reduction; the stated filtration supplies the
monodromy condition for $H$.
Consequently a universal B.4 index in height $h\ge2$ satisfies
$\nu_{\min}(p,e,h)\ge n$ in this range.
\end{theorem}
\begin{proof}
Take $G$ and its rational point $P$ from the arc lemma, and put
$C=\langle P\rangle\simeq\mathbf Z/p^n$ and $G'=G/C$.
The finite-flat closure of $C$ in the good model over $\mathcal O_L$
and its quotient show that $G'$ also acquires good reduction.
Equivalently this is the usual quotient of a BT group by a finite
flat subgroup; the quotient is isogenous to the original group.

The rank-one crystalline determinant argument in the full proof of
Proposition~\ref{prop:B-RD18}, using
\cite[Corollary~2.2.6]{Kisin2006Crystalline}, gives
\[
                      \det T_pG=\mathbf Z_p(1)\otimes\psi,
\]
where $\psi$ is trivial on $G_L$.  Indeed, over $L$ the determinant
has weight one and is the Tate representation of a height-one
multiplicative group; the residue field is algebraically closed,
so there is no residual unramified character.  Its finite inertia
value on $\sigma$ is
\[
                         \psi(\sigma)=a^{p+1}.
                                                        \tag{B-RD62.13}
\]
Thus $\psi$ is a character of order $p-1$, valued in the Teichmuller
units of $\mathbf Z_p$.

The multiplication-by-$p^n$ snake sequence for $G\to G'$ is
\[
 0\longrightarrow C\longrightarrow G[p^n]
  \longrightarrow G'[p^n]\longrightarrow C\longrightarrow0.
\]
The image is the cyclic quotient $G[p^n]/C$.  Because $C$ is
constant and primitive as a free $\mathbf Z/p^n$ submodule, the
determinant identifies this quotient with
$(\mathbf Z/p^n)(\chi_p\psi)$.  Set $M=G'\otimes\psi^{-1}$.
It follows that its actual truncation fits into
\[
 0\longrightarrow\mu_{p^n}\longrightarrow M[p^n]
       \longrightarrow(\mathbf Z/p^n)(\psi^{-1})
       \longrightarrow0.
\]
Lemma~\ref{lem:B-RD62-extension} constructs $H$ with precisely this
truncation.  Twisting by the finite character preserves potential
good reduction of $M$; $H$ has the explicitly stated good graded
pieces.

Isogeny preserves the special rational character, and the finite
twist acts on it with the same scalar.  Thus, at the chosen generator
of the finite tame descent quotient, the two eigenvalue multisets are
\[
 \begin{array}{c|c}
 M&\{a^{-p},a^{-1}\}=\{\zeta^p,\zeta\}\\
 H&\{1,a^{-(p+1)}\}=\{1,\zeta^{p+1}\}.
 \end{array}                                           \tag{B-RD62.14}
\]
Their trace difference is
\[
 \zeta+\zeta^p-1-\zeta^{p+1}
                  =-(\zeta-1)(\zeta^p-1)\ne0.
\]
The generator has order $N$ in the finite special-inertia quotient,
which is prime to $p$; no finite-order lift in absolute inertia is
asserted.  Adding the same constant etale group of height $h-2$
preserves the truncation isomorphism and the trace difference.
Since $p^n$ is the manuscript level $\nu=n-1$, failure there gives
the lower bound $\nu_{\min}\ge n$.
\end{proof}

\begin{corollary}[The first ramification step and an infinite sequence]
\label{cor:B-RD62}
\status{Proved; exact lower bounds only, with full independent audits.}
For every $p\ge5$ and $h\ge2$,
\[
                      \nu_{\min}(p,2(p-1),h)\ge2.
                                                        \tag{B-RD62.15}
\]
More generally, for $n\ge2$ put $e_n=p^{2n-3}+p-2$.  Then
$p-1\mid e_n$, $p\nmid e_n$, and
\[
                         \nu_{\min}(p,e_n,h)\ge n.
                                                        \tag{B-RD62.16}
\]
\end{corollary}
\begin{proof}
The integer $e_n$ is the least multiple of $p-1$ strictly greater
than $p^{2n-3}$.  It is strictly smaller than $p^{2n-2}$, and its
residue modulo $p$ is $p-2\ne0$.  Apply the theorem.  For $n=2$
it is $2(p-1)$.
\end{proof}

\begin{remark}[Scope and exact finite checks]
\label{rem:B-RD62-scope}
\status{Proved scope and limitation; no upper bound or exact next index.}
These examples compare a potentially good group with a group having
a two-step potentially good filtration.  They do not assert that
both entire groups are potentially good.  They are compatible with
Theorem~\ref{thm:B-RD61}: the smallest new ramification is $2(p-1)$,
whereas that theorem has $e=p-1$.  The formal group and its exact
descent arc are specified by (B-RD62.4)--(B-RD62.6); the filtered
comparison uses the explicit cohomological lifting argument.
No finite computation is being used to assert effectivity of either
construction.  The exact certificate
\path{code/verify_tame_cyclic_level_counterexamples.py}
checks $15{,}556$ arithmetic cases, $94{,}756$ strict iterate-weight
comparisons, $14{,}592$ logarithm monomials, $40$ independently
recursively computed multiplication coefficients, $24$ finite
character identities, and $264$ first-threshold calculations.
The range is primes $5\le p\le101$ and levels $2\le n\le12$:
all admissible $e$ are tested for $n=2$, and explicitly specified
first, last, and middle windows are tested at higher levels.
The direct multiplication coefficients through degree $p^2$ are
checked for $p=5,7,11,13$.  The actual author run on
19 November 2026 took $0.129167474$ seconds and passed; the complete
source and output are retained.  Full independent root and obstruction/crystal
audits passed, with author-preserving exact reruns taking
$0.137289279$ and $0.123158989$ seconds, respectively.
The upper bound at $e=2(p-1)$
and the general least-index problem remain open here.
\end{remark}


\subsection{The exact ramification threshold for a mixed cyclic comparison}

\begin{lemma}[A rational cyclic point at arbitrary level]
\label{lem:B-RD63-point}
\status{Proved; full root and obstruction/crystal audits passed.}
Let $p\ge5$, let $R$ be a complete discrete valuation ring of mixed
characteristic $(0,p)$ with algebraically closed residue field, and
normalize $v(\pi_R)=1$. Put $a=v(p)$ and assume $p\nmid a$.
If a one-dimensional formal group over $R$ of special height two
has an $R$-rational point of exact order $p^n$, where $n\ge2$, then
\[
                    a>(p^2-1)p^{2n-3}.                 \tag{B-RD63.1}
\]
\end{lemma}
\begin{proof}
Use the integral universal coordinate and coefficient bounds in
the full proof of Lemma~\ref{lem:B-RD49-point}; these come from
\cite[Proposition~5.7; (12.3)--(12.5); Proposition~12.10 and proof]
{HopkinsGross1994Bundles}. Put $q=p^2$, $N=q-1$, and $b=v(u)$,
allowing $b=+\infty$. The only possible vertices of the positive
root Newton polygon of $[p](T)$ are
\[
                         (1,a),\ (p,b),\ (q,0).
\]
All intermediate coefficients lie in $(p,u)$, coefficients in
degrees $2,\ldots,p-1$ vanish, and the degree-$q$ coefficient is a
unit. Later integral coefficients lie strictly above the relevant
negative-slope segments. Thus, if $b\ge ap/(p+1)$, every nonzero
$p$-torsion point has valuation $a/N$. Otherwise its two valuations
are
\[
       \alpha=\frac{a-b}{p-1},\qquad
       \beta=\frac{b}{p(p-1)},                           \tag{B-RD63.2}
\]
with multiplicities $p-1$ and $q-p$, respectively.

We also require the exact inverse-step polygon. For any nonzero
torsion point of valuation $Z<a$, the polynomial obtained by
Weierstrass preparation from $[p](T)-z$ has its possible lower
vertices at $(0,Z),(p,b),(q,0)$; the point $(1,a)$ is strictly
above them. Consequently its $p$-preimages have valuations
\[
 \begin{cases}
 Z/q &\text{if } b\ge Z(p-1)/p,\\
 (Z-b)/p\ \text{($p$ times)},\quad
 \beta\ \text{($q-p$ times)}&\text{if }b<Z(p-1)/p.
 \end{cases}                                           \tag{B-RD63.3}
\]
At equality both displayed slopes agree. Every inverse-step
valuation is strictly smaller than $Z$. Every nonzero $p$-torsion
valuation is at most $a/(p-1)<a$, so (B-RD63.3) applies at every
step of a rational chain. Once a valuation is $\beta$, all further
inverse steps divide it by $q$, because
$b=p(p-1)\beta>\beta(p-1)/p$. The same remains true after any
uniform inverse step whose input satisfies the first inequality
of (B-RD63.3).

Let $P$ have order $p^n$ and let $Z$ be the valuation of
$[p]^{n-1}P$. All points in this chain are rational, so their
positive valuations are integers. In the one-slope case all
inverse steps divide by $q$, giving
$v(P)=a/(Nq^{n-1})$. Its integrality contradicts $p\nmid a$.
In the two-slope case, if $Z=\beta$, then
\[
 b\ge p(p-1)q^{n-1},\qquad
 a>\frac{p+1}{p}b\ge Nq^{n-1},
\]
which is stronger than (B-RD63.1).

It remains that $Z=\alpha$. If $b\ge\alpha(p-1)/p$, all inverse
steps are uniform, and integrality gives $\alpha\ge q^{n-1}$.
Hence
\[
 a=b+(p-1)\alpha\ge\frac{N}{p}\alpha
       \ge Np^{2n-3}.
\]
Equality would make $p\mid a$, so the desired inequality is strict.
If instead $b<\alpha(p-1)/p$, a rational high inverse step of
valuation $A=(\alpha-b)/p$ would imply
\[
                   a=p\bigl(b+(p-1)A\bigr),
\]
again impossible. The first inverse step must therefore have
valuation $\beta$. The remaining $n-2$ steps divide by $q$;
thus $b\ge p(p-1)q^{n-2}$. The strict inequality for this case
is equivalent to $a>(p+1)b$, and gives
$a>Np^{2n-3}$. This exhausts both Newton polygons and every
rational inverse-step choice.
\end{proof}

\begin{proposition}[The mixed-case necessary threshold]
\label{prop:B-RD63-mixed}
\status{Proved; full root and obstruction/crystal audits passed.}
Let $K$ be complete discretely valued of mixed characteristic
$(0,p)$, $p\ge5$, with perfect residue field and absolute
ramification $e$ prime to $p$. Let $M$ be a height-two BT group
which acquires one-dimensional supersingular good reduction,
and let $H$ admit a filtration by two potentially good height-one
groups whose dimensions sum to one. For $n\ge2$,
\[
              M[p^n]\simeq H[p^n]\quad\Longrightarrow
                         e>p^{2n-3}.                   \tag{B-RD63.4}
\]
\end{proposition}
\begin{proof}
Pass to the completed maximal unramified extension, preserving
$e$ and the torsion comparison. Its residue field is algebraically
closed. The full quotient-and-twist proof of
Proposition~\ref{prop:B-RD49-mixed} works at level $p^n$ without
change: transport the first filtered subgroup to a primitive
cyclic subgroup of $M[p^n]$. If its character is a finite
$\epsilon$, twist by $\epsilon^{-1}$. If its character is
$\chi_p\epsilon$, first quotient by that cyclic subgroup and
then twist by the inverse finite quotient character. This produces
a potentially supersingular height-two group with a rational
point of exact order $p^n$. Finite-flat closure in a good model
justifies the quotient, and the determinant is $\chi_p$ times
a finite $\mu_{p-1}$-valued character, exactly as in that proof.

Lemma~\ref{lem:B-RD49-inertia} gives good formal reduction over
a tame extension of degree $m\mid p^2-1$. Its absolute index is
$a=me$, prime to $p$. The rational point extends to the finite
flat torsion model and lies in its formal group, since the special
fibre is connected. Lemma~\ref{lem:B-RD63-point} now gives
$me>(p^2-1)p^{2n-3}$. Since $m\le p^2-1$, (B-RD63.4) follows.
\end{proof}

\begin{theorem}[All indices above the cyclic threshold occur]
\label{thm:B-RD63-arc}
\status{Proved; full root and obstruction/crystal audits passed.}
Let $p\ge5$, $n\ge2$, and let $e$ be a positive multiple of
$p-1$ with $e>p^{2n-3}$. There is a complete mixed-characteristic
field $K$ with residue field $\overline{\mathbf F}_p$, absolute
ramification $e$, and $\mu_p\subset K$, and height-two BT groups
$M,H$ such that $M[p^n]\simeq H[p^n]$, while their requested
prime-to-$p$ special characters differ. Here $M$ is wholly
potentially supersingular, and $H$ has a two-step potentially
good height-one filtration of total dimension one.
\end{theorem}
\begin{proof}
Put $q=p^2$, $N=q-1$, $Q=q^{n-1}$, and $x_0=e/Q$.
Choose an integer $A$ with
\[
                  x_0<A<px_0,\qquad N\nmid A.          \tag{B-RD63.5}
\]
For completeness take $A_0=\lfloor x_0\rfloor+1$, and replace it
by $A_0+1$ only if $N\mid A_0$. If $x_0<1$, the hypothesis
$px_0>1$ makes $A_0=1$ valid. If $x_0\ge1$, the chosen value
satisfies $A\le x_0+2<px_0$ because $p\ge5$; two consecutive
integers cannot both be divisible by $N$.

Use the universal group and exact logarithm from the proof of
Lemma~\ref{lem:B-RD62-arc}, and define
\[
 \begin{gathered}
 d=Ne,\qquad b=d-(p-1)AQ,\qquad Y_0(X)=X^A,\\
 Y_{j+1}(X)=[p]_{X^b}(Y_j(X)),\qquad
 P_A(X)=\Phi_{X^b}(Y_{n-1}(X)).
 \end{gathered}                                        \tag{B-RD63.6}
\]
The two strict inequalities in (B-RD63.5) give exactly
\[
 \begin{gathered}
 b>(q-p)Aq^{n-2}>0,\qquad AQ<d<NAQ,\\
 b\equiv-(p-1)A\pmod N.
 \end{gathered}                                        \tag{B-RD63.7}
\]
Indeed the first difference is $N(e-Ap^{2n-3})$; also
$AQ<pe<Ne=d$. Every coefficient argument of the audited arc
now applies with all input weights multiplied by $A$:
\[
 \begin{gathered}
 Y_j(X)\equiv X^{Aq^j}+\text{higher terms}\pmod p,\\
 P_A(0)=p,\qquad
 P_A(X)\equiv X^d+\text{higher terms}\pmod p.
 \end{gathered}
\]
Thus $R=W(\overline{\mathbf F}_p)[[X]]/(P_A)$ is a complete
Eisenstein DVR with uniformizer $x$ and $v_R(p)=d$.
The strict bound $d>AQ$ prevents $p$-divisible coefficients from
changing any valuation $v_R(Y_j(x))=Aq^j$. Consequently $x^A$
has exact order $p^n$ in $F_{x^b}$.

Choose $\zeta$ of order $N$ and set $a=\zeta^{-A}$.
The exact homogeneous logarithm identities (B-RD62.9)--(B-RD62.10)
hold for this Teichmuller unit too. The substitution
$\sigma(x)=\zeta x$ sends $u=x^b$ to $a^{p-1}u$, sends every
$Y_j$ to $a^{-1}Y_j$, and fixes $P_A$. It has order $N$ on $R$.
The semilinear coordinate map $aT$ fixes $x^A$ and satisfies
the descent cocycle. Therefore the generic BT group and marked
point descend to $K=\operatorname{Frac}(R)^{\langle\sigma\rangle}$,
with $e_K=e$. This degree-$N$ good extension need not be faithful
on the special group; no faithfulness is required here.
The same tame cyclotomic argument as in B-RD62 gives $\mu_p\subset K$.

On the special isocrystal the two eigenvalues are $a,a^p$.
The determinant is $\chi_p\psi$, where
$\psi(\sigma)=a^{p+1}$ has order dividing $p-1$.
Apply the identical finite-flat cyclic quotient, twist by
$\psi^{-1}$, and compatible Kummer-extension lifting of
Theorem~\ref{thm:B-RD62} and Lemma~\ref{lem:B-RD62-extension}.
They produce $M,H$ with exactly the claimed common truncation.
Writing $\lambda=\zeta^A$, their special eigenvalue multisets are
\[
                       \{\lambda,\lambda^p\},\qquad
                       \{1,\lambda^{p+1}\}.
\]
Their trace difference is $-(\lambda-1)(\lambda^p-1)\ne0$, since
$N\nmid A$ and $p$ is coprime to $N$. This is a difference on
an element of prime-to-$p$ order in the finite special inertia
quotient, with no assertion of finite order in absolute inertia.
\end{proof}

\begin{corollary}[Sharp mixed threshold and a general lower bound]
\label{cor:B-RD63}
\status{Proved; full root and obstruction/crystal audits passed.}
For fixed $p\ge5$, $n\ge2$, and a positive multiple $e$ of $p-1$
prime to $p$, a comparison of the two classes in
Proposition~\ref{prop:B-RD63-mixed} can occur over some permitted
field of index $e$ if and only if $e>p^{2n-3}$.
For every positive multiple $e$ of $p-1$ and height $h\ge2$,
the universal B.4 index satisfies
\[
 \nu_{\min}(p,e,h)\ge
 \left\lfloor\frac{\lfloor\log_p(e-1)\rfloor+3}{2}\right\rfloor.
                                                               \tag{B-RD63.8}
\]
\end{corollary}
\begin{proof}
Necessity and existence are exactly the preceding proposition
and theorem; the latter is an existence assertion for some $K$,
not for every field with that index. A failure at actual level
$p^n$ is a failure at manuscript index $\nu=n-1$, so requires
$\nu_{\min}\ge n$. The largest $n\ge2$ with
$p^{2n-3}<e$ is the displayed integer, whenever that integer is
at least two. In the remaining range it equals one, and the
existing level-one height-two examples of
Proposition~\ref{prop:B-RD18}, with ramification a multiple of
$p-1$, give the same lower bound. Adding equal good height-one
factors proves every $h\ge2$ without changing either comparison.
\end{proof}

\begin{remark}[What the threshold does and does not settle]
\label{rem:B-RD63-scope}
\status{Proved; full root and obstruction/crystal audits passed.}
The necessary bound assumes $p\nmid e$ and compares a whole
potentially supersingular group to a filtered height-one group.
It is not an upper bound for comparisons between two potentially
supersingular groups, and does not determine the unrestricted
least B.4 index outside the earlier exact cases. The existence
construction requires $p-1\mid e$ but not $p\nmid e$; it removes
the upper restriction $e<p^{2n-2}$ from B-RD62. The finite-character
and filtration conventions are exactly those of that theorem.

The exact supplement
\path{code/reduction_descent/verify_cyclic_torsion_threshold.py}
checks $433130$ Newton polygons and $1299390$ rational chain steps
for six primes with $a\le400$, $943$ explicit arcs, $5706$ weight
inequalities, and $14360$ integer-index comparisons. It passed in
$1.3598105860000942$ seconds at 20:00:45 UTC. The finite ranges
and actual adjusted choices of $A$ are retained in the output;
no infinite-field or BT existence assertion is inferred from them.
\end{remark}


\subsection{Tame weights on the actual quotient and the next exact height-two index}

\begin{lemma}[An equivariant coefficient bound for the actual target]
\label{lem:B-RD64-weighted-kernel}
\status{Proved; complete independent proof, source and code audits passed.}
Let $p\ge5$, $N=p^2-1$, and let $L/K$ be a totally ramified
cyclic tame extension of degree $N$ of complete mixed-characteristic
$(0,p)$ discrete valuation fields with algebraically closed residue
field $k$. Put $R=\mathcal O_L$, normalize $v_L(\pi)=1$, and
write $a=v_L(p)$. Choose a generator $\tau$ with
$\tau\pi=\zeta\pi$, where $\zeta\in K$ has order $N$.
Let $G/R$ be a finite-height formal group with a semilinear
$\tau$-descent action. Assume $G_k$ is a product of one-dimensional
height-two formal groups, and its cotangent eigenvalues are
$\bar\zeta^{c_1},\ldots,\bar\zeta^{c_d}$ with
\[
                       c_i\equiv c_j\pmod{p-1}
                       \quad\text{for every }i,j.      \tag{B-RD64.1}
\]
Define
\[
 q_0=1,\qquad q_{j+1}=\min\{a+q_j,\ p q_j+p-1\}.
                                                               \tag{B-RD64.2}
\]
If $f:F\to G$ has no nonzero finite flat subgroup in its kernel,
has finite special-fibre kernel, and $q_u\ge a/(p-1)$, then
\[
                         F^u\ker(f_k)=0.              \tag{B-RD64.3}
\]
Here $F^u$ is iterated relative Frobenius. The assumptions and
coefficient bound concern this actual target $G$, not a source
or a group merely isogenous to $G$.
\end{lemma}
\begin{proof}
Lift an eigenbasis of the special cotangent space to parameters,
and average each parameter using
$N^{-1}\sum_{j=0}^{N-1}\zeta^{-jc_i}\tau^j$.
The reductions still form a basis, so the averaged parameters
are integral formal coordinates, with exact ring action
$\tau(X_i)=\zeta^{c_i}X_i$. This uses only that $N$ is a unit
and its Teichmuller roots lie in $K$. For a coefficient
$b_{j,\alpha}$ of $X^\alpha$ in the $j$th coordinate of $[p]_G$,
commutation with this action gives
\[
 \tau(b_{j,\alpha})=
       \zeta^{c_j-\sum_i\alpha_i c_i}b_{j,\alpha}.       \tag{B-RD64.4}
\]
The opposite point-action convention reverses both exponents
and gives the same ensuing divisibility. A nonzero eigen-coefficient
has valuation congruent to its exponent modulo $N$: divide it by
$\pi^{v_L(b_{j,\alpha})}$ and reduce the resulting unit, on which
$\tau$ acts trivially. For $|\alpha|=p$, (B-RD64.1) therefore gives
\[
                         v_L(b_{j,\alpha})\equiv0\pmod{p-1}.
\]
The product-height-two special-fibre condition makes every
coefficient of total degree less than $p^2$ vanish modulo $\pi$.
In particular every nonzero degree-$p$ coefficient has valuation
at least $p-1$.

There is a second coefficient constraint, independent of the
chosen eigenbasis. If $U(X)dX$ is the matrix of normalized
invariant differentials, then
\[
                         U([p]X)J_{[p]}(X)=pU(X).
\]
Thus a monomial with any exponent not divisible by $p$ has
coefficient divisible by $p$. The linear term is $pX$.
Apart from the degree-$p$ monomials, every nonlinear monomial
of total degree less than $2p$ has an exponent not divisible
by $p$. The remaining terms of degree $2p$ through $p^2-1$
have coefficients in $\pi R$, and all later coefficients are
integral. Consequently an augmentation vector of coefficient
valuation at least $q\ge1$ is carried by $[p]_G$ into one of
coefficient valuation at least
\[
 \min\{a+q,\ p-1+pq,\ 1+2pq,\ p^2q\}
             =\min\{a+q,\ p-1+pq\}.                   \tag{B-RD64.5}
\]
The last equality follows from $q\ge1$ and $p\ge5$.
The special order condition is preserved by the averaging
coordinate change: substitutions with invertible linear part
preserve the ideal of total degree at least $p^2$.
Iteration proves the valuations in (B-RD64.2) for
$[p]^j_G(\pi h)$, $h\in tR[[t]]^d$.

For clarity, we give the exact passage to the kernel estimate.
The coefficient $l_\alpha$ of a normalized logarithm, with
$n=|\alpha|>0$, satisfies
\[
 v_L(l_\alpha)\ge-a\lfloor\log_p n\rfloor,
 \qquad \lfloor\log_p n\rfloor\le\frac{n-1}{p-1}.
\]
This follows by integrating the invariant differentials, exactly
as in Lemma~\ref{lem:B-RD34-log}. If $q_u\ge a/(p-1)$, formal
substitution, coefficient by coefficient, therefore gives
\[
               p^u\log_G(\pi h)\in\pi^{q_u}R[[t]]^d.  \tag{B-RD64.6}
\]
Equality at the bound is permitted: each coefficient of an
augmentation-series substitution involves only finitely many terms.

Set $J=\pi^{\lceil a/(p-1)\rceil}R$ and
$\mathcal R=R[[\Delta]][1/p]$. The original Cartier identities are
\[
 \mathcal C_G=\Omega\pi D_{G_\pi},\qquad
 J[[\Delta]]^d\subset D_G,\qquad
 E(G_k)=D_G/(D_G\cap\mathcal R^d).
\]
They are Proposition~10(4), Proposition~9, and Lemma~3 of
\cite[pp.~674--676]{Bondarko2006Generic}; the bounded-projection
identity is Proposition~11(2). Integral Witt-coefficient action
of $\Omega$ preserves every coefficient-valuation ideal, so
(B-RD64.6) implies $p^u\mathcal C_G\subset J\{\{\Delta\}\}^d$.
For $z\in D_G$ and an integral matrix $A$ with $Az\in\mathcal R^b$,
the bounded-projection identity and
$\mathbf fz=p(z-z(0))\Delta^{-1}$ imply
\[
                         \mathbf f^uAz\in J[[\Delta]]^b.
\]
The discarded constant terms lie in $pR\subset J$, exactly as
in (B-RD34.8).

Finally choose an integral complement to
$df(0)(L^{\dim F})\cap R^d$ and its projection $\operatorname{pr}$.
Absence of a nonzero finite flat kernel gives
$df(0)D_F=D_G\cap df(0)(L^{\dim F})[[\Delta]]$.
For a class in the saturation of $E(F_k)$ in $E(G_k)$,
represented by $z\in D_G$, the projected series is bounded.
Thus $\mathbf f^u\operatorname{pr}z\in J[[\Delta]]^d\subset D_G$
is bounded and reduces to zero, while its difference from
$\mathbf f^uz$ belongs to $df(0)D_F$.
The saturation quotient is killed by $\mathbf f^u$.
Proposition~16 with its full proof, p.~681, and Theorem~2(8),
equation~(6), p.~668 of the cited source identify this quotient
with the Dieudonne module of $\ker f_k$, with $\mathbf f$
equal to relative Frobenius. This proves (B-RD64.3).
\end{proof}

\begin{proposition}[Common determinants control the quotient weights]
\label{prop:B-RD64-formal}
\status{Proved; complete independent proof, source and code audits passed.}
Let $p\ge5$ and let $K$ be complete discretely valued of mixed
characteristic $(0,p)$ with algebraically closed residue field
and absolute index $e$. Let $M,N$ be height-two groups with
potentially supersingular good reduction and common determinant
$\chi_p\psi$, where $\psi$ is a finite $\mathbf Z_p^\times$-valued character.
Put $a=e(p^2-1)$ and define $q_j$ by (B-RD64.2).
If $n\ge2$ and $q_{n-1}\ge a/(p-1)$, then an isomorphism
$M[p^n]\simeq N[p^n]$ determines their entire finite special
inertia characters.
\end{proposition}
\begin{proof}
Assume instead that the special characters differ. By
Lemma~\ref{lem:B-RD49-inertia}, both groups have good formal
models over a common tame field of degree dividing $N_0=p^2-1$.
Enlarge it to the unique tame extension $L/K$ of degree $N_0$;
this extension is not required to act faithfully on either
special group. Choose $\pi_L$ with $\pi_L^{N_0}=\pi_K$ and
$\tau\pi_L=\zeta\pi_L$.
Let $M_0,N_0'$ denote the two special formal groups.
Their special eigenvalue multisets are Frobenius orbits
$\{\lambda,\lambda^p\}$ and $\{\mu,\mu^p\}$ in
$\mu_{p^2-1}$, allowing repeated roots. Distinct orbits are
disjoint. The characteristic polynomials over $\mathbf Z_p$
therefore have relatively prime reductions, so an integral
polynomial in $\tau$ projects $M_0\times N_0'$ onto either factor.

Let $H$ be the finite-flat closure of the given generic graph
inside $\mathcal M[p^n]\times\mathcal N[p^n]$.
Over $L$ both Tate determinants are cyclotomic, because the finite
character is killed on the good field. The alternating generic pairings
and their inverse duality maps extend to perfect pairings on
the full good BT models by
\cite[Theorem~4 and Corollaries~1--2, p.~181]{Tate1967PDivisible}.
Scale one pairing by the unit multiplier of the torsion map.
Then $H$ is its own orthogonal: it is isotropic by flatness,
its orthogonal is finite flat of the same rank, and the
equal-rank closed immersion is an isomorphism. This identity
commutes with special-fibre base change by exact finite-flat
Cartier duality, as in the full proof of Proposition~\ref{prop:B-RD33}.

The integral tame projector preserves $H_k$. Hence
$H_k=J_M\times J_N$. The unique subgroup of rank $p^r$ in a
one-dimensional formal group over $k$ is $\ker F^r$:
its ideal is $(T^{p^r})$ in $k[[T]]$. Self-orthogonality in
the $p^n$-torsion, of rank $p^{2n}$ in each factor, forces
both subgroup exponents to equal $n$. Therefore
\[
                     H_k=\ker F_{M_0}^n\times
                                  \ker F_{N_0'}^n.      \tag{B-RD64.7}
\]

Set $Q=(\mathcal M\times\mathcal N)/H$. Quotient formation
commutes with base change, so its special fibre is the actual
product $M_0^{(p^n)}\times (N_0')^{(p^n)}$, with its induced
descent action. This is an isomorphism of special groups with
descent, not just an isogeny. Put $G=Q\times\mathcal M$.
The map
\[
 \mathcal M\times\mathcal N\longrightarrow G,
                \qquad z\longmapsto(qz,\operatorname{pr}_{\mathcal M}z)
                                                               \tag{B-RD64.8}
\]
has no nonzero finite flat subgroup in its kernel, since its
generic graph meets $0\times N$ trivially. Its special kernel
is exactly $0\times\ker F_{N_0'}^n$.

We now verify the crucial hypothesis on this actual target.
Write the special cotangent eigenvalues of $M_0,N_0'$ as
$\bar\zeta^{c_M},\bar\zeta^{c_N}$. Their reduced norms are the
same finite determinant value $\psi(\tau)$, by the covariant
dictionary and determinant calculation in
Proposition~\ref{prop:B-RD4-dictionary} and
Lemma~\ref{lem:B-RD49-inertia}. Thus
$c_M\equiv c_N\pmod{p-1}$. A simultaneous inversion convention
would give the same congruence.
Relative Frobenius of order $n$ raises cotangent characters
to their $p^n$th powers. Consequently the three eigen-exponents
of $G_k$ are
\[
                          p^n c_M,\quad p^n c_N,\quad c_M
                          \pmod{p^2-1}.
\]
They are all congruent modulo $p-1$. Hence
Lemma~\ref{lem:B-RD64-weighted-kernel} applies to (B-RD64.8)
with $u=n-1$. It would kill $\ker F_{N_0'}^n$ by $F^{n-1}$.
Its image instead has rank $p$, a contradiction. This excludes
different special characters under the stated bound. No
thickness of $Q$ was inferred from thickness of its sources;
the degree-$p$ coefficient bound was proved from $Q$'s own
semilinear descent action.
\end{proof}

\begin{theorem}[The next cyclotomic ramification has least index two]
\label{thm:B-RD64}
\status{Proved; complete independent proof, source and code audits passed.}
For every prime $p\ge5$, let $K$ range over complete discretely
valued mixed-characteristic $(0,p)$ fields in the full B.4
monodromy setting, with height two, $\mu_p\subset K$, perfect
residue field and
$e_K=2(p-1)$,
\[
                         \nu_{\min}(p,2(p-1),2)=2.     \tag{B-RD64.9}
\]
In particular the actual $p^3$-torsion always determines the
requested character, while actual $p^2$-torsion need not.
\end{theorem}
\begin{proof}
The lower bound is the already proved Theorem~\ref{thm:B-RD62}
with $n=2$, since $p<2(p-1)<p^2$.
For sufficiency suppose two allowed groups have isomorphic
$p^3$-torsion, and pass to the completed maximal unramified
base. Their generic determinants have the form
$\chi_p^d\epsilon$ with $d\in\{0,1,2\}$ and finite
$\mu_{p-1}$-valued $\epsilon$. The cyclotomic image modulo
$p^3$ has order $p^2$, because $\mu_p\subset K$ and
$p\nmid e_K$. Comparing determinants therefore forces the
same $d$ (their difference has absolute value at most two),
and then the same finite determinant character.

For $d=0$ the special character is the Brauer lift of the
semisimplified generic mod-$p$ inertia representation, with
additivity across the monodromy filtration. Cartier duality
reduces $d=2$ to this case. For $d=1$, each group is either
wholly potentially supersingular or is filtered by two
potentially good height-one groups: the connected--etale
filtration of a potentially ordinary graded piece descends,
as in the full proof of Theorem~\ref{thm:B-RD49}.
Two filtered groups have their finite special characters
determined by the semisimplified generic mod-$p$ representation,
since $\chi_p\equiv1\pmod p$ on inertia.

A mixed comparison is impossible by
Proposition~\ref{prop:B-RD63-mixed} at $n=3$, because
$p\nmid2(p-1)$ and $2(p-1)<p^3$.
It remains to compare two whole potentially supersingular
groups with common determinant. For their enlarged common
tame good field take
\[
 a=2(p-1)(p^2-1),\qquad q_0=1,\quad q_1=2p-1,
 \quad q_2=2p^2-1>2(p^2-1)=a/(p-1).
\]
These are exactly the first two values of (B-RD64.2).
Proposition~\ref{prop:B-RD64-formal} with $n=3$ excludes the
last possible character difference. Thus manuscript index
$\nu=2$ always suffices, and the lower example proves equality.
\end{proof}

\begin{remark}[Scope of the new target argument]
\label{rem:B-RD64-scope}
\status{Proved; complete independent proof, source and code audits passed.}
The coefficient lemma proves additional divisibility on the
actual quotient target by its tame action. It does not assume
that a ramified quotient preserves the Hasse-coefficient
valuation of a source. The complete exact-index assertion is
only for $p\ge5$, $e=2(p-1)$ and height two. The formal-pair
proposition has its own explicit recurrence bound and common
determinant hypothesis; it is not an unrestricted least-index
formula for other heights or ramification indices.

The exact supplement
\path{code/reduction_descent/verify_tame_quotient_weights.py}
checks $78770$ distinct equal-norm orbit pairs, $236310$ actual
quotient-weight triples, $2126790$ degree-$p$ coefficient bounds,
and $1740891$ low-degree monomials for nine primes. It also checks
$1800$ recurrence identities, $45000$ logarithm inequalities and
$1485$ orthogonality cases. The author run passed in
$1.6924098020026577$ seconds at 20:13:57 UTC. Its output retains
the exact scope; these finite checks prove no geometric existence
or universal formal-group theorem. The independent root run passed in
$1.6690704469947377$ seconds, and the independent companion run in
$1.7823184230001061$ seconds, with the author output preserved.
Both audits read the complete proof and original Cartier-module
source passages; the exact prepromotion source is retained.
\end{remark}


% Separate post-CP23 candidate; exact proof and certificate are audited separately.
\subsection{The canonical high quotient at \(p=5,e=12\)}

Throughout this subsection, after the explicitly indicated unramified
base change, $K$ is complete discretely valued of mixed characteristic
$(0,5)$ with algebraically closed residue field, $e_K=12$ and
$\mu_5\subset K$. Put $N=24$. For two wholly potentially supersingular
height-two groups, enlarge the common tame good field to $L/K$ of degree
$N$, and normalize $v_L$ to have value group $\mathbf Z$.
Then $v_L(5)=288$. In the nonsquare finite-determinant case a common
finite height-one twist makes that determinant primitive of order four.
Choose $\sigma\pi_L=\zeta\pi_L$ with $\zeta$ primitive of order $24$.
As in B53 and B60, the actual special tangent exponent $c$ satisfies
$c\equiv1\pmod4$, and its special character is the Frobenius orbit
$\{c,5c\}$ modulo $24$.

\begin{lemma}[Newton and point-field rules at index twelve]
\label{lem:B-RD65-Newton}
\status{Proved; full independent proof, source and code audits passed.}
In the nonsquare normalized case, a reducible residual module has a
unique high-valuation line $H$ and an integer $B$ satisfying
\[
 \begin{gathered}
 0<B<60,\qquad B\equiv c\pmod6,\\
 \alpha=v_L(T(Q))=72-B\quad(0\ne Q\in H),\\
 \beta=v_L(T(Q))=B/5\quad(Q\notin H).
 \end{gathered}                                      \tag{B-RD65.1}
\]
The high line is rational over $L$ and its tame character is
$\zeta^{c-B}$. Every point $P$ of order $25$ with $5P\notin H$
has valuation $B/125$.
For two such groups with isomorphic actual $25$-torsion:
\begin{enumerate}
\item If their residual wild image is nontrivial, their high lines
coincide. If neither $B_i$ is divisible by $5$, then $B_1=B_2$ and
$c_1=c_2$ modulo $24$. It is impossible that exactly one $B_i$ is
divisible by $5$.
\item If their residual wild image is trivial, write $B_i=5b_i$.
If both $b_i$ are prime to $5$, then $b_1=b_2$, their high lines
coincide, and $c_1=c_2$ modulo $24$.
\end{enumerate}
\end{lemma}
\begin{proof}
We repeat the calculations of B60--B61 with the actual new absolute
index, rather than applying their minimal-ramification statements
outside their scope. In the universal height-two coordinate,
$[5](T)$ has the three Newton vertices
$(1,288),(5,b),(25,0)$, where $b=v_L(u)$; the intermediate coefficient
bounds are exactly those established from Hopkins--Gross in
Lemma~\ref{lem:B-RD49-point}. If $b\ge240$, all nonzero $5$-torsion
points have valuation $12$. Their leading-coordinate map is injective,
so wild inertia acts trivially. The tame representation on its
leading plane is irreducible over $\mathbf F_5$ or scalar; the latter
would have square determinant, contrary to the normalization. Thus
reducibility implies $0<b<240$.

The high line has valuation $(288-b)/4$, and its character has values
in $\mathbf F_5^\times$, so it is killed by the tame good field $L$.
Its coordinate valuation is integral. Write $b=4B$; this gives
$\alpha=72-B$, $\beta=B/5$ and $0<B<60$. Tame averaging of the
formal parameter preserves the degree-five coefficient valuation
$b<288$. Compatibility with the exact action $T\mapsto\zeta^c\sigma(T)$
gives $b\equiv4c\pmod{24}$, hence $B\equiv c\pmod6$.
The high point's leading coefficient gives the character
$\zeta^{c+72-B}=\zeta^{c-B}$. In particular $B$ is odd.
For a target of valuation $\beta=B/5<12$, the inverse equation
$[5](T)=T(5P)$ has a single Newton segment of degree $25$, because
both $288$ and $4B$ exceed $\beta$. Thus $v_L(T(P))=B/125$.

Let $R=L(M[5])$. A reducible residual image over $L$ is either
trivial or cyclic unipotent of order five. In the latter case $H$
is its unique fixed line. For an element $g\ne1$ of
$\operatorname{Gal}(R/L)$ and a low point $Q$, one has
$v_R(T(Q))=B$ and
$v_R(T(gQ)-T(Q))=5\alpha=5(72-B)$.
The exact lower-break lemma~\ref{lem:B-RD60-break} therefore gives
\[
 r=360-6B\quad(5\nmid B),\qquad
 r<360-6B\quad(5\mid B),                              \tag{B-RD65.2}
\]
where $r$ is the single actual lower break of $R/L$.
Two prime-to-five $B_i$ are equal from the common break, and then
from the common high-line character their $c_i$ are equal.
If precisely one $B_i$ is divisible by five, index so that
$5\nmid B_1$ and $5\mid B_2$. Equation (B-RD65.2) forces $B_1>B_2$.

Choose corresponding $P_i$ with $5P_i$ outside the common high line.
Over $R$, their valuations are $B_i/25$. On the first side the
valuation denominator forces the point field $E=R(P_i)$ to have
degree at least $25$, and the elementary-abelian point-field
Lemma~\ref{lem:B-RD61-point-field} bounds it by $25$.
Thus the common field is Galois of degree $25$, every $5$-torsion
translation occurs, $[E:L]=125$, and $v_E(T_i(P_i))=B_i$.
Choose a translation outside $H$. Its difference has valuation
$25B_i$ in model $i$. The first coordinate has prime-to-five
valuation, so its actual lower break is $24B_1$. Applying only the
lower-bound part of the same lemma to the second coordinate gives
$25B_2\ge B_2+24B_1$, contradicting $B_1>B_2$.

In the tame residual case $R=L$ and $B_i=5b_i$ is necessary for
the low point's valuation to be integral. Choose a common residual
vector outside both high lines and lift it to $P_i$ of order $25$.
Their valuations are $b_i/25$. If both $b_i$ are prime to five,
the common point field has degree $25$ over $L$ and both normalized
coordinate valuations equal $b_i$. A translation outside both high
lines gives the same actual break $24b_1=24b_2$. Thus $b_1=b_2=b$.
If their high lines differed, a translation high only in the first
model would have respective difference valuations $25\alpha_1$
and $25b$, while the two coordinates have the same prime-to-five
valuation $b$. The equality part of the break lemma would force
$\alpha_1=b$, contradicting the strict high/low inequality.
Their high lines therefore agree, and their characters give $c_1=c_2$.
\end{proof}

\begin{lemma}[The exact canonical quotient transformation]
\label{lem:B-RD65-quotient}
\status{Proved; full independent proof, source and code audits passed.}
Let a group in Lemma~\ref{lem:B-RD65-Newton} have parameter $B\ne12$,
and let $M'=M/H$ be its quotient by the high line. It is again wholly
potentially supersingular. Its unequal Newton parameter is
\[
 B'=\begin{cases}5B,&B<12,\\72-B,&B>12.\end{cases}      \tag{B-RD65.3}
\]
Its special character is the same as that of $M$.
If two actual $125$-torsion modules are isomorphic and the isomorphism
preserves their high lines, then the actual $25$-torsion modules of
their high quotients are isomorphic.
\end{lemma}
\begin{proof}
Over the good DVR take the finite flat closure of $H$. Its quotient
is a good height-two formal group, isogenous to the original group.
Choose integral parameters for the source and quotient. The quotient
map has special degree five and, by Weierstrass preparation, its
parameter series is a unit times
$\prod_{Q\in H}(T-T(Q))$. The unit has valuation zero at every
positive-valuation point. In particular, if $v_L(T(P))<\alpha$,
then the quotient coordinate of $P$ has valuation $5v_L(T(P))$.
This uses the actual integral quotient, without a scaling of its
uniformizer or an identification of its Hasse coefficient with
that of the source.

The image of $M[5]/H$ is a line of $M'[5]$ with nonzero coordinate
valuation $5\beta=B$. All other points of $M'[5]$ come from points
$P$ with $5P$ a nonzero member of $H$.
If $B>12$, the inverse Newton equation with target valuation
$\alpha=72-B$ has one segment: the condition is
$4B>4\alpha/5$, equivalent to $B>12$.
Every preimage has valuation $\alpha/25<\alpha$, so the other
$20$ quotient points have valuation $\alpha/5$.
Since $B>\alpha/5$, the high and low valuations in the quotient
are $B$ and $(72-B)/5$, which gives $B'=72-B$.

If $B<12$, that inverse equation has five roots of valuation
$(\alpha-4B)/5=(72-5B)/5$ and twenty roots of valuation $B/5$.
Both valuations are less than $\alpha$. After quotienting by $H$,
the high preimages give four points of valuation $72-5B$, and
all other nonzero quotient points have valuation $B$.
(The latter include the four points of $M[5]/H$.) Hence $B'=5B$.
These counts follow by taking all four nonzero target points of $H$
and dividing the preimage counts by five. They also show that the
claimed unequal valuations are strict. The excluded value $B=12$
cannot in fact occur under $B\equiv c\pmod6$, $c\equiv1\pmod4$.

The special isogeny identifies the rational Dieudonne modules
with their descent actions, so the requested special characters
are unchanged. More concretely its special kernel is the unique
rank-five subgroup, and its target is the Frobenius twist of the
special group; the new tangent exponent is $5c$ modulo $24$, which
has the same unordered Frobenius orbit.
Finally the inverse image of $M'[25]$ is the actual finite subgroup
\[
       \{P\in M[125]:25P\in H\}.
\]
Dividing it by $H$ gives $M'[25]$ as a finite generic group scheme.
A $125$-torsion isomorphism preserving $H$ therefore induces the
asserted $25$-torsion isomorphism. This is the exact quotient
construction, not a comparison of isogeny classes of lattices.
\end{proof}

\begin{proposition}[All nonsquare formal comparisons at index twelve]
\label{prop:B-RD65-nonsquare}
\status{Proved; full independent proof, source and code audits passed.}
Two normalized potentially supersingular height-two groups over $K$
with isomorphic actual $125$-torsion have the same finite special
inertia character.
\end{proposition}
\begin{proof}
First consider an irreducible residual module. An unequal Newton
polygon would provide an invariant high line, so both polygons
have one slope. The leading-coordinate calculation gives residual
tame eigenvalue orbit $\{c+12,5(c+12)\}$ modulo $24$.
As $5\cdot12\equiv12\pmod{24}$, equality of these orbits is exactly
equality of the special orbits $\{c,5c\}$.
A one-slope reducible representation would be scalar and have
square determinant, already excluded. It remains to consider the
unequal polygons of Lemma~\ref{lem:B-RD65-Newton}.

If residual wild inertia is nontrivial, its fixed high line is
common. That lemma handles the cases where either $B_i$ is prime
to five. If both are divisible, then $B_i=5b_i$ with
\[
                       b_i\in\{1,3,5,7,9,11\}.        \tag{B-RD65.4}
\]
Quotient by the common high line. Their $25$-torsion modules remain
isomorphic by Lemma~\ref{lem:B-RD65-quotient}.
For every $B_i\ge15$, the new parameter $72-B_i$ is prime to five.
The only smaller possibility is $B_i=5$, for which $B_i'=25$.
If both original parameters are $5$, their common high-line
character already identifies $c_i$. If both are at least $15$,
both quotient parameters are prime to five; their nonintegral low
valuations force nontrivial residual wild inertia in the quotients.
Lemma~\ref{lem:B-RD65-Newton} makes their parameters and special
characters equal. If exactly one original parameter is $5$, the
quotients have one divisible and one prime-to-five parameter, so
the same lemma forbids their $25$-torsion isomorphism.
Isogeny invariance transfers the conclusion to the original pair.

Suppose instead residual inertia is tame. Again write $B_i=5b_i$
with the list (B-RD65.4). If neither $b_i$ is divisible by five,
Lemma~\ref{lem:B-RD65-Newton} identifies the characters.
The only exceptional value is $b=5$, namely $B=25$.
If both parameters have this value, their allowed tangent exponents
are $1$ or $13$ modulo $24$. Their high-line characters are respectively
$1$ or $-1$, and their common finite determinant is primitive of
order four, say $\eta$. The residual unordered characters are
therefore $\{1,\eta\}$ and $\{-1,-\eta\}$, respectively. These sets
are distinct, so a common residual module identifies the special
character.

Finally suppose $b_2=5$ and $b_1=b$ is prime to five. Choose a common
residual vector outside both high lines and lift it to $P_i$ of
order $25$. The common point field $E=L(P_i)$ has degree $25$,
forced by the first valuation $b/25$, and all translations occur.
The two coordinates have normalized valuations $b$ and $5$.
If the high lines differed, choose a translation high in model one
and low in model two. The first coordinate's prime-to-five valuation
gives lower break $25(72-5b)-b$, whereas the lower-bound inequality
for the second coordinate would require
\[
                 125\ge5+25(72-5b)-b.                 \tag{B-RD65.5}
\]
This is impossible for every $b\in\{1,3,7,9,11\}$: its right side
is at least $419$.
Thus the high lines agree. A translation outside that line has
break $24b$ in the first coordinate, so the second coordinate gives
$125\ge5+24b$. Hence $b\le5$, leaving $b=1$ or $3$.
Their high-line characters now give
\[
                        c_1-c_2\equiv5b-25\pmod{24}.
                                                        \tag{B-RD65.6}
\]
Since both $c_i\equiv1\pmod4$, this forces $b\equiv1\pmod4$,
so only $b=1$ remains. Then $c_2=1$ or $13$, while
$c_1=c_2+4$ is $5$ or $17$, respectively. These are the same
Frobenius orbits. This exhausts the tame case.
\end{proof}

\begin{lemma}[The square-determinant formal comparison at index twelve]
\label{lem:B-RD65-square}
\status{Proved; full independent proof, source and code audits passed.}
Two wholly potentially supersingular height-two groups over $K$
with common square finite determinant and isomorphic actual
$125$-torsion have the same finite special inertia character.
\end{lemma}
\begin{proof}
Twist both groups by the same finite square root to make their
full determinants cyclotomic. Their finite special eigenvalues
then have norm one, so their common faithful good field is tame
of degree $m\mid6$, as in Proposition~\ref{prop:B-RD50-square}.
Its absolute index is $a=12m\le72$.
If the special characters differ, their two Frobenius eigenvalue
orbits are disjoint. The integral tame projector and the actual
perfect pairings of Proposition~\ref{prop:B-RD64-formal} make the
special graph closure in the $125$-torsion equal to
$\ker F_M^3\times\ker F_N^3$. The actual quotient $Q$ has special
fibre $M_0^{(5^3)}\times N_0^{(5^3)}$.
The map $M\times N\to Q\times M$ used there has no nonzero finite
flat kernel and has special kernel $0\times\ker F_N^3$.

We check the required logarithm bound on its actual target
$G=Q\times M$, without any transferred source thickness.
The invariant-differential identity makes all nonlinear coefficients
of total degree below five divisible by five; the product-height-two
special fibre makes coefficients of degrees five through twenty-four
lie in the maximal ideal. Starting with an augmentation vector of
valuation at least one, two applications of $[5]$ therefore give
coefficient valuations at least
\[
 q_1\ge6,\qquad q_2\ge\min\{a+6,31\}\ge a/4,
                    \qquad 12\le a\le72.                  \tag{B-RD65.7}
\]
The degree-at-least-twenty-five terms are larger throughout.
The closed coefficientwise logarithm estimate and the exact
Cartier bounded-projection/saturation argument of
Lemma~\ref{lem:B-RD64-weighted-kernel} now apply with two
Cartier Frobenius operators. Only its coefficient estimate has
been replaced by (B-RD65.7); its integral module argument is
unchanged. They force $F^2\ker f_k=0$, whereas
$F^2\ker F_N^3$ has rank five. This contradiction proves the lemma.
\end{proof}

\begin{theorem}[The next five-adic least index]
\label{thm:B-RD65}
\status{Proved; full independent proof, source and code audits passed.}
In the full B.4 monodromy setting, for height two, complete mixed
characteristic $(0,5)$ fields with perfect residue, $\mu_5\subset K$
and $e_K=12$, the exact least manuscript index is
\[
                           \nu_{\min}(5,12,2)=2.       \tag{B-RD65.8}
\]
\end{theorem}
\begin{proof}
The counterexample of Theorem~\ref{thm:B-RD63-arc} at $n=2$
exists because $12>5$, so index one need not suffice.
For an actual $125$-torsion comparison, pass to the completed
maximal unramified base. As $5\nmid12$, the cyclotomic image
modulo $125$ has order $25$. The determinant argument of
Theorem~\ref{thm:B-RD64} identifies the total dimension and the
finite determinant. Dimensions zero and two, and comparisons of
two filtered height-one systems, are determined by generic
mod-five inertia as proved there. A mixed whole-supersingular
versus filtered comparison is impossible by
Proposition~\ref{prop:B-RD63-mixed}, since $12\le5^3$.
For two whole potentially supersingular groups, the square
finite-determinant case is Lemma~\ref{lem:B-RD65-square}.
In the nonsquare case a common finite twist gives exactly the
normalization used above, and Proposition~\ref{prop:B-RD65-nonsquare}
identifies the characters. Every allowed comparison is covered,
so index two suffices and is least.
\end{proof}

\begin{remark}[Scope and retained finite checks]
\label{rem:B-RD65-scope}
\status{Proved; full independent proof, source and code audits passed.}
This argument uses the actual high-line quotient and the exact
$125$-to-$25$ subgroup quotient. It does not infer a rank-one lift
from a small Hasse valuation, and does not transfer thickness to
an isogenous quotient. The Newton and point-field rules are proved
at the new absolute index $288$. The finite supplement checks
only their displayed arithmetic and case exhaustion; it is not
used as a substitute for the formal-group or Galois proofs.
No exact least index at $p=5,e=16$ or at other new ramification
indices is asserted here.
The exact supplement
\path{code/reduction_descent/verify_p5_e12_high_quotient.py}
checks $60$ normalized Newton records, $120$ inverse polygons,
$60$ quotient point counts, $4716$ common-break rows (including
$4524$ divisible-parameter rows), $72$ tame residual pairs,
$25$ exceptional inequalities, $36$ one-slope orbit comparisons
and four square-determinant logarithm bounds. The author run passed
in $0.008917200997530017$ seconds at 20:33:10 UTC.
The independent root and companion full audits passed; their
author-output-preserving reruns took $0.007177088$ and
$0.007459839995135553$ seconds, respectively. The exact audited
prepromotion source and all programs and outputs are retained.
\end{remark}


% Proved after full root and companion audits; exact candidate retained in notes/.
\subsection{A uniform middle-ramification quotient argument}

Fix $p\ge5$ and put
\[
 N=p^2-1,\qquad e_0=N/2,\qquad s=e_0(p+1).
                                                        \tag{B-RD66.1}
\]
After completed maximal unramified base change, $K$ is complete of
mixed characteristic $(0,p)$ with algebraically closed residue field,
$e_K=e_0$ and $\mu_p\subset K$. For a pair of whole potentially
supersingular height-two groups take a common tame good field
$L/K$ of degree $N$. Thus $v_L(p)=e_0N=(p-1)s$.
In the nonsquare finite-determinant case normalize by a common
finite height-one twist as in B50, so an actual special tangent
exponent $c$ satisfies $c\equiv1\pmod{p-1}$. Choose
$\sigma\pi_L=\zeta\pi_L$, with $\zeta$ primitive of order $N$.

\begin{lemma}[The unequal Newton comparison rules at middle ramification]
\label{lem:B-RD66-rules}
\status{Proved; full independent root and companion proof, source and exact-code audits passed.}
A reducible normalized residual module has a unique high line $H$
and an integer $B$ with
\[
 0<B<pe_0,\quad B\equiv c\pmod{p+1},\quad
 \alpha=s-B,\quad\beta=B/p,\quad\alpha>\beta.          \tag{B-RD66.2}
\]
Here nonzero points of $H$ have valuation $\alpha$, all other
nonzero $p$-torsion points have valuation $\beta$, and the high-line
character is $\zeta^{c-B}$. The integer $B$ is odd. If $P$ has
order $p^2$ and $pP\notin H$, then $v_L(T(P))=B/p^3$.
For two such groups with isomorphic actual $p^2$-torsion:
\begin{enumerate}
\item In the residual wild case their high lines agree. If both $B_i$
are prime to $p$, they are equal and the special characters agree.
Exactly one $B_i$ cannot be divisible by $p$.
\item In the residual tame case write $B_i=pb_i$. If both $b_i$ are
prime to $p$, then $b_1=b_2$, their high lines agree, and the special
characters agree.
\end{enumerate}
\end{lemma}
\begin{proof}
The three Newton vertices of $[p](T)$ are
$(1,(p-1)s),(p,b),(p^2,0)$ with $b=v_L(u)$, using exactly the
universal coefficient bounds proved in Lemma~\ref{lem:B-RD49-point}.
If the polygon has one slope, its leading-coordinate residual
representation is tame and either irreducible over $\mathbf F_p$
or scalar. The scalar case has square determinant, so is excluded.
Reducibility therefore gives two slopes. The high line is inertia
stable and its character takes values in $\mathbf F_p^\times$;
it is rational over $L$. Consequently its valuation is integral,
so $b=(p-1)B$ with $B\in\mathbf Z$. The two slopes give
$\alpha=s-B$ and $\beta=B/p$, and their strict order gives
$B<ps/(p+1)=pe_0$. Exact tame averaging gives
$b\equiv(p-1)c\pmod N$, hence $B\equiv c\pmod{p+1}$.
Since $s=((p+1)/2)N$ is a multiple of $N$, the high-line character
is $\zeta^{c+s-B}=\zeta^{c-B}$. In particular $B$ is odd.
For target valuation $\beta$, the inverse Newton polygon has one
segment of degree $p^2$, giving $B/p^3$ for every preimage.

For completeness the comparison rules are reproved at this index.
The common residual division field $R/L$ is trivial or cyclic of
order $p$. In the wild case its high line is its unique fixed line.
For a low point, $v_R(T(Q))=B$ and its difference under a nontrivial
wild element has valuation $p\alpha=p(s-B)$. The exact lower-break
lemma~\ref{lem:B-RD60-break} gives
\[
 r=ps-(p+1)B\ (p\nmid B),\qquad
 r<ps-(p+1)B\ (p\mid B).                              \tag{B-RD66.3}
\]
Two prime-to-$p$ parameters are equal from the same actual break;
their high-line characters then identify $c$ modulo $N$.
If only $B_2$ is divisible by $p$, the common break forces
$B_1>B_2$. Choose corresponding $P_i$ with $pP_i$ outside the
common high line. Over $R$, their valuations are $B_i/p^2$.
The prime-to-$p$ denominator on the first side and the elementary
point-field lemma force the common $E=R(P_i)$ to have degree $p^2$;
all translations in $M[p]$ occur. Thus $[E:L]=p^3$ and the point
coordinates have valuations $B_i$. A translation outside the high
line has difference valuations $p^2B_i$. Its actual break from the
first coordinate is $NB_1$, while the second coordinate would require
$p^2B_2\ge B_2+NB_1$, a contradiction.

In the tame case $R=L$, so $B_i=pb_i$. Choose $pP_i$ outside both
high lines. If both $b_i$ are prime to $p$, their common point field
has degree $p^2$, and their coordinate valuations are $b_i$.
A translation outside both high lines gives the same break
$Nb_1=Nb_2$. Thus $b_1=b_2=b$. A translation high only in one
model, if it existed, would force the strictly higher valuation
$\alpha$ to equal $\beta=b$ by the same prime-valuation equality.
The high lines therefore agree and their characters give $c_1=c_2$.
This proves both rules without a minimal-ramification assumption.
\end{proof}

\begin{lemma}[High and low cyclic quotient coordinates]
\label{lem:B-RD66-quotients}
\status{Proved; full independent root and companion proof, source and exact-code audits passed.}
For a group as in Lemma~\ref{lem:B-RD66-rules}, quotient by its
high line transforms its Newton parameter by
\[
 B_{\rm high}=\begin{cases}pB,&B<e_0,\\s-B,&B>e_0.\end{cases}
                                                        \tag{B-RD66.4}
\]
The equality case is impossible since $e_0$ is even and $B$ odd.
If the full residual module is rational over $L$ and $J$ is an
inertia-stable low line, quotient by $J$ transforms it by
\[
                         B_{\rm low}=B/p.             \tag{B-RD66.5}
\]
These are actual good formal quotients over $\mathcal O_L$;
their generic quotients descend to $K$. Their special tangent
exponent is $pc$ modulo $N$, so their finite special character
is unchanged. An actual $p^3$-torsion isomorphism preserving the
chosen lines induces an actual $p^2$-torsion isomorphism on the
quotients.
\end{lemma}
\begin{proof}
The finite flat closure of a chosen line has rank $p$. The good
formal quotient exists, is isogenous to the original height-two
group, and its parameter map is a unit times
$\prod_{Q\in J}(T-T(Q))$ by Weierstrass preparation. Thus the
quotient coordinate valuation is the sum of these difference
valuations. All statements use integral quotient coordinates.

For $J=H$, a point of valuation less than $\alpha$ acquires
valuation multiplied by $p$. The line $M[p]/H$ therefore has
valuation $B$. The other points arise from preimages of nonzero
high points. If $B>e_0$, the inverse Newton polygon has one segment,
with all roots of valuation $\alpha/p^2$, giving $p^2-p$ quotient
points of valuation $\alpha/p$. Since $B>\alpha/p$, this is
$B_{\rm high}=s-B$.
If $B<e_0$, the inverse polygon has $p$ roots of valuation
$(s-pB)/p$ and $p^2-p$ roots of valuation $B/p$.
Both are less than $\alpha$. Taking all $p-1$ nonzero high targets
and dividing the preimage counts by $p$ gives $p-1$ quotient points
of valuation $s-pB$ and $p^2-p$ of valuation $B$.
This gives $B_{\rm high}=pB$.

For a low line $J$, all its nonzero points have valuation $\beta$.
Every nonzero class in $M[p]/J$ has a unique high representative,
so its quotient valuation is
$\alpha+(p-1)\beta=s-B/p$. The other quotient $p$-torsion points
arise from preimages of nonzero low points. They have source
valuation $\beta/p^2<\beta$, hence quotient valuation
$\beta/p=B/p^2$. Their number is $p^2-p$, yielding (B-RD66.5).
Here $B/p$ is an integer because $M[p]$ is rational over $L$.

The special rank-$p$ subgroup of a one-dimensional connected
formal group is its Frobenius kernel. The special quotient is
therefore the Frobenius twist with inherited tangent action $pc$,
and its rational special character is unchanged. The generic
quotient descends because the chosen line was inertia stable.
Finally the inverse image of the quotient's $p^2$-torsion is
$\{P\in M[p^3]:p^2P\in J\}$; quotienting this actual finite group
by $J$ proves the asserted induced torsion isomorphism.
\end{proof}

\begin{proposition}[Nonsquare formal recovery at middle ramification]
\label{prop:B-RD66-nonsquare}
\status{Proved; full independent root and companion proof, source and exact-code audits passed.}
Two normalized whole potentially supersingular height-two groups
with isomorphic actual $p^3$-torsion have the same finite special
inertia character.
\end{proposition}
\begin{proof}
For an irreducible residual module both Newton polygons have one
slope. The leading character is obtained from the special tangent
exponent by adding $e_0=N/2$, a scalar sign in $\mathbf F_p^\times$.
Since $pe_0\equiv e_0\pmod N$, the residual Frobenius orbit
$\{c+e_0,p(c+e_0)\}$ determines $\{c,pc\}$.
A reducible one-slope module has square determinant, so the remaining
cases have the unequal parameters of Lemma~\ref{lem:B-RD66-rules}.

If residual wild inertia is nontrivial, the high line is common.
The rules already handle every case where at least one $B_i$ is
prime to $p$. Suppose both are divisible. If both $B_i<e_0$, put
$B_i=pb_i$. Then $0<b_i<e_0/p<p/2$. The common high-line character
and $c_i\equiv1\pmod{p-1}$ imply
$b_1\equiv b_2\pmod{p-1}$, hence $b_1=b_2$ and $c_1=c_2$ modulo $N$.
If both $B_i>e_0$, quotient by their common high line. The parameters
$s-B_i$ are prime to $p$ because $p\nmid s$, and the actual
$p^2$-torsion remains isomorphic. The rules make the quotient
parameters and characters equal. If one $B_i$ lies on each side
of $e_0$, the quotient parameters have opposite $p$-divisibility,
so the rules forbid their $p^2$-torsion isomorphism. This exhausts
the wild case.

Now let residual inertia be tame, and write $B_i=pb_i$ with
$0<b_i<e_0$. If both $b_i$ are prime to $p$, the rules identify the
characters. Suppose first that $b_1=b$ is prime to $p$ and
$b_2=pk$. Then $0<k<e_0/p<p/2$ and $B_2=p^2k>e_0$.
Choose corresponding order-$p^2$ points with residual targets
outside both high lines. Their common point field has degree $p^2$
over $L$, forced by the coordinate valuation $b/p^2$ in the first
model. All translations occur, and the normalized point valuations
are $b$ and $pk$.
If their high lines differed, choose a translation high in the first
model and low in the second. The prime-valuation coordinate gives
actual break $p^2\alpha_1-b$, while the lower-bound inequality for
the other would require
\[
                  p^2(\alpha_1-pk)\le b-pk.            \tag{B-RD66.6}
\]
But $\alpha_1=s-pb\ge e_0+p$, $pk<e_0$ and $b<e_0$;
the left side is greater than $p^3$, while the right side is less
than $e_0<p^2$. Thus the high lines coincide.
Quotient by that common line. The second parameter $s-p^2k$ is
prime to $p$. If $B_1<e_0$, the first quotient parameter $pB_1$
is divisible by $p$, impossible by the rules and the inherited
$p^2$-torsion isomorphism. If $B_1>e_0$, both new parameters are
prime to $p$; their equality forces $B_1=B_2$, contrary to
$p\nmid b$ and $b_2=pk$. Thus the mixed $b$-divisibility case is
impossible.

Finally suppose $b_i=pk_i$ for both groups. Then
$B_i=p^2k_i>e_0$. If their high lines agree, quotienting by that
line gives two prime-to-$p$ parameters $s-B_i$, which the rules
identify. If their high lines differ, quotient the first group
by its high line and the second by the corresponding low line.
Both chosen lines are inertia stable. The quotient parameters are
$s-B_1$, prime to $p$, and $B_2/p=pk_2$, divisible by $p$.
Their actual $p^2$-torsion isomorphism is forbidden by the rules.
These alternatives exhaust tame residual inertia. Isogeny invariance
in every quotient comparison proves the original character equality.
\end{proof}

\begin{theorem}[A uniform additional exact least-index family]
\label{thm:B-RD66}
\status{Proved; full independent root and companion proof, source and exact-code audits passed.}
For every $p\ge5$, in the full B.4 monodromy setting with height two,
perfect residue field and $\mu_p\subset K$,
\[
                    \nu_{\min}\bigl(p,(p^2-1)/2,2\bigr)=2.
                                                        \tag{B-RD66.7}
\]
\end{theorem}
\begin{proof}
The lower example is Theorem~\ref{thm:B-RD63-arc} with $n=2$,
since $e_0>p$ and $p-1$ divides $e_0$.
For sufficiency pass to the completed maximal unramified base.
Since $p\nmid e_0$, the cyclotomic image modulo $p^3$ has order $p^2$.
The determinant and dimension argument of Theorem~\ref{thm:B-RD64}
reduces to two whole potentially supersingular groups with a common
finite determinant: dimension zero/two and two filtered systems are
already determined; the mixed comparison is impossible by
Proposition~\ref{prop:B-RD63-mixed}, since $e_0<p^3$.
The nonsquare finite determinant case is the preceding proposition.

If the finite determinant is square, twist both groups by its finite
square root. Their common good field has tame degree $m\mid p+1$,
so $a=v_L(p)\le e_0(p+1)$. Different special characters again give
an integral tame projector and the self-orthogonal graph closure
$\ker F_M^3\times\ker F_N^3$, exactly as in
Lemma~\ref{lem:B-RD65-square}. The actual target $Q\times M$ of
the no-flat-kernel map there has product height-two special fibre.
Two applications of $[p]$ to a coefficientwise valuation-one vector
have valuations at least
\[
 q_1\ge p+1,\qquad q_2\ge\min\{a+p+1,p^2+p+1\}\ge a/(p-1).
                                                        \tag{B-RD66.8}
\]
Indeed $a\ge e_0\ge p$, and
\[
     a/(p-1)\le(p+1)^2/2\le p^2+p+1.
\]
The closed logarithm and actual Cartier saturation proof already
established in B64 forces $F^2$ to annihilate the special kernel
$\ker F_N^3$, which is impossible. This is an actual-target bound,
not inherited Hasse thickness. Thus square and nonsquare cases
both recover the character at actual $p^3$-torsion. Together with
the lower example this proves the exact manuscript index.
\end{proof}

\begin{remark}[Scope of the uniform quotient argument]
\label{rem:B-RD66-scope}
\status{Proved; full independent root and companion proof, source and exact-code audits passed.}
The family includes B65 at $p=5,e_0=12$, and for example gives
$e_0=24$ at $p=7$. It does not settle $p=5,e=16$. The proof uses
both high and low stable cyclic quotients, with their actual
$p^3$-to-$p^2$ comparison, and a strict bound on the small integer
$B/p$; no enumeration in finitely many primes proves the theorem.
The exact supplement
\path{code/reduction_descent/verify_middle_tame_quotients.py}
checks $20490$ Newton/high-quotient records, $834$ low quotients,
$40980$ tangent weight transformations, $761$ small-parameter
congruences, $1659$ wild divisible transformations, $18822$ mixed
tame comparisons, $761$ two-divisible tame comparisons, $3664$
one-slope orbit comparisons, and $52$ square-target bounds at
nine primes. Its author run passed in $0.8840610159968492$ seconds
at 20:40:55 UTC. The complete independent root and companion audits passed, with author-output-preserving reruns in $0.8864672490017256$ and $0.7995483489939943$ seconds, respectively. These checks are supplementary to the all-prime proof.
\end{remark}


% Proved after complete root and companion audits; exact candidate retained in notes/.
\subsection{Irreducible residual inertia at arbitrary prime-to-\(p\) ramification}

Throughout this subsection \(p\ge5\), \(q=p^2\), \(N=q-1\), and
\(K\) is complete discretely valued of mixed characteristic \((0,p)\)
with algebraically closed residue field \(k\), \(\mu_p\subset K\),
and absolute ramification \(e\) prime to \(p\). The two height-two
groups \(M_1,M_2\) acquire one-dimensional formal good reduction,
and their actual \(p^2\)-torsion is isomorphic. Their common residual
inertia representation is irreducible over \(\mathbf F_p\).
We prove equality of their finite special inertia characters.
No restriction \(e=p-1\), CM hypothesis, or change of Tate lattice
is imposed.

\begin{lemma}[Common fields, coordinates, and the new normalization]
\label{lem:B-RD67-normalization}
\status{Proved; full independent root and companion proof, original-source and exact-code audits passed.}
Take the tame extension \(L/K\) of degree \(N\), whether or not
it acts faithfully on either special fibre. Both actual groups
have good formal models over \(L\), and all their \(p\)-torsion
is \(L\)-rational. Their common second-level division field \(F/L\)
has elementary abelian group
\[
 W\subset M_2(\mathbf F_p),\qquad |W|=p^r,\quad r\in\{2,3,4\}.
 \tag{B-RD67.1}
\]
Here \(B\in W\) acts on the torsion module by \(1+pB\).
There are a tame complement \(\langle\sigma\rangle\) of order \(N\),
an exact eigen-uniformizer \(\varpi\) of \(F\), and integral
Honda-framed coordinates \(T_i\) such that
\[
 \begin{gathered}
 \sigma\varpi=\eta\varpi,\quad \eta^{p^r}=\zeta,\quad
 T_i(\sigma P)=\zeta^{c_i}\sigma(T_i(P)),\\
 A=v_F(T_i(P))=e p^{r-2}\quad(P\text{ primitive of order }p^2),\\
 v_F(T_i(Q))=qA\quad(0\ne Q\in M_i[p]),\qquad
 e_F=NqA .
 \end{gathered}
 \tag{B-RD67.2}
\]
The root \(\zeta\) is primitive of order \(N\), defined by the action
of \(\sigma\) on a tame uniformizer of \(L\). Put \(b_0=Ne\).
Primitive leading coefficients identify the two residual spaces
with \(d_i\mathbf F_q\subset k\). Their tame eigenvalue pairs are
\[
 \{\zeta^{c_i+e},\zeta^{p(c_i+e)}\}.
 \tag{B-RD67.3}
\]
If the special character pairs \(\{c_i,pc_i\}\) differ modulo \(N\),
the residual intertwiner in these coordinates is
\(z\mapsto a z^p\), with \(a\in\mathbf F_q^\times\).
\end{lemma}
\begin{proof}
Lemma~\ref{lem:B-RD49-inertia} gives actual good reduction over
a tame degree dividing \(N\). A normal \(p\)-subgroup of an
irreducible two-dimensional \(\mathbf F_p\)-representation acts
trivially: its nonzero fixed subspace is stable under the full
group. The remaining cyclic irreducible residual image lies in
\(\mathbf F_q^\times\), so \(L\) also contains the residual division
field. The second-level kernel is the additive matrix space.
After adjoining \(L\) there is no further prime-to-\(p\) quotient
in that kernel. Thus \(F/L\) is the indicated elementary abelian
extension.

The valuation filtration on \(M_i[p]\) is inertia stable.
Irreducibility makes all nonzero first-level coordinates have
one valuation. The product of the \(q-1\) nonzero torsion
coordinates, or equivalently the linear coefficient of the
Weierstrass polynomial of \([p]\), gives that valuation
\(e_L/N=e\). The one-segment Newton polygon has endpoints
\((1,eN),(q,0)\). Subtracting the coordinate of a nonzero
first-level point gives one inverse segment from \((0,e)\) to
\((q,0)\): all intermediate coefficients are strictly above it.
Every primitive second-level coordinate therefore has
\(v_L\)-valuation \(e/q\). Its integral valuation in \(F\) is
\(e p^{r-2}\); since \(p\nmid e\), integrality forces \(r\ge2\).
This proves all valuations in (B-RD67.2).

The extension of \(C_N\) by the additive group \(W\) splits.
Indeed, if a lift \(s\) of a generator has \(s^N=w\), then \(w\)
is fixed by \(s\); replacing \(s\) by \(s(-N^{-1}w)\) makes its
\(N\)-th power one. Averaging a uniformizer against its leading
character gives the exact eigen-uniformizer. Its eigenvalue
satisfies \(\eta^{p^r}=\zeta\), hence is primitive of order \(N\).

Tate full faithfulness extends the actual semilinear descent maps
and their inverses to the good models
\cite[Theorem 4 and Corollaries 1--2, p.~181]{Tate1967PDivisible}.
The finite special action can be put into a standard Honda
coordinate integrally. Explicitly, for a special unit \(v\) of
order \(m\mid N\) with residue equal to a standard root \(\lambda\),
\[
 x=\sum_{j=0}^{m-1}v^j\lambda^{-j}
 \quad\text{satisfies}\quad vx=x\lambda,\qquad
 x\equiv m\pmod\Pi .
\]
Thus \(x\) is a unit in the height-two maximal order, not merely
a rational conjugator. Averaging an integral lift of the resulting
Honda parameter gives the displayed exact semilinear coordinate;
its special \(p\)-series remains \(T^q\).

Write \(R_i=M_i[p^2](F)\), \(V_i=R_i/pR_i\). The leading
coefficient of \(T_i(P)/\varpi^A\) gives an injective
\(\mathbf F_p\)-linear map \(V_i\to k\): two primitive points
have the same leading coefficient exactly when their difference
is first-level. Multiplication by \(p\) raises this coefficient
to its \(q\)-th power, since the degree-\(q\) term of \([p]\)
is uniquely lowest at a primitive point. The first-level leading
space, acted on irreducibly by tame inertia, is a scalar multiple
of \(\mathbf F_q\), and so is the primitive leading space.
The uniformizer contribution is
\(\eta^A=\zeta^{p^{-r}e p^{r-2}}=\zeta^e\).
This proves (B-RD67.3). Its eigenvalues are not in \(\mathbf F_p\).
An intertwiner between the two multiplication actions is either
\(\mathbf F_q\)-linear or Frobenius-semilinear. The linear case
forces \(c_1=c_2\pmod N\), giving equal special characters.
Thus different special characters require the latter case.
Under multiplication by \(p\), normalized leading coordinates
are \(z\mapsto z^q=z\); the same intertwiner acts on the
first-level submodule. This fixes the orientation used below.
\end{proof}

\begin{lemma}[Exact coefficient calculus and the wild images of dimensions two and three]
\label{lem:B-RD67-small-wild}
\status{Proved; full independent root and companion proof, original-source and exact-code audits passed.}
Under the preceding hypotheses, a different-special-character
comparison cannot have \(r=2\) or \(r=3\).
\end{lemma}
\begin{proof}
Choose a coefficient ring in \(\mathcal O_K\), and representatives
of every residue class in that ring, by
\cite[Theorem 10.160.8 and proof]{Stacks032A}.
They are fixed by \(W\). Since \(e_F=NqA>qA\), all calculations
modulo \(\varpi^{qA+1}\) with these representatives are
characteristic-\(p\) calculations. For \(g\ne1\) define its lower
break and leading coefficient by
\[
 v_F(g\varpi-\varpi)=b_g+1,\qquad
 d_g=\overline{(g\varpi-\varpi)/\varpi^{b_g+1}} .
 \tag{B-RD67.4}
\]
A monomial of degree \(j\), with \(p^s\Vert j\), first changes
in degree \(j+p^s b_g\), with coefficient
\((j/p^s)d_g^{p^s}\) modulo \(p\). Terms divisible by \(p\)
in the mixed-characteristic expansion occur beyond \(qA\).
On a subgroup having one common break, the leading map
\(g\mapsto d_g\) is additive and injective. Tame conjugation
acts on it by \(\eta^{b_g}\).

If \(r=2\), then \(A=e\) is prime to \(p\). Every primitive
coordinate changes first in degree \(A+b_g\), with nonzero
coefficient \(\bar e\,d_i z d_g\). Thus every \(g\ne1\)
displaces every primitive point, and \(A+b_g=qA\). The matrix
\(B_i(g)\) on the residual space consequently acts by
\[
 B_i(g)z=\kappa_i d_g z,\qquad \kappa_i\ne0.
 \tag{B-RD67.5}
\]
Its image as \(g\) varies is the entire Cartan field
\(\mathbf F_q\), since \(W\) has order \(q\).
The different-character Frobenius intertwiner instead requires
the scalar in model 2 to be the \(p\)-th power of that in
model 1. Two fixed scalar multiples of the same \(q\)-element
space cannot have this property: it would make Frobenius scalar
on all of \(\mathbf F_q\). This excludes \(r=2\).

Suppose \(r=3\), so \(A=pe\). Identify the residual space with
\(\mathbf F_q\) and decompose
\[
 \operatorname{End}_{\mathbf F_p}(\mathbf F_q)=Z\oplus U,\qquad
 Z=\{z\mapsto az\},\quad U=\{z\mapsto az^p\}.
\]
If the residual projective order is \(h>2\), tame conjugation
is trivial on \(Z\) and irreducible on \(U\), with multiplier
\(\lambda^{1-p}\) of order \(h\mid p+1\). Thus
\(W=Z_1\oplus U\) for a Cartan line \(Z_1\), and \(U\) has
one lower break. If \(h=2\), conjugation is \(+1\) on \(Z\)
and \(-1\) on \(U\), so either \(W=Z_1\oplus U\) or
\(W=Z\oplus U_1\). The projective order cannot be one,
since the residual representation is irreducible.

For every primitive class \(z\), use its fixed expansion
\[
 T_i(P)=d_i z\varpi^A+
       \sum_{A<j<qA}f_{i,j}(z)\varpi^j\pmod{\varpi^{qA}}.
 \tag{B-RD67.6}
\]
The coefficients below \(qA\) are independent of the representative,
because changing it adds a first-level point. No additivity
of the other coefficient functions is assumed. Every nonzero
matrix in \(Z\) or \(U\) is invertible, so its displacement
has valuation exactly \(qA\) for every \(z\ne0\).
For an element of break \(b<b_0=Ne\), the leading term first
changes in degree \(A+pb<qA\). Every other degree divisible
by \(p\) changes later. Cancellation forces the first degree
prime to \(p\) to be
\[
 j=A+(p-1)b,\qquad
 j f_{i,j}(z)=-\bar e\,d_i d_g^{p-1}z .
 \tag{B-RD67.7}
\]
In particular \(p\nmid j\). For \(b>b_0\), all degrees divisible
by \(p\) change after \(qA\), so that first prime-to-\(p\)
degree is \(j=qA-b\). Its coefficient is a nonzero multiple
of \(z\) for a Cartan element, and of \(z^p\) for an anti-Cartan
element. These statements refer to the same expansion (B-RD67.6).

First suppose \(U\subset W\) has one break \(b\). Its tame action
is nontrivial, so \(\eta^b\ne1\), whence \(b\ne b_0\).
If \(b<b_0\), fix \(z\ne0\) in (B-RD67.7). It requires
\(d_g^{p-1}\) to be constant on the \(q-1\) nonzero elements
of the leading coefficient space of \(U\), impossible.
If \(b>b_0\), the displacement equation at \(j=qA-b\) is
\[
 d_i^q\beta_i(g)z^p=j f_{i,j}(z)d_g,\qquad g\in U,
\]
where \(\beta_i:U\to\mathbf F_q\) is an additive isomorphism.
Fixing one \(g\ne0\), then varying \(g\), gives
\(\beta_i(g)=\kappa_i d_g\). The residual intertwiner requires
\(\beta_2(g)=a^{1-p}\beta_1(g)^p\), again making Frobenius
scalar on all of \(\mathbf F_q\). This excludes a single break.

If \(U\subset W\) has two breaks \(b_1<b_2\), necessarily \(h=2\);
both satisfy \(\eta^{b_j}=-1\), so neither equals \(b_0\).
If \(b_1<b_0\), (B-RD67.7) gives the first prime-to-\(p\)
term with coefficient proportional to \(z\). At \(b_2\) its
change is uniquely earlier than all terms divisible by \(p\).
It must therefore be the required first-level displacement,
but would be Cartan-linear in \(z\), contrary to its nonzero
anti-Cartan type. If \(b_1>b_0\), that first term and all
other terms change after \(qA\) at \(b_2\), also impossible.

It remains that \(h=2\) and \(W=Z\oplus U_1\).
Write \(b\) for the break of \(U_1\). Its tame weight is \(-1\),
while every Cartan break \(c\) has \(\eta^c=1\); hence
\(b\ne b_0\) and \(c\ne b\).
If \(b>b_0\), the first ordinary coefficient is proportional
to \(z^p\). A Cartan break \(c<b_0\) would make that same
coefficient proportional to \(z\). If \(b_0\le c<b\),
its change is uniquely before \(qA\); if \(c>b\), all changes
are after \(qA\). All possibilities are excluded, including
the boundary \(c=b_0\). Thus \(b<b_0\).
The first ordinary degree is now \(j=A+(p-1)b\), with coefficient
proportional to \(z\). A Cartan break below \(b\) would force
an earlier first ordinary term. Every Cartan break is therefore
greater than \(b\), and its unique first change forces
\[
 c=qA-j=p b_0-(p-1)b>b_0 .
 \tag{B-RD67.8}
\]
All nonzero Cartan elements have this same break. Their leading
displacement makes their scalar coefficients
\(\beta_i(g)=\kappa_i d_g\) on the full plane \(Z\).
Frobenius-semilinearity requires
\(\beta_2(g)=\beta_1(g)^p\), the same contradiction as before.
This exhausts every three-dimensional \(W\).
\end{proof}

\begin{lemma}[The final plane for the full wild image]
\label{lem:B-RD67-final-plane}
\status{Proved; full independent root and companion proof, original-source and exact-code audits passed.}
If \(W=M_2(\mathbf F_p)\), then every primitive point field
\(L(P)/L\) is Galois of degree \(q\), with one lower and upper
break \(b_0=Ne\). The last nonzero upper group of \(F/L\) is
the full Cartan plane \(Z\). Its common lower break \(c\) satisfies
\(c>b_0\), and every anti-Cartan break is smaller than \(c\).
\end{lemma}
\begin{proof}
For a primitive residual vector \(v\ne0\), its stabilizer in
\(W\) is \(H_v=\{B:Bv=0\}\), a two-dimensional subspace.
Thus \(E=L(P)\) is Galois of degree \(q\). Its primitive
coordinate has valuation \(e\), prime to \(p\), and every
nonidentity automorphism displaces it by a first-level point
of valuation \(qe\). Fixed-coefficient uniformizer expansion
therefore gives the common lower break \(qe-e=b_0\).
A single lower break is also the single upper break.

Apply upper-numbering quotient compatibility
\cite[\S3.1, Proposition 3, p.~128]{Serre1961LocalClosed};
the full norm-polynomial and Herbrand proofs are also in
\cite[Proposition 3.7, Theorem 3.10 and Proposition 3.13,
pp.~14--16]{ColmezCorpsLocaux}.
The group \(S=W^{b_0}\) surjects under evaluation
\(B\mapsto Bv\) for every \(v\ne0\). There is no nonzero upper
group after \(b_0\), since \(\bigcap_v H_v=0\).

Over the good field \(L\), the determinant is cyclotomic by
the height-one determinant argument already proved in B18.
Its degree-\(p\) quotient is \(E_{\rm cyc}=L(\mu_{p^2})\).
This has degree \(p\), since \(\mu_p\subset L\) and \(p\nmid e_L\).
For \(x=\zeta_{p^2}-1\),
\[
 v_{E_{\rm cyc}}(x)=\frac{eN}{p-1},\qquad
 v_{E_{\rm cyc}}(gx-x)=\frac{peN}{p-1}\quad(g\ne1).
 \tag{B-RD67.9}
\]
The first integer is prime to \(p\), so this quotient also has
single break \(Ne=b_0\). Consequently trace is nonzero on \(S\).

At the last lower break \(c\), the leading ramification map
\(S\hookrightarrow k\) is additive and injective, with tame
weight \(\eta^c\). The nonzero invariant trace functional,
by averaging, supplies a nonzero tame-fixed vector of \(S\).
Hence \(\eta^c=1\), and injectivity makes all of \(S\) tame-fixed.
The centralizer of the irreducible residual tame element is
exactly \(Z\). Evaluation surjectivity forces dimension at least
two, so \(S=Z\). The Herbrand function has value \(b_0\) at
this last lower break. Since \(Z\) is a proper subgroup of \(W\),
its slope is strictly less than one on a nonempty interval before
\(c\), giving \(c>b_0\). This proves all assertions.
\end{proof}

\begin{lemma}[All two-dimensional anti-Cartan break patterns]
\label{lem:B-RD67-full-wild}
\status{Proved; full independent root and companion proof, original-source and exact-code audits passed.}
The full wild image \(W=M_2(\mathbf F_p)\) also cannot occur
in a different-special-character comparison.
\end{lemma}
\begin{proof}
Now \(A=qe\), with \(v_p(A)=2\), and the final Cartan plane
has the break \(c>b_0\) just proved. For a fixed primitive
coordinate \(z\ne0\), let \(J\) be the first nonzero degree
of exact \(p\)-valuation one, and \(K_0\) the first nonzero
degree prime to \(p\), in its expansion below \(qA\).
An absent degree is interpreted as infinity. All degrees
divisible by \(q\) change after \(qA\) at the break \(c\).
The \(q\)-element leading ramification space of \(Z\) cannot
be the root space of a nonzero polynomial of degree at most \(p\).
Applying this to a possible cancellation between the \(J\)
and \(K_0\) terms gives
\[
 J+pc\ge qA,\qquad K_0+c\ge qA,\qquad
 \min(J+pc,K_0+c)=qA .
 \tag{B-RD67.10}
\]
Every nonzero anti-Cartan matrix is invertible, so its displacement
also has exact valuation \(qA\).

First suppose the anti-Cartan plane has one break \(b<c\).
Its leading coefficient space \(D\) has dimension two over
\(\mathbf F_p\), but need not be a scalar multiple of \(\mathbf F_q\).
Its root polynomial is
\[
 P_D(X)=X^q+s_DX^p+t_DX,\qquad t_D\ne0.
 \tag{B-RD67.11}
\]
For example, with basis \(a,b\), put \(Y=X^p-a^{p-1}X\);
then \(P_D=Y^p-(b^p-a^{p-1}b)^{p-1}Y\).
The tame action on the anti-Cartan plane has no fixed vector,
so \(\eta^b\ne1\), and \(b\ne b_0\).
If \(b>b_0\), the leading degree \(A\) already changes after
\(qA\). An earlier cancellation between the \(J\) and \(K_0\)
terms would be a polynomial of degree at most \(p\) vanishing
on all of \(D\), impossible. Thus both of those changes are
at least \(qA\) at \(b\), contradicting (B-RD67.10) at \(c>b\).
Hence \(b<b_0\).

At this lower break, the degree-\(A\) change in degree \(A+qb\)
has nonzero coefficient \(\bar e\,d_i z d_g^q\).
No earlier change from the other two classes can cancel on
all of \(D\), by the same degree-\(p\) root bound.
Cancellation at \(A+qb\) must therefore be a multiple of
(B-RD67.11). Its nonzero linear coefficient forces
\[
 K_0=A+Nb,\qquad J\ge A+(q-p)b,
 \tag{B-RD67.12}
\]
and \(f_{i,K_0}(z)\) is one fixed nonzero multiple of \(z\).
The middle coefficient \(s_D\) is retained; equality in the
bound for \(J\) is allowed. At the later break \(c\),
\[
 J+pc-(K_0+c)\ge(p-1)(c-b)>0,\qquad
 A+qc-(K_0+c)=N(c-b)>0 .
 \tag{B-RD67.13}
\]
Thus the prime-to-\(p\) term is uniquely first and
\(K_0+c=qA\).

If the anti-Cartan plane has two breaks \(b_1<b_2<c\),
consider the three possible first-change degrees
\[
 A+qb,\qquad J+pb,\qquad K_0+b.
 \tag{B-RD67.14}
\]
At \(b_2\), (B-RD67.10) makes the minimum smaller than \(qA\);
it must be attained twice. It cannot involve the first term:
its larger slope would make it uniquely smallest at \(b_1\),
still below \(qA\). Hence
\[
 J+pb_2=K_0+b_2<A+qb_2.
 \tag{B-RD67.15}
\]
At \(b_1\), the middle term is strictly smaller than the last
and is below \(qA\). Cancellation forces equality with the first:
\[
 J=A+(q-p)b_1,\qquad
 K_0=J+(p-1)b_2,\qquad b_1<b_0.
 \tag{B-RD67.16}
\]
Choose fixed elements of the two breaks. The first cancellation
makes \(f_{i,J}(z)\) a fixed nonzero multiple of \(z\);
the second makes \(f_{i,K_0}(z)\) a fixed nonzero multiple
of that coefficient. These fixed elements work for every
\(z\ne0\); therefore the degrees and multiples are independent
of \(z\). At \(c\), the ordinary term is strictly first, since
\(J+pc-(K_0+c)=(p-1)(c-b_2)>0\); (B-RD67.10) again gives
\(K_0+c=qA\).

In either pattern, on the full final plane \(Z\) the displacement
equation is
\[
 d_i^q\beta_i(g)z=K_0 f_{i,K_0}(z)d_g .
 \tag{B-RD67.17}
\]
Thus \(\beta_i(g)=\kappa_i d_g\), and
\(\beta_i:Z\to\mathbf F_q\) is bijective.
The different-character intertwiner requires
\(\beta_2(g)=\beta_1(g)^p\), impossible for two fixed scalar
multiples of the same \(q\)-element space. This excludes \(r=4\).
\end{proof}

\begin{theorem}[Irreducible residual recovery for every prime-to-\(p\) index]
\label{thm:B-RD67}
\status{Proved; full independent root and companion proof, original-source and exact-code audits passed.}
Let \(p\ge5\), let \(K\) be complete discretely valued of mixed
characteristic \((0,p)\) with perfect residue field, \(\mu_p\subset K\)
and \(p\nmid e_K\). Let \(M_1,M_2\) be height-two BT groups
which acquire one-dimensional supersingular good reduction.
If their \(p^2\)-torsion is isomorphic and the common residual
inertia is irreducible over \(\mathbf F_p\), then their finite
special inertia characters are equal.
\end{theorem}
\begin{proof}
Pass to the completed maximal unramified extension. This preserves
the absolute ramification and the given torsion comparison, while
the residue field becomes algebraically closed. The preceding
lemmas exclude each possible wild image size \(p^2,p^3,p^4\)
under the assumption that the special characters differ.
They exhaust the possible sizes by (B-RD67.1). All calculations
refer to the actual common division field and the given torsion
isomorphism, so they prove equality on inertia over the original
field as well.
\end{proof}

\begin{remark}[Scope and the role of the ramification index]
\label{rem:B-RD67-scope}
\status{Proved; full independent root and companion proof, original-source and exact-code audits passed.}
The proof replaces the earlier minimal-index calculation by explicit
normalizations \(A=e p^{r-2}\), \(b_0=(p^2-1)e\), and
\(e_F=(p^2-1)p^2A\). The coefficient of the first changed leading
monomial is \(\bar e\), nonzero precisely because \(p\nmid e\).
No identification of lower and upper ramification numbers is made
before the point-field and cyclotomic quotient calculation.
The tame extension of degree \(p^2-1\) need not act faithfully on
the special groups. The theorem does not claim character recovery
for reducible residual inertia, ramification divisible by \(p\),
or general heights.
The retained exact supplement
\path{code/reduction_descent/verify_irreducible_arbitrary_index.py}
checks \(150\) varying prime-to-\(p\) indices, \(450\) field
normalizations, \(8{,}093\) direct binomial first terms, \(60\)
higher-divisibility terms, \(30{,}450\) leading-character cases,
\(38{,}689\) cubic Cartan comparisons, \(31{,}489\) single-break
and \(12{,}948\) two-break full-wild patterns. It also checks
all \(221\) two-planes in the cubic fields over
\(\mathbf F_5,\mathbf F_7,\mathbf F_{11}\), with
\(200{,}449\) exact root tests and \(806\) directly multiplied
polynomial coefficients; the middle coefficient is retained
and nonzero in these planes. The author run passed in
\(2.847146783999051\) seconds at 20:57:06 UTC.
The complete independent root and companion audits passed; their
author-output-preserving reruns took
\(2.895408538999618\) and \(2.5705544949887553\) seconds,
respectively. These are finite arithmetic checks of the proof,
not formal-group or pair constructions.
\end{remark}


% Proved after complete root and companion audits; B67 is proved and exact candidates are retained.
\subsection{The initial tame ramification sector in height two}

Fix \(p\ge5\) and
\[
 e=t(p-1),\qquad 2\le t\le p-1,\qquad
 q=p^2,\quad N=q-1,\quad s=e(p+1)=tN .
 \tag{B-RD68.1}
\]
In particular \(p\nmid e\), \(e<p^2\), and \(e\) is even.
After completed maximal unramified base change, the field \(K\)
is complete of mixed characteristic \((0,p)\), with algebraically
closed residue field, \(e_K=e\) and \(\mu_p\subset K\).
For whole potentially supersingular height-two groups we take
the common good tame extension \(L/K\) of degree \(N\), and normalize
\(v_L(\pi_L)=1\). Thus \(v_L(p)=eN=(p-1)s\).

\begin{lemma}[Unequal-Newton rules throughout the sector]
\label{lem:B-RD68-rules}
\status{Proved; full independent root and companion proof, source and exact-code audits passed.}
Suppose the common finite determinant character is nonsquare.
Make a common finite height-one twist so that every actual
special tangent exponent satisfies \(c\equiv1\pmod{p-1}\).
If the residual module is reducible, its Newton polygon has two
slopes, with unique high line \(H\) and integer \(B\) satisfying
\[
 \begin{gathered}
 0<B<pe,\qquad B\equiv c\pmod{p+1},\qquad B\ \text{odd},\\
 \alpha=s-B,\quad\beta=B/p,\quad\alpha>\beta,\\
 \text{high-line character } \zeta^{c-B}.
 \end{gathered}
 \tag{B-RD68.2}
\]
Every order-\(p^2\) point whose \(p\)-multiple is outside \(H\)
has \(v_L\)-valuation \(B/p^3\).
For two such groups with isomorphic actual \(p^2\)-torsion:

\begin{enumerate}
\item If their residual wild image is nontrivial, the high lines
agree. Two parameters prime to \(p\) are equal and give equal
special characters. Exactly one parameter cannot be divisible
by \(p\).
\item If the residual image is tame, write \(B_i=pb_i\).
When both \(b_i\) are prime to \(p\), they are equal, their
high lines agree, and the special characters agree.
\end{enumerate}
\end{lemma}
\begin{proof}
The universal height-two coefficient bounds established in
Lemma~\ref{lem:B-RD49-point} give the only possible negative
Newton vertices \((1,(p-1)s),(p,v_L(u)),(q,0)\).
If all nonzero first-level points have one valuation, their
leading-coefficient map is injective. Wild inertia acts trivially
on that map: its uniformizer multiplier and residue action are
both trivial, and the good descent is tame. Thus the residual
module is tame, with a scalar leading-coordinate action.
If it is reducible over \(\mathbf F_p\), that scalar is in
\(\mathbf F_p^\times\), and its determinant character is square.
Because \(\mu_p\subset K\), reduction of the cyclotomic determinant
is trivial, so this contradicts the nonsquare finite determinant.
The reducible case must therefore have two slopes.

The high line is stable under inertia, and its character has
values in \(\mathbf F_p^\times\). It is rational over \(L\).
Its valuation is integral, so writing
\(v_L(u)=(p-1)B\) gives the integer \(B\) and the slopes
in (B-RD68.2). Their strict order gives \(B<pe\).
Choose an exact tame coordinate and
\(\sigma\pi_L=\zeta\pi_L\), as in the integral averaging proof
of Lemma~\ref{lem:B-RD67-normalization}. On a nonzero high
point the leading character is \(\zeta^{c+\alpha}\).
Since it is \(\mathbf F_p\)-valued,
\(c+s-B\equiv0\pmod{p+1}\), giving the asserted congruence.
Here \(s=tN\), so its actual value is \(\zeta^{c-B}\).
As \(c\) is odd and \(p+1\) is even, \(B\) is odd.
The inverse Newton polygon at target valuation \(\beta\)
is one segment of degree \(q\), giving \(B/p^3\) for every
preimage.

Let \(R/L\) be the common residual division field. It is trivial
or cyclic of degree \(p\). In the latter case the high line
is the unique fixed line of the wild unipotent image, so it is
common. A low point has \(v_R\)-valuation \(B\), while its
nonzero wild displacement has valuation \(p\alpha\).
The fixed-coefficient lower-break calculation gives
\[
 r=ps-(p+1)B\quad(p\nmid B),\qquad
 r<ps-(p+1)B\quad(p\mid B).
 \tag{B-RD68.3}
\]
Indeed a coordinate of valuation \(B\) changes in valuation
\(B+r\) if \(p\nmid B\), and strictly later if \(p\mid B\);
every higher term changes later still. This is precisely the
proof of Lemma~\ref{lem:B-RD60-break}, with the present valuations.
Two prime-to-\(p\) parameters give the same \(r\), hence the
same \(B\); their common high-line character then gives the
same \(c\) modulo \(N\).

If \(p\nmid B_1\) and \(p\mid B_2\), (B-RD68.3) gives \(B_1>B_2\).
Choose corresponding primitive points \(P_i\) with \(pP_i\)
outside the common high line. Their \(v_R\)-valuations are
\(B_i/p^2\). Their point fields over \(R\) are equal, because
the given Galois-module isomorphism identifies their stabilizers.
The denominator on the first side forces their common point
field \(E\) to have degree \(p^2\): it is a torsor under
the \(R\)-rational first-level module, so degree is at most \(p^2\).
Thus every first-level translation occurs.
In \(E\), the two point coordinates have valuations \(B_i\).
A translation outside the common high line has displacement
valuation \(p^2B_i\). The first coordinate gives actual break
\(NB_1\); the universal lower bound on the second gives
\[
 p^2B_2\ge B_2+NB_1,
 \tag{B-RD68.4}
\]
contradicting \(B_1>B_2\). This proves the first rule.

In the tame case the residual module is rational over \(L\),
so \(B_i=pb_i\). Choose \(pP_i\) outside both high lines.
If both \(b_i\) are prime to \(p\), the common point field has
degree \(p^2\), and the point coordinates there have valuations
\(b_i\). A translation outside both high lines gives the same
break \(Nb_1=Nb_2\), hence \(b_1=b_2=b\).
If the high lines differed, a translation high only in one
model would have unequal displacement valuations, although
both coordinates have the same prime-to-\(p\) valuation \(b\)
and must change in valuation \(b+r\). This is impossible.
The high-line characters now identify \(c\) modulo \(N\).
\end{proof}

\begin{lemma}[The actual cyclic quotient transformations]
\label{lem:B-RD68-quotients}
\status{Proved; full independent root and companion proof, source and exact-code audits passed.}
The actual quotient by the high line has Newton parameter
\[
 B_H=\begin{cases}
       pB,&B<e,\\
       s-B,&B>e.
      \end{cases}
 \tag{B-RD68.5}
\]
There is no equality case, because \(B\) is odd and \(e\) even.
If the residual module is rational over \(L\) and \(J\) is an
inertia-stable low line, its actual quotient has parameter
\[
 B_J=B/p .
 \tag{B-RD68.6}
\]
Each quotient is good over \(\mathcal O_L\), descends generically
to \(K\), and has special tangent exponent \(pc\) modulo \(N\).
Its finite special character is unchanged. A \(p^3\)-torsion
isomorphism preserving the chosen lines gives a \(p^2\)-torsion
isomorphism of these quotients.
\end{lemma}
\begin{proof}
The finite flat closure of a chosen rank-\(p\) line has an
actual good formal quotient. Its coordinate map is a unit times
\(\prod_{Q\in J}(T-T(Q))\). Valuations are therefore computed
by this product, without assuming that a rational isogeny
preserves a chosen parameter thickness.

For \(J=H\), every point below valuation \(\alpha\) has its
valuation multiplied by \(p\). The image of \(M[p]/H\) has
valuation \(B\). For a nonzero high target, if \(B>e\) the
inverse polygon is one segment with \(q\) roots of valuation
\(\alpha/q\); the resulting \(q-p\) nonzero quotient points
have valuation \(\alpha/p\). The other \(p-1\) points have
valuation \(B>\alpha/p\), giving \(B_H=s-B\).
If \(B<e\), each nonzero high target has \(p\) inverse roots
of valuation \((s-pB)/p\) and \(q-p\) roots of valuation \(B/p\).
All lie below \(\alpha\). After quotienting the root counts
by \(p\), and including \(M[p]/H\), there are \(p-1\)
points of valuation \(s-pB\) and \(q-p\) of valuation \(B\).
Thus \(B_H=pB\).

For a low line \(J\), every nonzero class in \(M[p]/J\)
has one high representative. Its quotient valuation is
\(\alpha+(p-1)\beta=s-B/p\). All preimages of nonzero low
targets have valuation \(\beta/q<\beta\), and their quotient
valuations are \(\beta/p=B/p^2\). Their count is \(q-p\).
This gives (B-RD68.6); \(B/p\) is integral because the
first-level module is \(L\)-rational.

On the special fibre the unique subgroup of rank \(p\) is
the Frobenius kernel. Thus the special quotient is the
Frobenius twist and has inherited tangent exponent \(pc\).
The unordered pair \(\{c,pc\}\) is unchanged.
The chosen line is stable under the actual generic descent,
so the generic quotient is defined over \(K\).
Finally its \(p^2\)-torsion is exactly
\[
 \{P\in M[p^3]:p^2P\in J\}/J .
 \tag{B-RD68.7}
\]
This identity of actual finite generic groups proves the
claimed passage between the two torsion levels.
\end{proof}

\begin{proposition}[Nonsquare formal comparisons in the initial sector]
\label{prop:B-RD68-nonsquare}
\status{Proved; full independent root and companion proof, source and exact-code audits passed.}
Two whole potentially supersingular height-two groups in this
sector with common nonsquare finite determinant and isomorphic
actual \(p^3\)-torsion have equal finite special inertia characters.
\end{proposition}
\begin{proof}
Normalize the common finite determinant by a common finite
height-one twist. An irreducible residual comparison is already
settled at \(p^2\) by Theorem~\ref{thm:B-RD67}.
The other comparisons have the unequal parameters just proved.

In the residual wild case, their high line is common.
The rules handle any case with a prime-to-\(p\) parameter.
Suppose both \(B_i\) are divisible by \(p\).
If both are below \(e\), write \(B_i=pb_i\).
Their common high-line character and \(c_i\equiv1\pmod{p-1}\)
give \(b_1\equiv b_2\pmod{p-1}\). But
\[
 0<b_i<e/p\le (p-1)^2/p<p-1,
 \tag{B-RD68.8}
\]
so \(b_1=b_2\), and the high-line character identifies \(c_1=c_2\).
If both \(B_i>e\), quotient by the common high line.
The new parameters \(s-B_i\) are prime to \(p\), since \(p\nmid s\).
Their actual \(p^2\)-torsion is isomorphic, so the rules identify
the quotient special characters, and therefore the original ones.
If one parameter is below \(e\) and one above, the two quotient
parameters \(pB_i\) and \(s-B_j\) have opposite \(p\)-divisibility.
They cannot have isomorphic \(p^2\)-torsion by the same rules.

Now suppose residual inertia is tame, and write \(B_i=pb_i\),
where \(0<b_i<e\). Two \(b_i\) prime to \(p\) are already
handled. Suppose \(p\nmid b_1=b\) and \(b_2=pk\).
Then \(B_2=p^2k>e\), since \(e<(p^2)\).
Choose corresponding primitive points whose \(p\)-multiples
are outside both high lines. The common point field has degree
\(p^2\) over \(L\), forced by the first coordinate valuation
\(b/p^2\); its normalized coordinate valuations are \(b\) and \(pk\).
If the high lines differ, choose a translation high in model 1
and low in model 2. The first coordinate, prime to \(p\), fixes
its actual lower break as \(p^2\alpha_1-b\).
The universal lower bound in model 2 would require
\[
 p^2(\alpha_1-pk)\le b-pk .
 \tag{B-RD68.9}
\]
Yet \(\alpha_1=s-pb\ge e+p\), while \(pk<e\), \(b<e<p^2\).
The left side is greater than \(p^3\), and the right side less
than \(p^2\), a contradiction. Thus the high lines agree.
Quotient by this common line. The second parameter \(s-B_2\)
is prime to \(p\). If \(B_1<e\), the first parameter \(pB_1\)
is divisible, excluded by the rules. If \(B_1>e\), both new
parameters are prime to \(p\), so their equality forces
\(B_1=B_2\), contradicting \(p\nmid b\).

Finally let \(b_i=pk_i\) for both groups. Both
\(B_i=p^2k_i>e\). If the high lines agree, their high quotients
have prime-to-\(p\) parameters and equal characters by the rules.
If the high lines differ, quotient the first group by its
high line and the second by the corresponding stable low line.
The parameters are \(s-B_1\), prime to \(p\), and
\(B_2/p=pk_2\), divisible by \(p\). Their inherited actual
\(p^2\)-torsion isomorphism is impossible. These are all
residual and divisibility cases.
\end{proof}

\begin{lemma}[Square determinants in the same sector]
\label{lem:B-RD68-square}
\status{Proved; full independent root and companion proof, source and exact-code audits passed.}
The same conclusion holds when the common finite determinant
character is square.
\end{lemma}
\begin{proof}
Twist both groups by its common finite square root.
Their special inertia eigenvalues have norm one, so their
simultaneous good field has degree \(m\mid p+1\).
Put \(a=e m\le e(p+1)\).
If the special characters differ, their finite special inertia
eigenvalue orbits are disjoint. The integral Tate pairings and self-orthogonal
graph argument of Proposition~\ref{prop:B-RD64-formal} give
the special graph closure
\(\ker F_M^3\times\ker F_N^3\).
In detail the graph is isotropic generically, hence integrally
by flatness; equality with its orthogonal follows from equal
finite flat ranks. Integral tame polynomial projectors split
its special fibre into the two formal factors, and uniqueness
of rank-\(p^a\) subgroups of a one-dimensional connected formal
group gives exponent three on each factor.
The actual quotient \(Q=(M\times N)/H\) consequently has
special fibre \(M_0^{(p^3)}\times N_0^{(p^3)}\).
The map
\[
 M\times N\longrightarrow Q\times M,\qquad
 (x,y)\longmapsto(q(x,y),x)
\]
has no nonzero finite flat subgroup in its kernel, while its
special kernel is \(0\times\ker F_N^3\).

The actual target is a product of height-two formal groups
on the special fibre. For a coefficientwise valuation-one
vector its first two \(p\)-multiplications have valuations
\[
 q_1\ge p+1,\qquad
 q_2\ge\min\{a+p+1,p^2+p+1\}\ge a/(p-1).
 \tag{B-RD68.10}
\]
Here \(a\ge e\ge2(p-1)\ge p\).
The differential identity for the \(p\)-series makes any
coefficient with a non-\(p\)-divisible exponent divisible by \(p\);
all total degrees below \(p^2\) vanish specially, so the
remaining such coefficients are divisible by a uniformizer.
These are precisely the intrinsic coefficient estimates on
the actual target used in B34 and B64.
Finally
\[
 a/(p-1)\le e(p+1)/(p-1)\le(p-1)(p+1)=N .
 \tag{B-RD68.11}
\]
The closed coefficientwise logarithm and Cartier bounded-projection
and saturation proof in Lemma~\ref{lem:B-RD64-weighted-kernel},
using the just stated unweighted bounds, forces \(F^2\) to
kill this special kernel. It cannot kill \(\ker F_N^3\).
This contradiction uses the actual quotient target, without
transferring any Hasse-coefficient thickness from \(M\) or \(N\).
\end{proof}

\begin{theorem}[The exact initial tame sector]
\label{thm:B-RD68}
\status{Proved; full independent root and companion proof, source and exact-code audits passed.}
For every prime \(p\ge5\), in the full B.4 monodromy setting
with height two, perfect residue field and \(\mu_p\subset K\),
\[
 \boxed{\displaystyle
 \nu_{\min}\bigl(p,t(p-1),2\bigr)=2
 \quad\text{for every }2\le t\le p-1.}
 \tag{B-RD68.12}
\]
Thus actual \(p^3\)-torsion always determines the requested
character in this sector, whereas actual \(p^2\)-torsion
does not do so uniformly. In particular
\(\nu_{\min}(5,16,2)=2\).
\end{theorem}
\begin{proof}
For the lower bound, \(e=t(p-1)>p\) and \(p-1\mid e\).
The explicit actual \(p^2\)-torsion comparison in
Theorem~\ref{thm:B-RD63-arc} therefore applies at every
one of these indices.

For the upper bound, pass to the completed maximal unramified
base and compare actual \(p^3\)-torsion. The determinant of
an allowed height-two monodromy system is
\(\chi_p^d\epsilon\), \(d\in\{0,1,2\}\), with \(\epsilon\)
finite and \(\mu_{p-1}\)-valued. Since \(p\nmid e\), the
cyclotomic image modulo \(p^3\) has order \(p^2\).
Equality of determinants gives the same \(d\), because
\(|d_1-d_2|\le2<p\), and then the same \(\epsilon\).
For \(d=0\), the special character is the Brauer lift of the
semisimplified residual inertia character; \(d=2\) follows
by Cartier duality. For \(d=1\), each group is either whole
potentially supersingular or filtered by two potentially good
height-one groups. Two filtered systems have their finite
characters determined modulo \(p\), since \(\mu_p\subset K\).
A mixed comparison is impossible by
Proposition~\ref{prop:B-RD63-mixed}, because \(e<p^2<p^3\).
The two whole formal groups have the common finite determinant
just established. Propositions~\ref{prop:B-RD68-nonsquare}
and Lemma~\ref{lem:B-RD68-square} settle its two square classes.
This proves the upper bound, and hence the exact index.
\end{proof}

\begin{remark}[Boundary of the sector]
\label{rem:B-RD68-scope}
\status{Proved; full independent root and companion proof, source and exact-code audits passed.}
The proof uses \(e<p^2\) and \(e/p<p-1\) in the reducible
residual comparisons, and the actual-target estimate
\(e(p+1)/(p-1)\le p^2-1\) for square determinants.
It does not extrapolate those inequalities to \(t\ge p\)
or ramification divisible by \(p\).
The minimal index \(t=1\) has the separately proved value one
from B61. At \(p=5\), the new sector consists of \(e=8,12,16\);
the first two were already proved in B64 and B65.
The theorem concerns the full filtered monodromy condition,
not only comparisons of two whole potentially good groups.
The retained exact supplement
\path{code/reduction_descent/verify_initial_tame_sector.py}
checks all \(60\) sectors \(2\le t\le p-1\) at
\(p=5,7,11,13,17,19\). It independently computes
\(56{,}916\) Newton/high-quotient records, \(3{,}420\) low
quotients, \(113{,}832\) tangent weights, \(63{,}018\)
mixed-divisibility point-field inequalities and \(322\)
actual-target logarithm bounds, together with every stated
small-parameter and two-divisible congruence in those ranges.
The author run passed in \(2.3066381700045895\) seconds at
20:57:28 UTC. The complete independent root and companion audits passed, with
author-output-preserving reruns in \(1.9408804569975473\)
and \(2.1150716640113387\) seconds, respectively.
The all-prime assertion is the preceding proof and does not follow
from this finite enumeration.
\end{remark}


% Proved after full root and companion audits; exact candidate and outputs retained.
\subsection{One-slope quotients and the first tame boundary}

Throughout this subsection \(p\ge5\), \(q=p^2\), \(N=q-1\), and
\(K\) is complete discretely valued of mixed characteristic \((0,p)\),
with algebraically closed residue field, \(\mu_p\subset K\), and
absolute ramification \(e\) prime to \(p\). Write
\[
 p-1\mid e,\qquad s=e(p+1),\qquad
 L/K\text{ tame of degree }N,\qquad v_L(p)=eN=(p-1)s .
 \tag{B-RD69.1}
\]
The extension of degree \(N\) is used even when the special action
is not faithful. By Lemma~\ref{lem:B-RD49-inertia}, it is a good
field for every whole potentially supersingular height-two group.
Fix \(\sigma\pi_L=\zeta\pi_L\), with \(\zeta\) primitive of order \(N\).
An integral Honda framing and tame averaging, as in
Lemma~\ref{lem:B-RD67-normalization}, give an exact semilinear
coordinate whose special tangent multiplier is \(\zeta^c\).
The finite special character has eigenvalue exponents \(\{c,pc\}\).
No square-class normalization of its determinant is made here.

\begin{lemma}[Comparison rules conditional on unequal slopes]
\label{lem:B-RD69-rules}
\status{Proved; complete independent root and companion proof, source and exact-code audits passed.}
If the first-level Newton polygon has two slopes, it has a
unique high line \(H\) and an integer \(B\) with
\[
 \begin{gathered}
 0<B<pe,\qquad B\equiv c\pmod{p+1},\\
 \alpha=s-B,\qquad \beta=B/p,\qquad \alpha>\beta,\\
 \text{high-line character }\zeta^{c-B}.
 \end{gathered}
 \tag{B-RD69.2}
\]
There are \(p-1\) nonzero points of valuation \(\alpha\) and
\(q-p\) of valuation \(\beta\). Every order-\(p^2\) point whose
\(p\)-multiple is outside \(H\) has \(v_L\)-valuation \(B/p^3\).
For two such groups with isomorphic actual \(p^2\)-torsion:

\begin{enumerate}
\item If residual wild inertia is nontrivial, the high lines
agree. Two parameters prime to \(p\) are equal and give equal
special characters; exactly one parameter cannot be divisible
by \(p\).
\item If residual inertia is tame, the first-level modules are
rational over \(L\), and \(B_i=pb_i\). If both \(b_i\) are prime
to \(p\), they are equal, the high lines agree, and the special
characters agree.
\end{enumerate}
These statements impose no parity condition on \(B\) or \(c\).
\end{lemma}
\begin{proof}
The exact universal coefficient bounds and their original source
are given in Lemma~\ref{lem:B-RD49-point}:
\cite[Proposition~5.7; equations (12.3)--(12.5);
Proposition~12.10 and proof]{HopkinsGross1994Bundles}.
The only possible negative Newton vertices are
\((1,(p-1)s),(p,b),(q,0)\), where \(b=v_L(u)\).
The high line is inertia stable. Its character is
\(\mathbf F_p^\times\)-valued and hence trivial over \(L\);
its coordinates have integral valuation. Therefore
\(\alpha=s-b/(p-1)\) is integral, so \(b=(p-1)B\).
The strict slope inequality is exactly \(0<B<pe\).
Since \(s\) is a multiple of \(N\), the leading high-line
character is \(\zeta^{c+\alpha}=\zeta^{c-B}\).
It is \(\mathbf F_p\)-valued, giving the displayed congruence.
At a low target of valuation \(\beta\), the inverse Newton polygon
has one segment of degree \(q\), giving \(B/p^3\) for every
preimage.

For clarity, the common-point-field argument is rechecked without
the nonsquare hypothesis of Lemma~\ref{lem:B-RD68-rules}.
Let \(R/L\) be the common first-level division field. The high
line and the determinant are trivial over \(L\), so the residual
image there is a subgroup of the upper unipotent group of order
\(p\). Thus \(R=L\) or \([R:L]=p\).
In the latter case the unique wild fixed line is \(H\), so the
two high lines agree. A low coordinate has \(v_R\)-valuation
\(B\); its nonzero wild displacement has valuation \(p\alpha\).
If \(r\) is the actual lower break, the fixed-coefficient
uniformizer expansion of Lemma~\ref{lem:B-RD60-break} gives
\[
 r=ps-(p+1)B \quad(p\nmid B),\qquad
 r<ps-(p+1)B \quad(p\mid B).
 \tag{B-RD69.3}
\]
Explicitly, choose residue representatives in a coefficient
ring fixed by the automorphism. For a coordinate of valuation
\(B\), its first change has valuation \(B+r\) if \(p\nmid B\);
if \(p\mid B\), the first monomial changes strictly later,
as do all higher monomials. This proves the strict version too.
Two prime-to-\(p\) parameters give the same \(B\). Their common
high-line character then identifies their \(c\)'s modulo \(N\).

If \(p\nmid B_1\) and \(p\mid B_2\), (B-RD69.3) gives
\(B_1>B_2\). Choose corresponding primitive \(P_i\) whose
\(p\)-multiples are outside the common high line. Their
\(v_R\)-valuations are \(B_i/p^2\). Their stabilizers under
the actual torsion isomorphism agree, so their point fields
\(E/R\) agree. Such a point field is a torsor under the
\(R\)-rational first-level module, of degree at most \(p^2\).
The first valuation denominator forces degree \(p^2\).
Every first-level translation therefore occurs, and the
normalized valuations in \(E\) are \(B_i\).
A translation outside \(H\) has displacement valuation \(p^2B_i\).
The first coordinate gives its exact break \(NB_1\), while
the universal lower bound for the other gives
\[
 p^2B_2\ge B_2+NB_1,
 \tag{B-RD69.4}
\]
contradicting \(B_1>B_2\).

If residual inertia is tame, the residual division field over
\(L\) is trivial, hence \(B_i=pb_i\). Choose a common residual
vector outside the two high lines; such a vector exists because
two lines do not cover \(\mathbf F_p^2\).
When both \(b_i\) are prime to \(p\), the corresponding primitive
points have valuations \(b_i/p^2\), and their common point field
over \(L\) has degree \(p^2\). A translation outside both high
lines fixes the actual break as \(Nb_1=Nb_2\).
Thus \(b_1=b_2=b\). If the high lines differed, a translation
high in one model and low in the other would have different
displacement valuations. Both coordinates have the same
prime-to-\(p\) valuation \(b\), however, so each changes in
valuation \(b+r\) for that automorphism. This contradiction
identifies the high lines; their characters identify \(c\).
\end{proof}

\begin{lemma}[Actual quotients, including the equal-slope case]
\label{lem:B-RD69-quotients}
\status{Proved; complete independent root and companion proof, source and exact-code audits passed.}
The quotient of a two-slope group by its high line has parameter
\[
 B_H=pB\quad(B<e),\qquad B_H=s-B\quad(B>e).
 \tag{B-RD69.5}
\]
At \(B=e\), its quotient instead has one slope: all its nonzero
\(p\)-torsion points have valuation \(e\).
If the first-level module is \(L\)-rational and \(J\) is an
inertia-stable low line, its quotient has parameter \(B_J=B/p\).
If the original group has one slope, all its nonzero
\(p\)-torsion points have valuation \(e\), and its quotient by
any inertia-stable rank-\(p\) line has two slopes
\[
 \alpha'=pe,\qquad\beta'=e/p,\qquad B'=e.
 \tag{B-RD69.6}
\]
All these are actual good formal quotients over \(\mathcal O_L\).
Their generic groups descend to \(K\), their special tangent
exponents are \(pc\), and their finite special characters are
unchanged. An actual \(p^3\)-torsion isomorphism preserving the
chosen lines induces an actual \(p^2\)-torsion isomorphism of
the quotients.
\end{lemma}
\begin{proof}
Take the finite flat closure of the chosen line in the good
model. Its quotient is a good one-dimensional formal group;
the quotient coordinate map has the form
\[
 T\longmapsto U(T)\prod_{Q\in J}(T-T(Q)),\qquad U(T)\in
 \mathcal O_L[[T]]^\times .
 \tag{B-RD69.7}
\]
This is Weierstrass preparation for the actual finite-flat
isogeny. In particular all asserted valuations are sums of
actual difference valuations.

For the unequal polygon, quotienting by the high line multiplies
every valuation below \(\alpha\) by \(p\).
The image of \(M[p]/H\) has valuation \(B\).
If \(B>e\), the inverse polygon over a nonzero high point has
\(q\) roots of valuation \(\alpha/q\); their quotient images
give \(q-p\) points of valuation \(\alpha/p\), while the
remaining \(p-1\) points have valuation \(B>\alpha/p\).
This gives \(B_H=s-B\).
If \(B<e\), the inverse roots over each nonzero high target
have valuations \((s-pB)/p\), \(p\) times, and \(B/p\),
\(q-p\) times. All lie below \(\alpha\). Including the image
of \(M[p]/H\) and dividing preimage counts by \(p\) gives
\(p-1\) points of valuation \(s-pB\) and \(q-p\) of valuation
\(B\), hence \(B_H=pB\).
At the equality \(B=e\), the inverse polygon has one segment
of slope \(-e/p\), with \(q\) roots. Both contributions to
the quotient therefore have valuation \(e\).

For a low line every nonzero class of \(M[p]/J\) has a unique
high representative. Its quotient valuation is
\(\alpha+(p-1)\beta=s-B/p\).
Preimages of nonzero low targets have valuation \(\beta/q\);
after quotienting, the \(q-p\) resulting points have valuation
\(\beta/p=B/p^2\). This proves \(B_J=B/p\), an integer because
the first-level module is \(L\)-rational.

Now start from a one-slope polygon. The product of the \(q-1\)
nonzero \(p\)-torsion parameters has valuation \(v_L(p)=eN\),
so every such parameter has valuation \(e\). For a nonzero
class of \(M[p]/J\), all \(p\) differences in (B-RD69.7) are
nonzero \(p\)-torsion parameters and have valuation \(e\).
Its image therefore has valuation \(pe\), giving \(p-1\)
nonzero quotient points. A nonzero target in \(J\) has valuation
\(e\); the inverse Newton polygon is one segment, with \(q\)
roots of valuation \(e/q\). These values are strictly below
those of nonzero points of \(J\). Formula (B-RD69.7) then gives
valuation \(e/p\). There are \((p-1)q/p=q-p\) distinct such
quotient points. The two contributions are disjoint and exhaust
the \(q-1\) nonzero points, proving (B-RD69.6).

The unique rank-\(p\) subgroup of the one-dimensional connected
special formal group is its Frobenius kernel. Thus the special
quotient is its Frobenius twist; the inherited special tangent
exponent is \(pc\), and the unordered pair \(\{c,pc\}\) is
unchanged. Generic inertia stability of \(J\) gives actual
generic descent of the quotient to \(K\).
More explicitly, Tate full faithfulness extends the original
semilinear generic descent maps and their inverses to the good
models. They preserve the schematic closure of \(J\), hence
induce semilinear quotient maps satisfying the same cocycle.
Only the generic quotient is being descended through the
ramified extension; no integral good model over \(\mathcal O_K\)
is asserted.
Finally the actual generic identity
\[
 (M/J)[p^2]=\{P\in M[p^3]:p^2P\in J\}/J
 \tag{B-RD69.8}
\]
proves the claimed passage from a \(p^3\)-comparison to a
\(p^2\)-comparison. Here both sides are finite groups over
the characteristic-zero field, so the identity also identifies
their finite group schemes and their Galois actions.
\end{proof}

\begin{proposition}[Excluding a mixture of Newton polygons]
\label{prop:B-RD69-one-slope}
\status{Proved; complete independent root and companion proof, source and exact-code audits passed.}
For every \(e\) as in (B-RD69.1), two whole potentially
supersingular height-two groups with isomorphic actual
\(p^3\)-torsion cannot have one one-slope and one two-slope
first-level Newton polygon. If both have one slope, their finite
special characters agree.
\end{proposition}
\begin{proof}
For a one-slope group the map taking a \(p\)-torsion coordinate
to the leading coefficient of \(T/\pi_L^e\) is injective and
\(\mathbf F_p\)-linear. Indeed every nonzero difference is a
nonzero \(p\)-torsion point and has valuation \(e\); the higher
terms of the formal group law have strictly larger valuation.
Inertia over \(L\) fixes this map, because it fixes \(\pi_L\)
and the algebraically closed residue field. Thus all \(p\)-torsion
is \(L\)-rational. The chosen generator over \(K\) acts on the
two-dimensional leading-coefficient space by multiplication by
\[
 \lambda=\zeta^{c+e}\in\mathbf F_q^\times .
 \tag{B-RD69.9}
\]
If its \(\mathbf F_p\)-representation is reducible, an invariant
line gives \(\lambda\in\mathbf F_p^\times\); its action on the
entire module is then scalar. Otherwise the residual module
is irreducible over \(\mathbf F_p\).

Suppose the first polygon is one-slope and the second two-slope.
The latter has a stable high line, so their common residual
module is reducible, hence scalar by the preceding paragraph.
In particular every line is stable under the actual inertia
action, and the second first-level module is \(L\)-rational.
Write its parameter as \(B\); then \(p\mid B\).
Transport its high line across the actual torsion isomorphism
to a line in the first group, and quotient both groups.
Lemma~\ref{lem:B-RD69-quotients} gives an actual \(p^2\)-comparison
with first parameter \(e\), prime to \(p\). If \(B<e\), the
second parameter is \(pB\), divisible by \(p\). Such a comparison
is impossible by Lemma~\ref{lem:B-RD69-rules}: the first parameter
forces nontrivial residual wild inertia in the two quotients.
If \(B>e\), the second parameter is \(s-B\), prime to \(p\)
because \(p\nmid s\). Equality of the two parameters forces
\(s-B=e\), or \(B=pe\), contrary to the strict two-slope bound.
The equality \(B=e\) cannot occur here, because \(p\mid B\)
whereas \(p\nmid e\). This excludes mixed polygons.

If both polygons have one slope and their common residual
module is irreducible, Theorem~\ref{thm:B-RD67} already recovers
their special characters from their actual \(p^2\)-torsion.
If it is reducible, both residual actions are scalar.
An isomorphism of the actual residual representations equates
the two scalars in (B-RD69.9). With the same generator and the
same \(e\), this gives \(c_1=c_2\pmod N\), hence the same
special character. This uses the actual scalar action, not
merely its determinant.
\end{proof}

\begin{proposition}[Two unequal polygons below \(p^2\)]
\label{prop:B-RD69-two-slope}
\status{Proved; complete independent root and companion proof, source and exact-code audits passed.}
Suppose in addition \(e<p^2\). Two whole potentially
supersingular height-two groups with isomorphic actual
\(p^3\)-torsion and two-slope first-level polygons have the
same finite special characters, regardless of the square
class of their common finite determinant.
\end{proposition}
\begin{proof}
Write their special tangent exponents as \(c_i\).
Their generic determinant is \(\chi_p\epsilon_i\), with
\(\epsilon_i\) finite and \(\mu_{p-1}\)-valued.
The actual torsion comparison and \(\mu_p\subset K\) imply
\(\epsilon_1=\epsilon_2\); reduction modulo \(p\) is already
injective on these finite characters. Its special reduced-norm
description gives
\[
 c_1\equiv c_2\pmod{p-1}.
 \tag{B-RD69.10}
\]
This equality, rather than a normalization to a primitive norm,
is all that is needed below.

If residual wild inertia is nontrivial, the high line is common.
Lemma~\ref{lem:B-RD69-rules} handles every case in which at
least one \(B_i\) is prime to \(p\). This expressly includes
the possible critical value \(B_i=e\).
For the remaining cases both \(B_i\) are divisible by \(p\),
so neither equals \(e\). If both are below \(e\), write
\(B_i=pb_i\). Their common high-line character and (B-RD69.10)
give \(b_1\equiv b_2\pmod{p-1}\). But
\[
 0<b_i<e/p<p,
 \tag{B-RD69.11}
\]
so both are in the interval \(\{1,\ldots,p-1\}\).
Reduction modulo \(p-1\) is injective on this interval, and
\(b_1=b_2\). The common high-line character identifies \(c_1=c_2\).
If both \(B_i>e\), their high quotients have parameters
\(s-B_i\) prime to \(p\); their actual \(p^2\)-comparison and
Lemma~\ref{lem:B-RD69-rules} give equal quotient characters,
hence equal original ones.
If the two \(B_i\)'s lie on opposite sides of \(e\), the high
quotient parameters have opposite \(p\)-divisibility and their
actual \(p^2\)-comparison is impossible by the same lemma.

It remains that the residual module is tame. Write \(B_i=pb_i\),
where \(0<b_i<e\); again \(B_i=e\) is impossible.
Two \(b_i\) prime to \(p\) are settled by the lemma.
Suppose \(p\nmid b_1=b\) and \(b_2=pk\).
Then \(B_2=p^2k>e\), since \(e<p^2\).
Choose a common residual vector outside both high lines and
corresponding primitive lifts. Their common point field has
degree \(p^2\) over \(L\), forced by the first valuation
\(b/p^2\); its two normalized coordinate valuations are
\(b\) and \(pk\). If the high lines differ, choose a translation
high in the first model and low in the second.
The prime-to-\(p\) valuation of the first coordinate fixes its
exact lower break as \(p^2\alpha_1-b\). The universal lower
bound for the second coordinate would give
\[
 p^2(\alpha_1-pk)\le b-pk .
 \tag{B-RD69.12}
\]
However \(\alpha_1=s-pb\ge e+p\), since \(b<e\) is integral,
while \(pk<e\) and \(b<e<p^2\). The left side is greater than
\(p^3\) and the right side less than \(p^2\). Thus the high
lines agree. Quotient by that common line.
The second new parameter \(s-B_2\) is prime to \(p\).
If \(B_1<e\), the first new parameter \(pB_1\) is divisible
by \(p\), an impossibility. If \(B_1>e\), equality of the two
prime-to-\(p\) new parameters forces \(B_1=B_2\), contrary
to \(p\nmid b\) and \(p\mid b_2\).

Finally suppose \(b_i=pk_i\) for both groups.
Then \(B_i=p^2k_i>e\). If their high lines agree, quotienting
by that line gives two prime-to-\(p\) parameters and recovers
the characters. If their high lines differ, quotient the
first by its high line and the second by the corresponding
stable low line. Both lines are stable because they correspond
under the actual residual isomorphism.
The new parameters \(s-B_1\) and \(B_2/p=pk_2\) have opposite
\(p\)-divisibility, contradicting their inherited actual
\(p^2\)-torsion isomorphism.
These cases exhaust the residual actions and all divisibilities.
Only \(p\)-divisible original parameters are quotiented in this
case analysis, so the critical quotient at \(B=e\) is never
silently treated as a two-slope quotient.
\end{proof}

\begin{theorem}[The first tame boundary has exact index two]
\label{thm:B-RD69}
\status{Proved; complete independent root and companion proof, source and exact-code audits passed.}
For every prime \(p\ge5\), in the full B.4 monodromy setting
with height two, perfect residue field and \(\mu_p\subset K\),
\[
 \boxed{\displaystyle \nu_{\min}(p,p^2-1,2)=2.}
 \tag{B-RD69.13}
\]
In particular \(\nu_{\min}(5,24,2)=2\).
More generally, actual \(p^3\)-torsion determines the requested
character whenever \(p-1\mid e\), \(p\nmid e\), and \(e<p^2\).
\end{theorem}
\begin{proof}
Pass to the completed maximal unramified base, retaining the
absolute ramification and the actual torsion comparison.
The determinant of an allowed height-two monodromy system is
\(\chi_p^d\epsilon\), with \(d\in\{0,1,2\}\) and finite
\(\mu_{p-1}\)-valued \(\epsilon\), by the dimension and
filtration argument in Theorem~\ref{thm:B-RD68}.
As \(p\nmid e\), the cyclotomic image modulo \(p^3\) has
order \(p^2\). Equality of the actual determinants gives the
same \(d\), since \(|d_1-d_2|\le2<p\), and then the same
\(\epsilon\).
Dimensions zero and two are determined by residual inertia
and Cartier duality. In total dimension one, two systems
filtered by potentially good height-one pieces have their
finite characters determined modulo \(p\).
A comparison of a whole potentially supersingular group
with a filtered one is excluded by
Proposition~\ref{prop:B-RD63-mixed}, because \(e<p^2<p^3\).
For the two whole groups, the possible Newton-polygon pairs
are exhausted by Propositions~\ref{prop:B-RD69-one-slope}
and~\ref{prop:B-RD69-two-slope}. They prove the asserted
upper bound throughout the stated range.

At \(e=p^2-1>p\), the explicit comparison of
Theorem~\ref{thm:B-RD63-arc}, with \(n=2\), gives two permitted
height-two monodromy systems with equal actual \(p^2\)-torsion
and different requested characters. Thus manuscript index one
does not suffice, while index two does. This proves the exact
value, not only a mixed-case necessary bound.
\end{proof}

\begin{remark}[Scope and the critical value]
\label{rem:B-RD69-scope}
\status{Proved; complete independent root and companion proof, source and exact-code audits passed.}
The one-slope/two-slope exclusion is valid for every
prime-to-\(p\) absolute index divisible by \(p-1\).
The recovery proof for two unequal polygons uses \(e<p^2\)
in the small-parameter interval and in the tame divisibility
case split. It imposes no square-determinant hypothesis and
uses actual rank-\(p\) quotients, not transferred Hasse thickness.
The value \(B=e\), absent from the earlier odd-\(B\) normalized
proof, is included explicitly: it is prime to \(p\), is handled
before the quotient case analysis, and its actual high quotient
is correctly recorded as one-slope.
This result does not cover \(e=p(p-1)\), nor any other absolute
index divisible by \(p\).
The retained exact supplement
\path{code/reduction_descent/verify_one_slope_boundary.py}
reconstructs the Newton polygons with rational arithmetic and
counts quotient roots from the actual product formula.
For \(72\) sectors at \(p=5,7,11,13,17,19\), it checks
\(130{,}398\) unequal polygons and high quotients,
\(216\) one-slope product quotients, all \(72\) critical
\(B=e\) quotients, \(7{,}842\) low quotients,
\(2{,}099{,}574\) unrestricted tangent weights, and
\(76{,}764\) tame mixed-divisibility inequalities.
It also checks the one-/two-slope arithmetic at additional
indices above \(p^2\), where only that particular exclusion is
claimed. The author run passed in \(7.471134923995123\) seconds
at 21:28:53 UTC. The exact program, full output and pre-certificate
proof are retained. These finite checks supplement the all-prime
proofs; they do not establish existence or nonexistence of
formal-group comparisons by enumeration alone.
The complete independent root and companion audits passed.
Their author-output-preserving exact reruns took
\(6.916882791992975\) and \(7.527760954006226\) seconds,
respectively. The exact audited candidate remains preserved.
\end{remark}


% Proved after full root and companion audits; exact candidate and outputs retained.
\subsection{A Frobenius-orbit collision at the next tame boundary}

Fix \(p\ge5\), put \(q=p^2\), \(N=q-1\), and set
\[
 e=(p+2)(p-1)=q+p-2,\qquad s=e(p+1)=(p+2)N.
 \tag{B-RD70.1}
\]
The complete mixed-characteristic \((0,p)\) field \(K\) has
algebraically closed residue field, \(\mu_p\subset K\), and
absolute ramification \(e\). Thus \(p\nmid e\),
\(q<e<2q\), and \(e<p^3\).
We use the actual good tame extension \(L/K\) of degree \(N\),
and the tangent exponent \(c\), exactly as in B69.
For a two-slope group its high line, parameter and valuations are
\[
 H,\qquad 0<B<pe,\quad B\equiv c\pmod{p+1},\qquad
 \alpha=s-B,\quad\beta=B/p .
 \tag{B-RD70.2}
\]
The conditional comparison rules and the actual high/low quotient
identities of Lemmas~\ref{lem:B-RD69-rules}
and~\ref{lem:B-RD69-quotients} apply at this \(e\): neither
lemma requires \(e<p^2\).

\begin{lemma}[Recovery when the high line is common]
\label{lem:B-RD70-high}
\status{Proved; complete independent root and companion proof and exact-code audits passed.}
Let two whole potentially supersingular height-two groups have
isomorphic actual \(p^3\)-torsion, two-slope first-level polygons,
and corresponding high lines. Their finite special inertia
characters agree.
\end{lemma}
\begin{proof}
The actual comparison identifies their finite determinant
characters. As in (B-RD69.10), their tangent exponents therefore
satisfy \(c_1\equiv c_2\pmod{p-1}\).
Their common high-line character gives
\[
 c_1-B_1\equiv c_2-B_2\pmod N .
 \tag{B-RD70.3}
\]
If either \(B_i\) is prime to \(p\), the common residual action
over \(L\) is wild, since its low coordinates have nonintegral
\(v_L\)-valuation \(B_i/p\). The first rule of
Lemma~\ref{lem:B-RD69-rules} then proves equality of the
characters or excludes a mixed-divisibility comparison.
This includes the possible critical value \(B_i=e\).

For the remaining cases both \(B_i\) are divisible by \(p\);
in particular neither equals \(e\).
If both \(B_i<e\), write \(B_i=pb_i\).
The common norm and (B-RD70.3) imply
\(b_1\equiv b_2\pmod{p-1}\), while
\[
 0<b_i<e/p=p+1-2/p,\qquad b_i\in\{1,\ldots,p\}.
 \tag{B-RD70.4}
\]
Consequently either \(b_1=b_2\) or
\(\{b_1,b_2\}=\{1,p\}\).
Equality gives \(c_1=c_2\pmod N\).
In the unequal case interchange the two models if necessary
so that \(b_1=1\), \(b_2=p\). Then
\[
 c_2-c_1\equiv p(p-1)\pmod N,\qquad
 c_1\equiv B_1=p\pmod{p+1}.
 \tag{B-RD70.5}
\]
Multiplying the latter congruence by \(p-1\) shows
\((p-1)c_1\equiv p(p-1)\pmod N\), so \(c_2\equiv pc_1\pmod N\).
The two unordered special orbits \(\{c_i,pc_i\}\) are equal.
Here unequal numerical parameters give the same character;
one must not incorrectly infer equality of the parameters.

If both \(B_i>e\), quotient by the common high line.
The inherited actual \(p^2\)-comparison has parameters
\(s-B_i\), both prime to \(p\), because \(p\nmid s\).
Lemma~\ref{lem:B-RD69-rules} recovers the quotient characters,
and isogeny invariance recovers the original ones.
If one \(B_i\) is below \(e\) and the other above, the high
quotient parameters are respectively \(pB_i\) and \(s-B_j\).
Their opposite \(p\)-divisibility contradicts the same lemma.
These are all possibilities.
\end{proof}

\begin{lemma}[A constraint on exchanged tame high lines]
\label{lem:B-RD70-exchange}
\status{Proved; complete independent root and companion proof and exact-code audits passed.}
Suppose two two-slope groups have isomorphic residual modules
with tame inertia. If the isomorphism does not identify their
high lines, then
\[
                         p-1\mid B_1+B_2.
 \tag{B-RD70.6}
\]
In particular two such groups with \(B_1=B_2=p^2\) must have
corresponding high lines.
\end{lemma}
\begin{proof}
The high eigenvalue of the tame generator is
\(\zeta^{c_i-B_i}\). Its determinant is
\(\zeta^{(p+1)c_i}\), because \(\mu_p\subset K\).
Thus the other residual eigenvalue is \(\zeta^{pc_i+B_i}\).
These are characters of the actual residual representation,
not of the special isocrystal.
Tame residual inertia is semisimple. If its two eigenvalues
are distinct, every stable line other than its high line
is its other eigenspace. If the representation is scalar,
the other eigenvalue has the same value and the following
identities remain valid for every line. Hence distinct
corresponding high lines force
\[
 c_1-B_1\equiv pc_2+B_2,\qquad
 c_2-B_2\equiv pc_1+B_1\pmod N .
 \tag{B-RD70.7}
\]
Multiplying the first congruence by \(p\), adding the second,
and using \(p^2\equiv1\pmod N\), gives
\(N\mid(p+1)(B_1+B_2)\). This is precisely (B-RD70.6).
For \(B_1=B_2=p^2\), it would require \(p-1\mid2p^2\),
or \(p-1\mid2\), impossible for \(p\ge5\).
\end{proof}

\begin{proposition}[Tame reducible comparisons at the next boundary]
\label{prop:B-RD70-tame}
\status{Proved; complete independent root and companion proof and exact-code audits passed.}
Two whole potentially supersingular height-two groups at the
index (B-RD70.1), with two-slope first-level polygons, tame
residual inertia and isomorphic actual \(p^3\)-torsion, have
equal finite special inertia characters.
\end{proposition}
\begin{proof}
Their first-level modules are rational over \(L\), so write
\(B_i=pb_i\), with \(0<b_i<e\).
If both \(b_i\) are prime to \(p\), the second comparison
rule of Lemma~\ref{lem:B-RD69-rules} already recovers their
characters from actual \(p^2\)-torsion.
Suppose exactly one is prime to \(p\), say \(b_1=b\), while
\(b_2=pk\). Choose corresponding residual vectors outside
both high lines and corresponding primitive lifts.
The first lift has valuation \(b/p^2\). Its denominator
forces the common point field to have degree \(p^2\)
over \(L\); all translations by the \(L\)-rational first-level
module occur. Its normalized coordinate valuations in that
field are \(b\) and \(pk\).

If the high lines differ, choose a translation high in the
first group and low in the second. The first coordinate has
prime-to-\(p\) valuation, so the actual lower break is
\(p^2\alpha_1-b\). The lower bound for the second gives
\[
 p^2(\alpha_1-pk)\le b-pk .
 \tag{B-RD70.8}
\]
But \(\alpha_1=s-pb\ge e+p\), and \(pk<e\).
Consequently the left side is greater than \(p^3\),
whereas the right side is less than \(e<p^3\).
Thus the high lines agree, and
Lemma~\ref{lem:B-RD70-high} proves the result.
Unlike the earlier proof below \(p^2\), this step does not
assert \(p^2k>e\); it allows the new case \(k=1\).

It remains that \(b_i=pk_i\) for both groups, so
\(B_i=p^2k_i\). If the high lines agree, use the same lemma.
Suppose they do not. If either \(B_i>e\), label that group
first, quotient it by its high line, and quotient the other
by the corresponding low line. The lines are inertia stable
because they correspond under the actual Galois-module
isomorphism. The inherited actual \(p^2\)-comparison has
parameters
\[
                  s-B_1\quad\text{and}\quad B_2/p=pk_2.
 \tag{B-RD70.9}
\]
The first is prime to \(p\), the second divisible by \(p\);
Lemma~\ref{lem:B-RD69-rules} forbids this comparison.
Neither original \(B_i\) equals \(e\), since \(p\mid B_i\)
and \(p\nmid e\). The only remaining possibility is
\(B_1<e\) and \(B_2<e\). Since \(p^2<e<2p^2\), this
forces \(B_1=B_2=p^2\).
Lemma~\ref{lem:B-RD70-exchange} now says that their high
lines must agree, contradicting the supposition.
This exhausts every tame divisibility case.
\end{proof}

\begin{proposition}[All whole formal comparisons at this index]
\label{prop:B-RD70-formal}
\status{Proved; complete independent root and companion proof and exact-code audits passed.}
At \(e=(p+2)(p-1)\), actual \(p^3\)-torsion determines the
finite special inertia character among whole potentially
supersingular height-two groups.
\end{proposition}
\begin{proof}
Proposition~\ref{prop:B-RD69-one-slope}, which has no upper
bound on \(e\), excludes a mixture of one-slope and two-slope
polygons and recovers two one-slope characters.
For two two-slope groups with nontrivial residual wild inertia,
the high line is the unique wild fixed line, hence common;
Lemma~\ref{lem:B-RD70-high} applies.
Tame residual inertia is precisely the preceding proposition.
The argument covers all finite determinant characters without
a square-class normalization.
\end{proof}

\begin{theorem}[The next tame family has exact index two]
\label{thm:B-RD70}
\status{Proved; complete independent root and companion proof and exact-code audits passed.}
For every prime \(p\ge5\), in the full B.4 monodromy setting
with height two, perfect residue field and \(\mu_p\subset K\),
\[
 \boxed{\displaystyle
       \nu_{\min}\bigl(p,(p+2)(p-1),2\bigr)=2 .}
 \tag{B-RD70.10}
\]
In particular \(\nu_{\min}(5,28,2)=2\).
\end{theorem}
\begin{proof}
For the upper bound pass to the completed maximal unramified
base. Since \(p\nmid e\), the cyclotomic image modulo \(p^3\)
has order \(p^2\); the determinant argument of
Theorem~\ref{thm:B-RD69} identifies the total dimension
\(d\in\{0,1,2\}\) and the finite determinant.
Dimensions zero and two, and two systems filtered by
potentially good height-one pieces, are detected residually
as in that proof. In dimension one a mixed whole/filtered
comparison is excluded by
Proposition~\ref{prop:B-RD63-mixed}, since \(e<p^3\).
Two whole potentially supersingular groups are covered by
Proposition~\ref{prop:B-RD70-formal}.
Thus actual \(p^3\)-torsion always suffices.
As \(e>p\), Theorem~\ref{thm:B-RD63-arc} with \(n=2\)
provides permitted height-two systems with equal actual
\(p^2\)-torsion and different requested characters.
Index one therefore fails and index two is least.
\end{proof}

\begin{remark}[The orbit collision and the next limitation]
\label{rem:B-RD70-scope}
\status{Proved; complete independent root and companion proof and exact-code audits passed.}
The new small-parameter collision is \(\{b_1,b_2\}=\{1,p\}\).
It preserves the unordered special character by exchanging its
two Frobenius-conjugate eigenvalues. The proof does not rule out
all unequal Newton parameters.
The new tame case \(B_1=B_2=p^2\) is handled by actual residual
eigenvalues, not by a square-determinant quotient-thickness bound.
At \(e=(p+3)(p-1)\), the small interval first admits
\(\{2,p+1\}\), which is not automatically a Frobenius orbit;
no conclusion for that index follows from this proof.
No absolute ramification divisible by \(p\) is treated.
The retained exact supplement
\path{code/reduction_descent/verify_next_tame_boundary.py}
checks twelve primes from \(5\) through \(43\).
It verifies \(8{,}496\) common-high-line character orbits,
including \(528\) unequal-parameter Frobenius collisions,
\(7{,}704\) critical exchanged-line pairs,
\(280{,}368\) mixed point-field inequalities, and every
high/low quotient-divisibility case in the stated finite range.
The first unexcluded later arithmetic pattern
\((b_1,b_2)=(2,p+1)\) is retained explicitly; it is not
asserted to arise from formal groups with isomorphic torsion.
The author run passed in \(0.14006940000399482\) seconds at
21:40:22 UTC. An unused malformed draft block was removed
before the first execution, and that unexecuted draft is
preserved. No executed numerical check failed.
These finite checks supplement the preceding all-prime proofs.
The complete independent root and companion audits passed,
with author-output-preserving exact reruns in
\(0.16807877800602\) and \(0.1444874750013696\) seconds,
respectively. Both full audit records and the exact
prepromotion candidate are retained.
\end{remark}


% Proved after full root and companion audits; exact candidate and outputs retained.
\subsection{Marked low lines and a larger exact tame sector}

Throughout the local lemmas \(p\ge5\), \(q=p^2\), \(N=q-1\),
and \(K\) is complete discretely valued of mixed characteristic
\((0,p)\), with algebraically closed residue field, \(\mu_p\subset K\),
and absolute ramification \(e\) prime to \(p\). Put \(s=e(p+1)\).
The actual good field \(L/K\) is tame of degree \(N\), with
\(v_L(p)=eN=(p-1)s\), and \(s\) is a multiple of \(N\).
For a two-slope group use the notation of B69:
\[
 0<B<pe,\quad B\equiv c\pmod{p+1},\quad
 \alpha=s-B,\quad\beta=B/p,\quad
 \lambda_H=\zeta^{c-B}.
 \tag{B-RD71.1}
\]
All coordinate valuations below belong to actual integral good
models and their finite-flat quotients.

\begin{lemma}[Prime Newton parameters are determined at the first level]
\label{lem:B-RD71-first-level}
\status{Proved; complete independent root and companion proof and exact-code audits passed.}
Two two-slope groups with parameters \(B_1,B_2\) prime to \(p\)
and isomorphic actual \(p\)-torsion have \(B_1=B_2\) and the
same finite special inertia character.
\end{lemma}
\begin{proof}
The common first-level division field \(R/L\) is cyclic of
degree \(p\). Indeed the high character and determinant are
trivial over \(L\), so the residual action lies in the
upper-unipotent group of order \(p\); the low valuation \(B_i/p\)
is not integral over \(L\), so the action is nontrivial.
The high lines are the common unique wild fixed line.
In \(R\), a low point has coordinate valuation \(B_i\).
For any nonidentity \(g\in\operatorname{Gal}(R/L)\), its
displacement is a nonzero high point, of valuation
\(p(s-B_i)\). Formal subtraction has the same valuation as
coordinate subtraction, since the relevant parameters have
positive valuation.
The exact prime-valuation part of
Lemma~\ref{lem:B-RD60-break} therefore gives the actual lower break
\[
 r=p(s-B_i)-B_i=ps-(p+1)B_i.
 \tag{B-RD71.2}
\]
The field and break are common, so \(B_1=B_2\).
Their identified high-line characters in (B-RD71.1) then
give \(c_1=c_2\pmod N\). This argument uses only the first-level
torsion isomorphism; the second-level assumption in the broader
comparison lemma is unnecessary for this particular branch.
\end{proof}

\begin{lemma}[The preserved low line in a small high quotient]
\label{lem:B-RD71-marked-low}
\status{Proved; complete independent root and companion proof and exact-code audits passed.}
Let \(M\) have two-slope parameter \(B<e\) and high line \(H\).
The actual quotient \(M'=M/H\) has parameter \(D=pB\), and
\[
 J=\operatorname{im}\bigl(M[p]\longrightarrow M'[p]\bigr)
   \simeq M[p]/H
 \tag{B-RD71.3}
\]
is an inertia-stable low line of \(M'[p]\), of valuation \(B\).
The first-level module \(M'[p]\) is rational over \(L\).
An actual \(p^3\)-torsion isomorphism identifying two original
high lines induces an actual \(p^2\)-torsion isomorphism of
the quotients which also identifies these marked low lines.
The quotient's low/high eigenvalue ratio is
\(\zeta^{(p+1)B}\); it may be scalar.
\end{lemma}
\begin{proof}
The high-quotient product calculation in
Lemma~\ref{lem:B-RD69-quotients} gives \(D=pB\).
The \(p-1\) nonzero points of \(M[p]/H\) have quotient
valuation \(B\), whereas the high valuation in \(M'[p]\)
is \(s-pB>B\), since \(B<e\).
Thus their image is precisely a low line, and is stable
under the actual generic inertia action.
The high line of \(M'[p]\) is stable too, because it is
characterized by its strictly larger valuation.
These two distinct invariant \(\mathbf F_p\)-lines make
the entire residual representation diagonal with
\(\mathbf F_p^\times\)-valued characters. They are killed
by the tame extension of degree \(N\); consequently \(M'[p]\)
is \(L\)-rational, even when the original residual action is wild.

The actual quotient has tangent exponent \(pc\).
Its high eigenvalue is \(\zeta^{pc-pB}\), and its other
eigenvalue is \(\zeta^{c+pB}\), by its determinant and
\(p^2\equiv1\pmod N\). Their ratio has exponent
\((1-p)c+2pB\), congruent to \((p+1)B\), because
\((p-1)(c-B)\) is divisible by \(N\).
In particular scalar action can occur when \(p-1\mid B\).
No uniqueness of an unmarked complementary line is assumed.

The induced actual \(p^2\)-isomorphism is obtained from
\(\{P\in M[p^3]:p^2P\in H\}/H\).
It carries the image of the intrinsic subgroup \(M[p]\)
to the corresponding image in the other quotient. This
proves preservation of \(J\) as an actual Galois-module
subgroup. The finite-flat quotients and semilinear generic
descent are those of Lemma~\ref{lem:B-RD69-quotients}.
\end{proof}

\begin{proposition}[A scalar-safe recovery of the quotient high line]
\label{prop:B-RD71-separated-high}
\status{Proved; complete independent root and companion proof and exact-code audits passed.}
Suppose an actual \(p^3\)-torsion isomorphism identifies the
high lines of two two-slope groups, both with
\[
 B_i=pb_i<e.
 \tag{B-RD71.4}
\]
If \(q\min(b_1,b_2)>e\), then \(B_1=B_2\) and their special
characters agree. This assertion includes scalar residual
actions in the first quotients.
\end{proposition}
\begin{proof}
If \(b_1=b_2\), the common high-line character already
identifies \(c_1=c_2\). Otherwise interchange the models
so that \(b_1<b_2\).
Quotient by the common high line and write \(M_i'=M_i/H_i\).
By Lemma~\ref{lem:B-RD71-marked-low} these groups have an
actual \(p^2\)-torsion isomorphism \(\phi'\), rational
first-level modules over \(L\), and parameters
\[
 D_i=q b_i,\qquad D_i>e,\qquad
 A_i=s-D_i,\qquad C_i=D_i/p .
 \tag{B-RD71.5}
\]
Here \(A_i\) and \(C_i\) are their high and low first-level
valuations. In particular \(p\nmid A_i\).

We prove that \(\phi'\) identifies their high lines.
Suppose otherwise. Choose a primitive order-\(p^2\) point
\(X_1\) of \(M_1'\) whose \(p\)-multiple is a nonzero high
point, and put \(X_2=\phi'(X_1)\).
The \(p\)-multiple of \(X_2\) is then low in \(M_2'\).
Because \(D_1>e\), the high-target inverse Newton polygon
of \(M_1'\) has a single segment of degree \(q\):
the exact inverse calculation in Lemma~\ref{lem:B-RD69-quotients}
gives
\[
 v_L(X_1)=A_1/q,\qquad
 v_L(X_2)=D_2/p^3=C_2/q .
 \tag{B-RD71.6}
\]
The second identity is the low-target inverse calculation.
Since \(A_1=s-qb_1\) is prime to \(p\), the denominator of
the first valuation is exactly \(q\).

Both \(p\)-multiples are \(L\)-rational, and both first-level
modules are \(L\)-rational. Therefore each point field over
\(L\) is Galois with group a subgroup of its first-level
translation group, of order at most \(q\): every conjugate
is \(X_i+Q\), and all such sums lie in the point field.
The actual torsion isomorphism identifies their stabilizers,
so the two point fields are the same field \(E\).
The first valuation denominator forces its ramification
index to be at least \(q\). Consequently
\[
 [E:L]=e(E/L)=q,\qquad
 v_E(X_1)=A_1,\quad v_E(X_2)=C_2,
 \tag{B-RD71.7}
\]
and every first-level translation occurs.

Choose \(g\in\operatorname{Gal}(E/L)\) which translates \(X_1\)
by a nonzero high point of \(M_1'[p]\).
Its corresponding translation in \(M_2'[p]\) is low, by
the assumed failure to identify high lines.
Thus its coordinate displacements have valuations
\(qA_1\) and \(qC_2\), respectively.
Since \(p\nmid A_1\), the exact lower-break lemma gives
\(b(g)=(q-1)A_1\).
The unconditional lower bound on the second coordinate
then requires
\[
 qC_2\ge C_2+(q-1)A_1,\qquad\text{hence }C_2\ge A_1.
 \tag{B-RD71.8}
\]
Yet \(pb_2=B_2<e\), so
\[
 \begin{aligned}
 A_1-C_2
   &=s-qb_1-pb_2\\
   &=s-p(p+1)b_2+q(b_2-b_1)
     >q(b_2-b_1)>0.
 \end{aligned}
 \tag{B-RD71.9}
\]
This contradiction identifies the high lines, without
assuming distinct residual eigenvalues.

We may therefore quotient both \(M_i'\) by their high lines.
Since \(D_i>e\), their actual high quotients have parameters
\(s-D_i=A_i\), both prime to \(p\).
The induced comparison is now exactly at the first level:
\[
 (M_i'/H_i')[p]
   =\{Z\in M_i'[p^2]:pZ\in H_i'\}/H_i'.
 \tag{B-RD71.10}
\]
Lemma~\ref{lem:B-RD71-first-level} gives \(A_1=A_2\),
so \(D_1=D_2\), contrary to \(b_1<b_2\).
Thus only \(b_1=b_2\) was possible, proving the assertion.
\end{proof}

\begin{proposition}[Common-high comparisons in the enlarged sector]
\label{prop:B-RD71-common-high}
\status{Proved; complete independent root and companion proof and exact-code audits passed.}
Suppose now that
\[
 e=t(p-1),\qquad 2\le t\le2p+1,\qquad p\nmid t.
 \tag{B-RD71.11}
\]
An actual \(p^3\)-torsion isomorphism between two two-slope
groups which identifies their high lines determines their
finite special inertia characters.
\end{proposition}
\begin{proof}
Their common finite determinant and high-line character give
\[
 c_1\equiv c_2\pmod{p-1},\qquad
 c_1-B_1\equiv c_2-B_2\pmod N.
 \tag{B-RD71.12}
\]
Any case with a prime-to-\(p\) parameter is already settled
by Lemma~\ref{lem:B-RD69-rules}; it includes \(B_i=e\).
Assume both \(B_i\) are divisible by \(p\).
Neither equals \(e\). If both are above \(e\), their actual
high quotients have prime parameters \(s-B_i\), and the
inherited \(p^2\)-comparison recovers their characters.
If the parameters lie on opposite sides of \(e\), those
quotients have opposite \(p\)-divisibility and the same
comparison lemma excludes them.

It remains that \(B_i=pb_i<e\). The range (B-RD71.11) gives
\[
 e\le2p^2-p-1<2q,\qquad 1\le b_i\le2p-2.
 \tag{B-RD71.13}
\]
If the \(b_i\)'s are equal, (B-RD71.12) identifies \(c_1=c_2\).
Otherwise take \(b_1<b_2\).
If \(b_1\ge2\), then \(qb_1\ge2q>e\);
Proposition~\ref{prop:B-RD71-separated-high} forbids the
inequality \(b_1<b_2\).
Thus \(b_1=1\).
Equation (B-RD71.12) implies \(b_2\equiv1\pmod{p-1}\),
and the upper bound \(b_2\le2p-2\) forces \(b_2=p\).
Exactly as in (B-RD70.5),
\[
 c_2-c_1\equiv p(p-1),\qquad
 c_1\equiv p\pmod{p+1}
 \quad\Longrightarrow\quad c_2\equiv pc_1\pmod N.
 \tag{B-RD71.14}
\]
This exchanges the two Frobenius-conjugate eigenvalues
and preserves the special character. Every case is covered.
\end{proof}

\begin{proposition}[Tame residual comparisons in the enlarged sector]
\label{prop:B-RD71-tame}
\status{Proved; complete independent root and companion proof and exact-code audits passed.}
Under (B-RD71.11), two two-slope groups with tame residual
inertia and isomorphic actual \(p^3\)-torsion have equal
finite special inertia characters.
\end{proposition}
\begin{proof}
Write \(B_i=pb_i\), \(0<b_i<e\), since their first-level
modules are rational over \(L\).
If both \(b_i\) are prime to \(p\), the second rule of
Lemma~\ref{lem:B-RD69-rules} gives the conclusion.
If exactly one is prime, say \(b_1=b\), \(b_2=pk\), take
primitive points whose \(p\)-multiples are outside both
high lines. Their common point field over \(L\) has degree
\(q\), forced by the first valuation \(b/q\).
Their normalized coordinate valuations are \(b\) and \(pk\).
Distinct high lines would permit a translation high in
the first model and low in the second. The exact break
from the first coordinate and the lower bound for the second
would give
\[
 q(\alpha_1-pk)\le b-pk .
 \tag{B-RD71.15}
\]
But \(\alpha_1=s-pb\ge e+p\), \(pk<e\), and
\(e<2q<p^3\). The left side is greater than \(p^3\),
the right side less than \(p^3\). Thus their high lines
agree, and Proposition~\ref{prop:B-RD71-common-high} applies.

Finally let \(b_i=pk_i\), so \(B_i=qk_i\).
If the high lines agree, apply the same proposition.
If not, and one \(B_i>e\), quotient that model by its high
line and the other by the corresponding stable low line.
The inherited actual \(p^2\)-comparison has parameters
\(s-B_i\), prime to \(p\), and \(B_j/p=pk_j\),
divisible by \(p\). This is excluded by
Lemma~\ref{lem:B-RD69-rules}.
The equality \(B_i=e\) is impossible because \(p\nmid e\).
If both \(B_i<e<2q\), then \(B_1=B_2=q\).

For completeness the exchanged-line obstruction does not
depend on the particular index used in B70.
The residual high and low eigenvalue exponents in this
last case are \(c_i-q\) and \(pc_i+q\).
Semisimplicity of tame inertia and noncorresponding high
lines imply
\[
 c_1-q\equiv pc_2+q,\qquad
 c_2-q\equiv pc_1+q\pmod N.
 \tag{B-RD71.16}
\]
These equalities remain valid for scalar action as well.
Multiplying the first by \(p\) and adding the second would
give \(p-1\mid2q\), or \(p-1\mid2\), impossible for \(p\ge5\).
Hence the high lines must correspond after all.
This finishes the last case.
\end{proof}

\begin{theorem}[A larger exact tame sector in full height-two monodromy]
\label{thm:B-RD71}
\status{Proved; complete independent root and companion proof and exact-code audits passed.}
For every prime \(p\ge5\), in the full B.4 monodromy setting
with height two, perfect residue field and \(\mu_p\subset K\),
\[
 \boxed{\displaystyle
 \nu_{\min}\bigl(p,t(p-1),2\bigr)=2
 \quad(2\le t\le2p+1,\ p\nmid t).}
 \tag{B-RD71.17}
\]
For \(p=5\), the additionally settled prime-to-\(5\) indices
are \(e=32,36,44\).
\end{theorem}
\begin{proof}
After completed maximal unramified base change, compare
actual \(p^3\)-torsion. Since \(p\nmid e\), its cyclotomic
image has order \(p^2\); the determinant separates the
total dimensions \(0,1,2\) and their finite determinant
characters exactly as in Theorem~\ref{thm:B-RD69}.
Dimensions zero and two are detected residually and by
Cartier duality. Two systems filtered by potentially good
height-one pieces are also detected residually.
Mixed whole-supersingular/filtered comparisons are excluded
by Proposition~\ref{prop:B-RD63-mixed}, since \(e<2q<p^3\).

For two whole potentially supersingular groups,
Proposition~\ref{prop:B-RD69-one-slope} excludes mixed
Newton polygons and recovers two one-slope characters,
without an upper restriction on \(e\).
For two unequal polygons, residual wild inertia makes
their high lines the common unique fixed line, so
Proposition~\ref{prop:B-RD71-common-high} applies.
Residual tame inertia is covered by
Proposition~\ref{prop:B-RD71-tame}.
Thus actual \(p^3\)-torsion always determines the character.
For every index in (B-RD71.11), \(e\ge2(p-1)>p\).
Theorem~\ref{thm:B-RD63-arc}, at \(n=2\), supplies an actual
\(p^2\)-torsion comparison with different characters.
The manuscript index is therefore exactly two.
\end{proof}

\begin{remark}[Scalar quotients and the next arithmetic boundary]
\label{rem:B-RD71-scope}
\status{Proved; complete independent root and companion proof and exact-code audits passed.}
The new point-field argument deliberately permits scalar
residual actions after the first high quotient.
Its decisive inputs are the actual high-target inverse
valuation, its prime-to-\(p\) numerator, the common degree-\(q\)
point field, and the exact lower-break inequality.
It does not identify a high line from a scalar residual
representation alone.
The second quotient only preserves actual \(p\)-torsion;
the separate first-level lemma supplies precisely the
recovery statement needed at that level.

At the next permitted index \(e=2N=2(p^2-1)\), the small
interval can contain \(b_1=1\), \(b_2=2p-1\).
Their tangent exponents may satisfy the common-norm and
common-high-line congruences without being Frobenius conjugate,
and \(q b_1=q<e\) prevents use of
Proposition~\ref{prop:B-RD71-separated-high}.
This is an unexcluded arithmetic pattern, not a formal-group
existence assertion. No absolute index divisible by \(p\)
is treated in this subsection.
The retained exact supplement
\path{code/reduction_descent/verify_marked_quotient_sector.py}
reconstructs the inverse Newton polygons and the exact
point-field normalizations with rational arithmetic.
Across \(132\) sectors at \(p=5,7,11,13,17,19\), it checks
\(459{,}096\) first-level prime-parameter breaks,
\(29{,}376\) marked quotient polygons,
\(457{,}344\) low-line character weights, and
\(15{,}945\) separated-high point-field records, including
all six endpoint scalar pairs in these primes.
It also checks \(578{,}094\) tame mixed point-field inequalities
and the complete stated common-norm and high/low divisibility
case split. The author run passed in
\(1.3819726219953736\) seconds at 21:59:24 UTC.
The program, full output, scalar normalization samples and
pre-certificate proof are retained. No executed check failed.
The next arithmetic patterns are retained with their limited
scope; no formal-group existence conclusion is inferred.
These finite checks supplement the complete all-prime proofs.
The complete independent root and companion audits passed.
Their author-output-preserving exact reruns took
\(1.4025961040024413\) and \(1.4007752400066238\) seconds,
respectively. The exact audited candidate and all full audit
records are preserved.
\end{remark}


% Proved after full root and companion audits; kept outside frozen CP26 paper.
\subsection{A first-level mixed break and the doubled tame boundary}

For the local assertions let \(p\ge5\), \(q=p^2\), \(N=q-1\).
The field \(K\) is complete discretely valued of mixed characteristic
\((0,p)\), with algebraically closed residue field, \(\mu_p\subset K\),
and absolute ramification \(e\) prime to \(p\). Put \(s=e(p+1)\),
and use the actual tame good extension \(L/K\) of degree \(N\).
Thus \(v_L(p)=eN=(p-1)s\) and \(N\mid s\).
For every two-slope group, \(B\) denotes its integral parameter
\[
 0<B<pe,\qquad \alpha=s-B,\qquad \beta=B/p,\qquad
 B\equiv c\pmod{p+1},\qquad \lambda_H=\zeta^{c-B}.
 \tag{B-RD72.1}
\]
These are parameters of actual good formal groups and actual
finite-flat quotients, with the semilinear generic descent of
Lemma~\ref{lem:B-RD69-quotients}.

\begin{lemma}[A mixed-divisibility inequality already at the first level]
\label{lem:B-RD72-first-mixed}
\status{Proved; complete independent root and companion proof and exact-code audits passed.}
Suppose two two-slope groups have isomorphic actual \(p\)-torsion,
with parameters \(C_{\mathrm{pr}}\) and \(C_{\mathrm{div}}\), where
\(p\nmid C_{\mathrm{pr}}\) and \(p\mid C_{\mathrm{div}}\).
Then necessarily
\[
                 C_{\mathrm{pr}}>C_{\mathrm{div}}.
 \tag{B-RD72.2}
\]
\end{lemma}
\begin{proof}
The first-level division field \(R/L\) of the prime-parameter group
has degree \(p\). Its high line and determinant are trivial over
\(L\), so its residual action is upper unipotent of order at most
\(p\); the nonintegral low valuation \(C_{\mathrm{pr}}/p\) makes
the action nontrivial. The actual torsion isomorphism identifies
the division fields and their unique wild fixed lines.

Fix a nonidentity \(g\in\operatorname{Gal}(R/L)\) and corresponding
low points. In the normalized valuation of \(R\), their coordinate
valuations are \(C_{\mathrm{pr}}\) and \(C_{\mathrm{div}}\), and
their nonzero displacements are high points, of valuations
\(p(s-C_{\mathrm{pr}})\) and \(p(s-C_{\mathrm{div}})\).
Formal subtraction has the same valuation as coordinate subtraction.
Writing the actual lower break as
\(r=v_R(g\varpi-\varpi)-1\), the equality and strict-inequality
parts of Lemma~\ref{lem:B-RD60-break} give
\[
 \begin{aligned}
 r&=ps-(p+1)C_{\mathrm{pr}},\\
 r&<ps-(p+1)C_{\mathrm{div}}.
 \end{aligned}
 \tag{B-RD72.3}
\]
Subtracting proves (B-RD72.2). This uses only actual first-level
torsion. No second-level comparison is imported into this lemma.
\end{proof}

\begin{proposition}[A necessary inequality for mixed second quotients]
\label{prop:B-RD72-mixed-quotients}
\status{Proved; complete independent root and companion proof and exact-code audits passed.}
Suppose an actual \(p^3\)-torsion isomorphism identifies the high
lines of two two-slope groups with
\[
 B_i=pb_i<e,\qquad qb_1<e<qb_2,\qquad p-1\nmid b_1.
 \tag{B-RD72.4}
\]
Then it is necessary that
\[
                    e(p+1)>q(pb_1+b_2).
 \tag{B-RD72.5}
\]
\end{proposition}
\begin{proof}
The actual comparison identifies the finite determinant and the
original high-line characters. Consequently
\(b_1\equiv b_2\pmod{p-1}\), as in (B-RD71.12).
Take the actual quotients \(M_i'=M_i/H_i\).
Lemma~\ref{lem:B-RD71-marked-low} gives their parameters
\(D_i=q b_i\), an actual \(p^2\)-torsion comparison, and
corresponding marked stable low lines
\[
 J_i=\operatorname{im}\bigl(M_i[p]\longrightarrow M_i'[p]\bigr).
 \tag{B-RD72.6}
\]
Their residual actions are tame, and their first-level modules are
rational over \(L\). The quotient low/high eigenvalue ratio is
\(\zeta^{(p+1)B_i}\). It is nontrivial, since
\(p-1\nmid b_i\) and \(B_i=pb_i\).
Thus each residual representation has two distinct eigenspaces.
Preservation of the marked low eigenspace forces preservation of
the other eigenspace, which is precisely its high line \(H_i'\).
This argument uses the actual marked line; it does not choose a
complement in a scalar representation.

Now quotient both \(M_i'\) by the corresponding high lines. The
inherited comparison is exactly at the first level:
\[
 (M_i'/H_i')[p]
   =\{Z\in M_i'[p^2]:pZ\in H_i'\}/H_i'.
 \tag{B-RD72.7}
\]
Since \(D_1<e<D_2\), Lemma~\ref{lem:B-RD69-quotients} gives
their respective parameters
\[
 C_1=pD_1=p^3b_1,
 \qquad C_2=s-D_2=s-qb_2.
 \tag{B-RD72.8}
\]
Both are in \((0,pe)\): the first because \(D_1<e\), and the
second because \(e<D_2<pe\). Moreover \(p\mid C_1\) and
\(p\nmid C_2\), since \(p\nmid s\).
Lemma~\ref{lem:B-RD72-first-mixed} therefore requires
\(s-qb_2>p^3b_1\). This is exactly (B-RD72.5).
All quotient identities concern actual generic torsion with its
inertia action; none asserts descent of the integral good model
through a ramified quotient.
\end{proof}

\begin{proposition}[Common-high recovery at the doubled boundary]
\label{prop:B-RD72-common-high}
\status{Proved; complete independent root and companion proof and exact-code audits passed.}
Set \(e=2N=2(p^2-1)\). An actual \(p^3\)-torsion isomorphism
between two two-slope groups which identifies their high lines
determines their finite special inertia characters.
\end{proposition}
\begin{proof}
The common determinant and high character give
\[
 c_1\equiv c_2\pmod{p-1},\qquad
 c_1-B_1\equiv c_2-B_2\pmod N.
 \tag{B-RD72.9}
\]
If either \(B_i\) is prime to \(p\),
Lemma~\ref{lem:B-RD69-rules} recovers the characters or excludes
the mixed-divisibility comparison. This includes \(B_i=e\).
Otherwise both are divisible by \(p\), so neither equals \(e\).
If both exceed \(e\), their actual high quotients have prime
parameters \(s-B_i\); their inherited \(p^2\)-comparison recovers
the characters. If they lie on opposite sides of \(e\), these
quotients have opposite parameter divisibility, excluded by the
same comparison lemma.

It remains that \(B_i=pb_i<e\). Here
\[
 1\le b_i\le2p-1,
 \qquad b_1\equiv b_2\pmod{p-1}.
 \tag{B-RD72.10}
\]
Equality of the two integers identifies \(c_1=c_2\pmod N\).
If they differ, reorder so that \(b_1<b_2\).
When \(b_1\ge2\), one has \(qb_1\ge2q>e\), and
Proposition~\ref{prop:B-RD71-separated-high} excludes this
comparison. Thus \(b_1=1\), leaving only \(b_2=p\) or
\(b_2=2p-1\).
The pair \(1,p\) exchanges the two Frobenius-conjugate special
eigenvalues by (B-RD71.14), so gives the same character.

For \((b_1,b_2)=(1,2p-1)\), all hypotheses of
Proposition~\ref{prop:B-RD72-mixed-quotients} hold:
\(p-1\nmid1\) and \(q<2(q-1)<q(2p-1)\).
The second actual quotients have
\[
 C_1=p^3,
 \qquad C_2=3p^2-2p-2,
 \qquad C_1-C_2=p(p-1)(p-2)+2>0.
 \tag{B-RD72.11}
\]
This contradicts the necessary inequality \(C_2>C_1\).
The additional arithmetic pattern is therefore excluded by an
actual first-level break comparison, completing the proof.
\end{proof}

\begin{proposition}[All whole formal comparisons at the boundary]
\label{prop:B-RD72-formal}
\status{Proved; complete independent root and companion proof and exact-code audits passed.}
At \(e=2N\), actual \(p^3\)-torsion determines the finite special
inertia character among whole potentially supersingular
height-two groups, with arbitrary finite determinant.
\end{proposition}
\begin{proof}
The one-slope and mixed-polygon cases are covered by
Proposition~\ref{prop:B-RD69-one-slope}, with no upper restriction
on \(e\). For two two-slope groups, residual wild inertia makes
their high line the unique fixed line, so the preceding
proposition applies.

For completeness consider every residual tame case at this new
endpoint. Write \(B_i=pb_i\), \(0<b_i<e\), since the first-level
modules are rational over \(L\). If both \(b_i\) are prime to
\(p\), Lemma~\ref{lem:B-RD69-rules} recovers the characters.
If exactly one is prime, say \(b_1=b\), \(b_2=pk\), choose
corresponding primitive points whose \(p\)-multiples lie outside
both high lines. Their common point field has degree \(q\)
over \(L\), forced by the first valuation \(b/q\), and all
first-level translations occur. Its normalized coordinate
valuations are \(b\) and \(pk\).
Distinct high lines would allow a translation high in the first
model and low in the second. The exact and lower-bound parts
of Lemma~\ref{lem:B-RD60-break} would give
\[
 q(\alpha_1-pk)\le b-pk.
 \tag{B-RD72.12}
\]
But \(\alpha_1=s-pb\ge e+p\), \(pk<e\), and
\(e=2q-2<p^3\). The left side is greater than \(p^3\),
whereas the right side is less than \(p^3\). Thus the high
lines correspond, and the common-high proposition applies.

Finally, if both \(b_i=pk_i\), then \(B_i=qk_i\).
Assume their high lines do not correspond. If one \(B_i>e\),
quotient that model by its high line and the other by the
corresponding stable low line. The actual \(p^2\)-comparison
has parameters \(s-B_i\) and \(B_j/p=pk_j\), of opposite
divisibility, again impossible.
Neither \(B_i=e\), since \(p\nmid e\).
If both \(B_i<e=2q-2\), they must both equal \(q\).
Lemma~\ref{lem:B-RD70-exchange} would then require
\(p-1\mid2q\), equivalently \(p-1\mid2\), impossible.
This exhausts all tame cases as well.
\end{proof}

\begin{theorem}[Exact index at twice the full tame degree]
\label{thm:B-RD72}
\status{Proved; complete independent root and companion proof and exact-code audits passed.}
For every prime \(p\ge5\), in the full B.4 monodromy setting
with height two, perfect residue field and \(\mu_p\subset K\),
\[
 \boxed{\nu_{\min}\bigl(p,2(p^2-1),2\bigr)=2.}
 \tag{B-RD72.13}
\]
In particular \(\nu_{\min}(5,48,2)=2\).
Combined with Theorem~\ref{thm:B-RD71}, the exact formula is
\[
 \nu_{\min}\bigl(p,t(p-1),2\bigr)=2
 \quad\text{for }2\le t\le2p+2,\quad p\nmid t.
 \tag{B-RD72.14}
\]
\end{theorem}
\begin{proof}
Pass to the completed maximal unramified base to prove the
upper bound. Since \(p\nmid e\), the actual cyclotomic character
modulo \(p^3\) has order \(p^2\), separating the total dimensions
\(0,1,2\) and the finite determinant as in
Theorem~\ref{thm:B-RD69}. Dimensions zero and two, and comparisons
of two systems filtered by height-one potentially good pieces,
are determined residually, with Cartier duality for dimension two.
The mixed whole-supersingular/filtered comparison is excluded by
Proposition~\ref{prop:B-RD63-mixed}, because
\(e=2q-2<p^3\). Two whole groups are handled by
Proposition~\ref{prop:B-RD72-formal}. Thus actual \(p^3\)-torsion
always determines the requested character in the full setting.
Because \(e>p\), Theorem~\ref{thm:B-RD63-arc} at \(n=2\)
supplies permitted systems with equal actual \(p^2\)-torsion and
different characters. The manuscript's least index is therefore
two. The final range adds the single endpoint \(t=2p+2\) to B71.
\end{proof}

\begin{remark}[Precise scope of the first-level obstruction]
\label{rem:B-RD72-scope}
\status{Proved; complete independent root and companion proof and exact-code audits passed.}
The new strict inequality is a first-level necessary condition,
not a first-level character-recovery theorem for arbitrary
mixed-divisibility parameters. The second quotient in
(B-RD72.7) leaves only actual \(p\)-torsion; (B-RD72.2) is
precisely the assertion used there.
The intermediate residual eigenvalues are distinct in the new
boundary pair, so their preserved low lines determine the high
lines. Scalar intermediate actions are not treated by that
argument; the separate scalar-safe B71 proposition is used in
the cases where it applies.
No absolute ramification divisible by \(p\) is settled here.
The next prime-to-\(5\) index \(e=52\), and the wild index
\(e=20\), remain outside this theorem.
The retained exact supplement
\path{code/reduction_descent/verify_first_level_mixed_boundary.py}
reconstructs both actual high-quotient Newton polygons from the
universal vertices and counts quotient roots using the product map.
Across twelve primes from \(5\) through \(43\), it checks
\(24\) double-quotient root records, \(264\) corresponding
distinct-eigenvalue records, \(545{,}376\) first-level prime breaks,
and \(1{,}074{,}792\) tame mixed point-field inequalities.
It also checks the entire stated small common-norm case split,
the boundary strict-break contradiction at every tested prime,
and all remaining tame quotient-divisibility branches.
For \(p=5,e=48\), the actual second parameters are \(125\) and
\(63\); the prime model requires lower break \(1062\), whereas
the divisible model requires it to be strictly less than \(690\).
The author run passed in \(0.5185015480092261\) seconds at
22:14:11 UTC. The complete program, output, and pre-certificate
proof are retained. No executed numerical check failed.
These finite checks supplement the all-prime proof and do not
provide a formal-group existence test.
The complete independent root and companion audits passed.
Their author-output-preserving exact reruns took
\(0.5013780940062134\) and \(0.5564822399901459\) seconds,
respectively. The exact audited candidate and complete audit
records are preserved.
\end{remark}


% Proved after full independent root and companion audits; exact candidate preserved.
\subsection{Rational first-level torsion and the exact index-two range}

Throughout the local assertions, \(p\ge5\), \(q=p^2\), \(N=q-1\),
and \(K\) is complete discretely valued of mixed characteristic
\((0,p)\), with algebraically closed residue field and \(\mu_p\subset K\).
Its absolute ramification \(e\) is prime to \(p\). Set \(s=e(p+1)\).
We use the actual good tame extension \(L/K\) of degree \(N\),
so that \(v_L(p)=eN=(p-1)s\) and \(N\mid s\).
A two-slope formal group has parameter \(D\) and valuations
\[
 0<D<pe,\qquad \alpha=s-D,\qquad\beta=D/p,
 \qquad \alpha>\beta.
 \tag{B-RD73.1}
\]
Its canonical high line has \(p-1\) nonzero points of valuation
\(\alpha\), and all other nonzero first-level points have valuation
\(\beta\). We always use actual integral good formal models,
actual generic torsion, and their actual inertia action.

\begin{lemma}[Rationality and the small high inverse branch]
\label{lem:B-RD73-high-branch}
\status{Proved; complete independent root and companion proof and exact-code audits passed.}
Let \(M\) be such a two-slope group. If its residual inertia action
is tame, then \(M[p]\) is rational over \(L\).
This is also true when \(M[p]\) has two distinct inertia-stable
\(\mathbf F_p\)-lines, without any semisimplicity assumption made
in advance. Conversely rationality over \(L\) implies \(p\mid D\).

Suppose \(M[p]\) is \(L\)-rational and \(D<e\).
For any nonzero high point \(Q\), the inverse equation \([p]X=Q\)
has exactly \(p\) roots of valuation
\[
 v_L(X)=\frac{C}{p},\qquad C=s-pD,
 \tag{B-RD73.2}
\]
and \(q-p\) roots of valuation \(D/p\).
The \(p\) higher-valuation roots form one coset of the high line.
Each such root has a Galois point field \(E/L\) of degree \(p\),
whose translation group is exactly that high line.
For every nonidentity \(g\in\operatorname{Gal}(E/L)\), its actual
lower break is
\[
 b(g)=eN.
 \tag{B-RD73.3}
\]
\end{lemma}
\begin{proof}
The high line is inertia stable. If residual inertia is tame,
its finite cyclic image is semisimple. Thus this stable
\(\mathbf F_p\)-line has a stable complementary line, and the
two characters take values in \(\mathbf F_p^\times\).
The tame extension of degree \(N\), a multiple of \(p-1\),
kills both characters. If two distinct stable lines are given
initially, the representation is already their direct sum,
with the same conclusion. In particular wild inertia acts
trivially: a \(p\)-group has no nontrivial character into
\(\mathbf F_p^\times\). These arguments use the algebraically
closed residue field to identify the full absolute Galois
action with inertia. They prove actual \(L\)-rationality,
not merely rationality of a semisimplification.
If all first-level points are \(L\)-rational, their low coordinate
valuation \(D/p\) is integral, so \(p\mid D\).

For the inverse equation, the exact universal coefficient
bounds used in Lemmas~\ref{lem:B-RD69-quotients}
and~\ref{lem:B-RD71-marked-low} give lower Newton vertices
\[
 (0,s-D),\qquad (p,(p-1)D),\qquad(q,0).
 \tag{B-RD73.4}
\]
When \(D<e\), the middle vertex is strictly below the chord:
this is equivalent to \((p+1)D<s\).
The initial segment gives \(p\) roots of valuation \((s-pD)/p\),
and the final segment gives \(q-p\) roots of valuation \(D/p\).
The potential degree-one point \((1,eN)\) and the remaining
coefficients lie above this polygon by those same bounds.
Here
\[
 C>D>0,\qquad p\nmid C,
 \qquad \alpha>C/p>D/p.
 \tag{B-RD73.5}
\]
The middle assertion follows from \(p\nmid s\).

Translation by any high first-level point preserves the larger
inverse valuation, because that point has valuation \(\alpha>C/p\).
Translation by a low point changes the valuation to \(D/p<C/p\).
For unequal positive parameter valuations, an integral formal
group law has valuation equal to their minimum; this follows
from its linear terms and integral terms of degree at least two.
Thus the larger inverse branch is exactly one coset of the
high line, containing \(p\) points.

The target \(Q\) and all first-level translations are rational
over \(L\). Every conjugate of a chosen larger-valuation root
remains in the same branch and is a translation of it by a high
point. In particular it belongs to \(L(X)\), so that field is
Galois of degree at most \(p\). The denominator \(p\) in
\(v_L(X)=C/p\), with \(p\nmid C\), forces ramification degree
at least \(p\). Hence \([E:L]=e(E/L)=p\), and all high-line
translations occur. Normalize \(v_E\) to be integral.
Then \(v_E(X)=C\), and every nonidentity displacement is a
high point of valuation \(p(s-D)\).
The exact prime-valuation part of Lemma~\ref{lem:B-RD60-break}
gives
\[
 b(g)=p(s-D)-C=(p-1)s=eN,
\]
as asserted. No upper/lower numbering conversion is used.
\end{proof}

\begin{proposition}[The high line is recovered from the next level]
\label{prop:B-RD73-high-recovery}
\status{Proved; complete independent root and companion proof and exact-code audits passed.}
Let \(M_1,M_2\) be two two-slope groups whose actual
\(p\)-torsion modules are \(L\)-rational. Every isomorphism of
their actual \(p^2\)-torsion identifies the high lines.
There is no upper bound on \(e\) in this assertion.
\end{proposition}
\begin{proof}
Let their parameters be \(D_1,D_2\).
They are divisible by \(p\) by the preceding lemma, so neither
equals \(e\). Assume the actual isomorphism \(\phi\) does not
identify the high lines.

Suppose first one parameter is below \(e\); after exchanging
the models and replacing \(\phi\) by its inverse if necessary,
take \(D_1<e\). Choose a larger-valuation inverse \(X_1\) of a
nonzero high point in the preceding lemma, and put \(X_2=\phi(X_1)\).
Its multiple \(pX_2\) is low in the second group. The low-target
inverse polygon gives
\[
 v_L(X_1)=\frac{s-pD_1}{p},\qquad
 v_L(X_2)=\frac{D_2}{p^3}.
 \tag{B-RD73.6}
\]
The actual torsion isomorphism identifies stabilizers, so their
point fields are the same \(E/L\), of degree \(p\).
Consequently
\[
 v_E(X_1)=s-pD_1,\qquad v_E(X_2)=D_2/q.
 \tag{B-RD73.7}
\]
If the second value is not integral, this already contradicts
\(X_2\in E\). Otherwise choose any nonidentity \(g\) in the
common degree-\(p\) group. It translates \(X_1\) by a high
point and \(X_2\) by a low point, because the first translation
group is exactly its high line and \(\phi\) sends that line
to a line distinct from the second high line.
The second displacement has valuation \(D_2\) in \(E\),
whereas the actual lower break is \(eN\). The unconditional
lower bound of Lemma~\ref{lem:B-RD60-break} therefore requires
\[
 D_2\ge D_2/q+eN,
 \qquad\text{equivalently }D_2\ge qe.
 \tag{B-RD73.8}
\]
This is impossible because \(D_2<pe<qe\).
The argument does not use distinct tame eigenvalues or a chosen
complement in a scalar residual representation.

It remains that both \(D_i>e\). Reorder so \(D_1\le D_2\).
Put \(A_1=s-D_1\); it is prime to \(p\).
For a high target in the first model the inverse polygon now
has one segment with \(q\) roots of valuation \(A_1/q\).
Choose one such root \(X_1\) and again set \(X_2=\phi(X_1)\).
Its target is low, hence \(v_L(X_2)=D_2/p^3\).
Both first-level modules are \(L\)-rational, so each point
field is Galois with translation group of order at most \(q\).
The prime-to-\(p\) numerator \(A_1\) forces the common field
to have degree \(q\); all first-level translations occur.
Its coordinate valuations are \(A_1\) and \(D_2/p\).
Choose a translation high in the first model, necessarily low
in the second. Exactly as in the degree-\(q\) comparison in
Proposition~\ref{prop:B-RD71-separated-high}, the exact break
and lower bound imply
\[
 q(D_2/p)\ge D_2/p+N A_1,
 \qquad D_2/p\ge A_1.
 \tag{B-RD73.9}
\]
But
\[
 A_1=s-D_1\ge s-D_2>D_2/p,
\]
the strict inequality being equivalent to \(D_2<pe\).
This is the final contradiction. The two cases exhaust the
possible parameters, proving the proposition.
\end{proof}

\begin{lemma}[Common-high recovery below the next cyclic threshold]
\label{lem:B-RD73-common-high}
\status{Proved; complete independent root and companion proof and exact-code audits passed.}
Assume \(e<p^3\). Two two-slope groups with isomorphic actual
\(p^3\)-torsion and corresponding high lines have the same
finite special inertia character.
\end{lemma}
\begin{proof}
Their common determinant and high-line character give
\[
 c_1\equiv c_2\pmod{p-1},\qquad
 c_1-B_1\equiv c_2-B_2\pmod N.
 \tag{B-RD73.10}
\]
Any prime-to-\(p\) parameter is settled by
Lemma~\ref{lem:B-RD69-rules}; this includes \(B_i=e\).
Assume both parameters are divisible by \(p\).
If both exceed \(e\), the actual high quotients have prime
parameters \(s-B_i\), and their inherited \(p^2\)-comparison
recovers equality. If the parameters lie on opposite sides of
\(e\), the high quotients have opposite divisibility, forbidden
by the same second-level rule.

It remains that \(B_i=pb_i<e\). Let \(M_i'=M_i/H_i\).
Their parameters are \(D_i=q b_i\), and the actual
\(p^3\)-comparison induces an actual \(p^2\)-comparison of
these quotients. By Lemma~\ref{lem:B-RD71-marked-low}, each
\(M_i'[p]\) has an intrinsic stable low line
\[
 J_i=\operatorname{im}(M_i[p]\longrightarrow M_i'[p]),
\]
distinct from its canonical high line. The rationality part of
Lemma~\ref{lem:B-RD73-high-branch} therefore shows that both
first-level modules are \(L\)-rational.
Proposition~\ref{prop:B-RD73-high-recovery} identifies their
high lines \(H_i'\), including when the residual representation
is scalar.

Take these second actual quotients. The surviving comparison
is exactly at level \(p\):
\[
 (M_i'/H_i')[p]
   =\{Z\in M_i'[p^2]:pZ\in H_i'\}/H_i'.
 \tag{B-RD73.11}
\]
Equivalently, in the original groups it is the quotient of
the corresponding subgroup of \(M_i[p^3]\) by the full
preimage of \(H_i'\), whose order is \(p^2\).
There is no unasserted second-level comparison after this step.
The second parameters are
\[
 C_i=pD_i\quad(D_i<e),\qquad
 C_i=s-D_i\quad(D_i>e).
 \tag{B-RD73.12}
\]
The equality \(D_i=e\) is impossible because \(p\mid D_i\)
and \(p\nmid e\).

If both \(D_i>e\), both second parameters are prime to \(p\).
The first-level recovery in Lemma~\ref{lem:B-RD71-first-level}
then gives \(D_1=D_2\), hence \(B_1=B_2\).
If \(D_1<e<D_2\), the first second-quotient module has two
distinct stable lines: its canonical high line and the
intrinsic low line
\[
 \operatorname{im}\bigl(M_1'[p]\longrightarrow(M_1'/H_1')[p]\bigr).
 \tag{B-RD73.13}
\]
This is another direct application of
Lemma~\ref{lem:B-RD71-marked-low}, now to the actual group
\(M_1'\) with parameter \(D_1<e\).
Its residual wild inertia is trivial. In contrast, the second
quotient has prime parameter \(s-D_2\); its low valuation
\((s-D_2)/p\) is not integral over \(L\), so its residual
division field is wild of degree \(p\). This contradicts the
actual first-level isomorphism (B-RD73.11).
The other mixed ordering is identical.

Finally, if both \(D_i<e\), then
\[
 0<b_i<e/q<p,
 \qquad b_i\in\{1,\ldots,p-1\}.
 \tag{B-RD73.14}
\]
Equation (B-RD73.10) gives \(b_1\equiv b_2\pmod{p-1}\).
Reduction modulo \(p-1\) is injective on this precise interval,
so \(b_1=b_2\). Thus every surviving branch has \(B_1=B_2\),
and the common original high character identifies \(c_1=c_2\).
This proves the assertion without a determinant square-class
normalization or a nonscalar residual assumption.
\end{proof}

\begin{proposition}[Whole formal character recovery below \(p^3\)]
\label{prop:B-RD73-formal}
\status{Proved; complete independent root and companion proof and exact-code audits passed.}
If \(e<p^3\), actual \(p^3\)-torsion determines the finite
special inertia character among whole potentially supersingular
height-two groups.
\end{proposition}
\begin{proof}
Proposition~\ref{prop:B-RD69-one-slope} handles two one-slope
groups and excludes mixed one-slope/two-slope comparisons,
without an upper restriction on \(e\).
For two two-slope groups with nontrivial residual wild inertia,
their high line is the common unique wild fixed line.
If residual inertia is tame, the first-level rationality lemma
gives \(M_i[p](L)=M_i[p]\), and
Proposition~\ref{prop:B-RD73-high-recovery} recovers their high
lines from the actual \(p^2\)-isomorphism already available.
In either case Lemma~\ref{lem:B-RD73-common-high} applies.
These are all residual possibilities.
\end{proof}

\begin{theorem}[The exact index-two range at prime-to-\(p\) ramification]
\label{thm:B-RD73}
\status{Proved; complete independent root and companion proof and exact-code audits passed.}
In the full B.4 monodromy setting of height two, with perfect
residue field, \(p\ge5\), \(\mu_p\subset K\), and \(p\nmid e\),
\[
 \boxed{\nu_{\min}(p,e,2)=2\quad\Longleftrightarrow\quad p<e<p^3.}
 \tag{B-RD73.15}
\]
Equivalently the range where the least index is two is
\[
 e=t(p-1),\qquad 2\le t\le p^2+p+1,
 \qquad p\nmid t.
 \tag{B-RD73.16}
\]
For \(p=5\) this includes every multiple of four from \(8\)
through \(124\) which is prime to five. At the next permitted
prime-to-five index \(e=128\), the least index is at least three.
\end{theorem}
\begin{proof}
For \(p<e<p^3\), pass to the completed maximal unramified base.
The cyclotomic character modulo \(p^3\) has order \(p^2\),
because \(p\nmid e\). The determinant argument of
Theorem~\ref{thm:B-RD69} identifies total dimensions \(0,1,2\)
and the finite determinant characters. Dimensions zero and two
are determined residually, with Cartier duality for the latter;
two systems filtered by potentially good height-one pieces are
also determined residually. A mixed whole-supersingular/filtered
comparison is impossible by Proposition~\ref{prop:B-RD63-mixed},
because \(e<p^3\). The two whole-group case is precisely
Proposition~\ref{prop:B-RD73-formal}. Hence actual \(p^3\)-torsion
determines the requested character in the full monodromy setting.
Theorem~\ref{thm:B-RD63-arc} at \(n=2\) gives unequal-character
systems with isomorphic actual \(p^2\)-torsion whenever \(e>p\).
Thus the least manuscript index is exactly two.

For completeness, \(\mu_p\subset K\) forces \(p-1\mid e\).
The only positive possibility with \(e\le p\) is \(e=p-1\),
whose exact index is one by Theorem~\ref{thm:B-RD61}.
There is no permitted equality \(e=p^3\), since \(p\nmid e\).
For \(e>p^3\), the same actual lower construction B63 at
\(n=3\) gives systems with equal actual \(p^3\)-torsion and
different characters, so the least index is at least three.
This proves the reverse implication as well.
Finally \((p-1)(p^2+p+1)=p^3-1\), giving the equivalent range
and the stated numerical examples. No exact index at \(e>p^3\)
is inferred.
\end{proof}

\begin{remark}[Scope of the rational-branch argument]
\label{rem:B-RD73-scope}
\status{Proved; complete independent root and companion proof and exact-code audits passed.}
The high-line recovery proposition is independent of an upper bound
on \(e\), but assumes actual \(L\)-rational first-level modules.
That hypothesis is proved separately where it is used: from tame
residual inertia, or from the intrinsic stable low line of a small
actual high quotient. Scalar residual representations cause no
exception, because the inverse-branch point fields recover the
high line without selecting a complementary eigenspace.
The common-high proof uses \(e<p^3\) only in the final integer
interval after two quotients; the full-monodromy proof also uses
the exact mixed whole/filtered threshold.
The comparison left after the second quotient is actual
\(p\)-torsion, and every subsequent argument is at that level.
No statement for absolute ramification divisible by \(p\),
no exact least index above the displayed range, and no higher-height
classification is asserted.

The retained exact supplement is
\path{code/reduction_descent/verify_rational_high_line.py}.
Its author run passed in \(9.558507562003797\) seconds at
22:30:37 UTC on 19 November 2026. For \(p=5,7,11,13\), it
checks all \(360\) permitted sectors in the displayed range,
\(275{,}760\) rational first-level polygons, the
\(22{,}932\) degree-\(p\) and \(252{,}828\) degree-\(q\)
high inverse branches, and \(22{,}932\) actual second-high-quotient
root-count records. Exact rational arithmetic reconstructs the
Newton hulls and checks the displacement inequalities and all
remaining parameter branches. Another \(48\) records above
\(p^3\) test only the unbounded high-line recovery proposition;
they do not extend the exact-index theorem. This finite supplement
does not replace the all-prime proof or establish existence of
formal groups with any enumerated parameter.

A fresh targeted search using two engines was retained with its
actual queries and read scopes. It located Hattori's finite-flat
ramification-character work and Hoshi's tame-blind extension
bound. Only the explicitly logged introduction, theorem-statement
and short-corollary passages were newly read; no unexamined main
proof is claimed as audited, no new result is imported from a
snippet, and no priority conclusion is inferred. The proof here
uses the already audited Newton, actual-quotient, lower-break and
mixed-threshold dependencies indicated above.
The independent root and companion audits passed for all six
assertions and the complete program. Their author-output-preserving
reruns took \(9.449625776003813\) and \(9.306663159004529\)
seconds, respectively. Exact prepromotion bytes, audit scopes,
the superseded weaker outline, and all executed code/output
remain retained.
\end{remark}


% Proved after full independent root and companion audits; exact candidate preserved.
\subsection{Iterated actual high quotients and a third-index interval}

For the local assertions let \(K\) be complete discretely valued
of mixed characteristic \((0,p)\), with algebraically closed
residue field, \(p\ge5\), \(\mu_p\subset K\), and absolute
ramification \(e\) prime to \(p\). Set
\[
 q=p^2,\qquad N=q-1,\qquad s=e(p+1).
 \tag{B-RD74.1}
\]
Use the actual common good tame field \(L/K\) of degree \(N\),
so \(v_L(p)=eN\). The two-slope parameter and high-line
character are those of Lemma~\ref{lem:B-RD69-rules}:
\[
 0<B<pe,\qquad \alpha=s-B,\qquad\beta=B/p,
 \qquad \lambda_H=\zeta^{c-B}.
 \tag{B-RD74.2}
\]
There is no square-class restriction on the finite determinant.

\begin{lemma}[Actual levels in a sequence of high quotients]
\label{lem:B-RD74-actual-levels}
\status{Proved; complete independent root and companion proof and exact-code audits passed.}
Let \(n\ge2\) and \(0\le j\le n-1\). Suppose actual
\(p^n\)-torsion modules of
\(M_1,M_2\) are identified, and at each of the first \(j\)
stages the chosen high lines are identified. Write
\(M_i^{(0)}=M_i\) and
\(M_i^{(j+1)}=M_i^{(j)}/H_i^{(j)}\).
After \(j\) quotients the induced comparison is actual
\(p^{n-j}\)-torsion. Each quotient is the actual good formal
quotient over \(\mathcal O_L\), with its inherited generic
semilinear descent and the same finite special character.

If a stage has divisible parameter \(D\), then \(D\ne e\).
For \(D<e\) the next group has parameter \(pD\) and its
first-level module has two distinct stable lines, hence is
\(L\)-rational. For \(D>e\) the next parameter \(s-D\) is
prime to \(p\), and its first-level division field over \(L\)
is wild of degree \(p\).
\end{lemma}
\begin{proof}
If the current level is \(p^r\), with \(r\ge2\), the exact
generic finite-group construction is
\[
 (M/H)[p^{r-1}]
  =\{X\in M[p^r]:p^{r-1}X\in H\}/H.
 \tag{B-RD74.3}
\]
Indeed a lift of a point killed by \(p^{r-1}\) has its
\(p^{r-1}\)-multiple in \(H\), hence is killed by \(p^r\).
Conversely this condition is sufficient, and the kernel is
exactly \(H\). Applying the actual torsion isomorphism to this
subgroup and quotient gives the asserted comparison. Iteration
loses one level at every stage, including at the final stage,
which retains only \(p\)-torsion.

The finite-flat closure of each stable generic line is stable
under the actual semilinear descent action. The good formal
quotient, quotient coordinate, and effective generic descent
are exactly those proved in
Lemma~\ref{lem:B-RD69-quotients}; its special isogeny preserves
the rational finite special character. The total kernel after
\(j\) stages has rank \(p^j\), is killed by \(p^j\), and its
preimage construction is contained in the original
\(M[p^n]\). Thus no comparison of higher truncation than
\(p^{n-j}\) is being asserted.

When \(p\mid D\), equality \(D=e\) is excluded by \(p\nmid e\).
The exact actual high-quotient parameter is \(pD\) for \(D<e\)
and \(s-D\) for \(D>e\). In the first case the intrinsic
image of \(M[p]\) in \((M/H)[p]\) is a stable low line,
distinct from its stable high line, by
Lemma~\ref{lem:B-RD71-marked-low}. Both characters are
\(\mathbf F_p^\times\)-valued, so the degree-\(N\) tame field
kills them. This proves actual first-level rationality.
In the second case \(p\nmid(s-D)\). The low valuation
\((s-D)/p\) is nonintegral over \(L\), whereas the high
character and determinant are trivial there. The residual
action is therefore the nontrivial upper-unipotent group of
order \(p\), as in Lemma~\ref{lem:B-RD71-first-level}.
This proves the stated tame/wild alternatives.
\end{proof}

\begin{proposition}[Uniform iteration bound for whole formal groups]
\label{prop:B-RD74-formal}
\status{Proved; complete independent root and companion proof and exact-code audits passed.}
Let \(n\ge3\) and \(e<p^n\). An isomorphism of actual
\(p^n\)-torsion determines the finite special inertia character
among whole potentially supersingular height-two groups.
\end{proposition}
\begin{proof}
The actual \(p^3\)-comparison already available handles both
one-slope polygons and excludes a mixed one-slope/two-slope
comparison by Proposition~\ref{prop:B-RD69-one-slope}, whose
scope has no upper bound on \(e\).
For two two-slope groups, nontrivial residual wild inertia
identifies the high line as its unique fixed line. If the
residual action is tame, Lemma~\ref{lem:B-RD73-high-branch}
gives actual \(L\)-rational first-level modules, and
Proposition~\ref{prop:B-RD73-high-recovery} identifies the
high lines from the available actual \(p^2\)-comparison.
Thus in every two-slope case the original high lines agree.

Write the original parameters \(B_i\) and tangent exponents
\(c_i\). Their common finite determinant and high character
give
\[
 c_1\equiv c_2\pmod{p-1},\qquad
 c_1-B_1\equiv c_2-B_2\pmod N.
 \tag{B-RD74.4}
\]
If at least one \(B_i\) is prime to \(p\), the general
comparison rules of Lemma~\ref{lem:B-RD69-rules}, using only
the available \(p^2\)-torsion, recover the character or
exclude opposite divisibility. In particular this settles
the possible original equality \(B_i=e\).
Assume henceforth that both \(B_i\) are divisible by \(p\).
If both exceed \(e\), their first high quotients have prime
parameters \(s-B_i\). Their actual first-level comparison
recovers the character by Lemma~\ref{lem:B-RD71-first-level},
and the same holds for the original groups by isogeny
invariance. If they straddle \(e\), the next first-level
modules are tame and wild respectively by the preceding
lemma, which is impossible. Thus the only remaining case is
\[
 B_i=pb_i<e,\qquad b_i>0.
 \tag{B-RD74.5}
\]

We now give the complete induction, including its level count.
As long as the previous parameters were below \(e\), the
actual group after \(j\) high quotients has parameter
\[
 D_i^{(j)}=p^j B_i=p^{j+1}b_i,
 \qquad 0\le j\le n-2,
 \tag{B-RD74.6}
\]
and its available actual comparison has level \(p^{n-j}\).
For \(j=0\) the high lines were already identified. For
\(1\le j\le n-2\), the preceding small quotient supplies
the intrinsic stable low line in the current group. Its
first-level module is \(L\)-rational. Since \(n-j\ge2\),
Proposition~\ref{prop:B-RD73-high-recovery} identifies the
current high lines, even for scalar residual modules.
We may therefore take the next actual quotient, leaving
level \(p^{n-j-1}\), which is at least \(p\).

Each \(D_i^{(j)}\) is divisible by \(p\), so it cannot equal
\(e\). If both are greater than \(e\), the next parameters
\(s-D_i^{(j)}\) are prime to \(p\); their actual first-level
comparison recovers equality and the finite characters, by
Lemma~\ref{lem:B-RD71-first-level}. If they lie on opposite
sides of \(e\), the next first-level representations have
trivial and nontrivial wild inertia, by
Lemma~\ref{lem:B-RD74-actual-levels}. This is impossible.
These two conclusions require only the \(p\)-torsion that is
still present even at \(j=n-2\).

The only way to continue through every stage is that both
parameters remain below \(e\), including at \(j=n-2\).
In that case
\[
 0<b_i<\frac{e}{p^{n-1}}<p,
 \qquad b_i\in\{1,\ldots,p-1\}.
 \tag{B-RD74.7}
\]
The common characters in (B-RD74.4) imply
\(B_1\equiv B_2\pmod{p-1}\), hence
\(b_1\equiv b_2\pmod{p-1}\). Reduction modulo \(p-1\) is
injective on precisely \(\{1,\ldots,p-1\}\); thus
\(b_1=b_2\). The common original high character then gives
\(c_1=c_2\pmod N\), proving equality of the special
characters. Every branch has been exhausted, and at no point
was a second-level comparison assumed after only the first
level remained.
\end{proof}

\begin{theorem}[A logarithmic upper bound in the tame absolute-index case]
\label{thm:B-RD74-upper}
\status{Proved; complete independent root and companion proof and exact-code audits passed.}
In the full B.4 monodromy setting of height two, with perfect
residue field, \(p\ge5\), \(\mu_p\subset K\), and \(p\nmid e\),
every integer \(n\ge3\) such that \(e<p^n\) satisfies
\[
 \nu_{\min}(p,e,2)\le n-1.
 \tag{B-RD74.8}
\]
In particular, if \(e>p\) and \(k\) is the unique integer
with \(p^k<e<p^{k+1}\), then
\[
 \left\lfloor\frac{k+3}{2}\right\rfloor
 \le \nu_{\min}(p,e,2)\le\max\{2,k\}.
 \tag{B-RD74.9}
\]
\end{theorem}
\begin{proof}
Pass to the completed maximal unramified base as in the
established character-descent reduction. The cyclotomic
character modulo \(p^n\) has order \(p^{n-1}\), because
\(p\nmid e\); the determinant comparison distinguishes total
dimensions \(0,1,2\) and their finite determinant characters.
The height-one filtered cases are determined residually.
For a mixed whole-supersingular/filtered comparison,
Proposition~\ref{prop:B-RD63-mixed} requires
\(e>p^{2n-3}\). But \(n\ge3\) gives \(n\le2n-3\), so
\(e<p^n\le p^{2n-3}\), excluding this possibility.
Two whole formal groups are covered by the preceding
proposition. This proves (B-RD74.8) in the full setting,
rather than only for potentially good whole groups.

For the upper bound in (B-RD74.9), take
\(n=\max\{3,k+1\}\). For the lower bound set
\(r=\lfloor(k+3)/2\rfloor\). Then \(r\ge2\) and
\(2r-3\le k\), so \(e>p^k\ge p^{2r-3}\).
The actual \(p^r\)-torsion lower construction in
Theorem~\ref{thm:B-RD63-arc} therefore gives
\(\nu_{\min}\ge r\). The conversion uses the manuscript
convention that actual \(p^r\)-torsion is its index \(r-1\).
No equality of these two bounds is asserted when they differ.
\end{proof}

\begin{corollary}[A complete third-index interval]
\label{cor:B-RD74-third-index}
\status{Proved; complete independent root and companion proof and exact-code audits passed.}
Under the hypotheses of the preceding theorem,
\[
 p^3<e<p^4\quad\Longrightarrow\quad
 \boxed{\nu_{\min}(p,e,2)=3}.
 \tag{B-RD74.10}
\]
For \(p=5\), this includes every multiple of four from
\(128\) through \(624\) which is prime to five.
\end{corollary}
\begin{proof}
Take \(n=4\) in the upper bound. Since \(e>p^3\),
Theorem~\ref{thm:B-RD63-arc} with \(r=3\) supplies actual
\(p^3\)-torsion counterexamples, so the least index is at
least three. These bounds coincide. The numerical endpoints
are the first and last multiples of four in \((125,625)\).
\end{proof}

\begin{remark}[Scope and remaining interval]
\label{rem:B-RD74-scope}
\status{Proved; complete independent root and companion proof and exact-code audits passed.}
The iteration uses actual generic torsion, actual finite-flat
high quotients, and their generic semilinear descent. The
second line ensuring rationality is the intrinsic image of
the preceding first-level module; no arbitrary complement
or scalar eigenspace is chosen. The exact third-index result
is a sufficient interval, not an ``if and only if'' statement.
For example, \(p^4<e<p^5\) remains bounded only by
\(3\le\nu_{\min}\le4\) here. The arguments do not settle
absolute ramification divisible by \(p\), higher heights,
or the exact least index throughout all larger intervals.
The retained finite supplement is
\path{code/reduction_descent/verify_iterated_high_quotient.py}.
Its final author run passed in \(0.863016332994448\) seconds
at 22:44:19 UTC on 19 November 2026. It exhaustively verifies
four actual rank-two module quotients at \(p=5,7\), levels
two and three, and \(81{,}600\) induced-matrix point identities.
It checks all \(196\) permitted sectors for
\((p,n)=(5,3),(5,4),(7,3)\), including every compatible small
parameter pair, and \(406\) explicitly sampled deeper sectors
for \(p=5,7,11,13\) and \(3\le n\le8\).
The complete and sampled parameter-pair counts are respectively
\(88{,}870\) and \(161{,}887\); among the checked stages,
\(16{,}832\) stop when only the final first level remains.
The program also checks the unordered character orbit under
each isogeny and the precise lower/upper-bound arithmetic.
These are finite algebra checks, not a proof of formal-group
existence or a substitute for the induction.

The first executed program and its successful output are
preserved. A subsequent harmless rewrite replaced division by
multiplication in one integer inequality, and the final exact
integer version was rerun successfully. No executed check
failed. The source chain is the fully audited actual quotient,
Newton, lower-break and mixed-threshold proofs B69/B71/B73/B63;
no additional theorem is inferred from the bounded B73
literature search.
The complete independent root and companion proof/code audits
passed. Their author-output-preserving exact reruns took
\(0.8197192030056613\) and \(0.8659103850077372\) seconds,
respectively. All five assertions, exact candidate bytes,
executed versions, and audit records are retained.
\end{remark}


% Promoted after CP27 archive release; exact audited candidate retained.
\subsection{The terminal Frobenius pair in the quotient induction}

Retain the hypotheses and notation of B74: \(p\ge5\),
\(q=p^2\), \(N=q-1\), \(\mu_p\subset K\), and the
absolute ramification \(e\) is prime to \(p\).
Local arguments use algebraically closed residue field and
the actual common good tame field of degree \(N\).

\begin{lemma}[The enlarged terminal interval]
\label{lem:B-RD75-terminal}
\status{Proved; complete independent root and companion proof and exact-code audits passed.}
Suppose two original two-slope groups have parameters
\(B_i=pb_i\), common high-line character, and common finite
determinant. If \(1\le b_i\le p\), their finite special
inertia characters agree.
\end{lemma}
\begin{proof}
The same character equations as (B-RD74.4) give
\[
 b_1\equiv b_2\pmod{p-1},\qquad
 c_2-c_1\equiv p(b_2-b_1)\pmod N,
 \qquad c_i\equiv pb_i\pmod{p+1}.
 \tag{B-RD75.1}
\]
The last congruence in (B-RD75.1) is exactly (B-RD69.2):
the canonical high-line character is \(\mathbf F_p\)-valued,
which gives \(c_i\equiv B_i=pb_i\pmod{p+1}\).
Equal \(b_i\) give equal \(c_i\pmod N\). The only unequal
congruent pair in \(\{1,\ldots,p\}\) is \(\{1,p\}\).
Reorder to \(b_1=1,b_2=p\). Then
\[
 c_1\equiv p\pmod{p+1},\qquad
 c_2-c_1\equiv p(p-1)\equiv(p-1)c_1\pmod N.
 \tag{B-RD75.2}
\]
Thus \(c_2\equiv pc_1\pmod N\), and
\(\{c_2,pc_2\}=\{c_1,pc_1\}\) because \(p^2\equiv1\pmod N\).
These are precisely the finite special eigenvalue multisets.
The parameters need not be equal: only the requested character
is being recovered. This calculation is at the original
marked tangent exponents, so no change of the common inertia
generator or covariant convention is implicit.
\end{proof}

\begin{proposition}[A larger induction upper range]
\label{prop:B-RD75-formal}
\status{Proved; complete independent root and companion proof and exact-code audits passed.}
Let \(n\ge3\). Whole potentially supersingular height-two
groups have their finite special characters determined by
actual \(p^n\)-torsion whenever
\[
 e<p^{n-1}(p+1).
 \tag{B-RD75.3}
\]
\end{proposition}
\begin{proof}
Run the complete actual-quotient induction of
Proposition~\ref{prop:B-RD74-formal} unchanged through all
its stopping cases. The actual quotient level, semilinear
descent, high-line recovery, prime-parameter case and tame/wild
case do not use its upper bound on \(e\). That bound is used
only at the final all-small stage. Under the present bound,
that stage instead gives
\[
 0<b_i<\frac{e}{p^{n-1}}<p+1,
 \qquad b_i\in\{1,\ldots,p\}.
 \tag{B-RD75.4}
\]
The preceding lemma recovers the original finite character
also in the additional pair \(\{1,p\}\). Every other branch
was already settled before this terminal interval was used.
The full proof still retains only actual \(p\)-torsion after
the last quotient; no further quotient is taken.
\end{proof}

\begin{theorem}[A longer third-index interval]
\label{thm:B-RD75}
\status{Proved; complete independent root and companion proof and exact-code audits passed.}
In the full B.4 height-two monodromy setting with perfect
residue field, \(p\ge5\), \(\mu_p\subset K\), and \(p\nmid e\),
every \(n\ge4\) satisfies
\[
 e<p^{n-1}(p+1)\quad\Longrightarrow\quad
 \nu_{\min}(p,e,2)\le n-1.
 \tag{B-RD75.5}
\]
Consequently
\[
 p^3<e<p^3(p+1)\quad\Longrightarrow\quad
 \boxed{\nu_{\min}(p,e,2)=3}.
 \tag{B-RD75.6}
\]
For \(p=5\), every permitted multiple of four from \(128\)
through \(748\), prime to five, therefore has exact index three.
\end{theorem}
\begin{proof}
The determinant and height-one filtered reductions in
Theorem~\ref{thm:B-RD74-upper} are unchanged. Its mixed
whole/filtered comparison is excluded because, for \(n\ge4\),
\[
 e<p^{n-1}(p+1)<p^{2n-3}.
 \tag{B-RD75.7}
\]
The second strict inequality is equivalent to
\(p+1<p^{n-2}\), valid for \(p\ge5,n\ge4\).
The preceding proposition supplies the whole-group upper
bound, completing the full-monodromy argument.
Taking \(n=4\) and using the B63 actual \(p^3\)-torsion
lower construction whenever \(e>p^3\) makes the two bounds
equal to three. At \(p=5\), the interval is \((125,750)\),
whose first and last permitted multiples of four are \(128\)
and \(748\).
\end{proof}

\begin{remark}[The first nontrivial terminal collision remains open]
\label{rem:B-RD75-scope}
\status{Proved; complete independent root and companion proof and exact-code audits passed.}
The extra terminal pair is harmless because it preserves the
unordered finite character; no formal-group existence is
inferred from its arithmetic compatibility. At the next
terminal value \(p+1\), the pair \(\{2,p+1\}\) need not be
a Frobenius orbit. This is an unresolved arithmetic possibility,
not a constructed comparison of actual torsion groups.
For \(p=5,e=752\), the proved bounds would still be
\(3\le\nu_{\min}\le4\). The general exact third-index range,
wild absolute ramification and higher heights remain open.
The retained exact supplement is
\path{code/reduction_descent/verify_terminal_frobenius_pair.py}.
The author run passed in \(0.08401663199765608\) seconds
at 22:54:37 UTC on 19 November 2026. For all \(24\) primes
from \(5\) through \(101\), it checks the \(75{,}960\)
compatible terminal character pairs, including \(1{,}132\)
unequal Frobenius pairs, and \(144\) full-monodromy threshold
inequalities. The explicit next arithmetic pair is separately
marked as having no asserted realization. A \(p=5,e=128,n=3\)
scope check records that the formal bound applies while the
mixed full-monodromy obstruction does not, preventing an
incorrect index-two conclusion. These finite checks support
the formulas; the proof remains the all-prime calculation
and the already audited actual quotient induction. No
executed computation failed.
The complete independent root and companion proof/code audits
passed; their author-output-preserving exact reruns took
\(0.07519893000426237\) and \(0.0776550739974482\) seconds,
respectively. The exact candidate, all four assertion records,
program, output and audit provenance are retained.
\end{remark}


% Promoted after full root and companion audits; exact wild-index scope is retained.
\subsection{A wild absolute-index cyclic-point restriction}

\begin{lemma}[The number of rational high inverse steps]
\label{lem:B-RD76-point}
\status{Proved; complete independent root and companion proof and exact-code audits passed.}
Let \(p\ge5\), and let \(R\) be a complete discrete valuation
ring of mixed characteristic \((0,p)\) with algebraically closed
residue field. Normalize \(v_R(\pi_R)=1\), put
\(a=v_R(p)\) and \(v=v_p(a)\). If a one-dimensional formal
group of special height two has a rational point of exact order
\(p^n\), where \(n\ge2\), then
\[
             a>(p^2-1)p^{2n-3-v}.                 \tag{B-RD76.1}
\]
The right side is interpreted as a positive rational number if
the exponent is negative.
\end{lemma}
\begin{proof}
Use precisely the integral coordinate, coefficient bounds and
Weierstrass polygons in the proof of
Lemma~\ref{lem:B-RD63-point}, from
\cite[Proposition~5.7; (12.3)--(12.5); Proposition~12.10 and proof]
{HopkinsGross1994Bundles}. Those coefficient statements do not
require \(p\nmid a\). Set \(q=p^2\), \(N=q-1\), and
\(b=v_R(u)\), allowing \(b=+\infty\).
If \(b\ge ap/(p+1)\), all nonzero first-level valuations are
\(a/N\), and all subsequent inverse valuations divide by \(q\).
Rationality therefore gives \(a\ge Nq^{n-1}\), which is strictly
stronger than (B-RD76.1).

Otherwise \(b\) is a positive integer and the first-level
valuations are
\[
 \alpha=\frac{a-b}{p-1},\qquad
 \beta=\frac{b}{p(p-1)},\qquad \alpha>\beta>0.
                                                     \tag{B-RD76.2}
\]
For an input of valuation \(0<Z<a\), the exact inverse polygon
gives
\[
 \begin{cases}
 Z/q,& b\ge Z(p-1)/p,\\
 (Z-b)/p\text{ with multiplicity }p,
 \quad\beta\text{ with multiplicity }q-p,
      &b<Z(p-1)/p.
 \end{cases}                                         \tag{B-RD76.3}
\]
Equality belongs to the uniform first case. Each new valuation
is smaller than its input; all input valuations are below \(a\).
After either a uniform step or a low step to \(\beta\), every
remaining step is uniform, as shown in that same proof.

Let \(P\) be the asserted point. Every point in its successive
multiplication chain is rational, so every valuation is a
positive integer. If the first-level point has valuation
\(\beta\), then
\[
 b\ge p(p-1)q^{n-1},\qquad
 a>\frac{p+1}{p}b\ge Nq^{n-1},
\]
again stronger than required. It remains that the chain starts
at \(Z_0=\alpha\). After \(j\) consecutive high inverse steps,
the exact recurrence in (B-RD76.3) gives
\[
 Z_j=\frac{a/p^j-b}{p-1},\qquad
 a=p^j\bigl(b+(p-1)Z_j\bigr).
                                                     \tag{B-RD76.4}
\]
Integrality of \(Z_j\) and \(b\) forces \(p^j\mid a\), hence
\(j\le v\). This is the only new constraint replacing the
prime-to-\(p\) hypothesis in B63.

Suppose first that the first non-high step is uniform, after
\(j\le\min(v,n-2)\) high steps. There are
\(r=n-1-j\ge1\) uniform inverse steps left. Rationality of
the terminal valuation implies \(Z_j\ge q^r\), and the
uniform inequality implies \(b\ge Z_j(p-1)/p\). Therefore
\[
 a=p^j\bigl(b+(p-1)Z_j\bigr)
    \ge Np^{j-1}Z_j
    \ge Np^{2n-3-j}
    \ge Np^{2n-3-v}.                              \tag{B-RD76.5}
\]
If \(j<v\), the last inequality is strict. If \(j=v\),
equality throughout would give
\(v_p(a)=2n-3-v>v\), because \(v=j\le n-2\).
This contradicts the definition of \(v\). The bound is thus
strict in this case as well.

If the first non-high step is low, after \(j\) high steps,
then the remaining \(n-2-j\) inverse steps divide \(\beta\)
by \(q\). Integrality and the strict high-case inequality give
\[
 \begin{gathered}
 b\ge p(p-1)q^{n-2-j},\qquad
 a/p^j>(p+1)b,\\
 a>Np^{2n-3-j}\ge Np^{2n-3-v}.
 \end{gathered}                                    \tag{B-RD76.6}
\]
Finally, if all \(n-1\) inverse steps are high, then
\(v\ge n-1\) and \(Z_{n-1}\ge1\). Since \(b\ge1\),
\[
 a=p^{n-1}\bigl(b+(p-1)Z_{n-1}\bigr)
       \ge p^n>Np^{n-2}\ge Np^{2n-3-v}.
                                                     \tag{B-RD76.7}
\]
This exhausts every possible chain, including each equality
boundary of the Newton polygon.
\end{proof}

\begin{proposition}[A mixed comparison at wild absolute index]
\label{prop:B-RD76-mixed}
\status{Proved; complete independent root and companion proof and exact-code audits passed.}
Let \(K\) be complete discretely valued of mixed characteristic
\((0,p)\), with perfect residue field and \(p\ge5\). Put
\(e=v_K(p)\) and \(v=v_p(e)\). Let \(M\) be a height-two
BT group which acquires one-dimensional supersingular good
reduction, and let \(H\) have a two-step potentially good
height-one filtration of total dimension one. For \(n\ge2\),
\[
 M[p^n]\simeq H[p^n]\quad\Longrightarrow\quad
                     e>p^{2n-3-v}.                 \tag{B-RD76.8}
\]
\end{proposition}
\begin{proof}
Pass to the completed maximal unramified extension, preserving
\(e\), \(v\), and the comparison. The quotient-and-finite-twist
argument of Proposition~\ref{prop:B-RD63-mixed}, whose complete
construction is in Proposition~\ref{prop:B-RD49-mixed}, produces
a potentially supersingular height-two group with a rational
point of exact order \(p^n\). None of those finite-flat closure,
actual quotient, height-one character, or finite twist steps
assumes \(p\nmid e\).

Lemma~\ref{lem:B-RD49-inertia} supplies a good formal model
over a tame extension of degree \(m\mid N=p^2-1\), for every
absolute index \(e\). The rational point extends to its finite
flat torsion model and lies in the formal group, because the
special fibre is connected. Here \(a=me\) and \(v_p(a)=v\).
The preceding lemma gives \(me>Np^{2n-3-v}\). Since
\(m\le N\), (B-RD76.8) follows.
\end{proof}

\begin{lemma}[Dimension detection with a wild absolute index]
\label{lem:B-RD76-determinant}
\status{Proved; complete independent root and companion proof and exact-code audits passed.}
In the height-two monodromy setting of B.4, assume \(p\ge5\),
\(\mu_p\subset K\), perfect residue field, and write
\(v=v_p(e)\). If \(n\ge v+2\), an isomorphism of actual
\(p^n\)-torsion determines the total dimension of the potentially
good graded pieces and their finite determinant character.
\end{lemma}
\begin{proof}
Pass again to the completed maximal unramified base. The generic
determinant on inertia is \(\chi_p^d\epsilon\), where
\(d\in\{0,1,2\}\) and \(\epsilon\) is
\(\mu_{p-1}\)-valued, by the determinant proof in B49.
Since \(\mu_p\subset K\), the extension \(K(\mu_{p^n})/K\)
is cyclic of degree \(p^r\). Its residue field is unchanged.
In its normalized valuation,
\[
 v_{K(\mu_{p^n})}(\zeta_{p^n}-1)
     =\frac{ep^r}{p^{n-1}(p-1)}\in\mathbf Z_{>0},
 \qquad r\ge n-1-v.                                \tag{B-RD76.9}
\]
Thus the cyclotomic character modulo \(p^n\) has a nontrivial
\(p\)-part when \(n\ge v+2\). Equality of determinants forces
\(p\mid d_1-d_2\), whereas \(|d_1-d_2|\le2<p\), so
\(d_1=d_2\). Their remaining finite characters agree, since
reduction modulo \(p\) is injective on \(\mu_{p-1}\).
\end{proof}

\begin{corollary}[The remaining branch at \(p=5,e=20\)]
\label{cor:B-RD76}
\status{Proved; complete independent root and companion proof and exact-code audits passed.}
Under the full B.4 height-two hypotheses with \(p\ge5\),
\(\mu_p\subset K\), and perfect residue field, suppose
\[
 n\ge\max\{2,v_p(e)+2\},\qquad
 e\le p^{2n-3-v_p(e)}.                              \tag{B-RD76.10}
\]
Any failure of special-character determination by actual
\(p^n\)-torsion must compare two whole potentially
supersingular height-two groups. In particular, at \(p=5,e=20\),
any failure at actual \(125\)-torsion is confined to that branch;
the mixed whole/filtered branch is impossible.
\end{corollary}
\begin{proof}
The determinant lemma fixes the total dimension. Dimension zero
is determined by the Brauer lift of semisimplified residual
inertia; dimension two reduces to it by Cartier duality, exactly
as in B49. In dimension one, each allowed group is either whole
potentially supersingular or filtered by two potentially good
height-one pieces. The characters of two filtered groups are
determined residually because \(\chi_p\equiv1\pmod p\).
The mixed possibility is excluded by (B-RD76.8) and the stated
inequality. At \(p=5,e=20\) we have \(v_p(e)=1\), and
\(n=3\) satisfies \(20\le5^2\) and \(n=v_p(e)+2\).
\end{proof}

\begin{remark}[No wild least-index assertion]
\label{rem:B-RD76-scope}
\status{Proved; complete independent root and companion proof and exact-code audits passed.}
For \(v_p(e)=0\), the mixed necessary inequality is exactly B63.
For positive \(v_p(e)\), it is only a necessary inequality; no
matching existence construction at its new boundary is asserted.
The new \(p=5,e=20\) conclusion eliminates a comparison class
and does not prove that \(125\)-torsion determines every allowed
character. B63 already supplies failure at actual \(25\)-torsion
for this index, so its least manuscript index is at least two.
The whole-potentially-supersingular \(125\)-comparison remains
open. All claims here use the actual torsion level, not a rational
isogeny class or a count of numerical Newton possibilities.

The retained supplement is
\path{code/reduction_descent/verify_wild_cyclic_point.py}.
Its first author run passed in \(3.575578117000987\) seconds
at 23:06:45 UTC on 19 November 2026. It checks \(431{,}375\)
finite Newton parameter pairs, \(100{,}121\) inverse branches
and \(4{,}349\) integral torsion-valuation chains through level
six; \(2{,}734\) chains have wild absolute index. All five
cases of the proof occur in the retained finite data, including
three consecutive high inverse steps. The separate arithmetic
checks cover \(1{,}584\) uniform equality cases,
\(1{,}104{,}000\) good-degree transfers and \(396{,}000\)
cyclotomic denominators. These are exact rational calculations,
not constructions of points or formal groups. The full code,
output, finite ranges and pre-certificate proof are preserved;
no executed computation failed.
The complete independent root and companion proof/code audits
passed; their unchanged-program, author-output-preserving reruns
took \(3.5021978250006214\) and \(3.4459782960038865\) seconds,
respectively. All five assertion records and exact audited
candidate bytes are retained.
\end{remark}


% Proved after full root and companion audits; the nonsquare whole-formal branch remains open.
\subsection{The finite-character gate at the first wild cyclotomic index}

\begin{lemma}[An unweighted iterated product-height-two bound]
\label{lem:B-RD77-kernel}
\status{Proved; complete independent root and companion proof and exact-code audits passed.}
Let \(R\) be a complete discrete valuation ring of mixed
characteristic \((0,p)\), with algebraically closed residue
field and \(p\ge5\). Write \(a=v_R(p)\). Let \(f:F\to G\)
be a homomorphism of finite-height formal groups with no nonzero
finite flat subgroup in its kernel and with finite special-fibre
kernel. Assume that the actual special fibre \(G_k\) is a product
of one-dimensional height-two formal groups. Define
\[
 q_0=1,\qquad q_{j+1}=\min\{a+q_j,1+pq_j\}.
                                                        \tag{B-RD77.1}
\]
If \(q_u\ge a/(p-1)\), then \(F^u\ker f_k=0\).
\end{lemma}
\begin{proof}
The invariant-differential identity
\(U([p]X)J_{[p]}(X)=pU(X)\) shows that every monomial with
an exponent not divisible by \(p\) has coefficient divisible
by \(p\). The product-height-two special fibre makes every
coefficient of degree less than \(p^2\) vanish modulo the
uniformizer, except that the linear term is already \(pX\).
For an augmentation vector of coefficient valuation at least
\(q\ge1\), this gives the lower bound
\[
 \min\{a+q,1+pq,p^2q\}
       =\min\{a+q,1+pq\}.                             \tag{B-RD77.2}
\]
This is the unweighted counterpart of (B-RD64.5); it concerns
the actual target and is invariant under integral coordinate
changes. Iteration gives \(q_j\).

For precision, the normalized logarithm coefficient bound is
\(v(l_\alpha)\ge-a\lfloor\log_p|\alpha|\rfloor\).
Therefore, at the stated closed threshold, coefficientwise
substitution gives
\[
       p^u\log_G(\pi h)\in\pi^{q_u}R[[t]]^{\dim G}
       \quad(h\in tR[[t]]^{\dim G}).                  \tag{B-RD77.3}
\]
The exact integral Cartier action, bounded projection and
saturation argument in Lemma~\ref{lem:B-RD64-weighted-kernel}
uses only this last inclusion after its initial coefficient
calculation. It applies unchanged: for
\(J=\pi^{\lceil a/(p-1)\rceil}R\), the integral
\(\Omega\)-action preserves \(J\), the bounded projection
of \(\mathbf f^u z\) lies in \(J[[\Delta]]\), and the
absence of a nonzero finite flat kernel identifies the source
Cartier module with the target intersection. Thus the saturation
quotient is killed by \(\mathbf f^u\). The cited identification
of that quotient with the Dieudonne module of the special kernel
makes \(\mathbf f\) relative Frobenius. These are precisely
\cite[Propositions~9--12 and Lemma~3, pp.~674--676;
Proposition~16 and proof, p.~681; Theorem~2(8), p.~668]
{Bondarko2006Generic}, with the closed coefficientwise boundary
already proved in Lemma~\ref{lem:B-RD50-boundary}.
No target thickness is inferred from a source or an isogeny.
\end{proof}

\begin{proposition}[Square finite determinants at the first wild index]
\label{prop:B-RD77-square}
\status{Proved; complete independent root and companion proof and exact-code audits passed.}
Let \(p\ge5\), and let \(K\) be complete discretely valued
of mixed characteristic \((0,p)\), with algebraically closed
residue field and \(e=p(p-1)\). Let \(M,N\) be whole
potentially supersingular height-two groups with common
determinant \(\chi_p\psi\). If the finite
\(\mu_{p-1}\)-valued character \(\psi\) is a square in
\(\operatorname{Hom}_{\rm cont}(I_K,\mu_{p-1})\), then
\[
             M[p^3]\simeq N[p^3]
   \quad\Longrightarrow\quad
   \text{their finite special inertia characters agree}.
                                                        \tag{B-RD77.4}
\]
\end{proposition}
\begin{proof}
Twist both groups by the inverse of a square root of \(\psi\).
This is the actual finite twist of
Proposition~\ref{prop:B-RD50-square}; it preserves the given
torsion comparison and preserves inequality of the two special
characters. Their finite special determinants are now one.
The finite height-two inertia groups are cyclic subgroups of
\(\mathbf F_{p^2}^{\times}\), and their norm-one condition
makes their common faithful good degree \(m\) divide \(p+1\).
Thus, over that actual good field,
\[
             a=em=p(p-1)m\le p(p^2-1).               \tag{B-RD77.5}
\]

Suppose the special characters differ. Their norm-one
Frobenius eigenvalue orbits are disjoint, so an integral tame
polynomial projector separates the two special factors, just
as in B50. Let \(H\) be the finite-flat closure of the given
generic graph in \(\mathcal M[p^3]\times\mathcal N[p^3]\).
The perfect generic alternating pairings, and their inverse
duality maps, extend to perfect integral BT pairings by
\cite[Theorem~4 and Corollaries~1--2, p.~181]{Tate1967PDivisible}.
After scaling one by the unit multiplier of the torsion map,
the graph closure is self-orthogonal, including on the special
fibre. The integral projector splits \(H_k\) as a product.
Uniqueness of a rank-\(p^r\) subgroup in each one-dimensional
connected formal group and self-orthogonality force
\[
            H_k=\ker F_{M_k}^{3}\times\ker F_{N_k}^{3}.
                                                        \tag{B-RD77.6}
\]

The actual quotient \(Q=(\mathcal M\times\mathcal N)/H\)
has special fibre \(M_k^{(p^3)}\times N_k^{(p^3)}\).
The map
\[
 \mathcal M\times\mathcal N\longrightarrow Q\times\mathcal M,
 \qquad z\longmapsto(qz,\operatorname{pr}_{\mathcal M}z)
                                                        \tag{B-RD77.7}
\]
has no nonzero finite flat kernel, since its generic kernel is
zero, and has special kernel \(0\times\ker F_{N_k}^{3}\).
Its actual target is a product of three height-two special
groups. Apply the preceding lemma with \(u=2\). Here
\[
 q_1=p+1,\qquad
 q_2=\min\{a+p+1,p^2+p+1\}\ge a/(p-1),              \tag{B-RD77.8}
\]
because \(a\ge p(p-1)\), \(a\le p(p^2-1)\), and both
entries of the minimum are at least that threshold. The
resulting assertion that \(F^2\) kills \(\ker F^3\) is
impossible: its image has rank \(p\). This proves the proposition.
\end{proof}

\begin{corollary}[Two remaining patterns at \(p=5,e=20\)]
\label{cor:B-RD77}
\status{Proved; complete independent root and companion proof and exact-code audits passed.}
In the full B.4 height-two setting with perfect residue field,
\(p\ge5\), \(\mu_p\subset K\), and \(e=p(p-1)\),
\[
                    2\le\nu_{\min}(p,p(p-1),2)\le3.
                                                        \tag{B-RD77.9}
\]
Any actual \(p^3\)-torsion failure must compare two whole
potentially supersingular groups with a nonsquare common finite
determinant. After a simultaneous finite twist, its value may
be normalized to be primitive of order \(p-1\).

For \(p=5,e=20\), the joint faithful tame degree of such a
remaining pair is \(24\). If \(\zeta\) is a primitive
\(24\)th root whose sixth power is the prescribed primitive
norm, the special characters have distinct types among
\[
 A=\{\zeta,\zeta^5\},\qquad
 B=\{\zeta^9,\zeta^{21}\},\qquad
 C=\{\zeta^{13},\zeta^{17}\}.                         \tag{B-RD77.10}
\]
Up to interchange and a simultaneous quadratic twist, the only
remaining type pairs are \(\{A,B\}\) and \(\{A,C\}\).
This is a necessary finite-character list, not an existence
assertion for torsion-congruent groups.
\end{corollary}
\begin{proof}
The B63 actual \(p^2\)-torsion construction applies because
\(p(p-1)>p\), and gives the lower bound two. At actual level
\(p^3\), Lemma~\ref{lem:B-RD76-determinant} applies because
\(v_p(e)=1\), and Corollary~\ref{cor:B-RD76} removes every
case except the whole-formal comparison. The preceding proposition
removes its square finite determinant.

For the upper bound, use actual \(p^4\)-torsion. Dimension
detection again applies, while the mixed comparison would require
\(e>p^{2\cdot4-3-1}=p^4\), contrary to \(e=p(p-1)\).
For two whole formal groups, the common determinant is fixed.
In Proposition~\ref{prop:B-RD64-formal} take
\(a=e(p^2-1)=p(p-1)(p^2-1)\). Its actual-target weighted
recurrence gives
\[
 q_1=2p-1,\qquad q_2=2p^2-1,\qquad q_3=2p^3-1
                  >p(p^2-1)=a/(p-1).                \tag{B-RD77.11}
\]
Thus level \(p^4\), or manuscript index three, suffices.

The normalization and finite-character enumeration do not
depend on the absolute ramification. The procyclic tame
character argument of Corollary~\ref{cor:B-RD50-remaining}
changes any nonsquare norm to a primitive one. The arithmetic
proof of Lemma~\ref{lem:B-RD51-orbits} then applies unchanged.
For \(p=5\), it permits joint degree \(8\) or \(24\), but
the primitive norm fibre in \(\mu_8\) has only one Frobenius
orbit, so different special characters require \(24\).
The three orbits and their simultaneous quadratic-twist action
are exactly the elementary computation in
Corollary~\ref{cor:B-RD51-p5}, giving (B-RD77.10) and the two
stated patterns. No earlier least-index statement at \(e=p-1\)
is imported to this wild index.
\end{proof}

\begin{remark}[The remaining wild comparison is still actual]
\label{rem:B-RD77-scope}
\status{Proved; complete independent root and companion proof and exact-code audits passed.}
The square exclusion uses only the actual quotient special fibre
and the coefficientwise logarithm bound; it does not transport a
source Hasse thickness through an isogeny. The remaining
\(p=5,e=20\) question is whether either of the displayed
nonsquare patterns occurs for an isomorphism of actual
\(125\)-torsion. The finite-character list and the upper index
three do not resolve it or supply a counterexample.

The retained exact supplement is
\path{code/reduction_descent/verify_first_wild_finite_patterns.py}.
Its first author run passed in \(0.6461968240037095\) seconds
at 23:15:03 UTC on 19 November 2026. For the \(24\) primes
from \(5\) through \(101\), it checks \(174\) norm-one
good-degree rows, \(348{,}000\) closed coefficientwise logarithm
inequalities, \(2{,}560{,}584\) low-degree coefficient bounds,
and all \(24\) weighted three-step upper bounds. The explicit
\(p=5\) primitive-norm fibres and quadratic-twist pair orbits
are retained separately. These finite arithmetic checks do
not imply existence of a BT group or a torsion congruence.
The complete code, output, finite ranges and pre-certificate
proof are retained; no executed computation failed.
The complete independent root and companion proof/code audits
passed, with unchanged-program reruns taking
\(0.7250874890014529\) and \(0.630934568005614\) seconds,
respectively. The author output and exact audited candidate
were preserved. The audits reread the full internal B64 and B50
proofs; they reused the previously audited original Bondarko
and Tate passages without claiming a fresh external-source read.
\end{remark}


% Proved after full independent audits. Potential CM means finite-extension endomorphisms.
\subsection{The last Lubin--Tate layer and the first wild CM index}

Use the definition of potentially CM in B52: the rational
endomorphism algebra over some finite extension with good
reduction contains a quadratic field over \(\mathbf Q_p\).
This is not the endomorphism algebra of an abstract constant
generic etale group over an algebraic closure.

\begin{lemma}[The nonsquare integral CM model at arbitrary index]
\label{lem:B-RD78-model}
\status{Proved; complete independent root and companion proof, primary-source and exact-code audits passed.}
Let \(p\ge5\), and let \(K\) be complete discretely valued
of mixed characteristic \((0,p)\), with algebraically closed
residue field and \(\mu_p\subset K\). There is no restriction
on its absolute index \(e\). Let \(M\) be a potentially CM,
potentially supersingular height-two group, whose determinant
is \(\chi_p\psi\), with nonsquare finite
\(\mu_{p-1}\)-valued character \(\psi\).
Then its quadratic endomorphism field is the unramified field
\(E=\mathbf Q_{p^2}\); all finite-extension endomorphisms are
defined over \(K\), and its Tate lattice is a free rank-one
\(\mathcal O_E\)-module. After choosing its tangent orientation,
\[
        T_p(M)=\mathcal O_E(\theta\epsilon),          \tag{B-RD78.1}
\]
where \(\theta\) is the character of the standard formal
\(\mathcal O_E\)-module \([p](T)=pT+T^{p^2}\), base changed
to \(K\), and \(\epsilon\) is finite and tame with values
in \(\mu_{p^2-1}\).
\end{lemma}
\begin{proof}
We verify the dependence on the absolute index in the complete
proofs of Lemmas~\ref{lem:B-RD52-CM} and~\ref{lem:B-RD52-LT}.
The faithful good extension is tame of degree dividing
\(p^2-1\) by Lemma~\ref{lem:B-RD49-inertia}, at arbitrary
\(e\). Tate full faithfulness extends integral endomorphisms
to its good model; clearing a power of \(p\) does the same
for rational ones. Reduction is injective by
\cite[Proposition~4.2 and proof]{HopkinsGross1994Bundles}.
Since the residue field is algebraically closed, finite-extension
endomorphisms therefore descend to the good field. Their
rational algebra is a division algebra acting on a two-dimensional
\(\mathbf Q_p\)-space, and the CM assumption makes it a
quadratic field.

If that field is ramified, its special prime-to-\(p\) normalizer
has order dividing \(2(p-1)\). For \(p\equiv3\pmod4\)
its reduced norms are squares. For \(p\equiv1\pmod4\), the
generic normalizer consists of multiplication by a unit and
such multiplication followed by conjugation; its determinants
are \(\pm N_{E/\mathbf Q_p}(u)\), both squares modulo \(p\).
Here \(\chi_p\equiv1\pmod p\), so either case contradicts
the nonsquare \(\psi\). Thus \(E\) is unramified.
The noncentral coset of its normalizer in the special division
algebra consists of \(E^\times\Pi\), whose valuations are
odd, so has no units. Special finite inertia centralizes \(E\),
and reduction injectivity makes the generic action centralize
\(E\) as well. This also descends the endomorphisms to \(K\).

Choose a Galois element with nonsquare finite determinant.
Its action is multiplication by a unit of \(E\) with nonsquare
residue norm, so its residue does not lie in \(\mathbf F_p\).
The algebra generated by that action over \(\mathbf Z_p\)
is \(\mathcal O_E\), since its quadratic polynomial has
irreducible separable reduction. The stable Tate lattice is
therefore free of rank one over \(\mathcal O_E\).

The integral action extends to the good model by Tate full
faithfulness. The explicit unit conjugator and deformation
argument in Lemma~\ref{lem:B-RD52-LT} identify this model with
the standard oriented Lubin--Tate group: averaging the universal
parameter against a primitive \((p^2-1)\)st-root action gives
a parameter multiplied by \(\lambda^{p-1}\ne1\), and an
\(\mathcal O_E\)-linear deformation must send it to zero.
The quotient of its generic character by \(\theta\) is trivial
on the tame good field, hence is the finite tame character
\(\epsilon\) in (B-RD78.1). None of these steps uses
\(p\nmid e\). Only the subsequent first-wild-layer claim
of B52 used that hypothesis; we replace it below.
\end{proof}

\begin{lemma}[A full terminal layer after wild base change]
\label{lem:B-RD78-layer}
\status{Proved; complete independent root and companion proof, primary-source and exact-code audits passed.}
Retain the preceding notation, set \(v=v_p(e)\), and let
\(L/K\) be any common good tame extension. If \(n\ge v+2\),
the image of \(\theta|_{G_L}\) modulo \(p^n\) contains
\[
          (1+p^{n-1}\mathcal O_E)/(1+p^n\mathcal O_E).
                                                        \tag{B-RD78.2}
\]
\end{lemma}
\begin{proof}
Choose a coefficient ring \(W(k)\subset\mathcal O_K\)
by the Cohen structure theorem, and set \(K_0=W(k)[1/p]\).
This uses \cite[Theorem~10.160.8, Tag~032A]{Stacks032A}.
Because the residue field is perfect and \(K\) is a complete
discrete valuation field, its ring of integers is finite free
of rank \(e\) over this coefficient ring: a uniformizer and
Weierstrass preparation give an Eisenstein relation of degree
\(e\), with basis \(1,\pi,\ldots,\pi^{e-1}\).
Thus \([L:K_0]=e[L:K]\), whose \(p\)-adic valuation is \(v\).
No identification of \(k\) with \(\overline{\mathbf F}_p\)
is made.

Let \(D_n\) be the full \(p^n\)-division field of the standard
Lubin--Tate group over \(K_0\), and put \(q=p^2\).
The direct proof of
\cite[Theorem~2(b),(d),(e) and complete proof, pp.~383--384]
{LubinTate1965FormalCM} works after this residue extension.
Explicitly, the primitive division polynomial
\[
 \Phi_n(T)=\bigl([p]^{\circ(n-1)}(T)\bigr)^{q-1}+p
                                                        \tag{B-RD78.3}
\]
is Eisenstein of degree \((q-1)q^{n-1}\), and a primitive
point generates the free rank-one \(\mathcal O_E/p^n\)-module
of torsion. Every other point is its \(\mathcal O_E\)-multiple.
The faithful Galois injection into
\((\mathcal O_E/p^n)^\times\), together with that Eisenstein
degree, is therefore an equality, and \(D_n/K_0\) is totally
ramified. This is an elementary division-field argument, not
an appeal to local class field theory over the larger residue
field.

Set \(U=(1+p\mathcal O_E)/(1+p^n\mathcal O_E)\), and let
\(J\) be the image after base change to \(L\). Galois
intersection gives
\[
 [\operatorname{Gal}(D_n/K_0):J]=[D_n\cap L:K_0],
 \qquad [U:U\cap J]\le p^v.                         \tag{B-RD78.4}
\]
The second index is the \(p\)-part of the first, since the
unit group is abelian with prime-to-\(p\) Teichmuller factor;
the intersection degree divides \([L:K_0]\).
An abelian group of order at most \(p^v\) is killed by \(p^v\),
so \(U^{p^v}\subset J\). For odd \(p\), successive binomial
expansions give
\[
 U^{p^v}=(1+p^{v+1}\mathcal O_E)/(1+p^n\mathcal O_E).
                                                        \tag{B-RD78.5}
\]
Indeed the \(p\)th-power map has leading term
\((1+p^r b)^p=1+p^{r+1}b\pmod{p^{r+2}}\) for \(r\ge1\);
successive lifts give surjectivity onto the next layer.
For \(n\ge v+2\), (B-RD78.5) contains (B-RD78.2).
\end{proof}

\begin{proposition}[Nonsquare CM character recovery]
\label{prop:B-RD78-CM}
\status{Proved; complete independent root and companion proof, primary-source and exact-code audits passed.}
Let \(K\) have perfect residue field and the other hypotheses
above, and put \(v=v_p(e)\). Two potentially CM, potentially
supersingular height-two groups with common nonsquare finite
determinant and isomorphic actual \(p^n\)-torsion have equal
special inertia characters whenever \(n\ge v+2\).
\end{proposition}
\begin{proof}
Pass to the completed maximal unramified field and choose
compatible orientations. The first lemma gives
\(T_p(M_i)=\mathcal O_E(\theta\epsilon_i)\), with the same
\(\theta\). On a common good tame field, both finite twists
are trivial. An intertwiner modulo \(p^n\) commutes with every
operator \(1+p^{n-1}b\), for \(b\in\mathcal O_E\), by the
preceding lemma. Dividing the commutator by \(p^{n-1}\) shows
that its reduction modulo \(p\) centralizes all multiplication
operators from \(\mathbf F_{p^2}\). This centralizer consists
of field multiplication itself. The reduction of the intertwiner
is therefore multiplication by a nonzero field element.

The full intertwining identity for every \(g\in G_K\), reduced
modulo \(p\), now gives
\(\bar\epsilon_1(g)=\bar\epsilon_2(g)\), after cancelling
that scalar and the nonzero scalar \(\bar\theta(g)\).
Reduction is injective on \(\mu_{p^2-1}\), so the two finite
characters agree. The good factor has trivial finite special
inertia; the covariant dictionary of B52 therefore gives the
same requested character for the two groups.
\end{proof}

\begin{theorem}[An exact first-wild index in the potential-CM subclass]
\label{thm:B-RD78}
\status{Proved; complete independent root and companion proof, primary-source and exact-code audits passed.}
For \(p\ge5\), \(e=p(p-1)\), \(\mu_p\subset K\), and
perfect residue field, actual \(p^3\)-torsion determines the
special character of two potentially CM, potentially supersingular
height-two groups. Actual \(p^2\)-torsion does not always do so
within this same subclass. Its least manuscript index is thus
exactly two.

Consequently any remaining full B.4 counterexample at actual
\(125\)-torsion for \(p=5,e=20\) must have at least one
whole formal model without quadratic finite-extension
endomorphisms, in either normalized pattern of B77.
\end{theorem}
\begin{proof}
For the upper bound, \(v_p(e)=1\), and the finite determinants
agree by Lemma~\ref{lem:B-RD76-determinant}. A square finite
determinant is covered by Proposition~\ref{prop:B-RD77-square};
the nonsquare case is the preceding proposition with \(n=3\).

For a lower example take \(k=\overline{\mathbf F}_p\),
\(K_0=W(k)[1/p]\), and the standard Lubin--Tate group \(F_0\).
Let \(D_2/K_0\) be its full \(p^2\)-division field, with
\[
 \operatorname{Gal}(D_2/K_0)
 = (\mathcal O_E/p^2)^\times
 =\mu_N\times(1+p\mathcal O_E)/(1+p^2\mathcal O_E),
 \quad N=p^2-1.
                                                        \tag{B-RD78.6}
\]
In this actual Galois group take
\[
 H=\mu_{p+1}\times
       \{1+pa:a\in\mathbf F_p\},\qquad K=D_2^H.
                                                        \tag{B-RD78.7}
\]
The field is completely specified. Its absolute ramification is
\(Np^2/[p(p+1)]=p(p-1)\). The determinant of the good
height-two dimension-one group \(F_0\) is cyclotomic and equals
the norm of \(\theta\). The subgroup \(H\) has norm one
modulo \(p\), hence \(\mu_p\subset K\).

Choose a uniformizer \(\pi\) of \(K\), and let
\(\omega:G_K\to\mu_N\) be the tame character from
\(\pi^{1/N}\). The first division equation is \(T^N=-p\).
Since the unit in \(-p=\pi^e u\) has an \(N\)th root in
\(K\) by Hensel's lemma, its Teichmuller torsion character is
\(\bar\theta=\bar\omega^e\). Write \(\sigma\) for the
unramified automorphism of \(\mathcal O_E\).
On the exact mod-\(p^2\) image \(H\), the wild factor in
(B-RD78.7) is scalar and fixed by \(\sigma\). Therefore
\[
 \frac{\sigma(\theta)}{\theta}
       =\omega^{(p-1)e}\pmod{p^2}.
                                                        \tag{B-RD78.8}
\]
Define actual finite CM twists of \(F_0\) by
\[
 \epsilon_1=\omega,\qquad
 \epsilon_2=\omega^{3p-2}.
                                                        \tag{B-RD78.9}
\]
They are effective generic BT twists, and acquire the original
good CM model over a tame extension. Since
\(p+(p-1)e\equiv3p-2\pmod N\), equation (B-RD78.8) says
that \(\sigma\), as a \(\mathbf Z/p^2\)-linear map of
\(\mathcal O_E/p^2\), intertwines
\(\theta\epsilon_1\) with \(\theta\epsilon_2\).
This is an isomorphism of the actual \(p^2\)-torsion modules,
and hence of generic finite group schemes.

Their finite covariant special eigenvalue exponents are
\[
                 \{1,p\},\qquad \{3p-2,3-2p\}
                 \pmod N.                            \tag{B-RD78.10}
\]
Here the good factor has trivial finite descent and the finite
twist contributes \(\epsilon_i\), exactly as in the explicit
covariant convention of B52. The exponent \(3p-2\) is neither
\(1\) nor \(p\) modulo \(N\): their differences are
\(3(p-1)\) and \(2(p-1)\), and \(p+1\ge6\) divides
neither \(3\) nor \(2\). The characters therefore differ.
Both finite norms are the same primitive character, since
\(3p-2\equiv1\pmod{p-1}\). This supplies the lower example
inside the stated CM subclass. At \(p=5\) its finite types
are \(A\) and \(C\).
\end{proof}

\begin{remark}[The unrestricted wild problem is not replaced by CM]
\label{rem:B-RD78-scope}
\status{Proved; complete independent root and companion proof, primary-source and exact-code audits passed.}
The terminal-layer proposition uses an index bound on an actual
Lubin--Tate division image and integral endomorphisms; it does
not assert that an arbitrary height-two formal group is CM.
The lower pair is an actual finite extension and actual pair
of generic BT twists, not merely two compatible finite characters.
The first-wild full least index remains in \(\{2,3\}\) until
the remaining non-CM \(p^3\)-comparisons are resolved.


\smallskip
\noindent\emph{Retained exact supplement.}
The independently rerunnable program
\path{code/reduction_descent/verify_wild_CM_layer.py}
checks 15,696 actual mod-\(p^2\) twist identities in explicit
unramified quadratic rings, 57,906 bounded-index additive
subgroups and their 115,812 nonzero terminal-layer basis tests,
46,228 residual centralizer matrices, 24 primitive Eisenstein
reductions, and 24 prime-index character rows.
The corrected author run passed in
\(0.21962125301070046\) seconds.
The first executed version and output are retained separately:
its additional terminal test used the ambient zero class
\(p^{n-1}\), although its stronger \(p^v\)-containment test
was already correct. The corrected explicit coordinate is
\(p^{n-2}\), because logarithmic coordinates for
\((1+p\mathcal O_E)/(1+p^n\mathcal O_E)\) divide by \(p\).
No finite range replaces the field construction or the
all-prime proof above.
The independent root and companion audits read every proof and
the entire corrected program; unchanged-program reruns passed in
\(0.21194225999352057\) and \(0.22546726299333386\) seconds.
Both freshly read the original Lubin--Tate proofs on
pp.~381--384 and the complete Stacks Tag~032A proof, retaining
the finite-residue limitation of the original Lubin--Tate source.
The exact audited candidate and every executed earlier output
remain preserved.
\end{remark}


% Proved after full post-CP28 independent audits. Larger wild images remain open.
\subsection{The minimal third-level wild image at the first wild index}

Throughout, \(p\ge5\), \(q=p^2\), \(N=q-1\), and \(K\) is a
complete discretely valued field of mixed characteristic \((0,p)\),
with algebraically closed residue field, \(\mu_p\subset K\), and
absolute index \(e=p(p-1)\). All torsion modules below are actual
generic torsion of height-two groups with potentially supersingular
good reduction. Their finite special characters use the covariant
convention fixed in B53.

\begin{lemma}[Third-level valuations and the two wild layers]
\label{lem:B-RD79-layers}
\status{Proved; complete independent root and companion proof and exact-code audits passed.}
Suppose the residual inertia module of \(M\) is irreducible over
\(\mathbf F_p\). Let \(L/K\) be the good tame extension of degree
\(N\), and put \(F=L(M[p^3])\), \(W=\operatorname{Gal}(F/L)\).
Then \(M[p]\) is \(L\)-rational, and every point of exact order
\(p^j\), \(1\le j\le3\), has coordinate valuation
\[
             v_L(T(P))=\frac{e}{q^{j-1}}.
                                                     \tag{B-RD79.1}
\]
Consequently \(W\) has order \(p^r\), where \(3\le r\le8\).
If \(r=3\), its reduction image and terminal kernel are
\[
 \begin{split}
 W_1&=\{A\bmod p:1+pA+p^2B\in W\}\subset M_2(\mathbf F_p),\\
 W_2&=\{B\bmod p:1+p^2B\in W\}\subset M_2(\mathbf F_p),\\
 \dim_{\mathbf F_p}W_1&=1,\qquad
 \dim_{\mathbf F_p}W_2=2,\qquad W_1\subset W_2.
 \end{split}                                         \tag{B-RD79.2}
\]
The inclusion in the last line is in the same actual matrix
coordinates, induced by taking \(p\)th powers.
\end{lemma}
\begin{proof}
A proper nonzero valuation-filtration step in \(M[p]\) would be
inertia stable, contradicting irreducibility. Thus all its nonzero
points have one valuation. The product of the \(q-1\) nonzero
parameters has valuation \(v_L(p)=eN\), so this valuation is \(e\).
The normal wild subgroup fixes a nonzero vector in characteristic
\(p\); its fixed space is inertia stable. Irreducibility therefore
makes wild inertia trivial. The remaining cyclic irreducible
image embeds in \(\mathbf F_q^\times\), and is killed by the tame
extension of degree \(N\). The good-field assertion itself is
Lemma~\ref{lem:B-RD49-inertia}, which permits arbitrary \(e\).

The first-level Newton polygon is the segment from \((1,eN)\)
to \((q,0)\). Its intermediate degree-\(i\) coefficients have
valuation at least \((q-i)e\). For a target of valuation
\(0<Z\le e\), these coefficients lie strictly above the segment
from \((0,Z)\) to \((q,0)\): for \(0<i<q\),
\((q-i)e>(1-i/q)Z\). Hence all its \(q\) inverse points have
valuation \(Z/q\). Applying this twice proves (B-RD79.1).
This is the same exact inverse-polygon argument as
Lemma~\ref{lem:B-RD55-valuations}; no approximation by an abstract
formal module is made.

Since the residual module is rational over \(L\), the group \(W\)
lies in the kernel of
\(\GL_2(\mathbf Z/p^3)\to\GL_2(\mathbf F_p)\), of order \(p^8\).
All residue extensions are trivial. Every primitive third-level
point has valuation
\(e/q^2=(p-1)/p^3\), whose reduced denominator is \(p^3\).
Its point field therefore has degree at least \(p^3\) over \(L\).
This proves \(3\le r\le8\).

Both spaces in (B-RD79.2) are additive, and
\(r=\dim W_1+\dim W_2\). For \(p\ge3\), the exact matrix identity
\[
       (1+pA+p^2B)^p=1+p^2A\pmod{p^3}
                                                     \tag{B-RD79.3}
\]
shows \(W_1\subset W_2\). If \(W_1=0\), each primitive point
has at most \(p^2\) conjugates: a terminal-layer element moves
it by an element of its first-level module, of cardinality \(p^2\).
This contradicts its valuation denominator. Thus, when \(r=3\),
the nonnegative dimensions with sum three, with the first positive
and no greater than the second, are exactly \((1,2)\).
\end{proof}

\begin{proposition}[The minimal third-level image cannot conceal the character]
\label{prop:B-RD79-minimal}
\status{Proved; complete independent root and companion proof and exact-code audits passed.}
Let \(M_1[p^3]\simeq M_2[p^3]\) be an actual isomorphism with
common irreducible residual inertia module. If its wild image
has order \(p^3\), the two finite special characters agree.
There is no CM assumption.
\end{proposition}
\begin{proof}
Use the same actual good field \(L/K\) of degree \(N\) and
common full division field
\(F=L(M_1[p^3])=L(M_2[p^3])\). Put \(R_i=M_i[p^3](F)\) and
\(V_i=R_i/pR_i\). The preceding lemma and \([F:L]=p^3\) give
\[
 \begin{gathered}
 A=v_F(T_i(P))=p-1\quad(P\text{ primitive of order }p^3),\\
 v_F(T_i(Q))=qA\quad(Q\text{ of exact order }p^2),\qquad
 v_F(T_i(S))=q^2A\quad(S\ne0\text{ of order }p),\\
                           v_F(p)=Nq^2A.
 \end{gathered}                                        \tag{B-RD79.4}
\]
Choose a common uniformizer \(\varpi\) of \(F\) and special
Honda coordinates with \([p](T)\equiv T^q\) on the special
fibre. Leading coefficients give injective linear maps
\[
       \ell_i:V_i\longrightarrow k,\qquad
       \ell_i(P)=\overline{T_i(P)/\varpi^A}.
                                                     \tag{B-RD79.5}
\]
Indeed two primitive points have the same leading coefficient
exactly when their difference lies in \(pR_i\): the latter
points have valuation at least \(qA>A\), whereas every point
outside \(pR_i\) has valuation \(A\). Integral formal addition
is ordinary addition in the leading term. The inverse-polygon
calculation shows that the degree-\(q\) term is uniquely lowest
at every primitive third- or second-level point. Consequently
multiplication by \(p^2\) sends leading coordinates to
\(\ell_i(P)^{q^2}\).

The first-level leading space is \(d\mathbf F_q\) for a nonzero
\(d\). To see this with the present absolute index, let
\(\zeta\) be the chosen primitive tame \(N\)th root. If the
special tangent exponent is \(c_i\), the first-level leading
action is multiplication by \(\zeta^{c_i+e}\), using
\(v_L(T(S))=e\). It is not in \(\mathbf F_p\) by residual
irreducibility. A two-dimensional \(\mathbf F_p\)-space stable
under this scalar is a one-dimensional \(\mathbf F_q\)-space.
Taking the unique \(q^2\)th roots in the algebraically closed
residue field shows that the image of (B-RD79.5) is likewise
\(d_i\mathbf F_q\), for some \(d_i\ne0\). Normalize it by
\(\ell_i(v)=d_i z\); its first-level normalization is then
\(d_i^{q^2}z\), since \(z^{q^2}=z\) on \(\mathbf F_q\).

Write \(U=\ker(W\to W_1)\), so that \(|U|=q\). For a
nonidentity \(g\in U\), its action on \(R_i\) is
\(1+p^2 B_i(g)\). Let \(b_g\) be its actual lower break and
\[
 d_g=\overline{(g\varpi-\varpi)/\varpi^{b_g+1}}\ne0.
\]
Choose residue representatives in a coefficient ring inside
\(\mathcal O_K\), as in the complete fixed-coefficient proof
of Lemma~\ref{lem:B-RD60-break} and
\cite[Theorem~10.160.8 and proof]{Stacks032A}.
Every \(g\) fixes these representatives. Since \(p\nmid A\),
the same uniformizer expansion proves for every \(x\) of
valuation \(A\)
\[
 v_F(gx-x)=A+b_g,\qquad
 \overline{(gx-x)/\varpi^{A+b_g}}
       =\overline A\,d_g\,\overline{x/\varpi^A}.
                                                     \tag{B-RD79.6}
\]
Explicitly the first binomial change of its degree-\(A\) term
is nonzero; every higher degree changes in valuation at least
its degree plus \(b_g\), strictly later. This accounts for
all unit terms, rather than replacing a point by a monomial.

Thus \(g\) moves every primitive point. Its displacement is a
nonzero first-level point, so (B-RD79.4) and (B-RD79.6) force
\[
          b_g=(q^2-1)A\quad(g\in U\setminus\{1\}).
                                                     \tag{B-RD79.7}
\]
Formal and ordinary coordinate subtraction have the same
leading coefficient. Applying the \(p^2\)-series leading term
therefore gives
\[
   \ell_i(B_i(g)v)^{q^2}
               =\overline A\,d_g\,\ell_i(v).
                                                     \tag{B-RD79.8}
\]
In the normalized field coordinates this says that \(B_i(g)\)
is multiplication by
\[
        \lambda_i(g)=\kappa_i d_g\in\mathbf F_q,
             \qquad \kappa_i=\overline A d_i^{1-q^2}.
                                                     \tag{B-RD79.9}
\]
The representation of \(U\) by these matrices is injective,
so \(\lambda_i(U)=\mathbf F_q\), with zero assigned to the
identity. In particular this is the full Cartan plane, obtained
from the actual last-layer action, not assumed from its dimension.

Suppose the requested finite characters differed. The isomorphism
on the first-level leading spaces must be of the form
\(z\mapsto a z^p\), \(a\in\mathbf F_q^\times\).
Here is the exact reason, independent of the index restriction
in the earlier B53 statement. The common irreducible tame
representations have eigenvalue orbits
\(\{\zeta^{c_i+e},\zeta^{p(c_i+e)}\}\).
An intertwiner is either field-linear, identifying \(c_1=c_2\)
modulo \(N\), or Frobenius-semilinear. The field-linear case
would give equal special characters, contrary to assumption.
The isomorphism on \(V_i\) is the same one under multiplication
by \(p^2\), since that multiplication raises normalized leading
coordinates to \(z^{q^2}=z\).
It follows that
\[
            \lambda_2(g)=\lambda_1(g)^p
                        \quad\text{for every }g\in U.
                                                     \tag{B-RD79.10}
\]
But (B-RD79.9) simultaneously makes \(\lambda_2(g)\) a fixed
scalar multiple of \(\lambda_1(g)\). Since the latter run through
all of \(\mathbf F_q\), this would make its nontrivial Frobenius
a scalar map. Evaluation at one and at an element outside
\(\mathbf F_p\) is a contradiction. Hence the special characters
agree.
\end{proof}

\begin{corollary}[A stricter remaining irreducible-residual gate]
\label{cor:B-RD79}
\status{Proved; complete independent root and companion proof and exact-code audits passed.}
At the first wild absolute index \(e=p(p-1)\), any actual
\(p^3\)-congruent pair with irreducible residual inertia and
different finite special characters has wild image of order
\(p^r\) with \(4\le r\le8\).
In particular, at \(p=5,e=20\), its common actual \(125\)-torsion
wild image has order at least \(625\).
\end{corollary}
\begin{proof}
The first lemma gives \(3\le r\le8\), and the proposition
excludes \(r=3\). The conclusion concerns the actual image,
not a conjectural formal-group model with that image size.
\end{proof}

\begin{remark}[Scope of the additional image restriction]
\label{rem:B-RD79-scope}
\status{Proved; complete independent root and companion proof and exact-code audits passed.}
This argument imposes no potential-CM hypothesis. It uses the
terminal kernel of an actual third-level representation and
the actual common division field; it does not identify that
kernel with the entire wild image. Neither any of the orders
\(p^4,\ldots,p^8\), nor the reducible-residual comparisons, is
excluded. In particular the unrestricted first-wild least
index remains in \(\{2,3\}\). The A/C finite-type relation alone
is not treated as an actual quadratic self-twist.


\smallskip
\noindent\emph{Exact supplement and provenance.}
The retained standalone program
\path{code/reduction_descent/verify_first_wild_irreducible_image.py}
checks 216 valuation-denominator cases, 600 layer-dimension
cases, 24 all-binomial-coefficient divisibility rows, 960
higher uniformizer terms, 75,984 \(q^2\)-power field identities,
75,984 distinct scalar candidates for Frobenius, and 88,335
literal matrix \(p\)th-power identities modulo \(p^3\).
The author run passed in \(0.9368512359942542\) seconds.
It does not construct a larger image or a formal-group pair.
No executed computation failed. The full fixed-coefficient
uniformizer and leading-coordinate arguments of B55, and the
actual inverse-polygon argument, were reread rather than
importing a bound from an abstract group order. No new
external-source theorem is asserted; the already freshly read
Stacks coefficient-ring proof supplies the fixed representatives.
The full independent root and companion audits reread the actual
B55, B53 and B60 arguments, with the arbitrary-index leading
action rederived here. Their literal isolated reruns passed in
\(0.9384170099947369\) and \(1.0166636889916845\) seconds,
respectively, preserving the author output. No additional fresh
external-source read is claimed beyond the recorded Stacks
proof used for B78. Exact candidate and audit records are retained.
\end{remark}


% Proved after complete independent audits; the final abelian-image branch remains unresolved.
\subsection{The next wild image and its actual integral Cartan}

Retain the hypotheses of B79: \(p\ge5\), \(q=p^2\), \(N=q-1\),
\(e_K=p(p-1)\), algebraically closed residue field and \(\mu_p\subset K\).
The common residual inertia representation is irreducible, and the
comparison, when used, is an actual \(p^3\)-torsion isomorphism.
Let \(L/K\) have tame degree \(N\), let \(F/L\) be the common full
third-level division field, and suppose \(|W|=|\operatorname{Gal}(F/L)|=p^4\).
Use the actual matrix spaces \(W_1,W_2\) of (B-RD79.2).
For each formal model, the primitive, second-level and first-level
coordinate valuations in \(F\) are respectively
\[
 A=p(p-1),\qquad qA,\qquad q^2A,
                 \qquad v_F(p)=Nq^2A.                 \tag{B-RD80.1}
\]
A tame complement of order \(N\) exists in the actual finite
Galois group: a lift of a tame generator has order \(Np^a\),
and taking its \(p^a\)th power gives an order-\(N\) lift of
another tame generator. Write \(T\) for its action on the residual module,
and identify its centralizer with \(\mathbf F_q\). Denote the
Cartan and anti-Cartan planes in its matrix algebra by
\[
 Z=\{z\mapsto az:a\in\mathbf F_q\},\qquad
 U=\{z\mapsto az^p:a\in\mathbf F_q\}.
\]
Thus tame conjugation fixes \(Z\) and acts on \(U\) through
multiplication by a root of order \(h\mid p+1\), where \(h>1\)
is the residual projective order.

\begin{lemma}[A first-coefficient comparison at either displacement level]
\label{lem:B-RD80-first}
\status{Proved; full independent root and companion proof and exact-code audits passed.}
Fix a Honda coordinate in one model and fixed residue representatives
in a coefficient ring of \(K\). For each \(z\in\mathbf F_q^\times\),
choose one primitive third-level point with leading coefficient \(dz\),
and keep these choices fixed. Let \(g\in W\) have lower break \(b\),
and suppose its displacement on all these points has valuation
\(D\), where \(D=qA\) or \(D=q^2A\). Suppose its leading displacement
map is a nonzero scalar multiple of either \(z\) or \(z^p\).
Put \(b_0=(D-A)/p\).
If \(b<b_0\), the first nonzero coefficient of degree prime to \(p\)
in the fixed point expansions below \(D\) has degree
\[
             j=A+(p-1)b,
 \qquad f_j(z)=Cz\,d_g^{p-1},\qquad C\ne0,           \tag{B-RD80.2}
\]
where \(d_g\) is the leading uniformizer displacement, and the
constant \(C\) depends on the fixed break and model, not on \(g\).
In particular \(p\nmid b\).
If \(b>b_0\), that first degree is \(j=D-b\), and its coefficient
is proportional to the leading displacement map divided by \(d_g\).
These conclusions concern the same fixed point expansions for every
chosen group element; no independence from other lifts modulo \(p\)
is asserted.
\end{lemma}
\begin{proof}
Work through degree \(D\); the coefficient ring is of characteristic
\(p\) in this truncated calculation because \(v_F(p)>D\).
A monomial of degree \(j\) first changes at degree
\(j+p^{v_p(j)}b\). The initial degree is \(A\), with \(v_p(A)=1\).
Among all degrees divisible by \(p\), it is the unique one whose
first change can occur at \(A+pb\); every later such degree changes
strictly later. If \(b<b_0\), this initial change must cancel.
No earlier prime-to-\(p\) degree can occur, since it would change
alone even earlier; none later can cancel. The required first degree
is therefore \(j=A+(p-1)b\), prime to \(p\). The leading cancellation is
\[
       (A/p)dz\,d_g^p+jf_j(z)d_g=0.
\]
Dividing by \(jd_g\) proves (B-RD80.2), including \(p\nmid b\).
If \(b>b_0\), every degree divisible by \(p\) changes after \(D\).
The first prime degree must instead be \(D-b\), and its linear
binomial change gives the stated coefficient. Any impossible
index or zero coefficient would contradict the prescribed nonzero
displacement. These are the full fixed-coefficient calculations of
B56--B57 with a displayed arbitrary target \(D\); they also apply
to a group element of order \(p^2\).
\end{proof}

\begin{proposition}[The three-dimensional terminal layer is excluded]
\label{prop:B-RD80-three-terminal}
\status{Proved; full independent root and companion proof and exact-code audits passed.}
If \(\dim W_1=1\) and \(\dim W_2=3\), two actual compared models
cannot have different finite special characters.
\end{proposition}
\begin{proof}
For a nonzero matrix in either \(Z\) or \(U\), every primitive
point is moved by a nonzero first-level point. Apply the lemma
with \(D=q^2A\) to the corresponding terminal elements.
Here \(b_0=(q^2-1)(p-1)\) is a multiple of \(N\).
Choose an exact tame eigen-uniformizer of \(F\), by averaging
an arbitrary uniformizer. Its eigenvalue \(\eta\) has order \(N\).
On a nonzero ramification-graded piece of lower break \(b\),
tame conjugation acts by \(\eta^b\) (reversing both conventions
has no effect on the following tests). The leading coefficient
map on a subgroup with one common break is injective and additive.

If \(h>2\), the anti-Cartan plane is irreducible over \(\mathbf F_p\).
Any three-dimensional tame-invariant plane consequently contains
all of \(U\), with one common break \(b\); moreover \(\eta^b\ne1\),
so \(b\ne b_0\). If \(b<b_0\), fix any \(z\ne0\) in (B-RD80.2)
and vary the nonzero elements of \(U\). Its two-dimensional
leading-coefficient space would have all its nonzero elements
satisfy the same equation \(X^{p-1}=c\), which has at most
\(p-1\) roots. This contradicts its \(p^2-1\) nonzero elements.
If \(b>b_0\), the high-break formula and one fixed element show
that the displacement parameters \(\beta_i(g)\) of each model
are fixed nonzero multiples of the same \(d_g\). The parameters
run through \(\mathbf F_q\). The Frobenius change between the
two leading orientations, proved at this absolute index in B79,
requires \(\beta_2(g)=a^{1-p}\beta_1(g)^p\) for the fixed
nonzero scalar \(a\) of that intertwiner, again impossible on the
whole field.
This proves the claim for \(h>2\).

Suppose \(h=2\). Tame conjugation is \(-1\) on \(U\), so every
anti-Cartan break satisfies \(b\equiv N/2\pmod N\), whereas
Cartan breaks satisfy \(b\equiv0\pmod N\). If all of \(U\)
is present but has two breaks \(b_1<b_2\), then \(b_1\ne b_0\).
For \(b_1<b_0\), the first prime degree is
\(j=A+(p-1)b_1\), with coefficient a nonzero multiple of \(z\).
Under an element at \(b_2\), its change at \(j+b_2\) is uniquely
first: it precedes \(A+pb_2\) and every later prime degree.
This would give a Cartan-linear displacement, contrary to an
anti-Cartan matrix. For \(b_1>b_0\), the first prime degree is
\(D-b_1\), so all its changes at \(b_2\), and all divisible
ones, occur after \(D\), again impossible. Thus a full \(U\)
has one break, and the previous single-break argument excludes it.

The remaining three-dimensional shape is \(Z\oplus U_1\), with
\(U_1\) an anti-Cartan line. Let \(b\) be its nonzero break.
It is neither a Cartan break nor \(b_0\). If \(b>b_0\), the first
prime coefficient is a multiple of \(z^p\), at \(j=D-b\).
A Cartan break \(c<b_0\) would require that same first coefficient
to be a multiple of \(z\). For \(b_0\le c<b\), the degree-\(j\)
term changes strictly before \(D\), while the divisible terms
do not; for \(c>b\), every change is later than \(D\).
The equality \(c=b\) is forbidden by tame characters. Hence
\(b<b_0\). Its first prime degree is now \(j=A+(p-1)b\), with
coefficient proportional to \(z\). A Cartan break smaller than
\(b\) would force an earlier first prime degree; equality is
again impossible. Each Cartan element must therefore have
\[
                      c=D-j=pb_0-(p-1)b>b_0.
\]
All nonzero elements of \(Z\) have this same break. The high-break
formula then makes their multiplication parameters fixed multiples
of the common \(d_g\) in both models. Frobenius orientation on
this full field again contradicts scalar proportionality. This
exhausts every possible three-dimensional terminal plane.
\end{proof}

\begin{lemma}[A correct integral-lift test for the two-dimensional layers]
\label{lem:B-RD80-plane}
\status{Proved; full independent root and companion proof and exact-code audits passed.}
If \(\dim W_1=\dim W_2=2\), their common matrix plane is a
nonsplit Cartan over \(\mathbf F_p\). The full wild group is
abelian, isomorphic to \((\mathbf Z/p^2)^2\).
\end{lemma}
\begin{proof}
The inclusion from (B-RD79.3) gives \(W_1=W_2=V\).
The exact commutator formula modulo \(p^3\) shows that
\([V,V]\subset V\). Thus \(V\) is a two-dimensional matrix
Lie algebra normalized by the irreducible residual tame matrix.
If it is nonabelian, its characteristic derived line has a generator
\(E\ne0\) and one may choose \(H\) with \([H,E]=E\).
Then \(\operatorname{tr}E=0\), and
\([H,E^2]=2E^2\) gives \(\operatorname{tr}E^2=0\).
Since \(p\ne2\), \(\det E=-\operatorname{tr}(E^2)/2=0\).
So \(E\) is a nonzero nilpotent and its kernel is an
\(\mathbf F_p\)-line preserved by the tame normalizer,
contradicting irreducibility.

The plane is therefore abelian. Choose a nonscalar element.
Its two-dimensional centralizer equals \(V\), and is either
a nonsplit Cartan, a split Cartan, or the dual-number algebra.
In the last case the unique nilpotent line and its kernel are
again preserved by the tame normalizer, a contradiction.

For the split Cartan use an eigenbasis so that \(V\) is the
entire diagonal plane. Choose an actual element
\(g=1+pE_{22}+p^2B\in W\). Multiplication by an actual terminal
diagonal matrix permits \(B_{11}=0\) modulo \(p\).
Choose a primitive vector
\[
                 v=e_1+p w,\qquad w_2=-B_{21}\pmod p.
                                                     \tag{B-RD80.3}
\]
Then \((E_{22}+pB)v=0\pmod{p^2}\), so \(gv=v\pmod{p^3}\).
But \(g^p=1+p^2E_{22}\ne1\), so \(g\) has order \(p^2\).
Its stabilizer on this actual primitive point has order at least
\(p^2\), giving an orbit of size at most \(p^2\). This contradicts
its degree-\(p^3\) valuation denominator from B79.
The choice of the integral lift in (B-RD80.3) is essential: no
bound by a sum of residual evaluation ranks is used.

Only a nonsplit Cartan remains. Its matrices commute, so the
actual commutator formula gives \([W,W]=1\) modulo \(p^3\).
Taking \(p\)th powers maps onto the terminal plane of order
\(p^2\). An abelian group of order \(p^4\), exponent at most
\(p^2\), with a \(p\)th-power subgroup of order \(p^2\), is
\((\mathbf Z/p^2)^2\), as asserted.
\end{proof}

\begin{lemma}[The terminal Cartan is the original residual Cartan]
\label{lem:B-RD80-orientation}
\status{Proved; full independent root and companion proof and exact-code audits passed.}
In the situation of the preceding lemma, \(V\) equals the
centralizer of the original irreducible residual tame matrix.
\end{lemma}
\begin{proof}
Write \(E'\) for the nonsplit Cartan \(V\). The normalizer of
\(E'\) is \(E'^\times\rtimes\operatorname{Gal}(E'/\mathbf F_p)\):
conjugation restricts to a field automorphism, and its kernel is
field multiplication. If the tame matrix lies in \(E'^\times\),
its irreducibility identifies its centralizer with \(E'\), as desired.
Suppose instead it lies in the other coset. Its square is scalar,
and its conjugation on \(E'\) has eigenspaces \(\mathbf F_p I\)
and a trace-zero line. Relative to the original residual field,
these are a Cartan scalar line and an anti-Cartan line.

The actual tame action on \(W\simeq(\mathbf Z/p^2)^2\) has order
prime to \(p\). Its square is the identity modulo \(pW\).
The reduction kernel in its automorphism group is a \(p\)-group,
so its square is the identity on \(W\) itself. Since two is a
unit, the exact \(\pm1\) projectors give two rank-one summands.
Choose generators \(g_+\) and \(g_-\), each of order \(p^2\),
with tame conjugation respectively \(g_+\mapsto g_+\) and
\(g_-\mapsto g_-^{-1}\). Their first matrices are respectively
nonzero scalar and anti-Cartan matrices, both invertible.
Their displacements on every primitive point have valuation \(qA\).

Apply the first-coefficient lemma with \(D=qA\) and
\(a_0=(q-1)(p-1)\). Write \(a\) for the lower break of \(g_-\)
and \(c\) for that of \(g_+\). Tame characters give
\(a\equiv N/2\pmod N\), \(c\equiv0\pmod N\).
Thus \(a\ne c\) and \(a\ne a_0\).
If \(a>a_0\), the first prime coefficient has degree \(qA-a\)
and is a nonzero multiple of \(z^p\). A break \(c<a_0\)
would force that first coefficient to be a multiple of \(z\);
for \(a_0\le c<a\) it would change strictly before \(qA\),
with no divisible term available to cancel; for \(c>a\) all
changes would occur later than \(qA\). Thus \(a<a_0\).
The first prime degree is
\(j=A+(p-1)a\), and its coefficient is a nonzero multiple of \(z\).

Now use the actual element \(g_-^p\), whose nonzero terminal
matrix is the same anti-Cartan matrix. Its lower break \(b\)
is strictly greater than \(a\): the first ramification coefficient
of a product is additive, so the first coefficient of the \(p\)th
power vanishes in the residue field. Its displacement must have
valuation \(q^2A\) with leading coefficient proportional to \(z^p\).
But the already fixed degree-\(j\) coefficient changes at
\(j+b<A+pb\), before every divisible term and every later
prime term. It is uniquely first and is proportional to \(z\).
If its degree differs from \(q^2A\), valuations already contradict;
if it equals \(q^2A\), the independence of \(z\) and \(z^p\) as
functions on \(\mathbf F_q\) contradicts the required leading map.
This excludes the noncentral normalizer case and proves the lemma.
\end{proof}

\begin{theorem}[The fourth-order wild branch has an actual abelian image]
\label{thm:B-RD80}
\status{Proved; full independent root and companion proof and exact-code audits passed.}
If the two finite special characters differ and their common
third-level wild image has order \(p^4\), then
\[
 W_1=W_2=Z,
 \qquad W\simeq(\mathbf Z/p^2)^2,
                                                     \tag{B-RD80.4}
\]
and the full common third-level inertia image is abelian.
More precisely its Tate module modulo \(p^3\) is free of rank one
over the unramified quadratic algebra \(\mathcal O_E/p^3\),
its full image acts by multiplication, and its wild image is exactly
\[
                 (1+p\mathcal O_E)/(1+p^3\mathcal O_E).
                                                     \tag{B-RD80.5}
\]
This is a statement about the finite-level action, not a CM
endomorphism assertion for the full BT group.
\end{theorem}
\begin{proof}
As in B79, \(W_1\ne0\), \(W_1\subset W_2\), and their dimensions
sum to four. The only possibilities are \((1,3)\) and \((2,2)\).
The former is excluded for a different-character pair by the
three-dimensional terminal proposition. The next two lemmas give
(B-RD80.4) in the latter case.

The tame complement acts trivially on \(W/pW=W_1=Z\).
Its finite prime-to-\(p\) action on \(W\simeq(\mathbf Z/p^2)^2\)
therefore is trivial, again because the reduction kernel in its
automorphism group is a \(p\)-group. Thus the complement centralizes
\(W\), and the full finite inertia image is abelian.

Let \(\widetilde T\) be the actual third-level tame matrix.
Its residual quadratic polynomial is irreducible and separable.
The algebra it generates over \(\mathbf Z/p^3\) is therefore the
unramified quadratic algebra \(\mathcal O_E/p^3\), and the free
rank-two underlying module is free of rank one over that algebra.
Every matrix commuting with \(\widetilde T\) is multiplication
by an element of it, as follows either from the basis
\(v,\widetilde T v\) or the centralizer equations.
The entire image commutes with \(\widetilde T\), so it lies in
these units. Its wild subgroup lies in their principal units,
whose order is \(p^4\). Equality of orders proves (B-RD80.5).
\end{proof}

\begin{remark}[What the finite Cartan does not prove]
\label{rem:B-RD80-scope}
\status{Proved; full independent root and companion proof and exact-code audits passed.}
This narrows the \(p^4\) branch left by B79 to a full finite
Cartan image. It neither excludes that branch nor identifies a
finite-level quadratic action with finite-extension endomorphisms
of the whole group. B78's CM theorem therefore cannot be applied
without an additional argument. Wild orders \(p^5,\ldots,p^8\)
and reducible residual representations remain open as well.

The earlier exploratory orbit-rank shortcut was withdrawn and its
exact note retained. The split-Cartan argument here instead constructs
an actual primitive lift with an order-\(p^2\) stabilizer. No assertion
that all fibre motions come from \(W_2v\) is used.


\smallskip
\noindent\emph{Retained exact supplement.}
The standalone program
\path{code/reduction_descent/verify_first_wild_fourth_layer.py}
enumerates all 19,882 matrix planes in reduced row echelon form at
\(p=5,7,11\), identifying 439 Lie-closed planes and the 29
normalized by the retained irreducible tame matrices. It checks
all 17,667 split-Cartan lift coefficient choices at those primes,
29 actual commutator identities, 24 layer-shape rows, 240
first-prime-degree cases, 768 later-break comparisons, 1,108
anti-lift power clashes, 1,180 quadratic-field power fibres,
and 91,728 exact anti-Cartan conjugacy identities.
The expanded author run passed in \(0.9882168779877247\) seconds.
The initial executed program and output are retained; no executed
computation failed. The anti-Cartan conjugation formula explicitly
retains its fixed scalar \(a^{1-p}\), and the earlier proof wording
without that scalar is preserved. That scalar does not change the
Frobenius-versus-scalar contradiction.
The full B56 and B57 coefficient arguments were reread and expanded
at the actual target displacement level; no new original-source read
is claimed. These finite algebra checks supplement the written
all-prime argument and establish no BT-group existence.

The independent root and companion audits read the entire final proof
and certificate and reread the actual B79, B57 and B60 dependency
arguments. Their unchanged isolated reruns passed in
\(0.983579940002528\) and \(1.0418004910025047\) seconds,
respectively, preserving the author output. These audits make no
additional fresh original-source claim. The exact audited candidate,
initial and corrected programs and outputs, and the withdrawn
orbit-rank shortcut remain retained with the promotion record.
\end{remark}


% Proved after complete independent audits; larger and reducible wild images remain open.
\subsection{Excluding the actual fourth-order Cartan comparison}

Throughout this subsection, \(K\) is complete discretely valued of
mixed characteristic \((0,p)\), its residue field \(k\) is
algebraically closed, \(p\ge5\), \(e_K=p(p-1)\), and
\(\mu_p\subset K\). All comparisons concern actual generic
\(p^3\)-torsion of potentially supersingular height-two groups.
Put \(q=p^2\), \(N=q-1\), and use the covariant special-character
convention of B53.

\begin{lemma}[The two actual Cartan displacement levels]
\label{lem:B-RD81-levels}
\status{Proved; full independent root and monodromy proof and exact-code audits passed.}
Suppose that two such groups have an actual \(p^3\)-torsion
isomorphism, irreducible residual inertia, and distinct finite
special characters. If their common wild third-level image has
order \(p^4\), use the common fields
\[
 L/K\text{ of tame degree }N,\qquad
 F=L(M_1[p^3])=L(M_2[p^3]).
\]
Then \(W=\operatorname{Gal}(F/L)\simeq(\mathbf Z/p^2)^2\),
\(U=pW\) is its terminal kernel of order \(q\), and both actual
matrix layers are the original nonsplit Cartan. For
\[
 A=p(p-1),\qquad D_1=qA,\qquad D_2=q^2A,
 \qquad v_F(p)=ND_2,                                  \tag{B-RD81.1}
\]
each primitive coordinate has valuation \(A\); its displacement
by \(g\in W\setminus U\) has valuation \(D_1\), and by
\(h\in U\setminus\{1\}\) has valuation \(D_2\).

Let \(a\) be the minimum lower break of nonidentity elements
of \(W\), with the convention
\(b(g)=v_F(g\varpi-\varpi)-1\). Every nonidentity element
of \(U\) has break strictly greater than \(a\), so a
minimum-break element lies outside \(U\).
\end{lemma}
\begin{proof}
The actual Cartan conclusion, including \(W=(\mathbf Z/p^2)^2\)
and \(pW=U\), is Theorem~\ref{thm:B-RD80}. The valuation
statements are (B-RD80.1). A nonzero Cartan coefficient is an
invertible multiplication operator on the residual two-dimensional
space. Hence a nonterminal element moves every primitive point
by a point of exact order \(p^2\), and a nonidentity terminal
element moves it by a nonzero point of order \(p\).

For completeness the strict-break assertion uses the actual
group, not its associated graded dimensions. If \(g\) has
break \(b\), the coefficient
\[
 d_g=\overline{(g\varpi-\varpi)/\varpi^{b+1}}
\]
adds under composition of elements of break at least \(b\).
Every such element fixes the residue field, so its \(p\)th
power has zero coefficient in that degree. Thus
\(b(g^p)>b(g)\) whenever \(g^p\ne1\). Every nonidentity
\(h\in U=pW\) is the \(p\)th power of an order-\(p^2\)
element of \(W\), which proves \(b(h)>a\).
\end{proof}

\begin{proposition}[The full finite Cartan does not conceal a different character]
\label{prop:B-RD81-Cartan}
\status{Proved; full independent root and monodromy proof and exact-code audits passed.}
Under the field hypotheses above, an actual \(p^3\)-congruent
pair with irreducible residual inertia and wild image of order
\(p^4\) has equal finite special characters.
\end{proposition}
\begin{proof}
Suppose the characters differ and apply the preceding lemma.
Fix one uniformizer \(\varpi\) of the actual common field \(F\).
For each model fix Honda coordinates and, once and for all,
primitive point lifts indexed by \(z\in\mathbf F_q^\times\).
Their leading coordinates are \(d_i z\), where \(d_i\in k^\times\).
The coordinate expansions are taken using residue representatives
in a coefficient ring of \(K\), and are written
\[
 x_i(z)=\sum_{r\ge A}[f_{i,r}(z)]\varpi^r,
 \qquad f_{i,A}(z)=d_i z.                             \tag{B-RD81.2}
\]
The representatives are fixed by \(W\). Since \(v_F(p)>D_2\),
all terms relevant up to displacement valuation \(D_2\) are
calculated in characteristic \(p\), exactly as in
Lemma~\ref{lem:B-RD80-first}. These are fixed point lifts
throughout the argument; no independence of higher coefficients
from the choice of lift is asserted.

Write \(\lambda_i(g)\in\mathbf F_q^\times\) for the first
matrix coefficient of \(g\in W\setminus U\), and
\(\lambda_i(h)\in\mathbf F_q^\times\) for the terminal
coefficient of \(h\in U\setminus\{1\}\). The corresponding
coefficient maps on \(W/U\) and \(U\) are additive isomorphisms
onto \(\mathbf F_q\). The actual leading displacement
coefficients in the two levels are
\[
 d_i^q\lambda_i(g)z,
 \qquad d_i^{q^2}\lambda_i(h)z.                       \tag{B-RD81.3}
\]
Indeed the leading \(p\)-series at a primitive point is
\(T^q\), uniquely, and multiplication by \(p^2\) raises
leading coordinates to the \(q^2\)th power. Normalized field
coordinates satisfy \(z^q=z\). Since the two special
characters differ, the actual intertwiner has Frobenius
orientation on these leading spaces, as proved in B79.
Consequently, on both matrix layers,
\[
                   \lambda_2=\lambda_1^p.             \tag{B-RD81.4}
\]
The possible scalar in the intertwiner cancels in conjugating
Cartan multiplication operators.

Set
\[
 a_0=\frac{D_1-A}{p}=(q-1)(p-1),\qquad
 j_0=A+(p-1)a_0=D_1-a_0.                             \tag{B-RD81.5}
\]
Here \(a_0\equiv1\pmod p\), and hence
\(j_0\equiv-1\pmod p\). In particular the boundary degree
\(j_0\) is never divisible by \(p\).
Choose a minimum-break element \(g\in W\setminus U\), of
break \(a\). The degree-\(A\) term of (B-RD81.2) first
changes at degree \(A+pa\), with coefficient
\((A/p)d_i z d_g^p\). Any other degree divisible by \(p\)
changes no earlier than \(A+pa\). The first nonzero degree
\(j\) prime to \(p\) changes at \(j+a\), with coefficient
\(j f_{i,j}(z)d_g\); it is the unique earliest prime-degree
change.

First suppose \(a<a_0\). Cancellation below \(D_1\)
forces the same first prime degree in both models,
\[
                  j=A+(p-1)a,
          \qquad f_{i,j}(z)=c_i z,
          \qquad c_i\ne0.                            \tag{B-RD81.6}
\]
This is the low-break conclusion of
Lemma~\ref{lem:B-RD80-first}; it also forces \(p\nmid a\).
If instead \(a>a_0\), the first change at \(D_1\) has to
come from a prime degree, and (B-RD81.3) gives
\[
                 j=D_1-a,
          \qquad f_{i,j}(z)=c_i z,
          \qquad c_i\ne0.                            \tag{B-RD81.7}
\]
If that degree is not a positive permitted prime degree, the
case is already impossible.

In either case every nonidentity \(h\in U\) has break
\(b>a\), and this fixed first prime term is uniquely the
first term changed by \(h\). To check the strict comparison,
the difference \(A+pb-(j+b)\) equals respectively
\[
 (p-1)(b-a),\qquad
 (p-1)(b-a)+p(a-a_0),                                \tag{B-RD81.8}
\]
both positive. All later prime terms change later as well.
Thus \(j+b=D_2\) for every nonidentity terminal element.
Its leading coefficient yields
\[
      d_i^{q^2}\lambda_i(h)z=j c_i z d_h,
      \qquad \lambda_i(h)=\kappa_i d_h,
      \qquad \kappa_i\ne0.                           \tag{B-RD81.9}
\]
The common lower break makes \(d_h\) the same actual
uniformizer coefficient for the two models. Since
\(\lambda_1(U)=\mathbf F_q\), equations (B-RD81.4) and
(B-RD81.9) would make Frobenius on \(\mathbf F_q\) a fixed
scalar map. Evaluating at \(1\) and at an element outside
\(\mathbf F_p\) is a contradiction. Therefore \(a=a_0\).

We now treat the boundary, including a vanishing first prime
coefficient. No prime degree below \(j_0\) can occur in
either expansion: its change would precede \(D_1\), whereas
all divisible degrees first change at least at \(D_1\).
At degree \(D_1\) the only possible contributions are
the initial term and the degree-\(j_0\) term. Thus, for
the fixed minimum-break element,
\[
 d_i^q\lambda_i(g)z
   =(A/p)d_i z d_g^p+j_0 f_{i,j_0}(z)d_g.             \tag{B-RD81.10}
\]
Since \(j_0d_g\ne0\), this proves
\(f_{i,j_0}(z)=c_i z\), with \(c_i\) possibly zero.

If \(c_i\ne0\) for at least one model, then for every
terminal element its break \(b>a_0\) makes this term
uniquely first, because
\(A+pb-(j_0+b)=(p-1)(b-a_0)>0\). Hence every such
break is \(b_*=D_2-j_0\). In the other model the same
actual terminal element must also have a nonzero first
prime coefficient at degree \(j_0\): its divisible terms
all change after \(D_2\), and an earlier or later first
prime term would give a displacement valuation different
from \(D_2\). Equivalently,
\[
 b_*-\frac{D_2-A}{p}
       =\frac{p-1}{p}(D_2-D_1)>0,                    \tag{B-RD81.11}
\]
so the high-break clause of Lemma~\ref{lem:B-RD80-first}
forces that exact nonzero degree. We obtain (B-RD81.9)
for both models and the same contradiction. It follows
that \(c_1=c_2=0\).

There is consequently no prime degree at or below \(j_0\)
in either expansion. No nonterminal element can now have
break \(b>a_0\): all divisible terms would change after
\(A+pa_0=D_1\), and all prime terms after
\(j_0+a_0=D_1\). This contradicts its actual displacement
valuation \(D_1\). Every element of \(W\setminus U\)
therefore has break \(a_0\), and every nonidentity element
of \(U\) has greater break. The leading ramification map
is an injective additive map \(W/U\to k\) with image of
dimension two over \(\mathbf F_p\).

For all nonterminal elements equation (B-RD81.10) now has
only its first term. Consequently
\[
       \lambda_i(g)=(A/p)d_i^{1-q}d_g^p
                   \quad(g\in W\setminus U).          \tag{B-RD81.12}
\]
The two coefficient maps are again fixed scalar multiples
of one another, while (B-RD81.4) makes one the Frobenius
of the other. Their values fill \(\mathbf F_q\), so this
is impossible. The contradiction proves the proposition.
\end{proof}

\begin{corollary}[The first two possible wild orders are excluded]
\label{cor:B-RD81}
\status{Proved; full independent root and monodromy proof and exact-code audits passed.}
At \(e=p(p-1)\), an actual \(p^3\)-congruent pair with
irreducible residual inertia and different finite special
characters must have wild third-level image of order \(p^r\)
with \(5\le r\le8\). In particular an irreducible-residual
counterexample at \(p=5,e=20\) would have actual
\(125\)-torsion wild image of order at least \(3125\).
\end{corollary}
\begin{proof}
B79 excludes order \(p^3\), and B80 together with the
preceding proposition excludes order \(p^4\). The upper
bound \(p^8\) is the actual matrix-kernel bound of B79.
\end{proof}

\begin{remark}[Exact scope of the Cartan exclusion]
\label{rem:B-RD81-scope}
\status{Proved; full independent root and monodromy proof and exact-code audits passed.}
The argument excludes an actual finite-level Cartan comparison
using its point coordinates and its exact \(p\)th-power map.
It does not promote a finite-level quadratic action to a
finite-extension endomorphism of the whole group and does
not invoke the potentially CM theorem for a non-CM model.
Wild image orders \(p^5,\ldots,p^8\) and reducible residual
comparisons remain open. Thus the full first-wild least
index still lies in \(\{2,3\}\), with no counterexample or
unrestricted \(125\)-recovery assertion.

\smallskip
\noindent\emph{Exact supplement and provenance.}
The standalone program
\path{code/reduction_descent/verify_first_wild_Cartan_exclusion.py}
verifies eight symbolic identities, 24 initial-binomial rows,
96 vanishing-boundary comparisons, 14,109 exhaustive permitted
first-break rows and 1,584 impossible prime-degree rows at
\(p=5,7,11\), together with 56,436 later-break comparisons.
It checks all 17,667 principal-unit \(p\)th powers in the
quadratic rings modulo \(p^3\) at those three primes and
their three exact terminal-image fibre counts. Its 75,984
scalar tests use the two field basis vectors and every
prime from 5 through 101. The author run passed in
\(0.4355178660043748\) seconds, with no failed execution.
All executed source and output are retained. These finite
checks supplement the complete uniform proof; they construct
no formal-group comparison. The full B80 coefficient argument
and the B79 actual displacement and orientation argument were
reread. No new original-source read is claimed.
The complete independent root and monodromy audits passed,
including every boundary coefficient and the actual Cartan
dependencies. Their unchanged isolated certificate reruns took
\(0.45372015300381463\) and \(0.4650933469965821\) seconds,
respectively, preserving the author output. The exact audited
candidate and all proof/code audit records are retained.
\end{remark}


% Proved after CP29 archive release; exact audited candidates retained.
\subsection{Actual terminal matrices for a reducible first-wild comparison}

Let \(p\ge5\), \(q=p^2\), \(N=q-1\), and let \(K\) be complete
discretely valued of mixed characteristic \((0,p)\), with
algebraically closed residue field, \(\mu_p\subset K\), and
\(e=e_K=p(p-1)\). Put \(s=(p+1)e\) and take the common
tame good field \(L/K\) of degree \(N\). All groups below
are whole potentially supersingular height-two groups; the
comparisons are of their actual generic \(p^3\)-torsion.

\begin{lemma}[Unequal slopes at the first wild index]
\label{lem:B-RD82-Newton}
\status{Proved; full independent root and companion proof and exact-code audits passed.}
If the first-level Newton polygon has two slopes, there is a
canonical high line \(H\) and an integer \(B\) such that
\[
 \begin{gathered}
 0<B<pe,\qquad \alpha=s-B,\qquad \beta=B/p,\\
 v_L(T(Q))=\alpha\ (0\ne Q\in H),\qquad
 v_L(T(Q))=\beta\ (Q\notin H).
 \end{gathered}                                      \tag{B-RD82.1}
\]
The high line is \(L\)-rational. If \(c\) is the finite
special tangent exponent, its actual tame character is
\(\zeta^{c-B}\), where \(\zeta\) is the generator chosen
on a tame uniformizer of \(L\). A primitive third-level
point with \(p^2P\notin H\) has
\[
                    v_L(T(P))=B/p^5.                 \tag{B-RD82.2}
\]
The residual division field over \(L\) has degree one or
\(p\). If it has degree \(p\), its unique wild fixed line
is \(H\).
\end{lemma}
\begin{proof}
Here the coefficient proof of B60 and B69 is reapplied at the
present index, rather than importing their prime-to-\(p\)
field scope. In the actual universal height-two coordinate,
the only possible negative Newton vertices of \([p]\) are
\((1,eN),(p,u_0),(q,0)\). The coefficient estimates are
those proved in Lemma~\ref{lem:B-RD49-point}. The unequal
polygon has \(p-1\) points at valuation
\((eN-u_0)/(p-1)\) and \(q-p\) at valuation
\(u_0/(p(p-1))\). The first points and zero form a
unique inertia-stable line. Its character is
\(\mathbf F_p^\times\)-valued and hence is killed by \(L\).
Its coordinate valuation is therefore integral. As
\(eN=(p-1)s\), this gives \(u_0=(p-1)B\), with the
range and values (B-RD82.1).

Use an integral Honda framing and tame averaging as in
Lemma~\ref{lem:B-RD67-normalization}. The leading action
on a high point is \(\zeta^{c+\alpha}=\zeta^{c-B}\),
because \(s\) is a multiple of \(N\). This is an actual
line character, and therefore remains identified by any
torsion isomorphism identifying the two high lines.

At a low target of valuation \(B/p\), the inverse polygon
has one segment from \((0,B/p)\) to \((q,0)\).
Indeed \(u_0=(p-1)B>B/p\), the linear coefficient has
valuation \(eN>B/p\), and the intervening coefficients
have the same lower bounds as in B60. Every inverse has
valuation \(B/p^3\). The still smaller target of that
valuation gives the same one-segment argument and
valuation \(B/p^5\), proving (B-RD82.2).

The residual high line and residual determinant are trivial
over \(L\): the latter is cyclotomic, with \(\mu_p\subset L\).
Thus the residual image over \(L\) lies in the upper
unipotent group of order \(p\). If nontrivial, its unique
fixed line is the already invariant high line.
\end{proof}

\begin{lemma}[An actual terminal plane forced by a prime point valuation]
\label{lem:B-RD82-terminal}
\status{Proved; full independent root and companion proof and exact-code audits passed.}
Suppose one such group has two slopes with \(B=pb\),
\(p\nmid b\), and its full third-level division field
\(F=L(M[p^3])\) has \([F:L]=p^4\). In matrix coordinates
with \(H=\mathbf F_p e_1\), its terminal matrix space
\[
 V_2=\{C\in M_2(\mathbf F_p):1+p^2C\in
                \operatorname{Gal}(F/L)\}
\]
has dimension two. It contains \(E_{12}\) and a matrix
whose induced coefficient on the quotient by \(H\) is
nonzero. Every element of \(V_2\) is upper triangular,
and every element with zero lower diagonal coefficient
is a multiple of \(E_{12}\).
\end{lemma}
\begin{proof}
Write \(W=\operatorname{Gal}(F/L)\). Its residual image
has order \(p^\epsilon\), with \(\epsilon\in\{0,1\}\).
By (B-RD82.2), every primitive point whose residual class
is outside \(H\) has integral \(F\)-valuation \(b\),
which is prime to \(p\).

For a nonidentity terminal element \(g=1+p^2C\), the
fixed-coefficient uniformizer argument of
Lemma~\ref{lem:B-RD60-break} gives
\[
             v_F(gT(P)-T(P))=b+b(g)                 \tag{B-RD82.3}
\]
for every such primitive point. In particular \(Cv\ne0\)
for every residual vector \(v\notin H\). Moreover all
these images \(Cv\) have the same high or low status:
their displacement valuations are either \(p^4\alpha\)
or \(p^4b\), and (B-RD82.3) is independent of \(v\).
Formal subtraction and coordinate subtraction have the
same valuation, since their ratio is an integral unit.

Write the lower row of \(C\) as \((c,d)\). If \(c\ne0\),
the vectors \((x,1)\) include one with zero lower image
and one with nonzero lower image. The first image cannot
be zero, and hence would be high, while the other is low.
This contradicts the common status just proved. Thus
\(c=0\). If \(d=0\), the upper row \((a,b_0)\) cannot
have \(a\ne0\), since it would then kill one vector
\((x,1)\). Consequently \(C=b_0E_{12}\).
The kernel of the linear map \(V_2\to\mathbf F_p\)
taking the lower diagonal coefficient is therefore
contained in \(\mathbf F_pE_{12}\), so \(t=\dim V_2\le2\).

Let \(V_1\) be the image of
\(\ker(W\to\GL_2(\mathbf F_p))\) under the first
congruence coefficient, and put \(r_1=\dim V_1\).
The actual matrix identity
\((1+pA+p^2B)^p=1+p^2A\pmod{p^3}\) gives
\(V_1\subset V_2\). The exact group-order filtration
therefore reads
\[
             4=\epsilon+r_1+t,
                    \qquad r_1\le t\le2.             \tag{B-RD82.4}
\]
It forces \(t=2\), with \(r_1=2\) for \(\epsilon=0\)
and \(r_1=1\) for \(\epsilon=1\). The lower-diagonal
map is now surjective and its one-dimensional kernel
is exactly \(\mathbf F_pE_{12}\). This proves the
assertion using actual terminal elements; no bound on
integral orbit sizes is inferred from residual evaluation ranks.
\end{proof}

\begin{proposition}[Recovering the high line from two terminal breaks]
\label{prop:B-RD82-recovery}
\status{Proved; full independent root and companion proof and exact-code audits passed.}
Let two groups with two-slope first-level polygons have
isomorphic actual \(p^3\)-torsion and common division
field \(F/L\) of degree \(p^4\). If one parameter is
\(B_1=pb_1\) with \(p\nmid b_1\), then their high lines
agree under the actual isomorphism, \(B_1=B_2\), and
their finite special characters agree.
\end{proposition}
\begin{proof}
Choose a residual vector outside both high lines. Its
primitive lifts in model 2 have valuation \(B_2/p^5\)
over \(L\). Integrality in the actual common field of
degree \(p^4\) implies \(p\mid B_2\). Write
\(B_2=pb_2\); both integers satisfy \(0<b_i<e\).
The preceding lemma supplies an actual terminal element
with matrix \(E_{12}\) relative to \(H_1\). For every
vector outside \(H_1\), its displacement in model 1 is
high, so its exact lower break is
\[
                  r_N=p^4(s-pb_1)-b_1.               \tag{B-RD82.5}
\]
If \(H_1\ne H_2\), choose the residual vector outside
both. Its nonzero image under \(E_{12}\) lies in
\(H_1\setminus H_2\), so the corresponding displacement
in model 2 is low. The universal lower bound, which
does not require \(p\nmid b_2\), would give
\[
       p^4b_2\ge b_2+p^4(s-pb_1)-b_1.                \tag{B-RD82.6}
\]
This is impossible: \(s-pb_1>e\), whereas both
\(b_1,b_2<e\). More explicitly the left side minus
the right side is increasing in both integers and,
at \(b_1=b_2=e-1\), is already \(-p^4(p+1)<0\).
Thus \(H_1=H_2\).

Take now an actual terminal matrix having nonzero lower
diagonal coefficient. It sends every residual vector
outside the common line to a low vector in both models.
Its exact break in model 1 is \((p^4-1)b_1\).
The universal inequality in model 2 implies
\[
        (p^4-1)(b_2-b_1)\ge0.                        \tag{B-RD82.7}
\]
For the nilpotent terminal element (B-RD82.5), both
displacements are high; the same inequality instead gives
\[
                    (p^5+1)(b_1-b_2)\ge0.            \tag{B-RD82.8}
\]
Hence \(b_1=b_2\). The actual high-line characters
\(\zeta^{c_i-B_i}\) agree under the given isomorphism.
Equality of \(B_i\) therefore identifies \(c_1=c_2\)
modulo \(N\), proving the special-character assertion.
\end{proof}

\begin{theorem}[A lower image bound also for nontrivial residual wild inertia]
\label{thm:B-RD82}
\status{Proved; full independent root and companion proof and exact-code audits passed.}
At the stated first wild index, an actual \(p^3\)-congruent
pair with nontrivial residual wild inertia and different
finite special characters has full wild image of order
\(p^r\), with \(5\le r\le9\). In particular its wild
image has order at least \(3125\) when \(p=5,e=20\).
\end{theorem}
\begin{proof}
Their residual wild image over \(L\) has order \(p\).
The remaining two congruence kernels have dimensions at
most four each, giving \(r\le9\). A one-slope first-level
polygon would have all nonzero coordinates at valuation
\(e\), and its injective leading-coordinate map over \(L\)
would make the residual wild action trivial. Thus both
polygons have two slopes and the actual isomorphism
identifies their unique wild fixed high lines.

Suppose \(r\le4\). Equation (B-RD82.2) and integrality
force \(B_i=pb_i\). If one \(b_i\) is prime to \(p\),
the same valuation denominator forces \(r\ge4\), so
\(r=4\); the preceding proposition gives equal characters.

It remains that \(B_i=p^2k_i\). Their range gives
\(1\le k_i\le p-2\). The finite determinant characters
agree by the actual torsion determinant comparison: both
generic determinants are the same cyclotomic character
times a finite \(\mu_{p-1}\)-valued character, and reduction
is injective on these prime-to-\(p\) roots. Therefore
\(c_1\equiv c_2\pmod{p-1}\). The common high-line character
gives \(c_1-B_1\equiv c_2-B_2\pmod N\). Reducing this
congruence modulo \(p-1\), and using \(p^2\equiv1\)
there, proves
\[
                 k_1\equiv k_2\pmod{p-1}.             \tag{B-RD82.9}
\]
The indicated integer interval forces \(k_1=k_2\);
the same line-character equality gives \(c_1=c_2\pmod N\).
This excludes every case with \(r\le4\).
\end{proof}

\begin{remark}[The residual-tame exceptional sector is not removed]
\label{rem:B-RD82-scope}
\status{Proved; full independent root and companion proof and exact-code audits passed.}
The terminal-plane proposition also applies when the residual
representation is tame, provided that its hypotheses on the
two polygons, the common field degree, and one prime-to-\(p\)
integer \(b_i\) hold. The final theorem uses the unique
residual wild fixed line to handle the remaining case with
both \(b_i\) divisible by \(p\). It does not exclude a
residual-tame comparison with two such parameters and
exchanged high lines, nor larger wild images in any branch.
It makes no existence or unrestricted \(p^3\)-recovery
assertion, and the full first-wild index remains in \(\{2,3\}\).

\smallskip
\noindent\emph{Exact supplement and provenance.}
The retained standalone program
\path{code/reduction_descent/verify_first_wild_reducible_terminal.py}
checks three symbolic break identities, all 17,664 nonzero
matrices and 1,723,680 outside-line images at \(p=5,7,11\),
and the 23 possible maximal terminal planes at those primes.
It checks 24 layer-shape rows, 71,456 small divisible-parameter
character comparisons, 3,628 field-integrality rows, and
35,000 rows each for high-line mismatch and two-break recovery.
The author run passed in \(0.42271842100308277\) seconds.
No executed computation failed; all source and output are
retained. These finite checks do not construct a formal
group or replace the uniform arguments above. The full
internal B60 fixed-coefficient proof, B69 actual Newton
and line-character arguments, and B49 coefficient bounds
were reread. No new original-source read is claimed.
The complete independent root and companion audits passed;
their author-output-preserving reruns took
\(0.3949265940027544\) and \(0.39488881299621426\) seconds,
respectively. Both auditors also checked the exact final
determinant simplification, which removes an unnecessary
square-character case split and finite twist without changing
the statement or certificate. Both source versions, the exact
diff, and all audit records remain retained. No additional
source reading or redundant rerun is claimed for that delta.
\end{remark}


% Separate post-CP29 candidate; not part of the frozen paper or current gate.
\subsection{Two opposite actual quotients and the first wild matrix layer}

Let \(K\) be complete discretely valued of mixed characteristic
\((0,p)\), with algebraically closed residue field,
\(p\ge5\), \(e_K=e=p(p-1)\), and \(\mu_p\subset K\).
Put \(q=p^2\), \(N=q-1\), \(s=(p+1)e=pN\), and use
the common tame good field \(L/K\) of degree \(N\).
All groups considered are height-two Barsotti--Tate groups
acquiring one-dimensional supersingular formal good reduction.
All torsion comparisons and quotients below are actual ones.

\begin{lemma}[The tame second-level point-field test at this index]
\label{lem:B-RD83-tame}
\status{Proved; full independent root and companion proof and exact-code audits passed.}
Suppose two good formal height-two groups over \(L\) have
two-slope first-level polygons with parameters \(B_i=pb_i\),
where \(p\nmid b_i\), and have isomorphic actual
\(p^2\)-torsion. If both first-level modules are
\(L\)-rational, then \(b_1=b_2\).
\end{lemma}
\begin{proof}
This is the common-point-field proof underlying B69, written
again without its prime-to-\(p\) absolute-index assumption.
Choose corresponding residual vectors outside the two
canonical high lines, and corresponding primitive
second-level lifts \(P_i\). The actual inverse Newton
polygon gives \(v_L(T_i(P_i))=b_i/q\). The torsion
isomorphism identifies their stabilizers, hence their
point fields \(E=L(P_1)=L(P_2)\).

Because the first-level modules are \(L\)-rational and
\(pP_i\) is rational, each point field is Galois, with
its group a subgroup of the first-level translation
group. Indeed every conjugate is \(P_i+Q\) for an
\(L\)-rational \(Q\), so it is defined over the same
point field. Thus \([E:L]\le q\). The reduced denominator
of \(b_i/q\) is \(q\), so \([E:L]=q\), and every
first-level translation occurs.

Choose a translation outside both high lines. With the
integral valuation on \(E\), the two primitive coordinates
have valuations \(b_i\), and their low displacements have
valuations \(qb_i\). The exact prime-valuation equality
in Lemma~\ref{lem:B-RD60-break} gives its common lower
break as \((q-1)b_1=(q-1)b_2\). This proves the assertion.
\end{proof}

\begin{lemma}[Both exchanged-line quotients have residual wild action]
\label{lem:B-RD83-quotients}
\status{Proved; full independent root and companion proof and exact-code audits passed.}
Suppose two groups with isomorphic actual \(p^3\)-torsion
have tame reducible residual inertia and two-slope parameters
\[
                 B_i=p^2k_i,\qquad1\le k_i\le p-2.
                                                        \tag{B-RD83.1}
\]
Their canonical high lines are assumed to be distinct under
the comparison. Quotienting both groups by either of the two
common stable lines gives an actual \(p^2\)-torsion comparison
with nontrivial residual wild inertia over \(L\).
\end{lemma}
\begin{proof}
The original first-level modules are \(L\)-rational. Each
of the two lines is stable under the original inertia action:
one is the high line of the first model and the other is
the high line of the second. Form their actual schematic
closures in the good models and the corresponding finite-flat
quotients. The generic quotients descend to \(K\) by the
actual semilinear descent argument in B69; no good integral
model over \(K\) is asserted.

Every parameter (B-RD83.1) is greater than \(e\).
The exact high and low quotient parameters are therefore
\[
            B_{i,H}=s-B_i=p(N-pk_i),
       \qquad B_{i,J}=B_i/p=pk_i\quad(J\ne H_i).
                                                        \tag{B-RD83.2}
\]
For clarity, these formulas are the actual Weierstrass
product calculation of Lemma~\ref{lem:B-RD69-quotients},
which does not use \(p\nmid e\). For the high quotient,
each inverse of a high target has valuation \(\alpha/q\)
because \(B>e\); quotienting multiplies this by \(p\),
while the old first-level image has valuation \(B\).
The resulting parameter is \(\alpha=s-B\). For a low
quotient the old high representative has quotient valuation
\(\alpha+(p-1)\beta=s-B/p\); the inverse low points
give the other valuation \(B/p^2\), hence parameter \(B/p\).

The actual identity
\[
 (M/J)[p^2]=\{P\in M[p^3]:p^2P\in J\}/J              \tag{B-RD83.3}
\]
gives the asserted comparisons. For either chosen line,
one model uses a high quotient and the other a low quotient.
Their parameters divided by \(p\) are \(N-pk_i\) and
\(k_j\), both prime to \(p\), and
\[
                  N-pk_i\ge2p-1>p-2\ge k_j.          \tag{B-RD83.4}
\]
The original third-level image over \(L\) is a
\(p\)-group, since its first-level module is rational.
The actual quotient comparison (B-RD83.3) is a subquotient
of that module, so its image over \(L\) is also a
\(p\)-group. If its residual action were tame, it would
therefore be trivial. The preceding lemma would then
equate these unequal integers. Thus its residual wild
action is nontrivial. The actual old first-level image
\(M[p]/J\subset(M/J)[p]\) also supplies an invariant
line, so this is a nontrivial unipotent residual action.
\end{proof}

\begin{proposition}[Opposite quotient wild actions force five matrix dimensions]
\label{prop:B-RD83-Lie}
\status{Proved; full independent root and companion proof and exact-code audits passed.}
Assume in addition \(p\equiv1\pmod4\) and that the common
finite determinant is nonsquare. Under the hypotheses of
Lemma~\ref{lem:B-RD83-quotients}, the original actual
third-level wild image over \(L\) has order at least \(p^5\).
\end{proposition}
\begin{proof}
Identify the actual third-level torsion modules by the given
isomorphism, and put the two stable lines at
\(\mathbf F_pe_1\) and \(\mathbf F_pe_2\). The common
residual representation is tame and diagonal over
\(\mathbf F_p\). All its action is killed by \(L\), so
\[
 W=\operatorname{Gal}(L(M[p^3])/L)
        \subset1+pM_2(\mathbf Z/p^3).
\]
Let \(W_1\) and \(W_2\) be its actual first and terminal
matrix spaces, as in (B-RD79.2). The actual matrix identities
give
\[
 W_1\subset W_2,\qquad [W_1,W_1]\subset W_2,
       \qquad |W|=p^{\dim W_1+\dim W_2}.              \tag{B-RD83.5}
\]

We identify the residual actions of the two quotients at
the integral Tate-lattice level. If \(T=T_pM\), then under
the rational identification supplied by the isogeny the
quotient by the first line has lattice
\(T'=T+\mathbf Z_p(e_1/p)\). In the basis
\((e_1/p,e_2)\), a matrix \(\rho=(\rho_{ij})\) becomes
\[
                \begin{pmatrix}
                 \rho_{11}&p\rho_{12}\\
                 \rho_{21}/p&\rho_{22}
                \end{pmatrix}.                       \tag{B-RD83.6}
\]
For \(\rho(g)=1+pA+p^2B\), its residual lower entry is
\(A_{21}\); the other off-diagonal residual entry is zero.
Thus the nontrivial residual wild action of this quotient
forces a nonzero \(21\)-projection of \(W_1\). Quotienting
by the second line gives, in the same way, a nonzero
\(12\)-projection. These are coefficients of the original
actual representation, not independently chosen matrices.

Let a residual tame generator have eigenvalues
\(\alpha,\beta\in\mathbf F_p^\times\). Since
\(\chi_p\equiv1\pmod p\), their product is the residual
finite determinant and is nonsquare. Therefore
\(\rho_0=\alpha/\beta=\alpha^2/(\alpha\beta)\) is
nonsquare as well. When \(p\equiv1\pmod4\), both
\(1\) and \(-1\) are squares. Hence
\(1,\rho_0,\rho_0^{-1}\) are distinct.

Tame conjugation preserves \(W_1\). Its three eigenspaces
in \(M_2(\mathbf F_p)\) are the diagonal plane,
\(\mathbf F_pE_{12}\), and \(\mathbf F_pE_{21}\), with
those distinct weights. The polynomial spectral projectors,
whose denominators are nonzero elements of \(\mathbf F_p\),
preserve \(W_1\). The two nonzero off-diagonal projections
therefore imply \(E_{12},E_{21}\in W_1\). Their commutator
is \(\operatorname{diag}(1,-1)\), independent from both.
Equation (B-RD83.5) now gives
\(\dim W_1\ge2\) and \(\dim W_2\ge3\), proving the bound.
\end{proof}

\begin{corollary}[One lower bound across all residual types at \(p=5\)]
\label{cor:B-RD83}
\status{Proved; full independent root and companion proof and exact-code audits passed.}
For \(p\equiv1\pmod4\), a different-special-character
actual \(p^3\)-congruent pair at \(e=p(p-1)\) has
third-level wild image of order at least \(p^5\), regardless
of its residual type. In particular any remaining
\(125\)-torsion ambiguity at \(p=5,e=20\) has wild image
of order at least \(3125\).
\end{corollary}
\begin{proof}
The square finite-determinant case is excluded by B77, so
assume the common finite determinant is nonsquare.
Irreducible residual inertia is covered by B81, and
nontrivial residual wild inertia by B82.

For tame reducible residual inertia a one-slope polygon
would have a leading-coordinate plane with scalar action
over \(\mathbf F_p\), and hence square determinant.
Thus both polygons have two slopes. Their first-level
modules are rational over \(L\), so \(B_i=pb_i\).
If one \(b_i\) is prime to \(p\), its primitive third-level
point has valuation \(b_i/p^4\), forcing a wild field
degree divisible by \(p^4\). Equality of that degree
with \(p^4\) recovers the character by B82. Hence an
unequal-character pair in this case has order at least
\(p^5\).

If both \(b_i\) are divisible by \(p\), write
\(B_i=p^2k_i\) with \(1\le k_i\le p-2\). If the high
lines agree, equality of finite determinants gives
\(c_1\equiv c_2\pmod{p-1}\); equality of their high-line
characters gives \(B_1\equiv B_2\pmod{p-1}\). The stated
range then gives \(k_1=k_2\), and the line character gives
\(c_1=c_2\pmod N\). Otherwise the lines are distinct,
and the proposition applies. These alternatives exhaust
the residual types.
\end{proof}

\begin{remark}[The still-open larger-image comparisons]
\label{rem:B-RD83-scope}
\status{Proved; full independent root and companion proof and exact-code audits passed.}
No order at least \(p^5\) is excluded or constructed here.
The full first-wild least index remains in \(\{2,3\}\).
The condition \(p\equiv1\pmod4\) is used precisely to
separate the two opposite tame weights. When
\(p\equiv3\pmod4\), a nonsquare ratio can equal \(-1\),
and the two nonzero projections alone do not produce the
two separate matrix units. That case is not included in
the proposition or the final corollary.

\smallskip
\noindent\emph{Exact supplement and provenance.}
The retained standalone program
\path{code/reduction_descent/verify_first_wild_swapped_quotients.py}
checks the two symbolic lattice conjugations and the two
exact residual matrices. For every prime from 5 through 101
it checks all 71,456 opposite-quotient parameter pairs,
36,848 nonsquare residual eigenvalue pairs, 24 literal
commutator and power rows, 24 dimension-bound rows and
74,804 residual-tame parameter branches. The order-two
exceptions at primes congruent to 3 modulo 4 are retained
explicitly, rather than included in the theorem.
The author run passed in \(0.3601258019916713\) seconds,
with no failed execution. The initial and clarified proofs,
every executed program and its output are retained. The
complete internal B69 point-field and actual quotient
proofs were reread and their relevant steps rederived
at the present wild index. No new original-source read
is claimed. These finite checks establish no formal-group
existence and do not replace the uniform proofs.
The complete independent root and companion proof and code
audits both passed. Their author-output-preserving reruns
took \(0.3534236279956531\) and
\(0.3482199970021611\) seconds, respectively. Both
auditors checked the actual quotient lattice direction,
the common point fields, and the same original first
congruence layer before applying its tame weight
projectors. Their exact read scopes, all prior proof
versions and the complete outputs remain retained.
\end{remark}


% Proved after CP29; no unrestricted first-wild least-index claim.
\subsection{Two translation breaks recover a tame residual Newton parameter}

Let \(p\ge5\), let \(K\) be complete discretely valued of mixed
characteristic \((0,p)\) with algebraically closed residue field,
and suppose \(\mu_p\subset K\). Its absolute index \(e\) is
arbitrary; in particular it may be divisible by \(p\).
The cyclotomic uniformizer valuation gives \(p-1\mid e\).
Put \(q=p^2\), \(N=q-1\), \(s=(p+1)e\), and pass to
the common tame good field \(L/K\) of degree \(N\).
The groups considered are whole potentially supersingular
height-two groups.

\begin{proposition}[One prime parameter suffices in the tame residual comparison]
\label{prop:B-RD84-tame}
\status{Proved; full independent root and companion proof and exact-code audits passed.}
Suppose two such groups have two-slope first-level polygons,
\(L\)-rational first-level modules, and isomorphic actual
\(p^2\)-torsion. Write their parameters and valuations as
\[
 \begin{gathered}
 B_i=pb_i,\quad0<b_i<e,\qquad
 \alpha_i=s-pb_i,\qquad\beta_i=b_i,\\
 v_L(T_i(Q))=\alpha_i\quad(0\ne Q\in H_i),\qquad
 v_L(T_i(Q))=b_i\quad(Q\notin H_i).
 \end{gathered}                                      \tag{B-RD84.1}
\]
If at least one of \(b_1,b_2\) is prime to \(p\), then
the actual comparison identifies the high lines,
\(b_1=b_2\), and the finite special characters agree.
No bound on the full second- or third-level division-field
degree is assumed.
\end{proposition}
\begin{proof}
Index the models so that \(p\nmid b_1\). Choose corresponding
residual vectors outside both high lines and corresponding
primitive second-level lifts \(P_i\). The actual inverse
Newton polygon gives
\[
                     v_L(T_i(P_i))=b_i/q.             \tag{B-RD84.2}
\]
Here the same universal coefficient and line-integrality
proof as in B69 and B82 applies at arbitrary \(e\);
no coefficient argument requires \(p\nmid e\).

The actual torsion isomorphism identifies the point
stabilizers, hence the point fields \(E=L(P_1)=L(P_2)\).
Because both first-level modules and both \(pP_i\) are
\(L\)-rational, each field is Galois with its group a
subgroup of the first-level translation group. Indeed
every conjugate is \(P_i+Q\) for an \(L\)-rational \(Q\),
and so is defined over the same field. Its degree is at
most \(q\). The reduced denominator of \(b_1/q\) is
\(q\), forcing \([E:L]=q\). Thus every first-level
translation occurs in the same actual point field, even
if \(b_2\) is divisible by \(p\).

Normalize \(v_E\) to have value group \(\mathbf Z\).
The coordinates of \(P_i\) have valuations \(b_i\).
For every nonidentity \(g\in\operatorname{Gal}(E/L)\),
Lemma~\ref{lem:B-RD60-break} gives the exact equality
for model 1 and the unconditional inequality for model 2:
\[
 \begin{split}
 v_E(gT_1(P_1)-T_1(P_1))&=b_1+r(g),\\
 v_E(gT_2(P_2)-T_2(P_2))&\ge b_2+r(g).
 \end{split}                                         \tag{B-RD84.3}
\]
Both concern the same actual lower break \(r(g)\).

If the high lines differed, choose a translation
\(0\ne Q_1\in H_1\) corresponding to \(Q_2\notin H_2\).
It occurs because the full first-level translation group
occurs. Its high displacement in model 1 gives
\(r(g)=q\alpha_1-b_1\), whereas its low displacement
in model 2 would require
\[
                  qb_2\ge b_2+q(s-pb_1)-b_1.         \tag{B-RD84.4}
\]
But \(\alpha_1=s-pb_1>e>b_1,b_2\), contradicting this
inequality. Explicitly the difference of the left side
and right side is increasing in both integers and at
\(b_1=b_2=e-1\) equals \(-q(p+1)<0\).
The high lines therefore agree.

Choose a nonzero low translation in the common module.
Its break from model 1 is \((q-1)b_1\); (B-RD84.3)
in model 2 gives
\[
                     (q-1)(b_2-b_1)\ge0.             \tag{B-RD84.5}
\]
Choose instead a nonzero high translation. Its break is
\(q(s-pb_1)-b_1\), and the other displacement gives
\[
                     (p^3+1)(b_1-b_2)\ge0.           \tag{B-RD84.6}
\]
These force \(b_1=b_2\), with no prime-to-\(p\) assumption
on the second coordinate. The identified high-line
characters are \(\zeta^{c_i-B_i}\), in an exact tame
Honda coordinate. Since \(B_1=B_2\), their equality
gives \(c_1=c_2\pmod N\), and hence equal finite
special characters.
\end{proof}

\begin{corollary}[The remaining tame reducible first-wild parameters]
\label{cor:B-RD84-first-wild}
\status{Proved; full independent root and companion proof and exact-code audits passed.}
At \(e=p(p-1)\), a different-character actual
\(p^3\)-congruent pair with tame reducible residual inertia
must have two-slope parameters
\[
 B_i=p^2k_i,\qquad1\le k_i\le p-2,                  \tag{B-RD84.7}
\]
and its high lines are exchanged. This restriction is
independent of the size of its full wild torsion image.

At \(p=5,e=20\), after a common finite height-one twist
normalizing the nonsquare finite determinant, the unordered
parameter pair is exactly \(\{25,75\}\). Its finite
types are A/B up to common quadratic twist, with A, B,
C as in B77. This is only a necessary pattern and asserts
no existence of such a pair.
\end{corollary}
\begin{proof}
The square determinant case is already excluded at the
third level by B77. A one-slope residual representation,
when tame and reducible, acts by a scalar on its
two-dimensional leading-coordinate plane and has square
determinant. Thus both polygons have two slopes. Since
their residual action is tame, both first-level modules
are rational over \(L\), and \(B_i=pb_i\).

Their actual \(p^2\)-torsion is already isomorphic.
The proposition therefore forces both \(b_i\) to be
divisible by \(p\) if the finite characters differ.
The range \(0<b_i<e=p(p-1)\) gives (B-RD84.7).
If their high lines agreed, equality of finite determinants
would give \(c_1\equiv c_2\pmod{p-1}\), and the common
high-line character would give
\(p^2k_1\equiv p^2k_2\pmod{p-1}\). The range
\(1\le k_i\le p-2\) forces \(k_1=k_2\); the same
high-line character then gives \(c_1=c_2\pmod N\).
This is excluded, so the high lines must be exchanged.

For the explicit pair at \(p=5\), make the common
nonsquare determinant primitive by the common finite
twist already justified in B50 and B77. This leaves the
good models over \(L\), hence their \(B_i\), unchanged.
Now \(c_i\equiv1\pmod4\), while
\(B_i\equiv c_i\pmod6\). Thus \(k_i\) is odd in
\(\{1,2,3\}\), giving \(B_i\in\{25,75\}\).
Equality of the residual characters with exchanged high
lines forces the two distinct values. More explicitly,
the possible normalized ordered pairs are
\[
 (B_1,c_1;B_2,c_2)=(25,1;75,9)
             \quad\text{or}\quad(25,13;75,21),         \tag{B-RD84.8}
\]
up to interchanging the two models. Their residual
high/low exponents are respectively \((0,6;6,0)\)
and \((12,18;18,12)\) modulo 24. The finite special
orbits are A/B and C/B; a common quadratic twist
interchanges A and C and preserves B. This proves the
stated necessary pattern directly.
\end{proof}

\begin{remark}[What this parameter recovery does not decide]
\label{rem:B-RD84-scope}
\status{Proved; full independent root and companion proof and exact-code audits passed.}
The proposition uses the small common point field of a
single primitive second-level point. It does not assume
that the full division field has that degree. Both the
high and the low translations in that same field are
essential; using only the low inequality would leave
spurious one-divisible parameter candidates.
The corollary removes the earlier exploratory tame pairs
\(\{35,75\}\) and \(\{55,75\}\) at every image size.
The pair \(\{25,75\}\), all larger-image questions in
other residual sectors, and the full index interval
\(\{2,3\}\) remain open. No formal-group pair is
constructed by the finite parameter enumeration.
\smallskip
\noindent\emph{Exact supplement and provenance.}
The retained program
\path{code/reduction_descent/verify_tame_residual_one_prime_recovery.py}
checks the three symbolic displacement identities,
80,536 point-field parameter rows over 96 prime/index
samples, including 6,600 mixed prime/divisible exclusions,
and the exact two normalized \(p=5\) character rows and
their common quadratic-twist orbit. All parameter pairs
are checked for \(p=5,7,11,13\) at the four displayed
index factors; larger primes through 101 use explicitly
retained boundary samples. The author run passed in
\(0.12348706599732395\) seconds with no failed execution.
The uniform common-point-field and two-break proof is
independent of those finite ranges. All executed code,
output and prior source versions remain retained.
The exact internal B69 quotient and inverse-polygon proof,
B60 fixed-coefficient break lemma and B82/B83 common-field
arguments were reread; no new original-source read is
claimed.
The complete independent root and companion proof and code
audits passed; their author-output-preserving reruns took
\(0.11858891700103413\) and
\(0.11411567399045452\) seconds, respectively.
Both exact audit records and the prior candidate remain
retained. The earlier exploratory parameter rows are
preserved as superseded research, without a retroactive
claim that the current argument had already been known.
\end{remark}


% Proved after checkpoint30 archive release; exact candidate and full audits retained.
\subsection{A small point field forces both first-layer diagonal projections}

Let \(K\) be complete discretely valued of mixed characteristic
\((0,p)\), with algebraically closed residue field,
\(p\ge5\), \(e_K=p(p-1)\), and \(\mu_p\subset K\).
Put \(q=p^2\), \(N=q-1\), and \(s=pN\).
Use the common tame good field \(L/K\) of degree \(N\).
The groups below are whole potentially supersingular
height-two Barsotti--Tate groups. All comparisons are actual
isomorphisms of torsion Galois modules, and the matrix layers
are the images of those actual representations.

\begin{lemma}[Neither diagonal projection of the first layer can vanish]
\label{lem:B-RD85-diagonal}
\status{Proved; full independent root and companion proof and exact-code audits passed.}
Suppose two groups with isomorphic actual \(p^3\)-torsion
have \(L\)-rational first-level modules, parameters
\[
 B_i=p^2k_i,\qquad 1\le k_i\le p-2,
       \qquad H_1=\mathbf F_pe_1,\quad H_2=\mathbf F_pe_2.
                                                        \tag{B-RD85.1}
\]
Thus the high lines are exchanged. Let
\(W\subset1+pM_2(\mathbf Z/p^3)\) be the common actual
third-level image over \(L\), and let \(W_1\) be its
image under \(1+pA\longmapsto A\pmod p\).
Both linear maps
\[
        W_1\longrightarrow\mathbf F_p,
        \qquad A\longmapsto A_{11},\quad A\longmapsto A_{22}
                                                        \tag{B-RD85.2}
\]
are nonzero. No bound on the full division-field degree
is assumed.
\end{lemma}
\begin{proof}
The actual opposite-quotient proof of
Lemma~\ref{lem:B-RD83-quotients} and its lattice calculation
shows that the \(12\)- and \(21\)-projections of the same
\(W_1\) are nonzero. This part does not use
\(p\equiv1\pmod4\).

Suppose the \(11\)-projection vanishes. It follows that
\(W_1e_1=\mathbf F_pe_2\). Choose corresponding primitive
second-level points \(P_i\) with \(pP_i=e_1\) in the
common first-level module. This choice means the residual
first-level target is high in model 1 and low in model 2.
The actual action modulo \(p^2\) translates such a point
by \(Ae_1\). Hence its orbit has exactly \(p\) elements.
Because all first-level points are \(L\)-rational, its
point field is Galois, and the actual comparison gives
one common field
\[
       E=L(P_1)=L(P_2),\qquad [E:L]=p,
       \qquad\operatorname{Gal}(E/L)=\mathbf F_pe_2
                                                        \tag{B-RD85.3}
\]
as the translation group. This computation uses the
second-level point orbit, not the order of \(W\).

We check the high inverse polygon explicitly. In the
normalized good coordinate the first-level high target
in model 1 has valuation
\(\alpha_1=s-p^2k_1=p(N-pk_1)\). The coefficient of
\(T^p\) in \([p](T)\) has valuation
\((p-1)p^2k_1\), and the coefficient of \(T^q\) is a
unit. The intermediate universal coefficient bounds
are those in the complete inverse-polygon proof of B69.
The intermediate vertex is strictly above the chord
from \((0,\alpha_1)\) to \((q,0)\), since
\[
          p^2k_1>N-pk_1;
       \qquad p(p+1)k_1>N.                           \tag{B-RD85.4}
\]
The linear coefficient also lies above that chord.
All \(q\) inverse points therefore have valuation
\(\alpha_1/q=(N-pk_1)/p\). The low inverse polygon
in model 2 similarly gives valuation \(k_2/p\).
With the integral valuation on the degree-\(p\) field
\(E\), the two coordinate valuations are consequently
\[
               v_E(T_1(P_1))=N-pk_1,
             \qquad v_E(T_2(P_2))=k_2.               \tag{B-RD85.5}
\]
Both are positive and prime to \(p\).

Every nonidentity automorphism of (B-RD85.3) translates
by a nonzero point of \(\mathbf F_pe_2\), which is low
in model 1 and high in model 2. Formal addition preserves
the coordinate valuation of this displacement. Its two
displacement valuations in \(E\) are thus \(qk_1\) and
\(q(N-pk_2)\), respectively. Applying the exact
prime-valuation equality in
Lemma~\ref{lem:B-RD60-break} to the same automorphism
and the two prime valuations in (B-RD85.5) gives
\[
 \begin{split}
  r(g)&=p(p+1)k_1-N,\\
  r(g)&=qN-(p^3+1)k_2.
 \end{split}                                         \tag{B-RD85.6}
\]
Equality would require
\[
         p(p+1)k_1+(p^3+1)k_2=p^4-1.                \tag{B-RD85.7}
\]
Both coefficients are positive. At the largest permitted
values \(k_1=k_2=p-2\), the right side exceeds the left
side by
\[
                     (p+1)(p^2+1)>0.                \tag{B-RD85.8}
\]
This is a contradiction. Interchanging the models and
the two basis vectors excludes a vanishing
\(22\)-projection in the same way.
\end{proof}

\begin{corollary}[Six wild matrix dimensions in the tame residual sector]
\label{cor:B-RD85-tame}
\status{Proved; full independent root and companion proof and exact-code audits passed.}
Suppose \(p\equiv1\pmod4\). At \(e=p(p-1)\), any
different-special-character actual \(p^3\)-congruent
pair with tame reducible residual inertia has
\[
                \dim W_1\ge3,\qquad\dim W_2\ge3,
                        \qquad |W|\ge p^6.          \tag{B-RD85.9}
\]
In particular the remaining tame reducible sector of
the \(p=5,e=20\) problem has wild image of order at
least \(15625\), and still has the necessary parameter
pair \(\{25,75\}\), type A/B up to a common quadratic
twist.
\end{corollary}
\begin{proof}
B77 excludes square finite determinant, and
Corollary~\ref{cor:B-RD84-first-wild} gives exactly the
exchanged-line parameters (B-RD85.1) at every image size.
The tame spectral projectors in
Proposition~\ref{prop:B-RD83-Lie} put
\(E_{12},E_{21}\) in \(W_1\), and also preserve its
diagonal part. The lemma gives a nonzero diagonal
projection, so that diagonal part is nonzero and
independent of the two off-diagonal directions.
Thus \(\dim W_1\ge3\). Actual \(p\)-th powers give
\(W_1\subset W_2\), hence \(\dim W_2\ge3\).
Since the original residual action over \(L\) is trivial,
\(|W|=p^{\dim W_1+\dim W_2}\), proving the result.
The explicit \(p=5\) pattern is the already proved
corollary to B84; it is not a new existence assertion.
\end{proof}

\begin{remark}[Scope of the diagonal recovery]
\label{rem:B-RD85-scope}
\status{Proved; full independent root and companion proof and exact-code audits passed.}
The point-field lemma holds for every \(p\ge5\) under
its stated exchanged-line hypotheses. Only the
six-dimensional matrix consequence uses
\(p\equiv1\pmod4\) to separate opposite tame weights.
The irreducible and nontrivial residual-wild sectors
retain their earlier bounds; no lower bound \(p^6\)
across all residual types is asserted. No image of order
\(p^6\) or greater is excluded or constructed, and the
unrestricted first-wild least-index interval remains
\(\{2,3\}\).

The argument was suggested by the retained parent note
\path{notes/root_B85_point_field_lead.md}. It uses an
actual degree-\(p\) second-level point field, with the
same automorphism evaluated on two different prime-valued
coordinates. It does not estimate an orbit from only
the ranks of evaluations in two congruence layers.
Nor does a zero trace on the terminal matrix layer by
itself force the entire finite-level determinant to
vanish; such an additional inference is not used.
The full B83 actual quotient and lattice argument and
the relevant B69 inverse-polygon and B60 exact-break
proofs were reread. No new original-source reading is
claimed.

\smallskip
\noindent\emph{Exact supplement.}
The retained program
\path{code/reduction_descent/verify_first_wild_diagonal_projection.py}
checks three symbolic identities, all 71,456 parameter
pairs for every prime from 5 through 101, and 142,912
inverse-slope and prime-valuation rows. It also checks
1,132 literal degree-\(p\) translation images,
21,272 separated tame-weight pairs and 43,752 independent
diagonal extensions in the stated congruence class.
The author run passed in \(0.3171308919991134\) seconds,
with no failed execution. The actual point-field and
lower-break proofs do not depend on these finite ranges.
All executed code, output and prior source bytes are
retained; no formal-group realization follows from
the finite checks.

\smallskip
\noindent\emph{Independent audits.}
The complete independent root and companion audits passed.
Their author-output-preserving exact reruns took
\(0.3250291720032692\) and
\(0.315905258001294\) seconds, respectively.
Their exact proof and source read scopes, all prior
candidate bytes and every executed output remain retained.
\end{remark}


% Proved after checkpoint30 archive release; exact candidate and full audits retained.
\subsection{The cyclotomic quotient excludes a traceless six-dimensional image}

Use the complete discretely valued field, actual torsion
comparison and tame good extension \(L/K\) of B85:
\(p\ge5\), \(e_K=p(p-1)\), algebraically closed residue,
\(\mu_p\subset K\), \([L:K]=N=p^2-1\).
All valuations below have value group \(\mathbf Z\) on
the field named in their subscript. Put \(q=p^2\).

\begin{lemma}[The last cyclotomic break after the tame good extension]
\label{lem:B-RD86-cyclotomic}
\status{Proved; full independent root and companion proof and exact-code audits passed.}
If \(\mu_{p^2}\subset L\), then \(\mu_{p^2}\subset K\),
\([L(\mu_{p^3}):L]=p\), and every nonidentity element
of that cyclic group has lower break \(N^2\).
\end{lemma}
\begin{proof}
Because \(\mu_p\subset K\), the degree of
\(K(\mu_{p^2})/K\) divides \(p\). It also divides
the prime-to-\(p\) degree \(N\), so is one.
Choose a Cohen coefficient ring \(\Lambda\subset\mathcal O_K\)
and put \(K_0=\operatorname{Frac}(\Lambda)\);
its existence follows from the full Cohen structure
theorem \cite[Tag 032A]{Stacks032A}.
The residue field is perfect, so \(\Lambda\) is an
unramified complete discrete valuation ring with
uniformizer \(p\) and the same residue field.

For a primitive \(p^2\)-th root \(\zeta_2\), the
polynomial \(\Phi_{p^2}(1+T)\) is Eisenstein over
\(\Lambda\), of degree \(p(p-1)\). Its root
\(\zeta_2-1\) is therefore a uniformizer of \(K\).
The complete subring \(\Lambda[\zeta_2]\) contains
fixed representatives of every residue element and
that uniformizer. Successive uniformizer expansion
therefore gives \(\mathcal O_K=\Lambda[\zeta_2]\),
and \(K=K_0(\zeta_2)\). This explicitly justifies
the coefficient-field identification at the exact
absolute index; no canonical choice of \(\Lambda\)
is required.

Let \(K_3=K_0(\zeta_3)\), where \(\zeta_3\) has
order \(p^3\) and \(\zeta_3^p=\zeta_2\).
The analogous Eisenstein polynomial has degree
\(p^2(p-1)\), and \(K_3/K\) is cyclic of degree
\(p\). A uniformizer is \(\pi=\zeta_3-1\).
For a nonidentity automorphism \(\sigma\),
\[
   \sigma\pi-\pi=\zeta_3(\zeta_p^a-1),\quad
   p\nmid a,\qquad v_{K_3}(\sigma\pi-\pi)=p^2.
                                                        \tag{B-RD86.1}
\]
Thus its lower break is \(p^2-1=N\).

The extensions \(L/K\) and \(K_3/K\) are disjoint.
Their composite \(C=LK_3\) is the tame extension of
\(K_3\) of degree \(N\). Since the residue field is
algebraically closed, every unit has an \(N\)-th root
by Hensel's lemma. The tame extension consequently
admits a uniformizer \(t\) with \(t^N=\pi\);
this is also the direct tame uniformizer argument in
\cite[Proposition 2.12 and proof, p.~11]{ColmezCorpsLocaux}.
Extend \(\sigma\) by fixing \(L\). Since it is wild,
\(\sigma(t)/t\) has residue one. The identity
\[
       \left(\frac{\sigma(t)}t\right)^N
           =\frac{\sigma\pi}{\pi}
                                                        \tag{B-RD86.2}
\]
and the unit \(N\) show that
\[
  v_C\left(\frac{\sigma(t)}t-1\right)
   =v_C\left(\frac{\sigma\pi}{\pi}-1\right)
   =N\,v_{K_3}\left(\frac{\sigma\pi}{\pi}-1\right)
   =N^2.                                               \tag{B-RD86.3}
\]
This is the asserted lower break of \(C/L\).
\end{proof}

\begin{proposition}[The two actual congruence layers cannot both be traceless]
\label{prop:B-RD86-traceless}
\status{Proved; full independent root and companion proof and exact-code audits passed.}
Suppose the residual representation is tame reducible,
the finite determinant is nonsquare, the high lines
are exchanged, and the actual parameters are
\(B_i=p^2k_i\), \(1\le k_i\le p-2\).
Assume the residual eigenvalue ratio is different
from one. Then the common actual third-level image
over \(L\) cannot have
\[
                  W_1=W_2=\mathfrak{sl}_2(\mathbf F_p)
                                                        \tag{B-RD86.4}
\]
in the basis \(H_1=\mathbf F_pe_1\),
\(H_2=\mathbf F_pe_2\).
\end{proposition}
\begin{proof}
Suppose (B-RD86.4) holds. On the good field \(L\)
the actual determinant is the cyclotomic character
\(\chi_p\), as in B76 and B77; the finite determinant
character is killed by the good extension. The first
traceless layer makes this determinant trivial modulo
\(p^2\), and therefore \(\mu_{p^2}\subset L\).
The lemma shows that the remaining character modulo
\(p^3\) is nonzero of order \(p\).

For every terminal element,
\[
       \det(1+p^2A)=1+p^2\operatorname{tr}(A)=1
                        \pmod{p^3}.                  \tag{B-RD86.5}
\]
Hence the actual determinant induces a nonzero linear
functional
\[
 \ell:W_1=\mathfrak{sl}_2(\mathbf F_p)\longrightarrow
 \mathbf F_p,\qquad
 \det(g)=1+p^2\ell(A)\pmod{p^3},\quad g\mapsto A.
                                                        \tag{B-RD86.6}
\]
This does not assert that the determinant itself is
zero merely because its terminal restriction is zero.
Its invariance under actual tame conjugation implies
\(\ell(E_{12})=\ell(E_{21})=0\), since the two
off-diagonal weights are different from one.
Writing \(H=\operatorname{diag}(1,-1)\), one has
\(\ell(H)\ne0\).

Choose an actual primitive second-level point \(P\)
with \(pP=e_1\), and let \(E=L(P)\). The common
second-level action is translation by \(W_1e_1\),
which is all of \(\mathbf F_p^2\). All first-level
points are rational, so \(E/L\) is Galois of degree
\(q\), with
\[
 \operatorname{Gal}(E/L)=
      \mathfrak{sl}_2(\mathbf F_p)/\mathbf F_pE_{12}.
                                                        \tag{B-RD86.7}
\]
The actual point stabilizer also contains the terminal
layer, on which the determinant vanishes. By (B-RD86.6)
the entire point stabilizer is contained in the
cyclotomic kernel. Thus the cyclotomic field
\(C=L(\mu_{p^3})\) is a subfield of this actual point
field \(E\).

Every class in (B-RD86.7) has a representative
\(aH+cE_{21}\). It translates \(P\) by
\(ae_1+ce_2\), and it has nontrivial image in
\(\operatorname{Gal}(C/L)\) exactly when \(a\ne0\).
In model 2 such a translation is outside its high
line \(\mathbf F_pe_2\), hence has low valuation
\(pk_2\) over \(L\). The target \(pP=e_1\) is
also low in model 2, so its inverse coordinate has
valuation \(k_2/p\) over \(L\). In \(E\),
\[
       v_E(T_2(P))=pk_2,\qquad
       v_E(gT_2(P)-T_2(P))=p^3k_2\quad(a\ne0).
                                                        \tag{B-RD86.8}
\]
The unconditional lower-break inequality of B60 gives
\[
             b_E(g)\le p^3k_2-pk_2=pNk_2
                \le pN(p-2)<N^2                   \tag{B-RD86.9}
\]
for every lift of every nonidentity element of
\(\operatorname{Gal}(C/L)\).

For completeness the exact quotient formula for a
nonidentity \(\bar g\in\operatorname{Gal}(C/L)\) is
\[
 b_C(\bar g)+1
   =\frac1{[E:C]}\sum_{g\mapsto\bar g}(b_E(g)+1).
                                                        \tag{B-RD86.10}
\]
Here both fields are totally ramified and the fibre
has \([E:C]=p\) elements. Formula (B-RD86.10) is the
norm-polynomial formula, with its full two-divisibility
proof, in \cite[Proposition 3.7, pp.~14--15]{ColmezCorpsLocaux}.
It uses no finite-residue hypothesis. It implies from
(B-RD86.9) that \(b_C(\bar g)<N^2\), contradicting
the cyclotomic lemma. This proves the proposition.
\end{proof}

\begin{corollary}[Seven dimensions in the exchanged tame sector]
\label{cor:B-RD86-tame}
\status{Proved; full independent root and companion proof and exact-code audits passed.}
Assume \(p\equiv1\pmod4\). In the different-character
tame reducible sector at \(e=p(p-1)\), an actual
\(p^3\)-congruent pair must have wild image order at
least \(p^7\). At \(p=5,e=20\), the necessary
\(\{25,75\}\) tame parameter pattern can therefore
remain only at wild orders \(5^7\) or \(5^8\).
\end{corollary}
\begin{proof}
B85 gives \(\dim W_1\ge3\), and actual powers give
\(W_1\subset W_2\). If the total dimension were at
most six, both spaces would have dimension three and
would coincide. B83 puts \(E_{12},E_{21}\) in the
first space, and their actual commutator puts
\(H=\operatorname{diag}(1,-1)\) in the second.
The two equal three-dimensional spaces would therefore
both be \(\mathfrak{sl}_2(\mathbf F_p)\), excluded by
the proposition. The total dimension is at least seven;
the general two-layer upper bound is eight.
The parameter restriction is B84, with no existence
assertion.
\end{proof}

\begin{remark}[The remaining image orders and source scope]
\label{rem:B-RD86-scope}
\status{Proved; full independent root and companion proof and exact-code audits passed.}
This argument does not exclude the seven- or
eight-dimensional images and does not strengthen the
other two residual sectors. The full first-wild
least-index interval remains \(\{2,3\}\).
The new obstruction uses the actual cyclotomic
subfield of a degree-\(p^2\) point field. It is not a
claim that all traceless terminal images have zero
finite determinant.

The complete original Cohen structure proof, Tag 032A,
and Colmez's original Proposition 2.12 with proof on
p.~11 and Proposition 3.7 with its full proof on
pp.~14--15 were freshly opened and read for this
candidate. The elementary cyclotomic and tame-root
calculations are written above. The actual B83/B85
matrix layers and the full B60 displacement lemma
were reread. The separate B85 dependency and this full proof
both completed independent audits before promotion.
They were promoted after the immutable checkpoint30
archive and first enter the following checkpoint.

\smallskip
\noindent\emph{Exact supplement.}
The retained standalone program
\path{code/reduction_descent/verify_first_wild_cyclotomic_quotient.py}
checks the symbolic determinant, bracket, point-field
evaluation and strict-gap identities, every small
parameter for primes 5 through 101, and all diagonal
characters on the displayed additive point quotients
for primes through 31. It also computes every
coefficient of the six shifted cyclotomic polynomials
at \(p=5,7,11\), \(n=2,3\), and verifies their
Eisenstein conditions. An explicit nonzero determinant
with traceless terminal power is retained as a guard
against the invalid simpler inference. The author run
passed in \(0.4382091660081642\) seconds. No execution failed.
The finite ranges do not replace the actual field
inclusion and quotient-ramification proofs. Every
executed program, output and prior proof version is
retained.

\smallskip
\noindent\emph{Independent audits.}
The complete independent root and companion audits passed.
Their author-output-preserving exact reruns took
\(0.4253982070076745\) and
\(0.4311180119984783\) seconds, respectively.
Their exact proof and source read scopes, all prior
candidate bytes and every executed output remain retained.
\end{remark}


% Proved after complete root and companion audits; exact candidate retained.
\subsection{Two shared cyclic quotients force the full first layer at \(p=5\)}

\begin{lemma}[A coordinate whose high-displacement line is not the last line]
\label{lem:B-RD87-coefficient}
\status{Proved; full independent root and companion audits passed.}
Let \(E/L\) be a totally ramified Galois extension of
complete discretely valued mixed-characteristic fields
with group \(G\simeq\mathbf F_p^2\) and the same perfect
residue field. Suppose \(G\) has two lower breaks
\(t<u\), with last ramification line \(T\).
Let \(x\in\mathfrak m_E\), put \(A=v_E(x)\), and
suppose \(v_p(A)=1\). Assume a line \(H\subset G\)
and integers \(D_0<D_1<v_E(p)\) satisfy
\[
 v_E(gx-x)=
 \begin{cases}
  D_1,&0\ne g\in H,\\
  D_0,&g\notin H.
 \end{cases}                                         \tag{B-RD87.1}
\]
If \(H\ne T\), then
\[
     \frac{u+(p-1)t}{p}=\frac{D_0-A}{p}.              \tag{B-RD87.2}
\]
The left side is the lower break of any cyclic
degree-\(p\) quotient in which \(T\) has nonzero image.
\end{lemma}
\begin{proof}
Choose a uniformizer \(\varpi\) of \(E\) and fixed
residue representatives in \(\mathcal O_L\). Write
\(x=\sum_{j\ge A}a_j\varpi^j\), with \(a_A\) a unit
and all nonzero \(a_j\) units. Let \(J\) be the first
degree with \(p\nmid J\) and \(a_J\ne0\), allowing
\(J=+\infty\).

Up to valuation \(D_1<v_E(p)\), terms whose integer
binomial coefficient is divisible by \(p\) do not
contribute. If an automorphism has lower break \(b\),
the leading degree \(A\), whose \(p\)-adic order is
exactly one, first changes in degree \(A+pb\), with
nonzero coefficient. Every higher degree divisible
by \(p\) changes later. The first prime-to-\(p\)
degree first changes in degree \(J+b\). Therefore
the first possible change has degree
\[
                 m(b)=\min\{A+pb,J+b\},              \tag{B-RD87.3}
\]
and it is the actual change degree when the two
quantities are distinct and below \(v_E(p)\).
Only equality of the two quantities permits
cancellation. These assertions follow directly
by expanding \(g\varpi=\varpi(1+d\varpi^b+\cdots)\);
the same fixed-coefficient convention is used in B60
and B81.

Take \(0\ne g\in H\) and \(0\ne h\in T\).
Since the lines differ, their breaks are \(t,u\),
and their coordinate displacements are \(D_1,D_0\).
In particular \(m(u)\le D_0\). Both expressions in
(B-RD87.3) strictly increase with \(b\), so
\(m(t)<D_0<D_1\). The first change for \(g\) must
cancel. Consequently \(J\) is finite and
\[
                      J=A+(p-1)t.                    \tag{B-RD87.4}
\]
At the larger break \(u\), the degree \(J+u\) is
strictly less than \(A+pu\), hence is the unique
first change. Thus
\(D_0=J+u=A+(p-1)t+u\), proving the identity.

For the quotient assertion, every nonzero coset of
a kernel line different from \(T\) contains one
element of \(T\) and \(p-1\) elements outside it.
The exact norm-polynomial quotient formula
\cite[Proposition 3.7 and full proof,
pp.~14--15]{ColmezCorpsLocaux} averages their
\(b+1\) values and gives precisely
\((u+(p-1)t)/p\). The quotient killing \(T\) instead
has break \(t\). Thus the former is the larger
of these two cyclic quotient breaks.
\end{proof}

Now let \(K\) be complete discretely valued of mixed
characteristic \((0,5)\), with algebraically closed
residue field, \(e_K=20\), and \(\mu_5\subset K\).
Let \(L/K\) be the common tame good field of degree
24. Suppose two whole potentially supersingular
height-two groups have isomorphic actual
\(125\)-torsion, distinct finite special characters,
and tame reducible residual inertia.

\begin{lemma}[The two point fields share a diagonal cyclic quotient]
\label{lem:B-RD87-fields}
\status{Proved; full independent root and companion audits passed.}
If \(\dim W_1=3\), then there are two actual Galois
point fields \(E_1,E_2\) of degree 25 over \(L\),
sharing a diagonal cyclic degree-five quotient field
\(D\). Write \(d\) for the break of \(D/L\), and
\(o_i\) for that of the complementary off-diagonal
cyclic quotient of \(E_i/L\). Then
\[
     d>0,\quad d\equiv0\pmod{24},\qquad
        o_1\equiv6\pmod{24},\quad o_2\equiv18\pmod{24}.
                                                        \tag{B-RD87.5}
\]
Each \(E_i/L\) has two distinct lower breaks, with
last ramification line diagonal if \(d>o_i\), and
off-diagonal if \(o_i>d\).
\end{lemma}
\begin{proof}
By B84, after interchanging the groups and making
the justified common finite twist, the parameters
are \(B_1=25\), \(B_2=75\). Put the high lines at
\(H_1=\mathbf F_5e_1\), \(H_2=\mathbf F_5e_2\).
The residual tame eigenvalue exponents on these
two vectors are \((0,6)\), or \((12,18)\),
modulo 24. Their ratio is the same in both rows.
Thus tame conjugation has weights 6 on \(E_{21}\),
18 on \(E_{12}\), and zero on diagonal matrices.

B83 and B85 give
\[
 W_1=\mathbf F_5E_{12}\oplus\mathbf F_5E_{21}
                    \oplus\mathbf F_5Z,
             \qquad Z=\operatorname{diag}(a,b),\quad ab\ne0.
                                                        \tag{B-RD87.6}
\]
Indeed the tame spectral projectors separate all
three weights, and both diagonal-coordinate
projections are nonzero. The common second-level
division field over \(L\) has additive Galois group
exactly \(W_1\).

Choose a primitive second-level point above \(e_i\),
and call its field \(E_i\). First-level rationality
makes it Galois, with translation group \(W_1e_i\).
Since \(a,b\ne0\), both evaluations are surjective,
with kernels \(\mathbf F_5E_{12}\) for \(i=1\)
and \(\mathbf F_5E_{21}\) for \(i=2\).
Their common quotient by both off-diagonal axes
is the same degree-five field \(D\). In each
translation group the diagonal line is the image
of \(Z\), and the other line is the image of the
remaining off-diagonal matrix. All these kernel
spaces are stable under tame conjugation, so the
three cyclic quotient fields are Galois over \(K\)
as well as over \(L\).

We spell out the tame-break congruence. In a cyclic
degree-five quotient \(F/L\), let \(\tau\) lift
the tame generator of \(L/K\), whose leading
uniformizer multiplier is a primitive 24th root
\(\zeta\). On a uniformizer of \(F\), its leading
multiplier is \(\eta\), with \(\eta^5=\zeta\),
so \(\eta=\zeta^5\). If the nonidentity wild
elements have lower break \(b_F\), conjugating
\(g\varpi=\varpi+c\varpi^{b_F+1}+\cdots\)
shows that \(\tau\) acts on their leading
coefficients by \(\eta^{b_F}=\zeta^{5b_F}\).
The tame weights of the three quotient groups
are respectively \(0,6,18\). This gives
(B-RD87.5), since multiplication by five fixes
6 and 18 modulo 24.

Ramification groups are characteristic under
field automorphisms and hence tame-stable.
The two distinct tame characters in each
\(\operatorname{Gal}(E_i/L)\) imply that a proper
ramification line is one of these two axes.
There cannot be just one break: then the
norm-polynomial average would give equal breaks
for all cyclic quotients, contrary to (B-RD87.5).
For two breaks \(t<u\), the quotient killing the
last line has break \(t\), and the other displayed
quotient has the strictly larger break
\((u+4t)/5\). This proves the final assertion.
\end{proof}

\begin{proposition}[Three first-layer dimensions are impossible at \(p=5\)]
\label{prop:B-RD87-three}
\status{Proved; full independent root and companion audits passed.}
Under the stated actual \(125\)-comparison hypotheses,
one has \(\dim W_1\ne3\).
\end{proposition}
\begin{proof}
Suppose otherwise and apply the preceding lemma.
In the two good models the first-level low/high
valuations over \(L\) are \((5,95)\) and
\((15,45)\), respectively. Each parameter is
greater than \(e_K=20\), so all inverse roots over
a high target have valuation equal to that
target's valuation divided by 25, as in B85.
The same single-segment assertion holds over a
low target. Since \([E_i:L]=25\), their actual
primitive coordinates have the following
integral valuations and translation displacements:
\[
\begin{array}{c|c|r|r|r|c}
 \text{field}&\text{model}&A&D_0&D_1&H\\\hline
 E_1&1&95&125&2375&\text{diagonal}\\
 E_1&2&15&375&1125&\text{off-diagonal}\\
 E_2&1&5&125&2375&\text{off-diagonal}\\
 E_2&2&45&375&1125&\text{diagonal}
\end{array}                                             \tag{B-RD87.7}
\]
Here \(A\) is the point-coordinate valuation;
\(D_0,D_1\) are the low/high values of
\(v_{E_i}(gT(P)-T(P))\); and \(H\) is the line
of translations attaining \(D_1\). These are
actual point fields, shared under the given
torsion isomorphism. Formal addition multiplies
the translating coordinate by a unit, so the
displacement entries are exactly 25 times the
first-level valuations. Every \(A\) has
5-adic order one, and
\[
           v_{E_i}(5)=25\cdot24\cdot20=12000>2375.
                                                        \tag{B-RD87.8}
\]
Thus Lemma~\ref{lem:B-RD87-coefficient} applies to
each row without a mixed-characteristic equality
or truncation exception.

In \(E_1\), if the last ramification line were
off-diagonal, the first row would force its
larger cyclic quotient break to be
\(o_1=(125-95)/5=6\). But it would have to
exceed \(d\ge24\), a contradiction. Hence the
last line in \(E_1\) is diagonal. The second
row now gives
\[
                         d=(375-15)/5=72.             \tag{B-RD87.9}
\]
In \(E_2\), a diagonal last line would make
the third row give \(d=(125-5)/5=24\), contrary
to (B-RD87.9). An off-diagonal last line would
make the fourth row give
\(o_2=(375-45)/5=66\), which would have to
exceed \(d=72\), again impossible. These are
the only two ramification-line choices. The
assumed three-dimensional first layer therefore
cannot occur.
\end{proof}

\begin{corollary}[The tame residual ambiguity requires the full congruence kernel]
\label{cor:B-RD87-full}
\status{Proved; full independent root and companion audits passed.}
At \(p=5,e_K=20\), a different-special-character
actual \(125\)-congruent pair in the tame reducible
residual sector must have
\[
    W_1=W_2=M_2(\mathbf F_5),\qquad
    W=1+5M_2(\mathbf Z/125),\qquad |W|=5^8.           \tag{B-RD87.10}
\]
Its base field does not contain \(\mu_{25}\).
The necessary parameter pattern remains
\(\{25,75\}\), type A/B up to common quadratic twist.
\end{corollary}
\begin{proof}
B85 gives \(\dim W_1\ge3\), and the proposition
excludes equality. The ambient matrix space has
dimension four, so \(W_1=M_2(\mathbf F_5)\).
Actual fifth powers give \(W_1\subset W_2\),
so the terminal layer is the same full space.
Thus \(W\) has order \(5^8\), equal to that of
the full principal congruence kernel containing
it, and the groups are equal. The determinant
modulo 25 has nonzero trace on the first layer.
On \(L\) this is the actual cyclotomic character,
so \(\mu_{25}\) is not contained in \(L\), and
in particular not in \(K\). The parameter and
finite-type statement is precisely B84.
\end{proof}

\begin{remark}[Scope and the shared-quotient requirement]
\label{rem:B-RD87-scope}
\status{Proved; full independent root and companion audits passed.}
The statement concerns only the tame reducible
residual sector of the \(p=5,e_K=20\) problem.
It does not bound the other residual sectors,
construct the surviving full congruence image,
or settle the unrestricted least index, which
remains in \(\{2,3\}\).
The shared diagonal field \(D\) follows from
the actual same three-dimensional first layer.
Two arbitrary degree-25 extensions with
unrelated diagonal quotients would not give
this contradiction. No cyclotomic identification
of \(D\) is assumed.

The coefficient proof keeps the actual point
lifts fixed, retains the strict bound
\(D_1<v_{E_i}(5)\), and includes the possibility
that the first prime-to-five coefficient is
absent until cancellation forces it to exist.
The complete B60 and B81 fixed-coefficient
proofs and the B83--B85 actual field and
matrix arguments were reread. Serre's original
Section 3.1, pp.~126--128, was reopened and
read for the ramification definitions and
filtration conventions; its quotient theorem
refers to other proofs. The complete direct
norm-polynomial proof used here was freshly
read in Colmez, Proposition 3.7, pp.~14--15.
The tame leading-coefficient calculation is
written explicitly above. No unexamined
finite-residue assertion is imported.

The retained exact supplement is
\path{code/reduction_descent/verify_p5_tame_full_first_layer.py}.
The author run passed at 01:32:46.564572 UTC on
20 November 2026 in 0.03949103799823206 seconds:
three symbolic identities, 900 literal binomial
checks, eight strict inverse-polygon vertices,
all four coordinate rows, all 2,000 vectors in
the 16 possible diagonal first-layer spaces,
three exact tame congruences, 12 invariant-line
checks, and all four last-line alternatives.
The uniform coefficient argument and the actual
field constructions are proved above rather
than extrapolated from this finite supplement.
No failed mathematical computation preceded
this run; the exact executed code and output
are retained without replacement.

\smallskip
\noindent\emph{Independent audits.}
The complete root and companion proof, source and code audits passed.
Their literal author-output-preserving reruns took
\(0.037566576997051015\) and
\(0.034868966991780326\) seconds, respectively.
The exact audited candidate, per-auditor read scopes,
all code and outputs, and the promotion diff are retained.
\end{remark}


% Proved necessary profile after complete root and companion audits; no realization or exclusion claim.
\subsection{The necessary ramification profile of the surviving full image}

Keep the hypotheses and normalization of B87:
\(p=5\), \(e_K=20\), \(\mu_5\subset K\),
algebraically closed residue field, common tame good
field \(L/K\) of degree 24, and a different-special-character
actual \(125\)-comparison with tame reducible residual inertia.
The parameters are \(B_1=25,B_2=75\), with high lines
\(\mathbf F_5e_1,\mathbf F_5e_2\), respectively.
B87 gives \(W_1=M_2(\mathbf F_5)\). Let
\(F_2=L(M[25])\) be the actual common second division field,
so \(\operatorname{Gal}(F_2/L)=W_1\).
All breaks below use normalized valuations and the convention
\(b(g)=v(g\varpi-\varpi)-1\).

\begin{proposition}[Both actual point-field ramification filtrations]
\label{prop:B-RD88-point-fields}
\status{Proved; full independent root and companion audits passed.}
Let \(E_i=L(P_i)\), where \(P_i\) is any primitive
second-level point with \(5P_i=e_i\).
Each \(E_i/L\) is Galois of degree 25.
Its two lower breaks and its two coordinate cyclic
quotient breaks are necessarily
\[
\begin{array}{c|c|c|c}
 &\text{lower breaks}&\text{diagonal quotient}&\text{off-diagonal quotient}\\\hline
 E_1/L&(6,336)&72&6\\
 E_2/L&(24,234)&24&66.
\end{array}                                             \tag{B-RD88.1}
\]
The diagonal quotients here are different actual fields;
they are the \(A_{11}\) and \(A_{22}\) quotients of \(W_1\).
\end{proposition}
\begin{proof}
Evaluation on \(e_i\) is surjective. Since the first-level
module is \(L\)-rational, its kernel is the actual stabilizer
of every such lift. It is a normal subgroup of the additive
second-level image, proving the field assertion. In \(E_1\)
the translation coordinates are \((A_{11},A_{21})\);
in \(E_2\) they are \((A_{12},A_{22})\).
Their tame weights are respectively \((0,6)\) and
\((18,0)\). The tame congruence and norm-quotient arguments
of Lemma~\ref{lem:B-RD87-fields} apply separately to each
field, without identifying their diagonal quotients.

Write \(d_i\) and \(o_i\) for their diagonal and
off-diagonal cyclic quotient breaks. Thus
\[
 d_i\in24\mathbf Z_{>0},\qquad
 o_1\equiv6\pmod{24},\qquad o_2\equiv18\pmod{24}.
                                                        \tag{B-RD88.2}
\]
There are two distinct lower breaks \(t<u\). The last
ramification line is an eigenaxis, and the larger cyclic
quotient break is \((u+4t)/5\), while the smaller is \(t\).
The four actual coordinate rows are still (B-RD87.7),
with the same strict cutoff \(12000\).

We use explicitly the fixed-coordinate minimum from
(B-RD87.3). For a row with leading valuation \(A\),
let \(J\) be its first nonzero exponent prime to five.
At lower break \(b\) the first possible displacement
degree is \(\min(A+5b,J+b)\); a strict smaller minimum
than the actual displacement can be removed only by
the tie \(J=A+4b\). Every prime exponent present
satisfies \(J>A\), and \(5\nmid J\).

First consider \(E_1\). An off-diagonal last line
would make the first row of (B-RD87.7) give
\(o_1=6>d_1\), impossible by (B-RD88.2).
The last line is therefore diagonal, and the second
row gives \(d_1=72\). Consequently
\[
       t=o_1\in\{6,30,54\},\qquad u=360-4t.          \tag{B-RD88.3}
\]
For the second row, whose high-displacement line is not
the last line, the compulsory prime degree is
\(J=15+4t\). If \(t=30\), it is 135, divisible by
five, a contradiction. If \(t=54\), apply the first
row instead: its displacement outside the last line
is 125, whereas \(A+5t=95+270=365\).
It would require \(J=125-54=71<A=95\), impossible.
Only \(t=6,u=336\) remains.

Now consider \(E_2\). If the last line were diagonal,
the third row would give \(d_2=24\). Necessarily
\(o_2=t=18\) and \(u=120-72=48\).
In the fourth row the smaller-break elements have
displacement 375, while \(A+5t=45+90=135\).
Cancellation forces \(J=45+4\cdot18=117\).
At the last line's break 48, this prime term changes
uniquely first, at degree \(117+48=165\), less than
the other possible degree \(45+5\cdot48=285\).
Its actual displacement would have to be 1125,
a contradiction.

Thus the last line in \(E_2\) is off-diagonal.
The fourth row gives \(o_2=66\), and hence
\[
             t=d_2\in\{24,48\},\qquad u=330-4t.      \tag{B-RD88.4}
\]
If \(t=48,u=138\), the third row has displacement
125 outside its high line, whereas \(A+5t=245\).
It follows that \(J=125-48=77\). At the last line,
the unique first change would then be
\(77+138=215<5+5\cdot138=695\), contradicting its
actual displacement 2375. Therefore \(t=24,u=234\).
This proves every entry of (B-RD88.1). All first-change
comparisons were strictly below \(v_{E_i}(5)=12000\);
none assumes an unproved independence of higher coefficients.
\end{proof}

\begin{corollary}[The full second division field and its cyclotomic quotient]
\label{cor:B-RD88-cyclotomic}
\status{Proved; full independent root and companion audits passed.}
The upper breaks of \(F_2/L\), together with the
corresponding coordinate characters, are
\[
\begin{array}{c|rrrr}
 \text{character}&A_{21}&A_{22}&A_{12}&A_{11}\\\hline
 \text{upper break}&6&24&66&72.
\end{array}                                             \tag{B-RD88.5}
\]
Its lower breaks are \(6,96,1146,1896\), with successive
ramification ranks \(4,3,2,1\) over \(\mathbf F_5\).
Among its 156 cyclic degree-five subfields, the numbers
with breaks \(6,24,66,72\) are \(1,5,25,125\), respectively.
The actual cyclotomic quotient has
\[
 b\bigl(L(\mu_{25})/L\bigr)=72,
 \qquad b\bigl(K(\mu_{25})/K\bigr)=3.                \tag{B-RD88.6}
\]
Both displayed extensions have degree five.
\end{corollary}
\begin{proof}
The four coordinate cyclic quotients have exactly
the breaks in (B-RD88.5), by the proposition.
We recall why a nonzero sum of independent cyclic
characters with different breaks has the larger break.
Factor through the actual degree-25 quotient defined
by the two characters. Its two quotient breaks differ,
so it has two ramification breaks and a last line.
The smaller-break character vanishes on the last line,
whereas the larger-break one does not. Their sum also
does not, and the same exact coset-average formula
gives it the larger break. This proves the assertion
directly, without an equality of unrelated point fields.
Induction now shows that every nonzero combination of
the four coordinate characters has break equal to the
largest break of its nonzero coordinates.

The quotient compatibility of upper numbering therefore
gives (B-RD88.5) and the stated successive ranks. Here
the complete norm-polynomial/Herbrand argument is the
one already checked in
\cite[Proposition 3.7 and Theorem 3.10, full proofs,
pp.~14--16]{ColmezCorpsLocaux}; the normalization is
also that of Serre's original Section 3.1.
Integrating the indices \(1,5,25,125\) gives the lower
breaks explicitly:
\[
 6,\quad6+5(24-6)=96,\quad
 96+25(66-24)=1146,\quad
 1146+125(72-66)=1896.                                 \tag{B-RD88.7}
\]
The number of one-dimensional character spaces whose
last nonzero coordinate is the \(j\)th in that order
is \(5^{j-1}\). This gives all 156 cyclic subfields.

On \(L\), the determinant of the actual good height-two
representation is the cyclotomic character. Modulo 25,
\(\det(1+5A)=1+5\operatorname{tr}(A)\).
Thus \(L(\mu_{25})\) is precisely the quotient field
for \(A_{11}+A_{22}\), which has break 72 by the
unequal-break sum argument. It is nontrivial, and
\(\mu_5\subset K\), so it and \(K(\mu_{25})/K\)
have degree five.

For clarity, the final tame scaling is an actual
field calculation. Put \(C=K(\mu_{25})\) and
\(D=LC\). The extensions \(C/K\) and \(L/K\)
are Galois of coprime degrees 5 and 24. Hence
\(D/C\) is tame of degree 24 and
\(\operatorname{Gal}(D/K)\) is their direct product.
Choose a uniformizer \(x\) of \(C\). The tame
uniformizer/Hensel argument provides a uniformizer
\(z\) of \(D\) with \(z^{24}=x\); see the complete
argument of \cite[Proposition 2.12, p.~11]{ColmezCorpsLocaux}.
For a nonidentity \(g\in\operatorname{Gal}(D/L)\),
the ratio \(g(z)/z\) has residue one, and
\[
   (g(z)/z)^{24}=g(x)/x.
\]
As 24 is a unit, the valuations of these ratios minus
one agree. Thus
\(b(D/L)=24b(C/K)\). The left side is 72, giving
the asserted break three over \(K\).
\end{proof}

\begin{remark}[A necessary profile, not a final exclusion]
\label{rem:B-RD88-scope}
\status{Proved; full independent root and companion audits passed.}
This result constrains the actual remaining \((4,4)\)
image in the tame reducible sector at \(p=5,e_K=20\).
It neither constructs an ambiguous pair nor proves that
the required cyclotomic break three is impossible over
the allowed base fields. In particular, the fixed value
\(24^2\) from B86 concerned the different hypothesis
\(\mu_{25}\subset K\) and is not reused here.
The two diagonal quotient fields are distinct because
\(A_{11},A_{22}\) are independent characters on the
actual full first layer. No false shared-diagonal-field
assertion is imported from the three-dimensional case.

The only original-source inputs beyond the earlier
proved dependencies are the complete Colmez tame
uniformizer and ramification quotient/Herbrand proofs,
already freshly read and scoped in B86. No additional
finite-residue assumption or priority claim is made.

The exact retained supplement is
\path{code/reduction_descent/verify_p5_full_image_ramification.py}.
Its first author run passed at 01:46:08.611352 UTC
on 20 November 2026 in 0.001338374990154989 seconds.
It checks both initially excluded last-line cases,
all three remaining \(E_1\) residue choices and both
\(E_2\) choices, the four upper-to-lower intervals,
all 624 nonzero coordinate characters and their
156 projective classes, the trace quotient and tame
scaling, and the equality of the different computed
from the lower filtration and from the character
conductors. These are exact finite consequences of
the preceding proof; no bounded search for an
unknown ramification break replaces that proof.
No failed mathematical execution preceded this run;
its exact code and output remain retained.

\smallskip
\noindent\emph{Independent audits.}
The complete root and companion proof, source and code audits passed.
Their literal author-output-preserving reruns took
\(0.0014191469963407144\) and
\(0.0013393660046858713\) seconds, respectively.
The exact audited candidate, full read scopes, code,
outputs and promotion diff are retained.
\end{remark}


% Proved base-field example after complete audits; no BT comparison is constructed.
\subsection{The required cyclotomic break three is attainable}

\begin{lemma}[An explicit allowed base with the required cyclotomic quotient]
\label{lem:B-RD89-base}
\status{Proved; full independent root and companion audits passed.}
Let \(k\) be an algebraically closed field of characteristic five,
put \(K_0=W(k)[1/5]\), choose a primitive fifth root \(\zeta_5\),
and set
\[
 K_1=K_0(\zeta_5),\qquad \pi_1=\zeta_5-1.
\]
Let \(t\) be a root of
\[
       (1+T)^5-\zeta_5+\pi_1^4T^2=0,
       \qquad K=K_1(t).                              \tag{B-RD89.1}
\]
Then \(K\) is complete discretely valued, has residue field \(k\),
and satisfies \(e_K=20\) and \(\mu_5\subset K\).
The extension \(C=K(\mu_{25})/K\) is cyclic of degree five,
with lower break three. The extension \(K(\mu_{125})/K\)
is cyclic of degree 25. If \(L/K\) is the tame extension
of degree 24, then \(L(\mu_{25})/L\) has break 72.
\end{lemma}
\begin{proof}
The usual cyclotomic Eisenstein polynomial
\(\Phi_5(1+X)\) shows that \(K_1/K_0\) has degree four,
with uniformizer \(\pi_1\). Expanded over \(K_1\), the
polynomial in (B-RD89.1) is
\[
 T^5+5T^4+10T^3+(10+\pi_1^4)T^2+5T-\pi_1.
                                                        \tag{B-RD89.2}
\]
It is Eisenstein at \(\pi_1\), so \([K:K_1]=5\),
\(t\) is a uniformizer, \(v_K(\pi_1)=5\), and
\(v_K(5)=20\). Completeness and the residue-field assertion
follow from this finite totally ramified extension.

Choose \(\zeta_{25}\) with \(\zeta_{25}^5=\zeta_5\),
and put \(\alpha=\zeta_{25}-1\). The element
\(x=\alpha-t\) satisfies the exact shifted polynomial
\[
 -\pi_1^4t^2+
 \sum_{j=1}^5\binom5j(1+t)^{5-j}X^j=0.               \tag{B-RD89.3}
\]
Its constant coefficient has \(K\)-valuation 22,
its coefficients in degrees one through four have
valuation 20, and its leading coefficient is one.
Thus every intermediate point lies strictly above
the chord from \((0,22)\) to \((5,0)\). All its roots,
and in particular \(x\), have valuation \(22/5\).

Since \(\mu_5\subset K\), adjoining \(\zeta_{25}\)
gives a Galois extension of degree dividing five.
The denominator in \(v_K(x)=22/5\) forces the degree
to equal five. Normalizing \(v_C\), one has
\(v_C(x)=22\). For every nonidentity
\(g\in\operatorname{Gal}(C/K)\),
\[
  gx-x=(\zeta_5^j-1)\zeta_{25}
       \quad\text{for some }j\in\{1,2,3,4\},
  \qquad v_C(gx-x)=25.                               \tag{B-RD89.4}
\]
The leading valuation 22 is prime to five, so the
exact lower-break equality of
Lemma~\ref{lem:B-RD60-break} applies to this actual
coordinate and gives \(b(C/K)=25-22=3\).
This also proves \(\mu_{25}\not\subset K\).

For the next cyclotomic layer,
\(K_1(\mu_{125})/K_1\) is cyclic of degree 25,
with unique degree-five subfield \(K_1(\mu_{25})\).
Its intersection with the degree-five extension
\(K/K_1\) is therefore \(K_1\): a nontrivial intersection
would make \(K=K_1(\mu_{25})\), contrary to what was
just proved. Hence \(K(\mu_{125})/K\) is cyclic of
degree 25. Finally the degree-24 tame extension is
disjoint from every such cyclotomic five-extension.
The tame uniformizer/Hensel calculation written in
Corollary~\ref{cor:B-RD88-cyclotomic} gives
\(b(L(\mu_{25})/L)=24\cdot3=72\).
\end{proof}

\begin{remark}[Only the base and its cyclotomic character are realized]
\label{rem:B-RD89-scope}
\status{Proved; full independent root and companion audits passed.}
The example meets the cyclotomic part of the necessary
profile in B88 and gives the full order-25 cyclotomic
image at level 125. It does not construct the remaining
three independent first-layer characters, the required
height-two formal groups, or a common actual torsion
comparison. In particular it is neither a lower bound
example for the unresolved least index nor an existence
proof for the full \((4,4)\) image sector. Its role is
to rule out the unsupported shortcut that break three
could never occur over an allowed base field.

Only Eisenstein and Newton-polygon arithmetic, the
already proved exact B60 lower-break equality, and
the explicit B88 tame scaling are used. No additional
external classification theorem is imported.

The retained exact checker is
\path{code/reduction_descent/verify_p5_cyclotomic_break_three_example.py}.
The completed author run passed at 01:55:21.154128 UTC
on 20 November 2026 in 0.04096559699974023 seconds.
It verifies the two symbolic polynomial identities,
all four strict Newton vertices, and the Eisenstein
coefficients of the three cyclotomic polynomials used.
The initial execution stopped at a structural SymPy
equality comparing an expanded expression with an
unexpanded one. Its exact program, standard streams
and failure record are preserved; the corrected check
tests that the expanded polynomial difference is zero.
No mathematical assertion was inferred from the failed run.

\smallskip
\noindent\emph{Independent audits.}
The complete root and companion proof and code audits passed.
Their literal author-output-preserving reruns took
\(0.043666591998771764\) and
\(0.04821994699887\) seconds, respectively.
Both audits reread the complete B60 and B88 dependencies;
neither claims a new external-source opening. The exact
candidate, failed and successful programs, outputs, and
promotion diff remain retained.
\end{remark}


% Proved after full root and companion audits; exact actual-field candidate retained.
\subsection{Third-level point fields exclude the tame full-image comparison}

Throughout this subsection, assume for contradiction the remaining tame
comparison of B87--B88. Thus \(K\) is complete discretely valued of mixed
characteristic \((0,5)\), with algebraically closed residue field,
\(e_K=20\), and \(\mu_5\subset K\). Two whole potentially supersingular
height-two groups have isomorphic actual \(125\)-torsion, distinct finite
special inertia characters, and tame reducible residual inertia.
After the already justified common finite twist and interchange, take
\(B_1=25,B_2=75\), with high lines \(\mathbf F_5e_1\) and
\(\mathbf F_5e_2\). The common tame good field \(L/K\) has degree 24.
Put \(F=L(M[125])\), \(F_2=L(M[25])\). The actual comparison identifies
these fields for the two models. B87 gives
\[
 W=\operatorname{Gal}(F/L)=1+5M_2(\mathbf Z/125\mathbf Z),
 \qquad U=\operatorname{Gal}(F/F_2)=1+25M_2(\mathbf F_5).
 \tag{B-RD90.1}
\]
Every lower break is \(b(g)=v(g\varpi-\varpi)-1\), for the normalized
valuation on the top field of the indicated extension.

\begin{lemma}[The actual third-point towers and their two lower breaks]
\label{lem:B-RD90-point-towers}
\status{Proved; full independent root and companion audits passed.}
One can choose a common basis \(Q_1,Q_2\) of the actual \(125\)-torsion
whose images \(25Q_i=e_i\) lie on the two displayed residual axes and
whose point fields are stable under a tame lift. Put
\[
 P_i=5Q_i,\qquad E_i=L(P_i),\qquad E_i^+=L(Q_i).
\]
Then \(E_i^+/E_i\) is Galois of degree 25, and
\[
\begin{array}{c|c|c|c}
 &\text{lower breaks}&\text{diagonal cyclic quotient}&
       \text{off-diagonal cyclic quotient}\\\hline
 E_1^+/E_1&(606,6936)&1872&606\\
 E_2^+/E_2&(624,6834)&624&1866.
\end{array}                                                   \tag{B-RD90.2}
\]
The terms diagonal and off-diagonal refer to the actual translations
of \(Q_i\) by the first-level module, with respect to \(e_1,e_2\).
\end{lemma}
\begin{proof}
The finite extension \(F/K\) is Galois, with normal 5-group \(W\) and
cyclic quotient of order 24. A lift of a generator has order
\(24\cdot5^r\); a suitable power therefore gives a lift \(\tau\) of
order 24. Its action on the torsion lattice diagonalizes integrally:
the two distinct residual eigenvalues have unique Teichmuller lifts,
and the corresponding coprime factors of \(X^{24}-1\) give integral
idempotents. Choose \(Q_i\) on these two direct summands. Then
\(\tau Q_i\) is a unit multiple of \(Q_i\), so \(E_i,E_i^+\) are
stable under \(\tau\). This choice works in both models through the
actual torsion isomorphism.

Write a matrix in \(W\) as \(1+5X\), with
\(X\in M_2(\mathbf Z/25\mathbf Z)\). The stabilizers of \(P_i\)
and \(Q_i\) are respectively
\[
 H_i=\{1+5X:X e_i=0\pmod5\},\qquad
 S_i=\{1+5X:X e_i=0\pmod{25}\}.
 \tag{B-RD90.3}
\]
Their orders are \(5^6\) and \(5^4\). On \(H_i\), the map
\(g\mapsto gQ_i-Q_i\) takes values in the \(L\)-rational first-level
module. It is a homomorphism, and the terminal matrices \(U\) already
surject onto that module. Its kernel is \(S_i\). This proves that
\(E_i^+/E_i\) is Galois of degree 25, and that every nonzero translation
occurs. It does not assert that \(E_i^+/L\) is Galois.

The first-level low/high valuations over \(L\) are \((5,95)\) and
\((15,45)\). The exact inverse polygon (B-RD76.3), whose coefficient
bounds allow wild absolute index, applies twice. Already the first
inverse over a high target is uniform because \(B_i>e_K\); all later
inverses are uniform. Thus division by five twice divides each of
these valuations by \(25^2\). Since \([E_i^+:L]=625\), the coordinate
valuations in \(E_i^+\) equal the first-level valuations again.
Translation by a first-level point changes a formal coordinate by
that point's coordinate times an integral unit. The complete rows are
therefore
\[
\begin{array}{c|c|r|r|r|c}
 \text{field}&\text{model}&A&D_0&D_1&\mathcal H\\\hline
 E_1^+&1&95&3125&59375&\text{diagonal}\\
 E_1^+&2&15&9375&28125&\text{off-diagonal}\\
 E_2^+&1&5&3125&59375&\text{off-diagonal}\\
 E_2^+&2&45&9375&28125&\text{diagonal}.
\end{array}                                                   \tag{B-RD90.4}
\]
Here \(A=v_{E_i^+}(T(Q_i))\), while the displacement is \(D_1\) on
the line \(\mathcal H\) and \(D_0\) off it. All four values \(A\)
have 5-adic order one. The strict coefficient cutoff is
\(v_{E_i^+}(5)=625\cdot480=300000\), above every displacement in
the table.

Tame conjugation on these translation groups has the same two weights
as in B87: zero and 6 for \(i=1\), zero and 18 for \(i=2\), modulo
24. Indeed the eigenvalue of \(Q_i\) cancels in the cocycle
\(\tau g\tau^{-1}Q_i-Q_i\). On a cyclic degree-five quotient over
\(E_i\), the total ramification degree over \(L\) is 125, so its
leading tame uniformizer multiplier \(\eta\) satisfies
\(\eta^{125}=\zeta\). Thus \(\eta=\zeta^5\). The conjugation
calculation of B87 again gives a positive diagonal break divisible
by 24, and an off-diagonal break congruent to 6 or 18 respectively.
The two quotient breaks differ; the last ramification line is one
of the two eigenaxes. For lower breaks \(t<u\), their cyclic quotient
breaks are \(t\) and \((u+4t)/5\).

We give the coefficient exclusions explicitly. For any row of
(B-RD90.4), let \(J\) be the first nonzero coordinate exponent prime
to five. The proof of Lemma~\ref{lem:B-RD87-coefficient} gives the
first possible displacement degree
\[
                  \min\{A+5b,J+b\}.                 \tag{B-RD90.5}
\]
A strict minimum cannot cancel. If this minimum is strictly less
than the actual displacement, cancellation therefore requires
\(J=A+4b\); in every case \(J>A\) and \(5\nmid J\).
The value \(J=\infty\) is allowed before these constraints force
a finite exponent. All uses are strictly below the stated cutoff.

In \(E_1^+/E_1\), if the last line were off-diagonal, the first
row and Lemma~\ref{lem:B-RD87-coefficient} would give the larger
cyclic break \((3125-95)/5=606\). Hence \(t<606\) and
\(u=3030-4t\). In the second row,
\(15+5t<9375\), so cancellation forces \(J=15+4t\).
At \(u>t\) this prime term changes uniquely first, at
\(J+u=3045\), instead of 28125. This is impossible.
The last line is diagonal. The second row now gives its cyclic
break \((9375-15)/5=1872\), so
\[
             t\equiv6\pmod{24},\quad 0<t<1872,
             \qquad u=9360-4t.                      \tag{B-RD90.6}
\]
For the first row, if \(t<606\), then \(95+5t<3125\);
cancellation forces \(J=95+4t\), whose displacement at \(u\)
is 9455, instead of 59375. If \(t>606\), one must have
\(J=3125-t\). At \(u\), this term changes strictly before the
leading term, since
\[
 J-95-4u=15t-34410<0 \quad(t<1872).
\]
Its displacement is \(12485-5t<59375\), again impossible.
Only \(t=606,u=6936\) remains. The compulsory first prime
exponents in the two rows are then 27839 and 2439; these are
indeed prime to five. This proves the first row of (B-RD90.2).

In \(E_2^+/E_2\), a diagonal last line would give the larger
quotient break \((3125-5)/5=624\), so \(t<624\) and
\(u=3120-4t\). The fourth row forces \(J=45+4t\), whose
unique displacement at \(u\) is 3165, instead of 28125.
The last line must instead be off-diagonal. Its cyclic break
is \((9375-45)/5=1866\), so
\[
             t\equiv0\pmod{24},\quad 0<t<1866,
             \qquad u=9330-4t.                      \tag{B-RD90.7}
\]
For the third row, \(t<624\) forces \(J=5+4t\) and the
incorrect displacement 9335 at \(u\). If \(t>624\), it forces
\(J=3125-t\), and
\(J-5-4u=15t-34200<0\). The unique displacement is
\(12455-5t<59375\), again impossible. Thus \(t=624,u=6834\).
The two first prime exponents are 27341 and 2541. This proves
the second row of (B-RD90.2).
\end{proof}

\begin{proposition}[The forced terminal filtration in the actual full field]
\label{prop:B-RD90-terminal}
\status{Proved; full independent root and companion audits passed.}
Under the same contradictory hypotheses, the terminal extension
\(F/F_2\) would have lower breaks
\[
                7896,\quad11646,\quad769146,\quad937896,
 \tag{B-RD90.8}
\]
with the successive coordinate directions removed in the order
\(B_{22},B_{21},B_{12},B_{11}\), for \(1+25B\in U\).
For elements outside \(U\), the lower breaks in \(F/L\) would
equal their first-layer breaks in \(F_2/L\), namely
\(6,96,1146,1896\), in the coordinate order
\(A_{21},A_{22},A_{12},A_{11}\).
\end{proposition}
\begin{proof}
First compute the actual intersections, rather than imposing
independence of abstract extensions. The subgroups in (B-RD90.3)
satisfy \(US_i=H_i\). Consequently
\[
 F_2\cap E_i^+=E_i,\qquad D_i:=F_2E_i^+=F^{U\cap S_i}.
 \tag{B-RD90.9}
\]
The quotient \(H_i/(U\cap S_i)\) has order \(5^4\) and is
elementary abelian. For its first-layer matrices \(A,C\), both
\(Ae_i=Ce_i=0\), so \([A,C]e_i=0\). The exact commutator
and fifth-power identities modulo 125 are
\[
 [1+5A,1+5C]=1+25[A,C],\qquad
       (1+5A+25B)^5=1+25A.
 \tag{B-RD90.10}
\]
Both right sides lie in \(U\cap S_i\). Hence \(D_i/E_i\) is
the actual elementary-abelian compositum of the disjoint
degree-25 extensions \(F_2/E_i\) and \(E_i^+/E_i\).

The subgroup filtration of the already proved \(F_2/L\) profile
in B88 gives lower breaks \((96,1146)\) for \(F_2/E_1\), with
characters \(A_{22},A_{12}\), and \((6,1896)\) for \(F_2/E_2\),
with characters \(A_{21},A_{11}\). Their upper breaks are
\((96,306)\) and \((6,384)\), respectively. The new extensions
have upper breaks \((606,1872)\) and \((624,1866)\), by
(B-RD90.2). All four breaks in each compositum are distinct.
The unequal-break sum argument proved in B88 therefore gives
all its character conductors, without assuming any further
linear independence of coefficients. Exact upper-to-lower
integration yields
\[
\begin{array}{c|c|c}
 &\text{upper breaks of }D_i/E_i&\text{lower breaks of }D_i/E_i\\\hline
 i=1&(96,306,606,1872)&(96,1146,8646,166896)\\
 i=2&(6,384,624,1866)&(6,1896,7896,163146).
\end{array}                                                   \tag{B-RD90.11}
\]
The last two character directions in each row are precisely
those from \(E_i^+/E_i\). Restricting the lower filtration to
\(\operatorname{Gal}(D_i/F_2)\) and then passing to upper
numbering therefore gives
\[
\begin{array}{c|c|c|c}
 &\text{lower breaks}&\text{upper breaks}&\text{terminal characters}\\\hline
 D_1/F_2&(8646,166896)&(8646,40296)&B_{21},B_{11}\\
 D_2/F_2&(7896,163146)&(7896,38946)&B_{22},B_{12}.
\end{array}                                                   \tag{B-RD90.12}
\]
Here the characters are the restrictions of the actual
translation maps \(g\mapsto gQ_i-Q_i\) to \(U\), hence are
exactly the indicated matrix entries, not chosen abstract
identifications of isomorphic Galois groups.

Since \((U\cap S_1)\cap(U\cap S_2)=1\), one has
\(D_1D_2=F\); their degree-25 extensions over \(F_2\) are
linearly disjoint. Thus \(F/F_2\) is elementary abelian of
rank four with coordinate upper breaks
\[
       B_{22}:7896,\quad B_{21}:8646,\quad
       B_{12}:38946,\quad B_{11}:40296.
 \tag{B-RD90.13}
\]
Integrating indices \(1,5,25,125\) gives (B-RD90.8).
These are the actual lower subgroup breaks in \(F/L\),
because lower numbering is preserved on subgroups.

Every nonidentity element of \(U\) consequently has break
at least 7896. For any \(g\notin U\), the exact coset-average
formula for \(F_2/L\) shows that some lift of its image has
break less than 7896: otherwise the quotient break would
be at least 7896, contrary to B88. Multiplication by any
element of \(U\) preserves such a smaller break, by the
valuation inequality for composition of automorphisms.
All lifts therefore have the same break, equal to the
quotient break. This proves the final assertion and
completely identifies every nonidentity break in \(W\).

Only subgroup restriction and the exact norm-polynomial
quotient/Herbrand formulas are used here, with the full proofs
already read in \cite[Proposition 3.7, Theorem 3.10 and Proposition 3.13(ii), full proofs, pp.~14--16]
{ColmezCorpsLocaux}. No Hasse--Arf theorem or assertion about
unrelated extensions is required.
\end{proof}

\begin{theorem}[No different-special-character tame residual comparison at 125]
\label{thm:B-RD90-tame-exclusion}
\status{Proved; full independent root and companion audits passed.}
For \(p=5,e_K=20\) in the stated setting, two whole potentially
supersingular height-two groups with isomorphic actual
\(125\)-torsion and tame reducible residual inertia have the
same finite special inertia characters.
\end{theorem}
\begin{proof}
Otherwise the preceding reductions give the actual full image
(B-RD90.1) and its forced filtration. On the good field \(L\),
the determinant is the cyclotomic character. Its image is the
full group \(1+5\mathbf Z/125\mathbf Z\), of order 25.
The actual quotient field is therefore \(C=L(\mu_{125})\),
and the determinant kernel has order \(5^6=15625\).

Fix a nonidentity order-five element of this quotient, of
determinant \(1+25r\), \(r\in\mathbf F_5^\times\).
Write each element uniquely as \(g=1+5A+25B\), with entries
of \(A,B\) chosen from \(0,1,2,3,4\). A lift in this coset
must have \(\operatorname{tr}(A)=0\) modulo five. For each
such \(A\), exactly 125 matrices \(B\) have the specified
determinant: the second determinant coordinate is an affine
trace equation, including the carry in \(\operatorname{tr}(A)\)
and \(\det A\). No incorrect entrywise linearization of
the full determinant is used.

Among the 124 nonzero traceless first-layer matrices,
100 have \(A_{21}\ne0\), 20 have \(A_{21}=0,A_{22}\ne0\),
and four have only \(A_{12}\ne0\). Their breaks are
6, 96, and 1146. For \(A=0\), the equation is exactly
\(B_{11}+B_{22}=r\). The terminal filtration just proved
gives respectively 100, 20, four, and one solutions with
breaks 7896, 11646, 769146, and 937896. Hence the exact
coset counts are
\[
\begin{array}{c|rrrrrrr}
 b(g)&6&96&1146&7896&11646&769146&937896\\\hline
 \#g&12500&2500&500&100&20&4&1.
\end{array}                                                   \tag{B-RD90.14}
\]
Their sum is 15625 and their weighted break sum is 5925000.
For this nonidentity quotient element \(\bar g\), the exact
formula of \cite[Proposition 3.7 and full proof]{ColmezCorpsLocaux}
is
\[
 b_C(\bar g)+1
   =\frac{1}{15625}\sum_{g\mapsto\bar g}\bigl(b_F(g)+1\bigr).
\]
It would give
\[
                  b_C(\bar g)=\frac{5925000}{15625}
                              =\frac{1896}{5}.        \tag{B-RD90.15}
\]
This is impossible: \(b_C(\bar g)=v_C(\bar g\varpi_C-\varpi_C)-1\)
is an integer for the normalized discrete valuation of the
actual field \(C\). The contradiction proves the theorem.
\end{proof}

\begin{remark}[Scope and exact supplement]
\label{rem:B-RD90-scope}
\status{Proved; full independent root and companion audits passed.}
This excludes only the tame reducible residual branch left
by B87--B88. The irreducible-residual and nontrivial-residual-wild
branches retain their previously proved restrictions; no
actual comparison in those branches is constructed or excluded
here. In particular the unrestricted first-wild least-index
interval remains \(\{2,3\}\). The allowed base constructed
in B89 is unaffected: it never supplied the other division
fields or a height-two torsion comparison.

The retained exact program is
\path{code/reduction_descent/verify_p5_third_point_ramification.py}.
Its first run passed at 02:14:42.322512 UTC on 20 November 2026
in 0.08913616598874796 seconds. It checks all 206 lower-break
possibilities inside the four bounds proved above, obtaining
the two displayed survivors; it reconstructs every quotient,
subgroup, and compositum break by exact rational arithmetic.
It also enumerates all 62500 actual matrices in the four
nonidentity order-five determinant fibers, independently
evaluating their determinants modulo 125, and checks 1250
stabilizer commutators. These finite checks supplement the
field identifications and uniform coefficient argument;
they do not replace them or assert realization of a numerical
profile. No failed execution preceded the recorded run.
All code and output are retained. The complete B60, B76,
B87, and B88 arguments were reread; no new external-source
opening is claimed for this candidate.

\smallskip
\noindent\emph{Independent audits.}
The complete root and monodromy proof and exact-code audits passed.
Their literal author-output-preserving reruns took
\(0.08867679999093525\) and
\(0.08529232500586659\) seconds, respectively.
The root subsequently reread the original Colmez proofs through
Proposition 3.13 and retained a separate source addendum. The canonical
citation names Proposition 3.7, Theorem 3.10 and Proposition 3.13(ii)
precisely; this corrects only the candidate's bibliographic labels.
The exact candidate, source scopes, code, outputs and promotion diff
remain retained.
\end{remark}


% Proved after complete root and monodromy audits; exact candidate retained.
\subsection{High quotients in the remaining residual-wild comparison}

Let \(K\) be complete discretely valued of mixed characteristic
\((0,p)\), with algebraically closed residue field, \(p\ge5\),
\(\mu_p\subset K\), and \(e=e_K=p(p-1)\). Put
\(q=p^2\), \(N=q-1\), \(s=(p+1)e=pN\), and let
\(L/K\) be the common tame good extension of degree \(N\).
We compare actual generic \(p^3\)-torsion of two whole potentially
supersingular height-two groups. Assume their residual wild action
is nontrivial. Its unique fixed line is their common canonical
high line \(H\), by Lemma~\ref{lem:B-RD82-Newton}. Write
\[
 0<B_i<pe,\qquad \alpha_i=s-B_i,\qquad \beta_i=B_i/p
                                                        \tag{B-RD91.1}
\]
for the first-level parameters and their high and low coordinate
valuations over \(L\). The common first-level division field
\(R/L\) is cyclic of degree \(p\).

\begin{lemma}[The first-level strict comparison]
\label{lem:B-RD91-first-break}
\status{Proved; full independent root and companion audits passed.}
If \(p\nmid B_1\), then the actual lower break of \(R/L\) is
\[
              b(R/L)=ps-(p+1)B_1.                    \tag{B-RD91.2}
\]
For any second model with the same actual first-level module and
identified high line, one has \(B_1\ge B_2\). In particular,
if both parameters are prime to \(p\), they are equal; if
\(p\mid B_2\), the inequality is strict.
\end{lemma}
\begin{proof}
Choose corresponding nonzero first-level points outside \(H\).
Their stabilizers are the same, and every such point generates
\(R\), since each nonidentity residual unipotent translates it
by a nonzero point of \(H\). The normalized \(R\)-valuations
of their formal coordinates are \(B_i\). Every nonidentity
\(g\in\operatorname{Gal}(R/L)\) displaces the first coordinate
by valuation \(p\alpha_1\). The exact prime-valuation equality
of Lemma~\ref{lem:B-RD60-break} therefore gives (B-RD91.2).
The same lemma's unconditional inequality in the other coordinate
gives
\[
 p(s-B_2)\ge B_2+ps-(p+1)B_1,
 \qquad (p+1)(B_1-B_2)\ge0.                          \tag{B-RD91.3}
\]
If both parameters are prime to \(p\), interchange the models.
If only the first is prime, equality is impossible by divisibility.
These are actual coordinates of the same degree-\(p\) point field;
no bound on the full division field is used.
\end{proof}

\begin{proposition}[Both parameters must be divisible by \(p\)]
\label{prop:B-RD91-divisibility}
\status{Proved; full independent root and companion audits passed.}
If the finite special characters of the two groups differ, then
\(p\mid B_1\) and \(p\mid B_2\).
\end{proposition}
\begin{proof}
Suppose \(p\nmid B_1\). If the second parameter is also prime,
the lemma gives \(B_1=B_2\); equality of the actual high-line
characters \(\zeta^{c_i-B_i}\) then gives \(c_1=c_2\pmod N\).
Thus assume \(p\mid B_2\). The lemma gives \(B_1>B_2\).

Form the actual good integral quotients \(Q_i=M_i/H\), using
the schematic closures of the common line. The product calculation
of Lemma~\ref{lem:B-RD69-quotients} gives their parameters
\[
 D_i=\begin{cases}
       pB_i,&B_i<e,\\
       s-B_i,&B_i>e.
      \end{cases}                                    \tag{B-RD91.4}
\]
At \(B_i=e\) the quotient polygon has one slope, with every
nonzero first-level point at valuation \(e\). The quotient and
the actual semilinear descent are the same constructions as in
B69 and B83: no good integral model over \(K\) is asserted.
The comparison induces an actual \(p^2\)-torsion isomorphism,
because
\[
 Q_i[p^2]=\{P\in M_i[p^3]:p^2P\in H\}/H.           \tag{B-RD91.5}
\]
It also identifies the intrinsic images \(M_i[p]/H\) in
\(Q_i[p]\). Isogeny preserves the unordered finite special
character, so a recovery for these quotients recovers the original
character.

If \(B_i<e\), the intrinsic old image has valuation \(B_i\),
which is the low value \(D_i/p\), while the canonical high value
is \(s-pB_i>e\). These are two distinct stable lines. Over \(L\)
their characters are trivial: the residual image is a \(p\)-group,
whereas every line character is \(\mathbf F_p^\times\)-valued.
The \(p\)-group assertion follows directly because the quotient
torsion is a subquotient of the original \(p^3\)-torsion and the
original residual wild group has order \(p\). Thus \(Q_i[p]\)
is \(L\)-rational. At \(B_i=e\), the one-slope leading-coordinate
map is injective and wild inertia is trivial; the same \(p\)-group
observation again makes \(Q_i[p]\) \(L\)-rational.

If \(B_1<e\), both parameters are below \(e\). Both quotient
modules are therefore \(L\)-rational, and their parameters are
\(D_i=pB_i\). Proposition~\ref{prop:B-RD84-tame}, applied to
their actual \(p^2\)-comparison with \(p\nmid B_1\), recovers
the special characters, a contradiction.

Suppose next \(B_2\le e<B_1\). The second quotient has rational
first-level torsion. The first quotient has parameter
\(D_1=s-B_1\) prime to \(p\); its low points have nonintegral
valuation \(D_1/p\), so its residual division field over \(L\)
has nontrivial wild degree \(p\). This contradicts the actual
quotient first-level isomorphism. This includes the boundary
\(B_2=e\); the prime parameter \(B_1\) cannot equal \(e\).

Finally, if \(e<B_2<B_1\), both quotients have parameters
\(D_i=s-B_i\), with \(D_1\) prime and \(D_2\) divisible by
\(p\). Their identified old images are the canonical high lines,
of valuation \(B_i=s-D_i\). The first quotient has nontrivial
residual wild action, hence so does the second. The first-level
lemma applied to this actual quotient comparison gives
\(D_1>D_2\), opposite to \(B_1>B_2\). All cases are excluded.
\end{proof}

\begin{proposition}[The crossing parameters at \(p=5\)]
\label{prop:B-RD91-parameters}
\status{Proved; full independent root and companion audits passed.}
For \(p=5,e=20\), normalize the nonsquare common finite determinant
as in B77. A remaining different-character residual-wild comparison
has
\[
 B_i=5b_i,\qquad b_i\in\{1,3,5,7,9,11,13,15,17,19\},
 \qquad b_1\equiv b_2\pmod4.                         \tag{B-RD91.6}
\]
The two parameters cannot both be below \(20\). If they straddle
\(20\), their ordered values, with \(B_1<20<B_2\), are exactly
among the two necessary rows
\[
                 (B_1,B_2)=(5,45),\ (15,95).          \tag{B-RD91.7}
\]
Their common high quotients have rational first-level torsion and
exchanged canonical high lines, with unordered parameters
\(\{25,75\}\). If both original parameters exceed \(20\), their
common high quotients have nontrivial residual wild action.
\end{proposition}
\begin{proof}
The preceding proposition gives \(B_i=5b_i\), with
\(1\le b_i\le19\). Normalize \(c_i\equiv1\pmod4\).
The high-line character is \(\mathbf F_5^\times\)-valued, so
\(c_i-B_i\equiv0\pmod6\). Consequently \(B_i\), and hence
\(b_i\), is odd. Equality of the high-line characters and finite
determinants gives \(B_1\equiv B_2\pmod4\), hence the displayed
congruence. Equal parameters recover the finite character, so they
are distinct here. The only parameters below \(20\) are \(5,15\),
whose quotients by five are not congruent modulo four. The boundary
\(B_i=20\) is excluded by oddness.

In a crossing row the smaller quotient has rational first-level
torsion, so the other quotient does as well. Their parameters are
\(25b_1\) and \(5(24-b_2)\). Proposition B84 forces
\(5\mid(24-b_2)\), since otherwise one parameter divided by five
is prime to five and the quotient comparison recovers the character.
The odd possibilities \(5\le b_2\le19\) are \(9,19\). The
congruence modulo four then gives \((b_1,b_2)=(1,9),(3,19)\),
which yields (B-RD91.7). The old first-level image is low in the
smaller original model's quotient and high in the other, proving
the stated exchange and the actual quotient parameters.

If both original parameters exceed \(20\), their quotient
parameters are \(D_i=5(24-b_i)\), and their identified old images
are the common quotient high line. Suppose the quotient residual
action were tame, hence trivial over \(L\). B84 would force
both \(24-b_i\) divisible by five. Thus both \(b_i\) would lie
in \(\{9,19\}\), and their congruence modulo four would force
equality. This recovers the character, a contradiction. Therefore
the quotient residual wild action is nontrivial.
\end{proof}

\begin{theorem}[Only the two largest layer shapes remain above the index]
\label{thm:B-RD91-large}
\status{Proved; full independent root and companion audits passed.}
At \(p=5,e=20\), assume additionally that both \(B_i>20\).
Let \(W\) be the common actual \(125\)-division image over \(L\),
and let \(W_1,W_2\subset M_2(\mathbf F_5)\) be its first and
terminal congruence layers. In a basis with the common high line
\(\mathbf F_5e_1\),
\[
 \mathfrak{sl}_2(\mathbf F_5)\subset W_1\subset W_2,
 \qquad (\dim W_1,\dim W_2)\in\{(3,4),(4,4)\}.
                                                        \tag{B-RD91.8}
\]
Thus its full wild order is respectively \(5^8\) or \(5^9\).
\end{theorem}
\begin{proof}
The original residual image consists of matrices
\(1+uE_{12}\). For \(g\in W\) put
\[
 u(g)=\rho_{12}(g)\bmod5,\qquad
 v(g)=\rho_{21}(g)/5\bmod5.                          \tag{B-RD91.9}
\]
Both are additive homomorphisms \(W\to\mathbf F_5\): the
diagonal entries are one modulo five and the lower entry is
divisible by five. In the actual high-quotient Tate lattice,
with basis \((e_1/5,e_2)\), the residual lower entry is \(v(g)\),
as in (B-RD83.6). The preceding proposition says that \(v\ne0\).

Choose a prime-to-five order-24 lift of a tame generator in the
actual finite division group, and diagonalize its action on the
\(125\)-torsion by the coprime factors of \(X^{24}-1\), exactly
as in B90. Let its residual eigenvalues be \(\alpha,\beta\).
Its ratio \(r=\alpha/\beta\) is nonsquare in \(\mathbf F_5\),
so \(1,r,r^{-1}\) are distinct. Conjugation gives weights
\(r,r^{-1}\) on \(u,v\). If \(v\) vanished on
\(W^1=\ker u\), it would factor as a nonzero scalar multiple
of \(u\), contradicting those distinct weights. Hence \(W_1\)
has a nonzero lower off-diagonal entry. The tame spectral
projector gives \(E_{21}\in W_1\). Conjugating it by the
residual matrices \(1+tE_{12}\), and using \(t=1,-1\), gives
\(E_{12}\) and \(\operatorname{diag}(1,-1)\) as well.
Thus \(\mathfrak{sl}_2\subset W_1\). Actual fifth powers give
\(W_1\subset W_2\), and
\[
                  |W|=5^{1+\dim W_1+\dim W_2}.        \tag{B-RD91.10}
\]

It remains to exclude \(W_1=W_2=\mathfrak{sl}_2\). In that
case determinants on \(W^1\) are one modulo 25. The map
\(g\mapsto(\det g-1)/25\bmod5\) is additive, vanishes on the
terminal subgroup \(1+25W_2\), and therefore factors through a
linear functional \(\ell:W_1\to\mathbf F_5\). It is invariant
under both tame conjugation and residual unipotent conjugation.
The two nontrivial tame weights kill \(\ell(E_{12})\) and
\(\ell(E_{21})\); conjugation of \(E_{21}\) by \(1+E_{12}\)
then kills \(\ell(\operatorname{diag}(1,-1))\). Thus
\(\det(W^1)=1\pmod{125}\).

The determinant on \(W\) is the actual cyclotomic character,
since its finite tame factor is trivial on \(L\). Its image
therefore has order at most five. The exact cyclotomic tower
argument in B86 shows that it has order five, that
\(\mu_{25}\subset K\), and that, for a Cohen coefficient field
\(K_0\subset K\), one has \(K=K_0(\mu_{25})\). Indeed the
absolute index is exactly 20, while \(\mu_{125}\) requires
index 100; a tame degree-24 extension cannot change the remaining
five-primary cyclotomic degree. Since the determinant is trivial
on \(W^1\), its degree-five field is precisely
\(R=F^{W^1}\), the original residual division field.

The exact cyclotomic calculation of B86 now gives
\(b(R/L)=24^2=576\). On the other hand, take an actual low
first-level point. Its normalized \(R\)-coordinate has value
\(B_i\), and every nonidentity residual automorphism displaces
it by value \(5(s-B_i)\). The unconditional bound of B60 gives
\[
             b(R/L)\le5s-6B_i=600-6B_i<576,
                        \qquad B_i>20,                \tag{B-RD91.11}
\]
a contradiction. This excludes the only layer shape of total
order \(5^7\), and proves (B-RD91.8).
\end{proof}

\begin{remark}[Exact remaining scope]
\label{rem:B-RD91-scope}
\status{Proved; full independent root and companion audits passed.}
The crossing rows (B-RD91.7) are necessary patterns, not examples.
The theorem addresses only the both-large residual-wild branch;
it neither excludes nor constructs its two remaining layer shapes.
The irreducible residual branch is unchanged. The unrestricted
first-wild least index at \(p=5,e=20\) remains in \(\{2,3\}\).
All quotient comparisons have the actual levels displayed in
(B-RD91.5); in particular a quotient \(25\)-comparison is never
silently treated as an actual \(125\)-comparison.

The proof rereads the complete internal B60 fixed-coefficient
lemma, B69 quotient formula, B82 residual normalization, B84
point-field recovery, and B86 cyclotomic-field calculation.
No fresh original-source reading is claimed here. The retained
standalone program
\path{code/reduction_descent/verify_residual_wild_high_quotient.py}
reconstructs 78,761 inverse-Newton quotient distributions, including
all nine tested equality boundaries, and checks 210,956 complete
prime/divisible case rows over the four stated small primes.
Its independent finite-character enumeration retains 32 normalized
distinct-character rows, the four ordered crossing rows representing
(B-RD91.7), and the 20 both-large rows. It reconstructs all 250
linear functionals on the two tame forms of the trace-zero plane,
excludes all eight nonzero intertwiners between the opposite tame
characters, and checks 75,571 cyclotomic break inequalities.
The first author run passed at 03:02:51.567982 UTC on 20 November
2026 in 0.6828030579927145 seconds. No executed computation failed.
The program, full output, and pre-certificate source are retained.
These finite checks do not replace the uniform field and quotient
proofs, and no parameter row is treated as a formal-group existence
assertion.

\smallskip
\noindent\emph{Independent audits.}
The complete root and monodromy proof and exact-code audits passed.
Their literal author-output-preserving reruns took
\(0.69858235999709\) and
\(0.6949913400021615\) seconds, respectively.
Both reread the complete invoked B60, B69, B82, B84 and B86
dependency proofs. No additional original-source opening is claimed.
The exact audited candidate, pre-promotion copy, programs, outputs,
two full audit records and promotion diff remain retained.
\end{remark}


% Proved after complete root and monodromy audits; exact candidate retained.
\subsection{Actual diagonal fields remove one residual-wild crossing}

Throughout this subsection \(K\) is complete discretely valued of
mixed characteristic \((0,5)\), with algebraically closed residue
field, \(e_K=20\), and \(\mu_5\subset K\). The common tame good
field \(L/K\) has degree 24, so \(v_L(5)=480\). All point fields,
quotients and matrix images below arise from the stated actual
torsion comparisons. Breaks use \(b(g)=v(g\varpi-\varpi)-1\).

\begin{lemma}[The exchanged \(25\)-comparison already has full first image]
\label{lem:B-RD92-tame25}
\status{Proved; full independent root and companion audits passed.}
Suppose two whole potentially supersingular height-two groups have
isomorphic actual \(25\)-torsion, rational first-level torsion over
\(L\), parameters \(25,75\), and exchanged high lines
\(\mathbf F_5e_1,\mathbf F_5e_2\). Assume their normalized residual
tame exponents are \((0,6)\), or \((12,18)\), in this basis.
Their common \(25\)-division image over \(L\) is
\(1+5M_2(\mathbf F_5)\). Its four coordinate cyclic characters
have breaks
\[
 \begin{array}{c|rrrr}
 \text{character}&A_{21}&A_{22}&A_{12}&A_{11}\\\hline
 \text{break}&6&24&66&72.
 \end{array}                                           \tag{B-RD92.1}
\]
No actual \(125\)-comparison is assumed in this lemma.
\end{lemma}
\begin{proof}
Let \(V\subset M_2(\mathbf F_5)\) be the actual additive image.
Choose a second-level point above \(e_i\), with field \(E_i\).
First-level rationality makes \(E_i/L\) Galois with translation
group \(Ve_i\). Its degree is at most 25. The inverse Newton
polygons in both models are single segments: both parameters exceed
20, including the high-target inverse. The coordinate valuations
over \(L\), in the two models, are
\[
  (19/5,3/5)\quad(i=1),\qquad
  (1/5,9/5)\quad(i=2).                                \tag{B-RD92.2}
\]
Thus each degree is at least five. If it were five, all corresponding
normalized coordinate valuations would be prime to five, and the
exact B60 break formula would apply. There are three possible types
of its translation line: the first high line, the second high line,
or a line different from both. For \(i=1\), the two exact breaks
would respectively be
\[
          (456,72),\qquad(6,222),\qquad(6,72);
                                                        \tag{B-RD92.3}
\]
for \(i=2\), they would be
\[
          (474,66),\qquad(24,216),\qquad(24,66).
                                                        \tag{B-RD92.4}
\]
For example the translating high/low valuations over a degree-five
field are \((475,25)\) in the first model and \((225,75)\) in the
second. Subtracting the two coordinate values in (B-RD92.2), scaled
by five, gives these rows. None agrees. Hence both evaluations
\(V\to\mathbf F_5^2\) are surjective.

Tame conjugation has distinct weights 0, 6 and 18 on the diagonal
plane and the two off-diagonal axes. Surjectivity of both evaluations
therefore puts \(E_{12},E_{21}\) in \(V\), and makes both diagonal
projections nonzero. Consequently \(\dim V\ge3\). If equality
held, the diagonal part would be \(k_5\operatorname{diag}(a,b)\),
with \(ab\ne0\); the two degree-25 point fields would share its
actual diagonal cyclic quotient. Here \(k_5=\mathbf F_5\).

The complete shared-diagonal proof of Proposition B87 then applies
with no higher torsion level: its four coordinate rows are
\[
\begin{array}{c|c|r|r|r|c}
 \text{field}&\text{model}&A&D_0&D_1&H\\\hline
 E_1&1&95&125&2375&\text{diagonal}\\
 E_1&2&15&375&1125&\text{off-diagonal}\\
 E_2&1&5&125&2375&\text{off-diagonal}\\
 E_2&2&45&375&1125&\text{diagonal}.
\end{array}                                             \tag{B-RD92.5}
\]
All four displacements are below \(v_{E_i}(5)=12000\), and every
\(A\) has five-adic order one. The diagonal cyclic break is a
positive multiple of 24, while the off-diagonal breaks are 6 and
18 modulo 24. In \(E_1\), an off-diagonal last line would give its
larger cyclic break 6, impossible; hence its diagonal break is 72.
In \(E_2\), a diagonal last line gives diagonal break 24, whereas
an off-diagonal last line gives off-diagonal break 66, which could
not exceed the shared diagonal break 72. This excludes dimension
three, using exactly the fixed-coordinate lemma B87 and the same
actual quotient field. Thus \(V=M_2(\mathbf F_5)\).

Finally the complete point-field calculation in B88 uses only this
actual full \(25\)-image and (B-RD92.5). To spell out its finite
alternatives, in \(E_1\) the last line is diagonal, with cyclic
break 72; the smaller break is \(t\in\{6,30,54\}\), and the
larger lower break is \(360-4t\). At \(t=30\) the compulsory
prime exponent is 135; at \(t=54\) the other row would require
an exponent \(71<95\). Hence only \((t,u)=(6,336)\) remains.
In \(E_2\), a diagonal last line gives \((18,48)\), whose forced
prime term changes at 165 instead of 1125. An off-diagonal last
line gives cyclic break 66 and \(t\in\{24,48\}\). The latter
case forces a change at 215 instead of 2375. Thus
\((t,u)=(24,234)\). The corresponding two pairs of cyclic
breaks are \((A_{11},A_{21})=(72,6)\) and
\((A_{22},A_{12})=(24,66)\), proving (B-RD92.1).
\end{proof}

\begin{lemma}[The actual image in either residual-wild crossing]
\label{lem:B-RD92-parahoric}
\status{Proved; full independent root and companion audits passed.}
Suppose an actual \(125\)-comparison is in one of the B91
crossing rows \((5,45)\) or \((15,95)\), with common high line
\(H=\mathbf F_5e_1\). Its full wild image is
\[
 W=\left\{
 \begin{pmatrix}1+5a&b\\25c&1+5d\end{pmatrix}:
 a,d\in\mathbf Z/25,\ b\in\mathbf Z/125,
 c\in\mathbf F_5\right\},\qquad |W|=5^8.            \tag{B-RD92.6}
\]
In particular its first and terminal layers are respectively
the upper-triangular three-space and all of \(M_2(\mathbf F_5)\).
Let \(R/L\) be the original first-level division field, and let
\(D_{22}/L\) be the cyclic field defined by
\((\rho_{22}-1)/5\bmod5\). Their breaks are respectively
\[
 \begin{array}{c|cc}
 (B_1,B_2)&b(R/L)&b(D_{22}/L)\\\hline
 (5,45)&66&24\\
 (15,95)&6&72.
 \end{array}                                           \tag{B-RD92.7}
\]
The two fields in each row are distinct actual subfields.
\end{lemma}
\begin{proof}
Quotient both groups by their common high line. B91 gives an
actual \(25\)-comparison with rational first-level torsion and
exchanged parameters \(25,75\). The preceding lemma applies.
In the original integral basis the quotient lattice has basis
\((e_1/5,e_2)\), so its matrices are
\[
       \begin{pmatrix}\rho_{11}&5\rho_{12}\\
                   \rho_{21}/5&\rho_{22}\end{pmatrix}.
                                                        \tag{B-RD92.8}
\]
Rationality of the quotient first-level module therefore forces
\(25\mid\rho_{21}\) on the entire original wild image. It lies
in the displayed group (B-RD92.6).

The full quotient \(25\)-image makes the four additive characters
\[
 \rho_{12}\bmod5,\quad
 (\rho_{11}-1)/5\bmod5,\quad
 (\rho_{22}-1)/5\bmod5,\quad
 \rho_{21}/25\bmod5                                  \tag{B-RD92.9}
\]
jointly surjective. They are additive directly from the matrix
product and \(25\mid\rho_{21}\). Choose an integral eigenbasis
for an actual order-24 tame lift. The first congruence layer is
upper triangular, and its diagonal projection is surjective:
set the first character of (B-RD92.9) to zero while prescribing
the next two. Spectral projection gives the full diagonal plane.
Also, an element with nonzero residual upper entry has fifth power
\(1+5uE_{12}\pmod{25}\). This follows either by multiplying
the triangular matrices or by summing the five upper entries;
the terms from their diagonal corrections have coefficient ten
and vanish modulo 25. Thus the first layer is the full upper
triangular three-space.

If the last character in (B-RD92.9) vanished on the terminal layer,
its restriction to the first congruence subgroup would factor
through that upper-triangular first layer. Its tame weight is the
lower off-diagonal weight, distinct from both the diagonal and
upper weights, so that restriction would vanish. It would then
factor through the original residual upper character, again of
the other weight, and would be zero altogether. This contradicts
the joint surjectivity. Thus the terminal layer has a nonzero lower
entry. Tame projection supplies \(E_{21}\), and fifth powers
already put the whole first layer in the terminal layer. The latter
is therefore \(M_2(\mathbf F_5)\). The total order is
\(5^{1+3+4}=5^8\), the order of the containing group in
(B-RD92.6); equality follows.

It remains to identify the actual cyclic fields, including the
index order. In the \((5,45)\) row, the quotient with parameter25
is the first model and its high line is the first coordinate of
(B-RD92.8). Thus the original residual upper character is its
\(A_{12}\), and the original lower diagonal is its \(A_{22}\).
Their breaks in (B-RD92.1) are 66 and 24. In the \((15,95)\)
row, the parameter25 quotient is the second model: its high line
is the old first-level image, the second coordinate of
(B-RD92.8). Reordering the quotient basis to put this high line
first sends the original upper residual character to \(A_{21}\)
and the original lower diagonal to \(A_{11}\). Their breaks are
6 and 72. These are distinct independent characters on the actual
full additive quotient image, proving (B-RD92.7) and distinctness.
\end{proof}

\begin{lemma}[Two coordinates with the same high-translation line]
\label{lem:B-RD92-common-high}
\status{Proved; full independent root and companion audits passed.}
Let \(R/L\) and \(E/R\) be actual extensions in the crossing
setting, with \([R:L]=5\), \(\operatorname{Gal}(E/R)\simeq
\mathbf F_5^2\), and let \(b_1<b_2\) be prime to five,
\(1\le b_i\le19\). Suppose two formal coordinates in \(E\)
have the same high-displacement line and data
\[
 A_i=5b_i,\qquad D_{0,i}=125b_i,\qquad
 D_{1,i}=15000-625b_i,\qquad v_E(5)=60000.             \tag{B-RD92.10}
\]
Then the last ramification line is the common high line, the two
lower breaks are
\[
 t=24b_1,\qquad u=15000-630b_2-96b_1,                 \tag{B-RD92.11}
\]
and the cyclic quotient killing the high line has break \(24b_1\).
\end{lemma}
\begin{proof}
All displacement values are strictly below the mixed-characteristic
cutoff. Since \(v_5(A_i)=1\), the same fixed-coordinate expansion
as in B87 has first possible change at
\(\min(A_i+5r,J_i+r)\), where \(J_i>A_i\) is the first nonzero
prime-to-five exponent, possibly infinite. A strict smaller value
can cancel only when \(J_i=A_i+4r\).

There cannot be a single ramification break \(r\). Its leading
coefficient map on \(\mathbf F_5^2\) has a two-dimensional image.
A first change strictly below \(D_{0,i}\) would be a nonzero
additive polynomial of degree at most five vanishing on all 25
leading coefficients, impossible. If the first change were
\(D_{0,i}\) with a unique monomial, it could not vanish exactly
on the high line. Thus a tie at \(D_{0,i}\) is required, giving
\(r=(D_{0,i}-A_i)/5=24b_i\) in both models, a contradiction.

Write the two breaks as \(t<u\), with last line \(T\). If
\(T\) differed from the common high line, the coefficient lemma
B87 would give the same larger cyclic quotient break \(24b_i\)
in both models. Hence \(T\) is the common high line. Put
\(r_i=24b_i\), so \(r_1<r_2\). At the smaller break, every
element outside the high line has displacement \(D_{0,i}\).
Consequently
\[
 \begin{array}{c|c|c}
 &J_i&\text{forced relation at }u\\\hline
 t<r_i&A_i+4t&u+4t=15000-630b_i\\
 t>r_i&D_{0,i}-t&u-t=15000-750b_i.
 \end{array}                                           \tag{B-RD92.12}
\]
In both rows the prime term is uniquely first at \(u\). If both
models use the same strict row, their parameters are equal. If
\(r_1<t<r_2\), subtracting the two relations gives
\(t=150b_1-126b_2<24b_1\), impossible.

If \(t=r_2\), the first row's model, namely model1 in the second
strict case, gives \(u=15000-750b_1+24b_2\). In model2 one must
have \(J_2\ge101b_2\). At \(u\), its leading divisible term
exceeds the required displacement by
\[
 (5b_2+5u)-D_{1,2}=750(80-5b_1+b_2)>0.
\]
It therefore requires \(J_2=D_{1,2}-u=750b_1-649b_2\), which
would force \(b_1\ge b_2\). Thus \(t=r_2\) is impossible.
The only remaining value is \(t=r_1\); model2 is then in the
first strict row of (B-RD92.12), giving (B-RD92.11). The exact
cyclic quotient killing the last line has the smaller break \(t\),
as follows directly from the quotient-average formula. This proves
the assertion without assuming that every formal compatible profile
is realized by a field or a group.
\end{proof}

\begin{theorem}[The \((15,95)\) crossing is impossible]
\label{thm:B-RD92-crossing}
\status{Proved; full independent root and companion audits passed.}
The residual-wild crossing row \((B_1,B_2)=(15,95)\) cannot occur
in an actual different-special-character \(125\)-comparison.
Consequently a remaining crossing must have ordered parameters
\((5,45)\), and then has the full actual image (B-RD92.6),
with residual field break66 and diagonal field break24 as above.
\end{theorem}
\begin{proof}
Take an actual primitive \(25\)-point \(P\) with \(5P=e_2\)
outside the original common high line, and put \(E=L(P)\).
The full group (B-RD92.6) modulo25 has order \(5^4\). Its orbit
on this point has size \(5^3\): the upper entry modulo25 and
the lower diagonal correction modulo five vary independently.
Thus \([E:L]=125\), and \(R\subset E\). Over \(R\), the
congruence image is the additive upper triangular three-space;
evaluation on \(e_2\) is surjective with kernel
\(\mathbf F_5E_{11}\). Therefore \(E/R\) is an actual Galois
extension of degree25 with every first-level translation present.
The low inverse coordinates have \(v_L(T(P))=b_i/25\), while
the translating points have values \(b_i\) and \(120-5b_i\).
Normalization by \([E:L]=125\) gives exactly (B-RD92.10).

In the translation coordinates for \(E/R\), the high line is
the upper entry and the quotient killing it is the lower diagonal
character. Hence its actual cyclic field is \(D_{22}R/R\).
In the \((15,95)\) row one has \((b_1,b_2)=(3,19)\), so
the preceding lemma gives
\[
                         b(D_{22}R/R)=72.             \tag{B-RD92.13}
\]
But (B-RD92.7) gives two distinct cyclic fields over \(L\) with
breaks \(b(R/L)=6\) and \(b(D_{22}/L)=72\). Their actual
elementary abelian compositum has upper breaks6 and72, and hence
lower breaks6 and \(6+5(72-6)=336\). Restriction to its subgroup
fixing \(R\), which is the last line, gives
\[
                         b(D_{22}R/R)=336,             \tag{B-RD92.14}
\]
a contradiction. The conversions are precisely subgroup lower
numbering and quotient upper numbering, as in the full norm and
Herbrand proofs already used in B88--B90. No relation between
unrelated abstract fields is used.

For the other row the same test gives24 over \(R\), consistently
with breaks66 and24 over \(L\). The theorem makes no exclusion
or existence claim for that row.
\end{proof}

\begin{remark}[The remaining actual crossing]
\label{rem:B-RD92-scope}
\status{Proved; full independent root and companion audits passed.}
The lower row is excluded by an actual cyclic field comparison,
not by the size of a hypothetical matrix group. The row \((5,45)\)
remains open despite the full necessary parahoric image. The
both-large residual-wild alternatives of B91 and the irreducible
residual branch are unchanged. The unrestricted index at
\(p=5,e=20\) therefore remains in \(\{2,3\}\).

The proof explicitly reuses only those parts of B87 and B88 that
concern the actual \(25\)-division image, and supplies the missing
degree-five point-field exclusions needed to obtain that image
without a \(125\)-comparison of the high quotients. The old level
is not silently restored after quotienting. The retained standalone
program \path{code/reduction_descent/verify_residual_wild_crossing.py}
reconstructs all 12 degree-five translation lines, the four possible
three-dimensional diagonal lines, the five relevant point-field
coefficient cases and the separate excluded diagonal case. It checks
625 actual triangular fifth powers and 3,125 quotient matrix-digit
rows, then tests 2,138,160 necessary common-high break pairs. The
surviving necessary profiles are exactly \((24,9234)\) and
\((72,2742)\); the actual diagonal-field comparison excludes only
the second crossing. The first author run passed at
03:16:05.929610 UTC on 20 November 2026 in
1.205554521002341 seconds. The complete program, output and
pre-certificate source are retained; no executed computation failed.
These finite checks supplement the actual field and coefficient
proofs and imply no realization of a surviving numerical profile.
The complete internal B60, B69, B87, B88, B90 and B91 dependency
proofs were reread; no fresh external original-source opening is
claimed in this candidate.

\smallskip
\noindent\emph{Independent audits.}
The complete root and monodromy proof and exact-code audits passed.
Their literal author-output-preserving reruns took
\(1.2081115340115502\) and
\(1.17941884699394\) seconds, respectively.
Both reread the full actual B87 and B88 coefficient and point-field
proofs and verified that precisely their actual 25-level statements
are used here. The invoked earlier quotient and fixed-coefficient
arguments were also rechecked. No fresh external original-source
opening is claimed. The exact candidate, pre-promotion copy,
programs, outputs, full audit records and promotion diff are retained.
\end{remark}

